\documentclass[11pt]{article}
\usepackage[T1]{fontenc}
\usepackage{amsmath,amssymb,amsthm}
\usepackage[margin=1in]{geometry}
\usepackage{enumitem}
\usepackage[colorlinks=true,linkcolor=blue,urlcolor=blue]{hyperref}

\newtheorem{theorem}{Theorem}[section]
\newtheorem{proposition}[theorem]{Proposition}
\newtheorem{lemma}[theorem]{Lemma}
\newtheorem{corollary}[theorem]{Corollary}
\newtheorem{definition}[theorem]{Definition}
\newtheorem{warning}[theorem]{Warning}
\newtheorem{firewall}[theorem]{Claim firewall}
\newtheorem{decision}[theorem]{Next decision}
\newtheorem{assessment}[theorem]{Assessment}
\newtheorem{ledger}[theorem]{Acquisition ledger}
\theoremstyle{remark}
\newtheorem{remark}[theorem]{Remark}

\title{Renewal Standard Model--Einstein Continuum Research Notes\\
Twelfth Cycle, V1}
\author{}
\date{9 September 2026}

\begin{document}
\maketitle
\tableofcontents

\section{Context, programme status, and the new compact cycle}
\label{sec:v12v1}

These notes continue the attempt to construct a mathematically
controlled continuum limit of a regulated Standard Model coupled to
Einstein gravity.  The aim is a projective continuum of gauge,
fermion, Higgs, and gravitational fields with a nonzero normalised
functional, the correct gauge and diffeomorphism identities,
reflection positivity on the physical observable algebra, and the
required long-distance clustering or massless-sector structure.

The programme has not reached that goal.  The purpose of this note is
to keep the distinction between proved implications and missing
physical inputs explicit while advancing the first unresolved
construction.

The preceding active manuscript reached V100 and has been frozen as
\[
 \texttt{Renewal\_SM\_Einstein\_Continuum\_Research\_Notes\_V100\_f11.tex}.
\]
Its terminal record is
\[
\begin{array}{c|c}
\text{lines}&29662\\
\text{bytes}&969144\\
\text{SHA--256}&
\texttt{127FDC3C8C6CDFEE05820C03237619C8C814E119A1626F586BC77330272CCD45}.
\end{array}                                                     \tag{1.1}
\]
Earlier frozen cycles remain separate archives.  This compact V1
summarises the major results rather than reproducing their proofs.
Future versions append to this file, replace only their immediate
active predecessor, audit the newest Yang--Mills and Navier--Stokes
companions at every fifth version, and roll over again at V100.

\subsection{Major mathematical results acquired so far}

The archived notes establish the following rigorous components.

\begin{enumerate}[leftmargin=2.2em]
\item Finite-cutoff renewal maps, local slice carriers, Schur and
      Feshbach reductions, and determinant regularisations have been
      separated into explicit hypotheses with stability estimates.
\item Exact orbit--stabiliser weighting, finite-grid Ward telescopes,
      and full-density coordinate covariance prevent premature
      symmetry compression of partial diagrammatic expressions.
\item Gaussian reference laws have projective limits under covariance
      consistency and spectral-tail tightness.  Summable adjacent
      covariance or density defects give convergent cylinder laws.
\item An action cost exceeding sector entropy gives exponential
      topological-charge tails.  Gapped Hermitian index carriers give
      integer charges that remain stable under certified refinement
      homotopies.
\item Positive reference estimates transfer to interacting amplitudes
      through uniform \(L^r\) bounds and a lower partition-function
      anchor, with the exact Holder exponent loss recorded.
\item A finite phase anchor survives summable positive-measure,
      charge-admissibility, determinant-transport, and line-holonomy
      defects, yielding a nonzero normalised complex cylinder
      functional.
\item Relative trace-class or Hilbert--Schmidt operator increments
      control determinant phases on zero-free paths.  Exact local
      counterterms remove extensive bulk phase terms only when the
      corresponding one-form is exact; global cycle holonomy remains
      a separate obstruction.
\item A quantitative diffeomorphism Faddeev--Popov gap gives a local
      slice and a relative determinant bound.  Global quotient
      measures additionally require full-Jacobian overlap equality,
      Gribov multiplicity and stabiliser corrections, and equivariant
      refinement.
\item The bare Euclidean Einstein--Hilbert action is unbounded below
      along bounded, fixed-volume, fixed-topology conformal
      oscillations.  It cannot be used without qualification as the
      positive amplitude base.
\item A coercive real Gaussian conformal reference has a
      Sobolev-valued projective continuum, but a naive fourth-order
      regulator cannot be removed while preserving both uniform
      \(L^r\) Einstein reweighting and nondegenerate fixed modes.
\item Complex contours are controlled by real-part coercivity,
      singular-divisor avoidance, relative homology, normalised
      new-mode factors, and a denominator anchor.  Antilinear
      invariance gives reality; a positive Gram kernel is required
      for Osterwalder--Schrader positivity.
\item Theta phases preserve the reflection kernel when they are exact
      half-space coboundaries.  Pairing conjugate sectors or thimbles
      gives reality but can leave negative reflection eigenchannels.
\item Independent negative quadratic conformal modes have an exact
      imaginary-axis Gaussian contour.  Dividing each new-mode form
      by its orientation phase produces exact projective consistency
      and an infinite Sobolev-valued amplitude under a trace
      condition.
\end{enumerate}

\begin{ledger}[Leading unresolved inputs]
\label{led:v12v1-open}
The full continuum still requires:
\begin{enumerate}
\item a nonlinear coupled conformal contour or a different
      gravitational ultraviolet construction;
\item a singularity-free, refinement-compatible diffeomorphism
      quotient atlas on a high-probability or controlled contour set;
\item uniform gauge and gravitational sector energy--entropy
      estimates;
\item a globally transported chiral determinant line and nonzero
      phase anchor;
\item reflection factorisation for the full fermion, ghost, gauge,
      Higgs, and metric kernel;
\item vanishing-defect Ward identities, clustering estimates, and
      physical universality and parameter matching.
\end{enumerate}
\end{ledger}

\subsection{Current direction}

The first exact contour block is quadratic.  The next noncircular
question is whether an interacting quadratic Hessian remains a
well-defined projective complex Gaussian after it is measured relative
to that block.  This is answered below in trace-class form.  Nonlinear
interactions and reflection positivity are left as subsequent,
separate gates.

\section{Trace-class quadratic perturbations of the imaginary
Gaussian}
\label{sec:v12v1-quadratic}

Let \({\cal K}=\ell^2(\mathbb N;\mathbb C)\), let \(P_N\) project onto
the first \(N\) coordinates, and let \(x_1,x_2,\ldots\) be independent
standard real Gaussians.  The normalised imaginary conformal variables
are \(z_n=i\kappa_n^{-1/2}x_n\), with \(\kappa_n>0\).

A quadratic holomorphic interaction
\[
                              V_N(z)={1\over2}z^TB_Nz         \tag{1.2}
\]
pulls back to
\[
                              -V_N(i\kappa^{-1/2}x)
                              ={1\over2}x^TC_Nx,\qquad
 C_N=P_NCP_N,                                                \tag{1.3}
\]
where \(C=C^T\) is the interaction Hessian measured relative to the
diagonal reference stiffness.

Assume throughout this section that \(C\) is trace class and
\[
                              \|C\|\le q<1.                  \tag{1.4}
\]

\begin{theorem}[Exact finite-cutoff denominator]
\label{thm:v12v1-denominator}
For every \(N\),
\[
 D_N
 :=\mathbb E\exp\left({1\over2}x^TC_Nx\right)
 =\det(I-C_N)^{-1/2},                                       \tag{1.5}
\]
where the square-root branch is continued from \(C_N=0\).
The integral is absolutely convergent and \(D_N\ne0\).
Moreover,
\[
 e^{-L_N}\le|D_N|\le e^{L_N},\qquad
 L_N={\|C_N\|_1\over2(1-q)}.                                \tag{1.6}
\]
\end{theorem}

\begin{proof}
For real \(x\),
\[
 \operatorname{Re}\,x^T(I-C_N)x
 \ge(1-q)\|x\|^2,                                           \tag{1.7}
\]
so the complex Gaussian is absolutely integrable.  The usual
finite-dimensional Gaussian integral, first for real symmetric
\(C_N\) and then by holomorphic continuation in the open ball
\(\|C_N\|<1\), gives (1.5).  The determinant is nonzero because
\(I-C_N\) is invertible.

The trace expansion converges absolutely:
\[
 \log D_N
 ={1\over2}\sum_{m\ge1}{\operatorname{Tr}(C_N^m)\over m}.    \tag{1.8}
\]
Since
\[
 |\operatorname{Tr}(C_N^m)|
 \le\|C_N\|_1\|C_N\|^{m-1}
 \le\|C_N\|_1q^{m-1},
\]
one obtains
\[
 |\log D_N|
 \le{\|C_N\|_1\over2}
       \sum_{m\ge1}{q^{m-1}\over m}
 \le{\|C_N\|_1\over2(1-q)}.
\]
Taking real parts and exponentiating proves (1.6).
\end{proof}

\begin{theorem}[Fredholm limit and adjacent phase budget]
\label{thm:v12v1-fredholm}
The compressions satisfy
\[
                              \|C_N-C\|_1\longrightarrow0,   \tag{1.9}
\]
and
\[
 D_N\longrightarrow
 D_\infty=\det(I-C)^{-1/2}\ne0.                             \tag{1.10}
\]
For two trace-class perturbations \(C_0,C_1\) with
\(\|C_j\|\le q\),
\[
 |\log D(C_1)-\log D(C_0)|
 \le{\|C_1-C_0\|_1\over2(1-q)}.                             \tag{1.11}
\]
Consequently a refinement schedule with
\[
                              \sum_N\|C_{N+1}-C_N\|_1<\infty \tag{1.12}
\]
has a summable determinant modulus and phase increment.
\end{theorem}

\begin{proof}
Trace-class operators are approximated in trace norm by finite-rank
operators, and strong coordinate projections converge uniformly on a
fixed finite-rank range.  This proves (1.9).  Continuity of the
Fredholm determinant gives (1.10).

For (1.11), put \(C(t)=C_0+t(C_1-C_0)\).  The operator norm is at most
\(q\), and
\[
 {d\over dt}\log D(C(t))
 ={1\over2}\operatorname{Tr}
   \left[(I-C(t))^{-1}(C_1-C_0)\right].
\]
The inverse norm is at most \((1-q)^{-1}\).  Integrate and use the
trace duality bound.  Summation gives the final claim.
\end{proof}

Condition (1.12) is sufficient for an append-by-shell proof.  Direct
trace-norm convergence (1.9) already gives the endpoint limit and can
be used when the chosen adjacent blocking is not absolutely
summable.

\subsection{Normalised cylinder correlations}

Define the finite-cutoff normalised functional
\[
 \Omega_N(F)
 ={ \mathbb E[
       F(x)\exp(\frac12x^TC_Nx)]\over D_N}.                  \tag{1.13}
\]

\begin{theorem}[Complex Gaussian cylinder continuum]
\label{thm:v12v1-cylinder}
For a finitely supported real vector \(h\),
\[
 \Omega_N(e^{ih^Tx})
 =\exp\left[
 -{1\over2}h^T(I-C_N)^{-1}h
 \right]                                                     \tag{1.14}
\]
once \(N\) contains the support of \(h\).  Furthermore,
\[
 \|(I-C_N)^{-1}-(I-C)^{-1}\|
 \le{\|C_N-C\|\over(1-q)^2}\longrightarrow0.                \tag{1.15}
\]
Hence all polynomial cylinder correlations converge and obey Wick's
rule with complex covariance
\[
                                      G=(I-C)^{-1}.           \tag{1.16}
\]
\end{theorem}

\begin{proof}
Complete the square in the absolutely convergent Gaussian integral to
obtain (1.14), initially in a real neighbourhood and then by
holomorphic continuation under (1.7).  The resolvent identity gives
\[
 (I-C_N)^{-1}-(I-C)^{-1}
 =(I-C_N)^{-1}(C_N-C)(I-C)^{-1},
\]
which proves (1.15).  Differentiating (1.14) at \(h=0\) gives Wick's
rule and convergence of every fixed polynomial moment.
\end{proof}

\subsection{What trace class means for the conformal Hessian}

If \(B\) is the unscaled quadratic interaction in (1.2) and
\(K=\operatorname{diag}(\kappa_n)\), then
\[
                              C=K^{-1/2}BK^{-1/2}.            \tag{1.17}
\]
The new theorem applies when this relative Hessian is trace class and
has operator norm below one.  A local two-derivative Einstein Hessian
measured against an equally low-order reference generally fails this
condition by mode counting.  A stronger reference can make (1.17)
trace class, but recovering the intended infrared dynamics then
requires a separate renormalisation and universality argument.

\begin{warning}[Existence does not imply reflection positivity]
\label{warn:v12v1-reflection}
The bound (1.4) proves absolute convergence, a nonzero Fredholm
denominator, and cylinder convergence.  A complex symmetric
covariance \(G\) need not yield a positive reflected kernel.
Reflection positivity requires the half-space Gram or relative-form
criteria recorded in the frozen V96--V98 notes.
\end{warning}

\begin{firewall}[Quadratic scope]
\label{fire:v12v1-boundary}
No nonlinear Einstein interaction, metric-degeneracy divisor,
diffeomorphism quotient, fermion determinant, or gauge coupling has
been controlled by Theorems~\ref{thm:v12v1-denominator}--
\ref{thm:v12v1-cylinder}.  The trace-class hypothesis (1.4) is a
concrete acquisition target, not an established property of the full
regulated theory.
\end{firewall}

\begin{decision}[V2 acquisition]
\label{dec:v12v1-next}
V2 will treat nonlinear polynomial and analytic interactions under
the imaginary Gaussian amplitude.  It will seek a dimension-uniform
cluster or dependency-graph criterion that gives absolute
integrability, a nonzero denominator, and summable conditional
increments without reducing the problem to a global sup norm.
\end{decision}

\begin{assessment}[Progress after twelfth-cycle V1]
\label{ass:v12v1-status}
The rollover summary is complete and the first new result is
substantive: every relatively trace-class quadratic perturbation with
norm below one has an exact zero-free Fredholm denominator and a
projective complex Gaussian cylinder continuum.  The next obstacle is
the nonlinear local interaction and its reflection-compatible
dependency structure.
\end{assessment}

% =============================================================================
% BEGIN APPENDED TWELFTH-CYCLE V2 MATERIAL
% =============================================================================

\section{Conditional shell normalisation for nonlinear interactions}
\label{sec:v12v2}

The trace-class quadratic theorem in V1 is exact but does not cover a
general nonlinear local action.  A global uniform bound on the whole
interaction is usually extensive.  This version replaces it by a
conditional shell construction whose errors are measured one
refinement layer at a time.

Let \((\Omega,{\cal F},\nu)\) carry the projective imaginary-Gaussian
reference of V1, and let
\[
              {\cal F}_0\subset{\cal F}_1\subset\cdots
\]
be the filtration generated by successive mode or geometric shells.

\begin{theorem}[Complex multiplicative martingale]
\label{thm:v12v2-martingale}
Let \(R_j\) be complex, \({\cal F}_j\)-measurable random variables
such that
\[
                 \mathbb E_\nu[R_j\mid{\cal F}_{j-1}]=1      \tag{2.1}
\]
and
\[
 \mathbb E_\nu[|R_j-1|^2\mid{\cal F}_{j-1}]
                         \le\varepsilon_j^2                 \tag{2.2}
\]
almost surely.  Put
\[
                              W_N=\prod_{j=1}^NR_j.           \tag{2.3}
\]
If
\[
                              \sum_{j\ge1}\varepsilon_j^2<\infty,
                                                                    \tag{2.4}
\]
then \((W_N)\) is an \(L^2\)-bounded complex martingale and converges
both almost surely and in \(L^2(\nu)\) to \(W_\infty\).  Moreover,
\[
 \mathbb E_\nu W_N=\mathbb E_\nu W_\infty=1,                \tag{2.5}
\]
\[
 \sup_N\mathbb E_\nu|W_N|^2
 \le\prod_{j\ge1}(1+\varepsilon_j^2)
 \le\exp\left(\sum_{j\ge1}\varepsilon_j^2\right).            \tag{2.6}
\]
\end{theorem}

\begin{proof}
Equation (2.1) gives
\[
 \mathbb E[W_j\mid{\cal F}_{j-1}]=W_{j-1},
\]
so \(W_j\) is a martingale and (2.5) holds at finite \(j\).  Also,
conditionally on \({\cal F}_{j-1}\),
\[
 \begin{split}
 \mathbb E[|R_j|^2\mid{\cal F}_{j-1}]
 &=1+2\operatorname{Re}
       \mathbb E[R_j-1\mid{\cal F}_{j-1}]\\
 &\quad+\mathbb E[|R_j-1|^2\mid{\cal F}_{j-1}]\\
 &\le1+\varepsilon_j^2.
 \end{split}
\]
Hence
\[
 \mathbb E|W_j|^2
 \le(1+\varepsilon_j^2)\mathbb E|W_{j-1}|^2.
\]
Iteration gives (2.6).  The \(L^2\) martingale convergence theorem
gives almost-sure and \(L^2\) convergence.  Passing expectation
through the \(L^2\), hence \(L^1\), limit proves (2.5).
\end{proof}

\begin{corollary}[A countably additive complex continuum measure]
\label{cor:v12v2-measure}
Define
\[
                              {\mathfrak m}(A)
                              =\int_AW_\infty\,d\nu.          \tag{2.7}
\]
Then \({\mathfrak m}\) is a countably additive complex measure with
\[
 {\mathfrak m}(\Omega)=1,\qquad
 \|{\mathfrak m}\|_{\rm TV}
 \le\|W_\infty\|_{L^1}
 \le\exp\left({1\over2}\sum_j\varepsilon_j^2\right).         \tag{2.8}
\]
For every bounded cylinder \(F\),
\[
                              \mathbb E_\nu[FW_N]
                              \longrightarrow\int F\,d{\mathfrak m}.
                                                                    \tag{2.9}
\]
\end{corollary}

\begin{proof}
An \(L^1\) density defines a countably additive complex measure.
Cauchy--Schwarz and (2.6) give (2.8).  Equation (2.9) follows from
\[
 |\mathbb E[F(W_N-W_\infty)]|
 \le\|F\|_\infty\|W_N-W_\infty\|_{L^1}.
\]
\end{proof}

Unlike an anchored denominator estimate, (2.5) gives the denominator
exactly.  The price is the exact conditional normalisation (2.1).

\subsection{Obtaining the factors from shell interactions}

Let \(v_j\) be a complex \({\cal F}_j\)-measurable shell interaction
and define its conditional partition factor
\[
                    m_j=
                    \mathbb E_\nu[e^{-v_j}\mid{\cal F}_{j-1}].
                                                                    \tag{2.10}
\]
Assume
\[
                              |m_j-1|\le a_j<1               \tag{2.11}
\]
almost surely, and set
\[
                              R_j={e^{-v_j}\over m_j}.        \tag{2.12}
\]

\begin{proposition}[Zero-free local normalisation]
\label{prop:v12v2-local}
Equations (2.10)--(2.12) imply (2.1).  If
\[
 \mathbb E_\nu[
 |e^{-v_j}-m_j|^2\mid{\cal F}_{j-1}]
                              \le\sigma_j^2,                 \tag{2.13}
\]
then
\[
 \mathbb E_\nu[
 |R_j-1|^2\mid{\cal F}_{j-1}]
                              \le{\sigma_j^2\over(1-a_j)^2}. \tag{2.14}
\]
Thus Theorem~\ref{thm:v12v2-martingale} applies when
\[
                              \sum_j{\sigma_j^2\over(1-a_j)^2}
                              <\infty.                        \tag{2.15}
\]
\end{proposition}

\begin{proof}
Condition (2.11) gives \(|m_j|\ge1-a_j>0\), so (2.12) is defined and
\[
 \mathbb E[R_j\mid{\cal F}_{j-1}]
 ={m_j\over m_j}=1.
\]
Moreover,
\[
 R_j-1={e^{-v_j}-m_j\over m_j}.
\]
Take the conditional square and use the lower bound on \(|m_j|\).
\end{proof}

A convenient sufficient estimate for (2.11) is
\[
 \mathbb E_\nu[e^{|v_j|}-1\mid{\cal F}_{j-1}]\le a_j,        \tag{2.16}
\]
because
\[
 |m_j-1|
 \le\mathbb E[|e^{-v_j}-1|\mid{\cal F}_{j-1}]
 \le\mathbb E[e^{|v_j|}-1\mid{\cal F}_{j-1}].
\]
This is local in the shell and can decay with scale even when
\(\sum_j|v_j|\) has no deterministic global bound.

\subsection{Dependency locality}

Suppose the Gaussian reference variables are grouped into independent
shells \(\xi_j\), and
\[
                              v_j=v_j(\xi_j,\xi_{B_j}),       \tag{2.17}
\]
where \(B_j\subset\{1,\ldots,j-1\}\) is a finite boundary
neighbourhood in a dependency graph.

\begin{lemma}[Locality of the conditional counterterm]
\label{lem:v12v2-locality}
Under (2.17),
\[
                              m_j=m_j(\xi_{B_j}).             \tag{2.18}
\]
Hence the counterterm \(\log m_j\), using the branch selected by
(2.11), depends only on the past boundary \(B_j\).
\end{lemma}

\begin{proof}
Condition on all past variables.  Independence makes integration over
the new shell \(\xi_j\) insensitive to past variables outside \(B_j\),
so (2.10) is measurable with respect to \(\xi_{B_j}\).
The disk \(|m_j-1|<1\) avoids zero and is simply connected, giving the
logarithm continued from one.
\end{proof}

Finite dependency degree controls the spatial support of this
counterterm.  It does not by itself prove the numerical decay
(2.15); that estimate must come from scale decay, covariance
localisation, or cancellations in the complete interaction.

\subsection{Relation to the original action}

The replacement
\[
                              e^{-v_j}=m_jR_j                \tag{2.19}
\]
is an exact identity.  Constructing the continuum with density
\(\prod_jR_j\) moves
\[
                              -\sum_j\log m_j                \tag{2.20}
\]
into the running local action.  This is an admissible renewal step
only if the counterterms (2.20) respect gauge symmetry,
diffeomorphism transport, the determinant line, and the selected
physical renormalisation conditions.

\begin{warning}[Ordering and reflection]
\label{warn:v12v2-order}
A one-sided filtration can make (2.20) local while breaking Euclidean
reflection symmetry.  A reflection-compatible construction must add
shells in reflected pairs or prove that the accumulated conditional
counterterm has the full re-treeing and Gram-kernel covariance.
Exact martingale normalisation alone does not imply
Osterwalder--Schrader positivity.
\end{warning}

\subsection{Perturbative estimates}

If \(|v_j|\le b_j\) almost surely, then
\[
 \sigma_j^2
 \le\mathbb E[|e^{-v_j}-1|^2\mid{\cal F}_{j-1}]
 \le(e^{b_j}-1)^2.                                         \tag{2.21}
\]
Thus \(\sum_j b_j^2<\infty\) with \(\sup_jb_j<\infty\) is a simple
sufficient condition for (2.15), after ensuring
\(\sup_ja_j<1\).  More generally, Gaussian exponential moments can
replace the sup bound in (2.21).

\begin{proposition}[Conditional exponential-moment estimate]
\label{prop:v12v2-exponential}
If for some \(p>2\)
\[
 \mathbb E[e^{p|v_j|}\mid{\cal F}_{j-1}]\le M_j             \tag{2.22}
\]
and
\[
 \mathbb E[|v_j|^2e^{2|v_j|}
                 \mid{\cal F}_{j-1}]\le\tau_j^2,            \tag{2.23}
\]
then
\[
                              \sigma_j^2\le\tau_j^2.          \tag{2.24}
\]
In particular (2.15) follows if
\(\sup_ja_j<1\) and \(\sum_j\tau_j^2<\infty\).
\end{proposition}

\begin{proof}
Conditional variance minimises mean-square distance from a
conditional constant, so
\[
 \sigma_j^2
 \le\mathbb E[|e^{-v_j}-1|^2\mid{\cal F}_{j-1}].
\]
The elementary bound
\[
                              |e^{-z}-1|\le|z|e^{|z|}
\]
and (2.23) give (2.24).  Condition (2.22) records the uniform
integrability needed when (2.23) is obtained by Holder estimates.
\end{proof}

\subsection{Application target for the coupled contour}

For the imaginary conformal Gaussian, a shell \(v_j\) must include
all terms generated at that refinement:
\[
 v_j=v_j^{\rm grav}+v_j^{\rm gauge}+v_j^{\rm Higgs}
      +v_j^{\rm det}+v_j^{\rm ghost}+v_j^{\rm Jac}.          \tag{2.25}
\]
The conditional mean in (2.10) must be taken after this recombination.
The YM re-treeing obstruction in the frozen V95 and V100 audits shows
why normalising separate diagrams would not define a covariant
counterterm.

\begin{firewall}[Missing nonlinear estimates]
\label{fire:v12v2-boundary}
No current calculation proves (2.11), (2.13), or (2.23) for the
complete shell (2.25), nor that the conditional counterterms satisfy
the full gauge, diffeomorphism, determinant-line, and reflection
constraints.  V2 supplies an exact nonlinear continuum criterion,
not its SM--Einstein verification.
\end{firewall}

\begin{decision}[V3 acquisition]
\label{dec:v12v2-next}
V3 will derive (2.23) for polynomial shell interactions from Gaussian
hypercontractivity and a finite dependency degree.  It will track the
growth with polynomial degree and shell rank, and will reject any
bound whose constants are extensive in the total cutoff volume.
\end{decision}

\begin{assessment}[Progress after twelfth-cycle V2]
\label{ass:v12v2-status}
Nonlinear complex interactions now have a dimension-local continuum
mechanism.  Zero-free conditional shell factors produce an exactly
normalised multiplicative martingale; square-summable conditional
fluctuations yield an \(L^2\) limit and a countably additive complex
measure with controlled total variation.  The remaining task is to
obtain those fluctuation estimates from the coupled local action.
\end{assessment}

\subsection{Parent record}

This V2 note was formed by appending the preceding material to the
validated twelfth-cycle V1 source, whose record was
\[
\begin{array}{c|c}
\text{lines}&368\\
\text{bytes}&13970\\
\text{SHA--256}&
\texttt{22924F9A64FCC9ABB773021956D18B5A81791B984DE406F8297B11951E0630C8}.
\end{array}                                                     \tag{2.26}
\]

% =============================================================================
% END APPENDED TWELFTH-CYCLE V2 MATERIAL
% =============================================================================
% =============================================================================
% BEGIN APPENDED TWELFTH-CYCLE V3 MATERIAL
% =============================================================================

\section{Gaussian chaos bounds for nonlinear shell fluctuations}
\label{sec:v12v3}

V2 reduced nonlinear continuum construction to conditional estimates
on one shell factor.  Its exponential-moment condition is convenient
but too strong for a generic high-degree Gaussian polynomial.
This version proves the correct hypercontractive tail and a sharper
route to the variance condition that uses only a lower bound on the
real part of the shell action.

Let \(\xi=(\xi_1,\ldots,\xi_m)\) be a standard real Gaussian vector
under a conditional shell law.  A polynomial \(P(\xi)\) has Gaussian
chaos degree at most \(d\) when its Hermite expansion contains no
degree above \(d\).

\begin{theorem}[Gaussian polynomial hypercontractivity]
\label{thm:v12v3-hyper}
If \(P\) has chaos degree at most \(d\), then for every \(p\ge2\),
\[
                 \|P\|_{L^p}
                 \le(p-1)^{d/2}\|P\|_{L^2}.                 \tag{3.1}
\]
The same estimate holds conditionally when the polynomial
coefficients are measurable functions of past variables and the
\(L^p\) norms are taken only over the independent new shell.
\end{theorem}

\begin{proof}
Let \(T_\rho\) be the Ornstein--Uhlenbeck operator.  Gaussian
hypercontractivity states that
\[
                 \|T_\rho f\|_{L^p}\le\|f\|_{L^2}
                 \quad\text{when }\rho\le(p-1)^{-1/2}.       \tag{3.2}
\]
Write \(P=\sum_{k=0}^dP_k\) in homogeneous Wiener chaoses and set
\(Q=\sum_{k=0}^d\rho^{-k}P_k\).  Then \(T_\rho Q=P\).  Orthogonality
of the chaoses gives
\[
 \|Q\|_2^2
 =\sum_{k=0}^d\rho^{-2k}\|P_k\|_2^2
 \le\rho^{-2d}\|P\|_2^2.
\]
Choose \(\rho=(p-1)^{-1/2}\) in (3.2) to obtain (3.1).  With past
variables fixed, the same proof applies to the conditional polynomial;
measurability permits the resulting pointwise conditional estimate.
\end{proof}

\begin{corollary}[Stretched-exponential tail]
\label{cor:v12v3-tail}
There are constants \(c_d,C_d>0\), depending only on \(d\), such that
for \(t\ge C_d\|P\|_2\),
\[
 \mathbb P\{|P|\ge t\}
 \le2\exp\left[-c_d
       \left({t\over\|P\|_2}\right)^{2/d}\right].            \tag{3.3}
\]
Consequently
\[
 \mathbb E\exp\left[
 c'_d\left({|P|\over\|P\|_2}\right)^{2/d}\right]<\infty      \tag{3.4}
\]
for a sufficiently small \(c'_d>0\).
\end{corollary}

\begin{proof}
Markov's inequality and (3.1) give
\[
 \mathbb P\{|P|\ge t\}
 \le\left({(p-1)^{d/2}\|P\|_2\over t}\right)^p.             \tag{3.5}
\]
Choose \(p\) comparable to
\((t/(e\|P\|_2))^{2/d}\), rounded down but at least two.
This yields (3.3), after changing constants.  Integrating the tail
bound gives (3.4) for a smaller exponent constant.
\end{proof}

\begin{warning}[Hypercontractivity is not an ordinary exponential
moment]
\label{warn:v12v3-exponential}
For \(d>2\), (3.4) controls
\(\exp(c|P|^{2/d})\), not \(\exp(c|P|)\).
For example, if \(P(\xi)=\xi^d\), then
\[
 \int_{\mathbb R}e^{c|\xi|^d}e^{-\xi^2/2}\,d\xi=\infty
 \quad(c>0,\ d>2).                                          \tag{3.6}
\]
Thus V2 condition (2.22) cannot be obtained for an arbitrary
higher-degree polynomial by hypercontractivity alone.
\end{warning}

\subsection{One-sided real-part control}

The conditional variance in V2 concerns \(e^{-v_j}\), not
\(e^{|v_j|}\).  A lower bound on the real part is enough.

\begin{lemma}[Exponential increment under a sectorial bound]
\label{lem:v12v3-sectorial}
If \(z\in\mathbb C\) and
\[
                              \operatorname{Re}z\ge-a,
                  \qquad a\ge0,                             \tag{3.7}
\]
then
\[
                              |e^{-z}-1|\le e^a|z|.           \tag{3.8}
\]
\end{lemma}

\begin{proof}
Use
\[
 e^{-z}-1=-z\int_0^1e^{-tz}\,dt.
\]
For \(0\le t\le1\),
\(|e^{-tz}|=e^{-t\operatorname{Re}z}\le e^{ta}\le e^a\).
\end{proof}

\begin{theorem}[Conditional \(L^2\) shell criterion]
\label{thm:v12v3-L2}
Let \(v_j\) be the shell interaction of V2.  Suppose, almost surely
with respect to the past,
\[
                       \operatorname{Re}v_j\ge-a_j           \tag{3.9}
\]
and
\[
                       \mathbb E[|v_j|^2\mid{\cal F}_{j-1}]
                       \le s_j^2.                           \tag{3.10}
\]
Then the conditional fluctuation in (2.13) satisfies
\[
                              \sigma_j^2\le e^{2a_j}s_j^2.   \tag{3.11}
\]
Moreover,
\[
                              |m_j-1|\le e^{a_j}s_j.          \tag{3.12}
\]
Hence the V2 martingale construction applies if
\[
 \sup_j e^{a_j}s_j<1,\qquad
 \sum_j{e^{2a_j}s_j^2\over(1-e^{a_j}s_j)^2}<\infty.          \tag{3.13}
\]
\end{theorem}

\begin{proof}
Conditional variance is no larger than mean-square distance from the
constant one:
\[
 \begin{split}
 \sigma_j^2
 &=\mathbb E[|e^{-v_j}-m_j|^2\mid{\cal F}_{j-1}]\\
 &\le\mathbb E[|e^{-v_j}-1|^2\mid{\cal F}_{j-1}]
 \le e^{2a_j}s_j^2,
 \end{split}
\]
where Lemma~\ref{lem:v12v3-sectorial} was used.  Conditional
Cauchy--Schwarz similarly gives
\[
 |m_j-1|
 \le\mathbb E[|e^{-v_j}-1|\mid{\cal F}_{j-1}]
 \le e^{a_j}s_j.
\]
Insert these estimates into Proposition~\ref{prop:v12v2-local} to obtain (3.13).
\end{proof}

For a purely imaginary polynomial \(v_j=iP_j\), one may take
\(a_j=0\).  A positive real stabilising interaction also has
\(a_j=0\).  The dangerous case is a negative real polynomial direction
whose lower bound worsens with shell rank.

\subsection{Finite dependency graphs}

Let a shell interaction be a sum
\[
                              v_j=\sum_{\alpha\in A_j}P_{j,\alpha},
                                                                    \tag{3.14}
\]
where, conditionally on the past, the centred variables
\(P_{j,\alpha}\) have a dependency graph of maximum degree
\(\Delta_j\).  Thus nonadjacent distinct vertices have zero conditional
covariance.  Suppose
\[
 \mathbb E[P_{j,\alpha}\mid{\cal F}_{j-1}]=0,\qquad
 \mathbb E[|P_{j,\alpha}|^2\mid{\cal F}_{j-1}]
                              \le b_j^2.                     \tag{3.15}
\]

\begin{proposition}[Shell-rank variance bound]
\label{prop:v12v3-dependency}
Under (3.14)--(3.15),
\[
 \mathbb E[|v_j|^2\mid{\cal F}_{j-1}]
 \le |A_j|(\Delta_j+1)b_j^2.                                \tag{3.16}
\]
\end{proposition}

\begin{proof}
Expand the conditional square.  Every diagonal term is at most
\(b_j^2\).  A fixed vertex has at most \(\Delta_j\) neighbours with a
possibly nonzero covariance.  Conditional Cauchy--Schwarz bounds each
such covariance in modulus by \(b_j^2\).  Counting ordered pairs gives
(3.16).
\end{proof}

Combining Theorem~\ref{thm:v12v3-L2} and
Proposition~\ref{prop:v12v3-dependency} yields the sufficient
summability condition
\[
 \sum_j e^{2a_j}|A_j|(\Delta_j+1)b_j^2<\infty,               \tag{3.17}
\]
provided the corresponding local zero-free margins stay below one.
This counts the new shell, not the entire accumulated cutoff volume.

\begin{corollary}[Dyadic exponent test]
\label{cor:v12v3-dyadic}
If
\[
 |A_j|\le C_A2^{\rho j},\qquad
 \Delta_j+1\le C_\Delta2^{\delta j},\qquad
 b_j\le C_b2^{-\sigma j},\qquad
 \sup_ja_j<\infty,                                          \tag{3.18}
\]
then (3.17) holds when
\[
                              2\sigma>\rho+\delta.           \tag{3.19}
\]
\end{corollary}

\begin{proof}
The summand in (3.17) is bounded by a constant times
\(2^{-(2\sigma-\rho-\delta)j}\), which is geometrically summable under
(3.19).
\end{proof}

The exponent \(\rho\) is the number of genuinely new local
interaction cells.  If a purported shell rewrite changes every old
cell, then \(\rho\) becomes the bulk dimension exponent and the gain
is lost.  Exact counterterm cancellation and quasi-local refinement
must first isolate the new layer.

\subsection{Using higher moments for bad events}

Hypercontractivity is still useful when the sectorial bound (3.9)
holds only on a good event.  If \(P\) has degree \(d\), (3.3) gives
\[
 \mathbb P\{|P|>L_j\}
 \le2\exp\left[-c_d
       (L_j/\|P\|_2)^{2/d}\right].                           \tag{3.20}
\]
On the good event one applies (3.8); on its complement the phase chord
is at most two if \(\operatorname{Re}v_j=0\), or requires a separate
amplitude bound otherwise.  A summable choice of the right side of
(3.20) can therefore be added to the V2 defect budget.

\begin{warning}[Complex phase versus negative real action]
\label{warn:v12v3-phase}
For \(v_j=iP_j\), arbitrarily large \(P_j\) does not increase
\(|e^{-v_j}|\), and \(L^2\) smallness is enough for (3.11).
For \(\operatorname{Re}v_j=-|P_j|\), the same polynomial produces
\(e^{|P_j|}\), which is nonintegrable in the generic degree-\(d>2\)
case.  These two situations must not be estimated by the same
absolute-value exponential moment.
\end{warning}

\subsection{Application ledger}

To verify (3.17) for the coupled SM--Einstein shell one must:
\begin{enumerate}
\item express the fully recombined shell action in conditional
      Gaussian chaoses;
\item identify a real-part lower bound before applying phase
      estimates;
\item count only genuinely new cells and their dependency degree;
\item prove the local \(L^2\) coefficient decay \(b_j\);
\item include bad spectral, gauge-slice, and singular-clearance events;
\item preserve reflection by using paired shells or an exact
      half-space factorisation.
\end{enumerate}

\begin{firewall}[No coupled exponent acquired]
\label{fire:v12v3-boundary}
The present sources do not yet provide the values of
\(\rho,\delta,\sigma\) for the fully measured coupled shell, nor a
uniform sectorial lower bound (3.9).  Equation (3.19) is a sharp
bookkeeping gate for the proposed construction, not a verified
SM--Einstein estimate.
\end{firewall}

\begin{decision}[V4 acquisition]
\label{dec:v12v3-next}
V4 will derive \(b_j\) and the shell-count exponent for Wick-ordered
local monomials of a Gaussian field with a dyadic covariance shell.
It will distinguish pointwise variance, integrated cell sums, and
boundary-supported interactions, and will test (3.19) in dimension
four.
\end{decision}

\begin{assessment}[Progress after twelfth-cycle V3]
\label{ass:v12v3-status}
The nonlinear martingale criterion now has the correct Gaussian
input.  Hypercontractivity yields stretched-exponential polynomial
tails, while a one-sided real-part bound converts conditional
\(L^2\) shell size directly into the variance needed by V2.
Finite dependency degree gives the explicit summability threshold
\(2\sigma>\rho+\delta\).
\end{assessment}

\subsection{Parent record}

This V3 note was formed by appending the preceding material to the
validated twelfth-cycle V2 source, whose record was
\[
\begin{array}{c|c}
\text{lines}&708\\
\text{bytes}&25022\\
\text{SHA--256}&
\texttt{5F9DB41B51B0BFA752C367CF68640B61EC7598794A6AD5714E37F159E2918E72}.
\end{array}                                                     \tag{3.21}
\]

% =============================================================================
% END APPENDED TWELFTH-CYCLE V3 MATERIAL
% =============================================================================
% =============================================================================
% BEGIN APPENDED TWELFTH-CYCLE V4 MATERIAL
% =============================================================================

\section{Dyadic Wick-shell variance and the four-dimensional test}
\label{sec:v12v4}

V3 reduces nonlinear shell convergence to the conditional
\(L^2\) size of the fully recombined new interaction.  This version
computes that size for local Wick monomials.  The calculation shows
that a raw four-dimensional quartic shell is not a small martingale
increment; relevant and marginal local terms must first be extracted
by a Wilsonian renewal step.

Let \(\Lambda_j\) be a lattice discretisation of a fixed
\(D\)-dimensional volume with spacing
\[
                              a_j=2^{-j},\qquad
                              |\Lambda_j|\le C_\Lambda2^{Dj}. \tag{4.1}
\]
Let \(\psi_j(x)\) be a centred real Gaussian shell with covariance
\(C_j(x,y)\).  Assume, for every integer \(k\) under consideration,
\[
 \sup_{x\in\Lambda_j}
       \sum_{y\in\Lambda_j}|C_j(x,y)|^k
                              \le S_kv_j^k,                 \tag{4.2}
\]
where \(v_j\) bounds the point variance.  A quasi-local covariance
with correlation length \(O(a_j)\) and uniformly summable rescaled
decay satisfies (4.2).

\begin{lemma}[Wick covariance identity]
\label{lem:v12v4-wick}
For \(k,\ell\ge0\),
\[
 \mathbb E\!\left[
 :\!\psi_j(x)^k\!:\ :\!\psi_j(y)^\ell\!:
 \right]
 ={\bf1}_{k=\ell}\,k!\,C_j(x,y)^k.                          \tag{4.3}
\]
\end{lemma}

\begin{proof}
Use the generating identity
\[
 :e^{t\psi_j(x)}:
 =\exp\left(t\psi_j(x)-{t^2\over2}C_j(x,x)\right).
\]
The expectation of the product at \(x,y\) is
\(\exp(stC_j(x,y))\).  Comparing the coefficients of
\(s^kt^\ell\) gives (4.3).
\end{proof}

\begin{theorem}[Integrated Wick-monomial shell variance]
\label{thm:v12v4-variance}
For
\[
 V_{j,k}=g_{j,k}a_j^D
       \sum_{x\in\Lambda_j}:\!\psi_j(x)^k\!:,                \tag{4.4}
\]
one has
\[
 \mathbb E|V_{j,k}|^2
 \le k!C_\Lambda S_k|g_{j,k}|^2
             2^{-Dj}v_j^k.                                 \tag{4.5}
\]
Different Wick degrees are orthogonal.
\end{theorem}

\begin{proof}
By Lemma~\ref{lem:v12v4-wick},
\[
 \mathbb E|V_{j,k}|^2
 =k!|g_{j,k}|^2a_j^{2D}
       \sum_{x,y\in\Lambda_j}C_j(x,y)^k.
\]
Take absolute values, use (4.2), then (4.1):
\[
 \mathbb E|V_{j,k}|^2
 \le k!|g_{j,k}|^2 2^{-2Dj}
       |\Lambda_j|S_kv_j^k,
\]
which is (4.5).  Orthogonality of distinct degrees is (4.3).
\end{proof}

The factor \(2^{-Dj}\) is the averaging volume of one correlation
cell.  It is already the benefit of finite dependency range; no
additional independent-cell square root should be inserted.

\subsection{Conditional low-field coefficients}

After splitting a field into past and new shell,
a local polynomial contains coefficients depending on the low field:
\[
 V_j=a_j^D\sum_{x\in\Lambda_j}\sum_{k=0}^d
       c_{j,k}(x):\!\psi_j(x)^k\!:.                          \tag{4.6}
\]
Condition on the past, so \(c_{j,k}(x)\) are fixed.

\begin{theorem}[Conditional Schur bound]
\label{thm:v12v4-conditional}
Under (4.2),
\[
 \mathbb E[|V_j-\mathbb E(V_j\mid{\cal F}_{j-1})|^2
                 \mid{\cal F}_{j-1}]
 \le a_j^{2D}\sum_{k=1}^d
       k!S_kv_j^k\sum_{x\in\Lambda_j}|c_{j,k}(x)|^2.         \tag{4.7}
\]
\end{theorem}

\begin{proof}
The \(k=0\) term is conditionally deterministic.  Distinct positive
degrees are orthogonal by (4.3).  For fixed \(k\), let
\(A_{xy}=|C_j(x,y)|^k\).  Its row and column sums are at most
\(S_kv_j^k\).  The Schur test for the quadratic form gives
\[
 \sum_{x,y}|c_{j,k}(x)|A_{xy}|c_{j,k}(y)|
 \le S_kv_j^k\sum_x|c_{j,k}(x)|^2.
\]
Multiply by \(k!a_j^{2D}\) and sum over \(k\).
\end{proof}

The conditional mean is precisely the local term to be placed into
the renewal counterterm before applying the V2 martingale
normalisation.

\subsection{Canonical massless scaling}

For a two-derivative massless scalar shell in \(D>2\), the point
variance has canonical size
\[
                              v_j\le C_v2^{(D-2)j}.           \tag{4.8}
\]
If
\[
                              |g_{j,k}|\le C_g2^{-r_kj},     \tag{4.9}
\]
Theorem~\ref{thm:v12v4-variance} gives
\[
 \mathbb E|V_{j,k}|^2
 \le C_k
 2^{[\,k(D-2)-D-2r_k\,]j}.                                  \tag{4.10}
\]

\begin{corollary}[Dyadic square-summability threshold]
\label{cor:v12v4-threshold}
The pure \(k\)-th Wick shell contributions in (4.4) satisfy
\[
                              \sum_j\mathbb E|V_{j,k}|^2<\infty
                                                                    \tag{4.11}
\]
whenever
\[
                              r_k>{k(D-2)-D\over2}.           \tag{4.12}
\]
\end{corollary}

\begin{proof}
Under (4.12), the exponent in (4.10) is strictly negative, so the
series is geometric.
\end{proof}

In four dimensions, (4.12) becomes
\[
                              r_k>k-2.                       \tag{4.13}
\]
Thus a raw quartic shell requires \(r_4>2\), a raw cubic shell requires
\(r_3>1\), and even the quadratic shell requires \(r_2>0\) for strict
square summability by this direct argument.

\begin{warning}[A marginal coupling is not a small raw shell]
\label{warn:v12v4-marginal}
A dimensionless four-dimensional quartic coupling has \(r_4=0\).
Equation (4.10) then grows like \(2^{4j}\), rather than tending to
zero.  Hypercontractivity cannot repair this variance growth.
The V2 factors cannot be the unprocessed quartic action at every
scale.
\end{warning}

This does not say that a renormalised quartic theory is impossible.
It says that the relevant and marginal local pieces must be absorbed
into running finite-dimensional data before the remaining shell
factor is tested by V2.

\subsection{Boundary and defect support}

Let an interaction be supported on a codimension-\(c\) lattice layer
\(B_j\) with
\[
                              |B_j|\le C_B2^{(D-c)j},         \tag{4.14}
\]
and use its natural cell measure \(a_j^{D-c}\):
\[
 V_{j,k}^{B}=g_{j,k}a_j^{D-c}
        \sum_{x\in B_j}:\!\psi_j(x)^k\!:.                    \tag{4.15}
\]
The same proof gives
\[
 \mathbb E|V_{j,k}^{B}|^2
 \le C_{k,B}|g_{j,k}|^2
       2^{-(D-c)j}v_j^k.                                    \tag{4.16}
\]
With (4.8)--(4.9), square summability requires
\[
                              r_k>{k(D-2)-(D-c)\over2}.       \tag{4.17}
\]
The exponent is larger by \(c/2\) because (4.15) is an interaction
integrated in the lower-dimensional natural measure.  A boundary
remainder inherited from a bulk action often carries additional
powers of \(a_j\); those powers must be included in \(r_k\).

\subsection{Wilsonian extraction}

Write the fully recombined shell action schematically as
\[
                              v_j={\cal L}_j+{\cal R}_j,      \tag{4.18}
\]
where \({\cal L}_j\) lies in a declared finite local basis containing
all relevant and marginal operators, and \({\cal R}_j\) is the
irrelevant remainder.

\begin{proposition}[Required remainder exponent]
\label{prop:v12v4-remainder}
Suppose every Wick component of \({\cal R}_j\) has coefficient decay
\(r_k\) satisfying (4.12), its real part has the lower bound required
in V3, and only finitely many tensor species occur with uniform
constants.  Then the conditional variances of \({\cal R}_j\) are
square summable.  The V2 martingale theorem applies after the local
conditional mean and the finite basis \({\cal L}_j\) have been
incorporated into the running action.
\end{proposition}

\begin{proof}
Apply (4.10) to each Wick component.  Each resulting geometric series
converges by (4.12), and a finite sum remains convergent.  The V3
sectorial estimate converts the variance into the V2 shell-factor
bound.
\end{proof}

The proposition is conditional on a stable finite local basis.
Operator mixing can regenerate relevant or marginal terms, so their
coefficients must be updated at every renewal step.

\subsection{Gauge, fermion, and metric complications}

The scalar identity (4.3) extends componentwise to Gaussian tensor
fields after contracting their covariance tensors.  Fermions use
Grassmann rather than Gaussian hypercontractivity and should be
integrated into determinant or Pfaffian carriers before the complex
shell variance is formed.  Gauge and metric propagators also contain
longitudinal or null directions; their covariance shell must be
defined on the quotient or split by the low/high Feshbach construction
recorded in the frozen notes.

For the imaginary conformal field \(z=i\psi\), the positive amplitude
is the real Gaussian law of \(\psi\).  Wick variances therefore use
that positive covariance; powers of \(i\) enter the complex
coefficients \(c_{j,k}\) and do not change (4.7).

\begin{ledger}[Data needed for the first coupled shell]
\label{led:v12v4-data}
A concrete calculation must supply:
\begin{enumerate}
\item the dyadic covariance estimate (4.2) for every retained field;
\item the local Wick expansion after determinant, ghost, coordinate,
      and quotient recombination;
\item the finite relevant/marginal basis \({\cal L}_j\);
\item coefficient decay for the remainder \({\cal R}_j\);
\item the real-part lower bound and bad-event probability;
\item reflection-paired rather than one-sided shell extraction.
\end{enumerate}
\end{ledger}

\begin{firewall}[Power counting is not an RG construction]
\label{fire:v12v4-boundary}
Equations (4.10)--(4.17) are rigorous variance estimates under the
stated covariance bounds.  The current programme has not constructed
the coupled Wilsonian map that produces a remainder satisfying
(4.12), nor proved stability of its finite local basis and reflection
factorisation.  The four-dimensional quartic warning is an
obstruction to the raw-shell approach, not a proof of a continuum
interaction.
\end{firewall}

\begin{decision}[Mandatory V5 audit]
\label{dec:v12v4-next}
V5 will perform the compulsory YM/NS companion audit.  The following
substantive version will go upstream to the finite local
renormalisation map required by Proposition~\ref{prop:v12v4-remainder},
using an orthogonal projection onto relevant and marginal Wick
operators and a contractive estimate on the irrelevant remainder.
\end{decision}

\begin{assessment}[Progress after twelfth-cycle V4]
\label{ass:v12v4-status}
The shell-size exponent is now explicit.  Integrated Wick monomials
have variance \(O(2^{[k(D-2)-D-2r_k]j})\); in four dimensions the raw
quartic interaction grows and cannot be a V2 martingale increment.
The continuum route must first extract running local operators and
apply the square-summability test only to the irrelevant remainder.
\end{assessment}

\subsection{Parent record}

This V4 note was formed by appending the preceding material to the
validated twelfth-cycle V3 source, whose record was
\[
\begin{array}{c|c}
\text{lines}&1035\\
\text{bytes}&36164\\
\text{SHA--256}&
\texttt{FE721F853EEEE4F8BFD4EBD8276F532B4E86AA82D3D94A023C34E314DB64439A}.
\end{array}                                                     \tag{4.19}
\]

% =============================================================================
% END APPENDED TWELFTH-CYCLE V4 MATERIAL
% =============================================================================
% =============================================================================
% BEGIN APPENDED TWELFTH-CYCLE V5 MATERIAL
% =============================================================================

\section{Companion audit at V5: corrected YM channel transport}
\label{sec:v12v5}

The compulsory audit records
\[
\begin{array}{c|r|r|l}
\text{file}&\text{lines}&\text{bytes}&\text{SHA--256}\\ \hline
\texttt{Renewal\_Yang\_Mills\_Mass\_Gap\_Research\_Notes\_V42.tex}
&17164&586899&
\texttt{008480F2DD4049D17ABE7DB62E49C99E3E497022BF7DA022D2856362181605DC}\\
\texttt{Renewal\_NS\_Stability\_Research\_Notes\_V26.tex}
&5319&278283&
\texttt{82C887F8C2C69E4F28FAD7C3DC8DE5B9099CB5A4BF431EBA1FFF3E91935978AD}.
\end{array}                                                     \tag{5.1}
\]

YM V42 supersedes the interpretation, but not the exact arithmetic,
of YM V41.  The retained boundary list was ordered by lattice
position while the response code treated its slot as the physical
direction.  The actual slot order is the reverse physical-axis order.
Consequently the V41 corner pair compared different transported
polarisation channels and did not diagnose a missing density term.

V42 constructs the exact spectator re-tree map.  At all sixteen
corners it has block form
\[
                              {\cal U}=
 \begin{pmatrix}S&0\\0&P\end{pmatrix},                       \tag{5.2}
\]
with invertible vertical block \(S\), the correct physical boundary
permutation \(P\), exact involution under the paired corner, and
Hessian congruence
\[
                              A(k)=S(k)^TA(\tau k)S(k).       \tag{5.3}
\]
It also verifies covariance of every routed toron channel.  The next
YM calculation will repeat the corner test with the corrected channel
stabiliser.

\begin{corollary}[Correction to the frozen V95 reading]
\label{cor:v12v5-correction}
The general frozen conclusion that a physical symmetry must transport
the complete coordinate, channel, and density system remains valid.
The specific V41 residual is not evidence for a nontrivial quotient
Jacobian: V42 proves that its first dichotomy branch, channel
transport, explains the comparison.
\end{corollary}

\begin{proof}
Equation (5.2) has zero vertical--boundary off-diagonal blocks, while
the physical boundary permutation moves the slot used in V41.
Therefore the fixed-slot subtraction was not a covariance comparison.
Equation (5.3) supplies the correctly transported Hessian identity.
\end{proof}

For the present SM--Einstein programme, the useful import is positive:
the local operator basis and its renormalisation projection must carry
a representation of the physical re-treeing group.  Coordinate slots
cannot be declared invariant before the exact channel map is applied.

The NS V26 file is unchanged and supplies no new Gaussian shell,
renormalisation, gauge, or gravitational input.

\begin{decision}[V6 equivariant local extraction]
\label{dec:v12v5-next}
V6 will define the relevant/marginal local subspace as an invariant
subspace of the exact symmetry and re-treeing representation.  It
will prove that the weighted orthogonal projection commutes with that
representation and combine it with a contractive recursion for the
irrelevant remainder, producing the coefficient decay required by
V4 when the spectral contraction beats the canonical shell exponent.
\end{decision}

\begin{firewall}[Audit boundary]
\label{fire:v12v5-boundary}
YM V42 corrects a finite channel label and proves corner re-tree
covariance.  It does not supply the coupled Wilsonian map, the
irrelevant-remainder contraction, continuum measure, reflection
kernel, or SM--Einstein universality data.
\end{firewall}

\begin{assessment}[Progress after twelfth-cycle V5]
\label{ass:v12v5-status}
The required audit is complete and prevents a false density anomaly
from entering the SM programme.  The next local renormalisation step
will be formulated in physical equivariant coordinates, with channel
transport included in the representation.
\end{assessment}

\subsection{Parent record}

This V5 note was formed by appending the preceding material to the
validated twelfth-cycle V4 source, whose record was
\[
\begin{array}{c|c}
\text{lines}&1353\\
\text{bytes}&47373\\
\text{SHA--256}&
\texttt{06B4FE513E3F54222D8DB357ECD420243A86A9A0368BFA8B840B42640FC3EA57}.
\end{array}                                                     \tag{5.4}
\]

% =============================================================================
% END APPENDED TWELFTH-CYCLE V5 MATERIAL
% =============================================================================
% =============================================================================
% BEGIN APPENDED TWELFTH-CYCLE V6 MATERIAL
% =============================================================================

\section{Equivariant local extraction and irrelevant contraction}
\label{sec:v12v6}

V4 shows that the raw relevant and marginal shell terms are too large
for the V2 martingale.  They must be projected into running local
couplings while the remainder contracts.  The V5 audit adds a
structural requirement: this projection must commute with physical
channel transport and re-treeing.

Let \({\cal H}_{\rm loc}\) be a Hilbert space of cutoff-local
interaction densities, including their tensor and channel labels.
Let a finite or compact group \(G\) act continuously through
\(U_g\).  The action includes physical lattice symmetries together
with the exact coordinate re-treeing needed to return to the chosen
slice.

\begin{lemma}[Invariant inner product]
\label{lem:v12v6-average}
Starting from any positive inner product
\(\langle\cdot,\cdot\rangle_0\), define
\[
 \langle x,y\rangle_G
 =\int_G\langle U_gx,U_gy\rangle_0\,dg,                     \tag{6.1}
\]
where \(dg\) is normalised Haar measure.  If every \(U_g\) is
invertible, then (6.1) is positive definite and the representation is
unitary for \(\langle\cdot,\cdot\rangle_G\).
\end{lemma}

\begin{proof}
Positivity follows because the integrand is nonnegative and, for
\(x\ne0\), it is positive at every \(g\) when the starting inner
product is positive and \(U_g\) is invertible.  Haar invariance gives
\[
 \langle U_hx,U_hy\rangle_G
 =\int_G\langle U_{gh}x,U_{gh}y\rangle_0\,dg
 =\langle x,y\rangle_G.
\]
\end{proof}

Let \({\cal L}\subset{\cal H}_{\rm loc}\) be the finite-dimensional
space spanned by the declared relevant and marginal local operators.
Assume
\[
                              U_g{\cal L}={\cal L}
                         \quad(g\in G).                      \tag{6.2}
\]
Let \(\Pi\) be the orthogonal projection onto \({\cal L}\) in the
inner product (6.1), and put \(Q=I-\Pi\).

\begin{theorem}[Equivariant local projection]
\label{thm:v12v6-projection}
Under (6.2),
\[
                              \Pi U_g=U_g\Pi,\qquad
                              QU_g=U_gQ.                     \tag{6.3}
\]
Consequently the extracted local action and the irrelevant remainder
transform in their correct physical representations.
\end{theorem}

\begin{proof}
Unitarity and (6.2) imply that both \({\cal L}\) and
\({\cal L}^{\perp_G}\) are invariant under every \(U_g\).
The orthogonal decomposition
\(x=\Pi x+Qx\) is therefore carried to the orthogonal decomposition
of \(U_gx\), which proves (6.3).
\end{proof}

A fixed coordinate-slot projection need not satisfy (6.2).  YM V42
provides an exact example: the physical spectator symmetry permutes
the retained slots, even though a fixed numerical slot was previously
treated as invariant.

\subsection{Conditioning of the finite local basis}

Let \(O_1,\ldots,O_m\) span \({\cal L}\) and set
\[
                              G_{ab}=\langle O_a,O_b\rangle_G. \tag{6.4}
\]
Then
\[
                 \Pi V=\sum_{a,b=1}^m
                       O_a(G^{-1})_{ab}\langle O_b,V\rangle_G.
                                                                    \tag{6.5}
\]

\begin{proposition}[Uniform coefficient extraction]
\label{prop:v12v6-gram}
If
\[
                              \lambda_{\min}(G)\ge g_0>0     \tag{6.6}
\]
uniformly in the cutoff, then the coefficient vector \(c(V)\) in
(6.5) obeys
\[
                              \|c(V)\|_{\ell^2}
                              \le g_0^{-1/2}\|\Pi V\|_G
                              \le g_0^{-1/2}\|V\|_G.          \tag{6.7}
\]
\end{proposition}

\begin{proof}
If \(\Pi V=\sum_ac_aO_a\), then
\[
 \|\Pi V\|_G^2=c^*Gc\ge g_0\|c\|_{\ell^2}^2.
\]
Orthogonal projection is contractive, giving (6.7).
\end{proof}

A nearly dependent local basis can turn a small density remainder
into large running coefficients.  The Gram lower bound is therefore
part of the uniform renormalisation data.

\subsection{Equivariant renewal map}

Let \({\cal T}_j\) be one exact shell-integration and rescaling map.
Write
\[
                              V_j=\ell_j+r_j,\qquad
 \ell_j=\Pi V_j,\quad r_j=QV_j.                              \tag{6.8}
\]
Define
\[
 \ell_{j+1}=\Pi{\cal T}_j(\ell_j+r_j),\qquad
 r_{j+1}=Q{\cal T}_j(\ell_j+r_j).                            \tag{6.9}
\]

\begin{proposition}[Symmetry preservation]
\label{prop:v12v6-symmetry}
If
\[
                        {\cal T}_j(U_gV)=U_g{\cal T}_j(V)    \tag{6.10}
\]
and \(V_j\) is \(G\)-invariant, then \(\ell_{j+1}\) and
\(r_{j+1}\) are \(G\)-invariant.  More generally, every
representation channel is preserved by the decomposition (6.9).
\end{proposition}

\begin{proof}
Combine (6.3) and (6.10).  For example,
\[
 U_g\ell_{j+1}
 =\Pi U_g{\cal T}_j(V_j)
 =\Pi{\cal T}_j(U_gV_j)=\ell_{j+1}.
\]
The same calculation applies to \(Q\) and to covariant channels.
\end{proof}

Exact equivariance must be checked for the complete density.  An
approximate commutator \(\Pi U_g-U_g\Pi\) enters the forcing error
below and cannot be relabelled as a physical anomaly without first
checking channel transport.

\subsection{Linear contraction with forcing}

Suppose the irrelevant recursion has the form
\[
                              r_{j+1}=A_jr_j+e_j,\qquad
                              \|A_j\|\le\kappa<1.             \tag{6.11}
\]

\begin{theorem}[Forced irrelevant decay]
\label{thm:v12v6-linear}
If
\[
                              \|e_j\|\le E2^{-\alpha j},     \tag{6.12}
\]
then, for every
\[
                              0<\sigma<
                   \min\{\alpha,-\log_2\kappa\},             \tag{6.13}
\]
there is \(C_\sigma<\infty\) such that
\[
                              \|r_j\|\le C_\sigma2^{-\sigma j}.
                                                                    \tag{6.14}
\]
\end{theorem}

\begin{proof}
Iteration gives
\[
 \|r_j\|
 \le\kappa^j\|r_0\|
       +E\sum_{n=0}^{j-1}\kappa^{j-1-n}2^{-\alpha n}.       \tag{6.15}
\]
Choose \(\sigma\) as in (6.13).  Then
\(\kappa2^\sigma<1\) and \(2^{-(\alpha-\sigma)}<1\).
After multiplying (6.15) by \(2^{\sigma j}\), the convolution is
bounded by the convolution of two summable geometric sequences,
uniformly in \(j\).  This proves (6.14).
\end{proof}

\subsection{Nonlinear contraction}

The exact renewal map also has a nonlinear irrelevant term.  Assume
\[
 r_{j+1}=A_jr_j+N_j(r_j)+e_j,\qquad
 \|N_j(r)\|\le C_N\|r\|^2                                  \tag{6.16}
\]
inside a fixed ball.

\begin{theorem}[Small nonlinear stable cone]
\label{thm:v12v6-nonlinear}
Choose \(\sigma\) satisfying (6.13).  If there is \(M>0\) such that
\[
             \kappa M+C_NM^2+E\le M2^{-\sigma},             \tag{6.17}
\]
and
\[
                              \|r_0\|\le M,                  \tag{6.18}
\]
then
\[
                              \|r_j\|\le M2^{-\sigma j}
                         \quad(j\ge0).                       \tag{6.19}
\]
\end{theorem}

\begin{proof}
Assume (6.19) at \(j\).  Since \(\alpha>\sigma\),
\[
 \begin{split}
 \|r_{j+1}\|
 &\le\kappa M2^{-\sigma j}
       +C_NM^22^{-2\sigma j}+E2^{-\alpha j}\\
 &\le2^{-\sigma j}
       (\kappa M+C_NM^2+E)\\
 &\le M2^{-\sigma(j+1)}
 \end{split}
\]
by (6.17).  Induction starting from (6.18) proves the claim.
\end{proof}

Condition (6.17) is feasible when
\(\kappa<2^{-\sigma}\) and the initial remainder and forcing are
sufficiently small.  The local coupling vector \(\ell_j\) may be
large or marginally running, but the constants in (6.16) must remain
uniform along its permitted trajectory.

\subsection{Comparison with the Wick threshold}

Suppose the norm in (6.19) controls a degree-\(k\) remainder
coefficient \(g_{j,k}\).  Then V4 applies with \(r_k=\sigma\).

\begin{corollary}[Four-dimensional contraction gate]
\label{cor:v12v6-four}
In \(D=4\), the degree-\(k\) Wick remainder is square summable by V4
if one can choose
\[
 k-2<\sigma<
 \min\{\alpha,-\log_2\kappa\}.                              \tag{6.20}
\]
A sufficient spectral and forcing condition is
\[
                              \alpha>k-2,\qquad
                              \kappa<2^{-(k-2)},              \tag{6.21}
\]
with strict room for a \(\sigma\) and smallness as in (6.17).
For a quartic remainder this requires \(\kappa<1/4\) in the norm and
rescaling convention used here.
\end{corollary}

\begin{proof}
Choose \(\sigma\) between the two sides of (6.20).
Theorem~\ref{thm:v12v6-nonlinear} gives coefficient decay
\(2^{-\sigma j}\), and V4 condition \(r_k>k-2\) then gives
square summability.  Conditions (6.21) guarantee that such a
\(\sigma\) exists.
\end{proof}

The numerical threshold is norm-dependent.  Rescaling canonical
powers into the definition of \(r_j\) changes both \(\kappa\) and the
V4 exponent but not the invariant requirement that the renormalised
remainder beat its shell variance growth.

\subsection{Ward-compatible local space}

Group invariance alone does not identify the allowed counterterms.
Let \(s\) be the linearised BRST or combined gauge--diffeomorphism
Ward differential on local functionals.  A physical local basis
should be chosen in
\[
                              \ker s
\]
modulo \(s\)-exact terms and declared total derivatives.  For the
projection to preserve the Ward complex, one needs
\[
                              s\Pi=\Pi s                     \tag{6.22}
\]
on the regulated local space, or an explicitly controlled commutator.

\begin{warning}[Symmetry averaging is not anomaly cancellation]
\label{warn:v12v6-anomaly}
The group average (6.1) makes a finite symmetry representation
unitary.  It does not prove nilpotence of the regulated Ward
differential, cancel a local anomaly, or trivialise a global
determinant line.  Those are independent cohomological inputs.
\end{warning}

\subsection{Acquisition ledger}

A concrete coupled renewal step must now provide:
\begin{enumerate}
\item the exact complete-density map \({\cal T}_j\);
\item a finite local operator space closed under physical channel
      transport and the Ward differential;
\item a uniformly conditioned Gram matrix;
\item an irrelevant linear contraction \(\kappa\);
\item a quadratic or Lipschitz nonlinear remainder bound;
\item a decaying forcing term containing discretisation, contour,
      slice-atlas, and determinant errors;
\item the inequality (6.20) for every Wick degree left in the
      irrelevant remainder.
\end{enumerate}

\begin{firewall}[No coupled contraction acquired]
\label{fire:v12v6-boundary}
V6 proves the equivariant projection and stable-cone implications.
The present manuscripts do not yet construct the coupled
renormalisation map or establish numerical constants satisfying
(6.17) and (6.20).  In particular, the finite relevant/marginal
SM--Einstein Ward cohomology has not been closed under the proposed
contour and quotient regulator.
\end{firewall}

\begin{decision}[V7 acquisition]
\label{dec:v12v6-next}
V7 will analyse the Ward-compatible projection in a finite Hilbert
complex.  It will use a Hodge decomposition and a spectral gap of the
Ward Laplacian to separate physical cohomology, exact gauge
directions, and anomaly obstructions, with perturbation bounds that
can be inserted into the forcing term of (6.16).
\end{decision}

\begin{assessment}[Progress after twelfth-cycle V6]
\label{ass:v12v6-status}
The Wilsonian extraction required by V4 now has a rigorous abstract
mechanism.  A group-averaged orthogonal projection preserves exact
re-treeing and channel covariance, while a small nonlinear stable
cone gives explicit irrelevant decay.  The four-dimensional gate
(6.20) states exactly how much contraction is needed before the V2
complex martingale can be applied.
\end{assessment}

\subsection{Parent record}

This V6 note was formed by appending the preceding material to the
validated twelfth-cycle V5 source, whose record was
\[
\begin{array}{c|c}
\text{lines}&1462\\
\text{bytes}&51979\\
\text{SHA--256}&
\texttt{D01003C660846AB1AEE88EBA7DE0806351E2250118B996D2F89D3F5ADE1D8F99}.
\end{array}                                                     \tag{6.23}
\]

% =============================================================================
% END APPENDED TWELFTH-CYCLE V6 MATERIAL
% =============================================================================
% =============================================================================
% BEGIN APPENDED TWELFTH-CYCLE V7 MATERIAL
% =============================================================================

\section{Finite Ward complexes and quantitative Hodge projection}
\label{sec:v12v7}

V6 requires a local projection compatible with the regulated Ward or
BRST differential.  This version gives the finite-dimensional Hodge
decomposition, identifies the harmonic anomaly obstruction, and
bounds the counterterm needed to remove an exact defect.

Let
\[
 \cdots\longrightarrow C^{q-1}
 \mathop{\longrightarrow}^{d_{q-1}} C^q
 \mathop{\longrightarrow}^{d_q} C^{q+1}
 \longrightarrow\cdots                                      \tag{7.1}
\]
be a finite-dimensional Hilbert complex:
\[
                              d_qd_{q-1}=0.                  \tag{7.2}
\]
The degree-\(q\) Ward Laplacian is
\[
 \Delta_q=d_{q-1}d_{q-1}^*+d_q^*d_q.                       \tag{7.3}
\]
Write \(H_q\) for the orthogonal projection onto
\({\cal H}^q=\ker\Delta_q\).

\begin{theorem}[Finite Hodge decomposition]
\label{thm:v12v7-hodge}
One has the orthogonal decomposition
\[
 C^q=\operatorname{im}d_{q-1}
       \mathbin{\oplus^\perp}{\cal H}^q
       \mathbin{\oplus^\perp}\operatorname{im}d_q^*,         \tag{7.4}
\]
with
\[
 {\cal H}^q=\ker d_q\cap\ker d_{q-1}^*.                     \tag{7.5}
\]
Every cohomology class in
\(\ker d_q/\operatorname{im}d_{q-1}\) has a unique harmonic
representative.
\end{theorem}

\begin{proof}
Equation (7.3) gives
\[
 \langle x,\Delta_qx\rangle
 =\|d_{q-1}^*x\|^2+\|d_qx\|^2,
\]
which proves (7.5).  Nilpotence makes
\(\operatorname{im}d_{q-1}\) orthogonal to
\(\operatorname{im}d_q^*\).  Their joint orthogonal complement is
\[
 \ker d_{q-1}^*\cap\ker d_q={\cal H}^q,
\]
proving (7.4).  If a harmonic vector is exact, it is orthogonal to
itself and hence zero.  Decomposition of a closed vector has no
coexact component, so every cohomology class has exactly one harmonic
representative.
\end{proof}

At ghost number zero, \({\cal H}^0\) represents physical local
couplings after exact redundancies have been removed.  At ghost number
one, \({\cal H}^1\) is the finite-cutoff anomaly space.

\subsection{Spectral gap and Green operator}

Assume the positive spectrum of \(\Delta_q\) obeys
\[
 \Delta_q\ge\gamma_q(I-H_q),\qquad \gamma_q>0.              \tag{7.6}
\]
Define the Green operator
\[
 G_q=\Delta_q^{-1}(I-H_q).                                  \tag{7.7}
\]
Then
\[
                              \|G_q\|\le\gamma_q^{-1}.        \tag{7.8}
\]

\begin{theorem}[Exact anomaly removal with norm control]
\label{thm:v12v7-anomaly}
Let \(a\in C^1\) satisfy the Wess--Zumino consistency condition
\[
                              d_1a=0.                        \tag{7.9}
\]
Then
\[
                              a=H_1a+d_0b,\qquad
                              b=d_0^*G_1a.                  \tag{7.10}
\]
In particular, \(a\) is removable by a local counterterm precisely
when
\[
                              H_1a=0.                        \tag{7.11}
\]
For the minimal choice in (7.10),
\[
                              \|b\|\le\gamma_1^{-1/2}\|a\|.  \tag{7.12}
\]
\end{theorem}

\begin{proof}
The Hodge decomposition of the closed vector \(a\) contains only its
harmonic and exact parts.  Since \(\Delta_1G_1=I-H_1\),
\[
 (I-H_1)a=d_0d_0^*G_1a+d_1^*d_1G_1a.
\]
The second term vanishes because Green operators commute with the
differential on the positive spectral subspaces and \(d_1a=0\).
This proves (7.10).  Exact and harmonic subspaces are orthogonal, so
(7.11) is necessary and sufficient.

Finally,
\[
 \begin{split}
 \|d_0^*G_1a\|^2
 &\le\langle G_1a,\Delta_1G_1a\rangle\\
 &=\langle G_1a,(I-H_1)a\rangle
 \le\gamma_1^{-1}\|a\|^2,
 \end{split}
\]
which proves (7.12).
\end{proof}

The harmonic part cannot be removed by changing a local counterterm
inside the same complex.  For the Standard Model it must cancel after
all fermion representations and mixed gauge--gravitational
contributions are combined.

\subsection{Failure of consistency}

For a general \(a\in C^1\), Hodge decomposition gives
\[
 a=H_1a+d_0d_0^*G_1a+d_1^*d_1G_1a.                        \tag{7.13}
\]
The final coexact term measures failure of (7.9).

\begin{proposition}[Consistency-defect bound]
\label{prop:v12v7-consistency}
Assume (7.6) at degrees one and two.  Then
\[
 d_1^*d_1G_1a=d_1^*G_2(d_1a),                              \tag{7.14}
\]
and
\[
 \|d_1^*G_2(d_1a)\|
                              \le\gamma_2^{-1/2}\|d_1a\|.    \tag{7.15}
\]
\end{proposition}

\begin{proof}
The intertwining identity
\(d_1\Delta_1=\Delta_2d_1\) follows from (7.2)--(7.3).
Functional calculus on the positive spectral subspaces gives
\(d_1G_1=G_2d_1\), proving (7.14).  The estimate is the analogue of
(7.12):
\[
 \|d_1^*G_2r\|^2
 \le\langle G_2r,\Delta_2G_2r\rangle
 \le\gamma_2^{-1}\|r\|^2.
\]
\end{proof}

Thus a cutoff Ward defect has three distinct pieces:
\[
 \boxed{
 a=\underbrace{H_1a}_{\text{cohomological anomaly}}
 +\underbrace{d_0d_0^*G_1a}_{\text{removable counterterm}}
 +\underbrace{d_1^*G_2d_1a}_{\text{consistency failure}}.}   \tag{7.16}
\]
They must not be combined into one residual norm.

\subsection{Compatibility with physical symmetries}

Suppose the finite symmetry and re-treeing representation of V6 is
unitary and
\[
                              d_qU_g=U_gd_q.                 \tag{7.17}
\]

\begin{theorem}[Equivariance of Hodge data]
\label{thm:v12v7-equivariant}
Under (7.17),
\[
 U_g\Delta_q=\Delta_qU_g,\qquad
 U_gH_q=H_qU_g,\qquad
 U_gG_q=G_qU_g.                                             \tag{7.18}
\]
Hence the counterterm \(b=d_0^*G_1a\) transforms covariantly whenever
the anomaly carrier \(a\) does.
\end{theorem}

\begin{proof}
Taking adjoints in (7.17) and using unitarity gives the corresponding
identity for \(d_q^*\).  Equation (7.3) then commutes with \(U_g\).
Its kernel projection and its inverse on the orthogonal complement
are spectral functions of \(\Delta_q\), proving (7.18).
\end{proof}

The corrected channel permutation in YM V42 is part of \(U_g\).
Omitting it invalidates (7.17) even when the coordinate formula for
the differential appears unchanged.

\subsection{A Ward-compatible local renormalisation projection}

Let \({\cal L}^q\subset C^q\) be the finite relevant/marginal local
subspace at degree \(q\), with orthogonal projection \(\Pi_q\).
Assume it is a reducing subcomplex:
\[
 d_q\Pi_q=\Pi_{q+1}d_q,\qquad
 d_{q-1}^*\Pi_q=\Pi_{q-1}d_{q-1}^*.                         \tag{7.19}
\]

\begin{theorem}[Commuting local and Hodge projections]
\label{thm:v12v7-local}
Under (7.19),
\[
 \Pi_q\Delta_q=\Delta_q\Pi_q,\qquad
 \Pi_qH_q=H_q\Pi_q,\qquad
 \Pi_qG_q=G_q\Pi_q.                                        \tag{7.20}
\]
The local extraction therefore separates physical harmonic
couplings, exact counterterms, and consistency defects without mixing
them.
\end{theorem}

\begin{proof}
Insert (7.19) and its adjacent-degree versions into (7.3) to obtain
commutation with \(\Delta_q\).  Spectral calculus gives the remaining
identities.
\end{proof}

This is stronger than asking \({\cal L}^0\subset\ker d_0\).
The full local basis across ghost degrees must close under the
differential and its adjoint in the chosen regulator inner product.

\subsection{Stability of the harmonic projection}

Let \(\Delta\) be a Ward Laplacian with zero eigenspace projection
\(H\) and positive spectral gap \(\gamma\).  Let
\[
                              \widetilde\Delta=\Delta+E      \tag{7.21}
\]
be another positive Ward Laplacian, assume
\[
                              \|E\|=\varepsilon<{\gamma\over2},\tag{7.22}
\]
and assume \(\dim\ker\widetilde\Delta=\dim\ker\Delta\).

\begin{theorem}[Hodge-projector perturbation]
\label{thm:v12v7-perturb}
The positive spectrum of \(\widetilde\Delta\) lies above
\(\gamma-\varepsilon\), and its harmonic projection
\(\widetilde H\) satisfies
\[
                              \|\widetilde H-H\|
                              \le{2\varepsilon\over\gamma-2\varepsilon}.
                                                                    \tag{7.23}
\]
\end{theorem}

\begin{proof}
The min--max principle and equality of the nullities give the lower
positive eigenvalue bound.  Let \({\cal C}\) be the circle of radius
\(\gamma/2\) around zero.  It separates the zero spectrum of both
operators from their positive spectra.  The Riesz formula and the
resolvent identity give
\[
 \widetilde H-H
 ={1\over2\pi i}\int_{\cal C}
   (\zeta-\widetilde\Delta)^{-1}
       E(\zeta-\Delta)^{-1}\,d\zeta.
\]
On \({\cal C}\), the two resolvent norms are at most
\(2/(\gamma-2\varepsilon)\) and \(2/\gamma\), respectively.
Multiplying by the contour length \(\pi\gamma\) proves (7.23).
\end{proof}

The nullity assumption is essential.  If the regulated differential
changes cohomology dimension, an anomaly class can appear or
disappear through a zero eigenvalue, and no small projector estimate
with fixed rank applies.

\subsection{Insertion into the irrelevant forcing}

Let \(a_j\) be the degree-one Ward defect produced by one renewal
step.  After subtracting its exact part, V6 forcing receives at least
\[
 e_j^{\rm Ward}
 =H_{1,j}a_j+d_{1,j}^*G_{2,j}d_{1,j}a_j
             +e_j^{\rm proj},                               \tag{7.24}
\]
where \(e_j^{\rm proj}\) contains the drift of the local/Hodge
projectors.  Under uniform gaps,
\[
 \|e_j^{\rm Ward}\|
 \le\|H_{1,j}a_j\|
    +\gamma_{2,j}^{-1/2}\|d_{1,j}a_j\|
    +\|e_j^{\rm proj}\|.                                    \tag{7.25}
\]
For the V6 contraction, every term on the right must decay faster
than the target irrelevant exponent.  Exact cancellation of the
harmonic term is preferable to estimating it: a persistent anomaly
is not an irrelevant forcing.

\begin{warning}[Approximate nilpotence]
\label{warn:v12v7-nilpotence}
If \(d_{q+1}d_q\ne0\), the sequence is not a complex and
Theorem~\ref{thm:v12v7-hodge} does not apply.  One must either repair
the regulator so nilpotence is exact or enlarge the complex so the
defect is an explicit curvature insertion.  Calling the resulting
small eigenvectors cohomology would be circular.
\end{warning}

\begin{ledger}[Ward acquisition data]
\label{led:v12v7-data}
A cutoff implementation must provide:
\begin{enumerate}
\item an exactly nilpotent combined gauge--diffeomorphism
      differential;
\item a positive inner product including ghost, channel, and quotient
      transport;
\item a uniform positive Ward-Laplacian gap off cohomology;
\item the harmonic anomaly coefficients for the complete SM
      representation;
\item a reducing relevant/marginal local subcomplex;
\item adjacent bounds on the Hodge and local projections;
\item decay of the remaining consistency defect sufficient for V6.
\end{enumerate}
\end{ledger}

\begin{firewall}[No regulator Ward complex acquired]
\label{fire:v12v7-boundary}
V7 proves the finite-complex implications but does not construct the
combined lattice BRST--diffeomorphism complex, establish its uniform
gap, or calculate its harmonic anomaly coefficients on the
SM--Einstein contour.  Perturbative anomaly cancellation cannot be
substituted for these regulator-level data.
\end{firewall}

\begin{decision}[V8 acquisition]
\label{dec:v12v7-next}
V8 will construct adjacent cochain comparison maps and prove
stability of cohomology representatives under refinement.  It will
separate exact cochain-map failure, harmonic transport, and homotopy
defects, producing a summable Ward-projector contribution to the V6
forcing.
\end{decision}

\begin{assessment}[Progress after twelfth-cycle V7]
\label{ass:v12v7-status}
The local Ward projection now has a rigorous Hodge carrier.
Consistent defects split into harmonic anomalies and removable exact
terms, inconsistent defects occupy the coexact channel, and every
piece has a spectral-gap estimate.  Symmetry-equivariant Hodge data
and projector perturbation bounds connect this decomposition to the
Wilsonian contraction of V6.
\end{assessment}

\subsection{Parent record}

This V7 note was formed by appending the preceding material to the
validated twelfth-cycle V6 source, whose record was
\[
\begin{array}{c|c}
\text{lines}&1828\\
\text{bytes}&64352\\
\text{SHA--256}&
\texttt{853D875373893E06FB298680063A51C0AF726E94A86BBCFF7CA0B8511FD22155}.
\end{array}                                                     \tag{7.26}
\]

% =============================================================================
% END APPENDED TWELFTH-CYCLE V7 MATERIAL
% =============================================================================
% =============================================================================
% BEGIN APPENDED TWELFTH-CYCLE V8 MATERIAL
% =============================================================================

\section{Refinement transport of Ward cohomology}
\label{sec:v12v8}

A Hodge decomposition at each cutoff does not identify anomaly or
physical-coupling classes across cutoffs.  This version constructs
that identification from adjacent cochain maps and homotopies, then
quantifies the defects when the comparison is only approximate.

Let \((C_i^\bullet,d_i)\) and
\((C_{i+1}^\bullet,d_{i+1})\) be finite Hilbert complexes.  Suppose
\[
 R_i:C_i^\bullet\longrightarrow C_{i+1}^\bullet,\qquad
 P_i:C_{i+1}^\bullet\longrightarrow C_i^\bullet             \tag{8.1}
\]
are degree-zero cochain maps:
\[
 d_{i+1}R_i=R_id_i,\qquad
 d_iP_i=P_id_{i+1}.                                         \tag{8.2}
\]
Assume there are degree-minus-one maps \(K_i,L_i\) satisfying
\[
 P_iR_i-I=d_iK_i+K_id_i,                                    \tag{8.3}
\]
\[
 R_iP_i-I=d_{i+1}L_i+L_id_{i+1}.                            \tag{8.4}
\]
Thus \(R_i\) and \(P_i\) are cochain-homotopy inverses.

\begin{theorem}[Adjacent cohomology isomorphism]
\label{thm:v12v8-cohomology}
The maps \(R_i\) and \(P_i\) induce mutually inverse maps on every
cohomology degree.
\end{theorem}

\begin{proof}
If \(x\) is closed, (8.3) gives
\[
                              P_iR_ix-x=d_iK_ix,
\]
so \(P_iR_i\) induces the identity cohomology map.  Equation (8.4)
gives the corresponding statement for \(R_iP_i\).  Cochain-map
identities ensure that closed and exact vectors are mapped to closed
and exact vectors.
\end{proof}

Let \(H_i^q\) and \(H_{i+1}^q\) be the harmonic projections from V7.
Define
\[
 T_i^q=H_{i+1}^qR_iH_i^q:
       {\cal H}_i^q\longrightarrow{\cal H}_{i+1}^q,          \tag{8.5}
\]
\[
 S_i^q=H_i^qP_iH_{i+1}^q:
       {\cal H}_{i+1}^q\longrightarrow{\cal H}_i^q.          \tag{8.6}
\]

\begin{theorem}[Exact harmonic transport]
\label{thm:v12v8-harmonic}
Under (8.2)--(8.4),
\[
                              S_i^qT_i^q=I_{{\cal H}_i^q},
 \qquad
                              T_i^qS_i^q=I_{{\cal H}_{i+1}^q}.\tag{8.7}
\]
Hence the harmonic dimensions agree and
\[
                  s_{\min}(T_i^q)\ge\|S_i^q\|^{-1}
                                      \ge\|P_i\|^{-1},       \tag{8.8}
\]
with the convention that the estimate is void when the harmonic
space is zero.
\end{theorem}

\begin{proof}
For harmonic \(h\), \(R_ih\) is closed.  Its Hodge decomposition is
\[
                              R_ih=H_{i+1}^qR_ih+d_{i+1}b.
\]
Applying the cochain map \(P_i\), the exact part remains exact and is
annihilated by \(H_i^q\).  Equation (8.3) says
\(P_iR_ih-h\) is exact, also annihilated by \(H_i^q\).  Therefore
\(S_i^qT_i^qh=h\).  The other identity follows symmetrically from
(8.4).  The inverse bound is
\[
 \|h\|=\|S_i^qT_i^qh\|
 \le\|S_i^q\|\,\|T_i^qh\|,
\]
and \(\|S_i^q\|\le\|P_i\|\).
\end{proof}

The theorem transports the cohomology class rather than the raw
representative.  The exact component discarded by \(H_{i+1}^q\) is a
change of counterterm convention.

\subsection{Symmetry-compatible refinement}

Let \(U_{i,g}\) and \(U_{i+1,g}\) be the physical symmetry and
re-treeing representations.  Suppose
\[
                         R_iU_{i,g}=U_{i+1,g}R_i,\qquad
                         P_iU_{i+1,g}=U_{i,g}P_i.            \tag{8.9}
\]

\begin{proposition}[Equivariant cohomology transport]
\label{prop:v12v8-equivariant}
If the differentials commute with their symmetry representations,
then
\[
                         T_i^qU_{i,g}=U_{i+1,g}T_i^q,\qquad
                         S_i^qU_{i+1,g}=U_{i,g}S_i^q.        \tag{8.10}
\]
\end{proposition}

\begin{proof}
The Hodge projections commute with the representations by V7.
Insert (8.9) into (8.5)--(8.6).
\end{proof}

Thus a harmonic anomaly coefficient is compared only after its
physical channel has been transported.  A fixed coordinate component
does not define an adjacent cohomology class.

\subsection{Cochain-map failure}

Suppose \(R_i\) is only approximately a cochain map and define
\[
                              E_i^q=d_{i+1}R_i-R_id_i.        \tag{8.11}
\]
For \(h\in{\cal H}_i^q\),
\[
                              d_{i+1}R_ih=E_i^qh.            \tag{8.12}
\]

\begin{theorem}[Consistency cost of approximate transport]
\label{thm:v12v8-map-defect}
Assume the degree-\((q+1)\) Ward Laplacian at cutoff \(i+1\) has
positive gap \(\gamma_{i+1}^{q+1}\).  The coexact component of
\(R_ih\) obeys
\[
 \left\|
 d_{i+1}^*G_{i+1}^{q+1}d_{i+1}R_ih
 \right\|
 \le(\gamma_{i+1}^{q+1})^{-1/2}
       \|E_i^q\|\,\|h\|.                                    \tag{8.13}
\]
\end{theorem}

\begin{proof}
Apply the V7 consistency-defect estimate to \(a=R_ih\), then use
(8.12).
\end{proof}

The harmonic vector \(H_{i+1}^qR_ih\) always exists, but if (8.11) is
nonzero it is not the image of a genuine induced cohomology map.
Equation (8.13) measures the consistency repair; it does not make the
map exact.

\subsection{Approximate homotopy inverses}

Keep the exact cochain-map identities (8.2), but replace (8.3)--(8.4)
by
\[
 P_iR_i-I=d_iK_i+K_id_i+F_i,                                \tag{8.14}
\]
\[
 R_iP_i-I=d_{i+1}L_i+L_id_{i+1}+\widetilde F_i.             \tag{8.15}
\]

\begin{theorem}[Harmonic inverse defect]
\label{thm:v12v8-homotopy-defect}
The harmonic transports satisfy
\[
 \|S_i^qT_i^q-I\|\le\|H_i^qF_iH_i^q\|\le\|F_i\|,            \tag{8.16}
\]
\[
 \|T_i^qS_i^q-I\|
       \le\|H_{i+1}^q\widetilde F_iH_{i+1}^q\|
       \le\|\widetilde F_i\|.                               \tag{8.17}
\]
If \(\|F_i\|=f_i<1\), then
\[
                              s_{\min}(T_i^q)
                         \ge{1-f_i\over\|S_i^q\|}
                         \ge{1-f_i\over\|P_i\|}.             \tag{8.18}
\]
\end{theorem}

\begin{proof}
Repeat the proof of Theorem~\ref{thm:v12v8-harmonic}.  Exact terms are
annihilated by the harmonic projection, leaving
\(H_i^qF_iH_i^q\) and its fine analogue.  For harmonic \(h\),
\[
 (1-f_i)\|h\|
 \le\|S_i^qT_i^qh\|
 \le\|S_i^q\|\,\|T_i^qh\|,
\]
which proves (8.18).
\end{proof}

Only the harmonic compressions of \(F_i\) affect the induced class.
A large exact homotopy error can be harmless for cohomology but still
matter for locality or the action norm.

\subsection{Adjacent anomaly recursion}

Let \(a_i\in C_i^1\) be the complete regulated anomaly carrier and
put
\[
                              h_i=H_i^1a_i.                  \tag{8.19}
\]
Define the transported harmonic defect
\[
                              \delta_i
 =h_{i+1}-T_i^1h_i.                                         \tag{8.20}
\]

\begin{proposition}[Harmonic anomaly transport]
\label{prop:v12v8-anomaly}
If the transports are exact isomorphisms and
\[
                              \sum_i\|\delta_i\|<\infty,     \tag{8.21}
\]
then the harmonic anomaly coefficients form a Cauchy sequence after
transport to any fixed reference cohomology space.  The continuum
anomaly vanishes only if the transported limit is zero.
\end{proposition}

\begin{proof}
Use the inverse maps \(S_i^1\) to identify adjacent harmonic spaces
with a fixed one.  Equation (8.20) becomes a telescoping increment;
absolute summability makes the transported sequence Cauchy.  Its
limit represents the continuum cohomology class, which vanishes
exactly when that vector is zero.
\end{proof}

Summability alone does not force zero.  Starting from \(h_{i_0}=0\)
and adding nonzero summable \(\delta_i\) can leave a nonzero limiting
anomaly.  Exact species cancellation at every cutoff, or an
independent proof that the transported total is zero, is required.

\subsection{Combined refinement forcing}

For a harmonic vector \(h\), an adjacent Ward forcing budget may be
split as
\[
 \begin{split}
 {\cal E}_i^{\rm Ward}
 :=&\ \|H_{i+1}(a_{i+1}-R_ia_i)\|\\
 &+(\gamma_{i+1}^{2})^{-1/2}\|E_i^1\|\,\|h_i\|\\
 &+\|H_iF_iH_i\|\,\|h_i\|
   +\|\widetilde H_{i+1}-H_{i+1}\|\,\|a_{i+1}\|.            \tag{8.22}
 \end{split}
\]
The four rows are, respectively, physical anomaly generation,
cochain-map failure, homotopy failure, and Hodge-projector drift.
The last term can be estimated by the V7 Riesz bound when the
Laplacian nullity is constant.

If
\[
                              {\cal E}_i^{\rm Ward}
                              \le C2^{-\alpha i},             \tag{8.23}
\]
it can enter the forcing \(e_i\) of V6.  A nonzero harmonic anomaly
should not be hidden in this irrelevant forcing: its limiting class
must be shown to vanish as in
Proposition~\ref{prop:v12v8-anomaly}.

\subsection{Refinement of physical couplings}

At ghost number zero, the same maps transport the harmonic
relevant/marginal coupling space.  If
\[
                              T_i^0:{\cal H}_i^0
                              \longrightarrow{\cal H}_{i+1}^0
\]
is an isomorphism with the lower bound (8.18), running coupling
coordinates can be compared without an ill-conditioned change of
basis.  A loss of harmonic dimension signals that the regulator has
changed the physical local cohomology and cannot be treated as a
small irrelevant correction.

\begin{warning}[Cohomology dimension is a discrete gate]
\label{warn:v12v8-dimension}
Operator-norm convergence of Ward Laplacians does not by itself keep
their kernels fixed.  The nullity, cochain-homotopy equivalence, and
positive gap must be verified.  A newly appearing zero mode changes
the number of anomaly or physical-coupling classes discontinuously.
\end{warning}

\begin{ledger}[Adjacent Ward data]
\label{led:v12v8-data}
The coupled cutoff construction must provide:
\begin{enumerate}
\item cochain maps \(R_i,P_i\) transporting physical channels;
\item local homotopies \(K_i,L_i\) and their support bounds;
\item constant cohomology dimensions;
\item uniform Ward-Laplacian gaps;
\item summable cochain, homotopy, and projector defects;
\item an exact or independently vanishing transported harmonic
      anomaly.
\end{enumerate}
\end{ledger}

\begin{firewall}[No SM--Einstein cochain equivalence acquired]
\label{fire:v12v8-boundary}
No current regulator construction supplies the maps and homotopies
(8.1)--(8.4) for the combined gauge, ghost, fermion, and
diffeomorphism complex on the complex metric contour.  The continuum
harmonic anomaly has therefore not been shown to vanish by V8.
\end{firewall}

\begin{decision}[V9 acquisition]
\label{dec:v12v8-next}
Before attempting a new anomaly calculation, V9 will inspect the
frozen archive for any existing Standard Model representation-level
anomaly cancellation proof.  It will avoid duplication and then
address the first missing row: lifting algebraic species
cancellation to a cutoff-local cochain identity compatible with the
adjacent maps of V8.
\end{decision}

\begin{assessment}[Progress after twelfth-cycle V8]
\label{ass:v12v8-status}
Ward cohomology now has an exact refinement carrier.  Cochain-homotopy
equivalence yields invertible harmonic transport with a quantitative
lower bound; cochain-map, homotopy, and projector failures enter
separate estimates.  The remaining physical question is whether the
complete Standard Model anomaly cancels as a cutoff-local transported
class.
\end{assessment}

\subsection{Parent record}

This V8 note was formed by appending the preceding material to the
validated twelfth-cycle V7 source, whose record was
\[
\begin{array}{c|c}
\text{lines}&2208\\
\text{bytes}&76963\\
\text{SHA--256}&
\texttt{605A0BF38D479084673C043567A35BE4DEDBDBB5837C309B47D2F029D5D04E2B}.
\end{array}                                                     \tag{8.24}
\]

% =============================================================================
% END APPENDED TWELFTH-CYCLE V8 MATERIAL
% =============================================================================
% =============================================================================
% BEGIN APPENDED TWELFTH-CYCLE V9 MATERIAL
% =============================================================================

\section{Lifting Standard Model anomaly arithmetic to the cutoff complex}
\label{sec:v12v9}

Before writing this version, the frozen
\texttt{V100\_f11} archive was searched for an existing
representation-level anomaly proof.  Its V24 block already computes,
for one Standard Model generation, the cancellation of
\[
 SU(3)^3,\quad SU(3)^2Y,\quad SU(2)^2Y,\quad
 Y^3,\quad {\rm grav}^2Y,                                   \tag{9.1}
\]
proves vanishing of the degree-six invariant anomaly polynomial, and
counts four left-handed \(SU(2)\) doublets modulo two.  Those
calculations are not repeated here.

The missing step is analytic and regulator-local: every species
carrier must descend through one common map into the finite Ward
complex of V7--V8 before the representation sum is taken.

\subsection{Universal invariant-polynomial carrier}

Let \({\cal P}\) be the finite-dimensional space of degree-six
invariant anomaly polynomials relevant in four dimensions.  For each
left-handed species \(s\), let \(A_s\in{\cal P}\) be its
representation coefficient, including spectator multiplicities.
The frozen V24 result is
\[
                              \sum_sA_s=0.                   \tag{9.2}
\]

At cutoff \(i\), let \(a_{i,s}\in C_i^1\) be the regulated local Ward
anomaly carrier.  Suppose there is one linear map
\[
                              J_i:{\cal P}\longrightarrow C_i^1
                                                                    \tag{9.3}
\]
and a decomposition
\[
                              a_{i,s}
 =J_iA_s+d_{i,0}b_{i,s}+r_{i,s}.                            \tag{9.4}
\]
Here \(b_{i,s}\in C_i^0\) is a local counterterm and \(r_{i,s}\) is a
regulator remainder.

\begin{theorem}[Cutoff-local species cancellation]
\label{thm:v12v9-common}
Under (9.2)--(9.4), the complete species sum satisfies
\[
                              a_i^{\rm tot}:=\sum_sa_{i,s}
 =d_{i,0}b_i^{\rm tot}+r_i^{\rm tot},                       \tag{9.5}
\]
where
\[
                              b_i^{\rm tot}=\sum_sb_{i,s},
 \qquad
                              r_i^{\rm tot}=\sum_sr_{i,s}.   \tag{9.6}
\]
Consequently its harmonic anomaly is
\[
                              H_i^1a_i^{\rm tot}
                              =H_i^1r_i^{\rm tot}.           \tag{9.7}
\]
\end{theorem}

\begin{proof}
Sum (9.4).  Linearity of the common map and (9.2) give
\(J_i\sum_sA_s=0\), proving (9.5).  Harmonic projection annihilates
the exact term \(d_{i,0}b_i^{\rm tot}\), giving (9.7).
\end{proof}

Equation (9.7) is the precise regulator bridge.  Representation
arithmetic removes the universal cohomology source, but the full
finite-cutoff anomaly vanishes only after the summed remainder has
zero harmonic projection.

\subsection{Why the map must be common}

A species-dependent regulator can instead produce
\[
 a_{i,s}=J_{i,s}A_s+d_{i,0}b_{i,s}+r_{i,s}.                 \tag{9.8}
\]
Choose a reference \(J_i\) and write
\[
                              J_{i,s}=J_i+\delta J_{i,s}.    \tag{9.9}
\]

\begin{proposition}[Pre-sum transport defect]
\label{prop:v12v9-presum}
Under (9.2), the harmonic anomaly obeys
\[
 \|H_i^1a_i^{\rm tot}\|
 \le\sum_s\|H_i^1\delta J_{i,s}A_s\|
       +\|H_i^1r_i^{\rm tot}\|.                             \tag{9.10}
\]
In particular, cancellation of \(\sum_sA_s\) does not cancel
species-dependent transport errors.
\end{proposition}

\begin{proof}
Insert (9.9) into the sum of (9.8).  The common term vanishes by
(9.2), exact terms vanish under \(H_i^1\), and the triangle inequality
gives (9.10).
\end{proof}

Different hypercharges require different representation matrices but
not different identifications of the spacetime, gauge, ghost, and
metric variation.  Projecting species before transporting them to a
common Ward complex can therefore create a false residual or hide a
real one.

\subsection{A quantitative common-map criterion}

Let \(\widehat J_i:{\cal P}\to C_i^1\) be an explicitly constructed
local descent map.  Suppose
\[
 a_{i,s}-\widehat J_iA_s
                              =d_{i,0}b_{i,s}+r_{i,s},       \tag{9.11}
\]
and the degree-one Ward Laplacian has gap \(\gamma_i^1\).

\begin{theorem}[Hodge test for the descent remainder]
\label{thm:v12v9-hodge}
If each \(a_{i,s}-\widehat J_iA_s\) is \(d_{i,1}\)-closed, then
\[
 r_{i,s}\ \text{may be chosen harmonic},\qquad
 r_{i,s}=H_i^1(a_{i,s}-\widehat J_iA_s),                    \tag{9.12}
\]
with the minimal counterterm
\[
 b_{i,s}=d_{i,0}^*G_i^1
            (a_{i,s}-\widehat J_iA_s)                       \tag{9.13}
\]
satisfying
\[
 \|b_{i,s}\|
 \le(\gamma_i^1)^{-1/2}
       \|a_{i,s}-\widehat J_iA_s\|.                         \tag{9.14}
\]
\end{theorem}

\begin{proof}
Apply the exact-anomaly Hodge theorem of V7 to
\(a_{i,s}-\widehat J_iA_s\).  Its harmonic part is (9.12), and its
exact part has the primitive (9.13) with bound (9.14).
\end{proof}

If the difference in (9.11) is not closed, its coexact consistency
piece must be retained separately using V7; it cannot be assigned to
the harmonic remainder.

\subsection{Refinement of the universal descent}

Let \(R_i:C_i^1\to C_{i+1}^1\) be the cochain comparison of V8.
Define
\[
                              \delta_i^J
                    =R_iJ_i-J_{i+1}.                        \tag{9.15}
\]
For the complete representation,
\[
                              \delta_i^J\sum_sA_s=0          \tag{9.16}
\]
identically.  Thus common descent drift cancels before any norm
estimate.

\begin{proposition}[Remainder transport controls the total class]
\label{prop:v12v9-refine}
Assume exact cochain transport and the common decomposition (9.4).
Then
\[
 \begin{split}
 H_{i+1}^1a_{i+1}^{\rm tot}
 -T_i^1H_i^1a_i^{\rm tot}
 =H_{i+1}^1
   \left(r_{i+1}^{\rm tot}-R_ir_i^{\rm tot}\right).          \tag{9.17}
 \end{split}
\]
\end{proposition}

\begin{proof}
Use (9.7) at both cutoffs.  The harmonic transport is
\(T_i^1=H_{i+1}^1R_iH_i^1\).  For any closed vector, subtracting its
harmonic projection leaves an exact vector, which remains exact under
the cochain map and is annihilated by \(H_{i+1}^1\).  Hence
\[
 T_i^1H_i^1r_i^{\rm tot}=H_{i+1}^1R_ir_i^{\rm tot},
\]
which gives (9.17).
\end{proof}

If the right side of (9.17) is summable, the total harmonic anomaly
converges after transport.  To prove anomaly freedom one still needs
\[
                              \|H_i^1r_i^{\rm tot}\|
                              \longrightarrow0,              \tag{9.18}
\]
or an exact zero at every cutoff.  Adjacent summability without
(9.18) permits a nonzero limiting class.

\subsection{Locality requirement}

The map \(J_i\) must be local in the regulator geometry.  A useful
quantitative form is a decomposition
\[
                              J_iA=\sum_{x\in\Lambda_i}
                                      j_{i,x}(A),            \tag{9.19}
\]
where \(j_{i,x}\) has uniformly bounded support radius and norm after
canonical rescaling.  The primitives \(b_{i,s}\) in (9.4) must obey
the same quasi-local algebra or a summable exponential tail.

A purely algebraic Hodge inverse on the full finite lattice may be
nonlocal even when its norm satisfies (9.14).  Therefore the spectral
gap is not enough; one also needs quasi-locality of the Green operator
or a local descent construction that produces \(b_{i,s}\) directly.

\begin{warning}[Global anomalies remain separate]
\label{warn:v12v9-global}
The common local polynomial cancellation (9.2) does not evaluate
determinant-line holonomy around large gauge transformations or
diffeomorphism loops.  The frozen V24 doublet count removes the usual
fundamental \(SU(2)\) parity obstruction only after the regulator has
been shown to realise the same mod-two spectral-flow class.
\end{warning}

\subsection{Mod-two lift criterion}

Let \(\Gamma_i\) be a cutoff loop representing the intended nontrivial
\(SU(2)\) class, and let
\(\operatorname{sf}_2(D_{i,s};\Gamma_i)\in\mathbb Z_2\) be the
mod-two spectral flow of species \(s\).

\begin{proposition}[Regulator parity cancellation]
\label{prop:v12v9-parity}
Suppose every fundamental left-handed doublet has the same unit
mod-two spectral flow on \(\Gamma_i\), singlets have zero flow, and
spectator colour creates three identical copies of the quark
doublet.  Then one generation has total flow
\[
                              3+1=0\pmod2.                   \tag{9.20}
\]
This parity is invariant under any refinement homotopy that keeps the
endpoint operators invertible.
\end{proposition}

\begin{proof}
Mod-two spectral flow is additive under direct sums.  The four
doublet copies therefore give zero.  Homotopy invariance follows
because eigenvalue crossings can be created or annihilated only in
pairs while the endpoint gap remains open.
\end{proof}

This proposition records the analytic lift required to use the
already-proved count.  Construction of the loop and equality of the
species spectral-flow units remain regulator tasks.

\begin{ledger}[Nonduplicated anomaly acquisition]
\label{led:v12v9-data}
The representation sums are complete in frozen V24.  The remaining
work is:
\begin{enumerate}
\item construct one common local descent map \(J_i\);
\item prove exact consistency or bound its coexact defect;
\item show the summed harmonic remainder vanishes as in (9.18);
\item prove quasi-locality of the counterterm primitive;
\item transport the data through the V8 cochain homotopies;
\item realise and stabilise the global mod-two loop.
\end{enumerate}
\end{ledger}

\begin{firewall}[Arithmetic not repeated, bridge not assumed]
\label{fire:v12v9-boundary}
V9 relies on the archived representation identities but does not
claim that (9.4), locality (9.19), vanishing (9.18), or the mod-two
spectral-flow hypotheses have been proved for the coupled renewal
regulator.  Those are the remaining analytic anomaly inputs.
\end{firewall}

\begin{decision}[Mandatory V10 audit]
\label{dec:v12v9-next}
V10 will perform the compulsory companion audit.  The following
substantive version will construct a quasi-local Hodge inverse from a
gapped finite-range Ward Laplacian, providing the locality estimate
missing after (9.14).
\end{decision}

\begin{assessment}[Progress after twelfth-cycle V9]
\label{ass:v12v9-status}
The prior Standard Model anomaly arithmetic has been reused without
duplication.  Its exact cancellation now feeds a common finite-cutoff
descent theorem: only species-dependent transport and the summed
harmonic regulator remainder survive.  Refinement and mod-two
criteria state precisely what must be proved before the arithmetic
becomes a continuum Ward identity.
\end{assessment}

\subsection{Parent record}

This V9 note was formed by appending the preceding material to the
validated twelfth-cycle V8 source, whose record was
\[
\begin{array}{c|c}
\text{lines}&2560\\
\text{bytes}&88785\\
\text{SHA--256}&
\texttt{48441F455FFF09E67B18BC321EBE2CC483F638A4A55F7C9B4C356F70F5914167}.
\end{array}                                                     \tag{9.21}
\]

% =============================================================================
% END APPENDED TWELFTH-CYCLE V9 MATERIAL
% =============================================================================
% =============================================================================
% BEGIN APPENDED TWELFTH-CYCLE V10 MATERIAL
% =============================================================================

\section{Twelfth-cycle V10: compulsory companion audit and nonlinear
transport universality}
\label{sec:v12v10}

\subsection{Purpose of the audit}

This is the compulsory five-version comparison requested for the
parallel Standard Model--Einstein, Yang--Mills, and Navier--Stokes
programmes.  The newest Yang--Mills note has advanced from V42 to
V44.  The newest Navier--Stokes note remains V26.  The audit changes
the immediate Standard Model task in one precise way: the
species-common descent map introduced in V9 must be required to be
common through the full nonlinear background-dependent coordinate
transport, and not merely through its linear or pure-momentum jet.

The files read for this audit are
\[
\begin{array}{c|r|r|c}
\text{record}&\text{lines}&\text{bytes}&\text{SHA--256}\\ \hline
\text{YM V44}&17718&606281&
\texttt{A9C778335C04A2B23239066574643AE57383B34B802F456F2CA7AA3807217C6C}\\
\text{NS V26}&5319&278283&
\texttt{82C887F8C2C69E4F28FAD7C3DC8DE5B9099CB5A4BF431EBA1FFF3E91935978AD}.
\end{array}                                                   \tag{10.1}
\]

\subsection{What the new Yang--Mills calculation proves}

The YM V43 calculation uses the corrected spectator transposition
\(\tau_{24}\).  At the first corrected corner pair its exact
tadpole--bubble residual vanishes at external momentum orders zero
and one but has the nonzero second coefficient
\[
 -\frac{489372186744556308849}
        {89418228107593602350}.                              \tag{10.2}
\]
YM V44 then constructs the exact re-tree jet through order two.  It
proves the involution identity, the triviality of the first two
log-determinant jets, and covariance of all eight Wilson and endpoint
router columns through order two.  Thus the residual (10.2) cannot be
assigned to the retained-slot order, the pure-momentum linear
re-tree map, its determinant, or the linear background router.  The
unresolved alternatives are the mixed background dependence of the
nonlinear group-valued coordinate map and a defect in a second-jet
vertex compiler.

This sharpens the lesson already imported in V5.  Exact covariance
of a linearized coordinate map is not a proof of covariance of its
mixed background jets.  Conversely, failure of a diagrammatic
partial sum must not be named a measure-density anomaly before the
full coordinate change and an independent vertex regression have
been computed.

\subsection{Cancellation under a common nonlinear descent map}

We now make the V9 species-universality statement independent of
perturbative order.  Let \(\mathcal P\) be the finite-dimensional
space containing the local anomaly polynomials and let
\(\mathcal C_i^1\) be the regulator-\(i\) anomaly cochains.  The
background and external-momentum variables are denoted by \(z\) and
may contain several components.  For every fermion species \(s\),
let \(A_s\in\mathcal P\) be its signed anomaly coefficient.  The
archived representation calculation gives
\[
                         \sum_s A_s=0.                       \tag{10.3}
\]
Suppose the regulator realizes descent by analytic linear maps in
the polynomial argument,
\[
 {\cal J}_{i,s}(z):\mathcal P\longrightarrow\mathcal C_i^1.
                                                               \tag{10.4}
\]
No linearity in \(z\) is assumed.

\begin{theorem}[All-order cancellation of a common nonlinear
transport]
\label{thm:v12v10-common}
If
\[
 {\cal J}_{i,s}(z)={\cal J}_i(z)
 \quad\hbox{for every \(s\) and every \(z\) near zero},       \tag{10.5}
\]
then the polynomial part of the total regulated anomaly vanishes
identically:
\[
             \sum_s{\cal J}_{i,s}(z)A_s=0.                  \tag{10.6}
\]
Consequently every mixed Taylor coefficient in background amplitude,
external momentum, and any other coordinate parameter also
vanishes.
\end{theorem}

\begin{proof}
For fixed \(z\), linearity of \({\cal J}_i(z)\) in its argument and
(10.3) give
\[
 \sum_s{\cal J}_{i,s}(z)A_s
 ={\cal J}_i(z)\sum_s A_s=0.
\]
The left side is identically zero in a neighbourhood of the origin,
so every one of its partial derivatives at the origin is zero.
\end{proof}

The theorem does not require the common nonlinear coordinate
Jacobian to vanish.  It says that a common descent contribution,
including every mixed jet of that contribution, is applied to the
already vanishing total anomaly polynomial.  The analytic work is
therefore to prove commonality and locality of the regulator map.

\subsection{A quantitative nonuniversality budget}

Write
\[
 {\cal J}_{i,s}(z)={\cal J}_i(z)+\delta{\cal J}_{i,s}(z).
                                                               \tag{10.7}
\]
For a multi-index \(\alpha\), put
\[
 E_{i,\alpha}:=
 \sum_s
 \left\|
 \partial^\alpha\delta{\cal J}_{i,s}(0)
 \right\|_{\mathcal P\to\mathcal C_i^1}\,\|A_s\|_{\mathcal P}.
                                                               \tag{10.8}
\]

\begin{proposition}[Mixed-jet defect bound]
\label{prop:v12v10-jet-bound}
Under (10.3) and (10.7),
\[
 \left\|
 \partial^\alpha\left[
   \sum_s{\cal J}_{i,s}(z)A_s
 \right]_{z=0}
 \right\|_{\mathcal C_i^1}
 \le E_{i,\alpha}.                                          \tag{10.9}
\]
If \(H_i\) is the orthogonal harmonic projector of V7, the same
bound holds after applying \(H_i\).
\end{proposition}

\begin{proof}
The common term cancels by differentiating the identity in
Theorem~\ref{thm:v12v10-common}.  The triangle inequality and the
definition of the operator norm give (10.9).  Since an orthogonal
projector has norm at most one, applying \(H_i\) cannot increase the
right side.
\end{proof}

Thus a putative mixed background/momentum remainder must be entered
in the refinement ledger as a species-nonuniversality defect, a
descent remainder, or a consistency defect.  It cannot be inferred
from the failure of a partial diagrammatic row.  In particular, the
YM V44 dichotomy prescribes an upstream validation order for the SM
regulator: compute the nonlinear coordinate Jacobian and independently
regress the relevant vertex jet before estimating
\(\delta{\cal J}_{i,s}\).

\subsection{The unchanged Navier--Stokes record}

NS V26 remains the current companion file.  It computes exact
pointwise rank-one norms for the first two derivatives of the Oseen
kernel.  The first derivative gains nothing over the full symmetric
traceless norm.  The second derivative has a smaller rank-one norm
on an intermediate radial interval, with the squared norm reduced to
the maximum of a quadratic polynomial in the squared angle cosine.

The reusable methodological point is modest but exact.  A restricted
source class can improve a kernel constant only after the restriction
is propagated through the tensor contraction; a bound on the
ambient operator space can conceal a real but localized gain.  For
the SM problem, this supports decomposing local Ward sources into
exact, coexact, and harmonic sectors before estimating the Green
operator.  It does not by itself provide a Ward gap, exponential
locality, or any continuum construction.

\subsection{Adjusted acquisition order}

The audit produces the following order.

\begin{enumerate}
\item Prove a quasi-locality theorem for the inverse Ward Laplacian
      under an explicit uniform gap and finite-range hypothesis.
      This converts the algebraic primitive \(d_i^*G_i a_i\) from
      V7 into a controlled quasi-local counterterm.
\item Apply that theorem only after removing harmonic modes; the
      inverse on the full complex is undefined when cohomology is
      present.
\item Formulate the regulator universality target for the complete
      nonlinear coordinate map.  Mixed jets are included in
      (10.5), and failures are charged to the explicit budget
      (10.8).
\item Keep the anomaly endpoint separate: summability of transported
      harmonic increments gives a limit, whereas cancellation
      requires that limit to be zero.
\end{enumerate}

\begin{assessment}[Status after V10]
\label{ass:v12v10-status}
The audit preserves the V9 direction but strengthens its hypothesis.
Representation-level cancellation survives arbitrary nonlinear
background and momentum dependence whenever the descent map is
species-common.  The new YM calculation shows why commonality must
be checked at the full mixed-jet level.  The next missing bridge is
spatial locality of the Hodge primitive under a regulator-uniform
spectral gap.
\end{assessment}

\begin{firewall}[Claim boundary after V10]
\label{fire:v12v10-claim}
No nonlinear Standard Model regulator Jacobian has yet been
constructed, and no uniform Ward gap, quasi-local counterterm,
reflection positivity, Einstein constraint measure, or full
SM--Einstein continuum has been proved.  Equations (10.6) and (10.9)
are conditional algebraic consequences of the stated regulator
maps.
\end{firewall}

\subsection{Parent record}

This V10 note appends to the validated V9 file with record
\[
\begin{array}{c|c}
\text{lines}&2876\\
\text{bytes}&100149\\
\text{SHA--256}&
\texttt{92C8065FBECFA30CABE9C10B989ADB05B7426CAE928DEB74ED4C7BE1F0FF2471}.
\end{array}
\]

% =============================================================================
% END APPENDED TWELFTH-CYCLE V10 MATERIAL
% =============================================================================
% =============================================================================
% BEGIN APPENDED TWELFTH-CYCLE V11 MATERIAL
% =============================================================================

\section{Twelfth-cycle V11: exponential quasi-locality of the gapped
Ward inverse}
\label{sec:v12v11}

\subsection{The locality problem left by the Hodge construction}

V7 produced the canonical primitive
\[
                         b=d^*G a,                          \tag{11.1}
\]
where \(G\) is the Moore--Penrose inverse of the finite-regulator Ward
Laplacian.  That algebraic formula and the norm bound
\(\|G\|\le\gamma^{-1}\) do not by themselves prove that \(b\) is a
local counterterm.  V10 therefore made quasi-locality the next
acquisition.

This section gives a volume-uniform exponential kernel bound from
three explicit hypotheses: finite propagation, a uniform upper norm
bound, and a uniform gap above the harmonic subspace.  A polynomial
formula for the pseudoinverse avoids assuming that the harmonic
projector itself is local.

\subsection{Local block operators}

Let \((\Lambda,\operatorname{dist})\) be a finite metric space and
\[
                     {\cal K}=\bigoplus_{x\in\Lambda}E_x     \tag{11.2}
\]
an orthogonal sum of finite-dimensional fibres.  For
\(A\in{\cal L}({\cal K})\), write \(A_{xy}:E_y\to E_x\) for its block
kernel.  We say that \(A\) has propagation at most \(R\) when
\[
                    A_{xy}=0\quad\hbox{if }\operatorname{dist}(x,y)>R.
                                                               \tag{11.3}
\]
Products of operators with propagations \(R_1,R_2\) have propagation
at most \(R_1+R_2\), by the triangle inequality.

Let \(\Delta=\Delta^*\ge0\) satisfy
\[
 \operatorname{spec}(\Delta)\subset\{0\}\cup[\gamma,M],
 \qquad 0<\gamma\le M<\infty,                               \tag{11.4}
\]
and suppose that \(\Delta\) has propagation at most \(R>0\).  Denote
its Moore--Penrose inverse by \(G=\Delta^\dagger\).  Thus \(G\) is
zero on \(\ker\Delta\) and equals \(\Delta^{-1}\) on
\((\ker\Delta)^\perp\).

Put
\[
 \alpha=\frac{2}{M+\gamma},
 \qquad
 \rho=\frac{M-\gamma}{M+\gamma}.                            \tag{11.5}
\]
When \(M>\gamma\), one has \(0<\rho<1\).

\begin{theorem}[Local polynomial series for the Ward inverse]
\label{thm:v12v11-series}
Under (11.3)--(11.5),
\[
 \boxed{
 G=\alpha^2\sum_{n=0}^{\infty}(n+1)\,
       \Delta\bigl(I-\alpha\Delta\bigr)^n}
                                                               \tag{11.6}
\]
with convergence in operator norm.  Every summand indexed by \(n\)
has propagation at most \((n+1)R\).
\end{theorem}

\begin{proof}
Set \(S=I-\alpha\Delta\).  On \(\ker\Delta\), every summand in
(11.6) is zero.  If \(\lambda\in[\gamma,M]\), then
\[
 |1-\alpha\lambda|
 \le\frac{M-\gamma}{M+\gamma}=\rho.                         \tag{11.7}
\]
Consequently the series is uniformly convergent on the positive
spectrum.  The scalar value of its sum there is
\[
 \alpha^2\lambda
 \sum_{n=0}^{\infty}(n+1)(1-\alpha\lambda)^n
 =
 \frac{\alpha^2\lambda}
      {\bigl(1-(1-\alpha\lambda)\bigr)^2}
 =\frac1\lambda.                                            \tag{11.8}
\]
The spectral theorem now gives (11.6).

The operator \(S\) has propagation at most \(R\).  Hence \(S^n\)
has propagation at most \(nR\), and multiplication by \(\Delta\)
raises this to at most \((n+1)R\).
\end{proof}

The factor \(\Delta\) in every summand is essential.  It annihilates
the harmonic subspace before the Neumann-type series is summed.
There is therefore no nonlocal harmonic projector in the
construction.

\subsection{The exponential kernel estimate}

For \(r\ge0\), define
\[
 N_R(r):=\max\left\{0,\left\lceil\frac rR\right\rceil-1\right\}.
                                                               \tag{11.9}
\]

\begin{theorem}[Uniform exponential decay of the pseudoinverse]
\label{thm:v12v11-decay}
Assume \(M>\gamma\).  For all \(x,y\in\Lambda\),
\[
 \boxed{
 \|G_{xy}\|
 \le
 \alpha^2M\,
 \rho^{N_R(\operatorname{dist}(x,y))}
 \frac{N_R(\operatorname{dist}(x,y))+1
       -N_R(\operatorname{dist}(x,y))\rho}
      {(1-\rho)^2}.}                                        \tag{11.10}
\]
In particular, for every
\[
             0<\mu<\frac1R\log\frac1\rho,                   \tag{11.11}
\]
there is a constant \(C_{\mu}=C_{\mu}(\gamma,M,R)\), independent of
\(\Lambda\) and of the fibre dimensions, such that
\[
                         \|G_{xy}\|
 \le C_{\mu}e^{-\mu\operatorname{dist}(x,y)}.                \tag{11.12}
\]
If \(M=\gamma\), the series has only its zeroth nonzero spectral
term and \(G=\gamma^{-2}\Delta\), so \(G\) has propagation at most
\(R\).
\end{theorem}

\begin{proof}
Let \(r=\operatorname{dist}(x,y)\) and \(N=N_R(r)\).  The propagation
statement in Theorem~\ref{thm:v12v11-series} shows that the
\((x,y)\) block of the \(n\)-th summand vanishes whenever \(n<N\).
Moreover, by (11.7),
\[
 \left\|\Delta S^n\right\|
 \le M\rho^n.                                                \tag{11.13}
\]
Taking the \((x,y)\) block in (11.6) and using the operator norm bound
on each block therefore gives
\[
 \|G_{xy}\|
 \le\alpha^2M\sum_{n=N}^{\infty}(n+1)\rho^n.                \tag{11.14}
\]
Direct summation yields
\[
 \sum_{n=N}^{\infty}(n+1)\rho^n
 =
 \rho^N\frac{N+1-N\rho}{(1-\rho)^2},                        \tag{11.15}
\]
which proves (11.10).

For (11.12), choose
\(\theta=e^{-\mu R}\), so \(\rho<\theta<1\).  The elementary bounds
\[
 (N+1)\rho^N
 \le\frac{\theta^N}{(1-\rho/\theta)^2},
 \qquad
 \rho^{N+1}\le\rho\theta^N                                  \tag{11.16}
\]
and \(N\ge r/R-1\) imply
\(\theta^N\le e^{\mu R}e^{-\mu r}\).  Substitution in
(11.10) gives (11.12), with a constant depending only on the
displayed parameters.  If \(M=\gamma\), the positive spectrum is
the singleton \(\{\gamma\}\), so
\(\Delta=\gamma(I-H)\) and
\(G=\gamma^{-1}(I-H)=\gamma^{-2}\Delta\).
\end{proof}

\begin{remark}[No locality assumption on the harmonic projector]
\label{rem:v12v11-harmonic}
The projector \(H\) onto \(\ker\Delta\) may be global.  It never
appears as an input to the proof.  The gap in (11.4) permits the
function that equals zero at \(0\) and \(1/\lambda\) on
\([\gamma,M]\) to be represented by the uniformly convergent local
polynomial series (11.6).  This is the precise reason harmonic modes
do not obstruct the kernel estimate when the positive gap is
uniform.
\end{remark}

\subsection{Application to a finite Ward complex}

Consider a finite Hilbert cochain complex
\[
 {\cal C}^{q-1}\xrightarrow{d_{q-1}}
 {\cal C}^{q}\xrightarrow{d_q}{\cal C}^{q+1},
 \qquad d_qd_{q-1}=0,                                      \tag{11.17}
\]
with Ward Laplacian
\[
 \Delta_q=d_{q-1}d_{q-1}^*+d_q^*d_q.                       \tag{11.18}
\]
Suppose its cochain spaces are decomposed over the same finite metric
set.  Assume \(\Delta_q\) has propagation \(R\), norm at most \(M\),
and the gap (11.4).  Let \(H_q\) and \(G_q\) be its harmonic
projector and pseudoinverse.

\begin{theorem}[Quasi-local exact Ward counterterm]
\label{thm:v12v11-counterterm}
Let \(a\in{\cal C}^q\) satisfy
\[
                         d_q a=0,\qquad H_q a=0.             \tag{11.19}
\]
Then
\[
                         b=d_{q-1}^*G_q a                   \tag{11.20}
\]
obeys
\[
                         d_{q-1}b=a,\qquad
 \|b\|\le\frac{\|d_{q-1}\|}{\gamma}\|a\|.                  \tag{11.21}
\]
Suppose in addition that \(d_{q-1}^*\) has propagation at most
\(R_d\) and
\[
 K_d:=\sup_x\sum_z
       \|(d_{q-1}^*)_{xz}\|<\infty.                         \tag{11.22}
\]
If \(a\) is supported in \(S\subset\Lambda\), then for every \(\mu\)
satisfying (11.11),
\[
 \boxed{
 \|b_x\|
 \le K_d C_\mu e^{\mu R_d}|S|^{1/2}
      e^{-\mu\operatorname{dist}(x,S)}\|a\|.}               \tag{11.23}
\]
Thus the canonical anomaly-removing primitive is exponentially
quasi-local.
\end{theorem}

\begin{proof}
The Hodge identities imply
\[
 \Delta_qG_q=I-H_q.                                        \tag{11.24}
\]
They also give \(d_qG_q=G_{q+1}d_q\): this follows first for every
positive spectral subspace from
\(d_q\Delta_q=\Delta_{q+1}d_q\), and both sides vanish on the
harmonic subspace.  Hence (11.19) implies \(d_qG_q a=0\).
Using (11.18), (11.19), and (11.24),
\[
 a=\Delta_qG_q a
  =d_{q-1}d_{q-1}^*G_q a+d_q^*d_qG_q a
  =d_{q-1}b.                                                \tag{11.25}
\]
The spectral gap gives \(\|G_q\|\le\gamma^{-1}\), proving the norm
bound.

By Theorem~\ref{thm:v12v11-decay},
\[
 \|(G_qa)_z\|
 \le\sum_{y\in S}\|(G_q)_{zy}\|\,\|a_y\|
 \le C_\mu e^{-\mu\operatorname{dist}(z,S)}
            |S|^{1/2}\|a\|.                                \tag{11.26}
\]
Only \(z\) with \(\operatorname{dist}(x,z)\le R_d\) contribute to
\(b_x\).  For such \(z\),
\[
 \operatorname{dist}(z,S)
 \ge\operatorname{dist}(x,S)-R_d.                           \tag{11.27}
\]
Equations (11.22), (11.26), and (11.27) give (11.23).
\end{proof}

\begin{corollary}[A summable quasi-local shell ledger]
\label{cor:v12v11-shell}
Let the regulator index be \(i\), and suppose the constants
\(\gamma,M,R,R_d,K_d\) in
Theorem~\ref{thm:v12v11-counterterm} are uniform in \(i\).  If closed,
harmonic-free local anomaly remainders \(a_i\), supported in sets of
uniformly bounded cardinality, satisfy
\[
                         \sum_i\|a_i\|<\infty,               \tag{11.28}
\]
then their canonical primitives \(b_i=d_i^*G_i a_i\) converge
absolutely in norm.  If the supports lie in a common set \(S\), their
sum has a uniform exponential tail away from \(S\).
\end{corollary}

\begin{proof}
The norm estimate in (11.21) gives absolute convergence.  Summing
(11.23) gives the pointwise exponential tail.
\end{proof}

\subsection{Where the real gap problem now lies}

Theorem~\ref{thm:v12v11-decay} is a conditional locality theorem, but
its hypotheses are directly testable on every regulator.  Locality
of the Ward differential normally gives finite propagation and an
upper bound after the fields have been scaled to a fixed block.
The hard hypothesis is a positive lower bound \(\gamma\) independent
of volume and refinement.

A global gauge or diffeomorphism complex can have long-wavelength
positive modes whose eigenvalues tend to zero.  Declaring the exact
zero modes harmonic does not remove these small positive modes.
For such a sequence, \(\rho\) approaches one and the decay rate in
(11.11) collapses.  Therefore the next step cannot simply assume a
global gap.  It must obtain a blockwise relative Poincare estimate,
remove the finite-dimensional matching modes on each block, and
prove that the resulting local primitives patch with summable
overlap errors.

\begin{assessment}[Status after V11]
\label{ass:v12v11-status}
The algebraic Hodge counterterm of V7 is now proved quasi-local under
a regulator-uniform finite-range spectral gap.  The proof gives an
explicit local polynomial expansion of the Moore--Penrose inverse
and does not assume locality of the harmonic projector.  The
remaining bottleneck has moved upstream to a uniform blockwise Ward
Poincare inequality and a controlled patching theorem.
\end{assessment}

\begin{firewall}[Claim boundary after V11]
\label{fire:v12v11-claim}
No uniform gap for the Standard Model--Einstein Ward complex has
been established.  In particular, the theorem does not remove
long-wavelength gauge or diffeomorphism modes, prove that local block
primitives agree on overlaps, construct the interacting continuum
measure, or complete the SM--Einstein continuum.
\end{firewall}

\subsection{Parent record}

This V11 note appends to the validated V10 file with record
\[
\begin{array}{c|c}
\text{lines}&3113\\
\text{bytes}&109619\\
\text{SHA--256}&
\texttt{ABF0023C8E6E678CF7652AB582685F74EF870E6FD64083ECCA1D7697D4373A1A}.
\end{array}
\]

% =============================================================================
% END APPENDED TWELFTH-CYCLE V11 MATERIAL
% =============================================================================
% =============================================================================
% BEGIN APPENDED TWELFTH-CYCLE V12 MATERIAL
% =============================================================================

\section{Twelfth-cycle V12: local-to-global Ward contraction from a
contractive overlap defect}
\label{sec:v12v12}

\subsection{Why the argument moves upstream}

V11 proves exponential locality when the global Ward Laplacian has a
uniform positive gap above its kernel.  In a growing volume that
hypothesis can fail because of long-wavelength positive modes.  The
appropriate upstream object is then a blockwise contracting homotopy:
local Ward equations are solved on bounded blocks, the solutions are
patched, and only the overlap defect is required to contract.

This section proves the abstract patching step.  It turns any
finite-range approximate homotopy with defect norm below one into an
exact exponentially quasi-local homotopy.  The remaining
model-dependent task is reduced to constructing the block solvers
and verifying one explicit norm inequality.

\subsection{An approximate contraction relative to matching modes}

Let
\[
 \cdots\longrightarrow{\cal C}^{q-1}
 \xrightarrow{d_{q-1}}{\cal C}^{q}
 \xrightarrow{d_q}{\cal C}^{q+1}
 \longrightarrow\cdots,\qquad d_qd_{q-1}=0                 \tag{12.1}
\]
be a finite Hilbert cochain complex.  Every \({\cal C}^q\) is an
orthogonal sum of fibres over a common finite metric space
\(\Lambda\).  Propagation is defined as in (11.3).

Let \(P_q:{\cal C}^q\to{\cal C}^q\) be an orthogonal projection and
let \(K_q:{\cal C}^q\to{\cal C}^{q-1}\) have degree minus one.
Assume
\[
 P_{q+1}d_q=d_qP_q,\qquad
 K_qP_q=0,\qquad P_{q-1}K_q=0.                              \tag{12.2}
\]
Thus \(P\) is a cochain projection onto the retained matching modes,
and the provisional homotopy acts only on their complement.  Define
the overlap error by
\[
 E_q:=I-P_q-d_{q-1}K_q-K_{q+1}d_q.                         \tag{12.3}
\]

\begin{lemma}[Algebra of the overlap error]
\label{lem:v12v12-error}
Under (12.1)--(12.3),
\[
 d_qE_q=E_{q+1}d_q,\qquad
 P_qE_q=E_qP_q=0.                                          \tag{12.4}
\]
\end{lemma}

\begin{proof}
Using \(d_qd_{q-1}=0\) and the cochain identity for \(P\),
\[
\begin{split}
 d_qE_q
 &=d_q-P_{q+1}d_q-d_qK_{q+1}d_q,\\
 E_{q+1}d_q
 &=d_q-P_{q+1}d_q-d_qK_{q+1}d_q,
\end{split}
\]
which proves the first identity.  Multiplication of (12.3) by \(P_q\)
on the right gives
\[
 E_qP_q=P_q-P_q-d_{q-1}K_qP_q-K_{q+1}P_{q+1}d_q=0.
\]
Multiplication on the left and (12.2) give \(P_qE_q=0\) in the same
way.
\end{proof}

Assume now that, uniformly in degree,
\[
 \|E_q\|\le\kappa<1.                                       \tag{12.5}
\]
Then
\[
 S_q:=(I-E_q)^{-1}=\sum_{n=0}^{\infty}E_q^n                \tag{12.6}
\]
converges in operator norm.  Lemma~\ref{lem:v12v12-error} implies
\[
 d_qS_q=S_{q+1}d_q,\qquad P_qS_q=S_qP_q=P_q.               \tag{12.7}
\]

\begin{theorem}[Exact correction of a patched Ward homotopy]
\label{thm:v12v12-correction}
Under (12.1)--(12.6), define
\[
                         \widehat K_q:=K_qS_q.              \tag{12.8}
\]
Then
\[
 \boxed{
 d_{q-1}\widehat K_q+\widehat K_{q+1}d_q=I-P_q.}            \tag{12.9}
\]
Moreover,
\[
 \|\widehat K_q\|
 \le\frac{\|K_q\|}{1-\kappa}.                              \tag{12.10}
\]
Hence every closed cochain \(a\in{\cal C}^q\) with \(P_qa=0\) has
the exact primitive
\[
                         b=\widehat K_q a,\qquad d_{q-1}b=a.
                                                               \tag{12.11}
\]
\end{theorem}

\begin{proof}
By (12.7),
\[
\begin{split}
 d_{q-1}\widehat K_q+\widehat K_{q+1}d_q
 &=
 \bigl(d_{q-1}K_q+K_{q+1}d_q\bigr)S_q\\
 &=(I-P_q-E_q)S_q\\
 &=(I-E_q)S_q-P_qS_q=I-P_q.
\end{split}                                                  \tag{12.12}
\]
The Neumann series gives
\(\|S_q\|\le(1-\kappa)^{-1}\), proving (12.10).  If \(d_qa=0\) and
\(P_qa=0\), (12.9) applied to \(a\) gives
\(d_{q-1}\widehat K_qa=a\).
\end{proof}

\begin{corollary}[Cohomology is carried by the matching sector]
\label{cor:v12v12-cohomology}
The inclusion of \(\operatorname{ran}P\) into \({\cal C}\) and the
cochain projection \(P\) induce inverse maps in cohomology.
\end{corollary}

\begin{proof}
The restriction of \(P\) to its range is the identity.  Equation
(12.9) is a cochain homotopy from the identity on \({\cal C}\) to
\(P\).  Homotopic cochain maps induce the same map in cohomology.
\end{proof}

This corollary identifies the correct role of \(P\).  It retains the
finite-dimensional global or block-matching data; everything in its
closed complement is removed by an exact counterterm.

\subsection{Exponential locality of the corrected homotopy}

Assume that \(K_q\) has propagation at most \(R_K\), \(E_q\) has
propagation at most \(R_E>0\), and
\[
                         \|K_q\|\le K_0,\qquad
                         \|E_q\|\le\kappa<1.                \tag{12.13}
\]
For \(r\ge0\), put
\[
 N_{K,E}(r):=
 \max\left\{0,
 \left\lceil\frac{r-R_K}{R_E}\right\rceil\right\}.           \tag{12.14}
\]

\begin{theorem}[Quasi-local block patching]
\label{thm:v12v12-local}
For all \(x,y\in\Lambda\),
\[
 \boxed{
 \|(\widehat K_q)_{xy}\|
 \le\frac{K_0}{1-\kappa}\,
       \kappa^{N_{K,E}(\operatorname{dist}(x,y))}.}          \tag{12.15}
\]
Consequently, for every
\[
             0<\mu<\frac1{R_E}\log\frac1\kappa,             \tag{12.16}
\]
there is a constant
\(C=C(K_0,\kappa,R_K,R_E,\mu)\), independent of the volume, such
that
\[
                    \|(\widehat K_q)_{xy}\|
 \le Ce^{-\mu\operatorname{dist}(x,y)}.                     \tag{12.17}
\]
\end{theorem}

\begin{proof}
Expanding (12.8) gives
\[
                         \widehat K_q
 =\sum_{n=0}^{\infty}K_qE_q^n.                             \tag{12.18}
\]
The \(n\)-th term has propagation at most \(R_K+nR_E\) and norm at
most \(K_0\kappa^n\).  If
\(r=\operatorname{dist}(x,y)\), every term with
\(n<N_{K,E}(r)\) therefore has zero \((x,y)\) block.  Summing the
remaining operator-norm bounds gives
\[
 \|(\widehat K_q)_{xy}\|
 \le K_0\sum_{n=N_{K,E}(r)}^\infty\kappa^n,
\]
which is (12.15).  The inequalities
\[
 N_{K,E}(r)\ge\frac{r-R_K}{R_E},
 \qquad
 \kappa^{N_{K,E}(r)}
 \le
 \exp\left[-\frac{\log(1/\kappa)}{R_E}(r-R_K)\right]         \tag{12.19}
\]
give (12.17), and any smaller exponent in (12.16) can be absorbed
into a uniform constant.
\end{proof}

Unlike V11, this theorem makes no assumption about the smallest
positive eigenvalue of the global Laplacian.  Locality follows from
the geometric-series suppression of repeated overlap errors.

\subsection{A finite verification target for block constructions}

Suppose a cover by uniformly bounded blocks has supplied provisional
local solvers.  After extension, restriction, and a partition of
unity, let their sum be \(K\), choose the retained matching
projection \(P\), and compute
\[
                         E=I-P-dK-Kd.                       \tag{12.20}
\]
The theorem requires the following certificate:

\begin{enumerate}
\item \(P^2=P=P^*\) and \(Pd=dP\);
\item \(KP=PK=0\), after replacing \(K\) by
      \((I-P)K(I-P)\) if this preserves the intended local solver;
\item finite propagation bounds for \(K\) and \(E\);
\item one uniform estimate \(\|E\|\le\kappa<1\).
\end{enumerate}

The first three items are exact algebraic and support checks.  The
fourth is the analytic overlap estimate.  It may be obtained by a
Schur bound, a block Fourier symbol, or a direct singular-value
certificate on the finite list of block types.  Once those four
items hold, no separate global spectral-gap estimate is needed for
the exact primitive (12.11).

\begin{proposition}[Stability margin for a certified patch]
\label{prop:v12v12-stability}
Let \(E\) obey \(\|E\|\le\kappa<1\), and let a perturbed block
construction have error \(\widetilde E\) with
\[
                         \|\widetilde E-E\|\le\varepsilon
                         <1-\kappa.                         \tag{12.21}
\]
Then \(\|\widetilde E\|\le\kappa+\varepsilon<1\), and its corrected
homotopy satisfies
\[
 \|\widetilde K(I-\widetilde E)^{-1}\|
 \le\frac{\|\widetilde K\|}{1-\kappa-\varepsilon}.          \tag{12.22}
\]
\end{proposition}

\begin{proof}
The first assertion is the triangle inequality.  The second follows
from the norm-convergent Neumann series for
\((I-\widetilde E)^{-1}\).
\end{proof}

This margin is the appropriate refinement quantity: adjacent
regulators may vary as long as their block certificates remain
inside a common open contraction cone.

\subsection{Integration with the anomaly ledger}

At regulator \(i\), let \(a_i^{\rm tot}\) be the total local anomaly
cochain after the common polynomial contribution has cancelled as in
V10.  If
\[
 d_i a_i^{\rm tot}=0,\qquad P_i a_i^{\rm tot}=0,            \tag{12.23}
\]
then
\[
                  b_i=\widehat K_i a_i^{\rm tot}            \tag{12.24}
\]
is an exact exponentially quasi-local counterterm.  If the
construction is uniform and
\(\sum_i\|a_i^{\rm tot}\|<\infty\), then
\[
 \sum_i\|b_i\|
 \le\frac{K_0}{1-\kappa}
       \sum_i\|a_i^{\rm tot}\|<\infty.                      \tag{12.25}
\]
The surviving obstruction is now the finite matching component
\(P_i a_i^{\rm tot}\), not a diffuse failure of locality.

There are two noninterchangeable ways for (12.23) to fail.
A nonzero \(d_i a_i^{\rm tot}\) is a Ward consistency defect.
A nonzero \(P_i a_i^{\rm tot}\) is a retained cohomological or
matching obstruction.  Neither can be repaired by the Neumann
correction, and both must remain explicit in the refinement ledger.

\subsection{Next concrete acquisition}

For the coupled SM--Einstein complex, one must now choose bounded
reference blocks and define \(P_i\) to retain their gauge,
diffeomorphism, and overlap-matching modes.  The next calculation is
not a global eigenvalue search.  It is the construction of a
relative block homotopy \(K_i\) with boundary conditions compatible
with the exact re-tree representation from V6 and the anomaly
transport from V8--V10, followed by an estimate of the finite-range
operator (12.20).

\begin{assessment}[Status after V12]
\label{ass:v12v12-status}
A global uniform Ward gap is no longer logically necessary for the
local counterterm step.  A finite-range approximate block homotopy
with overlap norm below one has been corrected to an exact,
volume-uniform, exponentially quasi-local homotopy.  This reduces the
next bottleneck to a block construction and the quantitative
certificate \(\|I-P-dK-Kd\|<1\).
\end{assessment}

\begin{firewall}[Claim boundary after V12]
\label{fire:v12v12-claim}
No Standard Model--Einstein block homotopy or contraction constant
has yet been computed.  The theorem does not show that the physical
Ward complex satisfies (12.5), eliminate a nonzero matching
component, construct the interacting continuum measure, or prove the
full SM--Einstein continuum.
\end{firewall}

\subsection{Parent record}

This V12 note appends to the validated V11 file with record
\[
\begin{array}{c|c}
\text{lines}&3459\\
\text{bytes}&121598\\
\text{SHA--256}&
\texttt{8FD77D5F074998F8264C279A42C2C5B96FF270AEFE342E20A3A3E4D17947DF0E}.
\end{array}
\]

% =============================================================================
% END APPENDED TWELFTH-CYCLE V12 MATERIAL
% =============================================================================
% =============================================================================
% BEGIN APPENDED TWELFTH-CYCLE V13 MATERIAL
% =============================================================================

\section{Twelfth-cycle V13: a quantitative block-cover certificate
for the Ward contraction}
\label{sec:v12v13}

\subsection{From bounded blocks to the V12 overlap operator}

V12 reduced quasi-local anomaly removal to a finite-range approximate
homotopy \(K\) whose error
\[
                         E=I-P-dK-Kd                       \tag{13.1}
\]
has norm less than one.  We now derive a sufficient estimate for
that norm from local Hodge solvers and a partition of unity.  The
result separates three quantities that a regulator calculation can
measure: the local block-solver norm, the commutator of the Ward
differential with the partition, and the mismatch of the retained
matching modes.

\subsection{Local Hodge solvers on a finite block library}

Let \({\cal B}\) be a cover of the finite metric set \(\Lambda\) by
blocks of uniformly bounded diameter.  On every block \(B\), suppose
there is a finite Hilbert cochain complex
\[
 \cdots\longrightarrow{\cal C}^{q-1}_B
 \xrightarrow{d_{B,q-1}}{\cal C}^{q}_B
 \xrightarrow{d_{B,q}}{\cal C}^{q+1}_B
 \longrightarrow\cdots.                                    \tag{13.2}
\]
Write
\[
 \Delta_{B,q}
 =d_{B,q-1}d_{B,q-1}^*+d_{B,q}^*d_{B,q},                  \tag{13.3}
\]
and let \(p_{B,q}\) be the orthogonal projection onto
\(\ker\Delta_{B,q}\).  Assume
\[
 \operatorname{spec}(\Delta_{B,q})
 \subset\{0\}\cup[\gamma_B,M_B].                            \tag{13.4}
\]
Define
\[
 G_{B,q}=\Delta_{B,q}^\dagger,\qquad
 h_{B,q}=d_{B,q-1}^*G_{B,q}.                               \tag{13.5}
\]

\begin{proposition}[Exact block Hodge contraction]
\label{prop:v12v13-block}
For every block and degree,
\[
 d_{B,q-1}h_{B,q}+h_{B,q+1}d_{B,q}=I-p_{B,q},              \tag{13.6}
\]
and
\[
                  \|h_{B,q}\|
 \le\frac{\|d_{B,q-1}\|}{\gamma_B}.                        \tag{13.7}
\]
If the blocks come from a finite library of rescaled block types,
the quantities
\[
 K_0:=\sup_{B,q}\|h_{B,q}\|,\qquad
 R_h:=\sup_{B,q}\operatorname{prop}(h_{B,q})                \tag{13.8}
\]
are finite.
\end{proposition}

\begin{proof}
The Hodge decomposition on the finite block gives
\(\Delta_BG_B=I-p_B\).  The intertwining
\(d_BG_{B,q}=G_{B,q+1}d_B\), proved as in V11, yields
\[
\begin{split}
 d_{B,q-1}h_{B,q}+h_{B,q+1}d_{B,q}
 &=
 d_{B,q-1}d_{B,q-1}^*G_{B,q}
 +d_{B,q}^*G_{B,q+1}d_{B,q}\\
 &=\Delta_{B,q}G_{B,q}=I-p_{B,q}.
\end{split}
\]
The gap in (13.4) gives
\(\|G_{B,q}\|\le\gamma_B^{-1}\), proving (13.7).  A finite library
has a positive minimum of its positive gaps and finite maxima of all
matrix norms and diameters.  Extension of a block operator by zero
has propagation no larger than the block diameter, proving (13.8).
\end{proof}

The proposition is useful precisely because the blocks remain
bounded after rescaling.  It does not assert a gap for a growing
global complex.

\subsection{Partition commutators}

For each degree \(q\), let
\(\chi_{B,q}:{\cal C}^q\to{\cal C}^q\) be a self-adjoint diagonal
multiplication operator supported in \(B\), with
\[
             0\le\chi_{B,q}\le I,\qquad
             \sum_{B\in{\cal B}}\chi_{B,q}^2=I.             \tag{13.9}
\]
Define the degree-\(q\) commutator
\[
 C_{B,q}:=
 d_q\chi_{B,q}-\chi_{B,q+1}d_q.                            \tag{13.10}
\]
After extending the block homotopies to the global complex, set
\[
 K_q:=\sum_{B\in{\cal B}}
       \chi_{B,q-1}h_{B,q}\chi_{B,q},                       \tag{13.11}
\]
and define the patched local-mode operator
\[
 P^{\rm loc}_q:=
 \sum_{B\in{\cal B}}
 \chi_{B,q}p_{B,q}\chi_{B,q}.                              \tag{13.12}
\]

The following identity is the central bookkeeping step.

\begin{lemma}[Exact partition defect formula]
\label{lem:v12v13-partition}
Assume that the local identity (13.6) is valid on the supports of
the corresponding cutoffs.  Then
\[
 d_{q-1}K_q+K_{q+1}d_q
 =I-P^{\rm loc}_q+Q_q,                                     \tag{13.13}
\]
where
\[
 Q_q=
 \sum_{B\in{\cal B}}
 \left(
 C_{B,q-1}h_{B,q}\chi_{B,q}
 -\chi_{B,q}h_{B,q+1}C_{B,q}
 \right).                                                   \tag{13.14}
\]
\end{lemma}

\begin{proof}
For a single block, the definition of the commutator gives
\[
 d_{q-1}\chi_{B,q-1}
 =\chi_{B,q}d_{q-1}+C_{B,q-1}                              \tag{13.15}
\]
and
\[
 \chi_{B,q+1}d_q
 =d_q\chi_{B,q}-C_{B,q}.                                   \tag{13.16}
\]
Substitution in the two terms containing \(K\) gives
\[
\begin{split}
 &d_{q-1}\chi_{B,q-1}h_{B,q}\chi_{B,q}
 +\chi_{B,q}h_{B,q+1}\chi_{B,q+1}d_q\\
 &\quad=
 \chi_{B,q}(d_{B,q-1}h_{B,q}
           +h_{B,q+1}d_{B,q})\chi_{B,q}\\
 &\qquad\quad+
 C_{B,q-1}h_{B,q}\chi_{B,q}
 -\chi_{B,q}h_{B,q+1}C_{B,q}.
\end{split}                                                 \tag{13.17}
\]
Use (13.6), sum over \(B\), and apply (13.9).  The result is
(13.13)--(13.14).
\end{proof}

\subsection{A volume-uniform norm estimate}

Assume the cover can be split into \(L\) colours so that, within each
colour, the supports of all operators appearing in each sum in
(13.14) are pairwise disjoint.  Such sums are block diagonal, so
their norm is the maximum of the block norms rather than the sum over
the number of blocks.  Suppose
\[
 \sup_{B,q}\|C_{B,q}\|\le\varepsilon,\qquad
 \sup_{B,q}\|h_{B,q}\|\le K_0.                             \tag{13.18}
\]
Let \(P_q\) be the desired global matching projection and put
\[
                         \delta:=
 \sup_q\|P^{\rm loc}_q-P_q\|.                              \tag{13.19}
\]

\begin{theorem}[Finite block-cover contraction certificate]
\label{thm:v12v13-certificate}
Under the preceding assumptions,
\[
 \left\|I-P_q-d_{q-1}K_q-K_{q+1}d_q\right\|
 \le \delta+2L\varepsilon K_0.                             \tag{13.20}
\]
Assume also that \(P\) is a cochain projection,
\(K_qP_q=P_{q-1}K_q=0\), and
\[
                         \boxed{
 \delta+2L\varepsilon K_0\le\kappa<1.}                     \tag{13.21}
\]
Then the corrected homotopy
\[
 \widehat K_q
 =K_q\left(
 I-\bigl[I-P_q-d_{q-1}K_q-K_{q+1}d_q\bigr]
 \right)^{-1}                                               \tag{13.22}
\]
is exact:
\[
 d_{q-1}\widehat K_q+\widehat K_{q+1}d_q=I-P_q.            \tag{13.23}
\]
If \(d,\chi_B,h_B\), and \(P\) have uniform finite propagation, then
\(\widehat K\) has a volume-uniform exponentially decaying block
kernel.
\end{theorem}

\begin{proof}
By Lemma~\ref{lem:v12v13-partition},
\[
 I-P_q-dK-Kd=P^{\rm loc}_q-P_q-Q_q.                        \tag{13.24}
\]
Within each colour, the first sum in (13.14) has norm at most
\(\varepsilon K_0\), and the second has the same bound.
The triangle inequality over the \(L\) colours gives
\(\|Q_q\|\le2L\varepsilon K_0\).  Equations (13.19) and (13.24)
prove (13.20).  Condition (13.21) and the side identities place the
construction under Theorems~\ref{thm:v12v12-correction} and
\ref{thm:v12v12-local}, which give (13.23) and exponential locality.
\end{proof}

\begin{remark}[Meaning of the three factors]
\label{rem:v12v13-factors}
The quantity \(K_0\) is fixed by the worst local positive spectral
gap.  The commutator \(\varepsilon\) measures the variation of the
partition across one interaction range of the Ward differential.
The colour number \(L\) measures overlap geometry.  The term
\(\delta\) measures whether the locally retained modes assemble into
the chosen global matching sector.  None of these quantities depends
on the total number of blocks when the hypotheses are uniform.
\end{remark}

\subsection{Closed anomalies and the matching obstruction}

Let \(a\in{\cal C}^q\) be a closed total anomaly remainder.  Under
(13.21), (13.23) gives
\[
                    a=P_qa+d_{q-1}\widehat K_qa.            \tag{13.25}
\]
The second term is removed by the exponentially quasi-local
counterterm \(\widehat K_qa\).  The first is the finite matching
obstruction.

\begin{corollary}[Quantitative counterterm bound]
\label{cor:v12v13-counterterm}
Under the hypotheses of
Theorem~\ref{thm:v12v13-certificate}, if \(d_qa=0\) and \(P_qa=0\),
then
\[
 b=\widehat K_qa,\qquad d_{q-1}b=a,\qquad
 \|b\|\le
 \frac{LK_0}{1-\delta-2L\varepsilon K_0}\,\|a\|.            \tag{13.26}
\]
\end{corollary}

\begin{proof}
Exactness follows from (13.23).  Within one colour the summands in
(13.11) are block diagonal and have norm at most \(K_0\).  The
triangle inequality over the \(L\) colours therefore gives
\(\|K\|\le LK_0\).  The Neumann bound from V12 gives
\[
 \|\widehat K\|
 \le\frac{\|K\|}
          {1-\|I-P-dK-Kd\|}.
\]
Combining this estimate with (13.20) proves (13.26).
\end{proof}

\begin{remark}[Possible improvement of the numerator]
\label{rem:v12v13-numerator}
A separately verified normalized synthesis estimate may improve
\(\|K\|\le LK_0\) to \(\|K\|\le K_0\).  Such an improvement changes
only the numerator of (13.26); it must be certified and is not used
in the contraction condition (13.21).
\end{remark}

\subsection{Consequences for the next regulator calculation}

Theorem~\ref{thm:v12v13-certificate} converts the next stage into a
finite list of regulator data.

\begin{enumerate}
\item Specify the bounded gauge--diffeomorphism block complexes and
      their boundary conditions.
\item Enumerate the finite block-type library and compute the exact
      positive block gaps entering \(K_0\).
\item Choose a quadratic partition and calculate the commutator
      matrices (13.10).
\item Identify the local harmonic modes, their global matching
      projection, and the mismatch \(\delta\).
\item Verify (13.21) together with the normalized synthesis bound for
      \(K\).
\end{enumerate}

The exact re-tree representation from V6 must act on this entire
block library.  The nonlinear regulator commonality requirement from
V10 must hold before the resulting local anomaly sources are summed.
The block certificate then removes only the closed component
orthogonal to \(P\); it does not alter the harmonic anomaly ledger of
V8--V10.

\begin{assessment}[Status after V13]
\label{ass:v12v13-status}
The abstract contraction threshold of V12 has been reduced to the
explicit block-cover inequality
\(\delta+2L\varepsilon K_0<1\).  Bounded local Hodge solvers supply
\(K_0\), partition commutators supply \(\varepsilon\), cover geometry
supplies \(L\), and matching-mode assembly supplies \(\delta\).
This is a finite, volume-independent certificate once a concrete
SM--Einstein regulator block is fixed.
\end{assessment}

\begin{firewall}[Claim boundary after V13]
\label{fire:v12v13-claim}
No physical block library, boundary condition, matching projection,
or numerical contraction certificate has yet been supplied for the
coupled theory.  In particular, (13.21) is a sufficient criterion,
not an established property of the Standard Model--Einstein Ward
complex, and the full continuum remains open.
\end{firewall}

\subsection{Parent record}

This V13 note appends to the validated V12 file with record
\[
\begin{array}{c|c}
\text{lines}&3801\\
\text{bytes}&132968\\
\text{SHA--256}&
\texttt{A1E8CAA721244A5E8A16656CBBF26F7C9282D2664C315A2B60A717649BE9BC0D}.
\end{array}
\]

% =============================================================================
% END APPENDED TWELFTH-CYCLE V13 MATERIAL
% =============================================================================
% =============================================================================
% BEGIN APPENDED TWELFTH-CYCLE V14 MATERIAL
% =============================================================================

\section{Twelfth-cycle V14: symmetry-equivariant block contractions
without loss of the overlap margin}
\label{sec:v12v14}

\subsection{The covariance requirement}

The V13 block certificate is useful for the coupled theory only if
the block differential, retained matching modes, and counterterm
primitive transform under the exact physical re-tree and internal
symmetry actions.  An arbitrary choice of a block base point or a
local pseudoinverse can conceal this covariance.  We now prove that
a valid finite-range approximate contraction can be made equivariant
without worsening its contraction norm.

The statement applies to an exact unitary action on the
finite-regulator linear Ward complex.  It includes finite block
automorphism and re-tree groups, and compact internal groups after a
unitary finite-dimensional realization has been fixed.  It does not
replace the nonlinear mixed-jet validation required in V10.

\subsection{Averaging an approximate cochain homotopy}

Let \(\Gamma\) be a compact group with normalized Haar measure
\(dg\).  Suppose that every cochain space carries a unitary
representation
\[
                 U_{g,q}:{\cal C}^q\longrightarrow{\cal C}^q
                                                               \tag{14.1}
\]
such that
\[
             d_qU_{g,q}=U_{g,q+1}d_q.                       \tag{14.2}
\]
Let \(P_q\) be an orthogonal cochain projection satisfying
\[
                 U_{g,q}P_q=P_qU_{g,q}.                    \tag{14.3}
\]
Given a degree-minus-one operator \(K_q\), define
\[
 \overline K_q
 :=
 \int_\Gamma
 U_{g,q-1}^{-1}K_qU_{g,q}\,dg.                             \tag{14.4}
\]
All integrals are ordinary Bochner integrals in finite-dimensional
operator spaces.

\begin{theorem}[Equivariantization preserves the contraction
threshold]
\label{thm:v12v14-average}
Let
\[
 E_q=I-P_q-d_{q-1}K_q-K_{q+1}d_q                           \tag{14.5}
\]
and define \(\overline E_q\) by the same formula with
\(\overline K\).  Then
\[
 \overline E_q
 =
 \int_\Gamma U_{g,q}^{-1}E_qU_{g,q}\,dg,                   \tag{14.6}
\]
and
\[
 \|\overline K_q\|\le\|K_q\|,
 \qquad
 \|\overline E_q\|\le\|E_q\|.                              \tag{14.7}
\]
The operators \(\overline K\) and \(\overline E\) commute with the
group action in their respective graded senses.  If
\[
 K_qP_q=P_{q-1}K_q=0,                                      \tag{14.8}
\]
the same identities hold for \(\overline K\).

If \(\|\overline E_q\|\le\kappa<1\), then
\[
 \widehat K_q^\Gamma
 =\overline K_q(I-\overline E_q)^{-1}                       \tag{14.9}
\]
is equivariant and satisfies
\[
 d_{q-1}\widehat K_q^\Gamma+
 \widehat K_{q+1}^\Gamma d_q=I-P_q.                         \tag{14.10}
\]
\end{theorem}

\begin{proof}
Use (14.2) to move \(d\) through \(U_g\), and use (14.3) for \(P\).
For every \(g\),
\[
\begin{split}
 &I-P_q-d_{q-1}
   U_{g,q-1}^{-1}K_qU_{g,q}
 -U_{g,q}^{-1}K_{q+1}U_{g,q+1}d_q\\
 &\qquad=U_{g,q}^{-1}
 \bigl(I-P_q-d_{q-1}K_q-K_{q+1}d_q\bigr)U_{g,q}.
\end{split}                                                  \tag{14.11}
\]
Integration proves (14.6).  Unitary conjugation preserves the
operator norm, and the norm of an integral is at most the integral
of the norms, proving (14.7).

Left invariance of Haar measure shows
\[
 U_{h,q-1}^{-1}\overline K_qU_{h,q}=\overline K_q           \tag{14.12}
\]
for every \(h\), and similarly for \(\overline E\).
Averaging (14.8) proves the side identities for \(\overline K\).
The Neumann series for
\((I-\overline E)^{-1}\) is equivariant term by term.
Equation (14.10) is then
Theorem~\ref{thm:v12v12-correction}.
\end{proof}

\begin{corollary}[Symmetry does not consume overlap margin]
\label{cor:v12v14-margin}
If the unsymmetrized block construction satisfies
\(\sup_q\|E_q\|\le\kappa<1\), its averaged construction satisfies
the same bound.  It therefore has an exact equivariant correction
with norm
\[
 \|\widehat K_q^\Gamma\|
 \le\frac{\|K_q\|}{1-\kappa}.                              \tag{14.13}
\]
\end{corollary}

\begin{proof}
Apply (14.7) and (14.9).
\end{proof}

The estimate may improve strictly under averaging because
noninvariant pieces of the overlap defect can cancel.  No such
improvement is assumed.

\subsection{Preservation of spatial locality}

Assume \(\Gamma\) acts on \(\Lambda\) by isometries and that
\(U_{g,q}\) maps the fibre at \(x\) unitarily to the fibre at \(gx\).
If \(K_q\) has propagation at most \(R_K\), then every conjugate in
(14.4) has the same propagation.  Hence so does \(\overline K_q\).
The analogous statement holds for \(E_q\).

\begin{proposition}[Equivariant quasi-local correction]
\label{prop:v12v14-local}
Under the preceding geometric action, suppose
\[
 \operatorname{prop}(K_q)\le R_K,\qquad
 \operatorname{prop}(E_q)\le R_E,\qquad
 \|E_q\|\le\kappa<1.                                      \tag{14.14}
\]
Then \(\widehat K^\Gamma\) is equivariant and has the exponential
kernel estimate
\[
 \|(\widehat K_q^\Gamma)_{xy}\|
 \le
 \frac{\|K_q\|}{1-\kappa}\,
 \kappa^{\max\{0,\lceil
   (\operatorname{dist}(x,y)-R_K)/R_E\rceil\}}.             \tag{14.15}
\]
\end{proposition}

\begin{proof}
Averaging preserves the two propagation bounds and does not increase
the norms.  Apply Theorem~\ref{thm:v12v12-local} to the averaged
operators.
\end{proof}

\subsection{Construction from block-orbit representatives}

Let \(\Gamma\) act on the block cover \({\cal B}\).  Choose one
representative \(B\) in every block orbit and let
\(\Gamma_B=\{g:gB=B\}\) be its stabilizer.  Suppose \(h_{B,q}\) is a
local Hodge contraction and \(p_{B,q}\) its harmonic projection.
Average over the stabilizer:
\[
 h_{B,q}^{\rm st}
 =
 \int_{\Gamma_B}
 U_{g,q-1}^{-1}h_{B,q}U_{g,q}\,dg.                         \tag{14.16}
\]

\begin{proposition}[Orbit-representative block library]
\label{prop:v12v14-orbits}
Assume the block differential and \(p_B\) commute with
\(\Gamma_B\).  Then
\[
 d_{B,q-1}h_{B,q}^{\rm st}
 +h_{B,q+1}^{\rm st}d_{B,q}=I-p_{B,q},                     \tag{14.17}
\]
and \(\|h_{B,q}^{\rm st}\|\le\|h_{B,q}\|\).
For \(B'=gB\), define
\[
 h_{B',q}
 =U_{g,q-1}h_{B,q}^{\rm st}U_{g,q}^{-1}.                   \tag{14.18}
\]
This definition is independent of the choice of \(g\), and all
blocks in the orbit have identical gap, norm, and propagation
bounds.
\end{proposition}

\begin{proof}
Average the exact block identity over the stabilizer to obtain
(14.17).  The norm estimate follows as in (14.7).
If \(g_1B=g_2B\), then \(g_2^{-1}g_1\in\Gamma_B\).  Stabilizer
invariance of \(h_B^{\rm st}\) makes the two definitions in (14.18)
equal.  Unitary conjugation preserves spectra and norms, while the
isometric site action preserves propagation.
\end{proof}

Thus the finite certificate from V13 needs to be computed only on
block-orbit representatives, including one representative of every
inequivalent boundary condition and field-content type.

\subsection{The matching projection must be chosen before averaging}

If \(V_q\subset{\cal C}^q\) is a \(\Gamma\)-invariant subspace, its
orthogonal projector \(P_q\) commutes with the unitary action.
Indeed, both \(V_q\) and \(V_q^\perp\) are invariant.  This supplies
(14.3) whenever the retained matching space has already been proved
invariant.

By contrast, the average of orthogonal projectors onto
noninvariant subspaces is generally a positive contraction rather
than a projection.  Therefore one must first define the invariant
matching subcomplex and then take its orthogonal projector.  Merely
averaging a channel-dependent matching projector does not establish
the hypotheses of V12.

\subsection{Consequences for the coupled regulator}

The exact finite-dimensional certificate now has the following
symmetry-reduced form.

\begin{enumerate}
\item Construct one local Ward complex and Hodge solver for each
      orbit representative of the bounded block library.
\item Prove that the retained block modes form a subrepresentation
      of the exact physical re-tree and internal symmetry action.
\item Stabilizer-average each representative solver and transport it
      by (14.18).
\item Assemble the partitioned operator of V13 and verify
      \(\delta+2L\varepsilon K_0<1\).
\item Average the assembled approximate homotopy if necessary and
      apply (14.9).
\end{enumerate}

The resulting primitive is simultaneously exact, quasi-local, and
equivariant under the verified linear action.  V10 still requires a
separate calculation for nonlinear background-dependent transport.
The present averaging theorem cannot infer that nonlinear
commonality from its linearization.

\begin{assessment}[Status after V14]
\label{ass:v12v14-status}
Symmetry compatibility no longer costs contraction margin: compact
unitary averaging preserves the approximate homotopy identity,
cannot increase the overlap norm, preserves finite propagation under
an isometric site action, and yields an exact equivariant
quasi-local correction.  The block computation reduces to orbit
representatives and stabilizer-invariant solvers.
\end{assessment}

\begin{firewall}[Claim boundary after V14]
\label{fire:v12v14-claim}
No exact compact-group action for the full nonlinear
SM--Einstein regulator has been constructed here.  BRST
nilpotence, noncompact gauge fixing, diffeomorphism covariance,
mixed background jets, the block contraction inequality, and the
continuum measure remain open.
\end{firewall}

\subsection{Parent record}

This V14 note appends to the validated V13 file with record
\[
\begin{array}{c|c}
\text{lines}&4145\\
\text{bytes}&144410\\
\text{SHA--256}&
\texttt{2AFCFA4E0AF5A761E88B1FF4273B37039856FC438EF0E6D855F23E9C23DDD802}.
\end{array}
\]

% =============================================================================
% END APPENDED TWELFTH-CYCLE V14 MATERIAL
% =============================================================================
% =============================================================================
% BEGIN APPENDED TWELFTH-CYCLE V15 MATERIAL
% =============================================================================

\section{Twelfth-cycle V15: compulsory companion audit and exact
separation of the compact gauge-fibre measure}
\label{sec:v12v15}

\subsection{Audit record}

This is the compulsory five-version audit following V10.  The newest
complete Yang--Mills file read is V46; it contains the V45
product-Haar result and the V46 direct compiler regression.  The
Navier--Stokes companion remains V26.  Their records are
\[
\begin{array}{c|r|r|c}
\text{record}&\text{lines}&\text{bytes}&\text{SHA--256}\\ \hline
\text{YM V46}&18335&627990&
\texttt{8D00DD00E4F38552262670E57FE9EA3810EFF75F3A2A7DC115BCCDBCF3A97B39}\\
\text{NS V26}&5319&278283&
\texttt{82C887F8C2C69E4F28FAD7C3DC8DE5B9099CB5A4BF431EBA1FFF3E91935978AD}.
\end{array}                                                   \tag{15.1}
\]

No Navier--Stokes result has changed since the V10 audit.  Its
rank-one kernel estimates therefore do not change the current
SM--Einstein acquisition order.

\subsection{The new Yang--Mills conclusions}

YM V45 proves exact product-Haar disintegration for a compact lattice
gauge group in spanning-tree gauge.  It identifies the retained
boundary links with coarse holonomies and proves that the nonlinear
fixed-boundary re-tree map preserves vertical product Haar measure
for every boundary value.  Consequently the complete
background/momentum mixed response of the re-tree coordinate
Jacobian plus local Haar density is zero.

YM V46 independently realizes the relevant corner phases without
the Walsh batching used earlier.  It reproduces the same nonzero
second-jet residual, thereby clearing the Walsh transform and corner
extraction.  Combined with the exact Haar result, this locates the
residual in the shared upstream separated vertex grammar.  It is a
compiler consistency defect rather than a physical anomaly or a
missing measure counterterm.

These proofs are not repeated here.  We extract one new consequence
for a coupled measure by applying the disintegration identity before
fermion, Higgs, ghost, and gravitational densities are combined.

\subsection{A conditional-measure separation theorem}

Let \((X,\mu)\) and \((Y,\nu)\) be finite or standard Borel measure
spaces, with \(\mu\) a probability measure.  Let
\[
 T(x,y)=(\phi_y(x),\psi(y))                                 \tag{15.2}
\]
be an invertible measurable map.  Assume that every
\(\phi_y:X\to X\) preserves \(\mu\), and that \(\psi_*\nu\) is
absolutely continuous with respect to \(\nu\).  Write
\[
 r_\psi(y')=\frac{d(\psi_*\nu)}{d\nu}(y').                  \tag{15.3}
\]
Let an interacting measure be
\[
                         d{\mathbb M}(x,y)
 =w(x,y)\,d\mu(x)\,d\nu(y),                                \tag{15.4}
\]
where \(w>0\) almost everywhere.

\begin{theorem}[No coordinate Jacobian from a conditionally
invariant fibre]
\label{thm:v12v15-disintegration}
For almost every \((x',y')\),
\[
 \frac{d(T_*{\mathbb M})}{d(\mu\otimes\nu)}(x',y')
 =
 w\!\left(
 \phi_{\psi^{-1}y'}^{-1}(x'),\psi^{-1}y'
 \right)r_\psi(y').                                        \tag{15.5}
\]
Consequently,
\[
 \boxed{
 \frac{d(T_*{\mathbb M})}{d{\mathbb M}}(x',y')
 =
 \frac{
 w\!\left(\phi_{\psi^{-1}y'}^{-1}(x'),\psi^{-1}y'\right)}
 {w(x',y')}\,r_\psi(y').}                                  \tag{15.6}
\]
There is no additional Radon--Nikodym factor from the fibre maps
\(\phi_y\).
\end{theorem}

\begin{proof}
For every bounded measurable test function \(F\),
\[
 \int F\,d(T_*{\mathbb M})
 =
 \int_Y\int_X
 F(\phi_y(x),\psi(y))w(x,y)\,d\mu(x)\,d\nu(y).              \tag{15.7}
\]
For fixed \(y\), change variables \(x'=\phi_y(x)\).  Fibre invariance
gives \(d\mu(x')=d\mu(x)\).  Then put \(y'=\psi(y)\) and use (15.3).
The resulting integral has density (15.5) with respect to
\(\mu\otimes\nu\).  Division by the positive density \(w(x',y')\)
gives (15.6).
\end{proof}

\begin{corollary}[Compact gauge-fibre contribution to the anomaly
ledger]
\label{cor:v12v15-gauge-ledger}
Take \(X\) to be the vertical compact-gauge links in spanning-tree
coordinates and \(Y\) to contain retained holonomies and all
fermion, Higgs, ghost, and gravitational variables.  Under the exact
fixed-boundary re-tree map proved in YM V45, the vertical
product-Haar coordinate contribution to the logarithmic
Radon--Nikodym cocycle is identically zero.  The remaining cocycle is
the sum of
\[
 \log w(T^{-1}(x',y'))-\log w(x',y')
 \quad\hbox{and}\quad
 \log r_\psi(y').                                          \tag{15.8}
\]
\end{corollary}

\begin{proof}
Apply Theorem~\ref{thm:v12v15-disintegration} and take logarithms.
The YM V45 theorem supplies the conditional invariance of the fibre
measure.
\end{proof}

This removes one possible source from the coupled anomaly ledger at
every nonlinear mixed order.  It does not remove the transformation
of the fermion determinant contained in \(w\), the ghost density, or
the base transformation of gravitational and retained variables.

\subsection{Compatibility with the V10 universality theorem}

The compact gauge re-tree measure is species-independent.  Hence its
absence in (15.8) is stronger than cancellation after summing
species: the factor is exactly one before the representation
coefficients are summed.  The anomaly-polynomial contribution still
cancels by Theorem~\ref{thm:v12v10-common} only when all species use
one common regulator descent map.  The two statements address
different factors:
\[
\begin{array}{c|c}
\text{factor}&\text{required argument}\\ \hline
\text{vertical compact-gauge coordinate measure}&
\text{conditional Haar invariance}\\
\text{fermion determinant phase and local descent}&
\text{species-common regulator plus }\sum_sA_s=0\\
\text{retained/gravitational base measure}&
\text{a separate Radon--Nikodym calculation}.
\end{array}                                                   \tag{15.9}
\]

The YM V46 finding also fixes the validation order.  A separated
diagrammatic expansion may be used only after a monolithic
finite-matrix derivative has been frozen and checked against exact
coordinate covariance.  A residual that violates such an identity
belongs to the consistency ledger until the shared compiler has been
repaired.

\subsection{Adjustment of the SM--Einstein direction}

The V13--V14 block programme remains the correct next step, with two
refinements.

\begin{enumerate}
\item The compact gauge-fibre measure is installed as an exactly
      invariant factor and is not charged to the overlap error
      \(\delta+2L\varepsilon K_0\).
\item The coupled block complex should first be assembled from
      sectorwise contractions.  Its uncoupled tensor differential
      can be contracted exactly; gauge--matter and
      gravity--matter mixing should then be entered as a controlled
      differential perturbation.
\end{enumerate}

The next note carries out the first of these algebraic steps: an
explicit tensor-product contraction onto the tensor product of the
sector matching complexes, with quasi-locality and norm bounds.

\begin{assessment}[Status after V15]
\label{ass:v12v15-status}
The mandatory audit is complete.  Exact conditional Haar invariance
removes the compact vertical gauge-coordinate measure from the
coupled nonlinear anomaly ledger.  The YM residual is confirmed to
be an upstream compiler inconsistency and supplies a validation rule,
not a physical counterterm.  The next constructive step is the
sectorwise assembly of the coupled block contraction.
\end{assessment}

\begin{firewall}[Claim boundary after V15]
\label{fire:v12v15-claim}
The conditional separation theorem does not prove invariance of the
fermion determinant, ghost measure, Higgs density, retained
holonomies, or gravitational measure.  It does not construct the
coupled block differential, verify its overlap threshold, or produce
the full SM--Einstein continuum.
\end{firewall}

\subsection{Parent record}

This V15 note appends to the validated V14 file with record
\[
\begin{array}{c|c}
\text{lines}&4435\\
\text{bytes}&154459\\
\text{SHA--256}&
\texttt{FD014513CF811FA9C02885AC626B98F8A4A1D3CEE3F2D6A08807D0C941C7F721}.
\end{array}
\]

% =============================================================================
% END APPENDED TWELFTH-CYCLE V15 MATERIAL
% =============================================================================
% =============================================================================
% BEGIN APPENDED TWELFTH-CYCLE V16 MATERIAL
% =============================================================================

\section{Twelfth-cycle V16: exact tensor-product contraction for the
sectorwise reference complex}
\label{sec:v12v16}

\subsection{Purpose}

The V13 certificate is stated for one coupled Ward complex, while a
finite-regulator construction normally begins with separately
controlled gauge, matter, and gravitational block complexes.  This
section gives the exact graded tensor-product contraction that
assembles such sectorwise data.  It identifies the product matching
space, controls the homotopy norm, and preserves quasi-locality.

The theorem concerns the uncoupled reference differential.  Mixing
terms in the physical BRST/Ward differential will be treated as a
differential perturbation; they are not silently included in the
tensor formula.

\subsection{Two contracted cochain complexes}

Let
\[
 ({\cal C},d_{\cal C}),\qquad({\cal D},d_{\cal D})           \tag{16.1}
\]
be finite Hilbert cochain complexes.  Suppose there are orthogonal
cochain projections \(P_{\cal C},P_{\cal D}\) and degree-minus-one
operators \(h_{\cal C},h_{\cal D}\) satisfying
\[
\begin{split}
 d_{\cal C}h_{\cal C}+h_{\cal C}d_{\cal C}&=I-P_{\cal C},\\
 d_{\cal D}h_{\cal D}+h_{\cal D}d_{\cal D}&=I-P_{\cal D},
\end{split}                                                  \tag{16.2}
\]
with the normalized side conditions
\[
 h_{\cal C}P_{\cal C}=P_{\cal C}h_{\cal C}=0,\qquad
 h_{\cal D}P_{\cal D}=P_{\cal D}h_{\cal D}=0.               \tag{16.3}
\]
Let \(S_{\cal C}\) be the parity operator,
\[
                  S_{\cal C}c=(-1)^p c
 \quad(c\in{\cal C}^p).                                    \tag{16.4}
\]
Then
\[
 S_{\cal C}d_{\cal C}=-d_{\cal C}S_{\cal C},\qquad
 S_{\cal C}h_{\cal C}=-h_{\cal C}S_{\cal C},                \tag{16.5}
\]
and \(S_{\cal C}\) commutes with \(P_{\cal C}\).

On the total complex
\[
 {\cal T}={\cal C}\widehat\otimes{\cal D},                  \tag{16.6}
\]
use the graded differential
\[
 \partial=d_{\cal C}\otimes I+
          S_{\cal C}\otimes d_{\cal D}.                    \tag{16.7}
\]

\begin{theorem}[Tensor-product contraction]
\label{thm:v12v16-tensor}
Define
\[
 P_{\cal T}=P_{\cal C}\otimes P_{\cal D}                    \tag{16.8}
\]
and
\[
 \boxed{
 h_{\cal T}
 =
 h_{\cal C}\otimes I+
 S_{\cal C}P_{\cal C}\otimes h_{\cal D}.}                  \tag{16.9}
\]
Then
\[
 \boxed{
 \partial h_{\cal T}+h_{\cal T}\partial=I-P_{\cal T}.}      \tag{16.10}
\]
Moreover,
\[
 h_{\cal T}P_{\cal T}=P_{\cal T}h_{\cal T}=0,\qquad
 \|h_{\cal T}\|
 \le\|h_{\cal C}\|+\|h_{\cal D}\|.                         \tag{16.11}
\]
\end{theorem}

\begin{proof}
The anticommutator of
\(d_{\cal C}\otimes I\) with \(h_{\cal C}\otimes I\) is
\[
                    (I-P_{\cal C})\otimes I.                \tag{16.12}
\]
The two cross terms containing
\(S_{\cal C}\otimes d_{\cal D}\) and
\(h_{\cal C}\otimes I\) cancel because
\(S_{\cal C}h_{\cal C}=-h_{\cal C}S_{\cal C}\).
The two cross terms containing
\(d_{\cal C}\otimes I\) and
\(S_{\cal C}P_{\cal C}\otimes h_{\cal D}\) cancel because
\(d_{\cal C}S_{\cal C}=-S_{\cal C}d_{\cal C}\) and
\(d_{\cal C}P_{\cal C}=P_{\cal C}d_{\cal C}\).
The remaining anticommutator is
\[
 P_{\cal C}\otimes
 (d_{\cal D}h_{\cal D}+h_{\cal D}d_{\cal D})
 =
 P_{\cal C}\otimes(I-P_{\cal D}).                          \tag{16.13}
\]
Adding (16.12) and (16.13) gives
\[
 (I-P_{\cal C})\otimes I+
 P_{\cal C}\otimes(I-P_{\cal D})
 =I-P_{\cal C}\otimes P_{\cal D},
\]
which proves (16.10).

The side conditions follow by multiplying (16.9) by
\(P_{\cal C}\otimes P_{\cal D}\) and using (16.3).  Since parity and
orthogonal projections have norm at most one, the triangle
inequality gives the norm bound in (16.11).
\end{proof}

\begin{corollary}[Product cohomology representatives]
\label{cor:v12v16-cohomology}
The tensor matching complex
\(\operatorname{ran}P_{\cal C}\widehat\otimes
 \operatorname{ran}P_{\cal D}\)
is cochain-homotopy equivalent to the total reference complex.
In particular, if the projections select harmonic representatives,
the retained total representatives are their graded tensor
products.
\end{corollary}

\begin{proof}
Equation (16.10) is a homotopy from the identity to
\(P_{\cal T}\).  The inclusion and projection are inverse on the
range of \(P_{\cal T}\), so they induce inverse maps in cohomology.
\end{proof}

\subsection{Quasi-locality on the product block}

Suppose
\[
 {\cal C}=\bigoplus_{x\in\Lambda_{\cal C}}C_x,\qquad
 {\cal D}=\bigoplus_{u\in\Lambda_{\cal D}}D_u               \tag{16.14}
\]
and equip
\(\Lambda_{\cal C}\times\Lambda_{\cal D}\) with
\[
 \operatorname{dist}_{\cal T}((x,u),(y,v))
 =
 \operatorname{dist}_{\cal C}(x,y)+
 \operatorname{dist}_{\cal D}(u,v).                        \tag{16.15}
\]
Assume
\[
\begin{split}
 \|(h_{\cal C})_{xy}\|&\le A_{\cal C}
 e^{-\mu_{\cal C}\operatorname{dist}_{\cal C}(x,y)},\\
 \|(P_{\cal C})_{xy}\|&\le B_{\cal C}
 e^{-\nu_{\cal C}\operatorname{dist}_{\cal C}(x,y)},\\
 \|(h_{\cal D})_{uv}\|&\le A_{\cal D}
 e^{-\mu_{\cal D}\operatorname{dist}_{\cal D}(u,v)}.
\end{split}                                                  \tag{16.16}
\]

\begin{proposition}[Product quasi-locality]
\label{prop:v12v16-local}
With
\[
 \eta=\min\{\mu_{\cal C},\nu_{\cal C},\mu_{\cal D}\},       \tag{16.17}
\]
the tensor homotopy has
\[
 \|(h_{\cal T})_{(x,u),(y,v)}\|
 \le
 \bigl(A_{\cal C}+B_{\cal C}A_{\cal D}\bigr)
 e^{-\eta\operatorname{dist}_{\cal T}((x,u),(y,v))}.        \tag{16.18}
\]
\end{proposition}

\begin{proof}
The first term of (16.9) has block kernel
\[
 (h_{\cal C})_{xy}\,\delta_{uv}.                            \tag{16.19}
\]
When it is nonzero, the \({\cal D}\)-distance is zero, so the first
bound in (16.16) gives the required product-distance estimate.
The second term has block norm at most
\[
 B_{\cal C}A_{\cal D}
 e^{-\nu_{\cal C}\operatorname{dist}_{\cal C}(x,y)}
 e^{-\mu_{\cal D}\operatorname{dist}_{\cal D}(u,v)},        \tag{16.20}
\]
which is bounded by the second contribution in (16.18).  Add the two
bounds.
\end{proof}

The projector bound in (16.16) is automatic when \(d_{\cal C}\) has
finite propagation and \(h_{\cal C}\) is exponentially quasi-local,
because (16.2) gives
\[
 P_{\cal C}=I-d_{\cal C}h_{\cal C}-h_{\cal C}d_{\cal C}.    \tag{16.21}
\]
Finite-range multiplication shifts the exponential kernel by only a
fixed distance and changes its constant by a fixed factor.

\begin{corollary}[Finite propagation version]
\label{cor:v12v16-propagation}
If \(h_{\cal C},P_{\cal C},h_{\cal D}\) have propagations
\(R_{\cal C},R_P,R_{\cal D}\), then
\[
 \operatorname{prop}(h_{\cal T})
 \le\max\{R_{\cal C},R_P+R_{\cal D}\}.                      \tag{16.22}
\]
\end{corollary}

\begin{proof}
This follows directly from the two kernels in (16.9) and the product
metric (16.15).
\end{proof}

\subsection{Iteration over all reference sectors}

Let \(j=1,\ldots,m\) label finitely many sector complexes with
contractions \((P_j,h_j)\).  Iterating
Theorem~\ref{thm:v12v16-tensor} gives
\[
                         P^{(m)}=\bigotimes_{j=1}^mP_j      \tag{16.23}
\]
and an exact total homotopy.  Inductively,
\[
                         \|h^{(m)}\|
 \le\sum_{j=1}^m\|h_j\|.                                   \tag{16.24}
\]
If every sector contraction and projection is uniformly quasi-local,
the finite tensor product is uniformly quasi-local with decay rate
at least the minimum of the sector rates.

For a reference SM--Einstein block, the labels may include compact
gauge, Higgs, fermion, ghost, and linearized gravitational sectors.
The statement is algebraic and does not assert that these physical
sectors are independent in the full differential.

\subsection{Equivariance}

Suppose each sector carries an exact unitary cochain representation
and its \(P_j,h_j\) are equivariant.  The graded tensor
representation commutes with the total differential.  Every term in
the iterated version of (16.9) is equivariant, so the assembled
homotopy is equivariant.  Thus the V14 orbit-representative reduction
can be performed sectorwise before tensor assembly.

\subsection{The coupled anomaly equation at reference level}

Let \(a\) be a closed total cochain for the reference differential
\(\partial\).  Equation (16.10) gives
\[
                         a=P_{\cal T}a+
                           \partial h_{\cal T}a.             \tag{16.25}
\]
When \(P_{\cal T}a=0\), the quasi-local primitive
\(b=h_{\cal T}a\) removes the reference anomaly exactly.  When the
projection is nonzero, it is a product-sector matching obstruction
and must be tested by the anomaly transport ledger of V8--V10.

The physical coupled differential has the form
\[
                         D=\partial+V,                      \tag{16.26}
\]
where \(V\) includes gauge action on matter, Yukawa and Higgs
linearization where appropriate, ghost coupling, and
gravity--matter mixing.  Even if \(D^2=0\), the primitive
\(h_{\cal T}a\) need not solve the \(D\)-equation.  The next upstream
step is a homological perturbation theorem with an explicit smallness
condition on \(h_{\cal T}V\) or a finite exact inversion on the
retained block library.

\begin{assessment}[Status after V16]
\label{ass:v12v16-status}
Sectorwise block contractions now assemble into an exact contracted
tensor reference complex.  The total matching projection is the
tensor product of the sector projections, the homotopy norm is
bounded by the sum of sector norms, and quasi-locality and verified
symmetry covariance are preserved.  This supplies the correct
starting point for a controlled coupled differential perturbation.
\end{assessment}

\begin{firewall}[Claim boundary after V16]
\label{fire:v12v16-claim}
No physical coupled differential \(D\), nilpotence certificate,
interaction perturbation bound, or coupled matching differential has
yet been constructed.  The tensor theorem applies to the reference
complex and does not prove anomaly cancellation or the interacting
SM--Einstein continuum.
\end{firewall}

\subsection{Parent record}

This V16 note appends to the validated V15 file with record
\[
\begin{array}{c|c}
\text{lines}&4654\\
\text{bytes}&162942\\
\text{SHA--256}&
\texttt{CCDC54E1AD7C9578FFBA54C93523200D1CD85256A1D22AC2DCC29507A65B13D3}.
\end{array}
\]

% =============================================================================
% END APPENDED TWELFTH-CYCLE V16 MATERIAL
% =============================================================================
% =============================================================================
% BEGIN APPENDED TWELFTH-CYCLE V17 MATERIAL
% =============================================================================

\section{Twelfth-cycle V17: controlled transfer of the coupled Ward
differential}
\label{sec:v12v17}

\subsection{The interaction problem}

V16 contracts the tensor product of sectorwise reference complexes.
The physical finite-regulator differential is
\[
                         D=d+V,                             \tag{17.1}
\]
where \(d\) is the reference differential and \(V\) contains the
coupling between sectors.  Even when \(D^2=0\), the old homotopy is
not a contraction for \(D\).  We now give the exact transferred
differential, inclusion, projection, and homotopy under a quantitative
Neumann condition.

No theorem with these formulas occurs in the retained SM notes; the
archive search found only the V16 statement that this acquisition was
next.

\subsection{Contraction data}

Let \(({\cal C},d)\) be a finite Hilbert cochain complex and
\(({\cal R},d_{\cal R})\) a retained complex.  Let
\[
 i:{\cal R}\to{\cal C},\qquad
 p:{\cal C}\to{\cal R},\qquad
 h:{\cal C}\to{\cal C}[-1]                                 \tag{17.2}
\]
satisfy
\[
 pi=I_{\cal R},\qquad
 di=id_{\cal R},\qquad
 pd=d_{\cal R}p,                                            \tag{17.3}
\]
and
\[
 dh+hd=I_{\cal C}-ip.                                      \tag{17.4}
\]
Assume the normalized side conditions
\[
 hi=0,\qquad ph=0,\qquad h^2=0.                            \tag{17.5}
\]
These conditions can be imposed on the finite block contraction and
are retained here as explicit hypotheses.

Let \(V\) have degree one and suppose
\[
                         (d+V)^2=0.                         \tag{17.6}
\]
Assume that \(I+hV\) and \(I+Vh\) are invertible.  Define
\[
 A:=V(I+hV)^{-1}=(I+Vh)^{-1}V.                             \tag{17.7}
\]
The equality follows by multiplying
\((I+Vh)V=V(I+hV)\) by the two inverses.

\subsection{The finite perturbation identity}

\begin{lemma}[Resolvent and Maurer--Cartan identities]
\label{lem:v12v17-algebra}
With \(P=ip\), the operator \(A\) satisfies
\[
\begin{split}
 A&=V-VhA=V-AhV,\\
 I-hA&=(I+hV)^{-1},\\
 I-Ah&=(I+Vh)^{-1},\\
 dA+Ad+APA&=0.
\end{split}                                                  \tag{17.8}
\]
\end{lemma}

\begin{proof}
The first three identities follow directly from the two expressions
in (17.7) and the resolvent identity
\((I+X)^{-1}=I-X(I+X)^{-1}\).

For the last identity, substitute
\(P=I-dh-hd\) from (17.4), expand the left side, and use both
relations \(A=V-VhA=V-AhV\).  All terms containing \(h\) cancel in
pairs, leaving
\[
                         dV+Vd+V^2,
\]
which is zero by (17.6).
\end{proof}

Define
\[
\begin{split}
 d_{\cal R}'&:=d_{\cal R}+pAi,\\
 i'&:=i-hAi,\\
 p'&:=p-pAh,\\
 h'&:=h-hAh.
\end{split}                                                  \tag{17.9}
\]

\begin{theorem}[Quantitative basic perturbation theorem]
\label{thm:v12v17-transfer}
The data in (17.9) satisfy
\[
 (d_{\cal R}')^2=0,\qquad
 (d+V)i'=i'd_{\cal R}',\qquad
 p'(d+V)=d_{\cal R}'p',                                    \tag{17.10}
\]
and
\[
 p'i'=I_{\cal R},\qquad
 (d+V)h'+h'(d+V)=I_{\cal C}-i'p'.                          \tag{17.11}
\]
They also obey
\[
 h'i'=0,\qquad p'h'=0,\qquad (h')^2=0.                     \tag{17.12}
\]
Thus the perturbed ambient complex contracts exactly onto the
retained complex with transferred differential \(d_{\cal R}'\).
\end{theorem}

\begin{proof}
Use Lemma~\ref{lem:v12v17-algebra}.  For example,
\[
\begin{split}
 (d+V)i'
 &=di+Vi-dhAi-VhAi\\
 &=i(d_{\cal R}+pAi)+h\,dAi.
\end{split}                                                  \tag{17.13}
\]
Apply the last identity in (17.8) to \(i\), use
\(di=id_{\cal R}\), and multiply by \(h\).  This turns the final term
in (17.13) into
\[
 -hAi\,d_{\cal R}-hAi\,pAi,
\]
so (17.13) equals \(i'd_{\cal R}'\).  The proof that
\(p'(d+V)=d_{\cal R}'p'\) is the adjoint-sided algebraic calculation.

Apply \(p\) on the left and \(i\) on the right of
\(dA+Ad+APA=0\).  Using (17.3) gives
\[
 d_{\cal R}(pAi)+(pAi)d_{\cal R}+(pAi)^2=0,                 \tag{17.14}
\]
which is exactly \((d_{\cal R}')^2=0\).

Expanding \(p'i'\), every correction contains one of
\(hi,ph,h^2\), and therefore vanishes by (17.5); hence \(p'i'=I\).
A direct expansion of
\((d+V)h'+h'(d+V)\), using (17.8), (17.4), and (17.5), gives
\(I-i'p'\).  The same side conditions together with the resolvent
forms
\[
 h'=(I+hV)^{-1}h=h(I+Vh)^{-1}                              \tag{17.15}
\]
give (17.12).
\end{proof}

\subsection{Norm bounds}

Assume
\[
 \|h\|\,\|V\|\le\eta<1.                                    \tag{17.16}
\]
Then both Neumann series in (17.7) converge and
\[
 \|A\|\le\frac{\|V\|}{1-\eta},\qquad
 \boxed{\|h'\|\le\frac{\|h\|}{1-\eta}.}                    \tag{17.17}
\]
If \(i\) is an isometry and \(p=i^*\), then
\[
\begin{split}
 \|d_{\cal R}'-d_{\cal R}\|
 &\le\frac{\|V\|}{1-\eta},\\
 \|i'-i\|,\ \|p'-p\|
 &\le\frac{\eta}{1-\eta}.
\end{split}                                                  \tag{17.18}
\]

\begin{proof}
Use the norm-convergent Neumann series and (17.15).  Equations
(17.9) then give (17.18).
\end{proof}

The effective interaction on the retained matching modes is not
\(pVi\) alone.  It is the resummed expression
\[
 \boxed{
 pAi=pV(I+hV)^{-1}i
 =\sum_{n=0}^{\infty}(-1)^n pV(hV)^n i.}                   \tag{17.19}
\]
Every excursion from the retained space into the contracted sector
is therefore included.

\subsection{A quasi-local Banach algebra}

For a block operator \(B=(B_{xy})\) over a finite metric space, set
\[
\begin{split}
 \|B\|_\mu
 :=\max\bigg\{&
 \sup_x\sum_y e^{\mu\operatorname{dist}(x,y)}\|B_{xy}\|,\\
 &\sup_y\sum_x e^{\mu\operatorname{dist}(x,y)}\|B_{xy}\|
 \bigg\}.
\end{split}                                                  \tag{17.20}
\]

\begin{lemma}[Weighted Schur algebra]
\label{lem:v12v17-schur}
For \(\mu\ge0\),
\[
                         \|BC\|_\mu\le\|B\|_\mu\|C\|_\mu.   \tag{17.21}
\]
Moreover,
\[
                         \|B_{xy}\|
 \le\|B\|_\mu e^{-\mu\operatorname{dist}(x,y)}.             \tag{17.22}
\]
\end{lemma}

\begin{proof}
The triangle inequality for the metric gives
\[
 e^{\mu\operatorname{dist}(x,z)}
 \le
 e^{\mu\operatorname{dist}(x,y)}
 e^{\mu\operatorname{dist}(y,z)}.
\]
Insert this inequality in the row and column Schur sums for the
product kernel and take the two suprema.  This proves (17.21).
Each individual nonnegative summand is bounded by its row sum,
giving (17.22).
\end{proof}

\begin{theorem}[Quasi-local coupled contraction]
\label{thm:v12v17-local}
Suppose \(h,V\) have finite \(\mu\)-norm and
\[
 \eta_\mu:=
 \max\{\|hV\|_\mu,\|Vh\|_\mu\}<1.                          \tag{17.23}
\]
Then all Neumann series in (17.7)--(17.15) converge in the weighted
Schur algebra, and
\[
 \|A\|_\mu\le
 \frac{\|V\|_\mu}{1-\eta_\mu},\qquad
 \boxed{
 \|h'\|_\mu\le\frac{\|h\|_\mu}{1-\eta_\mu}.}                \tag{17.24}
\]
In particular, \(h'\) has a volume-uniform exponential kernel with
rate \(\mu\).
\end{theorem}

\begin{proof}
The weighted Schur operators form a Banach algebra by
Lemma~\ref{lem:v12v17-schur}.  Conditions (17.23) make both resolvent
series converge in that algebra.  Apply (17.21) to (17.7) and
(17.15), then use (17.22).
\end{proof}

\subsection{Equivariance and the anomaly obstruction}

If \(d,V,i,p,h\) intertwine a verified compact unitary symmetry, then
every term in the resolvent series is equivariant.  Hence
\(d_{\cal R}',i',p',h'\) are equivariant as well.

Let \(a\) be a \(D\)-closed anomaly cochain.  Equation (17.11) gives
\[
                         a=i'p'a+D h'a.                     \tag{17.25}
\]
Thus
\[
 p'a=0\quad\Longrightarrow\quad
 a=D(h'a),                                                  \tag{17.26}
\]
with a quasi-local primitive under (17.23).  The actual coupled
obstruction is the transferred class \(p'a\) in
\(({\cal R},d_{\cal R}')\), rather than the unperturbed component
\(pa\).

This distinction matters for the anomaly ledger.  Interaction
excursions can change both the retained differential and the
projection of a closed cochain even when the reference tensor
projection vanishes.

\subsection{Choice of the reference differential}

The condition \(\|hV\|_\mu<1\) is not expected to hold if all
marginal gauge--matter or gravity--matter couplings are placed in
\(V\).  The theorem instead gives a precise splitting rule:

\begin{enumerate}
\item include every nonperturbatively large local mixing term in the
      finite block reference differential \(d\);
\item verify \(d^2=0\) and contract that enlarged bounded block
      exactly;
\item reserve \(V\) for the remaining inter-block or irrelevant
      mixing and certify (17.23);
\item compute the finite retained differential (17.19) rather than
      truncating it without an error bound.
\end{enumerate}

If no useful split satisfies (17.23), one must enlarge the reference
blocks or invert the complete finite block complex directly.  The
smallness condition is a diagnostic, not a substitute for the
physical construction.

\begin{assessment}[Status after V17]
\label{ass:v12v17-status}
The interaction step is now algebraically closed under an explicit
condition.  A nilpotent coupled perturbation with
\(\max\{\|hV\|_\mu,\|Vh\|_\mu\}<1\) transfers to an exact retained
differential, exact inclusion and projection, and an exponentially
quasi-local coupled homotopy.  The physical obstruction is the
transferred class \(p'a\).
\end{assessment}

\begin{firewall}[Claim boundary after V17]
\label{fire:v12v17-claim}
No concrete SM--Einstein differential split, nilpotence proof, or
weighted smallness certificate has been supplied.  Marginal
interactions may have to be incorporated into the reference block.
The theorem does not prove that the transferred anomaly class
vanishes or construct the interacting continuum measure.
\end{firewall}

\subsection{Parent record}

This V17 note appends to the validated V16 file with record
\[
\begin{array}{c|c}
\text{lines}&4971\\
\text{bytes}&173522\\
\text{SHA--256}&
\texttt{4DB6B90A054AA6F215C54564CEAA77A783B70C2FCE703C9B5F75F27A812DABE4}.
\end{array}
\]

% =============================================================================
% END APPENDED TWELFTH-CYCLE V17 MATERIAL
% =============================================================================
% =============================================================================
% BEGIN APPENDED TWELFTH-CYCLE V18 MATERIAL
% =============================================================================

\section{Twelfth-cycle V18: a finite generator certificate for the
coupled BRST differential}
\label{sec:v12v18}

\subsection{Relation to the retained BRST results}

Earlier frozen notes contain case-specific BRST constructions: a
finite contour analysis with a later correction of its scalar
sector, and an algebraically contractible frozen gauge quartet.  They
do not supply the generator-and-relation certificate needed for the
coupled regulator differential \(D\) in V17.  This section gives that
certificate and records exactly what must be tested before the
homological perturbation theorem may be applied.

\subsection{The square of an odd derivation}

Let \({\cal F}\) be the free unital graded associative algebra on
homogeneous generators \(z_1,\ldots,z_m\).  Let \({\cal I}\) be the
homogeneous two-sided ideal generated by relations
\(r_1,\ldots,r_\ell\), and set
\[
                         {\cal A}={\cal F}/{\cal I}.         \tag{18.1}
\]
Let \(s:{\cal F}\to{\cal F}\) be a degree-one graded derivation:
\[
 s(ab)=s(a)b+(-1)^{|a|}a\,s(b)                             \tag{18.2}
\]
for homogeneous \(a\).

\begin{lemma}[The square is an even derivation]
\label{lem:v12v18-even}
For homogeneous \(a,b\),
\[
                         s^2(ab)=s^2(a)b+a\,s^2(b).         \tag{18.3}
\]
\end{lemma}

\begin{proof}
Apply (18.2) twice:
\[
\begin{split}
 s^2(ab)
 &=s^2(a)b+(-1)^{|a|+1}s(a)s(b)\\
 &\quad+(-1)^{|a|}s(a)s(b)+a\,s^2(b).
\end{split}
\]
The two middle terms have opposite signs and cancel.
\end{proof}

\begin{theorem}[Finite generator-and-relation nilpotence
certificate]
\label{thm:v12v18-certificate}
Assume
\[
                         s(r_k)\in{\cal I}
 \quad(1\le k\le\ell)                                      \tag{18.4}
\]
and
\[
                         s^2(z_j)\in{\cal I}
 \quad(1\le j\le m).                                       \tag{18.5}
\]
Then \(s\) descends to a degree-one differential on \({\cal A}\) and
\[
                         \boxed{s^2=0\quad\hbox{on }{\cal A}.}
                                                               \tag{18.6}
\]
\end{theorem}

\begin{proof}
Every element of \({\cal I}\) is a sum of terms \(a r_k b\).
Equation (18.2), the ideal property, and (18.4) show that
\(s({\cal I})\subset{\cal I}\), so \(s\) descends to the quotient.
By Lemma~\ref{lem:v12v18-even}, \(s^2\) is an ordinary even
derivation.  Its value on any word is the sum obtained by applying
\(s^2\) to one generator at a time.  Condition (18.5) puts every
such term in \({\cal I}\).  Hence \(s^2({\cal F})\subset{\cal I}\),
which proves (18.6).
\end{proof}

The theorem turns nilpotence into a finite exact audit whenever the
regulated field algebra has finitely many generators and defining
relations.  Checking \(s^2\) on a selected list of composite
observables is insufficient unless that list generates the quotient
algebra and the relation ideal has also been shown invariant.

\subsection{A quantitative defect ledger}

Suppose a submultiplicative norm is fixed on a truncated word algebra,
\(\|ab\|\le\|a\|\|b\|\), and
\[
 \|z_j\|\le M,\qquad \|s^2z_j\|\le\varepsilon.              \tag{18.7}
\]

\begin{proposition}[Propagation of a generator defect]
\label{prop:v12v18-defect}
For a word \(w=z_{i_1}\cdots z_{i_n}\),
\[
                         \|s^2w\|
 \le n\varepsilon M^{n-1}.                                 \tag{18.8}
\]
For a polynomial \(F=\sum_w c_w w\) of word length at most \(N\),
\[
 \|s^2F\|
 \le\varepsilon
 \sum_w |c_w|\,|w|\,M^{|w|-1}.                             \tag{18.9}
\]
\end{proposition}

\begin{proof}
Lemma~\ref{lem:v12v18-even} gives
\[
 s^2w
 =
 \sum_{k=1}^n
 z_{i_1}\cdots z_{i_{k-1}}
 (s^2z_{i_k})
 z_{i_{k+1}}\cdots z_{i_n}.                                \tag{18.10}
\]
The submultiplicative norm proves (18.8).  Apply the triangle
inequality to the polynomial expansion to obtain (18.9).
\end{proof}

This estimate is a defect bound, not a nilpotence proof.
The perturbation theorem in V17 requires exact \(D^2=0\).
A small nonzero value may be recorded for convergence diagnostics,
but it may not be absorbed into the overlap or anomaly remainder.

\subsection{The Lie-algebra BRST generator test}

Let \(\mathfrak g\) be a finite-dimensional Lie algebra with basis
\(T_a\) and
\[
                         [T_b,T_c]=f^a{}_{bc}T_a.           \tag{18.11}
\]
Let \(\rho:\mathfrak g\to\operatorname{Der}({\cal B})\) be a
representation on the regulated field algebra:
\[
                         [\rho(T_a),\rho(T_b)]
 =f^c{}_{ab}\rho(T_c).                                     \tag{18.12}
\]
Adjoin odd ghost generators \(c^a\) and define
\[
 sc^a=-\frac12 f^a{}_{bc}c^bc^c,\qquad
 s\Phi=-c^a\rho(T_a)\Phi                                  \tag{18.13}
\]
on every field generator \(\Phi\), extending \(s\) by (18.2).

\begin{theorem}[BRST nilpotence from Jacobi and representation]
\label{thm:v12v18-Lie}
If the structure constants obey the Jacobi identity and (18.12)
holds exactly, then
\[
                         s^2c^a=0,\qquad s^2\Phi=0.         \tag{18.14}
\]
Subject to preservation of the regulator relations, \(s^2=0\) on
the full regulated field algebra.
\end{theorem}

\begin{proof}
Substitution of the first formula in (18.13) into \(s^2c^a\) gives a
cubic ghost polynomial whose coefficient is the antisymmetrized
combination
\[
                         f^a{}_{e[b}f^e{}_{cd]},            \tag{18.15}
\]
which vanishes by Jacobi.

For a field generator, the odd Leibniz rule gives
\[
\begin{split}
 s^2\Phi
 &=
 \frac12 f^a{}_{bc}c^bc^c\rho(T_a)\Phi
 -c^ac^b\rho(T_a)\rho(T_b)\Phi\\
 &=
 \frac12 c^ac^b
 \left(
 f^c{}_{ab}\rho(T_c)-[\rho(T_a),\rho(T_b)]
 \right)\Phi=0                                             \tag{18.16}
\end{split}
\]
by (18.12).  Apply
Theorem~\ref{thm:v12v18-certificate} to conclude on the quotient.
\end{proof}

\subsection{Semidirect gauge--diffeomorphism structure}

At the continuum algebraic level, infinitesimal diffeomorphisms and
gauge transformations form a semidirect Lie algebra.  In schematic
notation,
\[
 [(\xi,\alpha),(\eta,\beta)]
 =
 \left(
 [\xi,\eta],
 {\cal L}_\xi\beta-{\cal L}_\eta\alpha+[\alpha,\beta]
 \right).                                                   \tag{18.17}
\]
Jacobi follows when Lie differentiation acts by derivations of the
gauge bracket.  Theorem~\ref{thm:v12v18-Lie} then treats the combined
ghost as one Lie-algebra ghost.

A finite regulator cannot inherit this conclusion by notation alone.
It must supply exact regulated versions of:
\[
\begin{array}{ll}
\text{(N1)}&\text{the Jacobi identity for the regulated bracket},\\
\text{(N2)}&\text{the representation identity (18.12) on every
field generator},\\
\text{(N3)}&\text{invariance of chirality, boundary, reality, and
incidence relations},\\
\text{(N4)}&\text{compatibility with the exact re-tree action}.
\end{array}                                                  \tag{18.18}
\]
These are finite polynomial or matrix identities on each block type.
Once they hold on the generators and relation set,
Theorem~\ref{thm:v12v18-certificate} proves exact nilpotence on all
regulated composite fields.

\subsection{Consequences for the V17 split}

Write the coupled differential as \(D=d+V\).  Its nilpotence expands
as
\[
                         dV+Vd+V^2=0.                       \tag{18.19}
\]
The V17 perturbation formulas may be used only after (18.19) has
been certified exactly.  The finite generator theorem provides the
shortest route: define \(D\) on the block generators, check the
relations, and check \(D^2\) on those generators.  There is no need
to enumerate every composite cochain.

The order of estimates is therefore:
\[
 \boxed{
 \text{exact nilpotence first}
 \ \longrightarrow\
 \text{weighted contraction bound second}
 \ \longrightarrow\
 \text{transferred anomaly class third}.}                  \tag{18.20}
\]
A compiler residual in the first arrow is a consistency defect and
cannot be reclassified as an anomaly coefficient.

\begin{assessment}[Status after V18]
\label{ass:v12v18-status}
Exact nilpotence of the finite coupled differential has been reduced
to a finite generator-and-relation certificate.  For a regulated Lie
action, Jacobi and the exact representation identity prove the
generator test.  This supplies the missing logical gate before the
V17 homological perturbation and its anomaly transfer can be used.
\end{assessment}

\begin{firewall}[Claim boundary after V18]
\label{fire:v12v18-claim}
No regulator realizing the full gauge--diffeomorphism semidirect
algebra has been constructed, and conditions (N1)--(N4) have not been
verified for a physical SM--Einstein block.  Exact lattice
diffeomorphism covariance, chiral determinant phases, the weighted
perturbation bound, and the continuum measure remain open.
\end{firewall}

\subsection{Parent record}

This V18 note appends to the validated V17 file with record
\[
\begin{array}{c|c}
\text{lines}&5314\\
\text{bytes}&183970\\
\text{SHA--256}&
\texttt{49FD4B533EA2F18291263668576CB342F9366D92507A3874E639D3548980ABA5}.
\end{array}
\]

% =============================================================================
% END APPENDED TWELFTH-CYCLE V18 MATERIAL
% =============================================================================
% =============================================================================
% BEGIN APPENDED TWELFTH-CYCLE V19 MATERIAL
% =============================================================================

\section{Twelfth-cycle V19: the exact retained anomaly criterion and
its finite Hodge test}
\label{sec:v12v19}

\subsection{Refinement of the V17 sufficient condition}

V17 observed that a \(D\)-closed anomaly cochain \(a\) is exact when
\(p'a=0\).  That condition is sufficient but stronger than necessary.
The correct obstruction is the cohomology class of \(p'a\) in the
transferred retained complex.  This section proves the equivalence
and reduces it to a finite harmonic projection.

Retain all notation and hypotheses of
Theorem~\ref{thm:v12v17-transfer}.  Thus
\[
 D=d+V,\qquad
 D h'+h'D=I-i'p',\qquad
 Di'=i'd_{\cal R}',\qquad
 p'D=d_{\cal R}'p'.                                       \tag{19.1}
\]

\subsection{Cohomological equivalence}

\begin{theorem}[The ambient anomaly class is exactly the retained
class]
\label{thm:v12v19-equivalence}
Let \(a\in{\cal C}^q\) satisfy \(Da=0\).  Then
\[
                         \omega:=p'a                       \tag{19.2}
\]
is \(d_{\cal R}'\)-closed, and
\[
 \boxed{
 [a]=0\text{ in }H^q({\cal C},D)
 \quad\Longleftrightarrow\quad
 [\omega]=0\text{ in }H^q({\cal R},d_{\cal R}').}           \tag{19.3}
\]
More precisely:
\[
\begin{array}{ll}
a=Db&\Longrightarrow\ \omega=d_{\cal R}'(p'b),\\
\omega=d_{\cal R}'c&
\Longrightarrow\
a=D\bigl(h'a+i'c\bigr).
\end{array}                                                  \tag{19.4}
\]
\end{theorem}

\begin{proof}
The chain-map identity in (19.1) gives
\[
                 d_{\cal R}'\omega=p'Da=0.                 \tag{19.5}
\]
If \(a=Db\), the same identity gives the first line of (19.4).

Conversely, apply the homotopy identity in (19.1) to the closed
cochain:
\[
                         a=i'p'a+D h'a.                     \tag{19.6}
\]
If \(p'a=d_{\cal R}'c\), then
\[
 i'p'a=i'd_{\cal R}'c=D i'c.
\]
Substitution in (19.6) proves the second line of (19.4), and hence
(19.3).
\end{proof}

The theorem corrects the obstruction language: the cochain \(p'a\)
itself need not vanish.  It may be nonzero and still removable if it
is exact for the transferred differential.

\subsection{Finite retained Hodge reduction}

Give the retained complex its finite-regulator Hilbert inner product.
At degree \(q\), define
\[
 \Delta_{{\cal R},q}'
 =
 d_{{\cal R},q-1}'(d_{{\cal R},q-1}')^*
 +(d_{{\cal R},q}')^*d_{{\cal R},q}',                      \tag{19.7}
\]
with harmonic projector \(H_{{\cal R},q}'\) and Moore--Penrose inverse
\(G_{{\cal R},q}'\).  Suppose
\[
 \operatorname{spec}(\Delta_{{\cal R},q}')
 \subset\{0\}\cup[\gamma_{\cal R},\infty),
 \qquad \gamma_{\cal R}>0.                                 \tag{19.8}
\]

\begin{theorem}[Finite harmonic anomaly test]
\label{thm:v12v19-Hodge}
For \(Da=0\), put \(\omega=p'a\).  Then
\[
 \boxed{
 [a]=0
 \quad\Longleftrightarrow\quad
 H_{{\cal R},q}'\omega=0.}                                 \tag{19.9}
\]
When the harmonic component vanishes, define
\[
 c=(d_{{\cal R},q-1}')^*G_{{\cal R},q}'\omega              \tag{19.10}
\]
and
\[
                         b=h'a+i'c.                         \tag{19.11}
\]
Then
\[
                         Db=a                              \tag{19.12}
\]
and
\[
 \|b\|
 \le
 \|h'\|\,\|a\|
 +\|i'\|\frac{\|d_{{\cal R},q-1}'\|}{\gamma_{\cal R}}
            \|p'a\|.                                       \tag{19.13}
\]
\end{theorem}

\begin{proof}
Equation (19.5) says that \(\omega\) is closed.  The finite Hodge
decomposition of a closed cochain has only an exact and a harmonic
part.  Therefore \(\omega\) is exact if and only if its harmonic
part is zero.  Combine this fact with
Theorem~\ref{thm:v12v19-equivalence} to obtain (19.9).

When the harmonic part vanishes, the finite Hodge identity gives
\[
 d_{{\cal R},q-1}'c=\omega.                                \tag{19.14}
\]
Equation (19.12) now follows from the second line of (19.4).  The
gap gives
\(\|G_{{\cal R},q}'\|\le\gamma_{\cal R}^{-1}\), and the
triangle inequality gives (19.13).
\end{proof}

\begin{corollary}[Bound in the perturbative contraction cone]
\label{cor:v12v19-bound}
Assume the unperturbed inclusion is isometric, \(p=i^*\), and
\(\|h\|\|V\|\le\eta<1\).  Then the primitive in (19.11) obeys
\[
 \|b\|
 \le
 \left[
 \frac{\|h\|}{1-\eta}
 +
 \frac{\|d_{{\cal R},q-1}'\|}
      {\gamma_{\cal R}(1-\eta)^2}
 \right]\|a\|.                                              \tag{19.15}
\]
\end{corollary}

\begin{proof}
V17 gives
\[
 \|h'\|\le\frac{\|h\|}{1-\eta},\qquad
 \|i'\|,\|p'\|\le\frac1{1-\eta}.
\]
Insert these estimates and
\(\|p'a\|\le\|p'\|\|a\|\) in (19.13).
\end{proof}

\subsection{Quasi-locality and the retained sector}

Under the weighted hypothesis (17.23), \(h'\) is exponentially
quasi-local.  If the original inclusion \(i\) is local and
\(h,A\) belong to the same weighted Schur algebra, then
\[
                         i'=i-hAi                           \tag{19.16}
\]
is also quasi-local.  Consequently the first term in (19.11) and the
lifting of a block-local retained primitive are quasi-local.

There is an unavoidable qualification.  If the retained complex
contains global topology or boundary matching data, its Hodge inverse
may be global.  Equation (19.11) then separates the two effects:
\[
\begin{array}{c|c}
h'a&\text{contracted bulk primitive, controlled quasi-locally},\\
i'c&\text{retained matching primitive, finite but potentially global}.
\end{array}                                                   \tag{19.17}
\]
A full local-counterterm theorem therefore requires the harmonic-free
retained component either to decompose into bounded matching blocks or
to have its own uniform quasi-local Hodge estimate.

\subsection{A finite executable acceptance test}

For every coupled block type and regulator step, the anomaly test is
now finite:

\begin{enumerate}
\item construct the exact generator differential \(D\) and certify
      \(D^2=0\) by V18;
\item construct \(h',i',p'\), and \(d_{\cal R}'\) by V17;
\item verify
      \((d_{\cal R}')^2=0\) and
      \(d_{\cal R}'p'a=0\);
\item compute an orthonormal basis of
      \(\ker\Delta_{\cal R}'\);
\item evaluate the finite vector
\[
                         \Omega_{\cal R}:=
 H_{\cal R}'p'a;                                            \tag{19.18}
\]
\item if \(\Omega_{\cal R}=0\), verify (19.10)--(19.12)
      directly; otherwise retain \(\Omega_{\cal R}\) as the
      anomaly obstruction.
\end{enumerate}

Exact rational or algebraic arithmetic is preferred when the
regulator entries permit it.  An interval singular-value certificate
may be used when transcendental weights occur, but a floating-point
near-zero is not a proof of (19.18).

\subsection{Refinement transport}

At regulator level \(j\), define
\[
                         \Omega_j
 =H_{{\cal R}_j}'p_j'a_j.                                  \tag{19.19}
\]
The cochain-homotopy transport theorems of V8 apply to the retained
complexes.  If adjacent retained maps are exact cochain-homotopy
equivalences, the harmonic obstruction is transported by their
induced cohomology isomorphism.  Approximate maps add the explicit
consistency and homotopy defects already bounded there.

Summability of
\(\Omega_{j+1}-T_j\Omega_j\) proves convergence of a transported
obstruction.  It does not prove its vanishing.  The continuum anomaly
endpoint requires
\[
                         \lim_j\Omega_j=0                   \tag{19.20}
\]
in the identified finite cohomology space, or exact zero at every
level.

\begin{assessment}[Status after V19]
\label{ass:v12v19-status}
The coupled anomaly question has been reduced exactly to a finite
retained harmonic vector.  A closed ambient anomaly is removable if
and only if \(H_{\cal R}'p'a=0\); when it vanishes, (19.11) gives an
explicit primitive with norm and locality bookkeeping.  This replaces
the overly strong sufficient condition \(p'a=0\) by the correct
cohomological criterion.
\end{assessment}

\begin{firewall}[Claim boundary after V19]
\label{fire:v12v19-claim}
No physical retained complex, transferred differential, harmonic
basis, or vector \(\Omega_{\cal R}\) has yet been computed for the
full SM--Einstein regulator.  The theorem reduces the obstruction but
does not prove that it vanishes, establish locality of global matching
modes, or construct the continuum measure.
\end{firewall}

\subsection{Parent record}

This V19 note appends to the validated V18 file with record
\[
\begin{array}{c|c}
\text{lines}&5596\\
\text{bytes}&193426\\
\text{SHA--256}&
\texttt{6E8CC327B4550C8D4C38005403120376F0C42280EBABB01F211261837F4FB4C1}.
\end{array}
\]

% =============================================================================
% END APPENDED TWELFTH-CYCLE V19 MATERIAL
% =============================================================================
% =============================================================================
% BEGIN APPENDED TWELFTH-CYCLE V20 MATERIAL
% =============================================================================

\section{Twelfth-cycle V20: compulsory companion audit by content
fingerprint}
\label{sec:v12v20}

\subsection{Audit record}

This is the compulsory audit five versions after V15.  The companion
folder currently contains a nominal YM V47 file, a YM V46 file, and
NS V26.  Exact content inspection gives
\[
\begin{array}{c|r|r|c}
\text{record}&\text{lines}&\text{bytes}&\text{SHA--256}\\ \hline
\text{YM nominal V47}&18335&627990&
\texttt{8D00DD00E4F38552262670E57FE9EA3810EFF75F3A2A7DC115BCCDBCF3A97B39}\\
\text{YM V46}&18335&627990&
\texttt{8D00DD00E4F38552262670E57FE9EA3810EFF75F3A2A7DC115BCCDBCF3A97B39}\\
\text{NS V26}&5319&278283&
\texttt{82C887F8C2C69E4F28FAD7C3DC8DE5B9099CB5A4BF431EBA1FFF3E91935978AD}.
\end{array}                                                   \tag{20.1}
\]
The two YM files are byte-for-byte identical.  The nominal V47 file
contains no V47 section or V47 section label and ends with the V46
parent record.  It is therefore an unchanged snapshot of V46, not a
new mathematical version.  No companion file is modified or removed
by this audit.

\subsection{Mathematical consequence}

The substantive companion conclusions remain those recorded in V15:

\begin{enumerate}
\item exact conditional product-Haar invariance removes the compact
      vertical gauge-coordinate Jacobian from the nonlinear coupled
      measure ledger;
\item the nonzero YM corner residual is localized to a shared
      upstream second-jet compiler grammar;
\item NS V26 supplies no new estimate beyond the rank-one Oseen
      kernel analysis already audited.
\end{enumerate}

Consequently there is no evidence requiring a change to V16--V19.
The current SM--Einstein direction remains:

\[
\begin{split}
&\text{sector contractions}
 \longrightarrow
 \text{coupled nilpotent differential}
 \longrightarrow
 \text{transferred retained complex}\\
&\hspace{35mm}\longrightarrow
 \text{finite harmonic anomaly vector}
 \longrightarrow
 \text{continuum Ward-compatible measure}.
\end{split}                                                   \tag{20.2}
\]

\begin{proposition}[Version names do not supersede content
certificates]
\label{prop:v12v20-fingerprint}
Let two purported append-only versions have identical byte length and
SHA--256 digest, and suppose their bytes have been directly compared
or the digest is used under the ordinary collision-resistance
assumption.  Then every theorem, label, parent record, and claim
boundary in the two files is identical.  The later filename supplies
no additional mathematical premise.
\end{proposition}

\begin{proof}
Under direct byte equality the statement is immediate.  Under the
stated digest convention, the content records are identified by the
programme's version protocol.  A filename is external metadata and
does not change the encoded text.
\end{proof}

The proposition is used only for research-note bookkeeping.  It is
not a mathematical assumption about the physical regulator.

\subsection{Adjusted next acquisition}

V19 reduces the coupled anomaly to the finite harmonic vector
\[
                         \Omega_{\cal R}=H_{\cal R}'p'a.     \tag{20.3}
\]
Analytic refinement estimates can make this vector small, but V8
correctly warned that convergence alone does not make it zero.  The
next note supplies the missing exactness mechanism under a discrete
cohomology-lattice hypothesis: a nonzero lattice class has a positive
systolic norm, so a certified bound below that norm forces exact
vanishing.

\begin{assessment}[Status after V20]
\label{ass:v12v20-status}
The mandatory audit is complete and yields no new companion input.
The apparent YM version increase is an identical V46 snapshot.
The retained-anomaly programme therefore continues unchanged, with
the next step aimed at converting quantitative smallness into exact
zero when the obstruction is quantized.
\end{assessment}

\begin{firewall}[Claim boundary after V20]
\label{fire:v12v20-claim}
No quantization lattice for the full regulated SM--Einstein anomaly
has yet been constructed.  The audit does not repair the YM compiler,
advance the NS regularity programme, prove anomaly cancellation, or
construct the continuum measure.
\end{firewall}

\subsection{Parent record}

This V20 note appends to the validated V19 file with record
\[
\begin{array}{c|c}
\text{lines}&5874\\
\text{bytes}&202353\\
\text{SHA--256}&
\texttt{AC3DDA31428E7807E4A1352810350ACAC1B796AF7C1632AA5FD5594A262BB847}.
\end{array}
\]

% =============================================================================
% END APPENDED TWELFTH-CYCLE V20 MATERIAL
% =============================================================================
% =============================================================================
% BEGIN APPENDED TWELFTH-CYCLE V21 MATERIAL
% =============================================================================

\section{Twelfth-cycle V21: a cohomology-lattice gap converting
smallness into exact anomaly cancellation}
\label{sec:v12v21}

\subsection{What is new relative to the frozen notes}

The frozen third-volume notes already warn that a real Hodge
resolution need not preserve an integral anomaly lattice, and the
first large volume distinguishes Gauss--Manin lattice transport from
metric Kato transport.  Those warnings are retained.  The new result
here is the quantitative consequence needed after V19: once a
specific lattice has been constructed and preserved, its shortest
nonzero vector gives an exact zero test for the retained harmonic
anomaly.

\subsection{The lattice systole}

Let \(H\) be a finite-dimensional real Hilbert space and let
\(\Lambda\subset H\) be a lattice, meaning a discrete subgroup of its
real linear span.  Define
\[
 \operatorname{sys}(\Lambda)
 :=
 \inf\{\|\lambda\|:\lambda\in\Lambda,\ \lambda\ne0\}.        \tag{21.1}
\]

\begin{lemma}[Positive lattice systole]
\label{lem:v12v21-systole}
If \(\Lambda\ne\{0\}\), then
\[
                         \operatorname{sys}(\Lambda)>0.     \tag{21.2}
\]
\end{lemma}

\begin{proof}
Discreteness gives a neighbourhood of zero containing no other
lattice point.  Hence there is \(r>0\) such that
\(\Lambda\cap B_r(0)=\{0\}\).  Every nonzero lattice vector has norm
at least \(r\), proving (21.2).
\end{proof}

\begin{theorem}[Sub-systolic anomaly vanishing]
\label{thm:v12v21-zero}
Let the retained harmonic anomaly vector from V19 satisfy
\[
                         \Omega_{\cal R}\in\Lambda.          \tag{21.3}
\]
If
\[
                         \|\Omega_{\cal R}\|
 <\operatorname{sys}(\Lambda),                             \tag{21.4}
\]
then
\[
                         \boxed{\Omega_{\cal R}=0.}         \tag{21.5}
\]
Consequently the full anomaly cochain is exact and the primitive
(19.11) applies.
\end{theorem}

\begin{proof}
A nonzero element of \(\Lambda\) has norm at least its systole, which
contradicts (21.4).  The final statement follows from
Theorem~\ref{thm:v12v19-Hodge}.
\end{proof}

\begin{corollary}[Certified numerical enclosure]
\label{cor:v12v21-interval}
Suppose \(\widetilde\Omega\) is a computed vector and a rigorous error
bound gives
\[
 \|\Omega_{\cal R}-\widetilde\Omega\|\le\epsilon.           \tag{21.6}
\]
If \(\Omega_{\cal R}\in\Lambda\) and
\[
 \|\widetilde\Omega\|+\epsilon
 <\operatorname{sys}(\Lambda),                             \tag{21.7}
\]
then \(\Omega_{\cal R}=0\).
\end{corollary}

\begin{proof}
The triangle inequality gives
\(\|\Omega_{\cal R}\|<\operatorname{sys}(\Lambda)\).
Apply Theorem~\ref{thm:v12v21-zero}.
\end{proof}

A floating-point value close to zero is therefore useful only after
both the enclosure (21.6) and the lattice membership (21.3) have
been proved.

\subsection{Insertion of the perturbative norm bound}

Under the V17 contraction cone, the harmonic projector has norm at
most one and
\[
 \|\Omega_{\cal R}\|
 =\|H_{\cal R}'p'a\|
 \le\|p'\|\,\|a\|
 \le\frac{\|a\|}{1-\eta}.                                  \tag{21.8}
\]

\begin{corollary}[A sufficient anomaly remainder threshold]
\label{cor:v12v21-threshold}
Assume \(\Omega_{\cal R}\in\Lambda\) and
\(\|h\|\|V\|\le\eta<1\).  If
\[
 \boxed{
 \|a\|<(1-\eta)\operatorname{sys}(\Lambda),}               \tag{21.9}
\]
then the retained harmonic anomaly vanishes exactly and \(a\) has
the primitive (19.11).
\end{corollary}

\begin{proof}
Equation (21.8) and (21.9) imply (21.4).  Apply
Theorem~\ref{thm:v12v21-zero}.
\end{proof}

This is an exact implication from an analytic estimate.  It does not
replace the proof that the regulated class lies in the chosen
integral lattice.

\subsection{Lattice-preserving refinement transport}

Let \((H_j,\Lambda_j)\) be the retained harmonic space and anomaly
lattice at refinement level \(j\).  Suppose
\[
 T_j:H_j\longrightarrow H_{j+1},\qquad
                         T_j\Lambda_j=\Lambda_{j+1}         \tag{21.10}
\]
is an exact lattice isomorphism.  Let
\[
                         e_j:=\Omega_{j+1}-T_j\Omega_j.      \tag{21.11}
\]
If both anomaly vectors lie in their lattices, then
\(e_j\in\Lambda_{j+1}\).

\begin{theorem}[Sub-systolic transport defects vanish]
\label{thm:v12v21-transport}
If
\[
                         \|e_j\|
 <\operatorname{sys}(\Lambda_{j+1}),                       \tag{21.12}
\]
then \(e_j=0\).  In particular, if \(\Omega_{j_0}=0\) and (21.12)
holds for every \(j\ge j_0\), then
\[
                         \Omega_j=0
 \quad\hbox{for every }j\ge j_0.                           \tag{21.13}
\]
\end{theorem}

\begin{proof}
The defect is a lattice vector, so (21.12) and
Theorem~\ref{thm:v12v21-zero} give \(e_j=0\).
Equation (21.13) follows by induction from
\(\Omega_{j+1}=T_j\Omega_j\).
\end{proof}

This strengthens the summability conclusion of V8 when quantization
is available.  A sufficiently small quantized adjacent defect is
not merely summable; it is exactly zero.

\begin{corollary}[Uniform identified lattice]
\label{cor:v12v21-uniform}
Suppose all \(H_j\) are identified with one Hilbert space \(H\), all
\(\Lambda_j\) with one lattice \(\Lambda\), and
\(\Omega_j\in\Lambda\).  If \(\Omega_j\to0\), then
\(\Omega_j=0\) for all sufficiently large \(j\).
\end{corollary}

\begin{proof}
Choose \(j_0\) so that
\(\|\Omega_j\|<\operatorname{sys}(\Lambda)\) for
\(j\ge j_0\), and apply Theorem~\ref{thm:v12v21-zero}.
\end{proof}

\subsection{Metric transport versus lattice transport}

The norm in (21.1) may be induced by a regulator-dependent Hodge
metric, while the lattice is transported topologically.  Let
\(G_j:\Lambda\to\Lambda_j\) be a lattice-preserving Gauss--Manin
identifier.  Pull the Hodge norm back to the fixed space and suppose
\[
                         c\|v\|_0\le\|G_jv\|_j
 \le C\|v\|_0                                               \tag{21.14}
\]
with \(c>0\) uniform.  Then
\[
                         \operatorname{sys}_j(\Lambda_j)
 \ge c\,\operatorname{sys}_0(\Lambda).                     \tag{21.15}
\]
Thus the exact zero threshold remains uniform even when the
orthonormal Kato frame rotates relative to the integral lattice.

\begin{proof}
For nonzero \(\lambda_j=G_j\lambda\),
(21.14) gives
\(\|\lambda_j\|_j\ge c\|\lambda\|_0\).
Take the infimum over nonzero lattice vectors.
\end{proof}

Using a Kato-orthonormal frame as though it preserved integer
coordinates would invalidate (21.3).  The lattice membership test
must be performed in a topological lattice frame, while the norm may
be evaluated through the pulled-back Hodge metric.

\subsection{Torsion is a separate finite test}

A real harmonic space detects the free part of cohomology and does
not detect torsion.  Let \(A_{\rm tor}\) be a finite Abelian anomaly
group and \(\widehat A_{\rm tor}\) its character group.

\begin{proposition}[Finite-character torsion test]
\label{prop:v12v21-torsion}
For \(x\in A_{\rm tor}\),
\[
 x=0
 \quad\Longleftrightarrow\quad
 \chi(x)=1\quad\hbox{for every }\chi\in\widehat A_{\rm tor}. \tag{21.16}
\]
\end{proposition}

\begin{proof}
Characters of a finite Abelian group separate points.  Equivalently,
write the group as a product of finite cyclic groups and choose the
standard character on any factor in which \(x\) is nonzero.
\end{proof}

For the familiar \(SU(2)\) mod-two obstruction, the archived
one-generation arithmetic and V9 spectral-flow statement provide the
candidate zero character.  A regulator proof must still identify its
finite spectral-flow class with the same torsion element and preserve
the endpoint gap along refinement.

\subsection{The remaining quantization acquisition}

To use Theorem~\ref{thm:v12v21-zero} for the physical regulator one
must construct:

\begin{enumerate}
\item an integral or differential-cohomological retained anomaly
      lattice \(\Lambda_j\);
\item exact membership of \(H_{\cal R}'p'a\) in that lattice;
\item lattice-preserving adjacent transport;
\item a uniform lower bound for the pulled-back lattice systole;
\item separate finite characters for every torsion component.
\end{enumerate}

The representation-level Standard Model cancellation supplies the
expected zero class.  It does not by itself establish these
regulator-level identifications.

\begin{assessment}[Status after V21]
\label{ass:v12v21-status}
A precise mechanism now converts analytic smallness into exact
vanishing: a retained anomaly known to be lattice-valued is zero
below the lattice systole.  The same argument makes sub-systolic
lattice-valued refinement defects vanish identically.  Free and
torsion anomaly sectors are kept separate.
\end{assessment}

\begin{firewall}[Claim boundary after V21]
\label{fire:v12v21-claim}
No integral anomaly lattice, uniform systole, lattice membership, or
torsion-character comparison has yet been constructed for the full
finite SM--Einstein regulator.  The theorem is conditional and does
not prove the retained anomaly zero or construct the continuum
measure.
\end{firewall}

\subsection{Parent record}

This V21 note appends to the validated V20 file with record
\[
\begin{array}{c|c}
\text{lines}&6001\\
\text{bytes}&207203\\
\text{SHA--256}&
\texttt{2532DA348F13C300601FF25174F6332C2667857377F95E05859FAF6556F0C947}.
\end{array}
\]

% =============================================================================
% END APPENDED TWELFTH-CYCLE V21 MATERIAL
% =============================================================================
% =============================================================================
% BEGIN APPENDED TWELFTH-CYCLE V22 MATERIAL
% =============================================================================

\section{Twelfth-cycle V22: Ward identities for the normalized
complex shell martingale}
\label{sec:v12v22}

\subsection{The remaining gap after the archived Ward-limit theorem}

The frozen fifth-volume notes prove equivariant one-shell integration
and a general closedness criterion for Ward surfaces.  The current
V2 proves a different analytic fact: conditionally normalized complex
shell factors form an \(L^2\)-bounded martingale and hence a
countably additive complex continuum measure.  The missing bridge is
to show that exact or vanishing finite Ward defects pass through that
specific complex martingale limit.  This section supplies the bridge.

\subsection{Filtered Ward data}

Let \((\Omega,{\cal F},\mu)\) be a probability space with increasing
filtration \(({\cal F}_N)\).  Let \({\cal A}_{\rm cyl}\) be a graded
algebra of bounded cylinder observables, with a degree-one derivation
\(s\).  All statements may be read coefficientwise after adjoining a
fixed finite Grassmann algebra.

Assume the reference expectation has the integration-by-parts
identity
\[
                         {\mathbb E}_\mu[sG]=0              \tag{22.1}
\]
whenever \(G\), \(sG\), and the products used below belong to the
declared cylinder domain.

Let \(R_j\) be even, complex, \({\cal F}_j\)-measurable shell factors
such that
\[
 {\mathbb E}_\mu[R_j\mid{\cal F}_{j-1}]=1.                  \tag{22.2}
\]
Put
\[
                         W_N=\prod_{j=1}^N R_j,\qquad W_0=1.
                                                               \tag{22.3}
\]
Assume the square-summable conditional variance hypotheses of V2, so
\[
 W_N\longrightarrow W_\infty
 \quad\hbox{in }L^2(\mu)\hbox{ and almost surely}.           \tag{22.4}
\]

\subsection{Exact finite Ward invariance}

\begin{theorem}[Ward-compatible martingale continuum]
\label{thm:v12v22-Ward}
Assume
\[
                         sR_j=0\quad\hbox{for every }j.      \tag{22.5}
\]
Then \(sW_N=0\) for every \(N\).  The formula
\[
                         \nu(A):=
 {\mathbb E}_\mu[W_\infty{\bf1}_A]                          \tag{22.6}
\]
defines a normalized countably additive complex measure with
\[
                         \|\nu\|_{\rm TV}
 \le{\mathbb E}_\mu|W_\infty|
 \le\|W_\infty\|_2.                                        \tag{22.7}
\]
For every cylinder observable \(F\) with \(sF\in L^2(\mu)\),
\[
                         \boxed{\int_\Omega sF\,d\nu=0.}    \tag{22.8}
\]
Moreover, for \(F\in{\cal F}_m\),
\[
 \int F\,d\nu
 =
 {\mathbb E}_\mu[W_mF],                                    \tag{22.9}
\]
so the continuum Ward functional is projectively consistent with all
finite shell levels.
\end{theorem}

\begin{proof}
Since the factors are even, the derivation rule and (22.5) give
\(sW_N=0\).  Apply (22.1) to \(W_NF\):
\[
 0={\mathbb E}_\mu[s(W_NF)]
  ={\mathbb E}_\mu[W_NsF].                                 \tag{22.10}
\]
By Cauchy--Schwarz and (22.4),
\[
 \left|
 {\mathbb E}_\mu[(W_\infty-W_N)sF]
 \right|
 \le
 \|W_\infty-W_N\|_2\|sF\|_2\longrightarrow0.               \tag{22.11}
\]
Passing to the limit in (22.10) proves (22.8).

Since \(W_\infty\in L^2\subset L^1\), (22.6) is a countably additive
complex measure and (22.7) is the standard density bound.  The
martingale has expectation one at every level; \(L^1\) convergence
therefore gives \(\nu(\Omega)=1\).  For \(N\ge m\), the martingale
property gives
\[
 {\mathbb E}_\mu[W_N\mid{\cal F}_m]=W_m.
\]
Take the \(L^1\) limit and integrate against \(F\) to obtain (22.9).
\end{proof}

This theorem does not use positivity of \(W_N\).  The \(L^2\) bound
simultaneously supplies uniform integrability, countable additivity,
and passage of the Ward pairing.

\subsection{Conditional normalization preserves closedness}

Let a raw shell weight be
\[
 X_j=e^{-v_j},\qquad
 m_j={\mathbb E}_\mu[X_j\mid{\cal F}_{j-1}],\qquad
 R_j=\frac{X_j}{m_j},                                      \tag{22.12}
\]
with \(m_j\ne0\) almost surely.

\begin{proposition}[Equivariant conditional normalizer]
\label{prop:v12v22-normalizer}
Suppose
\[
 s\,{\mathbb E}_\mu[X\mid{\cal F}_{j-1}]
 =
 {\mathbb E}_\mu[sX\mid{\cal F}_{j-1}]                     \tag{22.13}
\]
on the shell domain and \(sX_j=0\).  Then
\[
                         sm_j=0,\qquad sR_j=0.              \tag{22.14}
\]
\end{proposition}

\begin{proof}
Equation (22.13) gives
\(sm_j={\mathbb E}_\mu[sX_j\mid{\cal F}_{j-1}]=0\).
The derivation rule for an invertible even element gives
\(s(m_j^{-1})=-m_j^{-2}sm_j=0\).  Hence
\(sR_j=s(X_jm_j^{-1})=0\).
\end{proof}

The commutation (22.13) is the filtered form of an equivariant shell
split.  It must be verified for the complete gauge-fixed measure and
cycle; it does not follow merely from invariance of the unsplit
action.

\subsection{Exact counterterm repair before normalization}

Suppose the raw shell action has a closed local Ward defect
\(a_j\).  If the V19 retained harmonic test vanishes, let \(b_j\) be
the quasi-local primitive satisfying
\[
                         s b_j=a_j.                         \tag{22.15}
\]
Adding the counterterm with the sign fixed by the Ward convention
produces a repaired raw weight \(\widetilde X_j\) with
\[
                         s\widetilde X_j=0.                 \tag{22.16}
\]
If its conditional normalizer is nonzero and (22.13) holds, the
normalized repaired factor is closed by
Proposition~\ref{prop:v12v22-normalizer}.  Thus anomaly repair must
precede conditional normalization in the proof; otherwise the
variation of the normalizer is an untracked term.

\subsection{Vanishing finite defects}

Exact closure at every cutoff is the clean route, but a quantitative
limit statement is also useful.

\begin{theorem}[Ward limit with an \(L^2\)-vanishing density defect]
\label{thm:v12v22-defect}
Drop (22.5) and assume instead that \(sW_N\in L^2(\mu)\) and
\[
                         \|sW_N\|_2\longrightarrow0.        \tag{22.17}
\]
Then every bounded cylinder \(F\) with \(sF\in L^2(\mu)\) satisfies
\[
                         \int sF\,d\nu=0.                   \tag{22.18}
\]
More precisely,
\[
 \left|
 {\mathbb E}_\mu[W_\infty sF]
 \right|
 \le
 \|W_\infty-W_N\|_2\|sF\|_2
 +\|sW_N\|_2\|F\|_2.                                      \tag{22.19}
\]
\end{theorem}

\begin{proof}
Apply (22.1) and the derivation rule:
\[
 {\mathbb E}_\mu[W_NsF]
 =-{\mathbb E}_\mu[(sW_N)F].                               \tag{22.20}
\]
Subtract the left side from
\({\mathbb E}_\mu[W_\infty sF]\) and use Cauchy--Schwarz twice.
This gives (22.19), whose right side tends to zero by (22.4) and
(22.17).
\end{proof}

A summable anomaly transport budget does not automatically imply
(22.17).  The accumulated density and its Ward derivative must be
estimated in the same \(L^2\) space.

\subsection{A shellwise sufficient defect estimate}

Because the factors are even,
\[
 sW_N
 =
 \sum_{j=1}^N
 \left(\prod_{k<j}R_k\right)
 (sR_j)
 \left(\prod_{k>j}R_k\right).                              \tag{22.21}
\]
A direct triangle inequality can be too expensive.  A sufficient
condition that avoids an unproved cancellation is to exhibit numbers
\(\delta_{j,N}\) such that
\[
 \left\|
 \left(\prod_{k<j}R_k\right)
 (sR_j)
 \left(\prod_{k>j}R_k\right)
 \right\|_2
 \le\delta_{j,N}                                           \tag{22.22}
\]
and
\[
                         \sum_{j=1}^N\delta_{j,N}
 \longrightarrow0.                                        \tag{22.23}
\]
Then (22.17) follows.  Exact repaired factors make every
\(\delta_{j,N}\) zero and avoid this additional analytic problem.

\subsection{Continuum measure implication reached here}

Combining V2, V19, V21, and the present theorem gives the following
conditional chain:

\[
\begin{split}
&\text{finite nilpotent regulator and vanishing retained anomaly}\\
&\quad\Longrightarrow
\text{quasi-local symmetry-restoring shell counterterms}\\
&\quad\Longrightarrow
\text{closed conditionally normalized shell factors}\\
&\quad\Longrightarrow
\text{normalized countably additive complex measure}\\
&\quad\Longrightarrow
\text{Ward identities on the }L^2\text{ cylinder core}.
\end{split}                                                  \tag{22.24}
\]

Each arrow is now proved under explicit hypotheses.  The unresolved
analytic input is to construct one physical regulator satisfying all
of them while also proving the shell variance summability from V2--V4.

\begin{assessment}[Status after V22]
\label{ass:v12v22-status}
The normalized complex martingale now carries the exact finite Ward
identity to its countably additive continuum measure.  Equivariant
conditional expectation preserves closure of each normalized factor,
and an \(L^2\)-vanishing density defect is sufficient when exact
closure is unavailable.  This joins the anomaly-repair complex to the
measure construction without assuming positivity.
\end{assessment}

\begin{firewall}[Claim boundary after V22]
\label{fire:v12v22-claim}
No physical family of repaired shell factors has been shown to obey
the V2 variance summability, the normalizer nonvanishing, or the
equivariant conditional-expectation identity.  The resulting complex
measure is not yet a positive physical Schwinger measure, and
reflection positivity, Einstein constraints, clustering, and the
full continuum remain open.
\end{firewall}

\subsection{Parent record}

This V22 note appends to the validated V21 file with record
\[
\begin{array}{c|c}
\text{lines}&6292\\
\text{bytes}&216830\\
\text{SHA--256}&
\texttt{1AA14C3C6ABF8F5938A81D81A1D68B181B7832506F5AF7588990F51D3CCB8389}.
\end{array}
\]

% =============================================================================
% END APPENDED TWELFTH-CYCLE V22 MATERIAL
% =============================================================================
% =============================================================================
% BEGIN APPENDED TWELF-CYCLE V23 MATERIAL
% =============================================================================

\section{Twelfth-cycle V23: stability of the conditional martingale
under symmetry-restoring counterterms}
\label{sec:v12v23}

\subsection{The normalization problem}

V19 constructs a counterterm primitive when the retained anomaly
vanishes, and V22 uses the repaired weight in a conditionally
normalized shell martingale.  Two analytic facts still had to be
proved: the new conditional normalizer must stay nonzero, and the
new shell variance must remain square summable.  This section proves
both from an explicit small counterterm bound.

\subsection{One repaired shell}

Fix a sigma-algebra \({\cal G}\).  Let \(R\) be a complex random
variable satisfying
\[
 {\mathbb E}[R\mid{\cal G}]=1,\qquad
 \left({\mathbb E}[|R-1|^2\mid{\cal G}]\right)^{1/2}
 \le\varepsilon                                             \tag{23.1}
\]
almost surely.  Let \(C\) be a complex multiplicative counterterm
factor with
\[
                         |C-1|\le\delta                     \tag{23.2}
\]
almost surely, where \(\varepsilon,\delta\ge0\) are deterministic.
Define
\[
                         m:={\mathbb E}[RC\mid{\cal G}].     \tag{23.3}
\]

\begin{theorem}[Nonvanishing and variance after conditional repair]
\label{thm:v12v23-one-shell}
Put
\[
                         \kappa:=\delta\sqrt{1+\varepsilon^2}.
                                                               \tag{23.4}
\]
If \(\kappa<1\), then
\[
                         |m|\ge1-\kappa>0.                  \tag{23.5}
\]
The normalized repaired factor
\[
                         \widetilde R:=\frac{RC}{m}          \tag{23.6}
\]
is well defined,
\[
 {\mathbb E}[\widetilde R\mid{\cal G}]=1,                  \tag{23.7}
\]
and obeys
\[
 \boxed{
 \left(
 {\mathbb E}[|\widetilde R-1|^2\mid{\cal G}]
 \right)^{1/2}
 \le
 \widetilde\varepsilon:=
 \frac{(1+\delta)\varepsilon+\delta}
      {1-\delta\sqrt{1+\varepsilon^2}}.}                   \tag{23.8}
\]
\end{theorem}

\begin{proof}
Since the conditional mean of \(R-1\) is zero,
\[
 {\mathbb E}[|R|^2\mid{\cal G}]
 =
 1+{\mathbb E}[|R-1|^2\mid{\cal G}]
 \le1+\varepsilon^2.                                       \tag{23.9}
\]
Consequently conditional Cauchy--Schwarz gives
\[
\begin{split}
 |m-1|
 &=
 \left|
 {\mathbb E}[R(C-1)\mid{\cal G}]
 \right|\\
 &\le
 \left({\mathbb E}[|R|^2\mid{\cal G}]\right)^{1/2}
 \left({\mathbb E}[|C-1|^2\mid{\cal G}]\right)^{1/2}\\
 &\le\delta\sqrt{1+\varepsilon^2}=\kappa.
\end{split}                                                  \tag{23.10}
\]
The reverse triangle inequality proves (23.5), and (23.7) follows
from the definition.

Conditional expectation is the orthogonal projection in conditional
\(L^2\), so
\[
 \left(
 {\mathbb E}[|RC-m|^2\mid{\cal G}]
 \right)^{1/2}
 \le
 \left(
 {\mathbb E}[|RC-1|^2\mid{\cal G}]
 \right)^{1/2}.                                             \tag{23.11}
\]
Now
\[
                         RC-1=(R-1)C+(C-1),
\]
and \(|C|\le1+\delta\).  Conditional Minkowski therefore gives
\[
 \left(
 {\mathbb E}[|RC-1|^2\mid{\cal G}]
 \right)^{1/2}
 \le(1+\delta)\varepsilon+\delta.                           \tag{23.12}
\]
Finally,
\(\widetilde R-1=(RC-m)/m\).  Combine (23.5),
(23.11), and (23.12) to obtain (23.8).
\end{proof}

The theorem allows \(C\) to depend on both the retained variables and
the current shell.  No independence between \(R\) and \(C\) is used.

\subsection{Square-summable repaired shells}

At shell \(j\), take
\[
 {\cal G}={\cal F}_{j-1},\qquad
 R=R_j,\qquad C=C_j,                                      \tag{23.13}
\]
and denote the original and counterterm bounds by
\(\varepsilon_j,\delta_j\).

\begin{theorem}[Stability of the \(L^2\) martingale criterion]
\label{thm:v12v23-summability}
Assume
\[
 \sum_j\varepsilon_j^2<\infty,\qquad
 \sum_j\delta_j^2<\infty,                                  \tag{23.14}
\]
and
\[
 \sup_j\delta_j\sqrt{1+\varepsilon_j^2}
 \le\kappa_*<1.                                             \tag{23.15}
\]
Then the bounds \(\widetilde\varepsilon_j\) from (23.8) satisfy
\[
                         \sum_j\widetilde\varepsilon_j^2<\infty.
                                                               \tag{23.16}
\]
Hence
\[
                         \widetilde W_N
 :=\prod_{j=1}^N\widetilde R_j                              \tag{23.17}
\]
is an \(L^2\)-bounded complex martingale and converges in \(L^2\) and
almost surely to a normalized limit \(\widetilde W_\infty\).
\end{theorem}

\begin{proof}
Square summability implies
\(\delta_*:=\sup_j\delta_j<\infty\).  From (23.8) and (23.15),
\[
 \widetilde\varepsilon_j
 \le
 \frac{(1+\delta_*)\varepsilon_j+\delta_j}{1-\kappa_*}
 \le C_*(\varepsilon_j+\delta_j),                           \tag{23.18}
\]
where \(C_*\) is independent of \(j\).  Squaring, summing, and using
\((x+y)^2\le2x^2+2y^2\) proves (23.16).  Apply the conditional
martingale theorem of V2 to the normalized factors
\(\widetilde R_j\).
\end{proof}

Only finitely many early shells can threaten (23.15) when
\(\delta_j\to0\).  They may be evaluated directly; the tail is
controlled uniformly by the displayed estimate.

\subsection{Counterterms in exponential form}

Let
\[
                         C_j=e^{-b_j}                       \tag{23.19}
\]
and suppose
\[
                         \|b_j\|_\infty\le\beta_j.           \tag{23.20}
\]
Then
\[
                         |C_j-1|
 \le e^{\beta_j}-1=:\delta_j.                              \tag{23.21}
\]

\begin{corollary}[Square-summable bounded counterterms]
\label{cor:v12v23-exponential}
If
\[
 \sum_j\varepsilon_j^2<\infty,\qquad
 \sum_j\beta_j^2<\infty,\qquad
 \sup_j(e^{\beta_j}-1)\sqrt{1+\varepsilon_j^2}<1,           \tag{23.22}
\]
then the repaired normalized product converges in \(L^2\).
\end{corollary}

\begin{proof}
Let \(\beta_*=\sup_j\beta_j<\infty\).  The mean-value theorem gives
\[
                         e^{\beta_j}-1
 \le e^{\beta_*}\beta_j.
\]
Thus \(\sum_j\delta_j^2<\infty\), and
Theorem~\ref{thm:v12v23-summability} applies.
\end{proof}

The quasi-local norm estimates of V19 do not automatically imply
(23.20).  A regulator-specific embedding from the local counterterm
space into \(L^\infty\), with constants uniform in refinement, is an
additional analytic requirement.

\subsection{Ward compatibility of the repaired factor}

\begin{proposition}[Ward-compatible repaired normalization]
\label{prop:v12v23-Ward}
Suppose the counterterm is chosen so that the unnormalized product is
closed:
\[
                         s(R_jC_j)=0.                       \tag{23.23}
\]
If conditional expectation commutes with \(s\) as in (22.13), then
\[
                         sm_j=0,\qquad s\widetilde R_j=0.   \tag{23.24}
\]
Therefore Theorems~\ref{thm:v12v23-summability} and
\ref{thm:v12v22-Ward} give a normalized countably additive complex
measure satisfying the Ward identities on the \(L^2\) cylinder core.
\end{proposition}

\begin{proof}
Apply Proposition~\ref{prop:v12v22-normalizer} to the raw factor
\(R_jC_j\).  The convergence and Ward conclusions are the cited
theorems.
\end{proof}

\subsection{A combined sufficient shell card}

For each shell it is now sufficient to certify:

\[
\begin{array}{ll}
\text{(C1)}&
 {\mathbb E}[R_j\mid{\cal F}_{j-1}]=1,\quad
 \|R_j-1\|_{2\mid{\cal F}_{j-1}}\le\varepsilon_j,\\
\text{(C2)}&
 s(R_jC_j)=0,\quad C_j=e^{-b_j},\\
\text{(C3)}&
 \|b_j\|_\infty\le\beta_j,\\
\text{(C4)}&
 \sum_j(\varepsilon_j^2+\beta_j^2)<\infty,\\
\text{(C5)}&
 (e^{\beta_j}-1)\sqrt{1+\varepsilon_j^2}<1
 \text{ with a uniform margin},\\
\text{(C6)}&
 s\,{\mathbb E}[X\mid{\cal F}_{j-1}]
 =
 {\mathbb E}[sX\mid{\cal F}_{j-1}]
 \text{ on the repaired shell domain}.
\end{array}                                                   \tag{23.25}
\]

This card links the block-cohomology repair to the analytic shell
measure without assuming that a small counterterm has no effect on
normalization.

\begin{assessment}[Status after V23]
\label{ass:v12v23-status}
Small symmetry-restoring counterterms are now proved compatible with
the conditional \(L^2\) martingale construction.  The new normalizer
has an explicit nonzero lower bound, and the repaired conditional
variance remains square summable whenever the original variance and
counterterm sup norms are square summable.
\end{assessment}

\begin{firewall}[Claim boundary after V23]
\label{fire:v12v23-claim}
No uniform \(L^\infty\) embedding for the physical quasi-local
counterterms, no shellwise bounds \(\varepsilon_j,\beta_j\), and no
equivariant conditional expectation have yet been proved for the
coupled regulator.  The theorem does not establish positivity,
reflection positivity, Einstein constraint concentration, or the
full continuum.
\end{firewall}

\subsection{Parent record}

This V23 note appends to the validated V22 file with record
\[
\begin{array}{c|c}
\text{lines}&6593\\
\text{bytes}&226774\\
\text{SHA--256}&
\texttt{BCFD766E7BAD36902EDEA8A9ED8BB5FFC9EAE4CA50D0DFB3E90499D3ED2A69C1}.
\end{array}
\]

% =============================================================================
% END APPENDED TWELF-CYCLE V23 MATERIAL
% =============================================================================
% =============================================================================
% BEGIN APPENDED TWELFTH-CYCLE V24 MATERIAL
% =============================================================================

\section{Twelfth-cycle V24: the reciprocal-normalizer gate for
simultaneous Ward repair and reflection positivity}
\label{sec:v12v24}

\subsection{Source-guided change of emphasis}

The current grand-tensor source was inspected at its RPESM regulator
and quantum-continuum sections.  Its file record is
\[
\begin{array}{c|r|c}
\text{bytes}&\text{SHA--256}\\ \hline
2804854&
\texttt{FACD058EC8BB210AD8B296C9D7F4D78A229D08D7A2AE378808700655CB8BC0FE}.
\end{array}                                                   \tag{24.1}
\]
It already supplies a named finite positive RPESM sequence with exact
static and dynamic projectivity, finite reflection positivity, and a
uniform refresh gap.  It also gives a conditional
reflection-positive continuum theorem.  Those results are not
reproved here.

The newer available copy of the attached physical-source programme
is V35, with record
\[
\begin{array}{c|r|c}
\text{bytes}&\text{SHA--256}\\ \hline
1410598&
\texttt{AE1FBC0A8125C37465E38E7649E2992725F4579A86F0BB17483B9AE829CFB17C}.
\end{array}                                                   \tag{24.2}
\]
Its terminal status says that the selected physical RPESM
right-energy fibre or complete five-matrix Duhamel card has not been
supplied.  No coefficient from its rational control is transferred
to the physical regulator.

Together these sources expose a compatibility gate for V23.  A
small Ward-restoring counterterm can preserve conditional
nonvanishing and martingale convergence while its conditional
normalizer destroys reflection positivity.  The exact matrix
criterion is proved below.

\subsection{Reflected boundary kernels}

Let \(B=\{1,\ldots,m\}\) be a finite reflected boundary alphabet.
A reflection-positive cylinder has a Hermitian positive semidefinite
boundary kernel
\[
                         K=(K_{ab})_{a,b\in B}\succeq0.      \tag{24.3}
\]
Let \(C=(C_{ab})\) be the boundary kernel of an unnormalized
counterterm insertion.  Let \(m_{ab}\ne0\) be the conditional
normalizer evaluated on the reflected coarse boundary pair and put
\[
                         N_{ab}:=\frac1{m_{ab}}.             \tag{24.4}
\]
Up to one positive global scalar, the repaired boundary kernel is
\[
                         \widetilde K=K\circ C\circ N,       \tag{24.5}
\]
where \(\circ\) denotes entrywise product.

\begin{lemma}[Schur-product positivity]
\label{lem:v12v24-Schur}
If \(A\succeq0\) and \(B\succeq0\), then
\[
                         A\circ B\succeq0.                  \tag{24.6}
\]
\end{lemma}

\begin{proof}
Choose Gram representations
\(A_{ij}=\langle u_i,u_j\rangle\) and
\(B_{ij}=\langle v_i,v_j\rangle\).  Then
\[
 (A\circ B)_{ij}
 =
 \langle u_i\otimes v_i,u_j\otimes v_j\rangle,
\]
so \(A\circ B\) is a Gram matrix.
\end{proof}

\begin{theorem}[Joint Ward-normalization and reflection gate]
\label{thm:v12v24-joint}
Assume
\[
                         K\succeq0,\qquad C\succeq0,\qquad
                         N\succeq0.                         \tag{24.7}
\]
Then the repaired kernel in (24.5) is positive semidefinite.  If its
total mass is positive, division by that positive scalar preserves
reflection positivity.

If in addition the shell conditions of V23 hold and the repaired
unnormalized factor is Ward closed, then the shell simultaneously
has:
\begin{enumerate}
\item a nonzero conditional normalizer;
\item the square-summable conditional variance bound (23.8);
\item exact finite Ward closure; and
\item a nonnegative Osterwalder--Schrader form on the tested boundary
      algebra.
\end{enumerate}
\end{theorem}

\begin{proof}
Apply Lemma~\ref{lem:v12v24-Schur} twice to (24.5).
Multiplication by a positive scalar does not change the sign of a
quadratic form.  The first three additional conclusions are
Theorem~\ref{thm:v12v23-one-shell},
Theorem~\ref{thm:v12v23-summability}, and
Proposition~\ref{prop:v12v23-Ward}; the fourth is the boundary-kernel
criterion for reflection positivity.
\end{proof}

The condition \(C\succeq0\) says that the counterterm itself is a
reflection-positive insertion.  Pointwise positivity of
\(e^{-b}\) is not enough.  The additional condition \(N\succeq0\)
is the new reciprocal-normalizer gate.

\subsection{A factorized sufficient normalizer}

\begin{proposition}[Reflected-square conditional normalizer]
\label{prop:v12v24-factor}
Suppose there are nonzero numbers \(u_a\) such that
\[
                         m_{ab}=\overline{u_a}u_b.           \tag{24.8}
\]
Then
\[
                         N_{ab}
 =\overline{u_a^{-1}}u_b^{-1}                              \tag{24.9}
\]
is a rank-one positive semidefinite kernel.  Consequently every
\(K\succeq0\) and \(C\succeq0\) obeys the gate (24.7).
\end{proposition}

\begin{proof}
Equation (24.9) is the Gram matrix of the scalars \(u_a^{-1}\).
Apply Theorem~\ref{thm:v12v24-joint}.
\end{proof}

A boundary-independent normalizer is the special case in which all
\(u_a\) are equal.  More generally, (24.8) is the exact
past/future reflected-square factorization required of the
conditional partition function.

\subsection{Nonvanishing and smallness are not sufficient}

\begin{counterexample}[An arbitrarily small normalizer variation can
break reflection positivity]
\label{cth:v12v24-small}
For \(0<\epsilon<1\), take
\[
 K=C=
 \begin{pmatrix}1&1\\1&1\end{pmatrix},
 \qquad
 M=(m_{ab})=
 \begin{pmatrix}1&1-\epsilon\\1-\epsilon&1\end{pmatrix}.
                                                               \tag{24.10}
\]
Both \(K\) and \(C\) are positive semidefinite, and \(M\) is positive
definite.  Every entry of \(M\) is nonzero and converges uniformly to
one as \(\epsilon\downarrow0\).  Nevertheless
\[
 \widetilde K=
 \begin{pmatrix}
 1&(1-\epsilon)^{-1}\\
 (1-\epsilon)^{-1}&1
 \end{pmatrix}                                               \tag{24.11}
\]
has eigenvalues
\[
                         1+(1-\epsilon)^{-1}>0,\qquad
                         1-(1-\epsilon)^{-1}<0.              \tag{24.12}
\]
Thus conditional nonvanishing, entrywise positivity, positive
definiteness of the normalizer matrix, and arbitrarily small
entrywise deviation from a constant do not preserve reflection
positivity.
\end{counterexample}

\begin{proof}
The eigenvalues of \(M\) are \(2-\epsilon\) and \(\epsilon\), so
\(M\succ0\).  Entrywise inversion gives (24.11), whose displayed
eigenvalues have opposite signs.
\end{proof}

The obstruction appears for arbitrarily small \(\epsilon\) because
the original reflected kernel has a null direction.  Hence no
ordinary smallness argument protects an unquotiented
Osterwalder--Schrader null space.

\subsection{A strict-gap robustness alternative}

Put
\[
                         A:=K\circ C.                       \tag{24.13}
\]
Let \(J\) be the all-ones matrix, the identity element for entrywise
multiplication.

\begin{proposition}[Robustness on a strictly positive reflected
quotient]
\label{prop:v12v24-gap}
Assume
\[
                         A\succeq\lambda I,\qquad\lambda>0, \tag{24.14}
\]
and write \(N=J+E\).  If
\[
                         \|A\|_{\rm F}
 \max_{a,b}|E_{ab}|<\lambda,                               \tag{24.15}
\]
then
\[
                         A\circ N\succ0.                    \tag{24.16}
\]
\end{proposition}

\begin{proof}
Since \(A\circ J=A\),
\[
 \|A\circ N-A\|_{\rm op}
 \le\|A\circ E\|_{\rm F}
 \le\|A\|_{\rm F}\max_{a,b}|E_{ab}|.                       \tag{24.17}
\]
The least eigenvalue changes by no more than the operator norm of the
perturbation.  Conditions (24.14)--(24.15) therefore give a positive
least eigenvalue.
\end{proof}

For a semidefinite reflected form, one may first quotient its exact
null ideal and apply this proposition only if the counterterm and
normalizer descend to that quotient.  If they move a null vector, the
exact positive-kernel test in (24.7) is required.

\subsection{Application to the named RPESM sequence}

The RPESM source uses positive conditional kernels of the form
\[
 K_{n+1/n}(dy\mid x)
 =
 Z_{n+1/n}(x)^{-1}e^{-V_{n+1/n}(x,y)}
 \kappa_{n+1/n}(dy\mid x).                                 \tag{24.18}
\]
Its original finite reflection statement is part of the named
construction.  If the V19 counterterm is inserted into (24.18), V23
controls the new scalar conditional denominator and its variance,
but reflection positivity survives only after one of the following
is proved:

\begin{enumerate}
\item the counterterm and reciprocal normalizer kernels both pass the
      exact positive-semidefinite tests in (24.7);
\item the normalizer has the reflected-square factorization (24.8);
      or
\item the tested OS quotient has a strict floor and the quantitative
      perturbation bound (24.15) holds.
\end{enumerate}

This is a finite matrix certificate on every declared positive-time
boundary algebra.  It must be transported with the adjacent-cutoff
maps just like the Ward and determinant-line data.

\begin{assessment}[Status after V24]
\label{ass:v12v24-status}
The Ward-repaired martingale and the reflection-positive RPESM branch
now have a precise compatibility gate.  Conditional nonvanishing and
small variance do not preserve reflection positivity; the entrywise
reciprocal normalizer must itself be a positive kernel, factor as a
reflected square, or be controlled on a strictly positive OS
quotient.
\end{assessment}

\begin{firewall}[Claim boundary after V24]
\label{fire:v12v24-claim}
No physical RPESM counterterm kernel or reciprocal normalizer has
been tested against (24.7), and the physical Duhamel fibre remains
unpopulated in the available V35 source.  The theorem does not prove
a scalarized chiral measure, determinant-line compatibility,
continuum tightness, or the full SM--Einstein continuum.
\end{firewall}

\subsection{Parent record}

This V24 note appends to the validated V23 file with record
\[
\begin{array}{c|c}
\text{lines}&6902\\
\text{bytes}&236091\\
\text{SHA--256}&
\texttt{951D9E42B2E1C3B5EBDBAFC0CDF6CF75507C542010E3A7F8688D46DC9F259ADB}.
\end{array}
\]

% =============================================================================
% END APPENDED TWELFTH-CYCLE V24 MATERIAL
% =============================================================================
% =============================================================================
% BEGIN APPENDED TWELFTH-CYCLE V25 MATERIAL
% =============================================================================

\section{Twelfth-cycle V25: compulsory companion audit and correction
from cardwise covariance to full-trace covariance}
\label{sec:v12v25}

\subsection{Audit record}

The current companion files are
\[
\begin{array}{c|r|r|c}
\text{record}&\text{lines}&\text{bytes}&\text{SHA--256}\\ \hline
\text{YM V47}&18675&640193&
\texttt{0C68117F66560B6DD180B9C2BB8C4F0EA84280A5D949D1CB79404DEBB85ABDDB}\\
\text{NS V26}&5319&278283&
\texttt{82C887F8C2C69E4F28FAD7C3DC8DE5B9099CB5A4BF431EBA1FFF3E91935978AD}.
\end{array}                                                   \tag{25.1}
\]
NS V26 is unchanged and supplies no new input.

YM V47 is new and materially changes the interpretation imported in
V15.  Its automatic nilpotent-variable differentiation of the
Wilson plaquette agrees entrywise with the legacy cubic and quartic
writers at both tested corners.  The nonzero cardwise second-jet
difference persists.  The error was the assumption that nonlinear
background-dependent re-tree transport preserves each individual
loop-momentum card.  Its linear background derivatives shift loop
momentum, so determinant covariance applies to the complete momentum
trace, or to a genuinely invariant sum of cards.

\subsection{Formal correction to the present series}

\begin{correction}[Withdrawal of the inherited compiler diagnosis]
\label{corr:v12v25-compiler}
The following earlier statements in this twelfth-cycle series are
withdrawn:

\begin{enumerate}
\item the V15 assertion that the YM residual had been proved to be an
      upstream compiler consistency defect;
\item item 2 in the V20 list, which repeated that diagnosis.
\end{enumerate}

The exact product-Haar disintegration and conditional fibre-Haar
invariance imported in V15 remain valid.  Their valid conclusion is
invariance of the complete measured transformation.  They do not
imply equality of individual momentum-card contributions when the
background-dependent coordinate map mixes those cards.

V24 does not rely on the withdrawn diagnosis.  Its reciprocal
normalizer theorem remains unchanged.
\end{correction}

\subsection{The general trace decomposition}

Let
\[
                         {\cal H}=\bigoplus_{k\in K}{\cal H}_k
                                                               \tag{25.2}
\]
be a finite-dimensional Hilbert space with orthogonal projections
\(P_k\).  Let \(U(z)\) be an analytic family of invertible operators
and let \(X(z)\) be any analytic operator family.  Under similarity
transport,
\[
                         X'(z)=U(z)X(z)U(z)^{-1}.            \tag{25.3}
\]

\begin{theorem}[Full trace invariance and the cardwise criterion]
\label{thm:v12v25-trace}
For every \(z\),
\[
                         \operatorname{Tr}X'(z)
 =\operatorname{Tr}X(z).                                   \tag{25.4}
\]
The contribution of card \(j\) after transport is
\[
 \operatorname{Tr}(P_jX'(z))
 =
 \operatorname{Tr}\!\left(
 U(z)^{-1}P_jU(z)X(z)
 \right).                                                   \tag{25.5}
\]
Consequently a cardwise identity transporting card \(k\) to card
\(j\) is guaranteed for all \(X\) if and only if
\[
                         U(z)^{-1}P_jU(z)=P_k.              \tag{25.6}
\]
Equality in (25.6) only at \(z=0\) does not imply equality of any
positive-order card derivatives at zero.
\end{theorem}

\begin{proof}
Cyclicity of the full finite-dimensional trace gives (25.4) and
(25.5).  If (25.6) holds, substitution proves the cardwise identity.
Conversely, if
\[
 \operatorname{Tr}\!\left[
 (U^{-1}P_jU-P_k)X
 \right]=0
\]
for every \(X\), nondegeneracy of the trace pairing forces (25.6).
An equality at one parameter value imposes no condition on the
derivatives of \(U^{-1}P_jU\), proving the final assertion.
\end{proof}

\begin{corollary}[First projected transport defect]
\label{cor:v12v25-first}
Assume \(U(0)=U_0\) and
\(U_0^{-1}P_jU_0=P_k\).  Put
\[
                         A:=U_0^{-1}\dot U(0).              \tag{25.7}
\]
Then
\[
 \frac{d}{dz}
 \left[U(z)^{-1}P_jU(z)\right]_{z=0}
 =
                         [P_k,A].                          \tag{25.8}
\]
Thus the first cardwise transport identity requires
\([P_k,A]=0\).  If \(A\) shifts momentum, this commutator is generally
nonzero, while the sum over all cards still cancels because
\(\sum_kP_k=I\).
\end{corollary}

\begin{proof}
Differentiate \(U^{-1}P_jU\) and use
\(\dot U^{-1}=-U^{-1}\dot U\,U^{-1}\).  After identifying the
zeroth-order target projector with \(P_k\), the result is (25.8).
The sum of the commutators is \([I,A]=0\).
\end{proof}

For a Hessian congruence rather than a similarity, the complete
log-determinant changes by the determinant of \(U\); exact Haar
transport cancels the corresponding measured coordinate term.
The cardwise obstruction remains the same: inserting \(P_k\) before
taking the trace is licensed only when the transported projector
identity holds through the required background order.

\subsection{Revised projection rule for the SM block programme}

The equivariant projection theorem in V6 and the averaging theorem in
V14 remain correct, but their representation must now be interpreted
on the complete space on which the nonlinear or background-linearized
transport acts.  A projector onto one momentum card, one channel, or
one block label is admissible only after proving that it commutes with
every transport derivative used by the calculation.

Let \({\cal U}\) be the algebra generated by the verified transport
operators and their required background derivatives.  The admissible
local extraction space must be a reducing \({\cal U}\)-module:
\[
                         [P,A]=0\qquad(A\in{\cal U}).        \tag{25.9}
\]
If an individual card fails (25.9), it must be enlarged to its
\({\cal U}\)-invariant envelope before a Ward, anomaly, determinant,
or reflection residual is assigned to that card.

This changes the next finite block certificate in one practical way.
The V13 mismatch \(\delta\) must compare global matching projections
after closing the proposed block space under all background-induced
momentum shifts.  Testing only the zero-background re-tree matrix can
underestimate the matching space and manufacture a false residual.

\subsection{Next acquisition}

For a finite momentum group and a finite set of background shifts,
the minimal invariant envelope can be computed exactly.  V26 will
prove that its blocks are cosets of the subgroup generated by those
shifts, and that convolutional transport is block diagonal precisely
on those cosets.  This gives the smallest legal trace packets for the
SM and YM background-response calculations.

\begin{assessment}[Status after V25]
\label{ass:v12v25-status}
The compulsory audit is complete and corrects the series promptly.
The YM Wilson compiler is validated on the tested cards; the failed
inference was cardwise covariance under a momentum-mixing
background.  The SM block programme now requires reducing spaces for
the full transport algebra, not merely for its zero-background
linearization.
\end{assessment}

\begin{firewall}[Claim boundary after V25]
\label{fire:v12v25-claim}
The complete sixteen-corner YM trace has not yet been certified, and
no physical SM transport algebra or invariant block envelope has
been computed.  The correction does not establish anomaly
cancellation, reflection-positive normalization, continuum
tightness, or the full SM--Einstein continuum.
\end{firewall}

\subsection{Parent record}

This V25 note appends to the validated V24 file with record
\[
\begin{array}{c|c}
\text{lines}&7204\\
\text{bytes}&246450\\
\text{SHA--256}&
\texttt{CD36DC23A105AFE2573AA42E0E8EC0A40C001D7C711A011CE9558BA1D1ADCAE7}.
\end{array}
\]

% =============================================================================
% END APPENDED TWELFTH-CYCLE V25 MATERIAL
% =============================================================================
% =============================================================================
% BEGIN APPENDED TWELFTH-CYCLE V26 MATERIAL
% =============================================================================

\section{Twelfth-cycle V26: momentum-shift cosets as the minimal
legal trace packets}
\label{sec:v12v26}

\subsection{Finite momentum geometry}

Let \(K\) be a finite Abelian momentum group and
\[
                         {\cal H}=\bigoplus_{k\in K}{\cal H}_k,
 \qquad P_k:{\cal H}\to{\cal H}_k.                         \tag{26.1}
\]
Let \(S\subset K\) be the set of momentum shifts appearing in the
background derivatives of the regulator transport.  Include the
negatives and define
\[
                         H_S:=\langle S\rangle\le K.         \tag{26.2}
\]
For a coset \(c=k+H_S\), put
\[
                         Q_c:=\sum_{\ell\in c}P_\ell.        \tag{26.3}
\]

A shift operator of type \(s\in S\) is an operator \(A_s\) satisfying
\[
                         A_s{\cal H}_k\subseteq{\cal H}_{k+s}
 \quad(k\in K).                                             \tag{26.4}
\]
Diagonal momentum operators preserve every \({\cal H}_k\).

\begin{theorem}[Coset reduction of the transport algebra]
\label{thm:v12v26-cosets}
Every projector \(Q_c\) commutes with every diagonal momentum
operator and every shift operator \(A_s\) with \(s\in S\):
\[
                         Q_cA_s=A_sQ_c.                     \tag{26.5}
\]
Consequently the unital star-algebra generated by the diagonal
operators, the shifts \(A_s\), and their adjoints is block diagonal
on
\[
                         {\cal H}=\bigoplus_{c\in K/H_S}Q_c{\cal H}.
                                                               \tag{26.6}
\]
The cosets are the smallest subsets of momentum labels closed under
all translations by \(S\cup(-S)\).  If the shift maps connect every
successive fibre nondegenerately, \(Q_c{\cal H}\) is the smallest
reducing support envelope containing any full card
\({\cal H}_k\) with \(k\in c\).
\end{theorem}

\begin{proof}
If \(k\in c\), then \(k+s\in c\); if \(k\notin c\), then
\(k+s\notin c\), because addition by an element of \(H_S\) does not
change a coset.  Hence
\[
 Q_cA_sP_k=
 \begin{cases}
 A_sP_k,&k\in c,\\
 0,&k\notin c,
 \end{cases}
 =
 A_sQ_cP_k.
\]
Sum over \(k\) to obtain (26.5).  Diagonal operators commute with
every sum of card projectors.  Adjoints shift by \(-s\), so the
star-algebra preserves each coset block.

A subset containing \(k\) and closed under every \(\pm s\) contains
\(k+\sum_jn_js_j\) for all integers \(n_j\), hence contains
\(k+H_S\).  The coset itself is closed, proving minimality at the
support level.  Under the stated nondegeneracy, repeated shifts
generate every full fibre in the coset from the initial full fibre,
so no smaller union of full card fibres is reducing.
\end{proof}

\begin{corollary}[When only the full trace is legal]
\label{cor:v12v26-full}
If \(H_S=K\), the only nonzero support projector formed from complete
momentum cards and reducing every background shift is \(I\).
Therefore symmetry covariance can constrain only the full momentum
trace, unless additional internal invariant subspaces are proved.
\end{corollary}

\begin{proof}
There is one coset in \(K/H_S\).  Apply
Theorem~\ref{thm:v12v26-cosets}.
\end{proof}

\subsection{Cosetwise trace invariance}

Let \(U(z)\) be an invertible transport family belonging to the
coset-reduced algebra in Theorem~\ref{thm:v12v26-cosets}, and let
\[
                         X'(z)=U(z)X(z)U(z)^{-1}.            \tag{26.7}
\]

\begin{theorem}[Exact cosetwise trace packets]
\label{thm:v12v26-trace}
For every coset \(c\),
\[
 \boxed{
 \operatorname{Tr}\!\left(Q_cX'(z)\right)
 =
 \operatorname{Tr}\!\left(Q_cX(z)\right).}                 \tag{26.8}
\]
Thus background-response traces may be assigned to cosets of
\(H_S\).  Assignment to an individual card is licensed only when
that card is itself a coset, equivalently when every occurring shift
fixes its momentum label.
\end{theorem}

\begin{proof}
The reduction theorem gives \(Q_cU=UQ_c\), hence
\[
\begin{split}
 \operatorname{Tr}(Q_cUXU^{-1})
 &=
 \operatorname{Tr}(UQ_cXU^{-1})\\
 &=\operatorname{Tr}(Q_cX)
\end{split}
\]
by cyclicity.  A singleton \(\{k\}\) is a coset precisely when
\(H_S=\{0\}\) on that orbit, which is the stated fixed-label
condition.
\end{proof}

For a Hessian congruence, the same packet statement holds after the
restricted determinant of \(U|_{Q_c{\cal H}}\) and its measured
coordinate density are included.  The Haar argument of V15 controls
that complete measured factor; it does not split it further than the
reducing cosets.

\subsection{Why second derivatives require intermediate cards}

Suppose a background mode has shifts \(+s\) and \(-s\), with
linear transport derivatives \(A_+\) and \(A_-\).  Although the
second-order compositions
\[
                         A_-A_+,\qquad A_+A_-               \tag{26.9}
\]
return \({\cal H}_k\) to itself, their intermediate images lie in
\({\cal H}_{k+s}\) and \({\cal H}_{k-s}\).  Inserting \(P_k\) between
the factors changes the full trace unless those adjacent fibres have
already been included.  The coset projector includes every such
intermediate card and commutes with both shifts.

\begin{proposition}[Closed-walk trace rule]
\label{prop:v12v26-walk}
A nonzero diagonal block of a word
\(A_{s_r}\cdots A_{s_1}\) on \({\cal H}_k\) requires
\[
                         s_1+\cdots+s_r=0\quad\hbox{in }K.  \tag{26.10}
\]
Every intermediate momentum
\[
                         k+s_1+\cdots+s_j                  \tag{26.11}
\]
lies in the same coset \(k+H_S\).  Therefore the coset trace contains
all closed shift walks contributing to every perturbative order.
\end{proposition}

\begin{proof}
Successive use of (26.4) sends \({\cal H}_k\) to
\({\cal H}_{k+s_1+\cdots+s_r}\).  A diagonal block can be nonzero
only if the terminal label equals \(k\), proving (26.10).  Every
partial sum belongs to \(H_S\), proving the second statement.
\end{proof}

This is the algebraic form of the V47 correction: a background
derivative may leave a card and return at second order, so the
intermediate card cannot be omitted from the trace packet.

\subsection{Smith computation on a periodic grid}

Take
\[
                         K=(\mathbb Z/M\mathbb Z)^D         \tag{26.12}
\]
and choose integer representatives
\(s_1,\ldots,s_r\in\mathbb Z^D\) for the background shifts.  Form the
integer matrix
\[
                         {\cal S}
 =
 \begin{bmatrix}
 MI_D&s_1&\cdots&s_r
 \end{bmatrix}.                                             \tag{26.13}
\]
Let
\[
                         U{\cal S}V
 =
 \operatorname{diag}(d_1,\ldots,d_D,0,\ldots,0)             \tag{26.14}
\]
be its Smith normal form, with
\(d_1\mid\cdots\mid d_D\).

\begin{theorem}[Exact number and size of legal packets]
\label{thm:v12v26-Smith}
The quotient of momentum labels by the shift subgroup has the
presentation
\[
 K/H_S
 \cong
 \mathbb Z^D/
 \left(M\mathbb Z^D+\sum_{j=1}^r\mathbb Zs_j\right)
 \cong
 \bigoplus_{a=1}^D\mathbb Z/d_a\mathbb Z.                  \tag{26.15}
\]
Hence
\[
 \#(K/H_S)=\prod_{a=1}^Dd_a,\qquad
 |H_S|=\frac{M^D}{\prod_{a=1}^Dd_a}.                       \tag{26.16}
\]
\end{theorem}

\begin{proof}
The columns of (26.13) generate the lattice
\(M\mathbb Z^D+\sum_j\mathbb Zs_j\).  Quotienting
\(\mathbb Z^D\) by that lattice is exactly quotienting
\((\mathbb Z/M\mathbb Z)^D\) by the subgroup generated by the
classes of the \(s_j\).  Smith normal form gives the invariant-factor
decomposition and its order.  The second formula follows from
\(|K|=|H_S|\,|K/H_S|\).
\end{proof}

For one shift \(s=(s_1,\ldots,s_D)\), its order is
\[
 \operatorname{ord}(s)
 =
 \operatorname{lcm}_{a=1}^D
 \frac{M}{\gcd(M,s_a)},                                    \tag{26.17}
\]
with the zero component contributing one.  Every legal packet in
that case has exactly \(\operatorname{ord}(s)\) cards.

\subsection{Impact on the coupled block certificate}

For a coupled SM--Einstein block, the shift set must contain every
momentum displacement generated by the retained background jets:
gauge backgrounds, Higgs/Yukawa insertions, metric/coframe
insertions, ghosts, and any routed boundary source.  The union is
taken before forming \(H_S\).  Colour, spin, flavour, and tensor
indices remain internal fibre coordinates unless their transport
also moves momentum.

The finite certificate is:

\begin{enumerate}
\item enumerate the exact shift support of every required derivative;
\item compute \(H_S\) and its cosets, using (26.13)--(26.16);
\item replace card projectors by the coset projectors \(Q_c\);
\item verify the exact commutators with the complete transport
      matrices;
\item evaluate Ward, determinant, anomaly, and reflection traces only
      on these reducing packets;
\item use the resulting packet projections in the V13 matching
      mismatch \(\delta\).
\end{enumerate}

\begin{assessment}[Status after V26]
\label{ass:v12v26-status}
The correction imported at V25 now has an exact finite solution.
Background-shift support partitions a finite momentum regulator into
cosets of the generated subgroup, and those cosets are the minimal
universal reducing trace packets.  Smith normal form computes their
number and size without a numerical eigenspace guess.
\end{assessment}

\begin{firewall}[Claim boundary after V26]
\label{fire:v12v26-claim}
The physical SM--Einstein shift set has not been extracted from a
selected regulator, and the YM V47 full sixteen-corner trace has not
been evaluated here.  The theorem determines legal packets but does
not prove their Ward or anomaly residuals vanish, establish continuum
quadrature, or construct the full continuum.
\end{firewall}

\subsection{Parent record}

This V26 note appends to the validated V25 file with record
\[
\begin{array}{c|c}
\text{lines}&7414\\
\text{bytes}&254524\\
\text{SHA--256}&
\texttt{E94A4D8A465850BE6A662B86428A081A8D4E2E9B270BF910B5B25A376A9AABB5}.
\end{array}
\]

% =============================================================================
% END APPENDED TWELFTH-CYCLE V26 MATERIAL
% =============================================================================
% =============================================================================
% BEGIN APPENDED TWELFTH-CYCLE V27 MATERIAL
% =============================================================================

\section{Twelfth-cycle V27: rigidity of positive reciprocal
normalizer kernels}
\label{sec:v12v27}

\subsection{The question left by V24}

V24 gave a sufficient reflection condition for a conditionally
normalized counterterm: both the normalizer kernel
\(M=(m_{ab})\) and, more importantly, its entrywise reciprocal
\(N=(m_{ab}^{-1})\) enter different parts of the argument.  We now
show that simultaneous positive semidefiniteness of \(M\) and \(N\)
is possible only on the reflected-square, rank-one branch.

\subsection{Reciprocal positivity theorem}

Let \(M=(M_{ij})_{1\le i,j\le m}\) be Hermitian positive
semidefinite and assume every entry is nonzero.  In particular
\(M_{ii}>0\).  Define its entrywise reciprocal
\[
                         M^{\circ-1}_{ij}:=\frac1{M_{ij}}.   \tag{27.1}
\]
This is not the ordinary matrix inverse.

\begin{theorem}[Positive Hadamard reciprocal rigidity]
\label{thm:v12v27-rigidity}
The following are equivalent:
\begin{enumerate}
\item \(M^{\circ-1}\succeq0\);
\item every \(2\times2\) principal minor of \(M\) has determinant
      zero;
\item \(M\) has rank one;
\item there are nonzero scalars \(u_1,\ldots,u_m\) such that
\[
                         M_{ij}=\overline{u_i}u_j.           \tag{27.2}
\]
\end{enumerate}
On this branch,
\[
                         M^{\circ-1}_{ij}
 =\overline{u_i^{-1}}u_j^{-1},                             \tag{27.3}
\]
so the reciprocal is also rank-one positive semidefinite.
\end{theorem}

\begin{proof}
Positivity of \(M\) gives, for every pair \(i,j\),
\[
                         |M_{ij}|^2\le M_{ii}M_{jj}.         \tag{27.4}
\]
If \(M^{\circ-1}\succeq0\), positivity of its corresponding
\(2\times2\) principal minor gives
\[
 \frac1{M_{ii}M_{jj}}-\frac1{|M_{ij}|^2}\ge0,
\]
or
\[
                         |M_{ij}|^2\ge M_{ii}M_{jj}.         \tag{27.5}
\]
Equations (27.4)--(27.5) are equalities for every pair, proving item
2.

Choose a Gram representation
\(M_{ij}=\langle v_i,v_j\rangle\).  Equality in Cauchy--Schwarz for
every pair says that all nonzero \(v_i\) are collinear.  Since every
diagonal is positive, none is zero.  Hence the Gram matrix has rank
one.  Conversely, a Hermitian positive rank-one matrix with positive
diagonal has the form (27.2).  Entrywise inversion gives (27.3),
which is again a Gram matrix.  This proves all equivalences.
\end{proof}

\begin{corollary}[A positive definite normalizer cannot have a
positive reciprocal]
\label{cor:v12v27-pd}
If \(m>1\) and \(M\succ0\), then
\[
                         M^{\circ-1}\not\succeq0.            \tag{27.6}
\]
\end{corollary}

\begin{proof}
A positive definite \(m\times m\) matrix has rank \(m>1\), contrary
to Theorem~\ref{thm:v12v27-rigidity}.
\end{proof}

This is stronger than the counterexample in V24: failure of
reciprocal positivity is generic on every strictly positive
multi-boundary normalizer, not only on one chosen matrix.

\subsection{An explicit negative reflected witness}

For a pair \(i\ne j\), define its normalized correlation
\[
                         r_{ij}:=
 \frac{|M_{ij}|}{\sqrt{M_{ii}M_{jj}}}\le1.                 \tag{27.7}
\]
If \(r_{ij}<1\), the reciprocal \(2\times2\) block is
\[
 N^{(ij)}
 =
 \begin{pmatrix}
 a&b\\ \overline b&c
 \end{pmatrix},
 \qquad
 a=M_{ii}^{-1},\quad c=M_{jj}^{-1},\quad
 |b|=\frac{\sqrt{ac}}{r_{ij}}.                             \tag{27.8}
\]

\begin{proposition}[Quantitative negative eigenvalue]
\label{prop:v12v27-negative}
The least eigenvalue of (27.8) is
\[
 \boxed{
 \lambda_-^{(ij)}
 =
 \frac{a+c-
 \sqrt{(a-c)^2+4ac/r_{ij}^2}}{2}<0.}                       \tag{27.9}
\]
Therefore any Hermitian positive-semidefinite approximation
\(\widehat N\) to \(M^{\circ-1}\) satisfies
\[
                         \|\widehat N-M^{\circ-1}\|_{\rm op}
 \ge|\lambda_-^{(ij)}|.                                    \tag{27.10}
\]
\end{proposition}

\begin{proof}
Formula (27.9) is the smaller eigenvalue of the Hermitian block.
Since \(r_{ij}<1\), its determinant
\(ac-|b|^2=ac(1-r_{ij}^{-2})\) is negative, so the smaller eigenvalue
is negative.  Restrict any proposed approximation to this principal
two-dimensional subspace.  Weyl's eigenvalue bound shows that moving
the negative eigenvalue to the nonnegative half-line requires an
operator perturbation of size at least its magnitude.
\end{proof}

The associated eigenvector is an explicit positive-time observable
of negative reflected norm.  Thus the failed gate returns a witness,
not merely a scalar diagnostic.

\subsection{Universal positivity preservation}

For a Hermitian kernel \(N\), consider the Schur multiplier
\[
                         \Phi_N(A)=A\circ N.                 \tag{27.11}
\]

\begin{theorem}[Universal conditional normalization criterion]
\label{thm:v12v27-universal}
The map \(\Phi_N\) sends every positive semidefinite matrix \(A\) to
a positive semidefinite matrix if and only if \(N\succeq0\).
Consequently, if \(N=M^{\circ-1}\) for a positive semidefinite
nonzero-entry normalizer \(M\), universal preservation holds if and
only if \(M\) has the reflected-square form (27.2).
\end{theorem}

\begin{proof}
If \(N\succeq0\), Schur-product positivity from V24 gives the forward
implication.  Conversely, apply \(\Phi_N\) to the positive
all-ones matrix \(J\).  Since \(J\circ N=N\), universal positivity
forces \(N\succeq0\).  The final equivalence is
Theorem~\ref{thm:v12v27-rigidity}.
\end{proof}

For one fixed reflected kernel \(A=K\circ C\), the product
\(A\circ M^{\circ-1}\) can remain positive even when the reciprocal
is indefinite.  That is a special cancellation and must be checked
on the complete matrix; it is not a universal normalizer property.

\subsection{Interpretation for shell conditioning}

Suppose \(M_{ab}\) is itself a reflection-positive conditional
partition-function kernel.  Then V24's simple sufficient route
\[
                         M\succeq0,\qquad
                         M^{\circ-1}\succeq0                \tag{27.12}
\]
is equivalent to the exact factorization
\[
                         M_{ab}=\overline{u_a}u_b.           \tag{27.13}
\]
Hence a nontrivial multi-boundary conditional normalizer cannot be
both strictly positive definite and universally harmless under
entrywise division.

There are three legitimate routes:

\begin{enumerate}
\item prove the reflected-square factorization (27.13);
\item test the complete repaired kernel
      \(K\circ C\circ M^{\circ-1}\) directly, retaining every OS
      null vector and the negative-eigenvalue witness;
\item avoid treating the reciprocal normalizer as a separate
      reflected insertion by constructing a fine
      reflection-positive measure first and then taking its exact
      reflected marginal.
\end{enumerate}

The third route is the natural one for the projective RPESM sequence
and will be formalized next.  It mirrors V25--V26: a complete
invariant object can be positive even when an artificially isolated
card or factor is not.

\begin{assessment}[Status after V27]
\label{ass:v12v27-status}
The reciprocal-normalizer gate is now classified exactly on the
positive normalizer branch.  Positive semidefiniteness of both a
nonzero-entry normalizer and its entrywise reciprocal forces rank-one
reflected-square factorization.  Any failure produces an explicit
two-boundary negative eigenvalue and a minimum repair distance.
\end{assessment}

\begin{firewall}[Claim boundary after V27]
\label{fire:v12v27-claim}
No physical RPESM normalizer has been shown to factor as (27.13), and
no complete repaired reflected kernel has been evaluated.  The
theorem does not supply the missing physical Duhamel fibre,
determinant-line scalarization, tightness, or the full continuum.
\end{firewall}

\subsection{Parent record}

This V27 note appends to the validated V26 file with record
\[
\begin{array}{c|c}
\text{lines}&7711\\
\text{bytes}&264711\\
\text{SHA--256}&
\texttt{C5F0A73081B407593EE89987B01F425541EEF47AB842EB80DAA7373FACDB55A6}.
\end{array}
\]

% =============================================================================
% END APPENDED TWELFTH-CYCLE V27 MATERIAL
% =============================================================================
% =============================================================================
% BEGIN APPENDED TWELFTH-CYCLE V28 MATERIAL
% =============================================================================

\section{Twelfth-cycle V28: fine-measure repair followed by exact
reflected marginalization}
\label{sec:v12v28}

\subsection{Why the order of operations matters}

V27 shows that dividing a reflected boundary kernel by a nonfactorized
conditional normalizer is not universally positivity preserving.
There is a different route already compatible with the projective
regulator philosophy: apply the counterterm to the complete fine
measure, prove positivity there, and define the coarse measure by
exact marginalization.  The conditional normalizer then appears once
in the coarse marginal and once inversely in the conditional law;
the two factors cancel in the complete joint kernel.

\subsection{Boundary disintegration}

Let \(B\) be a coarse boundary alphabet and \(F\) a detail alphabet.
Write a nonnegative fine boundary kernel as
\[
 {\cal L}_{(a,\alpha),(b,\beta)}
 =
 K_{ab}\,Q_{ab}(\alpha,\beta),                             \tag{28.1}
\]
where
\[
                         \sum_{\alpha,\beta\in F}
 Q_{ab}(\alpha,\beta)=1                                    \tag{28.2}
\]
whenever \(K_{ab}\ne0\).  Assume
\[
                         {\cal L}\succeq0.                  \tag{28.3}
\]
Let \(C=((C_{a\alpha,b\beta}))\succeq0\) be the reflected boundary
kernel of a fine counterterm insertion.  Put
\[
 {\cal L}^C:={\cal L}\circ C.                              \tag{28.4}
\]
By the Schur-product theorem, \({\cal L}^C\succeq0\).

Define
\[
 m_{ab}:=
 \sum_{\alpha,\beta}
 Q_{ab}(\alpha,\beta)C_{a\alpha,b\beta}.                   \tag{28.5}
\]
When \(m_{ab}\ne0\), the normalized conditional detail kernel is
\[
 Q^C_{ab}(\alpha,\beta)
 :=
 \frac{
 Q_{ab}(\alpha,\beta)C_{a\alpha,b\beta}}
 {m_{ab}}.                                                  \tag{28.6}
\]

\begin{theorem}[Normalizer cancellation in the complete Wilsonian
kernel]
\label{thm:v12v28-cancel}
Define the updated coarse kernel
\[
                         K^C_{ab}:=K_{ab}m_{ab}.             \tag{28.7}
\]
Then
\[
 \boxed{
 K^C_{ab}Q^C_{ab}(\alpha,\beta)
 =
 {\cal L}^C_{(a,\alpha),(b,\beta)}.}                        \tag{28.8}
\]
Moreover,
\[
 \boxed{
 K^C_{ab}
 =
 \sum_{\alpha,\beta}
 {\cal L}^C_{(a,\alpha),(b,\beta)},\qquad K^C\succeq0.}     \tag{28.9}
\]
Thus the updated coarse cylinder and the complete repaired fine
cylinder are reflection positive without requiring the entrywise
reciprocal matrix \((m_{ab}^{-1})\) to be positive semidefinite.
\end{theorem}

\begin{proof}
Substitution of (28.6)--(28.7) proves (28.8) by cancellation of
\(m_{ab}\).  Summing (28.8) over the detail indices gives the first
identity in (28.9).

Let
\[
 R:\mathbb C^B\longrightarrow\mathbb C^{B\times F},
 \qquad
 (Rf)_{a,\alpha}=f_a.                                      \tag{28.10}
\]
Then the marginal kernel is
\[
                         K^C=R^*{\cal L}^CR.                \tag{28.11}
\]
For every \(f\),
\[
 f^*K^Cf=(Rf)^*{\cal L}^C(Rf)\ge0,
\]
which proves the second assertion in (28.9).  Equation (28.8) shows
that no reciprocal normalizer remains in the complete kernel.
\end{proof}

\subsection{Fixed old marginal versus Wilsonian marginal}

If one insists on retaining the old coarse kernel \(K\), the repaired
joint kernel becomes
\[
 {\cal L}^{\rm fix}_{(a,\alpha),(b,\beta)}
 =
 K_{ab}Q^C_{ab}(\alpha,\beta)
 =
 {\cal L}^C_{(a,\alpha),(b,\beta)}m_{ab}^{-1}.              \tag{28.12}
\]
This is exactly the reciprocal-normalizer situation of V24--V27.
It can fail reflection positivity.

By contrast, the Wilsonian choice (28.7) changes the coarse action by
the shell term encoded by \(m_{ab}\) and keeps the complete fine
kernel equal to the manifestly positive Schur product (28.4).

\begin{corollary}[Exact location of the tradeoff]
\label{cor:v12v28-tradeoff}
After a nontrivial counterterm insertion, the following three demands
are not automatic simultaneously:
\begin{enumerate}
\item preserve the pre-counterterm coarse marginal \(K\);
\item use the conditionally normalized detail law \(Q^C\);
\item preserve reflection positivity.
\end{enumerate}
The Wilsonian construction preserves items 2 and 3 by replacing item
1 with the exact updated marginal \(K^C\).  Keeping item 1 requires a
direct positivity test of (28.12), or one of the special
normalizer conditions in V24.
\end{corollary}

\begin{proof}
Equations (28.8) and (28.12) are the two possible joint kernels.
The first is positive by (28.4), while the second contains the
uncontrolled entrywise reciprocal.  The counterexample and rigidity
theorem of V24--V27 show that its positivity is not automatic.
\end{proof}

This is the measure-theoretic analogue of the V25 correction: a
complete invariant or positive object may lose its property when an
intermediate factor is isolated and tested as though it were
autonomous.

\subsection{Ward identities commute with exact marginalization}

Let \(s_F\) be the fine Ward derivation and \(s_C\) the coarse one.
Assume pullback of coarse observables intertwines them:
\[
                         s_F(Rf)=R(s_Cf).                   \tag{28.13}
\]
Let \({\mathbb E}_{{\cal L}^C}\) and
\({\mathbb E}_{K^C}\) denote the globally normalized fine and coarse
expectations.

\begin{theorem}[Ward-compatible reflected pushforward]
\label{thm:v12v28-Ward}
If
\[
                         {\mathbb E}_{{\cal L}^C}[s_FG]=0   \tag{28.14}
\]
for every fine observable in the declared domain, then
\[
                         \boxed{
 {\mathbb E}_{K^C}[s_Cf]=0}                                \tag{28.15}
\]
for every coarse observable whose pullback lies in that domain.
\end{theorem}

\begin{proof}
Exact marginalization gives
\[
 {\mathbb E}_{K^C}[s_Cf]
 =
 {\mathbb E}_{{\cal L}^C}[R(s_Cf)]
 =
 {\mathbb E}_{{\cal L}^C}[s_F(Rf)]
 =0.
\]
\end{proof}

Thus a fine counterterm that is both a reflection-positive insertion
and an exact Ward repair produces a coarse measure with both
properties.  No separate variation of the conditional denominator
is omitted; it is contained in the coarse marginal (28.7).

\subsection{Relation to the V23 martingale estimate}

The conditional law (28.6) remains useful for sampling and for the
\(L^2\) shell filtration.  V23 gives nonvanishing and conditional
variance bounds when \(C\) is close to one.  The present theorem says
how that conditional law must be paired with its coarse factor:
\[
 \boxed{
 \text{coarse update }K\mapsto Km
 \quad+\quad
 \text{conditional update }Q\mapsto QC/m.}                 \tag{28.16}
\]
Only the product is used in the reflected positivity test.

In an append-only projective construction, (28.16) is implemented by
recomputing the coarser effective cylinder from the repaired fine
cylinder.  It need not agree with a previously frozen historical
coarse law.  Historical reuse is an additional equality constraint,
not a condition for existence of the repaired projective family.

\subsection{Cofinal consistency}

For three cutoffs \(Z\to Y\to X\), exact marginalization is
associative:
\[
                         (\pi_{Z/X})_!
 =
 (\pi_{Y/X})_!(\pi_{Z/Y})_!.                               \tag{28.17}
\]
Therefore, if every fine counterterm is applied before pushforward,
the resulting Wilsonian coarse kernels satisfy the exact two-step
cocycle automatically.  Conditional normalizers may be computed in
one step or two, but their complete coarse-times-conditional product
is the same fine repaired kernel.

\begin{proposition}[No normalization-order anomaly]
\label{prop:v12v28-associative}
Suppose all sums are finite and every required conditional
normalizer is nonzero.  Direct elimination \(Z\to X\) and staged
elimination \(Z\to Y\to X\) give the same coarse kernel and the same
normalized conditional joint law.
\end{proposition}

\begin{proof}
Both coarse kernels are the sum of the same repaired \(Z\)-kernel
over all variables above \(X\); finite Fubini gives equality.
Dividing the same joint kernel by the same coarse marginal gives the
same conditional law.
\end{proof}

\subsection{Application to RPESM}

For the named RPESM sequence, the safe counterterm order is:

\begin{enumerate}
\item construct the complete fine reflected cylinder, including
      determinant/Pfaffian scalarization and every boundary factor;
\item insert the V19 Ward counterterm only after its fine reflected
      kernel passes the positive-semidefinite test;
\item take the exact reflected marginal to define the updated coarse
      cylinder;
\item use the conditional quotient for martingale estimates, keeping
      its coarse normalizer factor in the Wilsonian head;
\item transport the full line, Ward, reflection, and locality packet
      across the resulting updated projective family.
\end{enumerate}

This route avoids the reciprocal positivity obstruction but transfers
the burden to the fine counterterm kernel and to the exact
recomputation of all coarser effective cylinders.

\begin{assessment}[Status after V28]
\label{ass:v12v28-status}
Simultaneous Ward repair and reflection positivity now have a viable
order of operations.  Reweight the complete fine positive kernel,
then marginalize.  The conditional normalizer and its reciprocal
cancel in the joint cylinder, while the normalizer is retained as
the exact Wilsonian coarse action.
\end{assessment}

\begin{firewall}[Claim boundary after V28]
\label{fire:v12v28-claim}
No physical fine RPESM counterterm kernel has been proved reflection
positive, and no repaired coarse family has been recomputed.  The
theorem does not scalarize the chiral determinant line, establish
uniform locality or tightness, populate the missing Duhamel fibre, or
construct the full continuum.
\end{firewall}

\subsection{Parent record}

This V28 note appends to the validated V27 file with record
\[
\begin{array}{c|c}
\text{lines}&7949\\
\text{bytes}&273078\\
\text{SHA--256}&
\texttt{1D469E895A3A0C86B82F377D69B27A163B45DDC70C63896C35627C7CB161E50F}.
\end{array}
\]

% =============================================================================
% END APPENDED TWELFTH-CYCLE V28 MATERIAL
% =============================================================================
% =============================================================================
% BEGIN APPENDED TWELFTH-CYCLE V29 MATERIAL
% =============================================================================

\section{Twelfth-cycle V29: density martingales and positive
counterterm cocycles}
\label{sec:v12v29}

V28 proved that a counterterm compatible with reflection positivity
must be inserted into the complete fine reflected kernel before the
coarse marginal is recomputed.  There is a second compatibility
condition: recomputed marginals obtained from different terminal
cutoffs must define one projective family.  The old V79 theorem gives
a limit once adjacent measure errors are summable.  The present note
addresses an earlier algebraic question.  It characterises exact
projectivity of normalized counterterm densities and gives a
finite-matrix certificate by which positivity survives a cofinal
counterterm cocycle.

\subsection{Conditional transfer of a density}

Let
\[
 E_0\longleftarrow E_1\longleftarrow E_2\longleftarrow\cdots
\]
be standard Borel spaces, with coarse maps
\(p_{n+1,n}:E_{n+1}\to E_n\).  Let \(\mu_n\) be probability
measures satisfying
\begin{equation}
 (p_{n+1,n})_*\mu_{n+1}=\mu_n.                 \tag{29.1}
\end{equation}
For \(f\in L^1(\mu_{n+1})\), define the conditional transfer
\(T_nf\in L^1(\mu_n)\) by
\begin{equation}
 \int_A T_nf\,d\mu_n
 =
 \int_{p_{n+1,n}^{-1}(A)} f\,d\mu_{n+1},
 \qquad A\in{\cal B}(E_n).                     \tag{29.2}
\end{equation}
Thus \(T_nf\) is the conditional expectation of \(f\), expressed as
a function of the retained variable.  It obeys
\begin{equation}
 \int T_nf\,d\mu_n=\int f\,d\mu_{n+1},
 \qquad
 T_n\bigl((g\circ p_{n+1,n})f\bigr)=g\,T_nf       \tag{29.3}
\end{equation}
whenever the products are integrable.

Let \(z_n\in L^1(\mu_n)\), put
\[
 Z_n:=\int z_n\,d\mu_n,
 \qquad Z_n\ne0,
\]
and define the normalized finite complex measure
\begin{equation}
 d\nu_n={z_n\over Z_n}\,d\mu_n.                 \tag{29.4}
\end{equation}
For nonnegative \(z_n\), this is a probability measure.

\begin{theorem}[Exact density-martingale criterion]
\label{thm:v12v29-density-martingale}
The reweighted measures \((\nu_n)\) are exactly projective if and
only if
\begin{equation}
 {T_nz_{n+1}\over Z_{n+1}}
 =
 {z_n\over Z_n}
 \quad\text{\(\mu_n\)-almost everywhere}         \tag{29.5}
\end{equation}
for every \(n\).
\end{theorem}

\begin{proof}
For bounded measurable \(g:E_n\to\mathbb C\), equations
\eqref{eq:v12v29-test-transfer-a} below are obtained directly from
the definition of \(T_n\):
\begin{align}
 \int g\,d(p_{n+1,n})_*\nu_{n+1}
 &=
 {1\over Z_{n+1}}
 \int (g\circ p_{n+1,n})z_{n+1}\,d\mu_{n+1}
 \nonumber\\
 &=
 \int g\,{T_nz_{n+1}\over Z_{n+1}}\,d\mu_n.
                                                        \tag{29.6}
 \label{eq:v12v29-test-transfer-a}
\end{align}
The same test against \(\nu_n\) is
\[
 \int g\,{z_n\over Z_n}\,d\mu_n.
\]
Equality for every bounded measurable \(g\) is equivalent to equality
of the two Radon--Nikodym densities, which is precisely
\eqref{eq:v12v29-test-transfer-b}.
\begin{equation}
 {T_nz_{n+1}\over Z_{n+1}}
 =
 {z_n\over Z_n}.                               \tag{29.7}
 \label{eq:v12v29-test-transfer-b}
\end{equation}
\end{proof}

The criterion is valid for complex weights; positivity is not used.
It shows that nonvanishing of each scalar normalizer is insufficient
for projectivity.  The normalized density itself must be a reverse
conditional-expectation chain.

\subsection{Why a shell counterterm changes the coarse action}

Suppose
\begin{equation}
 z_{n+1}
 =
 (z_n\circ p_{n+1,n})c_{n+1},                  \tag{29.8}
\end{equation}
where \(c_{n+1}\) is an integrable shell counterterm factor.  Define
its conditional normalizer
\begin{equation}
 m_n:=T_nc_{n+1}.                               \tag{29.9}
\end{equation}
The module identity in \eqref{eq:v12v29-shell-module} gives
\begin{equation}
 T_nz_{n+1}=z_nm_n.                             \tag{29.10}
 \label{eq:v12v29-shell-module}
\end{equation}

\begin{proposition}[Historical-coarse-law obstruction]
\label{prop:v12v29-historical-obstruction}
Assume \(z_n\ne0\) almost everywhere.  The old reweighted coarse law
\(\nu_n\) is the pushforward of \(\nu_{n+1}\) if and only if
\begin{equation}
 m_n={Z_{n+1}\over Z_n}
 \quad\text{\(\mu_n\)-almost everywhere}.        \tag{29.11}
\end{equation}
For a nonconstant \(m_n\), the exact unnormalized coarse weight is
instead
\begin{equation}
 z_n^{\rm new}:=z_nm_n.                         \tag{29.12}
\end{equation}
\end{proposition}

\begin{proof}
Insert \eqref{eq:v12v29-shell-module} into the criterion
\eqref{eq:v12v29-test-transfer-b} and cancel \(z_n\) on its conull
nonzero set.  This proves the equivalence.  Without imposing the
constant condition, equation \eqref{eq:v12v29-shell-module} states
that the pushforward of the unnormalized fine weight has density
\(z_nm_n\) relative to \(\mu_n\), proving
\eqref{eq:v12v29-updated-weight}.
\begin{equation}
 (p_{n+1,n})_*(z_{n+1}\mu_{n+1})
 =
 z_n^{\rm new}\mu_n.                            \tag{29.13}
 \label{eq:v12v29-updated-weight}
\end{equation}
\end{proof}

This is the measure-theoretic form of the V28 instruction to retain
the conditional normalizer in the Wilsonian head.  Freezing the old
coarse action silently imposes the strong constant-normalizer
condition \eqref{eq:v12v29-constant-condition}.
\begin{equation}
 T_nc_{n+1}\ \text{is constant}.                \tag{29.14}
 \label{eq:v12v29-constant-condition}
\end{equation}

\subsection{The inverse-limit martingale}

Let \(\mu_\infty\) be the inverse-limit probability measure associated
with \eqref{eq:v12v29-base-projective}, and let
\(X_n:E_\infty\to E_n\) be the coordinate maps.  Set
\[
 h_n={z_n\over Z_n},
 \qquad
 {\cal F}_n=\sigma(X_n).
\]
Here
\begin{equation}
 (p_{n+1,n})_*\mu_{n+1}=\mu_n                  \tag{29.15}
 \label{eq:v12v29-base-projective}
\end{equation}
is recalled only to identify the joint inverse-limit law.

\begin{corollary}[Uniformly integrable terminal density]
\label{cor:v12v29-terminal-density}
Equation \eqref{eq:v12v29-test-transfer-b} holds for every \(n\) if
and only if
\begin{equation}
 {\mathbb E}_{\mu_\infty}
 \bigl[h_{n+1}(X_{n+1})\mid{\cal F}_n\bigr]
 =
 h_n(X_n).                                      \tag{29.16}
\end{equation}
If this martingale is uniformly integrable, there is
\(h_\infty\in L^1(\mu_\infty)\) such that
\begin{equation}
 h_n(X_n)
 =
 {\mathbb E}_{\mu_\infty}
 \bigl[h_\infty\mid{\cal F}_n\bigr],
 \qquad
 h_n(X_n)\longrightarrow h_\infty
 \quad\text{in \(L^1\)}.                         \tag{29.17}
\end{equation}
The finite measures are then exactly the cylinder marginals of
\begin{equation}
 d\nu_\infty=h_\infty\,d\mu_\infty.             \tag{29.18}
\end{equation}
For nonnegative densities, \(\nu_\infty\) is a probability measure.
\end{corollary}

\begin{proof}
The conditional-expectation identity is exactly the transfer identity
of Theorem \ref{thm:v12v29-density-martingale} pulled to the inverse
limit.  The uniformly integrable martingale convergence theorem gives
\(L^1\) convergence and the conditional representation in
\eqref{eq:v12v29-martingale-limit}.  Testing a cylinder function
against \(h_\infty\mu_\infty\) and conditioning on \({\cal F}_n\)
gives \(\nu_n\).
\begin{equation}
 h_n(X_n)\to h_\infty\quad\hbox{in }L^1.
                                                        \tag{29.19}
 \label{eq:v12v29-martingale-limit}
\end{equation}
If \(h_n\ge0\) and \(\int h_n\,d\mu_n=1\), uniform integrability
preserves unit mass in the limit.
\end{proof}

\begin{warning}[Mass loss without uniform integrability]
\label{warn:v12v29-mass-loss}
A nonnegative unit-mean martingale converges almost surely, but its
limit can have expectation strictly less than one.  Thus pointwise
counterterm convergence does not by itself give a normalized
continuum density.  A uniform-integrability estimate, such as a
uniform \(L^r\) bound for some \(r>1\), is a separate denominator and
tail gate.
\end{warning}

\subsection{A reflected-kernel cocycle certificate}

The density criterion controls projectivity.  Reflection positivity
is controlled at the level of the complete reflected boundary
kernel.  At a finite terminal cutoff \(N\), write the unmodified
kernel as a Hermitian matrix \(K_N\succeq0\) on the positive-half
boundary index set.  Suppose counterterm increments \(j=1,\ldots,N\)
lift to Hermitian matrices \(C_{j,N}\succeq0\) on the same index set.
Define
\begin{equation}
 L_N
 :=
 K_N\circ C_{1,N}\circ\cdots\circ C_{N,N},       \tag{29.20}
\end{equation}
where \(\circ\) is entrywise product.

For a coarser reflected boundary index set at level \(n\le N\), let
\(R_{N,n}\) be the linear marginalization map, including the sums and
positive square-root weights of eliminated boundary variables.  The
complete coarse kernel descended from terminal level \(N\) is
\begin{equation}
 L_{N\downarrow n}:=R_{N,n}^*L_NR_{N,n}.         \tag{29.21}
\end{equation}

\begin{theorem}[Positive counterterm cocycle]
\label{thm:v12v29-positive-cocycle}
Under the preceding finite-dimensional hypotheses:
\begin{enumerate}
\item \(L_N\succeq0\) for every \(N\);
\item \(L_{N\downarrow n}\succeq0\) for every \(n\le N\);
\item any entrywise limit
\[
 L_{\infty\downarrow n}
 =
 \lim_{N\to\infty}L_{N\downarrow n}
\]
is positive semidefinite; and
\item if the descended normalized kernels are compatible under all
further marginal maps, their inverse-limit cylinder functional obeys
the Osterwalder--Schrader inequality on every finite positive-half
cylinder observable.
\end{enumerate}
\end{theorem}

\begin{proof}
The Schur product theorem applied repeatedly gives the first claim.
For every vector \(v\),
\[
 v^*R_{N,n}^*L_NR_{N,n}v
 =
 (R_{N,n}v)^*L_N(R_{N,n}v)\ge0,
\]
which proves the second.  The cone of positive semidefinite matrices
of fixed finite size is closed, proving the third.  A finite cylinder
observable is represented at some level \(n\) by a vector \(v\).  Its
reflected quadratic expectation is
\(v^*L_{\infty\downarrow n}v\), up to a positive scalar
normalization.  This is nonnegative.  Compatibility makes the value
independent of the chosen finer representation.
\end{proof}

\begin{remark}[Pullback preserves the increment certificate]
\label{rem:v12v29-pullback}
If a counterterm kernel \(C_j\) is positive semidefinite on a coarser
index set and \(q\) maps fine indices to coarse indices, then its
pullback \(\widetilde C_{\alpha\beta}=C_{q(\alpha)q(\beta)}\) is
positive semidefinite.  Indeed, it is \(Q^*C_jQ\) for the incidence
matrix \(Q\).  Hence coarse positive counterterms may be lifted into
a common terminal kernel before the Schur product is taken.
\end{remark}

\subsection{Combined exact gate}

Theorems \ref{thm:v12v29-density-martingale} and
\ref{thm:v12v29-positive-cocycle} impose different conditions:

\begin{enumerate}
\item the normalized weights must satisfy the conditional
density-martingale equation \eqref{eq:v12v29-test-transfer-b};
\item each full terminal counterterm insertion must lie in the
reflected positive cone, for example by the Schur certificate
\eqref{eq:v12v29-positive-product}; and
\item uniform integrability must prevent terminal density mass from
escaping.
\end{enumerate}
Here the certificate is
\begin{equation}
 K_N\succeq0,\quad C_{j,N}\succeq0
 \quad\Longrightarrow\quad
 K_N\circ\prod_{j=1}^N{}^{\circ}C_{j,N}\succeq0.
                                                        \tag{29.22}
 \label{eq:v12v29-positive-product}
\end{equation}
None of these three rows implies either of the other two.

\begin{assessment}[Status after V29]
\label{ass:v12v29-status}
The coarse-normalizer update of V28 now has an exact projective
criterion.  A compatible normalized counterterm tower is precisely a
density martingale.  A uniformly integrable tower has a terminal
inverse-limit density, while a Schur-positive terminal counterterm
cocycle preserves every finite reflected quadratic form through
marginalization and cofinal limits.
\end{assessment}

\begin{firewall}[Claim boundary after V29]
\label{fire:v12v29-claim}
No selected RPESM counterterm has yet been shown to satisfy the
density-martingale equation or the reflected Schur certificate.
Uniform integrability of the full interacting density remains
unproved.  The theorem neither supplies the missing physical Duhamel
fibre nor constructs the full SM--Einstein continuum.
\end{firewall}

\subsection{Parent record}

This V29 note appends to the validated V28 file with record
\[
\begin{array}{c|c}
\text{lines}&8249\\
\text{bytes}&283353\\
\text{SHA--256}&
\texttt{A79BFE05D11A5DEF7A5155FD23F7151E46D286CC08AADCEBDC8BE622ABA6A5B0}.
\end{array}
\]

% =============================================================================
% END APPENDED TWELFTH-CYCLE V29 MATERIAL
% =============================================================================
% =============================================================================
% BEGIN APPENDED TWELFTH-CYCLE V30 MATERIAL
% =============================================================================

\section{Twelfth-cycle V30: compulsory companion audit at the
counterterm-projectivity gate}
\label{sec:v12v30}

\subsection{Files and fingerprints}

This is the compulsory five-version audit following V25.  The newest
active companion files available during the audit were read without
modification:
\[
\begin{array}{c|r|r}
\text{file}&\text{lines}&\text{bytes}\\ \hline
\texttt{Renewal\_Yang\_Mills\_Mass\_Gap\_Research\_Notes\_V49.tex}
 &19212&658460\\
\texttt{Renewal\_NS\_Stability\_Research\_Notes\_V26.tex}
 &5320&278283.
\end{array}                                                   \tag{30.1}
\]
Their SHA--256 fingerprints are, respectively,
\[
\begin{split}
 &\texttt{6A8F5EEFAAD1AB422D61CA5B82F221647437844929DBA7D57B1264E087A59FD7},
 \\
 &\texttt{82C887F8C2C69E4F28FAD7C3DC8DE5B9099CB5A4BF431EBA1FFF3E91935978AD}.
\end{split}                                                   \tag{30.2}
\]
The Yang--Mills file advanced from V47 to V49 after the preceding
audit.  The Navier--Stokes file is unchanged.

\subsection{Yang--Mills transfer: exact finite arithmetic does not
certify a continuum observable}

Yang--Mills V48 removes the false pointwise orbit compression,
evaluates the uncompressed sixteen-corner response, and obtains a
strictly negative second response at \(M=2\).  V49 then proves that
the inherited analytic-strip estimate cannot transport that sign to
the Brillouin integral.  Its rigorous necessary resolution bound is
already larger than \(17\,220\,408\), and the sharper diagnostic
threshold is larger still.  The missing input is a much wider
certified complex inverse domain together with populated response
coefficient norms.

The transferable result is a rule about proof scope.

\begin{proposition}[Finite response data do not close a continuum
certificate]
\label{prop:v12v30-finite-not-continuum}
Let \(f\) be a periodic response integrand, let
\(\langle f\rangle_M\) be a finite grid average, and suppose the only
proved continuum comparison is
\begin{equation}
 \left|\langle f\rangle_M-\langle f\rangle\right|
 \leq N C_M.                                      \tag{30.3}
\end{equation}
If \(NC_M\geq|\langle f\rangle_M|\), then the data determine no sign
of \(\langle f\rangle\).
\end{proposition}

\begin{proof}
The certified interval
\[
 [\,\langle f\rangle_M-NC_M,\,
    \langle f\rangle_M+NC_M\,]
\]
contains zero.  It therefore contains values with both signs unless
one endpoint is zero, in which case it still does not prove a strict
sign.
\end{proof}

For the SM--Einstein programme, a selected finite Duhamel fibre or a
finite table of physical cards may be used only after its quadrature,
aliasing, and complex-resolvent errors are smaller than the claimed
continuum margin.  In particular, the missing physical RPESM
right-energy fibre cannot be replaced by a sparse sample and a sign
guess.

The same point applies to the V29 reflected counterterm condition.
Positive semidefiniteness of a few momentum cards does not imply
positive semidefiniteness of the complete reflected boundary kernel.
A legal momentum packet from V26 must be tested as a whole, and any
passage to a continuum momentum variable requires a uniform spectral
error bound.

\subsection{Navier--Stokes transfer: use a structured norm only on a
proved structured cone}

Navier--Stokes V26 contracts the first and second derivatives of the
Oseen kernel against inputs \(a\otimes a\).  It proves an exact
one-variable angular maximization.  The first-derivative structured
norm gives no improvement, while the second-derivative norm improves
the ambient symmetric-traceless norm on a finite radial range.

The useful general statement is elementary but prevents a dangerous
misuse.

\begin{lemma}[Cone-restricted norm rule]
\label{lem:v12v30-cone-norm}
Let \(A:V\to W\) be linear and let \(S\subset V\).  The bound
\[
 \|As\|\leq C_S\|s\|,\qquad s\in S,
\]
extends to a vector \(v\) only after one proves \(v\in S\), or after
one represents \(v\) as a controlled combination of elements of
\(S\).  A smaller value of \(C_S\) alone gives no bound on
\(\|A\|_{V\to W}\).
\end{lemma}

\begin{proof}
By definition, \(C_S\) controls only the displayed subset.  For
example, take \(V=\mathbb R^2\), \(S=\operatorname{span}(e_1)\), and
\(A=\operatorname{diag}(0,R)\).  Then \(C_S=0\), while
\(\|A\|=R\) is arbitrary.
\end{proof}

In the reflected-kernel problem, the OS inequality is tested against
every vector in the positive-half observable space.  Rank-one or
field-generated test vectors are therefore insufficient unless a
proved cone-generation theorem reduces all OS tests to them.  The
full positive-semidefinite test in V29 remains mandatory.  Restricted
norms can still sharpen estimates for a separately proved structured
Ward source.

\subsection{Direction selected after the audit}

The audit rules out two shortcuts:

\begin{enumerate}
\item do not infer the missing continuum Duhamel sign, positivity, or
mass parameter from finitely many exact response cards without a
certified continuum enclosure;
\item do not replace full reflected-kernel positivity by a
rank-one-input estimate without a cone-generation proof.
\end{enumerate}

The next upstream question is finite and exact.  Given a singular
positive reflected kernel \(K\) and an affine Ward equation for a
counterterm exponent \(B\), what necessary condition must \(B\)
satisfy for
\[
                         K\circ e^{-tB}
\]
to remain positive semidefinite for small \(t>0\)?  A strict version
of that condition will give a local sufficient theorem, and a dual
matrix witness will certify infeasibility.  This test can be run
before one attempts a continuum response estimate.

\begin{assessment}[Status after V30]
\label{ass:v12v30-status}
The companion audit preserves the V29 density-martingale and
Schur-positive cocycle gates.  It adds two restrictions on their
use: finite momentum arithmetic requires a quantitative continuum
enclosure, and OS positivity remains a full-cone condition.  The
next note will attack the finite semidefinite intersection of the
Ward affine space with the tangent cone of the reflected positive
cone.
\end{assessment}

\begin{firewall}[Claim boundary after V30]
\label{fire:v12v30-claim}
No Yang--Mills continuum response sign or Navier--Stokes regularity
claim is imported.  No physical SM--Einstein Duhamel fibre,
reflection-positive Ward counterterm, uniform determinant-line
bound, or full continuum has been constructed.
\end{firewall}

\subsection{Parent record}

This V30 note appends to the validated V29 file with record
\[
\begin{array}{c|c}
\text{lines}&8616\\
\text{bytes}&296270\\
\text{SHA--256}&
\texttt{9F88F8C26AFF8B63F2C13E7C06BC4F974D59A0B241B8D39DDEB1E8D54D267E83}.
\end{array}
\]

% =============================================================================
% END APPENDED TWELFTH-CYCLE V30 MATERIAL
% =============================================================================
% =============================================================================
% BEGIN APPENDED TWELF-CYCLE V31 MATERIAL
% =============================================================================

\section{Twelfth-cycle V31: the Ward affine space inside the
reflection-positive tangent cone}
\label{sec:v12v31}

V24 gives perturbative protection when the reflected kernel has a
strict positive gap.  Physical reflected kernels can be singular:
null vectors arise from constraints, quotient directions, or exact
relations among boundary observables.  In that case an arbitrarily
small counterterm can leave the positive cone at first order.  This
note gives the exact finite-dimensional tangent test and turns it
into a semidefinite feasibility problem for the Ward primitive.

\subsection{The entrywise exponential path}

Let \(K\in{\rm Herm}_d\) satisfy \(K\succeq0\).  Let
\[
 P_0:{\mathbb C}^d\longrightarrow\ker K
\]
be the orthogonal projection, and let \(P_1=I-P_0\).  For a Hermitian
matrix \(B\), define the Hermitian path
\begin{equation}
 C_B(t)_{ij}:=\exp(-tB_{ij}),
 \qquad
 L_B(t):=K\circ C_B(t),\qquad t\geq0.             \tag{31.1}
\end{equation}
Here \(\circ\) denotes entrywise product.  Since
\(C_B(0)\) is the all-ones matrix,
\[
 L_B(0)=K,\qquad
 H_B:=L_B'(0)=-K\circ B.                          \tag{31.2}
\]

\begin{theorem}[Necessary null-space tangent condition]
\label{thm:v12v31-tangent-necessary}
If \(L_B(t)\succeq0\) for every sufficiently small \(t\geq0\), then
\begin{equation}
 {\cal L}_K(B):=-P_0(K\circ B)P_0\succeq0
 \quad\text{on }\ker K.                           \tag{31.3}
\end{equation}
The same conclusion holds for a differentiable counterterm path
whose exponent is \(tB+o(t)\).
\end{theorem}

\begin{proof}
Let \(v\in\ker K\).  Positivity gives
\[
 \phi_v(t):=v^*L_B(t)v\geq0,\qquad \phi_v(0)=0.
\]
The right derivative at a minimum on the half-line is nonnegative.
Using \eqref{eq:v12v31-derivative},
\begin{equation}
 \phi_v'(0)
 =
 v^*H_Bv
 =
 -v^*(K\circ B)v\geq0.                            \tag{31.4}
 \label{eq:v12v31-derivative}
\end{equation}
This holds for every \(v\in\ker K\), which is exactly
\eqref{eq:v12v31-null-lmi}.  An \(o(t)\) correction to the exponent
does not change the derivative.
\begin{equation}
 {\cal L}_K(B)\succeq0.                           \tag{31.5}
 \label{eq:v12v31-null-lmi}
\end{equation}
\end{proof}

When \(K\) is positive definite, \(P_0=0\) and the condition is
vacuous; continuity then protects positivity for a sufficiently
small step.  The new obstruction occurs precisely at null
directions.

\subsection{A strict tangent gives a genuine positive path}

The necessary condition has a useful strict converse.

\begin{theorem}[Strict null-space tangent is locally sufficient]
\label{thm:v12v31-tangent-sufficient}
Suppose either \(K\succ0\), or
\begin{equation}
 {\cal L}_K(B)\succ0\quad\text{on }\ker K.         \tag{31.6}
\end{equation}
Then there is \(\varepsilon>0\) such that
\begin{equation}
 L_B(t)\succeq0\qquad(0\leq t\leq\varepsilon).
                                                               \tag{31.7}
\end{equation}
In the second case \(L_B(t)\succ0\) for \(0<t\leq\varepsilon\).
\end{theorem}

\begin{proof}
If \(K\succ0\), the claim follows from continuity of the least
eigenvalue.  Suppose \(K\) is singular.  Decompose
\[
 {\mathbb C}^d=P_1{\mathbb C}^d\oplus P_0{\mathbb C}^d.
\]
On the first summand, \(K_1:=P_1KP_1\) is positive definite.  Since
the entrywise exponential is analytic,
\begin{equation}
 L_B(t)=K+tH_B+R(t),\qquad \|R(t)\|=O(t^2).        \tag{31.8}
\end{equation}
Relative to the displayed decomposition, write
\[
 L_B(t)=
 \begin{pmatrix}
 A(t)&G(t)\\G(t)^*&D(t)
 \end{pmatrix}.
\]
There are constants \(\lambda,\gamma>0\) such that
\[
 K_1\succeq\lambda P_1,
 \qquad
 P_0H_BP_0={\cal L}_K(B)\succeq\gamma P_0.
\]
Equation (31.8) gives, for all sufficiently small \(t>0\),
\begin{align}
 A(t)&\succeq{\lambda\over2}P_1,                 \tag{31.9}\\
 D(t)&\succeq{\gamma t\over2}P_0,                \tag{31.10}\\
 \|G(t)\|&=O(t).                                 \tag{31.11}
\end{align}
The Schur complement of \(A(t)\) is
\[
 D(t)-G(t)^*A(t)^{-1}G(t).
\]
The second term is \(O(t^2)\), while the first is bounded below by
\(\gamma t/2\).  Shrinking \(\varepsilon\) makes the Schur complement
strictly positive.  Since \(A(t)\succ0\), the block matrix is then
strictly positive.
\end{proof}

A non-strict tangent is only a necessary gate in this note.  When
\({\cal L}_K(B)\) has its own null space, higher derivatives must be
tested on the unresolved subspace.

\subsection{Intersection with the Ward equation}

Let \(Y\) be a finite-dimensional real inner-product space containing
the selected Ward-anomaly cards, and let
\[
 {\cal D}:{\rm Herm}_d\longrightarrow Y
\]
be the real-linear map taking a counterterm exponent to its
linearized Ward correction.  With the sign convention fixed in
\({\cal D}\), the required equation is
\begin{equation}
                         {\cal D}B=a.             \tag{31.12}
\end{equation}
Define the reflected tangent cone in exponent coordinates by
\begin{equation}
 {\cal T}_K
 :=
 \{B\in{\rm Herm}_d:{\cal L}_K(B)\succeq0\}.      \tag{31.13}
\end{equation}

\begin{corollary}[Finite Ward--OS feasibility problem]
\label{cor:v12v31-feasibility}
A necessary condition for a differentiable reflection-positive
counterterm path with first exponent \(B\) to cancel the linearized
anomaly \(a\) is
\begin{equation}
 \boxed{\quad {\cal D}B=a,\qquad
               {\cal L}_K(B)\succeq0.\quad}       \tag{31.14}
\end{equation}
This is a finite semidefinite feasibility problem.  If it has a
solution with
\({\cal L}_K(B)\succ0\) on \(\ker K\), that solution generates a
reflection-positive entrywise exponential path for all sufficiently
small positive \(t\).
\end{corollary}

\begin{proof}
Necessity combines the Ward equation with Theorem
\ref{thm:v12v31-tangent-necessary}.  The cone in (31.13) is the
inverse image of the positive semidefinite cone under the linear map
\({\cal L}_K\), so (31.14) is an affine linear matrix inequality.
The strict conclusion is Theorem
\ref{thm:v12v31-tangent-sufficient}.
\end{proof}

The qualifier \emph{linearized} is essential.  Passing (31.14) does
not prove the nonlinear finite Ward identity.  It proves that the
first Ward correction is not already incompatible with reflected
positivity.

\subsection{A dual obstruction certificate}

Equip Hermitian matrices with
\(\langle A,B\rangle=\operatorname{Re}\operatorname{tr}(AB)\).
Let \({\cal D}^*\) and \({\cal L}_K^*\) be the adjoints, the latter
defined on Hermitian operators on \(\ker K\).

\begin{proposition}[Dual certificate of incompatibility]
\label{prop:v12v31-dual}
Suppose there are \(y\in Y\) and \(Q\succeq0\) on \(\ker K\) such
that
\begin{equation}
 {\cal D}^*y={\cal L}_K^*Q,
 \qquad
 \langle y,a\rangle<0.                           \tag{31.15}
\end{equation}
Then the feasibility problem (31.14) has no solution.  Consequently
no differentiable entrywise counterterm path with that Ward tangent
can preserve reflection positivity.
\end{proposition}

\begin{proof}
If \(B\) were feasible, then
\begin{align*}
 \langle y,a\rangle
 &=
 \langle y,{\cal D}B\rangle
 =
 \langle{\cal D}^*y,B\rangle\\
 &=
 \langle{\cal L}_K^*Q,B\rangle
 =
 \langle Q,{\cal L}_K(B)\rangle
 \geq0.
\end{align*}
The last inequality holds because both matrices are positive
semidefinite.  This contradicts (31.15).
\end{proof}

This proposition is a one-sided certificate and needs no
constraint qualification.  Under standard semidefinite
regularity, conic duality can also make such certificates complete,
but that stronger assertion is not needed here.

\subsection{An exact reflection-square subspace}

There is a smaller linear family that preserves positivity for every
\(t\), without a null-space calculation.  For \(b\in\mathbb C^d\),
define
\begin{equation}
 ({\cal S}b)_{ij}:=\overline{b_i}+b_j.            \tag{31.16}
\end{equation}
Then \({\cal S}b\) is Hermitian.

\begin{proposition}[Reflection-split Ward primitive]
\label{prop:v12v31-reflection-split}
If the Ward equation has a solution \(B={\cal S}b\), then
\begin{equation}
 K\circ C_B(t)=U(t)^*KU(t)\succeq0
 \qquad\text{for every }t\geq0,                  \tag{31.17}
\end{equation}
where
\[
 U(t)=\operatorname{diag}
 \bigl(e^{-tb_1},\ldots,e^{-tb_d}\bigr).
\]
Thus the finite linear equation
\begin{equation}
                         {\cal D}{\cal S}b=a      \tag{31.18}
\end{equation}
is an exact sufficient reflection-positivity certificate for the
entrywise exponential counterterm.
\end{proposition}

\begin{proof}
For every pair \(i,j\),
\[
 C_B(t)_{ij}
 =
 e^{-t\overline{b_i}}e^{-tb_j}.
\]
Therefore
\[
 (K\circ C_B(t))_{ij}
 =
 \overline{U(t)_{ii}}K_{ij}U(t)_{jj},
\]
which is the \((i,j)\) entry of \(U(t)^*KU(t)\).
Congruence preserves positive semidefiniteness.
\end{proof}

Equation (31.18) tests whether the V19 anomaly primitive can be chosen
as a reflected half-space coboundary on the complete boundary index
set.  Failure of this restricted equation does not exclude a more
general solution of (31.14).

\subsection{Executable finite card}

For a selected RPESM cutoff, the finite test is now:

\begin{enumerate}
\item assemble the complete reflected kernel \(K\), including every
determinant, ghost, coarea, stabilizer, and boundary factor;
\item compute its null projection \(P_0\) with certified spectral
intervals;
\item assemble the full Ward map \({\cal D}\) on the legal momentum
packets of V26 and the anomaly card \(a\);
\item first solve the linear reflection-split equation (31.18);
\item if it fails, solve the semidefinite system (31.14), recording
either a strict margin on \(\ker K\) or the dual witness (31.15);
\item only after this tangent gate passes, test the exact nonlinear
Ward identity, the V29 density-martingale equation, and uniform
cofinal bounds.
\end{enumerate}

\begin{assessment}[Status after V31]
\label{ass:v12v31-status}
The first compatibility test between Ward repair and a singular
reflected kernel is now exact.  The Ward affine space must intersect
the pullback of the positive-semidefinite tangent cone.  Strict
intersection yields a genuine small positive path, a dual pair
\((y,Q)\) can rigorously reject intersection, and a
reflection-split primitive preserves positivity for every
counterterm strength.
\end{assessment}

\begin{firewall}[Claim boundary after V31]
\label{fire:v12v31-claim}
The finite RPESM matrices \(K\), \({\cal D}\), and \(a\) have not been
populated, and no feasibility or dual certificate has been computed.
The tangent test does not imply exact nonlinear Ward cancellation,
cofinal uniformity, density-martingale compatibility, determinant
line trivialization, or the full SM--Einstein continuum.
\end{firewall}

\subsection{Parent record}

This V31 note appends to the validated V30 file with record
\[
\begin{array}{c|c}
\text{lines}&8800\\
\text{bytes}&303407\\
\text{SHA--256}&
\texttt{55A0AAFD432F4DB01B5054714C7ABE119C83A97410DD1CB9B8CD56109182E80D}.
\end{array}
\]

% =============================================================================
% END APPENDED TWELFTH-CYCLE V31 MATERIAL
% =============================================================================
% =============================================================================
% BEGIN APPENDED TWELFTH-CYCLE V32 MATERIAL
% =============================================================================

\section{Twelfth-cycle V32: second-order reflected positivity on
unresolved null directions}
\label{sec:v12v32}

The V31 strict tangent theorem does not decide the common boundary
case in which the first tangent is positive semidefinite but has a
kernel.  This note derives the next coefficient.  The new term is not
merely the second derivative of the entrywise counterterm: it contains
a negative Schur-complement correction caused by mixing with the
positive spectral subspace of the original reflected kernel.

\subsection{General second-order Schur reduction}

Let \(K\succeq0\) be a Hermitian \(d\times d\) matrix.  Retain the
orthogonal projections \(P_0\) onto \(\ker K\) and
\(P_1=I-P_0\).  Let
\begin{equation}
 L(t)=K+tH+t^2G+O(t^3)                            \tag{32.1}
\end{equation}
be a Hermitian analytic path in operator norm.  Denote by
\(K^\dagger\) the Moore--Penrose inverse; it is the inverse of \(K\)
on \(P_1\mathbb C^d\) and zero on \(\ker K\).  Define, on
\(P_0\mathbb C^d\),
\begin{align}
 A&:=P_0HP_0,                                    \tag{32.2}\\
 E&:=P_0(G-HK^\dagger H)P_0.                     \tag{32.3}
\end{align}
When \(A\succeq0\), let \(P_2\) be the orthogonal projection from
\(\ker K\) onto \(\ker A\).

\begin{theorem}[Second-order null-space criterion]
\label{thm:v12v32-second-order}
If \(L(t)\succeq0\) for all sufficiently small \(t\geq0\), then
\begin{equation}
 A\succeq0,\qquad P_2EP_2\succeq0.                \tag{32.4}
\end{equation}
Conversely, if
\begin{equation}
 A\succeq0,\qquad
 P_2EP_2\succ0\quad\text{on }\ker A,              \tag{32.5}
\end{equation}
then \(L(t)\succ0\) for every sufficiently small \(t>0\).
The convention is that the second condition is vacuous when
\(\ker A=\{0\}\).
\end{theorem}

\begin{proof}
On \(P_1\mathbb C^d\), the matrix
\[
 L_{11}(t):=P_1L(t)P_1
\]
is positive definite for all sufficiently small \(t\), because
\(P_1KP_1\) is positive definite.  Thus \(L(t)\succeq0\) is equivalent
to positivity of \(L_{11}(t)\) and of the Schur complement
\begin{equation}
 S(t):=
 P_0L(t)P_0
 -
 P_0L(t)P_1L_{11}(t)^{-1}P_1L(t)P_0.             \tag{32.6}
\end{equation}
The inverse expansion on \(P_1\mathbb C^d\) begins with
\[
 L_{11}(t)^{-1}=K^\dagger+O(t).
\]
Since \(P_0KP_1=0\), the off-diagonal block begins at order \(t\):
\[
 P_0L(t)P_1=tP_0HP_1+O(t^2).
\]
Substitution into (32.6) gives
\begin{equation}
 S(t)=tA+t^2E+O(t^3).                            \tag{32.7}
 \label{eq:v12v32-schur-expansion}
\end{equation}

If \(S(t)\succeq0\), division by \(t>0\) and passage to the limit
gives \(A\succeq0\).  For \(v\in\ker A\),
\[
 v^*S(t)v=t^2v^*Ev+O(t^3)\geq0.
\]
Division by \(t^2\) and passage to the limit gives
\(P_2EP_2\succeq0\), proving necessity.

For sufficiency, write
\[
 {S(t)\over t}=A+tE+O(t^2).
\]
Regard this as a Hermitian path on \(\ker K\), with base matrix
\(A\succeq0\).  Its derivative has strictly positive compression
\(P_2EP_2\) to \(\ker A\).  The block Schur argument used in Theorem
\ref{thm:v12v31-tangent-sufficient} therefore makes \(S(t)/t\),
and hence \(S(t)\), positive definite for all sufficiently small
\(t>0\).  Together with \(L_{11}(t)\succ0\), the Schur-complement
criterion gives \(L(t)\succ0\).
\end{proof}

The term
\begin{equation}
 -P_0HK^\dagger HP_0                             \tag{32.8}
\end{equation}
is negative semidefinite.  Omitting it can turn a failed
second-order test into an apparent pass.

\subsection{Specialization to an entrywise exponential counterterm}

For the V31 path
\[
 L_B(t)=K\circ\exp_\circ(-tB),
\]
where the exponential is taken entrywise, Taylor expansion gives
\begin{align}
 H_B&=-K\circ B,                                 \tag{32.9}\\
 G_B&={1\over2}K\circ B^{\circ2},                \tag{32.10}
\end{align}
where \(B^{\circ2}=B\circ B\).  Define
\begin{align}
 A_B&:=-P_0(K\circ B)P_0,                        \tag{32.11}\\
 E_B&:=P_0\left\{
 {1\over2}K\circ B^{\circ2}
 -(K\circ B)K^\dagger(K\circ B)
 \right\}P_0.                                    \tag{32.12}
\end{align}

\begin{corollary}[Second-order Ward--OS card]
\label{cor:v12v32-Ward-OS}
Let \(B\) solve the linearized Ward equation
\({\cal D}B=a\) of V31.  If \(L_B(t)\succeq0\) for all sufficiently
small \(t\geq0\), then
\begin{equation}
 A_B\succeq0,\qquad
 P_2E_BP_2\succeq0,                              \tag{32.13}
\end{equation}
where \(P_2\) projects onto \(\ker A_B\) inside \(\ker K\).
If the second inequality is strict on that space, the counterterm
path is positive definite for all sufficiently small \(t>0\).
\end{corollary}

\begin{proof}
Insert (32.9)--(32.10) into Theorem
\ref{thm:v12v32-second-order}.
\end{proof}

Thus a boundary solution of the V31 semidefinite program is not ready
for use merely because \(A_B\succeq0\).  It must pass (32.13), or be
certified by an exact structural factorization.

\subsection{A reduced second-order feasibility problem}

Suppose the first-order problem
\[
 {\cal D}B=a,\qquad A_B\succeq0
\]
has solutions but no solution with \(A_B\succ0\).  On a face where
the kernel projection \(P_2\) is fixed, the second-order condition is
\begin{equation}
 P_2P_0\left\{
 {1\over2}K\circ B^{\circ2}
 -(K\circ B)K^\dagger(K\circ B)
 \right\}P_0P_2\succeq0.                         \tag{32.14}
\end{equation}
This is a quadratic matrix inequality in \(B\), rather than an SDP.
It becomes a finite exact polynomial feasibility problem after real
and imaginary parts of the independent Hermitian entries are chosen
as coordinates.

\begin{proposition}[A computable rejection witness]
\label{prop:v12v32-rejection}
Let \(B\) satisfy the first-order Ward equation and tangent condition.
If there is \(v\in\ker A_B\) such that
\begin{equation}
 {1\over2}v^*(K\circ B^{\circ2})v
 <
 \bigl\|K^{\dagger/2}(K\circ B)v\bigr\|^2,        \tag{32.15}
\end{equation}
then \(L_B(t)\) is not positive semidefinite for all sufficiently
small \(t>0\).
\end{proposition}

\begin{proof}
Because \(v\in\ker K\) and \(v\in\ker A_B\), it lies in the
second unresolved subspace.  The difference between the two sides
of (32.15) is
\[
 v^*E_Bv.
\]
It is strictly negative, contradicting the necessary second-order
condition in Corollary \ref{cor:v12v32-Ward-OS}.
\end{proof}

This witness requires only one certified null vector and scalar
interval bounds for its two sides.  It can therefore reject a
candidate before the full quadratic feasibility problem is solved.

\subsection{Why exact reflection splitting escapes strictness}

For the reflection-split primitive \(B={\cal S}b\) of V31,
\[
 L_B(t)=U(t)^*KU(t)\succeq0
\]
for every \(t\).  Its rank is constant, so the evolving kernel need
not open strictly at either first or second order.  Theorem
\ref{thm:v12v32-second-order} implies in this case
\[
 A_B\succeq0,\qquad P_2E_BP_2\succeq0,
\]
possibly with equality.  This is not a defect: exact congruence
factorization proves positivity without a spectral opening.

\begin{remark}[Two valid routes]
\label{rem:v12v32-two-routes}
A candidate Ward primitive may therefore pass the OS gate in either
of two ways:
\begin{enumerate}
\item an exact factorization such as the reflection-split congruence
of V31; or
\item a finite-order spectral opening, first by \(A_B\succ0\), or,
when \(A_B\) is singular, by strict positivity of \(P_2E_BP_2\).
\end{enumerate}
Failure of strict opening is not failure when an exact factorization
is available.
\end{remark}

\subsection{Updated finite card}

The finite RPESM counterterm test now has the following order:

\begin{enumerate}
\item solve the reflection-split Ward equation, which gives exact
positivity if successful;
\item otherwise solve the V31 first-order semidefinite feasibility
problem;
\item for every solution on a singular face, compute \(K^\dagger\),
\(P_2\), and the full effective coefficient \(E_B\);
\item accept a local spectral route only with a strict certified
second-order margin, or continue the Schur hierarchy on its remaining
kernel;
\item reject a candidate immediately when (32.15) has a certified
strict sign.
\end{enumerate}

\begin{assessment}[Status after V32]
\label{ass:v12v32-status}
The boundary case left by V31 has a complete second-order test.
The correct coefficient is the raw Hadamard second derivative minus
the positive-subspace mixing square.  Strict positivity on the
remaining null space yields a genuine positive path; a single
negative scalar witness rejects it.  Exact reflection splitting
remains the preferred route because it requires no spectral opening.
\end{assessment}

\begin{firewall}[Claim boundary after V32]
\label{fire:v12v32-claim}
No physical RPESM matrix has been tested.  If both first- and
second-order effective coefficients remain singular, higher-order
Schur reduction is still required unless an exact factorization is
found.  The note does not prove nonlinear Ward cancellation,
cofinal uniformity, determinant-line control, the missing Duhamel
fibre, or the full SM--Einstein continuum.
\end{firewall}

\subsection{Parent record}

This V32 note appends to the validated V31 file with record
\[
\begin{array}{c|c}
\text{lines}&9138\\
\text{bytes}&314689\\
\text{SHA--256}&
\texttt{9B72ACDF84971C5DC5DA67E9DD54C902FE38DB98F8B95B4D219B0B37E68B6E92}.
\end{array}
\]

% =============================================================================
% END APPENDED TWELFTH-CYCLE V32 MATERIAL
% =============================================================================
% =============================================================================
% BEGIN APPENDED TWELFTH-CYCLE V33 MATERIAL
% =============================================================================

\section{Twelfth-cycle V33: exact Ward cancellation by a flat
reflection-split congruence}
\label{sec:v12v33}

V31--V32 test an entrywise exponential counterterm perturbatively.
The reflection-split subspace is stronger: its Schur factor is rank
one and acts on the reflected kernel by diagonal congruence, so
positivity is exact at every strength.  The remaining question is
whether such a congruence can make the kernel Ward invariant.  This
note gives the necessary and sufficient local equation and its
integrability obstruction.

\subsection{Ward leaves and reflected kernels}

Let \({\cal O}\) be a connected smooth integral leaf of the
finite-regulator Ward distribution.  Points of \({\cal O}\) differ
only by the selected finite gauge or diffeomorphism transformations.
Let
\[
 K:{\cal O}\longrightarrow{\rm Herm}_d
\]
be a \(C^1\) reflected boundary kernel satisfying
\begin{equation}
                         K(q)\succeq0
 \qquad(q\in{\cal O}).                            \tag{33.1}
\end{equation}
All determinant, ghost, coarea, stabilizer, and boundary factors are
understood to have been included before \(K\) is formed.

Let \(U:{\cal O}\to{\rm GL}_d(\mathbb C)\) be an invertible diagonal
map.  Its reflection-split Schur kernel is
\[
 C_U(q)_{ij}:=\overline{U_{ii}(q)}\,U_{jj}(q),
\]
and the repaired kernel is
\begin{equation}
 \widetilde K
 :=
 K\circ C_U
 =
 U^*KU.                                          \tag{33.2}
\end{equation}
It is positive semidefinite at every \(q\).

For a tangent Ward vector field \(X\) on \({\cal O}\), define the
right logarithmic derivative
\begin{equation}
 A(X):=(XU)U^{-1}.                                \tag{33.3}
\end{equation}
This is a diagonal matrix-valued one-form on \({\cal O}\).

\begin{lemma}[Exact congruence derivative]
\label{lem:v12v33-congruence-derivative}
For every tangent Ward vector field \(X\),
\begin{equation}
 X\widetilde K
 =
 U^*\{XK+A(X)^*K+KA(X)\}U.                       \tag{33.4}
\end{equation}
\end{lemma}

\begin{proof}
Differentiate \(U^*KU\):
\[
 X(U^*KU)
 =
 (XU)^*KU+U^*(XK)U+U^*K(XU).
\]
Since \(XU=A(X)U\), the first term is
\(U^*A(X)^*KU\), and the last is \(U^*KA(X)U\).
Factoring \(U^*\) and \(U\) gives (33.4).
\end{proof}

\subsection{The exact diagonal-congruence equation}

\begin{theorem}[Exact reflection-split Ward repair]
\label{thm:v12v33-exact-repair}
A diagonal congruence \(U\) makes the repaired kernel Ward invariant,
\begin{equation}
                         X\widetilde K=0
 \quad\text{for every tangent Ward field }X,      \tag{33.5}
\end{equation}
if and only if its logarithmic one-form \(A\) satisfies
\begin{equation}
 XK+A(X)^*K+KA(X)=0                              \tag{33.6}
 \label{eq:v12v33-congruence-equation}
\end{equation}
for every \(X\).

Conversely, suppose \({\cal O}\) is simply connected and there is a
\(C^1\) diagonal matrix-valued one-form \(A\) satisfying (33.6) and
\begin{equation}
                         dA=0.                   \tag{33.7}
\end{equation}
Then there is an invertible diagonal \(U\), unique after its value at
one base point is fixed, such that
\[
 (dU)U^{-1}=A.
\]
For this \(U\), the kernel \(U^*KU\) is positive semidefinite and
Ward invariant everywhere on \({\cal O}\).
\end{theorem}

\begin{proof}
The first equivalence follows from Lemma
\ref{lem:v12v33-congruence-derivative} and invertibility of \(U\).

For the converse, write
\[
 A=\operatorname{diag}(\alpha_1,\ldots,\alpha_d).
\]
The diagonal matrices commute, so the flatness equation for the
right logarithmic derivative reduces to \(d\alpha_i=0\) for every
\(i\).  On a simply connected leaf, the closed one-forms are exact:
\(\alpha_i=db_i\).  Fix nonzero constants \(u_i^0\) and put
\[
 U_{ii}(q)=u_i^0\exp\{b_i(q)-b_i(q_0)\}.
\]
Then \(dU\,U^{-1}=A\), and \(U\) is invertible.  Equation (33.6) and
the lemma give Ward invariance.  Positivity follows from
\[
 v^*U^*KUv=(Uv)^*K(Uv)\geq0.
\]
Fixing \(U(q_0)\) removes the integration constants and gives
uniqueness.
\end{proof}

For a nonsimply connected Ward leaf, closedness must be supplemented
by the componentwise trivial-holonomy conditions
\begin{equation}
 \exp\left\{\oint_\gamma\alpha_i\right\}=1
 \quad\text{for every closed loop \(\gamma\) and every \(i\)}.
                                                        \tag{33.8}
\end{equation}
This is a genuine global obstruction.  Local anomaly cancellation
does not imply (33.8).

\subsection{Entrywise finite equations}

Write
\[
 A(X)=\operatorname{diag}(a_1(X),\ldots,a_d(X)).
\]
Equation \eqref{eq:v12v33-congruence-equation} is equivalent to
\begin{equation}
 XK_{ij}
 +
 \{\overline{a_i(X)}+a_j(X)\}K_{ij}=0
 \qquad(1\leq i,j\leq d).                        \tag{33.9}
\end{equation}
On an edge of the support graph where \(K_{ij}\ne0\), this becomes
\begin{equation}
 \overline{a_i(X)}+a_j(X)
 =
 -{XK_{ij}\over K_{ij}}.                         \tag{33.10}
\end{equation}
At a zero entry, division is illegal and the exact necessary
condition is
\begin{equation}
                         K_{ij}=0
 \quad\Longrightarrow\quad XK_{ij}=0.            \tag{33.11}
\end{equation}

\begin{proposition}[Fixed-point linear card]
\label{prop:v12v33-linear-card}
At a fixed \(q\in{\cal O}\) and for a fixed Ward direction \(X\),
the existence of a diagonal infinitesimal congruence is equivalent
to solvability of the real-linear system (33.9) in the real and
imaginary parts of \(a_1,\ldots,a_d\).  Equivalently, if
\[
 {\cal C}_{K(q)}:
 \mathbb C^d\longrightarrow{\rm Herm}_d,\qquad
 {\cal C}_{K(q)}(a)_{ij}
 =
 (\overline a_i+a_j)K_{ij},
\]
then
\begin{equation}
                         -XK(q)\in
 \operatorname{Ran}{\cal C}_{K(q)}.              \tag{33.12}
\end{equation}
It fails exactly when there is a Hermitian left-null witness \(Y\)
such that
\begin{equation}
 {\cal C}_{K(q)}^*Y=0,\qquad
 \langle Y,XK(q)\rangle\ne0.                     \tag{33.13}
\end{equation}
\end{proposition}

\begin{proof}
Equation (33.9) is the coordinate form of
\({\cal C}_{K(q)}(a)=-XK(q)\), proving the first two statements.
In finite-dimensional real inner-product spaces,
\[
 \operatorname{Ran}{\cal C}_{K(q)}
 =
 \bigl(\ker{\cal C}_{K(q)}^*\bigr)^\perp.
\]
Thus membership in the range is equivalent to orthogonality to every
left-null vector.  Its failure is witnessed by (33.13).
\end{proof}

This card is cheaper than the V31 semidefinite program.  It should be
attempted first because a solution that also integrates to a flat
one-form gives positivity for all strengths, including on singular
OS kernels.

\subsection{Cycle and zero-pattern obstructions}

The support graph has vertices \(1,\ldots,d\) and an edge \(ij\)
whenever \(K_{ij}\ne0\).  Equations (33.10) couple the vertex
potentials \(a_i\).  Values prescribed independently on graph edges
need not arise from vertex potentials; their alternating sums around
cycles must obey the corresponding compatibility relations.  The
left-null witnesses in (33.13) are an invariant way to test all such
relations without choosing a spanning tree.

The zero-pattern condition (33.11) is separate.  A diagonal
congruence preserves every zero entry:
\[
 (U^*KU)_{ij}=\overline{U_{ii}}K_{ij}U_{jj}.
\]
Therefore a Ward flow that instantaneously creates a previously zero
entry cannot be made invariant by a reflection-split counterterm.

\begin{corollary}[Rank and inertia are fixed]
\label{cor:v12v33-rank}
Every repaired kernel \(U^*KU\) has the same rank as \(K\), and the
same numbers of positive, negative, and zero eigenvalues.  In
particular, an exact reflection-split repair transports OS null
directions but cannot remove them.
\end{corollary}

\begin{proof}
Congruence by an invertible matrix preserves rank and inertia.
\end{proof}

This explains why the strict spectral-opening alternatives in
V31--V32 and the exact reflection-split route are complementary.  A
strict opening can lift null directions.  A diagonal congruence
preserves them while moving their subspace.

\subsection{Compatibility with projectivity}

At each cutoff \(n\), an integrated solution \(U_n\) produces a
positive, Ward-invariant reflected kernel
\[
                         \widetilde K_n=U_n^*K_nU_n.
\]
This does not automatically make the family projective.  Its
normalized density must still satisfy the V29 conditional
density-martingale equation.  In unnormalized form, exact
marginalization requires
\begin{equation}
 (p_{n+1,n})_!\widetilde K_{n+1}
 =
 \widetilde K_n,                                  \tag{33.14}
\end{equation}
including every coarse conditional normalizer.  Thus the complete
exact route has three independent rows:
\begin{enumerate}
\item solve the pointwise congruence equation (33.6);
\item prove flatness and trivial holonomy of \(A\);
\item prove the projective equation (33.14), or equivalently the
normalized density-martingale relation.
\end{enumerate}

\begin{assessment}[Status after V33]
\label{ass:v12v33-status}
Reflection-split Ward repair now has an exact nonlinear
characterization.  The kernel derivative must be a diagonal
congruence derivative; the resulting diagonal connection must be
flat and have trivial holonomy.  When these conditions hold, one
integrated counterterm makes the complete reflected kernel Ward
invariant while preserving positivity at every strength.
\end{assessment}

\begin{firewall}[Claim boundary after V33]
\label{fire:v12v33-claim}
The selected RPESM kernel and its Ward derivatives have not been
assembled, so the range card (33.12), flatness, holonomy, and
projective equation (33.14) are untested.  A general
reflection-positive repair may exist even when the diagonal
congruence route fails.  No missing Duhamel fibre, determinant-line
trivialization, uniform continuum estimate, or full SM--Einstein
continuum has been constructed.
\end{firewall}

\subsection{Parent record}

This V33 note appends to the validated V32 file with record
\[
\begin{array}{c|c}
\text{lines}&9416\\
\text{bytes}&324332\\
\text{SHA--256}&
\texttt{491B8314859D532591415CF3AD57EBE09D8265EAB098475322A27A1F750AB698}.
\end{array}
\]

% =============================================================================
% END APPENDED TWELFTH-CYCLE V33 MATERIAL
% =============================================================================
% =============================================================================
% BEGIN APPENDED TWELFTH-CYCLE V34 MATERIAL
% =============================================================================

\section{Twelfth-cycle V34: when reflection-split congruence commutes
with marginalization}
\label{sec:v12v34}

V33 constructs an exact positive Ward repair at one cutoff when the
Ward defect is a flat diagonal-congruence cocycle.  A projective
tower requires more: repairing at the fine level and then
marginalizing must agree with repairing the previously defined
coarse kernel.  This note identifies the exact intertwiner that makes
the two operations commute and shows why it is a genuine additional
condition.

\subsection{Fine and coarse reflected spaces}

Let \({\cal H}_{\rm c}\) and \({\cal H}_{\rm f}\) be finite-dimensional
coarse and fine positive-half boundary spaces.  Let
\[
 R:{\cal H}_{\rm c}\longrightarrow{\cal H}_{\rm f}
\]
be the linear lift whose adjoint performs the reflected marginal.
For a fine reflected kernel \(K_{\rm f}\succeq0\), define
\begin{equation}
 K_{\rm c}:=R^*K_{\rm f}R\succeq0.                \tag{34.1}
\end{equation}
Let \(U_{\rm f}\) and \(U_{\rm c}\) be invertible diagonal
counterterm matrices in the chosen fine and coarse boundary bases.

There are two coarse repaired kernels:
\begin{align}
 K_{\rm c}^{({\rm fine})}
 &:=
 R^*U_{\rm f}^*K_{\rm f}U_{\rm f}R,              \tag{34.2}\\
 K_{\rm c}^{({\rm coarse})}
 &:=
 U_{\rm c}^*R^*K_{\rm f}RU_{\rm c}.              \tag{34.3}
\end{align}
Both are positive semidefinite, but they need not be equal.

\begin{theorem}[Exact congruence--marginal intertwiner]
\label{thm:v12v34-intertwiner}
If
\begin{equation}
                         U_{\rm f}R=RU_{\rm c},   \tag{34.4}
\end{equation}
then
\begin{equation}
 K_{\rm c}^{({\rm fine})}
 =
 K_{\rm c}^{({\rm coarse})}                      \tag{34.5}
\end{equation}
for every Hermitian fine kernel \(K_{\rm f}\).  Consequently, if the
unrepaired kernels form an exact projective family and (34.4) holds
at every adjacent scale, the reflection-split repaired kernels also
form an exact projective family.
\end{theorem}

\begin{proof}
Under (34.4),
\[
 R^*U_{\rm f}^*K_{\rm f}U_{\rm f}R
 =
 (U_{\rm f}R)^*K_{\rm f}(U_{\rm f}R)
 =
 (RU_{\rm c})^*K_{\rm f}(RU_{\rm c}),
\]
which is (34.3).  Apply this identity at each adjacent scale and use
associativity of the coarse lifts to obtain exact projectivity.
\end{proof}

A common scalar phase multiplying one side of (34.4) is irrelevant
to (34.5).

\subsection{Universal necessity}

The intertwiner is also necessary if equality is required
independently of the fine kernel.

\begin{proposition}[Universal equality forces an intertwiner]
\label{prop:v12v34-universal-necessity}
Assume every column of \(R\) is nonzero.  If
\begin{equation}
 R^*U_{\rm f}^*HU_{\rm f}R
 =
 U_{\rm c}^*R^*HRU_{\rm c}
 \qquad\text{for every Hermitian }H,              \tag{34.6}
\end{equation}
for every Hermitian \(H\) on \({\cal H}_{\rm f}\), then there is a
scalar phase \(e^{i\theta}\) such that
\begin{equation}
 U_{\rm f}R=e^{i\theta}RU_{\rm c}.                \tag{34.7}
\end{equation}
\end{proposition}

\begin{proof}
Put \(A=U_{\rm f}R\) and \(B=RU_{\rm c}\), and denote their columns
by \(a_j,b_j\).  Equation (34.6) says
\[
 a_i^*Ha_j=b_i^*Hb_j
\]
for every Hermitian \(H\) and every \(i,j\).  Hermitian matrices
separate complex matrices under the trace pairing, so
\[
 a_ja_i^*=b_jb_i^*
 \qquad\text{for all }i,j.                       \tag{34.8}
\]
Taking \(i=j\) and using nonzero columns gives
\(a_i=e^{i\theta_i}b_i\).  Substitution into (34.8) for \(i\ne j\)
shows \(e^{i(\theta_j-\theta_i)}=1\).  All phases are equal, proving
(34.7).
\end{proof}


\subsection{Fibrewise meaning for an incidence lift}

Suppose the lift is an incidence matrix: each fine boundary label
\(\alpha\) belongs to one coarse label \(a=q(\alpha)\), and
\(R_{\alpha a}\ne0\) only when \(a=q(\alpha)\).  Write
\[
 (U_{\rm f})_{\alpha\alpha}=u_{\rm f}(\alpha),
 \qquad
 (U_{\rm c})_{aa}=u_{\rm c}(a).
\]

\begin{corollary}[Fibre-constancy criterion]
\label{cor:v12v34-fibre-constant}
After absorbing one global phase, the intertwiner (34.4) holds if
and only if
\begin{equation}
 u_{\rm f}(\alpha)=u_{\rm c}(q(\alpha))
 \quad\text{whenever }R_{\alpha,q(\alpha)}\ne0.   \tag{34.9}
\end{equation}
Thus the fine reflection-split factor must be constant on every
coarse marginalization fibre seen by \(R\).
\end{corollary}

\begin{proof}
The \((\alpha,a)\) entry of (34.4) is
\[
 u_{\rm f}(\alpha)R_{\alpha a}
 =
 R_{\alpha a}u_{\rm c}(a).
\]
For a nonzero incidence entry this is exactly (34.9); all other
entries vanish on both sides.
\end{proof}

A Ward primitive may therefore be reflection split at every cutoff
and still fail to define a projective reflection-split family because
its half-space potential varies inside an eliminated fibre.

\subsection{Counterexample: positivity survives but the ansatz does
not}

Let
\[
 K_{\rm f}=
 \begin{pmatrix}
 1&0&1\\
 0&1&1\\
 1&1&2
 \end{pmatrix},
 \qquad
 R=
 \begin{pmatrix}
 1&0\\
 1&0\\
 0&1
 \end{pmatrix},
 \qquad
 U_{\rm f}=\operatorname{diag}(1,2,1).            \tag{34.10}
\]
The fine kernel is the Gram matrix of
\[
 e_1,\quad e_2,\quad e_1+e_2
\]
and is therefore positive semidefinite.  Direct calculation gives
\begin{align}
 K_{\rm c}
 &=
 R^*K_{\rm f}R
 =
 \begin{pmatrix}2&2\\2&2\end{pmatrix},            \tag{34.11}\\
 K_{\rm c}^{({\rm fine})}
 &=
 R^*U_{\rm f}^*K_{\rm f}U_{\rm f}R
 =
 \begin{pmatrix}5&3\\3&2\end{pmatrix}.            \tag{34.12}
\end{align}
The first has rank one, while the second has determinant one and
rank two.

\begin{proposition}[Marginalization can leave the diagonal ansatz]
\label{prop:v12v34-counterexample}
For the data (34.10), there is no invertible diagonal \(U_{\rm c}\)
such that
\[
 K_{\rm c}^{({\rm fine})}=U_{\rm c}^*K_{\rm c}U_{\rm c}.
\]
Nevertheless both sides of (34.2) are positive semidefinite.
\end{proposition}

\begin{proof}
Congruence by an invertible matrix preserves rank.  Every
\(U_{\rm c}^*K_{\rm c}U_{\rm c}\) has rank one, whereas (34.12) has
rank two.  Positivity of (34.12) follows either from its positive
leading minor and determinant, or directly from its form
\(R^*U_{\rm f}^*K_{\rm f}U_{\rm f}R\).
\end{proof}

The failure is caused by the two fine labels in the first coarse
fibre receiving unequal factors \(1\) and \(2\).  Recomputing the
coarse kernel remains safe, as V28 requires, but the result cannot
be encoded by a diagonal congruence of the historical coarse kernel.

\subsection{Approximate intertwining budget}

Exact fibre constancy may be stronger than necessary if one uses the
summable projective-repair theorem.  Put
\[
 A:=U_{\rm f}R,\qquad B:=RU_{\rm c},\qquad
 E:=A-B.
\]

\begin{proposition}[Quantitative marginal defect]
\label{prop:v12v34-defect}
For every fine Hermitian kernel \(K_{\rm f}\),
\begin{equation}
 \left\|
 A^*K_{\rm f}A-B^*K_{\rm f}B
 \right\|
 \leq
 \|K_{\rm f}\|
 \left(2\|B\|\|E\|+\|E\|^2\right).               \tag{34.13}
\end{equation}
The same estimate holds with any submultiplicative operator norm
compatible with adjoints.
\end{proposition}

\begin{proof}
Since \(A=B+E\),
\[
 A^*K_{\rm f}A-B^*K_{\rm f}B
 =
 E^*K_{\rm f}B+B^*K_{\rm f}E+E^*K_{\rm f}E.
\]
Submultiplicativity and the triangle inequality give (34.13).
\end{proof}

This is a kernel-norm estimate, not yet a probability
total-variation estimate.  To feed V79 or the cylinder-local V83
repair theorem, it must be combined with a lower normalization
anchor and an explicit comparison between kernel norm and the
selected cylinder-measure norm.

\subsection{Projective Ward connection card}

For adjacent cutoffs, an exact reflection-split projective repair
must now satisfy all of the following:

\begin{enumerate}
\item the V33 diagonal-congruence Ward equation on both scales;
\item flatness and trivial holonomy of both diagonal Ward
connections;
\item the interscale intertwiner \(U_{\rm f}R=RU_{\rm c}\);
\item the V29 conditional-normalizer equation after all eliminated
variables are included.
\end{enumerate}

If the third row fails, one may still push forward the complete
positive repaired fine kernel.  The coarse result stays positive but
must be treated as a newly recomputed Wilsonian kernel, potentially
outside the diagonal-congruence ansatz.

\begin{assessment}[Status after V34]
\label{ass:v12v34-status}
The exact reflection-split Ward construction now has its interscale
compatibility condition.  Diagonal repair commutes with
marginalization when its half-space factor intertwines the lift,
which for incidence maps means fibre constancy.  Without that
condition positivity survives exact pushforward, but even the rank
of the coarse kernel can differ from every diagonal repair of the
historical coarse kernel.  An explicit norm bound quantifies an
approximate intertwiner.
\end{assessment}

\begin{firewall}[Claim boundary after V34]
\label{fire:v12v34-claim}
No RPESM lift \(R\) or diagonal Ward connection has been populated,
and fibre constancy or summable intertwining defect has not been
proved.  Kernel-norm control alone does not give normalized measure
control.  The missing Duhamel fibre, determinant-line
trivialization, nonlinear gravitational estimates, and the full
SM--Einstein continuum remain open.
\end{firewall}

\subsection{Parent record}

This V34 note appends to the validated V33 file with record
\[
\begin{array}{c|c}
\text{lines}&9723\\
\text{bytes}&334792\\
\text{SHA--256}&
\texttt{7AF98084BFAD5BCC45F6EC565CB3349E86ACC71032007637DD8A61B90A719B99}.
\end{array}
\]

% =============================================================================
% END APPENDED TWELFTH-CYCLE V34 MATERIAL
% =============================================================================
% =============================================================================
% BEGIN APPENDED TWELFTH-CYCLE V35 MATERIAL
% =============================================================================

\section{Twelfth-cycle V35: compulsory companion audit and the
local-certificate correction}
\label{sec:v12v35}

\subsection{Audit identity}

This is the compulsory five-version audit following V30.  The newest
active companion files available at the audit time were
\[
\begin{array}{c|r|r}
\text{file}&\text{counted lines}&\text{bytes}\\ \hline
\texttt{Renewal\_Yang\_Mills\_Mass\_Gap\_Research\_Notes\_V52.tex}
 &20062&687538\\
\texttt{Renewal\_NS\_Stability\_Research\_Notes\_V26.tex}
 &5320&278283.
\end{array}                                                   \tag{35.1}
\]
Their SHA--256 fingerprints are
\[
\begin{split}
 &\texttt{301F7BB8BD639CE2B0317437CD35ED0588267CD6B91F1A1FDB539250A2243108},
 \\
 &\texttt{82C887F8C2C69E4F28FAD7C3DC8DE5B9099CB5A4BF431EBA1FFF3E91935978AD}.
\end{split}                                                   \tag{35.2}
\]
Both files were read without modification.

\subsection{Yang--Mills V50: independent confirmation of legal
momentum packets}

Yang--Mills V50 imports the SM V26 shift-subgroup theorem and
specializes it to one axial external momentum on
\((\mathbb Z/M\mathbb Z)^4\).  Its minimal invariant packets contain
all \(M\) values of the shifted coordinate, and all previously used
two-card spectator pairs cross those packets.  This independently
confirms the correction already recorded in SM V25--V26: covariance
belongs to the full shift-invariant trace packet, not to an
individual momentum card.

No additional theorem is imported here.  The useful evidence is that
the packet correction changes an actual downstream calculation and
is not merely formal bookkeeping.  Accordingly every future
SM--Einstein Duhamel or Ward matrix must retain the subgroup generated
by all background shifts before any trace, positivity, or anomaly
test is performed.

\subsection{Yang--Mills V51--V52: exact rays and the multivariable
boundary}

The Yang--Mills programme next restricts its exact Laurent Hessians
to a complex coordinate ray.  For the endpoint and Wilson actions it
proves
\begin{equation}
 \det A_e(e^{i\zeta},1,1,1)\ne0
 \quad\text{when}\quad|\operatorname{Im}\zeta|\leq\log2,
 \qquad e\in\{{\rm W},*\}.                         \tag{35.3}
\end{equation}
The endpoint proof uses exact interpolation and Rouch\'e dominance.
The Wilson proof obtains a reciprocal polynomial and factors it
exactly.  These results improve the former common axis width by more
than six orders of magnitude.

The companion note also states the decisive limitation: a
one-coordinate annulus does not imply a multivariable polystrip.
The planned next step is a cover of the real torus by boxes, with an
inverse residual or smallest-singular-value lower bound on every
box.

\begin{proposition}[Restriction certificates do not tensorize]
\label{prop:v12v35-no-tensorization}
Let \(P(z_1,z_2)\) be a polynomial.  Nonvanishing of
\(P(z,1)\) and \(P(1,z)\) on a domain \(D\) does not imply
nonvanishing of \(P\) on \(D^2\).
\end{proposition}

\begin{proof}
Take
\[
 P(z_1,z_2)=z_1+z_2-3.
\]
On the closed unit disk,
\(P(z,1)=z-2\) and \(P(1,z)=z-2\) do not vanish.  But
\(P(3/2,3/2)=0\).  Choosing any domain containing the unit disk and
\(3/2\) gives the stated failure.  More generally, restrictions to
coordinate slices do not detect mixed-variable zero sets.
\end{proof}

The particular example is only logical; it is not a claim about the
Yang--Mills or RPESM determinant.

\subsection{Transfer to the SM--Einstein programme}

The transferable method has two parts.

First, global coefficient-sum perturbation can be far weaker than
the actual inverse geometry.  A local cover may certify a useful
spectral gap even when one perturbation from a single reference
point is useless.  For a matrix family \(A(x)\), a box \(Q\) with
centre \(x_Q\) is certified whenever
\begin{equation}
 \|A(x_Q)^{-1}\|\,
 \sup_{x\in Q}\|A(x)-A(x_Q)\|<1.                \tag{35.4}
\end{equation}
The Neumann series then proves invertibility throughout \(Q\) and
gives the explicit bound
\begin{equation}
 \sup_{x\in Q}\|A(x)^{-1}\|
 \leq
 {\|A(x_Q)^{-1}\|\over
  1-\|A(x_Q)^{-1}\|
       \sup_{x\in Q}\|A(x)-A(x_Q)\|}.             \tag{35.5}
\end{equation}

Second, the certificate must cover the complete multivariable
domain.  Axis restrictions are useful regression faces and zero
locators, but they do not certify the physical configuration box or
the full Duhamel polystrip.

For V34, this changes the quantitative route from one global
normalizer floor to a finite or countable certified cover.  On each
box one may combine:
\begin{enumerate}
\item a local lower trace or smallest-eigenvalue anchor for the
positive reflected kernel;
\item the V34 approximate-intertwining defect;
\item overlap agreement of normalized coarse states.
\end{enumerate}
A uniform or summably weighted collection of such local estimates
can then feed projective repair.  None of the numerical Yang--Mills
widths transfers to the different SM--Einstein matrices.

\subsection{Navier--Stokes status}

Navier--Stokes V26 remains unchanged.  Its exact rank-one contraction
of the Oseen derivative continues to support the V30 rule that a
structured norm may be used only on a proved structured input cone.
It supplies no determinant, reflected-kernel, conditional-normalizer,
or projective estimate for the present SM--Einstein step.

\subsection{Post-audit acquisition}

V34 controls the difference between two unnormalized coarse
reflected kernels in operator norm.  That estimate is not yet a
distance between normalized states.  The next note will prove the
missing bridge: a lower trace anchor converts a trace-norm kernel
defect into a trace-distance bound for normalized positive states.
Combined with the V34 intertwining estimate and a dimension or local
rank bound, it will produce an explicit adjacent projective error
that can be summed by the earlier repair theorem.

\begin{assessment}[Status after V35]
\label{ass:v12v35-status}
The audit confirms the legal momentum-packet correction and imports
a method, rather than foreign numerical constants: replace crude
global inverse bounds by a fully covering family of certified local
boxes.  One-dimensional exact restrictions remain insufficient.
The immediate SM task is now the precise conversion from positive
kernel error to normalized projective-state error.
\end{assessment}

\begin{firewall}[Claim boundary after V35]
\label{fire:v12v35-claim}
No Yang--Mills polystrip, continuum response, or mass gap is imported.
No Navier--Stokes regularity statement is imported.  A local-box
cover, a physical RPESM lower normalizer, the missing Duhamel fibre,
and the full SM--Einstein continuum remain unconstructed.
\end{firewall}

\subsection{Parent record}

This V35 note appends to the validated V34 file with record
\[
\begin{array}{c|c}
\text{lines}&10037\\
\text{bytes}&344610\\
\text{SHA--256}&
\texttt{E862A21D88B13AA03A79FEEC4E4723E0A05D071722A9E30D24703E5C26F6A357}.
\end{array}
\]

% =============================================================================
% END APPENDED TWELFTH-CYCLE V35 MATERIAL
% =============================================================================
% =============================================================================
% BEGIN APPENDED TWELFTH-CYCLE V36 MATERIAL
% =============================================================================

\section{Twelfth-cycle V36: from positive-kernel defects to
normalized projective-state errors}
\label{sec:v12v36}

V34 bounds the operator-norm defect produced when a
reflection-split fine counterterm fails to intertwine the coarse
lift.  Projective repair, however, requires a bound after
normalization.  This note supplies that conversion first for
trace-normalized positive kernels and then for a general positive
partition functional.

\subsection{Trace-normalized positive kernels}

For a nonzero positive semidefinite matrix \(A\), put
\[
 \rho_A:={A\over\operatorname{Tr}A}.
\]
For two density matrices define their trace distance by
\[
 d_{\rm tr}(\rho,\sigma):={1\over2}\|\rho-\sigma\|_1.
\]

\begin{theorem}[Normalization stability in trace norm]
\label{thm:v12v36-trace-normalization}
Let \(A,B\succeq0\) be nonzero and let
\(\Delta=\|A-B\|_1\).  Then
\begin{equation}
 d_{\rm tr}(\rho_A,\rho_B)
 \leq
 {\Delta\over\operatorname{Tr}A}.                \tag{36.1}
\end{equation}
The same assertion with \(A\) and \(B\) exchanged also holds.
Consequently, if \(\operatorname{Tr}A\geq\zeta>0\), then
\begin{equation}
 d_{\rm tr}(\rho_A,\rho_B)\leq{\Delta\over\zeta}.
                                                               \tag{36.2}
\end{equation}
\end{theorem}

\begin{proof}
Write \(\alpha=\operatorname{Tr}A\) and
\(\beta=\operatorname{Tr}B\).  Positivity gives
\(\|B\|_1=\beta\), while trace duality gives
\[
 |\alpha-\beta|
 =
 |\operatorname{Tr}(A-B)|
 \leq\Delta.
\]
Now
\begin{align*}
 \left\|{A\over\alpha}-{B\over\beta}\right\|_1
 &\leq
 {1\over\alpha}\|A-B\|_1
 +
 \|B\|_1\left|{1\over\alpha}-{1\over\beta}\right|\\
 &=
 {\Delta\over\alpha}
 +
 {|\alpha-\beta|\over\alpha}
 \leq {2\Delta\over\alpha}.
\end{align*}
Division by two proves (36.1).  Interchanging the matrices proves the
symmetric statement, and the trace floor gives (36.2).
\end{proof}

For every bounded boundary observable \(O\),
\begin{equation}
 \left|
 \operatorname{Tr}\{O(\rho_A-\rho_B)\}
 \right|
 \leq
 2\|O\|\,d_{\rm tr}(\rho_A,\rho_B).               \tag{36.3}
\end{equation}
Thus the estimate controls all normalized finite boundary
expectations at once.

\subsection{Application to the V34 intertwining defect}

Retain the V34 notation and put
\[
 F:=U_{\rm f}R,\qquad
 G:=RU_{\rm c},\qquad
 E:=F-G.
\]
Define
\begin{align}
 A&:=G^*K_{\rm f}G,                              \tag{36.4}\\
 B&:=F^*K_{\rm f}F.                              \tag{36.5}
\end{align}
Both are positive semidefinite on the \(d_{\rm c}\)-dimensional
coarse boundary space.

\begin{corollary}[Normalized approximate-intertwining bound]
\label{cor:v12v36-intertwining}
Assume
\[
 \operatorname{Tr}(G^*K_{\rm f}G)\geq\zeta>0.
\]
Then
\begin{equation}
 d_{\rm tr}(\rho_A,\rho_B)
 \leq
 {d_{\rm c}\|K_{\rm f}\|\over\zeta}
 \left(2\|G\|\|E\|+\|E\|^2\right).               \tag{36.6}
\end{equation}
\end{corollary}

\begin{proof}
Proposition \ref{prop:v12v34-defect} gives
\[
 \|A-B\|
 \leq
 \|K_{\rm f}\|
 \left(2\|G\|\|E\|+\|E\|^2\right).
\]
For a matrix on a \(d_{\rm c}\)-dimensional space,
\(\|A-B\|_1\leq d_{\rm c}\|A-B\|\).
Insert this inequality into Theorem
\ref{thm:v12v36-trace-normalization}.
\end{proof}

If a rank bound
\(\operatorname{rank}(A-B)\leq r_{\rm def}\) is known, the factor
\(d_{\rm c}\) in (36.6) may be replaced by \(r_{\rm def}\).  Such a
replacement requires an actual rank certificate; locality alone
does not imply it.

\subsection{A general partition functional}

The physical cylinder normalizer need not be the matrix trace.  Let
\(\ell:{\rm Herm}_{d}\to\mathbb R\) be a positive linear functional,
and denote its dual norm relative to trace norm by
\[
 c_\ell:=\sup_{\|H\|_1\leq1}|\ell(H)|.
\]
For positive kernels \(A,B\), put
\[
 Z_A:=\ell(A),\qquad Z_B:=\ell(B).
\]

\begin{theorem}[Normalizer-aware kernel comparison]
\label{thm:v12v36-general-normalizer}
Assume \(A,B\succeq0\),
\[
 Z_A\geq\zeta>0,\qquad Z_B>0,\qquad
 {\operatorname{Tr}B\over Z_B}\leq M.
\]
Then, with \(\Delta=\|A-B\|_1\),
\begin{equation}
 \left\|
 {A\over Z_A}-{B\over Z_B}
 \right\|_1
 \leq
 {1+Mc_\ell\over\zeta}\,\Delta.                  \tag{36.7}
\end{equation}
\end{theorem}

\begin{proof}
Linearity and duality give
\[
 |Z_A-Z_B|=|\ell(A-B)|\leq c_\ell\Delta.
\]
Using \(B\) in the normalization-difference term,
\begin{align*}
 \left\|
 {A\over Z_A}-{B\over Z_B}
 \right\|_1
 &\leq
 {\Delta\over Z_A}
 +
 \operatorname{Tr}B\,
 { |Z_A-Z_B|\over Z_AZ_B}\\
 &\leq
 {\Delta\over\zeta}
 +
 {M c_\ell\Delta\over\zeta},
\end{align*}
which is (36.7).
\end{proof}

For a vector partition functional
\(\ell_w(H)=w^*Hw\), one has
\(c_{\ell_w}=\|w\|^2\).  The additional stock
\(\operatorname{Tr}B/Z_B\) measures how much positive kernel mass is
invisible to the chosen partition vector.  A lower value of \(Z_B\)
alone does not control that ratio.

\begin{warning}[The normalizer must match the physical cylinder]
\label{warn:v12v36-physical-normalizer}
Trace normalization is sufficient for a finite boundary density
matrix, but it may differ from the scalar normalization of the
original reflected cylinder.  Equation (36.7), with the actual
positive functional \(\ell\), is the applicable estimate in that
case.  Replacing \(\ell\) by the trace without proving equality would
change the theory.
\end{warning}

\subsection{Local boxes and integrated state errors}

Let \(q\) range over a measurable parameter region and suppose each
\(q\) has positive kernels \(A(q),B(q)\).  If a box cover proves
\[
 Z_A(q)\geq\zeta_Q,\qquad
 {\operatorname{Tr}B(q)\over Z_B(q)}\leq M_Q,\qquad
 \|A(q)-B(q)\|_1\leq\Delta_Q
 \quad(q\in Q),
\]
then (36.7) gives the boxwise error
\begin{equation}
 \varepsilon_Q
 :=
 {1+M_Qc_\ell\over\zeta_Q}\Delta_Q.              \tag{36.8}
\end{equation}
On overlaps the normalized kernels are the same pointwise objects,
so no partition-of-unity modification is required.

If \(\lambda\) is a probability law for \(q\), convexity of trace
norm gives
\begin{equation}
 \left\|
 \int {A(q)\over Z_A(q)}\,d\lambda(q)
 -
 \int {B(q)\over Z_B(q)}\,d\lambda(q)
 \right\|_1
 \leq
 \int\varepsilon(q)\,d\lambda(q),                \tag{36.9}
\end{equation}
where \(\varepsilon(q)\) is any measurable local bound from (36.7).

\subsection{Summable projective consequence}

Let \(\Phi_{n+1,n}\) be the normalized coarse channel between finite
boundary state spaces, and assume it contracts trace distance.  This
holds for completely positive trace-preserving maps and for
classical pushforwards.  Suppose positive normalized states
\(\rho_n\) satisfy
\begin{equation}
 d_{\rm tr}\bigl(
 \Phi_{n+1,n}(\rho_{n+1}),\rho_n
 \bigr)\leq\varepsilon_n,
 \qquad
 \sum_n\varepsilon_n<\infty.                    \tag{36.10}
\end{equation}

\begin{corollary}[Trace-state projective repair]
\label{cor:v12v36-projective-repair}
For every fixed \(k\), the descended states
\[
 \Phi_{N,k}(\rho_N),\qquad N\geq k,
\]
converge in trace norm to a state \(\bar\rho_k\).  The limit states
are exactly compatible, and
\begin{equation}
 d_{\rm tr}(\bar\rho_k,\rho_k)
 \leq
 \sum_{n=k}^\infty\varepsilon_n.                 \tag{36.11}
\end{equation}
\end{corollary}

\begin{proof}
Trace-distance contraction gives
\[
 d_{\rm tr}\bigl(
 \Phi_{N+1,k}(\rho_{N+1}),
 \Phi_{N,k}(\rho_N)
 \bigr)
 \leq\varepsilon_N.
\]
The summable tail makes the sequence Cauchy in the complete
finite-dimensional trace class and proves (36.11) by telescoping.
Positivity and unit trace pass to the limit.  Continuity of
\(\Phi_{k+1,k}\) gives exact compatibility.
\end{proof}

This is the finite positive-kernel specialization of the earlier
summable projective-repair mechanism.  Its new input is the explicit
error supplied by (36.6) or (36.7).

\subsection{Acquisition card}

For each adjacent RPESM scale one must now populate:

\begin{enumerate}
\item the complete positive fine kernel \(K_{{\rm f},n}\);
\item the exact or approximate intertwining defect
\(E_n=U_{{\rm f},n}R_n-R_nU_{{\rm c},n}\);
\item the actual physical normalizer functional \(\ell_n\);
\item local certified floors \(\zeta_{n,Q}\) and visibility stocks
\(M_{n,Q}\);
\item a dimension, rank, or direct trace-norm bound for the kernel
defect; and
\item a summable integrated or cylinder-local sequence
\(\varepsilon_n\).
\end{enumerate}

\begin{assessment}[Status after V36]
\label{ass:v12v36-status}
The gap between V34 kernel control and normalized projective control
is closed conditionally and quantitatively.  A trace floor converts
the approximate intertwiner into trace distance.  For the physical
partition functional, one additional visibility ratio is necessary.
Summability of the resulting adjacent errors produces an exactly
compatible limiting family of finite boundary states.
\end{assessment}

\begin{firewall}[Claim boundary after V36]
\label{fire:v12v36-claim}
The RPESM normalizer functional, local floors, visibility stocks, and
intertwining defects have not been populated.  Trace-state
projectivity does not by itself establish a Radon field measure,
tightness, clustering, determinant-line trivialization, the missing
Duhamel fibre, or the full SM--Einstein continuum.
\end{firewall}

\subsection{Parent record}

This V36 note appends to the validated V35 file with record
\[
\begin{array}{c|c}
\text{lines}&10224\\
\text{bytes}&351968\\
\text{SHA--256}&
\texttt{2A90E8D135128AEF84CC39633868A4D2FF66AD306A13D707F0AFA5E5283141FA}.
\end{array}
\]

% =============================================================================
% END APPENDED TWELFTH-CYCLE V36 MATERIAL
% =============================================================================
% =============================================================================
% BEGIN APPENDED TWELFTH-CYCLE V37 MATERIAL
% =============================================================================

\section{Twelfth-cycle V37: local box certificates for physical
normalizers and visibility}
\label{sec:v12v37}

V36 reduces normalized projective control to three local quantities:
a lower bound for the physical normalizer, an upper visibility ratio,
and a trace-norm kernel defect.  This note shows how to certify all
three from data at the centre of a parameter box plus operator-norm
variation bounds.  It also records the tail term required when the
certified boxes cover only a good set.

\subsection{Positive representation of the normalizer}

Every positive linear functional on \({\rm Herm}_d\) has the form
\begin{equation}
 \ell(H)=\operatorname{Tr}(SH),\qquad S\succeq0.             \tag{37.1}
\end{equation}
Two norms of \(S\) enter:
\[
 c_\ell:=\|S\|,
 \qquad
 s_\ell:=\operatorname{Tr}S.
\]
The first is the dual norm of \(\ell\) relative to trace norm; the
second is the dual norm relative to operator norm:
\begin{align}
 |\ell(H)|&\leq c_\ell\|H\|_1,                              \tag{37.2}\\
 |\ell(H)|&\leq s_\ell\|H\|.                                \tag{37.3}
\end{align}
For \(\ell(H)=\operatorname{Tr}H\),
\[
 c_\ell=1,\qquad s_\ell=d.
\]
For \(\ell_w(H)=w^*Hw\), both constants equal \(\|w\|^2\).

\subsection{One certified box}

Let \(Q\) be a parameter box with centre \(q_0\).  Let
\(A(q),B(q)\succeq0\) be the two positive coarse kernels to be
compared.  Suppose interval or analytic estimates prove
\begin{align}
 \sup_{q\in Q}\|A(q)-A(q_0)\|&\leq\eta_A,          \tag{37.4}\\
 \sup_{q\in Q}\|B(q)-B(q_0)\|&\leq\eta_B,          \tag{37.5}\\
 \sup_{q\in Q}\|A(q)-B(q)\|&\leq\delta.            \tag{37.6}
\end{align}
At the centre set
\[
 a_0:=\ell(A(q_0)),\qquad
 b_0:=\ell(B(q_0)),\qquad
 t_0:=\operatorname{Tr}B(q_0).
\]

\begin{theorem}[Box-center normalizer certificate]
\label{thm:v12v37-box}
Define
\begin{align}
 \zeta_Q&:=a_0-s_\ell\eta_A,                     \tag{37.7}\\
 \beta_Q&:=b_0-s_\ell\eta_B,                     \tag{37.8}\\
 T_Q&:=t_0+d\eta_B.                              \tag{37.9}
\end{align}
If \(\zeta_Q>0\) and \(\beta_Q>0\), then for every \(q\in Q\),
\begin{align}
 \ell(A(q))&\geq\zeta_Q,                         \tag{37.10}\\
 \ell(B(q))&\geq\beta_Q,                         \tag{37.11}\\
 {\operatorname{Tr}B(q)\over\ell(B(q))}
 &\leq {T_Q\over\beta_Q},                        \tag{37.12}\\
 \|A(q)-B(q)\|_1&\leq d\delta.                   \tag{37.13}
\end{align}
Consequently,
\begin{equation}
 \left\|
 {A(q)\over\ell(A(q))}
 -
 {B(q)\over\ell(B(q))}
 \right\|_1
 \leq
 \boxed{
 {d\delta\over\zeta_Q}
 \left(1+c_\ell{T_Q\over\beta_Q}\right)}
 =:\varepsilon_Q.                                \tag{37.14}
\end{equation}
\end{theorem}

\begin{proof}
By (37.3) and (37.4),
\[
 \ell(A(q))
 \geq
 \ell(A(q_0))-s_\ell\|A(q)-A(q_0)\|
 \geq\zeta_Q.
\]
The same argument proves (37.11).  Trace duality with the identity
matrix gives
\[
 \operatorname{Tr}B(q)
 \leq
 \operatorname{Tr}B(q_0)+d\|B(q)-B(q_0)\|
 \leq T_Q.
\]
Division by the positive lower bound \(\beta_Q\) proves (37.12).
The dimension inequality
\(\|H\|_1\leq d\|H\|\) gives (37.13).  Insert
(37.10)--(37.13) into Theorem
\ref{thm:v12v36-general-normalizer} to obtain (37.14).
\end{proof}

The box can be certified without a positive lower bound for every
eigenvalue.  A positive physical normalizer at the centre and a
sufficiently small operator variation are enough.  A smallest
eigenvalue bound may still be preferable because it can control
additional inverse and locality estimates.

\begin{corollary}[Rank-refined box error]
\label{cor:v12v37-rank}
If
\[
 \operatorname{rank}\{A(q)-B(q)\}\leq r
 \quad(q\in Q),
\]
then the factor \(d\delta\) in (37.14) may be replaced by
\(r\delta\).
\end{corollary}

\begin{proof}
For a Hermitian matrix of rank at most \(r\),
\[
 \|H\|_1\leq r\|H\|.
\]
Use this in place of (37.13).
\end{proof}

\subsection{Insertion of the V34 defect}

For the approximate congruence intertwiner, take
\[
 A=G^*K_{\rm f}G,\qquad B=F^*K_{\rm f}F,
 \qquad
 F=U_{\rm f}R,\quad G=RU_{\rm c}.
\]
Proposition \ref{prop:v12v34-defect} gives the pointwise choice
\begin{equation}
 \delta_Q
 =
 \sup_{q\in Q}
 \|K_{\rm f}(q)\|
 \left[
 2\|G(q)\|\|F(q)-G(q)\|
 +\|F(q)-G(q)\|^2
 \right].                                         \tag{37.15}
\end{equation}
Any certified upper bound for the right side may be inserted as
\(\delta\) in Theorem \ref{thm:v12v37-box}.  This produces a complete
boxwise normalized projective error using only:
\begin{enumerate}
\item centre values of two positive coarse kernels;
\item their variation from the centre;
\item the interscale congruence defect; and
\item the chosen physical normalizer matrix \(S\).
\end{enumerate}

\subsection{Subdivision terminates under a strict local margin}

A local certificate is useful only if refinement can eventually
verify it.

\begin{proposition}[Termination on a compact strict-margin set]
\label{prop:v12v37-termination}
Let \(G\) be compact.  Suppose \(A,B\) are uniformly continuous on
\(G\), and
\begin{equation}
 \inf_{q\in G}\ell(A(q))>0,\qquad
 \inf_{q\in G}\ell(B(q))>0.                       \tag{37.16}
\end{equation}
Then every sufficiently fine finite box cover of \(G\) has
\(\zeta_Q>0\) and \(\beta_Q>0\) on every box.  If the centre
quantities are enclosed rigorously and the variation moduli tend to
zero with box diameter, recursive bisection terminates for the two
normalizer tests.
\end{proposition}

\begin{proof}
Let the two infima in (37.16) be \(m_A,m_B>0\).
Uniform continuity gives a scale at which
\[
 s_\ell\eta_A<m_A/2,\qquad s_\ell\eta_B<m_B/2.
\]
Every centre belongs to \(G\), so \(a_0\geq m_A\) and
\(b_0\geq m_B\).  Hence
\(\zeta_Q\geq m_A/2\) and \(\beta_Q\geq m_B/2\).
Compactness gives a finite cover at that scale.  A bisection scheme
whose diameters tend to zero eventually reaches it.
\end{proof}

This proposition is conditional on the strict positive margins.
Subdivision cannot manufacture a margin when the physical
normalizer actually approaches zero.

\subsection{Certified good set plus uncovered tail}

A noncompact or singular configuration region may not admit one
finite strict-margin cover.  Let \(G\) be the union of certified
boxes and \(T=G^c\).  Let \(\lambda\) be the common parameter
probability law.  Put
\[
 \widehat A(q):={A(q)\over\ell(A(q))},
 \qquad
 \widehat B(q):={B(q)\over\ell(B(q))}.
\]

\begin{theorem}[Good-set comparison with tail stocks]
\label{thm:v12v37-tail}
Suppose
\begin{equation}
 \|\widehat A(q)-\widehat B(q)\|_1
 \leq\varepsilon(q)
 \quad(q\in G),                                   \tag{37.17}
\end{equation}
and
\begin{equation}
 \int_T\|\widehat A(q)\|_1\,d\lambda(q)\leq\tau_A,
 \qquad
 \int_T\|\widehat B(q)\|_1\,d\lambda(q)\leq\tau_B. \tag{37.18}
\end{equation}
Then
\begin{equation}
 \left\|
 \int\widehat A\,d\lambda-
 \int\widehat B\,d\lambda
 \right\|_1
 \leq
 \int_G\varepsilon(q)\,d\lambda(q)+\tau_A+\tau_B.
                                                        \tag{37.19}
\end{equation}
\end{theorem}

\begin{proof}
Split both integrals over \(G\) and \(T\).  Minkowski's inequality
and (37.17) bound the good part by
\(\int_G\varepsilon\,d\lambda\).  The triangle inequality and
(37.18) bound the two tail integrals by \(\tau_A+\tau_B\).
\end{proof}

The two tail stocks cannot be replaced merely by
\(\lambda(T)\) unless uniform bounds on
\(\|\widehat A\|_1\) and \(\|\widehat B\|_1\) are also proved.
Near a small normalizer, those quantities can be large even when the
uncovered set has small reference probability.

\subsection{Scale-summable local certificate}

At adjacent scale \(n\), let a certified cover give box errors
\(\varepsilon_{n,Q}\) and tails
\(\tau_{A,n},\tau_{B,n}\).  Define
\begin{equation}
 e_n:=
 \int_{G_n}\varepsilon_n(q)\,d\lambda_n(q)
 +\tau_{A,n}+\tau_{B,n}.                          \tag{37.20}
\end{equation}
If
\begin{equation}
                         \sum_ne_n<\infty,        \tag{37.21}
\end{equation}
then Corollary \ref{cor:v12v36-projective-repair} applies to the
integrated adjacent boundary states.  For local cylinder families,
the same construction may be performed with cylinder-dependent
dimension or rank and the cylinder-local repair theorem.

\begin{assessment}[Status after V37]
\label{ass:v12v37-status}
The local quantities required by V36 are now reducible to certified
box-centre arithmetic and variation bounds.  Strict positive margins
make subdivision terminate on compact good sets.  Uncovered regions
enter through explicit normalized trace-mass stocks, and the sum of
box and tail errors is the adjacent projective budget.
\end{assessment}

\begin{firewall}[Claim boundary after V37]
\label{fire:v12v37-claim}
No physical RPESM box centre, variation modulus, good-set cover, or
tail stock has been computed.  Compactness and strict normalizer
margins are hypotheses, not conclusions.  The missing Duhamel fibre,
determinant-line gluing, nonlinear gravitational tightness, and the
full SM--Einstein continuum remain open.
\end{firewall}

\subsection{Parent record}

This V37 note appends to the validated V36 file with record
\[
\begin{array}{c|c}
\text{lines}&10558\\
\text{bytes}&361767\\
\text{SHA--256}&
\texttt{82B723DC6E43516334E4188F8B1E7C66DD108993D6277948654E1C806BD26992}.
\end{array}
\]

% =============================================================================
% END APPENDED TWELFTH-CYCLE V37 MATERIAL
% =============================================================================
% =============================================================================
% BEGIN APPENDED TWELFTH-CYCLE V38 MATERIAL
% =============================================================================

\section{Twelfth-cycle V38: negative normalizer moments close the
uncovered-tail estimate}
\label{sec:v12v38}

V37 shows that an uncertified region must be controlled by the
normalized trace mass
\[
 \int_T{\operatorname{Tr}A\over\ell(A)}\,d\lambda,
\]
rather than by its reference probability alone.  This note gives a
sharp elementary criterion for that tail: a positive moment of the
kernel mass, a negative moment of the physical normalizer, and a
small probability for the uncovered set.

\subsection{Three-factor tail inequality}

Let \((\Omega,{\cal F},\lambda)\) be a probability space.  Let
\(A(\omega)\succeq0\) and let
\[
 X(\omega):=\operatorname{Tr}A(\omega),\qquad
 Z(\omega):=\ell(A(\omega)).
\]
Assume \(Z>0\) almost surely.  The trace norm of the physically
normalized kernel is
\begin{equation}
 \left\|{A\over Z}\right\|_1={X\over Z}.          \tag{38.1}
\end{equation}

\begin{theorem}[Moment-to-tail conversion]
\label{thm:v12v38-moment-tail}
Let \(p,q>1\) satisfy
\begin{equation}
                         {1\over p}+{1\over q}<1,
 \qquad
 \theta:=1-{1\over p}-{1\over q}>0.              \tag{38.2}
\end{equation}
If
\begin{equation}
 \|X\|_{L^p(\lambda)}\leq M_p,\qquad
 \|Z^{-1}\|_{L^q(\lambda)}\leq N_q,              \tag{38.3}
\end{equation}
then every measurable \(T\subset\Omega\) obeys
\begin{equation}
 \int_T\left\|{A\over Z}\right\|_1\,d\lambda
 \leq
 M_pN_q\,\lambda(T)^\theta.                      \tag{38.4}
\end{equation}
\end{theorem}

\begin{proof}
Choose \(r=\theta^{-1}\), so
\(1/p+1/q+1/r=1\).  Apply the three-factor H\"older inequality to
\(X\), \(Z^{-1}\), and the indicator \(1_T\):
\[
 \int_T{X\over Z}\,d\lambda
 \leq
 \|X\|_p\|Z^{-1}\|_q\|1_T\|_r
 \leq
 M_pN_q\,\lambda(T)^{1/r}.
\]
Use (38.1) and \(1/r=\theta\).
\end{proof}

The strict inequality in (38.2) is necessary for this argument to
gain a positive power of the uncovered probability.  Bounds exactly
on the H\"older line \(1/p+1/q=1\) give integrability but no uniform
decay with \(\lambda(T)\).

\begin{corollary}[Two-kernel tail stock]
\label{cor:v12v38-two-kernel}
Let \(A,B\succeq0\) have physical normalizers
\(Z_A=\ell(A)>0\) and \(Z_B=\ell(B)>0\).  Suppose
\[
 \|\operatorname{Tr}A\|_{p_A}\leq M_A,\quad
 \|Z_A^{-1}\|_{q_A}\leq N_A,
\]
and analogously \(M_B,N_B,p_B,q_B\), with
\(\theta_A,\theta_B>0\) as in (38.2).  Then the V37 tail stocks may
be chosen as
\begin{equation}
 \tau_A=M_AN_A\lambda(T)^{\theta_A},\qquad
 \tau_B=M_BN_B\lambda(T)^{\theta_B}.              \tag{38.5}
\end{equation}
\end{corollary}

\begin{proof}
Apply Theorem \ref{thm:v12v38-moment-tail} separately to the two
kernels.
\end{proof}

\subsection{Producing the bad-set probability}

Let \(Y\geq0\) be a control variable and let
\(T_R=\{Y>R\}\).  If
\[
 \mathbb E Y^s\leq C_s,
\]
then Markov's inequality and (38.4) give
\begin{equation}
 \int_{T_R}\left\|{A\over Z}\right\|_1d\lambda
 \leq
 M_pN_q C_s^\theta R^{-s\theta}.                 \tag{38.6}
\end{equation}
If instead
\[
 \mathbb E e^{cY^\alpha}\leq C_{\exp},
\]
then
\begin{equation}
 \int_{T_R}\left\|{A\over Z}\right\|_1d\lambda
 \leq
 M_pN_q C_{\exp}^{\theta}
 e^{-c\theta R^\alpha}.                          \tag{38.7}
\end{equation}

\begin{proof}
The probability bounds are
\[
 \lambda(T_R)\leq C_sR^{-s},
 \qquad
 \lambda(T_R)\leq C_{\exp}e^{-cR^\alpha},
\]
respectively.  Substitute them into (38.4).
\end{proof}

Thus a reference good-set probability becomes useful only after
being multiplied by the moment stocks that control amplification by
the inverse normalizer.

\subsection{Small-ball estimates and negative moments}

The negative moment in (38.3) is equivalent, up to exponent loss, to
quantitative exclusion of small normalizers.

\begin{theorem}[Small-ball bound implies inverse moments]
\label{thm:v12v38-small-ball}
Let \(Z>0\) almost surely.  Suppose that for some
\(\alpha,C,s_0>0\),
\begin{equation}
 \lambda\{Z\leq s\}\leq Cs^\alpha
 \qquad(0<s\leq s_0).                            \tag{38.8}
\end{equation}
Then for every \(0<q<\alpha\),
\begin{equation}
 \mathbb E Z^{-q}
 \leq
 s_0^{-q}
 +
 {Cq\over\alpha-q}s_0^{\alpha-q}.               \tag{38.9}
\end{equation}
Conversely, if \(\mathbb E Z^{-q}\leq N\), then
\begin{equation}
 \lambda\{Z\leq s\}\leq Ns^q
 \qquad(s>0).                                    \tag{38.10}
\end{equation}
\end{theorem}

\begin{proof}
For the first direction, use the layer-cake formula:
\[
 \mathbb E Z^{-q}
 =
 \int_0^\infty
 \lambda\{Z^{-q}>t\}\,dt.
\]
The part \(0<t<s_0^{-q}\) is at most \(s_0^{-q}\).  For
\(t\geq s_0^{-q}\), equation (38.8) gives
\[
 \lambda\{Z^{-q}>t\}
 =
 \lambda\{Z<t^{-1/q}\}
 \leq Ct^{-\alpha/q}.
\]
Because \(q<\alpha\),
\[
 \int_{s_0^{-q}}^\infty Ct^{-\alpha/q}\,dt
 =
 {Cq\over\alpha-q}s_0^{\alpha-q},
\]
proving (38.9).  The converse is Markov's inequality applied to
\(Z^{-q}\):
\[
 \lambda\{Z\leq s\}
 =
 \lambda\{Z^{-q}\geq s^{-q}\}
 \leq s^q\mathbb E Z^{-q}.
\]
\end{proof}

A proof that \(Z\ne0\) almost surely is much weaker than (38.8).
The exponent \(\alpha\) measures how rapidly probability leaves the
small-normalizer region and determines the available inverse
moments.

\subsection{Scale allocation}

At scale \(n\), suppose the uncovered region \(T_n\) satisfies
\[
 \lambda_n(T_n)\leq b_n
\]
and the two normalized-kernel moment stocks give constants
\(C_{A,n},C_{B,n}\) and exponents
\(\theta_{A,n},\theta_{B,n}>0\).  Then V37 may take
\begin{equation}
 e_n
 \leq
 \int_{G_n}\varepsilon_n(q)\,d\lambda_n(q)
 +
 C_{A,n}b_n^{\theta_{A,n}}
 +
 C_{B,n}b_n^{\theta_{B,n}}.                     \tag{38.11}
\end{equation}

\begin{proposition}[Geometric closure criterion]
\label{prop:v12v38-geometric}
Assume uniformly
\[
 C_{A,n}+C_{B,n}\leq C\Lambda^{\gamma n},
 \qquad
 b_n\leq C_0\Lambda^{-\beta n},
 \qquad
 \theta_{A,n},\theta_{B,n}\geq\theta>0
\]
for some \(\Lambda>1\).  If
\begin{equation}
                         \beta\theta>\gamma       \tag{38.12}
\end{equation}
and the good-set errors are summable, then
\(\sum_ne_n<\infty\).
\end{proposition}

\begin{proof}
Each tail term is bounded by a constant times
\[
 \Lambda^{\gamma n}\Lambda^{-\beta\theta n}
 =
 \Lambda^{-(\beta\theta-\gamma)n},
\]
which is a summable geometric sequence under (38.12).  Add the
summable good-set errors.
\end{proof}

This inequality exposes the competition hidden by a bare
high-probability statement: the good-set complement must shrink
faster than the normalized visibility stocks grow.

\subsection{What must produce the small-ball exponent}

For a positive reflected body, the physical normalizer is
nonnegative.  A usable estimate (38.8) might come from:
\begin{enumerate}
\item a deterministic lower bound on the certified good set;
\item transversality of the normalizer to its zero locus combined
with a bounded-density or coarea estimate;
\item a conditional small-ball inequality for the shell normalizer;
or
\item an exact reflected-square factorization with a quantitative
lower bound on its half-space factor.
\end{enumerate}
These are possible sources, not results already established for
RPESM.  The next upstream step is to prove the finite-dimensional
transversality-to-small-ball implication with all density and
tubular constants explicit.

\begin{assessment}[Status after V38]
\label{ass:v12v38-status}
The uncovered-tail term of V37 now has a precise sufficient
criterion.  Positive moments of the kernel trace and negative
moments of the physical normalizer yield a power of the bad-set
probability.  Negative moments in turn follow from a quantitative
small-ball exponent.  The geometric condition (38.12) states exactly
when shrinking bad sets overcome growing visibility stocks along the
regulator tower.
\end{assessment}

\begin{firewall}[Claim boundary after V38]
\label{fire:v12v38-claim}
No RPESM trace moment, inverse-normalizer moment, small-ball exponent,
or bad-set rate has been proved.  Almost-sure nonvanishing alone is
insufficient.  The note does not provide the missing Duhamel fibre,
determinant-line gluing, gravitational tightness, or the full
SM--Einstein continuum.
\end{firewall}

\subsection{Parent record}

This V38 note appends to the validated V37 file with record
\[
\begin{array}{c|c}
\text{lines}&10862\\
\text{bytes}&371436\\
\text{SHA--256}&
\texttt{C47A3B581D7F3F1F8357352D11F4CB4C81503939390EB13224D60D12030E34A3}.
\end{array}
\]

% =============================================================================
% END APPENDED TWELFTH-CYCLE V38 MATERIAL
% =============================================================================
% =============================================================================
% BEGIN APPENDED TWELFTH-CYCLE V39 MATERIAL
% =============================================================================

\section{Twelfth-cycle V39: coarea small-ball exponents and the
generic reflected-square obstruction}
\label{sec:v12v39}

V38 requires a quantitative small-ball exponent for the physical
normalizer.  An archived result propagates such an exponent to
coarea and inverse-Hessian tails but assumes the exponent as input.
This note derives it from finite-dimensional transversality.  The
calculation also shows that a generic complex reflected square gives
only the borderline exponent one, which is too weak for the
separate negative-moment factorization used in V38.

\subsection{A quantitative coarea theorem}

Let \(\Omega\subset\mathbb R^m\) be measurable, let
\(1\leq k\leq m\), and let
\[
 F:\Omega\longrightarrow\mathbb R^k
\]
be Lipschitz.  Write
\[
 J_kF(x):=
 \sqrt{\det\{DF(x)DF(x)^T\}}
\]
at differentiability points.  Let \(\lambda\) have Lebesgue density
\(h\).

\begin{theorem}[Transversality implies a small-ball exponent]
\label{thm:v12v39-coarea}
Fix \(r_0>0\) and suppose:
\begin{align}
 J_kF(x)&\geq j_0>0
 &&\text{for almost every }x\in\Omega
                 \text{ with }|F(x)|\leq r_0,    \tag{39.1}\\
 {\cal H}^{m-k}\bigl(\Omega\cap F^{-1}(y)\bigr)
 &\leq L
 &&\text{for almost every }|y|\leq r_0.          \tag{39.2}
\end{align}
If \(h\in L^\infty\), then
\begin{equation}
 \lambda\{|F|\leq r\}
 \leq
 {\|h\|_\infty L\omega_k\over j_0}\,r^k,
 \qquad 0<r\leq r_0,                             \tag{39.3}
\end{equation}
where \(\omega_k\) is the volume of the unit ball in
\(\mathbb R^k\).

More generally, if \(h\in L^s\) for \(1<s<\infty\), with conjugate
exponent \(s'=s/(s-1)\), then
\begin{equation}
 \lambda\{|F|\leq r\}
 \leq
 \|h\|_{L^s}
 \left({L\omega_k\over j_0}\right)^{1/s'}
 r^{k/s'}.                                       \tag{39.4}
\end{equation}
\end{theorem}

\begin{proof}
Let \(E_r=\{x\in\Omega:|F(x)|\leq r\}\).  The coarea formula and
(39.2) give
\begin{align*}
 \int_{E_r}J_kF(x)\,dx
 &=
 \int_{|y|\leq r}
 {\cal H}^{m-k}\bigl(\Omega\cap F^{-1}(y)\bigr)\,dy\\
 &\leq L\omega_kr^k.
\end{align*}
The lower bound (39.1) therefore implies
\[
 |E_r|\leq {L\omega_k\over j_0}r^k.
\]
Multiplication by the \(L^\infty\) density bound proves (39.3).
For \(h\in L^s\), H\"older's inequality gives
\[
 \lambda(E_r)=\int_{E_r}h\,dx
 \leq\|h\|_s|E_r|^{1/s'},
\]
which yields (39.4).
\end{proof}

Uniform control of the fibre volume in (39.2) is essential.  A
Jacobian lower bound alone does not prevent infinitely many
disjoint sheets from crossing the same target ball.

\subsection{Normalizers that are powers of an amplitude}

Let
\begin{equation}
                         Z(x)=|F(x)|^a,\qquad a>0.           \tag{39.5}
\end{equation}

\begin{corollary}[Codimension divided by vanishing order]
\label{cor:v12v39-exponent}
Under Theorem \ref{thm:v12v39-coarea}, the small-ball bound for \(Z\)
has exponent
\begin{equation}
 \alpha={k\over a}
 \quad\text{when }h\in L^\infty,                 \tag{39.6}
\end{equation}
and
\begin{equation}
 \alpha={k\over as'}
 \quad\text{when }h\in L^s.                      \tag{39.7}
\end{equation}
Consequently Theorem \ref{thm:v12v38-small-ball} supplies
\(\mathbb E Z^{-q}<\infty\) for every \(q<\alpha\).
\end{corollary}

\begin{proof}
The event \(Z\leq u\) is
\(|F|\leq u^{1/a}\).  Substitute
\(r=u^{1/a}\) in (39.3) or (39.4), then apply the V38 theorem.
\end{proof}

For a complex scalar amplitude \(f\), take
\(F=(\operatorname{Re}f,\operatorname{Im}f)\), so \(k=2\).
With bounded density:
\[
 Z=|f| \quad\Longrightarrow\quad\alpha=2,
 \qquad
 Z=|f|^2 \quad\Longrightarrow\quad\alpha=1.       \tag{39.8}
\]
The reflected square is therefore exactly at the first-inverse-moment
threshold.

\subsection{A tubular alternative}

Transversality may fail at a higher-order zero.  Let
\(N=\{Z=0\}\subset\Omega\).  Suppose for some
\(k,a,c,T,r_0,z_0>0\),
\begin{align}
 |\{x\in\Omega:\operatorname{dist}(x,N)\leq r\}|
 &\leq Tr^k
 &&(0<r\leq r_0),                                \tag{39.9}\\
 Z(x)&\geq c\,\operatorname{dist}(x,N)^a
 &&(\operatorname{dist}(x,N)\leq r_0),           \tag{39.10}\\
 Z(x)&\geq z_0
 &&(\operatorname{dist}(x,N)>r_0).               \tag{39.11}
\end{align}

\begin{proposition}[Tubular small-ball certificate]
\label{prop:v12v39-tubular}
If \(h\in L^\infty\), then for
\[
 0<u\leq\min\{z_0,cr_0^a\},
\]
\begin{equation}
 \lambda\{Z\leq u\}
 \leq
 \|h\|_\infty T c^{-k/a}u^{k/a}.                 \tag{39.12}
\end{equation}
For \(h\in L^s\), the exponent becomes \(k/(as')\).
\end{proposition}

\begin{proof}
Equations (39.10)--(39.11) imply
\[
 \{Z\leq u\}
 \subset
 \left\{
 \operatorname{dist}(x,N)\leq(u/c)^{1/a}
 \right\}.
\]
Apply (39.9) and the density bound.  The \(L^s\) version follows by
H\"older exactly as in Theorem \ref{thm:v12v39-coarea}.
\end{proof}

This formulation separates the codimension \(k\) of the zero set
from the vanishing order \(a\) of the normalizer.

\subsection{Sharp local obstruction at a regular zero}

The exponent \(k/a\) is not merely an artifact of the upper bound.

\begin{proposition}[Divergence at a regular zero]
\label{prop:v12v39-divergence}
Suppose \(F:\Omega\to\mathbb R^k\) is \(C^1\), \(F(x_0)=0\), and
\(DF(x_0)\) has rank \(k\).  Suppose the probability density is
bounded below by a positive constant on a neighbourhood of \(x_0\).
For \(Z=|F|^a\),
\begin{equation}
                         \mathbb E Z^{-q}=\infty
 \qquad\text{whenever }aq\geq k.                 \tag{39.13}
\end{equation}
\end{proposition}

\begin{proof}
By the constant-rank theorem, there are local coordinates
\((y,z)\in\mathbb R^k\times\mathbb R^{m-k}\) near \(x_0\) in which
\(F=y\), up to smooth coordinate changes with Jacobians bounded above
and below by positive constants.  Restricting the \(z\)-variables to
a fixed small box and using the lower density bound, the expectation
dominates a positive constant times
\[
 \int_{|y|<\varepsilon}|y|^{-aq}\,dy
 =
 \omega_{k-1}
 \int_0^\varepsilon r^{k-1-aq}\,dr.
\]
The last integral diverges exactly when \(k-aq\leq0\).
\end{proof}

\begin{corollary}[Generic complex reflected square is borderline]
\label{cor:v12v39-reflected-square}
At a regular zero of a complex scalar \(f\), with a locally positive
bounded density,
\[
                         \mathbb E|f|^{-2q}=\infty
 \qquad(q\geq1).                                  \tag{39.14}
\]
Hence the separate-factor hypothesis \(q>1\) in V38 cannot hold
across such a zero for the normalizer \(Z=|f|^2\).
\end{corollary}

\begin{proof}
Use Proposition \ref{prop:v12v39-divergence} with \(k=2\) and
\(a=2\).
\end{proof}

This does not say that the normalized kernel mass
\(\operatorname{Tr}A/Z\) diverges.  Its numerator can vanish at the
same order and cancel the small denominator.  It says that separating
the numerator and inverse denominator by H\"older is too strong in
the generic reflected-square case.

\subsection{Revised route for the RPESM normalizer}

There are now three distinct possibilities.

\begin{enumerate}
\item \textbf{Uniform floor.}  If \(Z\geq\zeta>0\) on the physical
support, V37 applies directly.
\item \textbf{High-codimension or repelled zeros.}  If the small-ball
exponent exceeds one by enough to choose \(p,q\) in (38.2), V38
controls the uncovered tail.
\item \textbf{Correlated cancellation.}  If \(Z\) is a generic
reflected square with exponent one, estimate
\(\operatorname{Tr}A/Z\) directly, using a common factorization of
the numerator and normalizer.  Do not demand a nonexistent
\(L^q\), \(q>1\), inverse-normalizer moment.
\end{enumerate}

For example, if
\[
 A=f\overline f\,A_0,\qquad
 Z=\ell(A)=|f|^2\ell(A_0),
\]
then
\[
 {\operatorname{Tr}A\over Z}
 =
 {\operatorname{Tr}A_0\over\ell(A_0)}
\]
wherever \(f\ne0\), and the apparent singularity cancels.  A rigorous
version must handle matrix-valued common factors, changing rank, and
extension across the zero set.

\begin{assessment}[Status after V39]
\label{ass:v12v39-status}
The missing small-ball exponent now has two geometric sources:
coarea transversality and tubular volume plus a vanishing-order
bound.  Their exponent is codimension divided by vanishing order,
with an additional density-integrability loss.  A generic complex
reflected-square zero makes every inverse moment of order at least
one diverge, so the programme must use correlated numerator-
denominator cancellation unless a uniform floor or stronger zero
repulsion is proved.
\end{assessment}

\begin{firewall}[Claim boundary after V39]
\label{fire:v12v39-claim}
No RPESM normalizer zero set, codimension, vanishing order, fibre
volume, or density bound has been identified.  The generic
reflected-square obstruction is conditional on a regular zero with
positive local density.  It neither proves that such a zero occurs
nor constructs the required correlated factorization.  The full
SM--Einstein continuum remains open.
\end{firewall}

\subsection{Parent record}

This V39 note appends to the validated V38 file with record
\[
\begin{array}{c|c}
\text{lines}&11165\\
\text{bytes}&380176\\
\text{SHA--256}&
\texttt{DC36C62B7D172E146478A9F015E6E2AC405231FF8DE056D7BAE3CBE01BF3D787}.
\end{array}
\]

% =============================================================================
% END APPENDED TWELFTH-CYCLE V39 MATERIAL
% =============================================================================
% =============================================================================
% BEGIN APPENDED TWELFTH-CYCLE V40 MATERIAL
% =============================================================================

\section{Twelfth-cycle V40: compulsory companion audit and the
cancellation-preserving route}
\label{sec:v12v40}

\subsection{Audit identity}

This is the compulsory five-version audit following V35.  The newest
active companion files at the audit time were
\[
\begin{array}{c|r|r}
\text{file}&\text{counted lines}&\text{bytes}\\ \hline
\texttt{Renewal\_Yang\_Mills\_Mass\_Gap\_Research\_Notes\_V54.tex}
 &20667&704814\\
\texttt{Renewal\_NS\_Stability\_Research\_Notes\_V26.tex}
 &5320&278283.
\end{array}                                                   \tag{40.1}
\]
Their SHA--256 fingerprints are
\[
\begin{split}
 &\texttt{D96369C9E12B42623FEB374B916DF0165CA03FBBEF5BEB95BE8B74F4362218A2},
 \\
 &\texttt{82C887F8C2C69E4F28FAD7C3DC8DE5B9099CB5A4BF431EBA1FFF3E91935978AD}.
\end{split}                                                   \tag{40.2}
\]
They were read without modification.

\subsection{Yang--Mills V53: preserving induced matrix structure}

Yang--Mills V53 replaces a sum of absolute values over all matrix
entries by induced row and column Laurent stocks.  It proves a
uniform four-complex-variable inverse strip nearly forty times wider
than its preceding global scalar estimate.  This closes the logical
multivariable gap left by its exact coordinate-ray results.

The transferable principle is precise.  If a matrix perturbation
\(\Delta\) has certified induced one- and infinity-norm bounds, then
\begin{equation}
 \|\Delta\|_2
 \leq
 \sqrt{\|\Delta\|_1\|\Delta\|_\infty}.            \tag{40.3}
\end{equation}
Summing every entry before taking an induced norm can multiply the
bound by the matrix dimension and destroy a usable inverse margin.
For the SM reflected kernels, V34--V37 should therefore retain block
row, column, rank, and support structure before converting to trace
norm.

No numerical Yang--Mills strip transfers to RPESM; its Hessian and
coefficient arrays are different.

\subsection{Yang--Mills V54: an absolute-value no-go theorem}

After propagating the wider strip, Yang--Mills V54 proves that the
existing absolute Fourier-tail lemma still cannot certify its
\(M=2\) continuum response.  More strongly, an exact Wilson
coordinate-line zero bounds every possible common strip, and that
upper bound is incompatible with the strip width the absolute alias
lemma would require at \(M=2\).  Further refinement of the same
absolute majorant cannot close the desired conclusion.  The selected
next route retains signed Fourier residue classes and packet
cancellation.

This is directly parallel to V39.  For a generic reflected-square
normalizer, separate absolute bounds on the numerator and
\(Z^{-1}\) can demand a negative moment that diverges, even when the
normalized ratio is regular because the same vanishing factor occurs
in both terms.

\begin{proposition}[Absolute factorization can lose an exact
cancellation]
\label{prop:v12v40-absolute-loss}
Let \(g\) be measurable, let \(a,b\geq0\), and set
\[
 X=|g|^2a,\qquad Z=|g|^2b.
\]
On \(\{g\ne0,b>0\}\),
\begin{equation}
                         {X\over Z}={a\over b}.   \tag{40.4}
\end{equation}
It is possible that \(a/b\) is bounded while
\(\|Z^{-1}\|_{L^q}=\infty\) for every \(q\geq1\).
\end{proposition}

\begin{proof}
The identity is cancellation.  For the example, take the unit disk
in \(\mathbb R^2\) with bounded positive density,
\(g(x_1,x_2)=x_1+ix_2\), and \(a=b=1\).  Then \(X/Z=1\) away from
the origin, while
\[
 \int_{|x|<1}Z^{-q}\,dx
 =
 2\pi\int_0^1r^{1-2q}\,dr
 =
 \infty
 \qquad(q\geq1).
\]
\end{proof}

Thus the failure of the separate V38 moment hypothesis need not
signal divergence of the physical normalized kernel.

\subsection{Navier--Stokes status}

Navier--Stokes V26 is unchanged.  Its rank-one norm calculation
continues to support structure-preserving estimates, but it supplies
no new normalizer, cancellation, determinant, or projective input.
No regularity conclusion is imported.

\subsection{Direction after the audit}

The next SM note will replace the separate inverse-normalizer demand
where a common factor is available.  It will prove:

\begin{enumerate}
\item a range-visibility inequality
\(\operatorname{Tr}A/\ell(A)\leq\chi^{-1}\) based on the angle between
the kernel range and the positive normalizer matrix;
\item exact removal of a common scalar vanishing factor from a
normalized positive kernel;
\item stability under an approximate common factor; and
\item insertion of the resulting bound into V36--V37 without any
negative moment of the canceled scalar.
\end{enumerate}

\begin{assessment}[Status after V40]
\label{ass:v12v40-status}
The audit validates the change from absolute factor separation to
cancellation-preserving estimates.  Induced matrix structure can
improve a valid bound substantially, but an exact analytic
obstruction may still rule out the entire absolute method.  The
next acquisition will control normalized reflected kernels by common
factor cancellation and range visibility.
\end{assessment}

\begin{firewall}[Claim boundary after V40]
\label{fire:v12v40-claim}
No Yang--Mills response sign or mass gap and no Navier--Stokes
regularity theorem is imported.  No RPESM common factor, range angle,
or physical visibility bound has yet been proved.  The missing
Duhamel fibre and the full SM--Einstein continuum remain open.
\end{firewall}

\subsection{Parent record}

This V40 note appends to the validated V39 file with record
\[
\begin{array}{c|c}
\text{lines}&11465\\
\text{bytes}&389677\\
\text{SHA--256}&
\texttt{ABD0947B30306687F0F48A4C919A3E70D11B73F8C0F723A63581C27D90D3EF97}.
\end{array}
\]

% =============================================================================
% END APPENDED TWELFTH-CYCLE V40 MATERIAL
% =============================================================================
% =============================================================================
% BEGIN APPENDED TWELFTH-CYCLE V41 MATERIAL
% =============================================================================

\section{Twelfth-cycle V41: range visibility and cancellation of
common normalizer zeros}
\label{sec:v12v41}

V39 shows that separate inverse-normalizer moments can fail at a
generic reflected-square zero.  V40 confirms, through the
Yang--Mills companion, that an absolute factorization can be
intrinsically incapable of proving a conclusion that survives signed
or correlated cancellation.  This note estimates the normalized
positive kernel without separating its numerator from its
normalizer.

\subsection{The range-visibility constant}

Let \(S\succeq0\) represent the physical positive functional
\[
 \ell(H)=\operatorname{Tr}(SH).
\]
For a nonzero \(A\succeq0\), let \(P_A\) be the orthogonal projection
onto \(\operatorname{Ran}A\), and define
\begin{equation}
 \chi(A;S):=
 \lambda_{\min}
 \left(P_ASP_A\big|_{\operatorname{Ran}A}\right).           \tag{41.1}
\end{equation}
This number measures how much of every direction carrying kernel
mass is seen by the physical normalizer.

\begin{theorem}[Range visibility controls normalized trace mass]
\label{thm:v12v41-visibility}
If \(\chi(A;S)>0\), then
\begin{equation}
 \ell(A)\geq\chi(A;S)\operatorname{Tr}A,           \tag{41.2}
\end{equation}
and hence
\begin{equation}
 \left\|{A\over\ell(A)}\right\|_1
 =
 {\operatorname{Tr}A\over\ell(A)}
 \leq {1\over\chi(A;S)}.                          \tag{41.3}
\end{equation}
\end{theorem}

\begin{proof}
The range of \(A^{1/2}\) is \(\operatorname{Ran}A\).  Therefore
\[
 A^{1/2}SA^{1/2}
 \succeq
 \chi(A;S)A.
\]
Taking traces gives
\[
 \ell(A)
 =
 \operatorname{Tr}(SA)
 =
 \operatorname{Tr}(A^{1/2}SA^{1/2})
 \geq
 \chi(A;S)\operatorname{Tr}A.
\]
Positivity gives \(\|A\|_1=\operatorname{Tr}A\), proving (41.3).
\end{proof}

\begin{corollary}[Tail control without an inverse moment]
\label{cor:v12v41-tail}
If
\[
 \chi(A(\omega);S)\geq\chi_0>0
\]
for almost every \(\omega\) in a measurable set \(T\), then
\begin{equation}
 \int_T
 \left\|{A(\omega)\over\ell(A(\omega))}\right\|_1
 d\lambda(\omega)
 \leq{\lambda(T)\over\chi_0}.                    \tag{41.4}
\end{equation}
\end{corollary}

\begin{proof}
Apply (41.3) pointwise and integrate.
\end{proof}

This estimate remains valid when \(\ell(A)\) tends to zero, provided
the range angle does not collapse.

\subsection{Full-rank and rank-one normalizers}

If \(S\succ0\), then
\[
 \chi(A;S)\geq\lambda_{\min}(S)
\]
for every nonzero positive \(A\).  Hence a full-rank normalizer
matrix gives uniform visibility automatically.

At the other extreme, let \(S=ww^*\) have rank one.

\begin{proposition}[Rank-one angle criterion]
\label{prop:v12v41-angle}
For \(S=ww^*\), one has \(\chi(A;S)>0\) only if
\(\operatorname{rank}A\leq1\).  If \(A=vv^*\ne0\), then
\begin{equation}
 \chi(A;ww^*)
 =
 {|w^*v|^2\over\|v\|^2},                         \tag{41.5}
\end{equation}
and
\begin{equation}
 {\operatorname{Tr}A\over\ell(A)}
 =
 {\|v\|^2\over|w^*v|^2}
 =
 {1\over\|w\|^2\cos^2\angle(w,v)}.               \tag{41.6}
\end{equation}
\end{proposition}

\begin{proof}
The restriction of a rank-one operator to a space of dimension at
least two has a nontrivial kernel, so its least eigenvalue there is
zero.  If \(\operatorname{Ran}A=\operatorname{span}(v)\), evaluation
on \(v/\|v\|\) gives (41.5).  Also
\[
 \operatorname{Tr}(vv^*)=\|v\|^2,\qquad
 \ell(vv^*)=\operatorname{Tr}(ww^*vv^*)=|w^*v|^2,
\]
which proves (41.6).
\end{proof}

Thus a partition vector sees a rank-one reflected square uniformly
exactly when the generating vector stays at a uniform nonzero angle
from the orthogonal complement of the partition vector.

\subsection{Exact common-factor cancellation}

Let \(\tau\geq0\) and let \(C\succeq0\) with \(\ell(C)>0\).  Put
\[
                         A=\tau C.               \tag{41.7}
\]

\begin{theorem}[Common scalar zero cancels before normalization]
\label{thm:v12v41-common-factor}
On \(\{\tau>0\}\),
\begin{equation}
 {A\over\ell(A)}={C\over\ell(C)}.                 \tag{41.8}
\end{equation}
Consequently the size, vanishing order, and inverse moments of
\(\tau\) do not enter the normalized kernel.  If
\(C/\ell(C)\) has a continuous limit at a point where \(\tau=0\),
equation (41.8) defines the unique continuous extension there.
\end{theorem}

\begin{proof}
Linearity gives \(\ell(A)=\tau\ell(C)\).  Cancel the positive scalar
\(\tau\).  Continuity gives the final assertion.
\end{proof}

For two kernels with the same scalar factor,
\[
 A=\tau C_A,\qquad B=\tau C_B,
\]
one has the exact identity
\begin{equation}
 \left\|
 {A\over\ell(A)}-{B\over\ell(B)}
 \right\|_1
 =
 \left\|
 {C_A\over\ell(C_A)}-{C_B\over\ell(C_B)}
 \right\|_1                                      \tag{41.9}
\end{equation}
where \(\tau>0\).  A small physical normalizer caused solely by
\(\tau\) therefore does not amplify the comparison error.

\subsection{Relative errors around a common factor}

Exact divisibility may hold only up to a controlled remainder.
Write
\begin{equation}
 A=\tau(C+E),\qquad C\succeq0,\quad C+E\succeq0.  \tag{41.10}
\end{equation}
The useful error is relative to \(\tau\): the trace norm of \(E\),
not the absolute trace norm of \(\tau E\).

\begin{theorem}[Stable cancellation under a relative remainder]
\label{thm:v12v41-relative}
Assume
\begin{equation}
 \ell(C)\geq\zeta>0,\qquad
 \|E\|_1\leq\eta,\qquad
 \chi(C+E;S)\geq\chi_0>0.                         \tag{41.11}
\end{equation}
Then on \(\{\tau>0\}\),
\begin{equation}
 \left\|
 {A\over\ell(A)}
 -
 {C\over\ell(C)}
 \right\|_1
 \leq
 {1+c_\ell/\chi_0\over\zeta}\,\eta,
 \qquad c_\ell=\|S\|.                             \tag{41.12}
\end{equation}
The bound is independent of \(\tau\).
\end{theorem}

\begin{proof}
First cancel \(\tau\):
\[
 {A\over\ell(A)}
 =
 {C+E\over\ell(C+E)}.
\]
Apply Theorem \ref{thm:v12v36-general-normalizer} with its first
kernel equal to \(C\) and its second equal to \(C+E\).
Theorem \ref{thm:v12v41-visibility} gives
\[
 {\operatorname{Tr}(C+E)\over\ell(C+E)}
 \leq\chi_0^{-1}.
\]
The remaining inputs to V36 are exactly
\(\ell(C)\geq\zeta\) and \(\|E\|_1\leq\eta\), which yield
(41.12).
\end{proof}

\begin{corollary}[Two relatively factored kernels]
\label{cor:v12v41-two-factor}
Let
\[
 A=\tau(C_A+E_A),\qquad
 B=\tau(C_B+E_B),
\]
and suppose the hypotheses of Theorem
\ref{thm:v12v41-relative} hold with constants
\((\zeta_A,\eta_A,\chi_A)\) and
\((\zeta_B,\eta_B,\chi_B)\).  Then
\begin{align}
 \left\|
 {A\over\ell(A)}-{B\over\ell(B)}
 \right\|_1
 \leq{}&
 {1+c_\ell/\chi_A\over\zeta_A}\eta_A
 +
 \left\|
 {C_A\over\ell(C_A)}-{C_B\over\ell(C_B)}
 \right\|_1
 \nonumber\\
 &+
 {1+c_\ell/\chi_B\over\zeta_B}\eta_B.            \tag{41.13}
\end{align}
\end{corollary}

\begin{proof}
Insert the two normalized core kernels between the left-hand terms
and use the triangle inequality together with (41.12).
\end{proof}

No negative moment of the common factor appears in (41.12)--(41.13).

\subsection{Matrix factors and a limitation}

A common scalar commutes with every kernel and with \(\ell\).
A common matrix congruence does not generally cancel:
\[
 A=T^*C_AT,\qquad B=T^*C_BT.
\]
The normalizers
\(\operatorname{Tr}(ST^*C_AT)\) and
\(\operatorname{Tr}(ST^*C_BT)\) depend on the orientation of \(T\)
relative to \(S\).  Controlling them requires the range-visibility
constant for \(T^*C_iT\), or an intertwining relation between \(T\)
and \(S\).  Treating a matrix factor as if it were scalar would lose
exactly the angle information measured by \(\chi\).

\subsection{Insertion into the projective budget}

At adjacent scale \(n\), suppose the unnormalized coarse kernels
share a scalar factor \(\tau_n\) and have cores satisfying
Corollary \ref{cor:v12v41-two-factor}.  Define the right side of
(41.13) to be \(e_n^{\rm core}\).  On an uncovered set \(T_n\),
suppose the actual kernels obey uniform range visibility
\(\chi_{A,n},\chi_{B,n}>0\).  Then the V37 total adjacent error may
be taken as
\begin{equation}
 e_n
 \leq
 \int_{G_n}e_n^{\rm core}(q)\,d\lambda_n(q)
 +
 {\lambda_n(T_n)\over\chi_{A,n}}
 +
 {\lambda_n(T_n)\over\chi_{B,n}}.                \tag{41.14}
\end{equation}
If \(\sum_ne_n<\infty\), the V36 projective repair applies without
an inverse moment of \(\tau_n\).

This route is strictly adapted to common zeros.  If the two kernels
vanish at different rates or in different directions, no scalar
cancellation is available and the relative remainders in (41.10)
will fail to stay bounded.

\begin{assessment}[Status after V41]
\label{ass:v12v41-status}
The generic reflected-square obstruction of V39 no longer forces a
negative-moment estimate when numerator and normalizer share their
zero.  Uniform visibility of the kernel range directly bounds
normalized trace mass.  A common scalar factor cancels exactly, and
relative trace-class remainders produce a scale-free normalized
comparison.  These estimates feed the V36--V37 projective budget.
\end{assessment}

\begin{firewall}[Claim boundary after V41]
\label{fire:v12v41-claim}
No RPESM common factor, range projection, visibility angle, or
relative remainder has been computed.  The existence of a
reflection-square representation does not by itself prove uniform
visibility of the physical partition functional.  The missing
Duhamel fibre, determinant-line gluing, nonlinear gravitational
tightness, and the full SM--Einstein continuum remain open.
\end{firewall}

\subsection{Parent record}

This V41 note appends to the validated V40 file with record
\[
\begin{array}{c|c}
\text{lines}&11627\\
\text{bytes}&395660\\
\text{SHA--256}&
\texttt{BFA8DD625282503CF74FF670037C84B74B12DA9EC0D3F262D887FEA21ED30AD4}.
\end{array}
\]

% =============================================================================
% END APPENDED TWELFTH-CYCLE V41 MATERIAL
% =============================================================================
% =============================================================================
% BEGIN APPENDED TWELFTH-CYCLE V42 MATERIAL
% =============================================================================

\section{Twelfth-cycle V42: canonical mass--shape decomposition and
the invisible face}
\label{sec:v12v42}

V41 cancels a common scalar factor when one is identified.  For a
positive kernel there is a canonical factorization that always
exists: total trace mass times a trace-one shape.  This separates
harmless amplitude zeros from the genuine obstruction, namely that
the normalized shape approaches the face invisible to the physical
partition functional.

\subsection{Canonical factorization}

Let \(A\succeq0\), \(A\ne0\), and define
\begin{equation}
 m_A:=\operatorname{Tr}A,\qquad
 \sigma_A:={A\over m_A}.                           \tag{42.1}
\end{equation}
Then \(\sigma_A\succeq0\) and
\(\operatorname{Tr}\sigma_A=1\).  Let
\[
 \ell(H)=\operatorname{Tr}(SH),\qquad S\succeq0,
\]
and define the shape visibility
\begin{equation}
 v_A:=\ell(\sigma_A)=\operatorname{Tr}(S\sigma_A).          \tag{42.2}
\end{equation}

\begin{theorem}[Mass cancels; shape visibility remains]
\label{thm:v12v42-mass-shape}
For every \(A\succeq0\), \(A\ne0\),
\begin{equation}
 \ell(A)=m_Av_A.                                  \tag{42.3}
\end{equation}
If \(v_A>0\), then
\begin{equation}
 {A\over\ell(A)}
 =
 {\sigma_A\over v_A},
 \qquad
 \left\|{A\over\ell(A)}\right\|_1
 =
 {1\over v_A}.                                   \tag{42.4}
\end{equation}
Thus small total mass \(m_A\) never causes divergence of the
physically normalized positive kernel.  Divergence can occur only
through loss of shape visibility \(v_A\).
\end{theorem}

\begin{proof}
Linearity gives
\[
 \ell(A)=\ell(m_A\sigma_A)=m_A\ell(\sigma_A)=m_Av_A.
\]
Cancel \(m_A>0\) to obtain the first identity in (42.4).
Since \(\sigma_A\succeq0\) and has trace one,
\(\|\sigma_A\|_1=1\), proving the second.
\end{proof}

The separate inverse moment of the full normalizer,
\[
 \ell(A)^{-q}=m_A^{-q}v_A^{-q},
\]
can therefore diverge solely because \(m_A\to0\), even though the
normalized kernel is completely regular.

\subsection{The invisible face}

Let \(P_S\) be the support projection of \(S\), and let
\[
 {\cal D}_d
 :=
 \{\sigma\succeq0:\operatorname{Tr}\sigma=1\}
\]
be the compact density-matrix simplex.  Define
\begin{equation}
 {\cal F}_S:=
 \{\sigma\in{\cal D}_d:
       \operatorname{Ran}\sigma\subset\ker S\}.    \tag{42.5}
\end{equation}

\begin{proposition}[Zero visibility is an exposed face]
\label{prop:v12v42-invisible-face}
For \(\sigma\in{\cal D}_d\),
\begin{equation}
 \operatorname{Tr}(S\sigma)=0
 \quad\Longleftrightarrow\quad
 \sigma\in{\cal F}_S.                             \tag{42.6}
\end{equation}
The set \({\cal F}_S\) is a closed exposed face of
\({\cal D}_d\).
\end{proposition}

\begin{proof}
The matrix \(S^{1/2}\sigma S^{1/2}\) is positive semidefinite and
has trace \(\operatorname{Tr}(S\sigma)\).  Zero trace therefore
implies
\(S^{1/2}\sigma S^{1/2}=0\).  Equivalently,
\(\sigma^{1/2}S^{1/2}=0\), so
\(\operatorname{Ran}\sigma\subset\ker S\).  The converse is
immediate.  The zero set of the nonnegative linear functional
\(\sigma\mapsto\operatorname{Tr}(S\sigma)\) is a closed exposed
face.
\end{proof}

Let \(\mu_S>0\) be the smallest positive eigenvalue of \(S\) on
\(\operatorname{Ran}S\).

\begin{corollary}[Visible-mass certificate]
\label{cor:v12v42-visible-mass}
For every \(\sigma\in{\cal D}_d\),
\begin{equation}
 \operatorname{Tr}(S\sigma)
 \geq
 \mu_S\operatorname{Tr}(P_S\sigma).              \tag{42.7}
\end{equation}
Hence
\[
 \operatorname{Tr}(P_S\sigma)\geq\eta>0
 \quad\Longrightarrow\quad
 v(\sigma)\geq\mu_S\eta.
\]
\end{corollary}

\begin{proof}
The spectral theorem gives \(S\succeq\mu_SP_S\).
Pair this operator inequality with \(\sigma\succeq0\).
\end{proof}

This average visible-mass condition is weaker than the V41
range-visibility condition.  V41 asks every vector in
\(\operatorname{Ran}A\) to be visible.  Equation (42.7) requires
only a positive amount of the trace-one shape to lie in the visible
subspace, which is exactly what the scalar normalization ratio uses.

\subsection{Stability of normalized shapes}

\begin{theorem}[Shape comparison at fixed visibility]
\label{thm:v12v42-shape-comparison}
Let \(A,B\succeq0\) be nonzero, and let
\(\sigma_A,\sigma_B\) be their trace-one shapes.  Put
\[
 \delta_{\rm sh}:=\|\sigma_A-\sigma_B\|_1.
\]
If
\begin{equation}
 v_A=\operatorname{Tr}(S\sigma_A)\geq\nu>0,\qquad
 v_B=\operatorname{Tr}(S\sigma_B)\geq\nu>0,       \tag{42.8}
\end{equation}
then
\begin{equation}
 \left\|
 {A\over\ell(A)}-{B\over\ell(B)}
 \right\|_1
 \leq
 \left({1\over\nu}+{\|S\|\over\nu^2}\right)
 \delta_{\rm sh}.                                 \tag{42.9}
\end{equation}
No lower bound on \(\operatorname{Tr}A\) or
\(\operatorname{Tr}B\) is required.
\end{theorem}

\begin{proof}
By Theorem \ref{thm:v12v42-mass-shape}, the left side is
\[
 \left\|{\sigma_A\over v_A}
              -{\sigma_B\over v_B}\right\|_1.
\]
Split the difference:
\begin{align*}
 \left\|{\sigma_A\over v_A}
              -{\sigma_B\over v_B}\right\|_1
 &\leq
 {\|\sigma_A-\sigma_B\|_1\over v_A}
 +
 \|\sigma_B\|_1\,{|v_A-v_B|\over v_Av_B}.
\end{align*}
Here \(\|\sigma_B\|_1=1\), and trace duality gives
\[
 |v_A-v_B|
 =
 |\operatorname{Tr}\{S(\sigma_A-\sigma_B)\}|
 \leq\|S\|\delta_{\rm sh}.
\]
Use (42.8) to obtain (42.9).
\end{proof}

\begin{corollary}[Relative unnormalized defect]
\label{cor:v12v42-relative-defect}
If
\begin{equation}
 \|A-B\|_1\leq\varepsilon\,\operatorname{Tr}A,
 \qquad 0\leq\varepsilon<1,                      \tag{42.10}
\end{equation}
then
\begin{equation}
 \|\sigma_A-\sigma_B\|_1\leq2\varepsilon.         \tag{42.11}
\end{equation}
Under (42.8), the right side of (42.9) is therefore at most
\[
 2\varepsilon
 \left({1\over\nu}+{\|S\|\over\nu^2}\right).      \tag{42.12}
\]
\end{corollary}

\begin{proof}
Theorem \ref{thm:v12v36-trace-normalization}, before division by two,
gives
\[
 \left\|{A\over\operatorname{Tr}A}
       -{B\over\operatorname{Tr}B}\right\|_1
 \leq
 {2\|A-B\|_1\over\operatorname{Tr}A}.
\]
Insert (42.10), then use Theorem
\ref{thm:v12v42-shape-comparison}.
\end{proof}

The correct smallness parameter near a vanishing kernel is therefore
a relative trace defect, not an absolute defect divided by a fixed
global normalizer floor.

\subsection{Rank-one amplitude versus angle}

Let \(A=vv^*\), \(v\ne0\), and let \(S=ww^*\).  Write
\[
 v=r u,\qquad r=\|v\|,\quad\|u\|=1.
\]
Then
\begin{equation}
 m_A=r^2,\qquad
 \sigma_A=uu^*,\qquad
 v_A=|w^*u|^2.                                    \tag{42.13}
\end{equation}
Hence
\begin{equation}
 {A\over\ell(A)}
 =
 {uu^*\over|w^*u|^2}.                            \tag{42.14}
\end{equation}

\begin{proposition}[Amplitude zeros cancel; angle zeros do not]
\label{prop:v12v42-amplitude-angle}
If \(r\to0\) while \(u\to u_0\) with \(w^*u_0\ne0\), the normalized
kernel (42.14) has a finite limit.  If
\(u\) approaches the hyperplane \(w^\perp\), its trace norm diverges
like \(|w^*u|^{-2}\), independently of \(r\).
\end{proposition}

\begin{proof}
Equation (42.14) contains no \(r\).  Continuity gives the first
claim, while its trace norm is exactly \(|w^*u|^{-2}\), giving the
second.
\end{proof}

Thus the regular-zero calculation of V39 must be applied to the
shape overlap when the zero is angular.  Applying it to the total
amplitude can impose a spurious inverse-mass requirement.

\subsection{Compact shape sets}

\begin{proposition}[Compact separation from the invisible face]
\label{prop:v12v42-compact}
Let \({\cal K}\subset{\cal D}_d\) be compact and
\({\cal K}\cap{\cal F}_S=\varnothing\).  Then
\begin{equation}
 \inf_{\sigma\in{\cal K}}\operatorname{Tr}(S\sigma)>0.
                                                        \tag{42.15}
\end{equation}
\end{proposition}

\begin{proof}
The visibility function is continuous on the compact set
\({\cal K}\), so it attains a minimum.  Proposition
\ref{prop:v12v42-invisible-face} says that a zero minimum would put a
point of \({\cal K}\) in \({\cal F}_S\).
\end{proof}

This gives the correct finite-box target: cover the trace-one kernel
shapes and prove that their image stays away from the invisible face.
A lower bound on the unnormalized trace mass is unnecessary.

\subsection{Projective shape budget}

For adjacent repaired coarse kernels \(A_n,B_n\), suppose
\[
 v_{A,n},v_{B,n}\geq\nu_n>0,\qquad
 \|\sigma_{A,n}-\sigma_{B,n}\|_1\leq\delta_{{\rm sh},n}.
\]
Then the adjacent normalized physical error is
\begin{equation}
 e_n
 \leq
 \left({1\over\nu_n}
       +{\|S_n\|\over\nu_n^2}\right)
 \delta_{{\rm sh},n}.                            \tag{42.16}
\end{equation}
If \(\sum_ne_n<\infty\), the V36 projective repair applies.  This
criterion remains meaningful when both unnormalized kernels tend to
zero with the regulator.

\begin{assessment}[Status after V42]
\label{ass:v12v42-status}
Every positive kernel now has a canonical harmless mass factor and
a trace-one shape.  The physical normalization cancels the mass
exactly; only approach of the shape to the exposed invisible face
can cause amplification.  A uniform visible-mass bound and a
relative trace defect give an explicit projective error without any
absolute normalizer floor or inverse total-mass moment.
\end{assessment}

\begin{firewall}[Claim boundary after V42]
\label{fire:v12v42-claim}
No RPESM trace-one shape family, invisible-face separation, or
relative adjacent defect has been computed.  Rank changes can make
shape and support control delicate.  The missing Duhamel fibre,
determinant-line gluing, gravitational tightness, and the full
SM--Einstein continuum remain open.
\end{firewall}

\subsection{Parent record}

This V42 note appends to the validated V41 file with record
\[
\begin{array}{c|c}
\text{lines}&11964\\
\text{bytes}&405651\\
\text{SHA--256}&
\texttt{190C3E2DA10C1BFAEA09C03C2CC65FBF690F8B27AD80D4EB2B794C654AD0E2B1}.
\end{array}
\]

% =============================================================================
% END APPENDED TWELFTH-CYCLE V42 MATERIAL
% =============================================================================
% =============================================================================
% BEGIN APPENDED TWELFTH-CYCLE V43 MATERIAL
% =============================================================================

\section{Twelfth-cycle V43: moving-normalizer shape boxes and
scale-free angular defects}
\label{sec:v12v43}

V42 identifies the trace-one kernel shape as the correct object for
physical normalization.  In the coupled SM--Einstein system the
partition functional can itself depend on the retained metric,
boundary frame, or coarse identification.  This note gives local
visibility and comparison certificates when both the shape and the
normalizer matrix move.

\subsection{Joint perturbation of shape and normalizer}

Let \(S,S_0\succeq0\) and let
\(\sigma,\sigma_0\succeq0\) have trace one.  Define
\[
 v(S,\sigma):=\operatorname{Tr}(S\sigma).
\]

\begin{lemma}[Visibility perturbation]
\label{lem:v12v43-visibility}
One has
\begin{equation}
 |v(S,\sigma)-v(S_0,\sigma_0)|
 \leq
 \|S-S_0\|
 +
 \|S_0\|\,\|\sigma-\sigma_0\|_1.                \tag{43.1}
\end{equation}
\end{lemma}

\begin{proof}
Insert \(v(S_0,\sigma)\):
\begin{align*}
 |v(S,\sigma)-v(S_0,\sigma_0)|
 &\leq
 |\operatorname{Tr}\{(S-S_0)\sigma\}|\\
 &\quad+
 |\operatorname{Tr}\{S_0(\sigma-\sigma_0)\}|.
\end{align*}
The first term is at most
\(\|S-S_0\|\|\sigma\|_1=\|S-S_0\|\).
Trace duality bounds the second by
\(\|S_0\|\|\sigma-\sigma_0\|_1\).
\end{proof}

\begin{corollary}[One-box visibility certificate]
\label{cor:v12v43-box-visibility}
Let \(q\) range over a box \(Q\) with centre \(q_0\).  Suppose
\begin{align}
 \sup_{q\in Q}\|S(q)-S(q_0)\|&\leq\eta_S,         \tag{43.2}\\
 \sup_{q\in Q}\|\sigma(q)-\sigma(q_0)\|_1
 &\leq\eta_\sigma.                               \tag{43.3}
\end{align}
Put
\begin{equation}
 \nu_Q:=
 \operatorname{Tr}\{S(q_0)\sigma(q_0)\}
 -\eta_S-\|S(q_0)\|\eta_\sigma.                  \tag{43.4}
\end{equation}
If \(\nu_Q>0\), then
\begin{equation}
 \operatorname{Tr}\{S(q)\sigma(q)\}\geq\nu_Q
 \qquad(q\in Q).                                 \tag{43.5}
\end{equation}
\end{corollary}

\begin{proof}
Apply Lemma \ref{lem:v12v43-visibility} at every \(q\in Q\).
\end{proof}

Unlike the V37 box floor, equation (43.4) is invariant under
multiplication of the unnormalized kernel by an arbitrary positive
scalar.

\subsection{A scale-free shape variation identity}

Let \(A(q),A(q_0)\succeq0\) be nonzero.  Write
\[
 m(q)=\operatorname{Tr}A(q),\qquad
 \sigma(q)={A(q)\over m(q)}.
\]

\begin{proposition}[Radial rescaling removes amplitude variation]
\label{prop:v12v43-radial}
The exact identity
\begin{equation}
 \|\sigma(q)-\sigma(q_0)\|_1
 =
 {1\over m(q)}
 \left\|
 A(q)-{m(q)\over m(q_0)}A(q_0)
 \right\|_1                                      \tag{43.6}
\end{equation}
holds.  Hence the relative angular defect
\begin{equation}
 \left\|
 A(q)-{m(q)\over m(q_0)}A(q_0)
 \right\|_1
 \leq\delta_{\rm ang}\,m(q)                     \tag{43.7}
\end{equation}
implies
\(\|\sigma(q)-\sigma(q_0)\|_1\leq\delta_{\rm ang}\), with no
lower bound on \(m(q)\).
\end{proposition}

\begin{proof}
Since \(A(q)=m(q)\sigma(q)\) and
\(A(q_0)=m(q_0)\sigma(q_0)\), the matrix inside the norm in
(43.6) is
\[
 m(q)\{\sigma(q)-\sigma(q_0)\}.
\]
Divide its trace norm by \(m(q)>0\).
\end{proof}

This is the correct box variation near an amplitude zero.  The
unscaled difference \(\|A(q)-A(q_0)\|_1\) confounds radial mass
change with angular shape change.

\subsection{Comparison when the physical functional also changes}

Let \(A,B\succeq0\) be nonzero, with trace-one shapes
\(\sigma_A,\sigma_B\).  Let their positive normalizer matrices be
\(S_A,S_B\), and put
\[
 v_A=\operatorname{Tr}(S_A\sigma_A),\qquad
 v_B=\operatorname{Tr}(S_B\sigma_B).
\]

\begin{theorem}[Joint shape--functional comparison]
\label{thm:v12v43-joint}
Assume
\[
 v_A,v_B\geq\nu>0.
\]
Set
\[
 \delta_\sigma:=\|\sigma_A-\sigma_B\|_1,\qquad
 \delta_S:=\|S_A-S_B\|,\qquad
 s_*:=\max\{\|S_A\|,\|S_B\|\}.
\]
Then
\begin{equation}
 \left\|
 {A\over\operatorname{Tr}(S_AA)}
 -
 {B\over\operatorname{Tr}(S_BB)}
 \right\|_1
 \leq
 {\delta_\sigma\over\nu}
 +
 {\delta_S+s_*\delta_\sigma\over\nu^2}.          \tag{43.8}
\end{equation}
\end{theorem}

\begin{proof}
The mass--shape identity gives
\[
 {A\over\operatorname{Tr}(S_AA)}
 =
 {\sigma_A\over v_A},
 \qquad
 {B\over\operatorname{Tr}(S_BB)}
 =
 {\sigma_B\over v_B}.
\]
As in V42,
\begin{equation}
 \left\|{\sigma_A\over v_A}
       -{\sigma_B\over v_B}\right\|_1
 \leq
 {\delta_\sigma\over v_A}
 +
 {|v_A-v_B|\over v_Av_B}.                       \tag{43.9}
\end{equation}
Furthermore,
\begin{align*}
 |v_A-v_B|
 &\leq
 |\operatorname{Tr}\{(S_A-S_B)\sigma_A\}|\\
 &\quad+
 |\operatorname{Tr}\{S_B(\sigma_A-\sigma_B)\}|\\
 &\leq\delta_S+s_*\delta_\sigma.
\end{align*}
Insert this and the lower visibility \(\nu\) in (43.9).
\end{proof}

Equation (43.8) separates two interscale owners: change of the
normalized reflected kernel shape and change of the physical
partition functional.  Neither can be omitted.

\subsection{A complete local adjacent certificate}

Let \(A_n(q),B_n(q)\) be the two adjacent coarse kernels after
transport to a common boundary space, and let their normalizer
matrices be \(S_{A,n}(q),S_{B,n}(q)\).  On a box \(Q\), certify:

\begin{enumerate}
\item angular defects from (43.6)--(43.7) for both kernel families;
\item lower visibilities
\(\nu_{A,Q},\nu_{B,Q}>0\) from (43.4);
\item a direct or triangular bound
\[
 \|\sigma_{A,n}(q)-\sigma_{B,n}(q)\|_1
 \leq\delta_{\sigma,Q};
\]
\item the functional transport error
\[
 \|S_{A,n}(q)-S_{B,n}(q)\|\leq\delta_{S,Q};
\]
and
\item \(s_Q\geq\max\{\|S_{A,n}(q)\|,\|S_{B,n}(q)\|\}\).
\end{enumerate}

With
\[
 \nu_Q=\min\{\nu_{A,Q},\nu_{B,Q}\},
\]
Theorem \ref{thm:v12v43-joint} yields
\begin{equation}
 e_{n,Q}
 =
 {\delta_{\sigma,Q}\over\nu_Q}
 +
 {\delta_{S,Q}+s_Q\delta_{\sigma,Q}\over\nu_Q^2}.
                                                        \tag{43.10}
\end{equation}
This replaces the absolute normalizer floor and unnormalized
operator defect in the V37 certificate.

\subsection{Termination on a compact visible shape image}

\begin{proposition}[Shape-box subdivision terminates]
\label{prop:v12v43-termination}
Let \(G\) be compact.  Suppose \(q\mapsto S(q)\) and
\(q\mapsto\sigma(q)\) are continuous on \(G\), and
\begin{equation}
 \inf_{q\in G}\operatorname{Tr}\{S(q)\sigma(q)\}
 =\nu_*>0.                                       \tag{43.11}
\end{equation}
Suppose the certified variation bounds in (43.2)--(43.3) tend to
zero with box diameter.  Then sufficiently fine recursive
subdivision certifies \(\nu_Q\geq\nu_*/2\) on every box of a finite
cover of \(G\).
\end{proposition}

\begin{proof}
Continuity on a compact set is uniform.  Choose the box diameter so
that
\[
 \eta_S+\sup_{q\in G}\|S(q)\|\eta_\sigma<\nu_*/2.
\]
Every centre has visibility at least \(\nu_*\), so (43.4) gives
\(\nu_Q\geq\nu_*/2\).  Compactness supplies a finite cover, and
recursive bisection eventually reaches the chosen diameter.
\end{proof}

The substantive hypothesis is (43.11).  Subdivision detects a strict
margin but cannot prove one when the shape image meets the invisible
face.

\subsection{Cofinal shape budget}

If the integrated or cylinder-local versions of (43.10) give
adjacent errors \(e_n\) satisfying
\[
                         \sum_ne_n<\infty,        \tag{43.12}
\]
then V36 gives an exactly compatible limiting family of normalized
positive boundary states.  Amplitude decay of the unnormalized
kernels is irrelevant to this sum.  The required inputs are instead:
\begin{enumerate}
\item summable angular shape drift;
\item summable transport error of the physical normalizer matrix; and
\item a noncollapsing visibility margin, or a separately controlled
small-visibility tail.
\end{enumerate}

\begin{assessment}[Status after V43]
\label{ass:v12v43-status}
Shape visibility is now locally certifiable when the physical
partition functional moves.  The exact radial-rescaling identity
isolates angular kernel drift without a trace-mass floor, and the
joint comparison theorem adds the previously missing functional
transport error.  Compact visible shape images admit terminating
box covers and a direct summable projective budget.
\end{assessment}

\begin{firewall}[Claim boundary after V43]
\label{fire:v12v43-claim}
No RPESM angular defect, moving normalizer matrix, visibility margin,
or functional transport error has been populated.  Continuity of
the trace-one shape through rank-changing zeros is not automatic.
The missing Duhamel fibre, determinant-line gluing, gravitational
tightness, and the full SM--Einstein continuum remain open.
\end{firewall}

\subsection{Parent record}

This V43 note appends to the validated V42 file with record
\[
\begin{array}{c|c}
\text{lines}&12303\\
\text{bytes}&415865\\
\text{SHA--256}&
\texttt{0C17731604FBB0A8AD3E4172AB378756427066FD5C7130904872DDE5E314C187}.
\end{array}
\]

% =============================================================================
% END APPENDED TWELFTH-CYCLE V43 MATERIAL
% =============================================================================
% =============================================================================
% BEGIN APPENDED TWELFTH-CYCLE V44 MATERIAL
% =============================================================================

\section{Twelfth-cycle V44: correction and shape extension through a
zero positive kernel}
\label{sec:v12v44}

\subsection{Correction to the V39 slice counterexample}

The proof of Proposition \ref{prop:v12v35-no-tensorization} uses
\[
 P(z_1,z_2)=z_1+z_2-3.
\]
Its final sentence said that any domain containing the unit disk and
\(3/2\) gives the counterexample.  That statement is too broad
because such a domain might also contain \(2\), where the coordinate
slice \(P(z,1)=z-2\) vanishes.  The required precise choice is
\begin{equation}
                         D=\{z\in\mathbb C:|z|\leq7/4\}.    \tag{44.1}
\end{equation}
Then \(P(z,1)=P(1,z)=z-2\) is nonzero on \(D\), while
\((3/2,3/2)\in D^2\) and \(P(3/2,3/2)=0\).  With this correction,
the proposition and its intended conclusion remain valid.

\subsection{The zero-kernel shape problem}

Let \(X\) be a metric parameter space, let \(x_0\in X\), and let
\[
 A:X\longrightarrow{\rm Herm}_d^+
\]
be continuous.  Assume \(A(x_0)=0\) and \(A(x)\ne0\) on a punctured
neighbourhood.  Put
\[
 m(x)=\operatorname{Tr}A(x),\qquad
 \sigma(x)={A(x)\over m(x)}
 \quad(x\ne x_0).
\]
The mass tends to zero automatically.  The shape need not have a
limit.

\begin{theorem}[Exact shape-extension criterion]
\label{thm:v12v44-extension}
A trace-one matrix \(\sigma_0\succeq0\) is the limit of
\(\sigma(x)\) as \(x\to x_0\) if and only if
\begin{equation}
 \|A(x)-m(x)\sigma_0\|_1=o(m(x))
 \qquad(x\to x_0,\ x\ne x_0).                    \tag{44.2}
\end{equation}
When it exists, the limiting shape is unique.
\end{theorem}

\begin{proof}
The exact identity
\[
 A(x)-m(x)\sigma_0
 =
 m(x)\{\sigma(x)-\sigma_0\}
\]
gives
\[
 {\|A(x)-m(x)\sigma_0\|_1\over m(x)}
 =
 \|\sigma(x)-\sigma_0\|_1.
\]
Thus (44.2) is equivalent to trace-norm convergence of the shapes.
Uniqueness follows from uniqueness of limits in a metric space.
\end{proof}

This is a relative leading-direction condition.  Absolute continuity
of \(A\) at its zero gives no information about it.

\subsection{Directional leading forms}

Assume now that \(X\) is a neighbourhood of the origin in
\(\mathbb R^n\), \(x_0=0\), and that for some integer \(k\geq1\)
there is a continuous positive matrix field \(H\) on the unit sphere
such that
\begin{equation}
 A(r\omega)
 =
 r^kH(\omega)+R(r,\omega),                       \tag{44.3}
\end{equation}
with
\begin{equation}
 \sup_\omega{\|R(r,\omega)\|_1\over r^k}\to0,
 \qquad
 \inf_\omega\operatorname{Tr}H(\omega)>0          \tag{44.4}
\end{equation}
as \(r\downarrow0\), over all directions belonging to the domain.

\begin{theorem}[Normalized leading-form criterion]
\label{thm:v12v44-leading}
For every fixed direction \(\omega\),
\begin{equation}
 \sigma(r\omega)
 \longrightarrow
 {H(\omega)\over\operatorname{Tr}H(\omega)}.
                                                        \tag{44.5}
\end{equation}
The shape extends continuously to the zero kernel if and only if
there is one density matrix \(\sigma_0\) such that
\begin{equation}
 {H(\omega)\over\operatorname{Tr}H(\omega)}
 =
 \sigma_0
 \quad\text{for every admissible }\omega.         \tag{44.6}
\end{equation}
\end{theorem}

\begin{proof}
Taking traces in (44.3) and using
\(|\operatorname{Tr}R|\leq\|R\|_1\) gives
\[
 m(r\omega)
 =
 r^k\operatorname{Tr}H(\omega)+o(r^k),
\]
uniformly in \(\omega\).  The positive lower trace in (44.4) permits
division, yielding (44.5) uniformly.

If a continuous extension exists, every directional limit equals
its value at the origin, so (44.6) is necessary.  If (44.6) holds,
the uniform version of (44.5) proves convergence to \(\sigma_0\)
along every approach, hence sufficiency.
\end{proof}

\begin{example}[Direction-dependent failure]
\label{ex:v12v44-direction-failure}
For
\[
 A(x,y)=
 \begin{pmatrix}x^2&0\\0&y^2\end{pmatrix},
\]
the shape tends to
\(\operatorname{diag}(1,0)\) along the \(x\)-axis and to
\(\operatorname{diag}(0,1)\) along the \(y\)-axis.  Although \(A\)
is analytic, positive semidefinite, and vanishes quadratically, its
trace-one shape has no limit at the origin.
\end{example}

This example shows why vanishing order alone cannot certify the V43
shape boxes.

\subsection{One analytic parameter}

Let \(t\mapsto A(t)\succeq0\) be real analytic for
\(|t|<\varepsilon\), with \(A(0)=0\) and not identically zero.
Let \(k\) be the first nonzero Taylor order:
\begin{equation}
 A(t)=t^kA_k+O(t^{k+1}).                          \tag{44.7}
\end{equation}

\begin{proposition}[Analytic one-parameter extension]
\label{prop:v12v44-analytic}
If positivity holds for both signs of \(t\), then \(k\) is even,
\(A_k\succeq0\), \(A_k\ne0\), and
\begin{equation}
 {A(t)\over\operatorname{Tr}A(t)}
 \longrightarrow
 {A_k\over\operatorname{Tr}A_k}.                 \tag{44.8}
\end{equation}
For a one-sided parameter \(t\downarrow0\), the same conclusion
holds without the parity assertion, provided the leading sign makes
\(A_k\succeq0\).
\end{proposition}

\begin{proof}
For every vector \(v\), the scalar analytic function
\(v^*A(t)v\) is nonnegative for both signs.  At least one such
function has first nonzero order \(k\).  A nonzero odd leading term
changes sign, so \(k\) is even.  Division of
\(v^*A(t)v\geq0\) by \(t^k>0\) and passage to the limit gives
\(v^*A_kv\geq0\); hence \(A_k\succeq0\).
It is nonzero by definition of \(k\), so its trace is positive.
Taking traces in (44.7) and dividing gives (44.8).  The one-sided
proof is identical once the leading coefficient has the positive
sign.
\end{proof}

Thus an analytic one-parameter shell path has a canonical limiting
shape.  A multidimensional background family can fail because its
normalized leading form depends on direction.

\subsection{Visibility through the zero}

Let \(S(x)\succeq0\) be continuous at \(x_0\).  If the shape extends
to \(\sigma_0\), then
\begin{equation}
 \operatorname{Tr}\{S(x)\sigma(x)\}
 \longrightarrow
 \operatorname{Tr}\{S(x_0)\sigma_0\}.             \tag{44.9}
\end{equation}

\begin{corollary}[Positive limiting visibility]
\label{cor:v12v44-visibility}
If
\begin{equation}
 \operatorname{Tr}\{S(x_0)\sigma_0\}>0,           \tag{44.10}
\end{equation}
then a punctured neighbourhood of \(x_0\) has a uniform positive
shape-visibility bound.  The normalized physical kernels
\[
 {A(x)\over\operatorname{Tr}\{S(x)A(x)\}}
 =
 {\sigma(x)\over\operatorname{Tr}\{S(x)\sigma(x)\}}
\]
extend continuously across \(x_0\).
\end{corollary}

\begin{proof}
Equation (44.9) follows from
\[
 |\operatorname{Tr}\{S(x)\sigma(x)\}
   -\operatorname{Tr}\{S(x_0)\sigma_0\}|
 \leq
 \|S(x)-S(x_0)\|
 +
 \|S(x_0)\|\|\sigma(x)-\sigma_0\|_1.
\]
Both terms tend to zero.  Under (44.10), the denominator is bounded
below nearby.  Continuity of the quotient follows.
\end{proof}

If the limiting shape lies in the invisible face of \(S(x_0)\), the
mass cancellation is insufficient and an angular small-ball or a
finer leading-order calculation is required.

\subsection{Cofinal implication}

At regulator scale \(n\), a zero kernel can be admitted into the V43
box system when one proves:
\begin{enumerate}
\item a relative leading expansion of the form (44.3);
\item direction independence of the normalized leading form;
\item compatibility of that limiting shape with the coarse lift;
\item positive limiting visibility; and
\item summable interscale drift of the extended shapes and moving
normalizer matrices.
\end{enumerate}
These conditions allow amplitude-zero strata to remain in the good
set rather than being discarded into a tail estimate.

\begin{assessment}[Status after V44]
\label{ass:v12v44-status}
The trace-one shape now has an exact extension criterion at a zero
positive kernel.  A one-parameter analytic positive path always has
a normalized leading shape, while a multidimensional family extends
only when every normalized directional leading form agrees.
Positive limiting visibility then extends the physical normalized
kernel without any inverse mass estimate.
\end{assessment}

\begin{firewall}[Claim boundary after V44]
\label{fire:v12v44-claim}
No RPESM zero-kernel stratum, leading matrix, or direction-independent
shape has been computed.  The correction at the start of V44 narrows
only the illustrative domain in V39 and changes no programme
conclusion.  The missing Duhamel fibre, determinant-line gluing,
gravitational tightness, and the full SM--Einstein continuum remain
open.
\end{firewall}

\subsection{Parent record}

This V44 note appends to the validated V43 file with record
\[
\begin{array}{c|c}
\text{lines}&12620\\
\text{bytes}&424990\\
\text{SHA--256}&
\texttt{BEFB58A4303D7E8E854F9A55DC5DE39926DC20D58F381229D023A11742432E67}.
\end{array}
\]

% =============================================================================
% END APPENDED TWELFTH-CYCLE V44 MATERIAL
% =============================================================================
% =============================================================================
% BEGIN APPENDED TWELFTH-CYCLE V45 MATERIAL
% =============================================================================

\section{Twelfth-cycle V45: compulsory companion audit and a
packet-first shape route}
\label{sec:v12v45}

\subsection{Audit identity}

This is the compulsory five-version audit following V40.  The newest
active companion files at the audit time were
\[
\begin{array}{c|r|r}
\text{file}&\text{counted lines}&\text{bytes}\\ \hline
\texttt{Renewal\_Yang\_Mills\_Mass\_Gap\_Research\_Notes\_V55.tex}
 &20960&714665\\
\texttt{Renewal\_NS\_Stability\_Research\_Notes\_V26.tex}
 &5320&278283.
\end{array}                                                   \tag{45.1}
\]
Their SHA--256 fingerprints are
\[
\begin{split}
 &\texttt{9D7B7A75BCF6713BCC41A1F224C6EDD96AA6384758C1C8DD3274701A9C1EF207},
 \\
 &\texttt{82C887F8C2C69E4F28FAD7C3DC8DE5B9099CB5A4BF431EBA1FFF3E91935978AD}.
\end{split}                                                   \tag{45.2}
\]
They were read without modification.

\subsection{Yang--Mills V55: signed core before absolute tail}

Yang--Mills V55 performs its own mandatory audit, retains the legal
shift-packet trace, and proves the exact Poisson identity
\[
 G_M(F)-I(F)
 =
 \sum_{\ell\ne0}\widehat F(M\ell).
\]
It then keeps a finite signed set of alias coefficients and applies
an absolute Wiener bound only to the remaining tail.  This evades
the method-specific no-go theorem in its V54, which applied absolute
values to the whole alias sum.

The result supplies no numerical SM--Einstein estimate, but it
confirms the correct proof order:
\begin{equation}
 \boxed{
 \text{complete legal packet}
 \longrightarrow
 \text{matrix or scalar coefficient}
 \longrightarrow
 \text{finite signed core}
 \longrightarrow
 \text{absolute residual tail}.}                \tag{45.3}
\end{equation}
Taking norms before the legal packet sum can erase exact cancellation
and can also violate the covariance correction of V25--V26.

\subsection{Transfer to reflected-kernel shapes}

The object required by V42--V44 is matrix valued: a trace-one
positive reflected-kernel shape.  The scalar Yang--Mills theorem
cannot simply be cited for it because:
\begin{enumerate}
\item convergence must hold in trace norm;
\item positivity and unit trace of the continuum and grid averages
must be retained;
\item the finite alias core is a Hermitian matrix whose cancellations
occur before its trace norm is taken; and
\item a moving physical normalizer matrix contributes the additional
V43 functional-transport term.
\end{enumerate}

These are the targets of V46.  The proof will use Bochner integration
in the trace class and a weighted trace-norm Fourier series.

\subsection{Navier--Stokes status}

Navier--Stokes V26 is unchanged.  Its restricted Oseen norm again
supports the rule that structure-sensitive estimates require a
proved structural input.  It gives no Fourier, reflected-kernel,
normalizer, or continuum-measure result for the present step.

\subsection{Direction after the audit}

Let
\[
 \Sigma(k)\in{\mathfrak S}_1({\cal H}),\qquad
 k\in{\mathbb T}^D,
\]
be a positive trace-one reflected shape after the full legal momentum
packet has been assembled.  V46 will prove:
\begin{enumerate}
\item exact matrix-valued Poisson congruence in trace class;
\item a finite Hermitian signed core and exponential trace-norm tail;
\item positivity and unit trace of the grid and continuum averages;
and
\item insertion of the resulting shape error into the V43
moving-normalizer estimate.
\end{enumerate}

\begin{assessment}[Status after V45]
\label{ass:v12v45-status}
The companion audit confirms that cancellation must be retained
through the complete invariant packet and finite core.  The next SM
step is a trace-class version for positive kernel shapes, which will
connect exact momentum arithmetic to the mass-free projective error
of V42--V43.
\end{assessment}

\begin{firewall}[Claim boundary after V45]
\label{fire:v12v45-claim}
No Yang--Mills continuum response or mass gap and no Navier--Stokes
regularity theorem is imported.  No RPESM trace-class Fourier stock
or signed shape core has been computed.  The missing Duhamel fibre
and the full SM--Einstein continuum remain open.
\end{firewall}

\subsection{Parent record}

This V45 note appends to the validated V44 file with record
\[
\begin{array}{c|c}
\text{lines}&12898\\
\text{bytes}&434038\\
\text{SHA--256}&
\texttt{81F1121708B9CDA7069DC16F7D924C110A419E5481DA240A0ADC1747FEB69785}.
\end{array}
\]

% =============================================================================
% END APPENDED TWELFTH-CYCLE V45 MATERIAL
% =============================================================================
% =============================================================================
% BEGIN APPENDED TWELFTH-CYCLE V46 MATERIAL
% =============================================================================

\section{Twelfth-cycle V46: signed Poisson cores for trace-class
reflected shapes}
\label{sec:v12v46}

V45 imports the cancellation-preserving order of the Yang--Mills
V55 argument but not its scalar conclusion.  The SM--Einstein object
is a positive trace-one boundary-kernel shape.  This note proves the
Poisson identity in the trace class, retains a finite Hermitian alias
core before taking a norm, and couples the result to the moving
physical normalizer of V43.

\subsection{Trace-class Fourier data}

Let \({\cal H}\) be a separable Hilbert space and let
\[
 \Sigma:{\mathbb T}^D\longrightarrow{\mathfrak S}_1({\cal H})
\]
be Bochner measurable.  Normalize Haar measure on the torus to mass
one.  Its trace-class Fourier coefficients are
\begin{equation}
 \widehat\Sigma(n)
 :=
 \int_{{\mathbb T}^D}
 e^{-in\cdot k}\Sigma(k)\,dk,
 \qquad n\in\mathbb Z^D.                          \tag{46.1}
\end{equation}
Assume the weighted Wiener stock
\begin{equation}
 {\cal W}_\rho(\Sigma)
 :=
 \sum_{n\in\mathbb Z^D}
 e^{\rho|n|_1}\|\widehat\Sigma(n)\|_1
 <\infty                                         \tag{46.2}
\end{equation}
for some \(\rho>0\).  Then the Fourier series converges absolutely
and uniformly in trace norm.

Define the continuum and \(M^D\)-grid averages by
\begin{align}
 I(\Sigma)&:=\int_{{\mathbb T}^D}\Sigma(k)\,dk,    \tag{46.3}\\
 G_M(\Sigma)&:=
 {1\over M^D}
 \sum_{j\in\{0,\ldots,M-1\}^D}
 \Sigma(2\pi j/M).                               \tag{46.4}
\end{align}

\begin{theorem}[Trace-class Poisson congruence]
\label{thm:v12v46-Poisson}
Under (46.2),
\begin{align}
 I(\Sigma)&=\widehat\Sigma(0),                   \tag{46.5}\\
 G_M(\Sigma)&=
 \sum_{\ell\in\mathbb Z^D}\widehat\Sigma(M\ell), \tag{46.6}\\
 G_M(\Sigma)-I(\Sigma)&=
 \sum_{\ell\ne0}\widehat\Sigma(M\ell).           \tag{46.7}
\end{align}
Every series converges absolutely in trace norm.
\end{theorem}

\begin{proof}
Absolute uniform trace-norm convergence permits termwise Bochner
integration and finite grid summation.  The torus integral kills
every nonzero character, proving (46.5).  On the grid,
\[
 {1\over M^D}
 \sum_j e^{2\pi i n\cdot j/M}
 =
 \prod_{r=1}^D
 \left\{
 {1\over M}\sum_{j_r=0}^{M-1}
 e^{2\pi i n_rj_r/M}
 \right\}.
\]
The product equals one exactly when every component of \(n\) is
divisible by \(M\), and zero otherwise.  This proves (46.6);
subtracting (46.5) gives (46.7).
\end{proof}

\subsection{Positivity and the Hermitian core}

Assume now that
\begin{equation}
 \Sigma(k)\succeq0,\qquad
 \operatorname{Tr}\Sigma(k)=1
 \quad\text{for every }k.                        \tag{46.8}
\end{equation}

\begin{proposition}[The two averages are kernel shapes]
\label{prop:v12v46-positive-averages}
Both \(I(\Sigma)\) and \(G_M(\Sigma)\) are positive trace-class
operators of trace one.  Moreover,
\begin{equation}
 \widehat\Sigma(-n)=\widehat\Sigma(n)^*,\qquad
 \operatorname{Tr}\widehat\Sigma(n)=0
 \quad(n\ne0).                                   \tag{46.9}
\end{equation}
\end{proposition}

\begin{proof}
Positive trace-class operators form a closed convex cone, so the
Bochner integral and finite average in (46.3)--(46.4) are positive.
Linearity of the trace and (46.8) give unit trace.  Taking adjoints
in (46.1) proves the first relation in (46.9).  Taking traces in
(46.1) reduces the second relation to the integral of a nonzero
torus character.
\end{proof}

For an integer \(L\geq0\), define the finite signed matrix core
\begin{equation}
 C_{M,L}(\Sigma)
 :=
 \sum_{0<|\ell|_1\leq L}\widehat\Sigma(M\ell).    \tag{46.10}
\end{equation}
The summation set is symmetric under \(\ell\mapsto-\ell\).

\begin{corollary}[Hermitian traceless core]
\label{cor:v12v46-core}
Under (46.8), \(C_{M,L}(\Sigma)\) is Hermitian and traceless.
\end{corollary}

\begin{proof}
Pair the terms at \(\ell\) and \(-\ell\) using (46.9).  Each
nonzero Fourier coefficient is traceless.
\end{proof}

The core need not be positive.  Positivity belongs to the complete
averages, not to a truncated alias correction.

\subsection{Cancellation-preserving trace-norm tail}

\begin{theorem}[Trace-class core--tail estimate]
\label{thm:v12v46-core-tail}
Under (46.2),
\begin{equation}
 \left\|
 G_M(\Sigma)-I(\Sigma)-C_{M,L}(\Sigma)
 \right\|_1
 \leq
 e^{-\rho M(L+1)}{\cal W}_\rho(\Sigma).           \tag{46.11}
\end{equation}
Consequently,
\begin{equation}
 \|G_M(\Sigma)-I(\Sigma)\|_1
 \leq
 \|C_{M,L}(\Sigma)\|_1
 +
 e^{-\rho M(L+1)}{\cal W}_\rho(\Sigma).           \tag{46.12}
\end{equation}
\end{theorem}

\begin{proof}
By Theorem \ref{thm:v12v46-Poisson}, the expression before the norm
in (46.11) is the sum over \(|\ell|_1\geq L+1\).  Hence
\begin{align*}
 \left\|
 \sum_{|\ell|_1\geq L+1}\widehat\Sigma(M\ell)
 \right\|_1
 &\leq
 \sum_{|\ell|_1\geq L+1}
 \|\widehat\Sigma(M\ell)\|_1\\
 &\leq
 e^{-\rho M(L+1)}
 \sum_{|\ell|_1\geq L+1}
 e^{\rho|M\ell|_1}\|\widehat\Sigma(M\ell)\|_1\\
 &\leq
 e^{-\rho M(L+1)}{\cal W}_\rho(\Sigma).
\end{align*}
This proves (46.11); the triangle inequality gives (46.12).
\end{proof}

The trace norm is taken only after the finite matrices in (46.10)
have been summed.  Thus cancellations among signs, tensor channels,
and opposite frequencies inside the core are retained exactly.

\begin{corollary}[Interval-computed core]
\label{cor:v12v46-interval-core}
Suppose a Hermitian finite matrix \(\widetilde C_{M,L}\) satisfies
\[
 \|C_{M,L}(\Sigma)-\widetilde C_{M,L}\|_1
 \leq\eta_{\rm core}.
\]
Then
\begin{equation}
 \|G_M(\Sigma)-I(\Sigma)\|_1
 \leq
 \|\widetilde C_{M,L}\|_1+\eta_{\rm core}
 +
 e^{-\rho M(L+1)}{\cal W}_\rho(\Sigma).           \tag{46.13}
\end{equation}
\end{corollary}

\begin{proof}
Insert \(\widetilde C_{M,L}\) in (46.12) and apply the triangle
inequality once to the complete finite core.
\end{proof}

Entrywise absolute values taken before assembling
\(\widetilde C_{M,L}\) would discard the cancellation this
corollary is designed to preserve.

\subsection{Physical normalization of the two averages}

Let \(S_I,S_G\succeq0\) be the continuum and grid physical
normalizer matrices.  Put
\[
 v_I=\operatorname{Tr}\{S_II(\Sigma)\},\qquad
 v_G=\operatorname{Tr}\{S_GG_M(\Sigma)\}.
\]
Let
\begin{equation}
 \delta_{M,L}:=
 \|C_{M,L}(\Sigma)\|_1
 +
 e^{-\rho M(L+1)}{\cal W}_\rho(\Sigma).           \tag{46.14}
\end{equation}
Then (46.12) gives
\(\|G_M(\Sigma)-I(\Sigma)\|_1\leq\delta_{M,L}\).

\begin{theorem}[Visibility-preserving physical comparison]
\label{thm:v12v46-physical}
Assume
\begin{equation}
 v_G\geq v_0>0,\qquad
 \|S_I-S_G\|\leq\delta_S,\qquad
 s_*:=\max\{\|S_I\|,\|S_G\|\}.                    \tag{46.15}
\end{equation}
If
\begin{equation}
 \nu:=
 v_0-\delta_S-s_*\delta_{M,L}>0,                 \tag{46.16}
\end{equation}
then \(v_I\geq\nu\) and
\begin{align}
 \left\|
 {I(\Sigma)\over v_I}
 -
 {G_M(\Sigma)\over v_G}
 \right\|_1
 \leq{}&
 {\delta_{M,L}\over\nu}
 +
 {\delta_S+s_*\delta_{M,L}\over\nu^2}.           \tag{46.17}
\end{align}
\end{theorem}

\begin{proof}
The visibility perturbation estimate used in V43 gives
\[
 |v_I-v_G|
 \leq
 \delta_S+s_*\|I(\Sigma)-G_M(\Sigma)\|_1
 \leq
 \delta_S+s_*\delta_{M,L}.
\]
Equations (46.15)--(46.16) yield \(v_I\geq\nu\), while
\(v_G\geq v_0\geq\nu\).  Apply Theorem
\ref{thm:v12v43-joint} with shape difference
\(\delta_{M,L}\).
\end{proof}

This theorem uses no unnormalized trace-mass floor.  The inputs are
a signed matrix core, a residual trace-class Wiener stock, and
visibility of the trace-one averages.

\subsection{Legal packet requirement}

The function \(\Sigma(k)\) must be formed only after all momentum
cards connected by the background-shift subgroup have been retained
and traced in the legal invariant packet.  If an individual card is
Fourier transformed first, its coefficients need not transform
covariantly, and cancellations in (46.10) have no invariant meaning.
This is the matrix-valued form of the ordering rule confirmed in
V25--V26 and YM V50--V55.

\subsection{Cofinal shape criterion}

At regulator scale \(n\), choose \(M_n,L_n,\rho_n\) and define the
physical comparison error by the right side of (46.17), including
any interval-core error from (46.13).  If
\begin{equation}
 \sum_n
 \left[
 {\delta_{n}\over\nu_n}
 +
 {\delta_{S,n}+s_n\delta_n\over\nu_n^2}
 \right]
 <\infty,                                        \tag{46.18}
\end{equation}
then the V36 trace-state repair produces an exactly compatible
limiting family.  This condition allows the grid sizes to stay
small when finitely many signed alias coefficients capture the
dominant discrepancy and the residual weighted tail is summable.

\begin{assessment}[Status after V46]
\label{ass:v12v46-status}
The cancellation-preserving Poisson method now applies to
trace-class reflected-kernel shapes.  Continuum and grid averages
remain positive and trace one; the finite alias core is Hermitian
and traceless; absolute values enter only in the residual tail.
The resulting shape error combines with a moving physical
normalizer through an explicit visibility margin and a cofinal
summability criterion.
\end{assessment}

\begin{firewall}[Claim boundary after V46]
\label{fire:v12v46-claim}
No RPESM packet-first shape function, Fourier coefficient, signed
matrix core, trace-class Wiener stock, or visibility margin has been
computed.  Absolute convergence in (46.2) is a hypothesis.  The note
does not supply the missing physical Duhamel fibre, determinant-line
gluing, gravitational tightness, or the full SM--Einstein continuum.
\end{firewall}

\subsection{Parent record}

This V46 note appends to the validated V45 file with record
\[
\begin{array}{c|c}
\text{lines}&13033\\
\text{bytes}&438816\\
\text{SHA--256}&
\texttt{1CE2788CE33D281C12A3031943CA977EA35BBFE8EA51DD886B23050FC54ABE61}.
\end{array}
\]

% =============================================================================
% END APPENDED TWELFTH-CYCLE V46 MATERIAL
% =============================================================================
% =============================================================================
% BEGIN APPENDED TWELFTH-CYCLE V47 MATERIAL
% =============================================================================

\section{Twelfth-cycle V47: polystrip control of the trace-class
Wiener stock}
\label{sec:v12v47}

V46 assumes a weighted trace-class Wiener stock for the packet-first
kernel shape.  This note derives that stock from a wider
Banach-holomorphic polystrip.  A separate statement removes a common
analytic amplitude factor before division, so zeros of the
unnormalized kernel mass do not falsely collapse the strip.

\subsection{Banach-holomorphic Fourier decay}

Let \({\cal B}\) be a complex Banach space and define the closed
periodic polystrip
\[
 {\cal P}_R
 :=
 \left\{
 z=k+iy\in\mathbb C^D/(2\pi\mathbb Z)^D:
 |y_j|\leq R,\ 1\leq j\leq D
 \right\}.
\]
Let \(F:{\cal P}_R\to{\cal B}\) be continuous on the closed
polystrip, holomorphic in its interior, and \(2\pi\)-periodic in
each coordinate.  Put
\[
 M_R:=\sup_{z\in{\cal P}_R}\|F(z)\|_{\cal B}.
\]
Its Bochner Fourier coefficients on the real torus are
\[
 \widehat F(n)
 =
 \int_{{\mathbb T}^D}e^{-in\cdot k}F(k)\,dk.
\]

\begin{theorem}[Banach-valued polystrip coefficient bound]
\label{thm:v12v47-coefficients}
For every \(n\in\mathbb Z^D\),
\begin{equation}
 \|\widehat F(n)\|_{\cal B}
 \leq
 M_Re^{-R|n|_1}.                                  \tag{47.1}
\end{equation}
Consequently, for every \(0\leq\rho<R\),
\begin{equation}
 \sum_{n\in\mathbb Z^D}
 e^{\rho|n|_1}\|\widehat F(n)\|_{\cal B}
 \leq
 M_R
 \coth^D\left({R-\rho\over2}\right).              \tag{47.2}
\end{equation}
\end{theorem}

\begin{proof}
Fix \(n\).  For each coordinate with \(n_j\ne0\), shift the
integration contour from the real circle to
\[
 \operatorname{Im}z_j=-R\,\operatorname{sgn}(n_j).
\]
Banach-valued Cauchy theory is justified by applying every continuous
linear functional on \({\cal B}\), using the scalar contour theorem,
and then invoking Hahn--Banach.  Periodicity cancels the vertical
sides.  On the shifted torus,
\[
 |e^{-in\cdot z}|=e^{-R|n|_1}.
\]
The normalized Haar integral therefore gives (47.1).

Let \(\delta=R-\rho>0\).  Summing (47.1) yields
\begin{align*}
 \sum_ne^{\rho|n|_1}\|\widehat F(n)\|_{\cal B}
 &\leq
 M_R\sum_{n\in\mathbb Z^D}e^{-\delta|n|_1}\\
 &=
 M_R
 \left(
 1+2\sum_{m=1}^\infty e^{-\delta m}
 \right)^D\\
 &=
 M_R
 \left({1+e^{-\delta}\over1-e^{-\delta}}\right)^D
 =
 M_R\coth^D(\delta/2).
\end{align*}
\end{proof}

If only the open strip is controlled, the same proof applies with
every \(R'<R\); no bound on the unattained boundary is implied.

\subsection{Trace-class specialization}

Take
\({\cal B}={\mathfrak S}_1({\cal H})\).  If
\(\Sigma(z)\) satisfies the hypotheses of Theorem
\ref{thm:v12v47-coefficients}, then the V46 Wiener stock obeys
\begin{equation}
 {\cal W}_\rho(\Sigma)
 \leq
 M_R^{(1)}
 \coth^D\left({R-\rho\over2}\right),
 \qquad
 M_R^{(1)}:=\sup_{{\cal P}_R}\|\Sigma(z)\|_1.      \tag{47.3}
\end{equation}

\begin{corollary}[Explicit V46 residual tail]
\label{cor:v12v47-tail}
For grid size \(M\), signed-core radius \(L\), and
\(0<\rho<R\),
\begin{equation}
 \left\|
 G_M(\Sigma)-I(\Sigma)-C_{M,L}(\Sigma)
 \right\|_1
 \leq
 M_R^{(1)}
 e^{-\rho M(L+1)}
 \coth^D\left({R-\rho\over2}\right).              \tag{47.4}
\end{equation}
The convenient symmetric choice \(\rho=R/2\) gives
\begin{equation}
 \left\|
 G_M-I-C_{M,L}
 \right\|_1
 \leq
 M_R^{(1)}
 e^{-RM(L+1)/2}
 \coth^D(R/4).                                   \tag{47.5}
\end{equation}
\end{corollary}

\begin{proof}
Insert (47.3) into Theorem
\ref{thm:v12v46-core-tail}.  The second formula is the stated choice
of \(\rho\).
\end{proof}

If every \(\Sigma(z)\) has rank at most \(r\) and an operator-norm
bound \(M_R^{(\infty)}\), then
\[
 M_R^{(1)}\leq rM_R^{(\infty)}.
\]
A rank bound must be valid on the complex polystrip; real-axis rank
alone does not control complex trace norm.

\subsection{Analytic normalization away from mass zeros}

Let \(A:{\cal P}_R\to{\mathfrak S}_1({\cal H})\) be holomorphic and
suppose
\begin{equation}
 \sup_{{\cal P}_R}\|A(z)\|_1\leq M_A,\qquad
 \inf_{{\cal P}_R}|\operatorname{Tr}A(z)|\geq m_A>0.
                                                        \tag{47.6}
\end{equation}
Then
\[
 \Sigma(z):={A(z)\over\operatorname{Tr}A(z)}
\]
is holomorphic and
\begin{equation}
                         M_R^{(1)}\leq {M_A\over m_A}.      \tag{47.7}
\end{equation}

A sufficient perturbative denominator test is the following.  If
\[
 \inf_{k\in{\mathbb T}^D}\operatorname{Tr}A(k)\geq m_0>0
\]
on the positive real axis and
\begin{equation}
 \sup_{\substack{k\in{\mathbb T}^D\\|y_j|\leq R}}
 \|A(k+iy)-A(k)\|_1\leq\eta<m_0,                 \tag{47.8}
\end{equation}
then
\[
 |\operatorname{Tr}A(k+iy)|
 \geq m_0-\eta.
\]
Thus (47.6) holds with \(m_A=m_0-\eta\), provided the numerator
stock is also bounded.

This route is unsuitable when the trace mass genuinely vanishes on
the real domain.  The next factorization handles removable zeros.

\subsection{Removal of a common analytic amplitude}

Suppose
\begin{equation}
                         A(z)=g(z)C(z),           \tag{47.9}
\end{equation}
where \(g\) is scalar holomorphic and \(C\) is trace-class
holomorphic on \({\cal P}_R\).  The function \(g\) may vanish.
Assume
\begin{equation}
 \sup_{{\cal P}_R}\|C(z)\|_1\leq M_C,\qquad
 \inf_{{\cal P}_R}|\operatorname{Tr}C(z)|\geq m_C>0.
                                                        \tag{47.10}
\end{equation}

\begin{theorem}[Analytic common-factor cancellation]
\label{thm:v12v47-common-factor}
On the set where \(g\ne0\),
\begin{equation}
 {A(z)\over\operatorname{Tr}A(z)}
 =
 {C(z)\over\operatorname{Tr}C(z)}.               \tag{47.11}
\end{equation}
The right side is the unique holomorphic extension of the kernel
shape through every zero of \(g\).  Its weighted Wiener stock
satisfies
\begin{equation}
 {\cal W}_\rho(\Sigma)
 \leq
 {M_C\over m_C}
 \coth^D\left({R-\rho\over2}\right),
 \qquad 0\leq\rho<R.                             \tag{47.12}
\end{equation}
\end{theorem}

\begin{proof}
Equation (47.11) follows by canceling \(g(z)\).  The denominator in
the right side is zero-free by (47.10), so the quotient is
holomorphic on the entire polystrip and bounded in trace norm by
\(M_C/m_C\).  The identity theorem gives uniqueness on every
connected component meeting \(\{g\ne0\}\).  Apply Theorem
\ref{thm:v12v47-coefficients} in the trace class to obtain (47.12).
\end{proof}

The factor \(g\) can contain the entire amplitude vanishing isolated
in V41--V44.  No inverse moment or complex lower bound for \(g\) is
required.

\subsection{Direction-dependent zeros remain obstructive}

If an unnormalized kernel has a multidirectional leading form whose
normalized shape depends on direction, as in V44, there is no single
holomorphic \(C/\operatorname{Tr}C\) extending the real shape.
Equation (47.9) with a zero-free core trace would force such an
extension and therefore cannot hold.  The polystrip method detects
the same obstruction in analytic language: removable scalar zeros
are harmless, while angular shape zeros are not.

\subsection{Cofinal allocation}

At scale \(n\), let \(R_n\) be a proved shape polystrip width,
let \(B_n\) bound the analytic trace norm of the normalized core, and
choose grid size \(M_n\) and core radius \(L_n\).  With
\(\rho_n=R_n/2\), the residual tail is at most
\begin{equation}
 t_n:=
 B_n
 \coth^D(R_n/4)
 \exp\{-R_nM_n(L_n+1)/2\}.                       \tag{47.13}
\end{equation}
If the signed finite cores, interval errors, functional-transport
errors, and \(t_n\) produce a summable V46 physical comparison
series, the trace-state projective repair applies.

A sufficient asymptotic condition for the tail alone is
\begin{equation}
 {R_nM_n(L_n+1)\over2}
 -
 \log B_n
 -
 D\log\coth(R_n/4)
 \geq(1+\epsilon)\log(n+2)                       \tag{47.14}
\end{equation}
for some \(\epsilon>0\).  Then \(t_n\leq(n+2)^{-1-\epsilon}\).

\begin{assessment}[Status after V47]
\label{ass:v12v47-status}
The weighted trace-class Wiener stock assumed in V46 now follows
from a wider holomorphic shape polystrip, with an explicit
strip-loss factor.  Common analytic amplitude zeros cancel before
normalization and do not reduce the usable strip.  The cofinal
condition (47.14) states how grid size, signed-core depth, strip
width, and analytic stock must balance to make residual alias tails
summable.
\end{assessment}

\begin{firewall}[Claim boundary after V47]
\label{fire:v12v47-claim}
No RPESM holomorphic shape polystrip, common analytic factor,
trace-class supremum, or cofinal allocation has been constructed.
The theorem applies only after the legal momentum packet is
assembled.  It does not supply the missing physical Duhamel fibre,
determinant-line gluing, gravitational tightness, or the full
SM--Einstein continuum.
\end{firewall}

\subsection{Parent record}

This V47 note appends to the validated V46 file with record
\[
\begin{array}{c|c}
\text{lines}&13365\\
\text{bytes}&448933\\
\text{SHA--256}&
\texttt{4D6741B0DDCE0312C1CF880D0807F4982654EEEB54074A0DECC8FBBB13E61F50}.
\end{array}
\]

% =============================================================================
% END APPENDED TWELFTH-CYCLE V47 MATERIAL
% =============================================================================
% =============================================================================
% BEGIN APPENDED TWELFTH-CYCLE V48 MATERIAL
% =============================================================================

\section{Twelfth-cycle V48: a finite-regularity trace-class
alternative to the analytic Wiener tail}
\label{sec:v12v48}

V47 supplies an exponential residual tail when the packet-first
kernel shape extends holomorphically to a uniform polystrip.  Such an
extension is a strong requirement for nonlinear gravitational
variables.  This note proves a polynomial alternative from finitely
many trace-class momentum derivatives.  The finite signed core of
V46 is unchanged; only the unresolved tail estimate is replaced.

\subsection{Trace-class Sobolev stock}

Let
\[
 \Sigma:{\mathbb T}^D\longrightarrow{\mathfrak S}_1({\cal H})
\]
be Bochner integrable.  Fix an integer \(m\geq1\), and assume the
distributional derivatives in
\[
 (I-\Delta)^m\Sigma
\]
form a Bochner \(L^1\) trace-class function.  Define
\begin{equation}
 {\cal B}_m(\Sigma)
 :=
 \int_{{\mathbb T}^D}
 \|(I-\Delta)^m\Sigma(k)\|_1\,dk.                \tag{48.1}
\end{equation}
Normalized Haar measure is used.

\begin{lemma}[Polynomial Fourier coefficient decay]
\label{lem:v12v48-coefficients}
For every \(n\in\mathbb Z^D\),
\begin{equation}
 \|\widehat\Sigma(n)\|_1
 \leq
 {{\cal B}_m(\Sigma)\over(1+|n|_2^2)^m}.          \tag{48.2}
\end{equation}
\end{lemma}

\begin{proof}
Fourier transformation of the distributional identity gives
\[
 \widehat{(I-\Delta)^m\Sigma}(n)
 =
 (1+|n|_2^2)^m\widehat\Sigma(n).
\]
The trace norm of a Bochner Fourier coefficient is bounded by the
\(L^1\) norm of its integrand.  Divide by
\((1+|n|_2^2)^m\).
\end{proof}

\subsection{Counting the unresolved alias lattice}

For \(r\geq1\), let
\[
 N_D(r):=\#\{\ell\in\mathbb Z^D:|\ell|_1=r\}.
\]

\begin{lemma}[Shell count]
\label{lem:v12v48-shell-count}
For every \(r\geq1\),
\begin{equation}
 N_D(r)
 \leq
 {2^DD^{D-1}\over(D-1)!}\,r^{D-1}.               \tag{48.3}
\end{equation}
\end{lemma}

\begin{proof}
Ignoring the overcount created by zero coordinates, choose signs in
at most \(2^D\) ways and choose nonnegative absolute values summing
to \(r\) in
\(\binom{r+D-1}{D-1}\) ways.  Since
\(r+D-1\leq Dr\) for \(r\geq1\),
\[
 \binom{r+D-1}{D-1}
 \leq{(Dr)^{D-1}\over(D-1)!}.
\]
Multiplication gives (48.3).
\end{proof}

\subsection{Polynomial signed-core tail}

Retain the V46 finite core
\[
 C_{M,L}(\Sigma)
 =
 \sum_{0<|\ell|_1\leq L}\widehat\Sigma(M\ell),
\]
with \(M,L\geq1\).

\begin{theorem}[Sobolev core--tail estimate]
\label{thm:v12v48-core-tail}
If \(2m>D\), then
\begin{equation}
 \left\|
 G_M(\Sigma)-I(\Sigma)-C_{M,L}(\Sigma)
 \right\|_1
 \leq
 C_{D,m}\,
 {{\cal B}_m(\Sigma)\over
  M^{2m}L^{\,2m-D}},                              \tag{48.4}
\end{equation}
where
\begin{equation}
 C_{D,m}:=
 {2^DD^{D-1+m}\over
  (D-1)!(2m-D)}.                                  \tag{48.5}
\end{equation}
Consequently,
\begin{equation}
 \|G_M(\Sigma)-I(\Sigma)\|_1
 \leq
 \|C_{M,L}(\Sigma)\|_1
 +
 C_{D,m}
 {{\cal B}_m(\Sigma)\over
  M^{2m}L^{\,2m-D}}.                              \tag{48.6}
\end{equation}
\end{theorem}

\begin{proof}
The trace-class Poisson identity of V46 makes the expression in
(48.4) equal to the sum over
\(|\ell|_1\geq L+1\).  Since
\[
 |M\ell|_2^2
 \geq {M^2|\ell|_1^2\over D},
\]
Lemma \ref{lem:v12v48-coefficients} gives
\[
 \|\widehat\Sigma(M\ell)\|_1
 \leq
 {\cal B}_m(\Sigma)
 {D^m\over M^{2m}|\ell|_1^{2m}}.
\]
Group the terms by \(r=|\ell|_1\) and apply Lemma
\ref{lem:v12v48-shell-count}:
\begin{align*}
 \sum_{|\ell|_1\geq L+1}
 \|\widehat\Sigma(M\ell)\|_1
 &\leq
 {{\cal B}_m(\Sigma)2^DD^{D-1+m}\over
   (D-1)!M^{2m}}
 \sum_{r=L+1}^\infty r^{D-1-2m}.
\end{align*}
Because \(2m>D\), the integral test yields
\[
 \sum_{r=L+1}^\infty r^{D-1-2m}
 \leq
 \int_L^\infty x^{D-1-2m}\,dx
 =
 {L^{D-2m}\over2m-D}.
\]
This proves (48.4).  The triangle inequality gives (48.6).
\end{proof}

The threshold \(2m>D\) is the point at which this absolute residual
shell sum becomes finite.  Cancellations are still retained inside
the finite core.

\subsection{Positivity and physical visibility}

If
\[
 \Sigma(k)\succeq0,\qquad\operatorname{Tr}\Sigma(k)=1
\]
on the real torus, then V46 already proves that \(I(\Sigma)\) and
\(G_M(\Sigma)\) are positive trace-one shapes.  Define
\begin{equation}
 \delta_{M,L}^{(m)}
 :=
 \|C_{M,L}(\Sigma)\|_1
 +
 C_{D,m}
 {{\cal B}_m(\Sigma)\over
  M^{2m}L^{\,2m-D}}.                              \tag{48.7}
\end{equation}
For physical normalizer matrices \(S_I,S_G\), suppose
\[
 v_G\geq v_0>0,\qquad
 \|S_I-S_G\|\leq\delta_S,\qquad
 s_*=\max\{\|S_I\|,\|S_G\|\}.
\]
If
\begin{equation}
 \nu:=v_0-\delta_S-s_*\delta_{M,L}^{(m)}>0,        \tag{48.8}
\end{equation}
then the proof of Theorem \ref{thm:v12v46-physical} gives
\begin{equation}
 \left\|
 {I(\Sigma)\over v_I}
 -
 {G_M(\Sigma)\over v_G}
 \right\|_1
 \leq
 {\delta_{M,L}^{(m)}\over\nu}
 +
 {\delta_S+s_*\delta_{M,L}^{(m)}\over\nu^2}.      \tag{48.9}
\end{equation}

Thus the analytic and finite-regularity routes feed exactly the same
visibility and projective machinery.

\subsection{Comparison with the analytic route}

The V47 tail is exponential in \(M(L+1)\) but contains a
polystrip-loss factor.  The V48 tail is polynomial and requires
\(2m>D\), but it uses only real-torus derivatives.  Neither dominates
universally.

A mixed sector may use:
\begin{enumerate}
\item the analytic bound for finite gauge, Higgs, or fermionic blocks
with a certified complex inverse strip;
\item the Sobolev bound for metric-dependent shapes with only a
finite differentiability budget; and
\item the same exact signed core wherever low alias coefficients are
computed.
\end{enumerate}
The complete legal shift packet must be assembled before either
regularity stock is taken.

\subsection{Cofinal polynomial allocation}

At scale \(n\), suppose
\[
 {\cal B}_{m,n}\leq C\Lambda^{\gamma n},\qquad
 M_n\geq\Lambda^{an},\qquad
 L_n\geq\Lambda^{bn},
 \qquad\Lambda>1.
\]
Then the residual tail in (48.4) is bounded by a constant times
\[
 \Lambda^{[
 \gamma-2ma-(2m-D)b]n}.
\]

\begin{corollary}[Geometric Sobolev-tail summability]
\label{cor:v12v48-cofinal}
If
\begin{equation}
                         2ma+(2m-D)b>\gamma,      \tag{48.10}
\end{equation}
then the V48 residual tails are summable in \(n\).
\end{corollary}

\begin{proof}
Under (48.10), the displayed exponent is strictly negative, so the
tail is bounded by a convergent geometric series.
\end{proof}

This condition quantifies the price of avoiding complex analyticity:
grid and core growth must beat the growth of the finite derivative
stock.

\begin{warning}[Weaker Schatten control]
\label{warn:v12v48-Schatten}
A Hilbert--Schmidt Sobolev estimate does not imply (48.1) uniformly
when the effective rank grows.  It may be converted to trace norm
only with a certified rank or singular-value tail bound.  Hiding that
factor would invalidate the cofinal estimate.
\end{warning}

\begin{assessment}[Status after V48]
\label{ass:v12v48-status}
The trace-class signed-core method no longer depends exclusively on
a complex polystrip.  More than \(D/2\) powers of the real-torus
Laplacian give an explicit polynomial residual alias bound.  The
same shape-visibility theorem then yields physical comparison, and
condition (48.10) gives a cofinal summability allocation.
\end{assessment}

\begin{firewall}[Claim boundary after V48]
\label{fire:v12v48-claim}
No RPESM trace-class derivative stock, signed core, grid allocation,
or visibility margin has been populated.  Smoothness of finite
regulators alone does not give uniform cofinal bounds.  The missing
physical Duhamel fibre, determinant-line gluing, gravitational
tightness, and the full SM--Einstein continuum remain open.
\end{firewall}

\subsection{Parent record}

This V48 note appends to the validated V47 file with record
\[
\begin{array}{c|c}
\text{lines}&13671\\
\text{bytes}&458290\\
\text{SHA--256}&
\texttt{6A5B21B8299E6DD8142F0DEE73339CDB6D9B6AA4D82D89B807D648201699FD8C}.
\end{array}
\]

% =============================================================================
% END APPENDED TWELFTH-CYCLE V48 MATERIAL
% =============================================================================
% =============================================================================
% BEGIN APPENDED TWELFTH-CYCLE V49 MATERIAL
% =============================================================================

\section{Twelfth-cycle V49: derivative recursion for a normalized
trace-class kernel shape}
\label{sec:v12v49}

V48 reduces the polynomial alias tail to the real-torus stock
\({\cal B}_m(\Sigma)\).  The normalized shape
\(\Sigma=A/\operatorname{Tr}A\) is a quotient, so differentiating it
naively can reintroduce the vanishing total mass removed in
V41--V44.  This note first cancels any common amplitude factor and
then gives a finite recursive bound for every shape derivative.

\subsection{Core normalization}

Let
\[
 C:{\mathbb T}^D\longrightarrow{\mathfrak S}_1({\cal H})
\]
be \(C^N\) in trace norm, with
\begin{equation}
 m(k):=\operatorname{Tr}C(k)\geq m_0>0             \tag{49.1}
\end{equation}
on the real torus.  Define
\begin{equation}
                         \Sigma(k):={C(k)\over m(k)}.       \tag{49.2}
\end{equation}
For \(0\leq j\leq N\), let
\begin{equation}
 c_j:=
 \sup_{k\in{\mathbb T}^D}
 \|D^jC(k)\|_{\operatorname{Mult}^j
  (\mathbb R^D;{\mathfrak S}_1)}.                 \tag{49.3}
\end{equation}
The trace functional has norm one on the trace class, so
\begin{equation}
 \|D^jm(k)\|_{\operatorname{Mult}^j(\mathbb R^D;\mathbb C)}
 \leq c_j.                                       \tag{49.4}
\end{equation}

If an unnormalized kernel has the form
\[
 A(k)=\tau(k)C(k),\qquad\tau(k)\geq0,
\]
then \(A/\operatorname{Tr}A=\Sigma\) wherever \(\tau>0\).
All following derivative bounds depend on the core \(C\), not on
\(\tau\).

\subsection{Recursive derivative theorem}

Define nonnegative numbers \(s_0,\ldots,s_N\) recursively by
\begin{align}
 s_0&:={c_0\over m_0},                            \tag{49.5}\\
 s_n&:={1\over m_0}
 \left[
 c_n+\sum_{j=1}^n
 \binom nj c_j s_{n-j}
 \right],
 \qquad1\leq n\leq N.                            \tag{49.6}
\end{align}

\begin{theorem}[Normalized-shape derivative recursion]
\label{thm:v12v49-recursion}
For every \(0\leq n\leq N\),
\begin{equation}
 \sup_{k\in{\mathbb T}^D}
 \|D^n\Sigma(k)\|_{\operatorname{Mult}^n
  (\mathbb R^D;{\mathfrak S}_1)}
 \leq s_n.                                       \tag{49.7}
\end{equation}
\end{theorem}

\begin{proof}
The case \(n=0\) follows from
\[
 \|\Sigma(k)\|_1\leq{\|C(k)\|_1\over m_0}\leq s_0.
\]
For \(n\geq1\), differentiate the identity
\[
                         m\Sigma=C.
\]
The symmetric Fr\'echet Leibniz rule gives
\begin{equation}
 D^nC
 =
 mD^n\Sigma
 +
 \sum_{\varnothing\ne J\subset\{1,\ldots,n\}}
 D^{|J|}m\,D^{n-|J|}\Sigma,                      \tag{49.8}
\end{equation}
where each term is evaluated on the arguments indexed by \(J\) and
its complement.  Grouping subsets by \(j=|J|\), taking multilinear
operator norms, using (49.1), (49.4), and the induction hypothesis,
gives
\[
 m_0\|D^n\Sigma\|
 \leq
 c_n+\sum_{j=1}^n\binom njc_js_{n-j}.
\]
Equation (49.6) proves the induction step.
\end{proof}

The recursion is finite, explicit, and contains every quotient-rule
cross term.  Retaining only \(c_n/m_0\) would miss derivatives of the
normalizer.

\subsection{Scale invariance of the recursion}

\begin{proposition}[Common constant rescaling cancels]
\label{prop:v12v49-scale}
Replace \(C\) by \(aC\) for a constant \(a>0\).  Then \(m_0\) and
every \(c_j\) scale by \(a\), while the recursively defined
\(s_j\) remain unchanged.
\end{proposition}

\begin{proof}
The ratio \(c_0/m_0\) is unchanged.  If
\(s_0,\ldots,s_{n-1}\) are unchanged, every term inside the brackets
in (49.6) scales by \(a\), as does its denominator.  Induction proves
the claim.
\end{proof}

This is the derivative-level form of mass cancellation.  A
scale-dependent amplitude must be factored before the \(c_j\) are
formed; otherwise its derivatives can contaminate the stock even
though it cancels from the shape.

\subsection{From derivative recursion to the V48 stock}

Take \(N=2m\).  Expanding the elliptic differential operator gives
\[
 (I-\Delta)^m
 =
 \sum_{r=0}^m\binom mr(-\Delta)^r.
\]
The operator \(\Delta^r\) is a sum of \(D^r\) ordered partial
derivatives of total order \(2r\).  Hence Theorem
\ref{thm:v12v49-recursion} implies
\begin{equation}
 \sup_k\|(I-\Delta)^m\Sigma(k)\|_1
 \leq
 \sum_{r=0}^m
 \binom mrD^rs_{2r}.                              \tag{49.9}
\end{equation}

\begin{corollary}[Explicit Sobolev stock]
\label{cor:v12v49-Sobolev-stock}
With normalized Haar measure,
\begin{equation}
 {\cal B}_m(\Sigma)
 \leq
 \boxed{
 B_m^{\rm rec}:=
 \sum_{r=0}^m\binom mrD^rs_{2r}.}                 \tag{49.10}
\end{equation}
Therefore the V48 residual tail satisfies
\begin{equation}
 \left\|G_M-I-C_{M,L}\right\|_1
 \leq
 C_{D,m}
 {B_m^{\rm rec}\over M^{2m}L^{2m-D}}
 \qquad(2m>D).                                   \tag{49.11}
\end{equation}
\end{corollary}

\begin{proof}
Integrate (49.9) over the probability torus to obtain (49.10), then
insert it into Theorem \ref{thm:v12v48-core-tail}.
\end{proof}

\subsection{Relative derivative stocks near an amplitude zero}

Suppose \(A=\tau C\), where \(\tau\) vanishes.  Direct derivative
stocks of \(A\) cannot in general recover the core stocks \(c_j\):
Leibniz expansion mixes derivatives of \(\tau\) and \(C\), and
division by \(\tau\) is singular at its zeros.  A valid certificate
must provide one of the following:
\begin{enumerate}
\item an explicit analytic or smooth factorization \(A=\tau C\);
\item a removable shape extension with direct derivative bounds for
\(\Sigma\); or
\item relative jet identities determining the core derivatives
without dividing interval enclosures by a set containing zero.
\end{enumerate}

\begin{warning}[Pointwise cancellation is not a derivative
certificate]
\label{warn:v12v49-pointwise}
The equality
\(A/\operatorname{Tr}A=C/\operatorname{Tr}C\) on
\(\{\tau>0\}\) proves pointwise cancellation.  It does not by itself
prove that the extended quotient is \(C^{2m}\) across
\(\{\tau=0\}\).  V44 direction independence or an explicit core
factorization is still required.
\end{warning}

\subsection{Moving physical normalization}

The derivative recursion concerns the trace-one shape \(\Sigma\).
The physical normalized kernel is
\[
 {\Sigma(k)\over v(k)},\qquad
 v(k)=\operatorname{Tr}\{S(k)\Sigma(k)\}.
\]
If Fourier analysis is applied after physical normalization, a
second recursion of the same form is needed for the scalar \(v\),
using derivatives of both \(S\) and \(\Sigma\).  The V46 route avoids
that extra quotient: it averages the trace-one shapes first and then
uses the finite-dimensional moving-normalizer comparison of V43.
These two operations are not interchangeable without a separate
commutation theorem.

\subsection{Cofinal derivative card}

At scale \(n\), the finite-regularity route is populated by:
\begin{enumerate}
\item a common amplitude factor and a \(C^{2m}\) trace-class core
\(C_n\);
\item a core trace floor \(m_{0,n}\);
\item derivative stocks \(c_{j,n}\), \(0\leq j\leq2m\);
\item the recursion (49.5)--(49.6);
\item the resulting \(B_{m,n}^{\rm rec}\);
\item a signed core and a mesh/core allocation satisfying the V48
summability condition; and
\item the V43 visibility and physical-functional transport bounds.
\end{enumerate}

Every row is invariant under constant rescaling of the core.  A
poorly chosen core factorization can nevertheless move variation
from \(\tau_n\) into \(C_n\) and enlarge the derivative stocks.

\begin{assessment}[Status after V49]
\label{ass:v12v49-status}
The real-regularity tail of V48 now has a finite computable input.
All normalized-shape derivatives through order \(2m\) are bounded
recursively by core derivatives and a core trace floor, after the
harmless amplitude factor is removed.  Equation (49.10) supplies the
trace-class Sobolev stock required by the signed-core tail.
\end{assessment}

\begin{firewall}[Claim boundary after V49]
\label{fire:v12v49-claim}
No RPESM common amplitude core, trace floor, derivative stock, or
signed Fourier core has been computed.  The recursion does not prove
smooth extension across an unidentified zero factor.  The missing
physical Duhamel fibre, determinant-line gluing, gravitational
tightness, and the full SM--Einstein continuum remain open.
\end{firewall}

\subsection{Parent record}

This V49 note appends to the validated V48 file with record
\[
\begin{array}{c|c}
\text{lines}&13974\\
\text{bytes}&466563\\
\text{SHA--256}&
\texttt{8ECCA77EA52287E392FB3C2F3A2C267AB53BB5DCA3201985BE69D96E4A516818}.
\end{array}
\]

% =============================================================================
% END APPENDED TWELFTH-CYCLE V49 MATERIAL
% =============================================================================
% =============================================================================
% BEGIN APPENDED TWELFTH-CYCLE V50 MATERIAL
% =============================================================================

\section{Twelfth-cycle V50: compulsory companion audit and the
Duhamel-stock decision}
\label{sec:v12v50}

\subsection{Audit identity}

This is the compulsory five-version audit following V45.  The
companion state is unchanged:
\[
\begin{array}{c|r|r}
\text{file}&\text{counted lines}&\text{bytes}\\ \hline
\texttt{Renewal\_Yang\_Mills\_Mass\_Gap\_Research\_Notes\_V55.tex}
 &20960&714665\\
\texttt{Renewal\_NS\_Stability\_Research\_Notes\_V26.tex}
 &5320&278283.
\end{array}                                                   \tag{50.1}
\]
The SHA--256 fingerprints remain
\[
\begin{split}
 &\texttt{9D7B7A75BCF6713BCC41A1F224C6EDD96AA6384758C1C8DD3274701A9C1EF207},
 \\
 &\texttt{82C887F8C2C69E4F28FAD7C3DC8DE5B9099CB5A4BF431EBA1FFF3E91935978AD}.
\end{split}                                                   \tag{50.2}
\]
Both files were read without modification.

\subsection{Audit conclusion}

Yang--Mills V55 still provides the signed packet-first Poisson
identity already transferred and upgraded in SM V46.  Its finite
Fourier core and residual Wiener stock have not yet been populated
for the Yang--Mills response.  It therefore supplies no new
coefficient, convergence rate, or physical conclusion.

Navier--Stokes V26 remains the same rank-one Oseen-kernel calculation.
It supplies no input to the trace-class analytic or Sobolev shape
stocks of V47--V49.

No theorem is repeated at this checkpoint.  The results since V45
have already addressed the available companion lesson:
\begin{enumerate}
\item V46 lifted the signed core--tail identity to trace-class
positive shapes;
\item V47 derived the exponential tail from a holomorphic polystrip;
\item V48 supplied a finite-regularity polynomial alternative; and
\item V49 reduced the Sobolev stock to derivatives of an
unnormalized core after common-factor cancellation.
\end{enumerate}

\subsection{Next upstream acquisition}

The remaining finite input in V49 is the derivative list
\[
 c_j=
 \sup_k\|D^jC(k)\|_1.
\]
For a selected finite Duhamel fibre, a natural positive core is
\[
                         C(k)=e^{-\beta H(k)}
\]
or a finite sum of such heat factors, where \(H(k)\) is the complete
Hermitian fibre Hamiltonian after the legal packet has been
assembled.

V51 will prove the exact time-ordered formula and trace-norm bound
for all Fr\'echet derivatives when \(H(k)\) is affine in the
background variables and has a certified spectral interval.  It
will then feed those constants through the V49 quotient recursion
and V48 signed-core tail.  Nonlinear Hamiltonian dependence will
remain a later Bell-polynomial extension.

\begin{assessment}[Status after V50]
\label{ass:v12v50-status}
The mandatory audit is complete and changes no current result.  The
next actionable bridge is from a finite Hermitian Duhamel fibre with
spectral bounds to the trace-class derivative stocks required by the
normalized-shape Fourier estimates.
\end{assessment}

\begin{firewall}[Claim boundary after V50]
\label{fire:v12v50-claim}
The physical RPESM fibre Hamiltonian and its matrices remain missing.
V50 records a proof route, not a populated Duhamel card.  No
Yang--Mills mass gap, Navier--Stokes regularity result, determinant-
line trivialization, gravitational tightness, or full SM--Einstein
continuum is claimed.
\end{firewall}

\subsection{Parent record}

This V50 note appends to the validated V49 file with record
\[
\begin{array}{c|c}
\text{lines}&14240\\
\text{bytes}&475330\\
\text{SHA--256}&
\texttt{5B3E838AE555F0915AFB9FE7E772FB23CBC3806A91012A20AF3811175A40C0CC}.
\end{array}
\]

% =============================================================================
% END APPENDED TWELFTH-CYCLE V50 MATERIAL
% =============================================================================
% =============================================================================
% BEGIN APPENDED TWELFTH-CYCLE V51 MATERIAL
% =============================================================================

\section{Twelfth-cycle V51: Duhamel simplex bounds for Gibbs kernel
shapes}
\label{sec:v12v51}

V49 reduces the real-regularity shape stock to derivatives of a
positive trace-class core.  For a finite physical fibre the natural
core is a heat factor \(e^{-\beta H}\).  This note derives all
Fr\'echet derivatives for an affine Hermitian Hamiltonian, bounds
them in trace norm by a time-ordered Schatten argument, and feeds
the resulting constants into the normalized-shape recursion.

\subsection{Affine Hamiltonian family}

Let \(U\subset\mathbb R^D\) be open and let
\[
 H(x)=H_0+V(x)
\]
be a self-adjoint family on a common dense domain, where
\(V:\mathbb R^D\to{\cal B}({\cal H})_{\rm sa}\) is linear and
bounded.  Assume
\begin{equation}
 \|V(h)\|\leq L\|h\|.                              \tag{51.1}
\end{equation}
Fix \(\beta>0\), and suppose \(e^{-\beta H(x)}\) is trace class for
every \(x\in U\).  Put
\begin{equation}
 C(x):=e^{-\beta H(x)},\qquad
 Z(x):=\operatorname{Tr}C(x).                    \tag{51.2}
\end{equation}

Let
\[
 \Delta_n(\beta)
 =
 \{(t_0,\ldots,t_n):t_j\geq0,\ \sum_{j=0}^nt_j=\beta\}
\]
carry the standard \(n\)-dimensional simplex measure, whose volume
is \(\beta^n/n!\).

\begin{theorem}[Time-ordered Duhamel derivative]
\label{thm:v12v51-Duhamel}
For \(n\geq1\),
\begin{align}
 D^nC(x)[h_1,\ldots,h_n]
 ={}&(-1)^n
 \sum_{\pi\in{\mathfrak S}_n}
 \int_{\Delta_n(\beta)}
 e^{-t_0H(x)}V(h_{\pi(1)})e^{-t_1H(x)}
 \nonumber\\
 &\qquad\cdots
 V(h_{\pi(n)})e^{-t_nH(x)}\,dt.                  \tag{51.3}
\end{align}
The formula is symmetric in the directions and converges in trace
norm.
\end{theorem}

\begin{proof}
For \(n=1\), the standard Duhamel identity is
\[
 DC(x)[h]
 =
 -\int_0^\beta
 e^{-(\beta-s)H(x)}V(h)e^{-sH(x)}\,ds,
\]
which is (51.3).  Differentiate an \(n\)-fold ordered integral once
more.  The derivative may enter any one of its \(n+1\) heat
segments.  Applying the one-derivative identity to that segment
inserts the new bounded operator at every possible ordered time.
The resulting regions partition
\(\Delta_{n+1}(\beta)\), while all placements of the labeled
directions give the permutations in
\({\mathfrak S}_{n+1}\).  Induction proves (51.3).

Trace-norm convergence follows from the uniform integrand estimate
proved in the next lemma, which also justifies differentiation under
the simplex integral.
\end{proof}

\subsection{The Schatten simplex bound}

\begin{lemma}[Time-ordered trace norm]
\label{lem:v12v51-Schatten}
Let \(H\) be self-adjoint with \(e^{-\beta H}\) trace class.  For
\(t_j\geq0\), \(\sum_{j=0}^nt_j=\beta\), and bounded
\(B_1,\ldots,B_n\),
\begin{equation}
 \left\|
 e^{-t_0H}B_1e^{-t_1H}\cdots B_ne^{-t_nH}
 \right\|_1
 \leq
 \operatorname{Tr}(e^{-\beta H})
 \prod_{j=1}^n\|B_j\|.                            \tag{51.4}
\end{equation}
\end{lemma}

\begin{proof}
For \(t_j>0\), choose the Schatten exponent
\(p_j=\beta/t_j\); if \(t_j=0\), use \(p_j=\infty\).
Then
\[
 \sum_{j=0}^n{1\over p_j}=1
\]
and
\[
 \|e^{-t_jH}\|_{p_j}
 =
 \{\operatorname{Tr}(e^{-\beta H})\}^{t_j/\beta}.
\]
Generalized Schatten H\"older inequality, with each \(B_j\) in
operator norm, gives (51.4), because the powers of the heat trace
multiply to one full power.  Zero-time factors follow by continuity.
\end{proof}

\begin{corollary}[All affine heat derivatives]
\label{cor:v12v51-heat-derivatives}
Suppose on a parameter region \(G\subset U\),
\begin{equation}
                         Z(x)\leq Z_+<\infty.      \tag{51.5}
\end{equation}
Then
\begin{equation}
 \sup_{x\in G}
 \|D^nC(x)\|_{\operatorname{Mult}^n
  (\mathbb R^D;{\mathfrak S}_1)}
 \leq
 Z_+(\beta L)^n
 \qquad(n\geq0).                                 \tag{51.6}
\end{equation}
\end{corollary}

\begin{proof}
Apply Lemma \ref{lem:v12v51-Schatten} to every integrand in (51.3)
and use (51.1).  There are \(n!\) permutations and the simplex has
volume \(\beta^n/n!\), so those factors cancel.  Equation (51.5)
gives (51.6).  The case \(n=0\) is \(\|C\|_1=Z\).
\end{proof}

\subsection{Normalized Gibbs shape}

Assume in addition
\begin{equation}
                         0<Z_-\leq Z(x)\leq Z_+
 \qquad(x\in G),                                  \tag{51.7}
\end{equation}
and put
\[
 \rho(x):={e^{-\beta H(x)}\over Z(x)}.
\]
This is a positive trace-one operator.  Define
\[
 q:={Z_+\over Z_-},\qquad a:=\beta L,
\]
and recursively
\begin{align}
 g_0&:=1,                                         \tag{51.8}\\
 g_n&:=
 q\left[
 a^n+\sum_{j=1}^n\binom nj a^jg_{n-j}
 \right],
 \qquad n\geq1.                                   \tag{51.9}
\end{align}

\begin{theorem}[Gibbs-shape derivative stock]
\label{thm:v12v51-Gibbs-stock}
For every \(n\geq0\),
\begin{equation}
 \sup_{x\in G}
 \|D^n\rho(x)\|_{\operatorname{Mult}^n
  (\mathbb R^D;{\mathfrak S}_1)}
 \leq g_n.                                       \tag{51.10}
\end{equation}
Consequently, for any integer \(m\) with \(2m>D\),
\begin{equation}
 {\cal B}_m(\rho)
 \leq
 \sum_{r=0}^m\binom mrD^rg_{2r}                 \tag{51.11}
\end{equation}
when \(G={\mathbb T}^D\).
\end{theorem}

\begin{proof}
Use the V49 quotient recursion with core
\(C=e^{-\beta H}\).  Corollary
\ref{cor:v12v51-heat-derivatives} gives
\(c_j\leq Z_+a^j\), and division by the trace floor \(Z_-\) gives
\(c_j/Z_-\leq qa^j\).  The order-zero bound is sharpened from the
generic quotient estimate to the exact identity
\(\|\rho\|_1=\operatorname{Tr}\rho=1\).  Induction in the proof of
Theorem \ref{thm:v12v49-recursion} then gives precisely
(51.8)--(51.10).  Corollary
\ref{cor:v12v49-Sobolev-stock} gives (51.11).
\end{proof}

Combining (51.11) with V48 gives a fully explicit polynomial
signed-core tail in terms of \(\beta,L,Z_+/Z_-\).

\subsection{Finite-dimensional spectral interval}

If \({\cal H}\) has dimension \(d\) and
\begin{equation}
 \mu I\preceq H(x)\preceq\Lambda I
 \qquad(x\in G),                                  \tag{51.12}
\end{equation}
then
\begin{equation}
 d e^{-\beta\Lambda}
 \leq Z(x)\leq
 d e^{-\beta\mu},                                 \tag{51.13}
\end{equation}
so
\begin{equation}
                         q\leq e^{\beta(\Lambda-\mu)}.      \tag{51.14}
\end{equation}
The dimension cancels from the partition ratio.  The relevant
spectral cost is the energy width, not the absolute energy origin.

\begin{proof}
Apply the decreasing scalar function \(e^{-\beta t}\) to the
eigenvalues in (51.12) and sum the resulting bounds.
\end{proof}

\subsection{Scalar-energy cancellation}

Let \(r:\mathbb R^D\to\mathbb R\) be affine and write
\begin{equation}
 H(x)=r(x)I+\widetilde H(x).                      \tag{51.15}
\end{equation}
Then
\[
 e^{-\beta H(x)}
 =
 e^{-\beta r(x)}e^{-\beta\widetilde H(x)}
\]
and hence
\begin{equation}
 {e^{-\beta H(x)}\over\operatorname{Tr}e^{-\beta H(x)}}
 =
 {e^{-\beta\widetilde H(x)}
  \over\operatorname{Tr}e^{-\beta\widetilde H(x)}}.        \tag{51.16}
\end{equation}

\begin{proposition}[Remove scalar energy before taking stocks]
\label{prop:v12v51-scalar-energy}
The bounds in Theorem \ref{thm:v12v51-Gibbs-stock} remain valid with
\(H\) replaced by \(\widetilde H\).  Therefore \(L\) may be replaced
by the norm of the centered derivative
\[
 V(h)-dr[h]I,
\]
and \(q\) by the partition ratio of the centered Hamiltonian.
\end{proposition}

\begin{proof}
Equation (51.16) is an exact identity of normalized shapes.
Because \(r\) is affine, \(\widetilde H\) remains an affine
self-adjoint family with bounded derivative.  Apply the preceding
theorems to it.
\end{proof}

This removes a background-dependent scalar energy that would
otherwise enlarge derivative stocks even though it has no effect on
the Gibbs state.

\subsection{Finite Duhamel acquisition card}

A selected physical fibre supplies the V49--V48 input once it
provides:
\begin{enumerate}
\item the complete legal-packet Hermitian matrix \(H(k)\);
\item a common domain and bounded affine derivative map \(V\);
\item a centered scalar energy \(r(k)\), when available;
\item uniform heat-trace bounds \(Z_-,Z_+\), or a finite spectral
interval;
\item the recursion (51.8)--(51.11);
\item the signed low Fourier core of the resulting Gibbs shape; and
\item the V43 physical visibility and normalizer-transport bounds.
\end{enumerate}

\begin{assessment}[Status after V51]
\label{ass:v12v51-status}
An affine Hermitian Duhamel fibre with heat-trace bounds now
generates every trace-class derivative stock required by V49.
The time-ordered simplex and Schatten H\"older inequality remove
factorial growth from the raw heat derivative bound.  A scalar
energy branch cancels exactly, leaving only centered Hamiltonian
variation and the partition ratio.
\end{assessment}

\begin{firewall}[Claim boundary after V51]
\label{fire:v12v51-claim}
The selected physical RPESM Hamiltonian, its legal packet, spectral
interval, and Fourier core remain missing.  Actual fibre dependence
may be nonlinear and then requires additional higher derivative
terms.  No determinant-line gluing, gravitational tightness, or
full SM--Einstein continuum is constructed.
\end{firewall}

\subsection{Parent record}

This V51 note appends to the validated V50 file with record
\[
\begin{array}{c|c}
\text{lines}&14348\\
\text{bytes}&479299\\
\text{SHA--256}&
\texttt{A39B469D12E138EFCB7DA2F40CA9BE42555359A47721C2547DBEADCBA4413CD1}.
\end{array}
\]

% =============================================================================
% END APPENDED TWELFTH-CYCLE V51 MATERIAL
% =============================================================================
% =============================================================================
% BEGIN APPENDED TWELFTH-CYCLE V52 MATERIAL
% =============================================================================

\section{Twelfth-cycle V52: nonlinear Duhamel jets and complete Bell
stocks}
\label{sec:v12v52}

V51 treats Hamiltonians affine in the background variables.  The
selected SM--Einstein fibre can depend nonlinearly on the metric,
connection, Higgs field, and retained boundary data.  This note
combines the operator Duhamel derivative with the Banach-valued
Fa\`a di Bruno formula.  The resulting constants are complete Bell
polynomials in the Hamiltonian jet norms.

\subsection{Nonlinear Hamiltonian hypotheses}

Let \(U\subset\mathbb R^D\) be open and let
\[
 H:U\longrightarrow{\rm SA}({\cal H})
\]
be a \(C^N\) self-adjoint family on a common domain.  Assume every
parameter derivative of positive order is bounded and
\begin{equation}
 \sup_{x\in G}
 \|D^jH(x)\|_{\operatorname{Mult}^j
  (\mathbb R^D;{\cal B}({\cal H}))}
 \leq L_j,
 \qquad1\leq j\leq N,                            \tag{52.1}
\end{equation}
on a parameter region \(G\).  Fix \(\beta>0\), assume
\(e^{-\beta H(x)}\) is trace class, and set
\[
 C(x)=e^{-\beta H(x)},\qquad
 Z(x)=\operatorname{Tr}C(x).
\]

\subsection{The exponential derivative as a multilinear map}

For fixed \(H\) and bounded self-adjoint
\(B_1,\ldots,B_p\), the \(p\)-th Fr\'echet derivative of
\(E(H)=e^{-\beta H}\) is
\begin{align}
 D^pE(H)[B_1,\ldots,B_p]
 ={}&(-1)^p
 \sum_{\pi\in{\mathfrak S}_p}
 \int_{\Delta_p(\beta)}
 e^{-t_0H}B_{\pi(1)}e^{-t_1H}
 \nonumber\\
 &\qquad\cdots
 B_{\pi(p)}e^{-t_pH}\,dt.                        \tag{52.2}
\end{align}
This is the fixed-base version of V51.

\begin{lemma}[Multilinear heat derivative bound]
\label{lem:v12v52-exponential}
For every \(p\geq1\),
\begin{equation}
 \|D^pE(H)[B_1,\ldots,B_p]\|_1
 \leq
 Z(H)\beta^p\prod_{j=1}^p\|B_j\|.                \tag{52.3}
\end{equation}
\end{lemma}

\begin{proof}
Lemma \ref{lem:v12v51-Schatten} bounds each ordered integrand by
\(Z(H)\prod_j\|B_j\|\).  The simplex volume is
\(\beta^p/p!\), and the permutation sum contains \(p!\) terms.
Multiplication gives (52.3).
\end{proof}

\subsection{Set partitions and Bell stocks}

Let \({\cal P}_n\) be the set of partitions of
\(\{1,\ldots,n\}\) into nonempty unordered blocks.  Define the
complete Bell polynomial by
\begin{equation}
 {\mathfrak B}_n(x_1,\ldots,x_n)
 :=
 \sum_{\Pi\in{\cal P}_n}
 \prod_{B\in\Pi}x_{|B|},
 \qquad {\mathfrak B}_0:=1.                      \tag{52.4}
\end{equation}
Put
\begin{equation}
                         a_n:=
 {\mathfrak B}_n(\beta L_1,\ldots,\beta L_n).
                                                        \tag{52.5}
\end{equation}
The first values are
\begin{align}
 a_1&=\beta L_1,                                  \tag{52.6}\\
 a_2&=\beta L_2+\beta^2L_1^2,                    \tag{52.7}\\
 a_3&=\beta L_3+3\beta^2L_1L_2+\beta^3L_1^3.     \tag{52.8}
\end{align}

\begin{theorem}[Nonlinear heat-kernel derivative stock]
\label{thm:v12v52-heat-stock}
If
\begin{equation}
                         Z(x)\leq Z_+
 \qquad(x\in G),                                  \tag{52.9}
\end{equation}
then for \(0\leq n\leq N\),
\begin{equation}
 \sup_{x\in G}
 \|D^nC(x)\|_{\operatorname{Mult}^n
  (\mathbb R^D;{\mathfrak S}_1)}
 \leq Z_+a_n.                                    \tag{52.10}
\end{equation}
\end{theorem}

\begin{proof}
The Banach-valued Fa\`a di Bruno formula for the composition
\(E\circ H\) is
\begin{equation}
 D^n(E\circ H)[h_1,\ldots,h_n]
 =
 \sum_{\Pi\in{\cal P}_n}
 D^{|\Pi|}E(H)
 \left[
 D^{|B|}H[h_i:i\in B]
 \right]_{B\in\Pi}.                              \tag{52.11}
\end{equation}
Apply Lemma \ref{lem:v12v52-exponential} to a term indexed by
\(\Pi\).  Its trace norm is at most
\[
 Z(x)\beta^{|\Pi|}
 \prod_{B\in\Pi}L_{|B|}
 =
 Z(x)\prod_{B\in\Pi}(\beta L_{|B|}).
\]
Sum over all set partitions and use (52.9), obtaining
\(Z_+a_n\).  The case \(n=0\) is
\(\|C(x)\|_1=Z(x)\leq Z_+\).
\end{proof}

The set-partition formula retains all mixed background legs.  Using
only pure coordinate derivatives would not control the Fr\'echet
multilinear norm required by V49.

\subsection{Normalized nonlinear Gibbs shape}

Assume
\begin{equation}
                         0<Z_-\leq Z(x)\leq Z_+
 \qquad(x\in G),                                  \tag{52.12}
\end{equation}
put \(q=Z_+/Z_-\), and define
\begin{align}
 g_0&:=1,                                         \tag{52.13}\\
 g_n&:=
 q\left[
 a_n+\sum_{j=1}^n\binom nj a_jg_{n-j}
 \right],
 \qquad1\leq n\leq N.                            \tag{52.14}
\end{align}

\begin{theorem}[Nonlinear Gibbs-shape stock]
\label{thm:v12v52-Gibbs}
For
\[
 \rho(x)={e^{-\beta H(x)}
          \over\operatorname{Tr}e^{-\beta H(x)}},
\]
one has
\begin{equation}
 \sup_{x\in G}
 \|D^n\rho(x)\|_{\operatorname{Mult}^n
  (\mathbb R^D;{\mathfrak S}_1)}
 \leq g_n,
 \qquad0\leq n\leq N.                            \tag{52.15}
\end{equation}
If \(N\geq2m\), then
\begin{equation}
 {\cal B}_m(\rho)
 \leq
 \sum_{r=0}^m\binom mrD^rg_{2r}                 \tag{52.16}
\end{equation}
on the probability torus.
\end{theorem}

\begin{proof}
Theorem \ref{thm:v12v52-heat-stock} supplies the V49 core derivative
bounds \(c_j\leq Z_+a_j\).  Divide by \(Z_-\) in the normalized-shape
recursion, use the exact order-zero value
\(\|\rho\|_1=1\), and obtain (52.13)--(52.15).
The Laplacian expansion in V49 gives (52.16).
\end{proof}

For an affine family, \(L_j=0\) for \(j\geq2\), so
\(a_n=(\beta L_1)^n\), and V51 is recovered.

\subsection{Nonlinear scalar-energy removal}

Let \(r:U\to\mathbb R\) be \(C^N\) and write
\begin{equation}
                         H(x)=r(x)I+\widetilde H(x).        \tag{52.17}
\end{equation}
No affinity assumption is needed.  Since the scalar commutes with
\(\widetilde H\),
\begin{equation}
 {e^{-\beta H(x)}
  \over\operatorname{Tr}e^{-\beta H(x)}}
 =
 {e^{-\beta\widetilde H(x)}
  \over\operatorname{Tr}e^{-\beta\widetilde H(x)}}.        \tag{52.18}
\end{equation}

\begin{corollary}[Centered nonlinear jet]
\label{cor:v12v52-centered}
Theorems \ref{thm:v12v52-heat-stock} and
\ref{thm:v12v52-Gibbs} may be applied to \(\widetilde H\), using
\begin{equation}
 \widetilde L_j
 :=
 \sup_{x\in G}
 \|D^jH(x)-D^jr(x)I\|,
 \qquad1\leq j\leq N,                            \tag{52.19}
\end{equation}
in the Bell stocks.  Derivatives of the removed scalar branch do
not appear elsewhere.
\end{corollary}

\begin{proof}
Equation (52.18) is exact, and
\(D^j\widetilde H=D^jH-D^jr\,I\).  Apply the theorems to the centered
family.
\end{proof}

The scalar jets \(D^jr\) must arise from one actual \(C^N\) function
\(r\).  Minimizing every jet independently without the integrability
relations of a scalar function would not define a valid
factorization.

\subsection{From nonlinear fibre jets to the signed tail}

For \(2m>D\), define
\[
 B_m^{\rm Gibbs}:=
 \sum_{r=0}^m\binom mrD^rg_{2r},
\]
using centered jet stocks when available.  Then V48 gives
\begin{equation}
 \left\|G_M-I-C_{M,L}\right\|_1
 \leq
 C_{D,m}
 {B_m^{\rm Gibbs}\over M^{2m}L^{2m-D}}.           \tag{52.20}
\end{equation}
Together with the finite signed core and the V43 visibility
estimate, this is the complete conditional chain from nonlinear
Hamiltonian jets to an adjacent normalized projective error.

\subsection{Exact acquisition order}

For a physical finite fibre:
\begin{enumerate}
\item assemble the legal momentum packet and its full Hamiltonian;
\item subtract one integrable scalar energy branch;
\item certify heat-trace upper and lower bounds;
\item bound every mixed Hamiltonian derivative through order \(2m\);
\item form the Bell stocks \(a_j\), quotient stocks \(g_j\), and
\(B_m^{\rm Gibbs}\);
\item compute the finite signed Fourier core; and
\item combine (52.20) with physical visibility and interscale
transport.
\end{enumerate}

Taking norms before the packet assembly or centering step can destroy
the cancellations on which the bounds depend.

\begin{assessment}[Status after V52]
\label{ass:v12v52-status}
The affine restriction in V51 has been removed.  Complete Bell
polynomials of the bounded Hamiltonian jets control every
trace-class heat derivative, and the normalized quotient recursion
then supplies the finite Sobolev stock.  An arbitrary smooth scalar
energy branch cancels exactly before these stocks are formed.
\end{assessment}

\begin{firewall}[Claim boundary after V52]
\label{fire:v12v52-claim}
No physical RPESM Hamiltonian jet, heat-trace bound, centered scalar
branch, or signed Fourier core has been populated.  The Bell formula
is conditional on common-domain differentiability and bounded
parameter derivatives.  Determinant-line gluing, nonlinear
gravitational tightness, and the full SM--Einstein continuum remain
open.
\end{firewall}

\subsection{Parent record}

This V52 note appends to the validated V51 file with record
\[
\begin{array}{c|c}
\text{lines}&14664\\
\text{bytes}&488998\\
\text{SHA--256}&
\texttt{A373C927D9F657DF472BB5063F8083821902D85C4C6F015B6CBDC21868D5AB50}.
\end{array}
\]

% =============================================================================
% END APPENDED TWELFTH-CYCLE V52 MATERIAL
% =============================================================================
% =============================================================================
% BEGIN APPENDED TWELFTH-CYCLE V53 MATERIAL
% =============================================================================

\section{Twelfth-cycle V53: scalar-quotient stability of Gibbs shapes}
\label{sec:v12v53}

V51--V52 control derivatives of one Gibbs shape as the background
varies.  Projective consistency also requires a comparison of two
nearby finite-fibre constructions, for example an exact packet
Hamiltonian and a local or truncated surrogate.  Estimating their
unnormalized heat kernels would introduce partition-function scales
that cancel from the physical state.  This note performs the
comparison after quotienting Hamiltonians by scalar energies.

\subsection{The scalar quotient norm}

For a bounded self-adjoint operator \(A\), define
\begin{equation}
 \|[A]\|_{\rm q}
 :=
 \inf_{c\in\mathbb R}\|A-cI\|.                    \tag{53.1}
\end{equation}
If
\[
 a_-:=\inf\operatorname{spec}A,\qquad
 a_+:=\sup\operatorname{spec}A,
\]
then
\begin{equation}
 \|[A]\|_{\rm q}={a_+-a_-\over2}.                 \tag{53.2}
\end{equation}

\begin{proof}
For every \(c\in\mathbb R\),
\[
 \|A-cI\|
 =
 \max\{|a_--c|,|a_+-c|\}.
\]
The maximum is minimized at \(c=(a_-+a_+)/2\), where it equals
\((a_+-a_-)/2\).
\end{proof}

Thus (53.1) measures the oscillating part of an energy perturbation.
It vanishes precisely for scalar perturbations, which leave a
normalized Gibbs state unchanged.

\subsection{A trace-norm path estimate}

Fix \(\beta>0\).  Let \(H\) be self-adjoint and suppose
\(e^{-\beta H}\) is trace class.  Let \(A\) be bounded and
self-adjoint, and put
\[
 H_t:=H+tA,\qquad
 C_t:=e^{-\beta H_t},\qquad
 Z_t:=\operatorname{Tr}C_t,\qquad
 \rho_t:={C_t\over Z_t},
 \qquad0\leq t\leq1.
\]
A bounded perturbation preserves the common domain and heat
trace-class property.  Indeed the min--max principle gives
\[
 \lambda_j(H)-t\|A\|
 \leq\lambda_j(H_t)
 \leq\lambda_j(H)+t\|A\|,
\]
with eigenvalues counted with multiplicity, and hence
\begin{equation}
 e^{-\beta t\|A\|}Z_0
 \leq Z_t\leq
 e^{\beta t\|A\|}Z_0.                             \tag{53.3}
\end{equation}
Here the discreteness needed by min--max follows from the
trace-class heat kernel.

\begin{lemma}[Normalized Duhamel velocity]
\label{lem:v12v53-velocity}
The path \(t\mapsto\rho_t\) is continuously differentiable in trace
norm and
\begin{equation}
 \|\dot\rho_t\|_1\leq2\beta\|A\|.
                                                            \tag{53.4}
\end{equation}
\end{lemma}

\begin{proof}
The bounded-perturbation Duhamel formula gives, in trace norm,
\begin{equation}
 \dot C_t
 =
 -\int_0^\beta
 e^{-(\beta-s)H_t}A e^{-sH_t}\,ds.                \tag{53.5}
\end{equation}
Lemma \ref{lem:v12v51-Schatten}, with one insertion, implies
\begin{equation}
 \|\dot C_t\|_1\leq\beta Z_t\|A\|.                \tag{53.6}
\end{equation}
It also supplies an integrable trace-norm majorant locally uniformly
in \(t\), using (53.3), and therefore justifies the asserted
trace-norm differentiation.

Cyclicity of the trace in (53.5) yields
\begin{equation}
 \dot Z_t
 =
 -\beta\operatorname{Tr}(A e^{-\beta H_t}),
 \qquad
 |\dot Z_t|\leq\beta Z_t\|A\|.                    \tag{53.7}
\end{equation}
Differentiate \(\rho_t=C_t/Z_t\).  Since
\(\|C_t\|_1=Z_t\),
\[
 \|\dot\rho_t\|_1
 \leq
 {\|\dot C_t\|_1\over Z_t}
 +
 {|\dot Z_t|\|C_t\|_1\over Z_t^2}.
\]
Equations (53.6)--(53.7) give (53.4).
\end{proof}

No lower bound on \(Z_t\) appears in (53.4): the two occurrences of
\(Z_t\) cancel exactly at each point of the path.

\subsection{The quotient-Lipschitz theorem}

For a self-adjoint \(L\) with trace-class heat kernel, write
\[
 \Gamma_\beta(L)
 :=
 {e^{-\beta L}\over\operatorname{Tr}e^{-\beta L}}.
\]

\begin{theorem}[Gibbs shapes are Lipschitz modulo scalars]
\label{thm:v12v53-quotient-Lipschitz}
Let \(H\) and \(K\) be self-adjoint on a common domain, suppose
\(e^{-\beta H}\) is trace class, and assume
\[
                         A:=K-H
\]
is bounded and self-adjoint.  Then \(e^{-\beta K}\) is trace class
and
\begin{equation}
 \|\Gamma_\beta(K)-\Gamma_\beta(H)\|_1
 \leq
 2\beta\|[K-H]\|_{\rm q}.                         \tag{53.8}
\end{equation}
Equivalently,
\begin{equation}
 \|\Gamma_\beta(K)-\Gamma_\beta(H)\|_1
 \leq
 \beta\{\sup\operatorname{spec}(K-H)
       -\inf\operatorname{spec}(K-H)\}.           \tag{53.9}
\end{equation}
The right side may harmlessly be replaced by its minimum with \(2\).
\end{theorem}

\begin{proof}
Fix \(c\in\mathbb R\).  Apply Lemma
\ref{lem:v12v53-velocity} to the path
\[
 H_t=H+t(A-cI).
\]
Its endpoint is \(K-cI\), and
\[
 \Gamma_\beta(K-cI)=\Gamma_\beta(K)
\]
because the scalar heat factor cancels in the normalization.
The fundamental theorem of calculus in the trace-class Banach space
therefore gives
\[
 \|\Gamma_\beta(K)-\Gamma_\beta(H)\|_1
 \leq
 \int_0^1\|\dot\rho_t\|_1\,dt
 \leq2\beta\|A-cI\|.
\]
Take the infimum over \(c\) and use (53.2).  Finally, the trace
distance between two density operators is at most \(2\).
\end{proof}

The estimate is uniform in the absolute energy origin, the Hilbert
space dimension, and the partition functions.  It requires only a
bounded difference of the two Hamiltonians; either Hamiltonian may
itself be unbounded.

\subsection{Passage through a retained-fibre channel}

Let
\[
 {\cal R}:{\mathfrak S}_1({\cal H})\longrightarrow
          {\mathfrak S}_1({\cal K})
\]
be positive and trace preserving.  Such a map contracts the trace
norm on self-adjoint inputs:
\begin{equation}
 \|{\cal R}(X)\|_1\leq\|X\|_1,
 \qquad X=X^*.                                    \tag{53.10}
\end{equation}

\begin{proof}
The Banach adjoint \({\cal R}^*\) is positive and unital.  Hence it
maps the order interval \([-I,I]\) into itself.  Trace-norm duality
for self-adjoint \(X\) gives
\[
 \|{\cal R}(X)\|_1
 =
 \sup_{-I\preceq B\preceq I}
 |\operatorname{Tr}(B{\cal R}(X))|
 =
 \sup_{-I\preceq B\preceq I}
 |\operatorname{Tr}({\cal R}^*(B)X)|
 \leq\|X\|_1.
\]
\end{proof}

\begin{corollary}[Adjacent retained-shape defect]
\label{cor:v12v53-retained}
Under the hypotheses of Theorem
\ref{thm:v12v53-quotient-Lipschitz},
\begin{equation}
 \|{\cal R}\Gamma_\beta(K)
      -{\cal R}\Gamma_\beta(H)\|_1
 \leq
 2\beta\|[K-H]\|_{\rm q}.                         \tag{53.11}
\end{equation}
In particular, if a surrogate fine Hamiltonian \(K\) has the exact
retained shape
\[
 {\cal R}\Gamma_\beta(K)=\sigma_{\rm coarse},
\]
then
\begin{equation}
 \|{\cal R}\Gamma_\beta(H)-\sigma_{\rm coarse}\|_1
 \leq
 2\beta\|[H-K]\|_{\rm q}.                         \tag{53.12}
\end{equation}
\end{corollary}

\begin{proof}
Apply (53.10) to the difference of the two Gibbs states and then use
(53.8).  The final assertion is substitution.
\end{proof}

Complete positivity is available for physical partial traces and
conditional expectations, but positivity and trace preservation
already suffice for this estimate.

\subsection{Cofinal application card}

At refinement level \(n\), suppose a legal fine packet supplies
self-adjoint Hamiltonians \(H_n,K_n\) on the same fibre and a
retained-fibre channel \({\cal R}_n\).  Put
\begin{equation}
 \varepsilon_n^{\rm Ham}
 :=
 2\beta_n\|[H_n-K_n]\|_{\rm q}.                   \tag{53.13}
\end{equation}
Then
\begin{equation}
 \|{\cal R}_n\Gamma_{\beta_n}(H_n)
      -{\cal R}_n\Gamma_{\beta_n}(K_n)\|_1
 \leq\min\{2,\varepsilon_n^{\rm Ham}\}.           \tag{53.14}
\end{equation}
Consequently, summability
\begin{equation}
 \sum_{n\geq n_0}\varepsilon_n^{\rm Ham}<\infty   \tag{53.15}
\end{equation}
is a sufficient Hamiltonian contribution to the telescoping
projective-error budget.  It must be combined with, rather than
substituted for, the signed Fourier-tail, visibility, Ward-repair,
and interscale-transport errors already isolated in V24--V52.

The correct acquisition is the scalar-quotient oscillation of the
full legal packet difference.  Bounding individual species blocks
before packet assembly can retain scalar pieces that cancel only
after the physical packet is complete.

\begin{assessment}[Status after V53]
\label{ass:v12v53-status}
Normalized Gibbs shapes now obey a dimension-free trace-norm
stability estimate controlled by the Hamiltonian difference modulo
scalar energies.  The estimate survives every positive
trace-preserving retained-fibre map and introduces no
partition-function floor.  It supplies a direct, summable
adjacent-scale defect once a physical surrogate pair \(H_n,K_n\) is
constructed.
\end{assessment}

\begin{firewall}[Claim boundary after V53]
\label{fire:v12v53-claim}
No physical RPESM Hamiltonian pair or estimate of
\(\|[H_n-K_n]\|_{\rm q}\) has been produced.  Equation (53.12)
requires a surrogate whose retained Gibbs shape is known exactly;
it does not assert that such a surrogate exists.  The result does
not by itself establish Ward compatibility, reflection positivity,
determinant-line gluing, gravitational tightness, or the full
SM--Einstein continuum.
\end{firewall}

\subsection{Parent record}

This V53 note appends to the validated V52 file with record
\[
\begin{array}{c|c}
\text{lines}&14969\\
\text{bytes}&498293\\
\text{SHA--256}&
\texttt{56D0848E725AF6A5497F18ABF3FB8E9B283094A8670971C0EE79A2DA4A1E4A7D}.
\end{array}
\]

% =============================================================================
% END APPENDED TWELFTH-CYCLE V53 MATERIAL
% =============================================================================
% =============================================================================
% BEGIN APPENDED TWELFTH-CYCLE V54 MATERIAL
% =============================================================================

\section{Twelfth-cycle V54: quotient-metric cell martingales for Gibbs
shapes}
\label{sec:v12v54}

V53 controls two normalized Gibbs states by the Hamiltonian
difference modulo scalars.  The present note turns that pointwise
estimate into an exact nested cell-average system and a summable
error bound for centre-sampled approximations.  This gives a
first-derivative route to convergence of a parameter-dependent
finite fibre, complementary to the higher-derivative Fourier tail
of V47--V52.

\subsection{The quotient differential length}

Let
\[
 H:{\mathbb T}^D\longrightarrow{\rm SA}({\cal H})
\]
be a \(C^1\) family on a common domain.  Assume one, hence every,
\(e^{-\beta H(x)}\) is trace class and that
\(H(y)-H(x)\) is bounded along coordinate charts.  For
\(v\in\mathbb R^D\), set
\[
 \|[DH(x)[v]]\|_{\rm q}
 :=
 \inf_{c\in\mathbb R}\|DH(x)[v]-cI\|
\]
and define
\begin{equation}
 L_{\rm q}
 :=
 \sup_{x\in{\mathbb T}^D}
 \sup_{|v|_2\leq1}
 \|[DH(x)[v]]\|_{\rm q}<\infty.                  \tag{54.1}
\end{equation}
Write
\[
 \rho(x):=\Gamma_\beta(H(x)).
\]

\begin{lemma}[Quotient-length Lipschitz bound]
\label{lem:v12v54-length}
For the flat torus distance \(d_{\mathbb T}\),
\begin{equation}
 \|\rho(y)-\rho(x)\|_1
 \leq2\beta L_{\rm q}d_{\mathbb T}(x,y).
                                                            \tag{54.2}
\end{equation}
\end{lemma}

\begin{proof}
Choose a shortest unit-speed flat geodesic
\(\gamma:[0,\ell]\to{\mathbb T}^D\) from \(x\) to \(y\), where
\(\ell=d_{\mathbb T}(x,y)\).  In the Banach quotient
\({\cal B}({\cal H})_{\rm sa}/\mathbb RI\), the fundamental theorem
of calculus and the quotient-norm triangle inequality give
\begin{align}
 \|[H(y)-H(x)]\|_{\rm q}
 &=
 \left\|\left[
 \int_0^\ell DH(\gamma(s))[\dot\gamma(s)]\,ds
 \right]\right\|_{\rm q}                         \nonumber\\
 &\leq
 \int_0^\ell
 \|[DH(\gamma(s))[\dot\gamma(s)]]\|_{\rm q}\,ds
 \leq L_{\rm q}\ell.                              \tag{54.3}
\end{align}
Theorem \ref{thm:v12v53-quotient-Lipschitz}, applied to the two
endpoints, now yields (54.2).
\end{proof}

This argument does not require choosing a differentiable scalar
centering \(r(x)\).  The quotient norm already performs the
pointwise scalar removal in a way compatible with path integration.

\subsection{Nested cell averages}

Identify \({\mathbb T}^D\) with \([0,1)^D\) with normalized Haar
measure \(dx\).  Let \({\cal P}_n\) be the half-open dyadic partition
into cubes of side \(2^{-n}\), and let \({\cal F}_n\) be its finite
sigma algebra.  Its cell diameter obeys
\begin{equation}
                         \delta_n:=\sqrt D\,2^{-n}.           \tag{54.4}
\end{equation}
For a Bochner-integrable
\(F:{\mathbb T}^D\to{\mathfrak S}_1({\cal H})\), let
\(E_nF\) be conditional expectation onto \({\cal F}_n\):
\begin{equation}
 (E_nF)(x)
 =
 {|Q|}^{-1}\int_QF(y)\,dy,
 \qquad x\in Q\in{\cal P}_n.                     \tag{54.5}
\end{equation}
The maps \(E_n\) are contractions on
\(L^1({\mathbb T}^D;{\mathfrak S}_1)\), and nesting gives
\begin{equation}
                         E_nE_{n+1}=E_n.          \tag{54.6}
\end{equation}

\begin{theorem}[Exact positive cell martingale]
\label{thm:v12v54-cell-martingale}
Put
\[
                         \rho_n:=E_n\rho.
\]
Then for every \(n\):
\begin{enumerate}
\item \(\rho_n(x)\) is positive and has trace one for almost every
\(x\);
\item the family is exactly projective,
\begin{equation}
                         E_n\rho_{n+1}=\rho_n;     \tag{54.7}
\end{equation}
\item it converges to \(\rho\) with the explicit estimate
\begin{equation}
 \|\rho-\rho_n\|_{L^1({\mathbb T}^D;{\mathfrak S}_1)}
 \leq
 2\beta L_{\rm q}\delta_n
 =
 2\beta L_{\rm q}\sqrt D\,2^{-n}.                 \tag{54.8}
\end{equation}
The same right side bounds the essential supremum of the pointwise
trace norm.
\end{enumerate}
\end{theorem}

\begin{proof}
The positive cone of \({\mathfrak S}_1\) is closed and convex, so its
Bochner averages are positive.  Continuity of the trace functional
and \(\operatorname{Tr}\rho(y)=1\) give
\(\operatorname{Tr}\rho_n(x)=1\).  Equation (54.7) is the tower
property (54.6).

For \(x\in Q\), Bochner's inequality and Lemma
\ref{lem:v12v54-length} give
\begin{align}
 \|\rho(x)-\rho_n(x)\|_1
 &\leq
 {|Q|}^{-1}\int_Q\|\rho(x)-\rho(y)\|_1\,dy
                                                        \nonumber\\
 &\leq
 2\beta L_{\rm q}\delta_n.                       \tag{54.9}
\end{align}
Take the essential supremum or integrate over the probability torus
to obtain (54.8).
\end{proof}

Thus a continuous Gibbs-shape field has an exact finite
positive-state reverse martingale at every dyadic resolution.  No
reciprocal normalizer enters the averaging operation.

\subsection{Centre samples and adjacent defects}

For each \(Q\in{\cal P}_n\), choose its centre \(x_Q\) on the flat
torus and define the computable sample field
\begin{equation}
 (S_n\rho)(x):=\rho(x_Q),\qquad x\in Q.            \tag{54.10}
\end{equation}
It is again positive and trace one.  Lemma
\ref{lem:v12v54-length} gives
\begin{equation}
 \|S_n\rho-\rho\|_{L^1}
 \leq2\beta L_{\rm q}\delta_n.                    \tag{54.11}
\end{equation}

\begin{proposition}[Summable sampled projective error]
\label{prop:v12v54-sampled}
The centre-sampled fields obey
\begin{equation}
 \|E_n(S_{n+1}\rho)-S_n\rho\|_{L^1}
 \leq
 2\beta L_{\rm q}(\delta_n+\delta_{n+1}).         \tag{54.12}
\end{equation}
Consequently,
\begin{equation}
 \sum_{n\geq n_0}
 \|E_n(S_{n+1}\rho)-S_n\rho\|_{L^1}
 \leq
 6\beta L_{\rm q}\sqrt D\,2^{-n_0}.               \tag{54.13}
\end{equation}
\end{proposition}

\begin{proof}
Insert the exact martingale:
\begin{align}
 \|E_n(S_{n+1}\rho)-S_n\rho\|_{L^1}
 &\leq
 \|E_n(S_{n+1}\rho-\rho_{n+1})\|_{L^1}
 +\|\rho_n-S_n\rho\|_{L^1}                       \nonumber\\
 &\leq
 \|S_{n+1}\rho-\rho\|_{L^1}
 +\|\rho-S_n\rho\|_{L^1}.                         \tag{54.14}
\end{align}
Here contraction of \(E_n\), the identity
\(E_n\rho_{n+1}=\rho_n\), and the triangle inequality were used.
The same cell-average calculation as (54.9), now with the fixed
point \(x_Q\) at the cell centre, gives
\(\|S_j\rho-\rho_j\|_{L^1}\leq
2\beta L_{\rm q}\delta_j\).  Apply this at levels \(j=n+1\) and
\(j=n\) in (54.14) to prove (54.12).
Since
\[
 \delta_n+\delta_{n+1}
 ={3\over2}\sqrt D\,2^{-n},
\]
summing the geometric series proves (54.13).
\end{proof}

The estimate also permits approximate packet Hamiltonians.  If the
sample at \(x_Q\) uses \(K_{n,Q}\) in place of \(H(x_Q)\), define
\begin{equation}
 \eta_n
 :=
 2\beta\max_{Q\in{\cal P}_n}
 \|[K_{n,Q}-H(x_Q)]\|_{\rm q}.                   \tag{54.15}
\end{equation}
Theorem \ref{thm:v12v53-quotient-Lipschitz} and the triangle
inequality add at most \(\eta_n+\eta_{n+1}\) to the right side of
(54.12).  Hence
\begin{equation}
 \sum_n\eta_n<\infty                              \tag{54.16}
\end{equation}
is sufficient to retain a summable adjacent sampled error.

\subsection{Relation to the Fourier route}

The V47--V52 route produces a finite signed Fourier core with a
uniform trace-norm tail.  It is adapted to translation-covariant
Poisson contraction and requires derivatives through order
\(2m>D\), or a holomorphic strip.  The cell route above uses only the
first quotient derivative and preserves positivity exactly, but its
finite approximants are piecewise constant and do not by themselves
supply a local Fourier/Wilson symbol.  The physical construction may
use either estimate where its coarse map requires it; their error
budgets should not be counted twice.

\subsection{Finite acquisition card}

A selected continuum fibre can use this route after certifying:
\begin{enumerate}
\item a common-domain periodic Hamiltonian \(H(x)\);
\item heat trace class at one base point;
\item the scalar-quotient differential stock \(L_{\rm q}\);
\item a nested physical partition whose mesh diameters tend to zero;
\item compatibility of its retained-fibre map with the cell
conditional expectations; and
\item summable surrogate packet errors \(\eta_n\), if exact cell
values are not used.
\end{enumerate}

\begin{assessment}[Status after V54]
\label{ass:v12v54-status}
The V53 pointwise comparison now yields an exact positive,
trace-one, nested cell-average martingale with explicit
\(O(2^{-n})\) trace-norm convergence.  Centre samples have a
geometrically summable adjacent projective defect, and approximate
finite packets add only their scalar-quotient Hamiltonian error.
\end{assessment}

\begin{firewall}[Claim boundary after V54]
\label{fire:v12v54-claim}
The theorem starts from a continuum Hamiltonian field and therefore
does not construct the physical RPESM Hamiltonian or its limiting
background measure.  Compatibility between dyadic parameter cells
and the actual Wilson/retained-field maps remains unproved.
Piecewise constant cell states do not establish locality,
reflection positivity of the full cylinder, determinant-line
gluing, gravitational tightness, or the full SM--Einstein
continuum.
\end{firewall}

\subsection{Parent record}

This V54 note appends to the validated V53 file with record
\[
\begin{array}{c|c}
\text{lines}&15284\\
\text{bytes}&507753\\
\text{SHA--256}&
\texttt{A3047DA6D6D3540575806FC8B85BA6233F9F824A68A52795FCF2273B735B2760}.
\end{array}
\]

% =============================================================================
% END APPENDED TWELFTH-CYCLE V54 MATERIAL
% =============================================================================
% =============================================================================
% BEGIN APPENDED TWELFTH-CYCLE V55 MATERIAL
% =============================================================================

\section{Twelfth-cycle V55: mandatory companion audit}
\label{sec:v12v55}

This is the required five-version review of the newest active
Yang--Mills and Navier--Stokes notes.  It is performed after the V53
scalar-quotient Gibbs estimate and the V54 cell-martingale
construction, and before the next substantive SM step.

\subsection{Files inspected}

The active companion records are:
\[
\begin{array}{c|r|r}
\text{file}&\text{lines}&\text{bytes}\\ \hline
\texttt{Renewal\_Yang\_Mills\_Mass\_Gap\_Research\_Notes\_V55.tex}
 &20959&714665\\
\texttt{Renewal\_NS\_Stability\_Research\_Notes\_V26.tex}
 &5319&278283.
\end{array}                                      \tag{55.1}
\]
Their SHA--256 digests at this audit are
\[
\begin{split}
 &\texttt{9D7B7A75BCF6713BCC41A1F224C6EDD96AA6384758C1C8DD3274701A9C1EF207},\\
 &\texttt{82C887F8C2C69E4F28FAD7C3DC8DE5B9099CB5A4BF431EBA1FFF3E91935978AD},
\end{split}                                      \tag{55.2}
\]
respectively.

\subsection{Yang--Mills transfer}

YM V55 requires the complete legal shift packet to be traced before
Fourier coefficients or cancellation-sensitive norms are formed.
It proves an exact signed Poisson alias identity and bounds only the
tail outside a finite computed signed core.  This agrees with the
operator-valued core-tail order already adopted in SM V46--V48.

The new consequence for V53--V54 is more specific.  The quotient
\[
                         \|[H-K]\|_{\rm q}
\]
must be taken after the complete legal SM packet has been assembled.
A scalar shift common to the full packet cancels from its normalized
Gibbs state.  Independent scalar shifts of species or momentum cards
generally change their relative Gibbs weights and therefore cannot
be discarded.  Thus a blockwise quotient stock is not automatically
a legal bound for the full packet.

YM V55 supplies no physical SM Hamiltonian, Fourier coefficient, or
numerical constant.  Its contribution is the packet-first proof
order:
\begin{equation}
 \boxed{
 \text{legal packet assembly}
 \longrightarrow
 \text{one global scalar quotient}
 \longrightarrow
 \text{retained shape or signed Fourier core}.}   \tag{55.3}
\end{equation}

\subsection{Navier--Stokes transfer}

NS V26 computes smaller rank-one norms for part of an Oseen-kernel
derivative, but explicitly permits those norms only when the actual
input remains in the rank-one cone.  The transferable proof rule is
that a structured norm is valid only after membership in its
structured input class has been proved.

For the SM programme, the scalar-quotient norm is a genuine quotient
norm on the entire self-adjoint packet difference, so V53 itself
needs no additional cone hypothesis.  A proposed reduction to
species blocks, symmetry orbits, tensor factors, or retained
subspaces does need a theorem showing how that structure sits inside
the full legal packet.  The next note will prove the direct-sum and
tensor-sum formulas needed for this check.

NS V26 contains no estimate for a Gibbs state, Ward complex,
reflection-positive kernel, or gravitational measure, and no
numerical constant is imported.

\subsection{Audit decision}

\begin{decision}[Direction after the V55 audit]
\label{dec:v12v55-direction}
Retain both V53 and V54 unchanged.  Next, compute the scalar quotient
of a complete block packet exactly.  Isolate the relative
block-centre spread that is lost by independent blockwise centering,
and prove the tensor-sum rule for species perturbations.  Only after
those formulas may a physical packet estimate use reduced block
data.
\end{decision}

\begin{assessment}[Status after V55]
\label{ass:v12v55-status}
The mandatory audit is complete.  YM V55 confirms the signed
core-tail and packet-first order; NS V26 confirms that every
structured norm needs a proved structural membership statement.
Together they expose one precise upstream issue for the SM
programme: global scalar cancellation is weaker than independent
blockwise cancellation.
\end{assessment}

\begin{firewall}[Claim boundary after V55]
\label{fire:v12v55-claim}
No companion theorem is treated as an SM theorem, and no companion
numerical value is imported.  The audit does not populate the
physical SM packet Hamiltonian or prove compatibility of V54 cells
with Wilson restriction.  The full SM--Einstein continuum remains
open.
\end{firewall}

\subsection{Parent record}

This V55 note appends to the validated V54 file with record
\[
\begin{array}{c|c}
\text{lines}&15574\\
\text{bytes}&517423\\
\text{SHA--256}&
\texttt{EC0FBF25B0AAB4D3CF9DCFEB6C60BA2F16D04F613C458C60C9EE67CBC786EC40}.
\end{array}
\]

% =============================================================================
% END APPENDED TWELFTH-CYCLE V55 MATERIAL
% =============================================================================
% =============================================================================
% BEGIN APPENDED TWELFTH-CYCLE V56 MATERIAL
% =============================================================================

\section{Twelfth-cycle V56: global scalar quotients of legal block
packets}
\label{sec:v12v56}

The V55 audit identified a gap between a quotient estimate for a
complete Gibbs packet and estimates performed separately on its
blocks.  This note computes that gap exactly.  Direct-sum sectors
retain their relative scalar offsets, whereas tensor-sum species
perturbations have an additive quotient norm.  These formulas state
precisely when a reduced packet calculation can feed the V53
stability theorem.

\subsection{A finite direct-sum formula}

Let
\[
 {\cal H}=\bigoplus_{\alpha=1}^N{\cal H}_\alpha,
 \qquad
 A=\bigoplus_{\alpha=1}^N A_\alpha,
\]
where every \(A_\alpha\) is bounded and self-adjoint.  Define
\begin{equation}
 a_\alpha^-:=\inf\operatorname{spec}A_\alpha,
 \quad
 a_\alpha^+:=\sup\operatorname{spec}A_\alpha,
 \quad
 q_\alpha:={a_\alpha^+-a_\alpha^-\over2},
 \quad
 m_\alpha:={a_\alpha^++a_\alpha^-\over2}.         \tag{56.1}
\end{equation}

\begin{theorem}[Exact packet quotient]
\label{thm:v12v56-direct-sum}
The scalar quotient of the complete direct-sum packet is
\begin{align}
 \|[A]\|_{\rm q}
 &=
 {1\over2}
 \left\{
 \max_\alpha a_\alpha^+
 -
 \min_\alpha a_\alpha^-
 \right\}                                         \tag{56.2}\\
 &=
 {1\over2}
 \left\{
 \max_\alpha(m_\alpha+q_\alpha)
 -
 \min_\alpha(m_\alpha-q_\alpha)
 \right\}.                                        \tag{56.3}
\end{align}
If
\[
 q_*:=\max_\alpha q_\alpha,
 \qquad
 \Delta_m:=\max_\alpha m_\alpha-\min_\alpha m_\alpha,
\]
then
\begin{equation}
 \max\{q_*,\Delta_m/2\}
 \leq\|[A]\|_{\rm q}
 \leq q_*+\Delta_m/2.                             \tag{56.4}
\end{equation}
\end{theorem}

\begin{proof}
The spectrum of a finite orthogonal direct sum is the union of the
block spectra.  Hence
\[
 \inf\operatorname{spec}A=\min_\alpha a_\alpha^-,
 \qquad
 \sup\operatorname{spec}A=\max_\alpha a_\alpha^+.
\]
Equations (56.2)--(56.3) follow from (53.2).

The global spectral interval contains every block interval, so its
half-width is at least \(q_*\).  Also
\[
 \max_\alpha a_\alpha^+\geq\max_\alpha m_\alpha,
 \qquad
 \min_\alpha a_\alpha^-\leq\min_\alpha m_\alpha,
\]
which gives the lower bound \(\Delta_m/2\).  Finally,
\[
 \max_\alpha(m_\alpha+q_\alpha)
 \leq\max_\alpha m_\alpha+q_*,
\]
and
\[
 \min_\alpha(m_\alpha-q_\alpha)
 \geq\min_\alpha m_\alpha-q_*.
\]
Insert these two inequalities in (56.3) to obtain the upper bound.
\end{proof}

The term \(\Delta_m/2\) is the relative block-centre stock.  It is
absent from the independently centered values \(q_\alpha\) and is
therefore the exact datum that a blockwise quotient calculation can
lose.

\begin{example}[Independent centering can erase the whole error]
\label{ex:v12v56-erasure}
Let \(A_\alpha=s_\alpha I_{{\cal H}_\alpha}\).  Then every
\(q_\alpha=0\), while
\begin{equation}
 \|[A]\|_{\rm q}
 =
 {\max_\alpha s_\alpha-\min_\alpha s_\alpha\over2}.           \tag{56.5}
\end{equation}
Thus independently subtracting \(s_\alpha\) from each block reports
zero even though the complete packet changes unless all
\(s_\alpha\) coincide.
\end{example}

\subsection{Why relative scalar offsets are physical}

Let
\[
 H=\bigoplus_{\alpha=1}^N H_\alpha
\]
have trace-class heat kernel in every block, and put
\[
 Z_\alpha:=\operatorname{Tr}e^{-\beta H_\alpha},
 \qquad
 p_\alpha:={Z_\alpha\over\sum_\gamma Z_\gamma}.
\]
Then
\begin{equation}
 \Gamma_\beta(H)
 =
 \bigoplus_{\alpha=1}^N
 p_\alpha\,\Gamma_\beta(H_\alpha).                \tag{56.6}
\end{equation}
If the blocks are shifted by \(s_\alpha I\), their new weights are
\begin{equation}
 p_\alpha'
 =
 {e^{-\beta s_\alpha}Z_\alpha
  \over\sum_\gamma e^{-\beta s_\gamma}Z_\gamma}.  \tag{56.7}
\end{equation}
Consequently, all block shifts cancel from the complete normalized
state if and only if they are common.

\begin{proof}
The heat exponential respects orthogonal direct sums, so its trace is
\(\sum_\alpha Z_\alpha\), which proves (56.6).  A scalar shift in
block \(\alpha\) multiplies that block heat kernel and partition
function by \(e^{-\beta s_\alpha}\), proving (56.7).
\end{proof}

This is the normalized-state reason for using one global quotient in
V53.  Sector-dependent scalar removal is a change of the sector
probabilities, not a gauge choice.

\subsection{Symmetry-orbit packets}

Suppose there is a reference block \(A_0\), unitaries
\(U_\alpha:{\cal H}_0\to{\cal H}_\alpha\), and real offsets
\(s_\alpha\) such that
\begin{equation}
 A_\alpha=U_\alpha A_0U_\alpha^*+s_\alpha I.       \tag{56.8}
\end{equation}

\begin{corollary}[Orbit reduction with an offset stock]
\label{cor:v12v56-orbit}
Under (56.8),
\begin{equation}
 \|[A]\|_{\rm q}
 =
 \|[A_0]\|_{\rm q}
 +
 {\max_\alpha s_\alpha-\min_\alpha s_\alpha\over2}.           \tag{56.9}
\end{equation}
In particular, a unitary orbit reduction preserves the complete
packet quotient exactly when the offsets are common.
\end{corollary}

\begin{proof}
Unitary conjugation preserves spectral endpoints.  Thus
\[
 a_\alpha^\pm=a_0^\pm+s_\alpha.
\]
Substitute in (56.2).
\end{proof}

Equation (56.9) is the certificate required before using a
symmetry-reduced packet norm.  A permutation of blocks alone does not
control the scalar offsets.

\subsection{Tensor-sum species perturbations}

Let
\[
 {\cal H}=\bigotimes_{j=1}^r{\cal H}_j
\]
and let \(A_j\) be bounded self-adjoint on \({\cal H}_j\).  On the
full tensor product define
\begin{equation}
 A
 :=
 \sum_{j=1}^r
 I_1\otimes\cdots\otimes I_{j-1}\otimes A_j
 \otimes I_{j+1}\otimes\cdots\otimes I_r.         \tag{56.10}
\end{equation}

\begin{theorem}[Exact tensor-sum quotient rule]
\label{thm:v12v56-tensor-sum}
For (56.10),
\begin{equation}
                         \|[A]\|_{\rm q}
 =
 \sum_{j=1}^r\|[A_j]\|_{\rm q}.                   \tag{56.11}
\end{equation}
\end{theorem}

\begin{proof}
The summands in (56.10) strongly commute.  If
\(a_j^\pm\) are the spectral endpoints of \(A_j\), the spectral
endpoints of their tensor sum are
\begin{equation}
 \inf\operatorname{spec}A=\sum_j a_j^-,
 \qquad
 \sup\operatorname{spec}A=\sum_j a_j^+.           \tag{56.12}
\end{equation}
For completeness, the operator inequalities
\(a_j^-I\preceq A_j\preceq a_j^+I\) give the corresponding two
bounds in (56.12).  Conversely, choose unit vectors in each factor
whose spectral expectations approach \(a_j^-\), respectively
\(a_j^+\); their tensor products show that both bounds are sharp.
Taking half the difference of the endpoints proves (56.11).
\end{proof}

Thus species-wise quotient stocks add exactly when the perturbation
is a genuine tensor sum on the complete tensor-product packet.  This
is a proved structural membership statement of the kind required by
the NS V26 audit.  Interactions, constraints, and direct-sum
superselection sectors fall outside (56.10) and must be included
before the quotient is evaluated.

\subsection{Compression monotonicity}

Let \(P\) be an orthogonal projection and regard \(PAP\) as an
operator on \(P{\cal H}\).

\begin{proposition}[Compression cannot increase scalar oscillation]
\label{prop:v12v56-compression}
For bounded self-adjoint \(A\),
\begin{equation}
                         \|[PAP]\|_{\rm q}
 \leq\|[A]\|_{\rm q}.                             \tag{56.13}
\end{equation}
\end{proposition}

\begin{proof}
For every \(c\in\mathbb R\),
\[
 \|PAP-cP\|
 =
 \|P(A-cI)P\|
 \leq\|A-cI\|.
\]
Take the infimum over \(c\).
\end{proof}

This proposition permits a retained-subspace estimate only when the
retained Hamiltonian difference is the actual compression
\(P(K-H)P\).  A Schur complement or an integrated effective
Hamiltonian includes additional terms and needs its own comparison.

\subsection{Packet-level Gibbs bound}

Let
\[
 H=\bigoplus_\alpha H_\alpha,\qquad
 K=\bigoplus_\alpha K_\alpha,
 \qquad
 A_\alpha:=K_\alpha-H_\alpha,
\]
on common block domains, with each \(A_\alpha\) bounded.  Combining
Theorems \ref{thm:v12v53-quotient-Lipschitz} and
\ref{thm:v12v56-direct-sum} gives
\begin{equation}
 \|\Gamma_\beta(K)-\Gamma_\beta(H)\|_1
 \leq
 \beta
 \left\{
 \max_\alpha\sup\operatorname{spec}A_\alpha
 -
 \min_\alpha\inf\operatorname{spec}A_\alpha
 \right\}.                                        \tag{56.14}
\end{equation}
This is the legal blockwise acquisition formula: compute both
spectral endpoints of every complete block difference, then take one
maximum and one minimum across the packet.  No per-block centering is
inserted before those two global extrema.

For a mixed packet containing tensor species inside direct-sum
sectors, apply (56.11) within each certified tensor sector, retain
its centre and half-width, and then apply (56.3) across the
direct-sum sectors.

\begin{assessment}[Status after V56]
\label{ass:v12v56-status}
The packet-first scalar quotient required by V53 is now explicit for
finite direct sums, unitary symmetry orbits, tensor sums, and genuine
compressions.  The relative direct-sum block centres are a necessary
stock; tensor-factor half-widths instead add exactly.  Equation
(56.14) converts certified block spectral endpoints into a
dimension-free Gibbs-shape error.
\end{assessment}

\begin{firewall}[Claim boundary after V56]
\label{fire:v12v56-claim}
No decomposition of the physical RPESM packet into the structures
above has been proved, and no physical block spectral endpoint has
been computed.  Gauge constraints, Yukawa mixing, chiral
determinants, and Schur-complement interactions may prevent the
simple tensor or direct-sum formulas from applying before the full
legal packet is assembled.  Reflection positivity, gravitational
tightness, and the full SM--Einstein continuum remain open.
\end{firewall}

\subsection{Parent record}

This V56 note appends to the validated V55 file with record
\[
\begin{array}{c|c}
\text{lines}&15704\\
\text{bytes}&522423\\
\text{SHA--256}&
\texttt{FEFD4A095EDD45E6FBF7CF93E22E0438B7AB4B8A5DC1FBEE87126B765E65462E}.
\end{array}
\]

% =============================================================================
% END APPENDED TWELFTH-CYCLE V56 MATERIAL
% =============================================================================
% =============================================================================
% BEGIN APPENDED TWELFTH-CYCLE V57 MATERIAL
% =============================================================================

\section{Twelfth-cycle V57: the sharp entropy bound for Gibbs-shape
stability}
\label{sec:v12v57}

V53 obtained
\(\|\Gamma_\beta(K)-\Gamma_\beta(H)\|_1
 \leq2\beta\|[K-H]\|_{\rm q}\)
by differentiating the normalized heat kernel and estimating its two
quotient-rule terms separately.  Those two terms are correlated.
Symmetrized quantum relative entropy retains their cancellation and
improves the constant by a factor of two.  This note proves the
sharper estimate and records its consequences for V54 and V56.

\subsection{Relative entropy and Pinsker's inequality}

For density operators \(\rho,\sigma\), with the usual support
condition, write
\[
 D(\rho\|\sigma)
 :=
 \operatorname{Tr}\rho(\log\rho-\log\sigma).
\]

\begin{lemma}[Quantum Pinsker inequality]
\label{lem:v12v57-Pinsker}
With natural logarithms,
\begin{equation}
 D(\rho\|\sigma)
 \geq{1\over2}\|\rho-\sigma\|_1^2.                \tag{57.1}
\end{equation}
\end{lemma}

\begin{proof}
Let \(X=\rho-\sigma\), let \(P\) be the positive spectral projection
of \(X\), and apply the two-outcome measurement
\[
 {\cal M}(\tau)
 =
 \bigl(\operatorname{Tr}(P\tau),
       \operatorname{Tr}((I-P)\tau)\bigr).
\]
Data processing for relative entropy gives
\[
 D(\rho\|\sigma)\geq
 D_{\rm cl}({\cal M}\rho\|{\cal M}\sigma).
\]
Because \(\operatorname{Tr}X=0\),
\[
 \operatorname{Tr}(PX)={1\over2}\|X\|_1.
\]
Thus the \(\ell^1\) distance between the two binary probability
vectors is \(\|X\|_1\).  The classical Pinsker inequality
\(D_{\rm cl}(p\|q)\geq\frac12\|p-q\|_1^2\) proves (57.1).
\end{proof}

The proof also covers nonfaithful states by the standard
infinite-relative-entropy convention and approximation on their
joint support.

\subsection{The symmetrized Gibbs identity}

Let \(H,K\) be self-adjoint on a common domain, assume
\(e^{-\beta H}\) is trace class, and suppose
\[
                         A:=K-H
\]
is bounded and self-adjoint.  Put
\[
 \rho:=\Gamma_\beta(H),\qquad
 \sigma:=\Gamma_\beta(K).
\]
As in V53, bounded perturbation implies that \(e^{-\beta K}\) is
trace class.

\begin{lemma}[Bounded-perturbation entropy identity]
\label{lem:v12v57-entropy-identity}
The two relative entropies are finite and
\begin{align}
 D(\rho\|\sigma)
 &=
 \beta\operatorname{Tr}(\rho A)
 +\log{Z_K\over Z_H},                             \tag{57.2}\\
 D(\sigma\|\rho)
 &=
 -\beta\operatorname{Tr}(\sigma A)
 +\log{Z_H\over Z_K},                             \tag{57.3}
\end{align}
where \(Z_H=\operatorname{Tr}e^{-\beta H}\) and
\(Z_K=\operatorname{Tr}e^{-\beta K}\).  Consequently,
\begin{equation}
 D(\rho\|\sigma)+D(\sigma\|\rho)
 =
 \beta\operatorname{Tr}\{A(\rho-\sigma)\}.        \tag{57.4}
\end{equation}
\end{lemma}

\begin{proof}
The Gibbs densities are faithful on \({\cal H}\) and functional
calculus gives
\[
 \log\rho=-\beta H-(\log Z_H)I,\qquad
 \log\sigma=-\beta K-(\log Z_K)I
\]
on their common form domain.  Their logarithmic difference is the
bounded operator
\[
 \log\rho-\log\sigma
 =
 \beta A+\log(Z_K/Z_H)I.
\]
Its expectation in \(\rho\) is finite and gives (57.2).  Reversing
the two states gives (57.3).  Equivalently, these identities follow
first on finite spectral cutoffs and then by monotone approximation;
boundedness of \(A\) and the partition comparison (53.3) provide a
uniform finite majorant.  Adding (57.2) and (57.3) cancels the two
log-partition terms and proves (57.4).
\end{proof}

The cancellation in (57.4) is exactly the correlation discarded by
the separate estimates of \(\dot C_t\) and \(\dot Z_t\) in V53.

\subsection{Sharp scalar-quotient stability}

\begin{theorem}[Sharp Gibbs quotient bound]
\label{thm:v12v57-sharp-Gibbs}
Under the preceding hypotheses,
\begin{equation}
 \boxed{
 \|\Gamma_\beta(K)-\Gamma_\beta(H)\|_1
 \leq
 \beta\|[K-H]\|_{\rm q}.}                         \tag{57.5}
\end{equation}
Equivalently,
\begin{equation}
 \|\Gamma_\beta(K)-\Gamma_\beta(H)\|_1
 \leq
 {\beta\over2}
 \left\{
 \sup\operatorname{spec}(K-H)
 -
 \inf\operatorname{spec}(K-H)
 \right\}.                                        \tag{57.6}
\end{equation}
The right side may be replaced by its minimum with \(2\).
The coefficient one in (57.5) is optimal at first order.
\end{theorem}

\begin{proof}
Set
\[
                         \delta:=\|\rho-\sigma\|_1.
\]
Apply Lemma \ref{lem:v12v57-Pinsker} in both directions and add:
\begin{equation}
 \delta^2
 \leq
 D(\rho\|\sigma)+D(\sigma\|\rho).                 \tag{57.7}
\end{equation}
For every \(c\in\mathbb R\), equations (57.4) and
\(\operatorname{Tr}(\rho-\sigma)=0\) give
\begin{align}
 \delta^2
 &\leq
 \beta\operatorname{Tr}\{(A-cI)(\rho-\sigma)\}
                                                        \nonumber\\
 &\leq
 \beta\|A-cI\|\,\delta.                            \tag{57.8}
\end{align}
If \(\delta=0\), the result is immediate.  Otherwise divide by
\(\delta\) and take the infimum over \(c\), proving (57.5).
Equation (53.2) gives (57.6).

For optimality, take a two-dimensional commuting example with
\(H=0\) and
\[
 K_t=t\begin{pmatrix}a&0\\0&-a\end{pmatrix}.
\]
Then \(\|[K_t-H]\|_{\rm q}=|t|a\), while direct expansion of the two
Gibbs probabilities gives
\[
 \|\Gamma_\beta(K_t)-\Gamma_\beta(H)\|_1
 =
 \beta|t|a+O(t^2).
\]
Hence the coefficient multiplying
\(\beta\|[K-H]\|_{\rm q}\) cannot be smaller than one.
\end{proof}

The V53 bound remains correct, but (57.5) supersedes it wherever the
same bounded-difference hypotheses hold.

\subsection{Retained channels and background paths}

A positive trace-preserving map is trace-norm contractive on
self-adjoint inputs by (53.10).  Therefore
\begin{equation}
 \|{\cal R}\Gamma_\beta(K)-{\cal R}\Gamma_\beta(H)\|_1
 \leq
 \beta\|[K-H]\|_{\rm q}.                           \tag{57.9}
\end{equation}
If \({\cal R}\Gamma_\beta(K)=\sigma_{\rm coarse}\), the same right
side bounds
\(\|{\cal R}\Gamma_\beta(H)-\sigma_{\rm coarse}\|_1\).

For the \(C^1\) torus family of V54, its quotient-length estimate
(54.3) and (57.5) give the improved Lipschitz bound
\begin{equation}
 \|\rho(y)-\rho(x)\|_1
 \leq
 \beta L_{\rm q}d_{\mathbb T}(x,y).               \tag{57.10}
\end{equation}
Repeating the unchanged Bochner-average proofs in V54 yields
\begin{align}
 \|\rho-E_n\rho\|_{L^1}
 &\leq
 \beta L_{\rm q}\sqrt D\,2^{-n},                  \tag{57.11}\\
 \|E_n(S_{n+1}\rho)-S_n\rho\|_{L^1}
 &\leq
 \beta L_{\rm q}(\delta_n+\delta_{n+1}),           \tag{57.12}\\
 \sum_{n\geq n_0}
 \|E_n(S_{n+1}\rho)-S_n\rho\|_{L^1}
 &\leq
 3\beta L_{\rm q}\sqrt D\,2^{-n_0}.               \tag{57.13}
\end{align}
For approximate cell Hamiltonians, the V54 stock is sharpened to
\begin{equation}
 \eta_n^{\rm sharp}
 :=
 \beta\max_{Q\in{\cal P}_n}
 \|[K_{n,Q}-H(x_Q)]\|_{\rm q}.                    \tag{57.14}
\end{equation}

\subsection{Sharpened packet endpoint formula}

For the direct-sum packet of V56, combine (57.6) with (56.2):
\begin{equation}
 \|\Gamma_\beta(K)-\Gamma_\beta(H)\|_1
 \leq
 {\beta\over2}
 \left\{
 \max_\alpha\sup\operatorname{spec}A_\alpha
 -
 \min_\alpha\inf\operatorname{spec}A_\alpha
 \right\},
 \qquad A_\alpha=K_\alpha-H_\alpha.               \tag{57.15}
\end{equation}
Thus (57.15) supersedes the coefficient in (56.14).  The structural
content of V56 is unchanged: one must still retain the direct-sum
centre spread, tensor-sum quotient stocks still add exactly, and
independent centering of superselection blocks remains illegal.

\subsection{Cofinal budget replacement}

In any later telescoping budget, replace the V53 Hamiltonian term
\[
 2\beta_n\|[H_n-K_n]\|_{\rm q}
\]
by
\begin{equation}
 \varepsilon_n^{\rm Ham,sharp}
 :=
 \beta_n\|[H_n-K_n]\|_{\rm q}.                    \tag{57.16}
\end{equation}
Summability of (57.16) is sufficient for the Hamiltonian comparison
part of the adjacent projective error.  This improvement does not
alter the other Ward, signed-tail, visibility, transport, or
gravitational contributions.

\begin{assessment}[Status after V57]
\label{ass:v12v57-status}
Symmetrized quantum relative entropy removes the factor two in the
normalized Duhamel estimate.  The sharp dimension-free bound is
\(\beta\) times the scalar-quotient norm of the complete Hamiltonian
difference, and its coefficient is locally optimal.  All V54 cell
rates and the V56 packet endpoint bound are correspondingly
sharpened.
\end{assessment}

\begin{firewall}[Claim boundary after V57]
\label{fire:v12v57-claim}
The theorem still requires a bounded difference of common-domain
Hamiltonians.  Metric or connection changes that alter an unbounded
principal part are not covered without a new form or resolvent
argument.  No physical RPESM packet endpoint, relative-entropy
locality estimate, determinant-line gluing, gravitational
tightness, or full SM--Einstein continuum has been established.
\end{firewall}

\subsection{Parent record}

This V57 note appends to the validated V56 file with record
\[
\begin{array}{c|c}
\text{lines}&16045\\
\text{bytes}&532728\\
\text{SHA--256}&
\texttt{4EF1A3A358FBC34C6053DCBBD95A66DDEE913F6BA0AF9270D1F259697F5A02D0}.
\end{array}
\]

% =============================================================================
% END APPENDED TWELFTH-CYCLE V57 MATERIAL
% =============================================================================
% =============================================================================
% BEGIN APPENDED TWELFTH-CYCLE V58 MATERIAL
% =============================================================================

\section{Twelfth-cycle V58: thermal-tail stability for unbounded
Gibbs perturbations}
\label{sec:v12v58}

V57 gives the sharp scalar-quotient estimate when \(K-H\) is
bounded.  Variations of a metric or connection can change the
principal part of a Dirac or Laplace-type operator, so that
hypothesis is too restrictive for the eventual SM--Einstein fibre.
This note isolates a rigorous extension: clip an entropy-admissible
unbounded perturbation to a bounded oscillating core and pay only
its thermal tail.

\subsection{Entropy-admissible Gibbs pairs}

Let \(H,K\) be self-adjoint with trace-class heat kernels at inverse
temperature \(\beta\), and put
\[
 \rho=\Gamma_\beta(H),\qquad
 \sigma=\Gamma_\beta(K).
\]
For this note, a self-adjoint operator \(A\) is called an
\emph{entropy-admissible difference} for the pair if:
\begin{enumerate}
\item \(A\) represents the quadratic-form difference \(K-H\) on a
common dense form core;
\item the expectations of \(A\) in \(\rho\) and \(\sigma\) are
well-defined and finite;
\item both relative entropies are finite; and
\item the symmetrized Gibbs identity holds:
\begin{equation}
 D(\rho\|\sigma)+D(\sigma\|\rho)
 =
 \beta\operatorname{Tr}\{A(\rho-\sigma)\}.        \tag{58.1}
\end{equation}
\end{enumerate}
For the bounded common-domain perturbation \(A=K-H\), Lemma
\ref{lem:v12v57-entropy-identity} proves all three conditions.  For
an unbounded form difference, (58.1) is an acquisition condition,
not an automatic assertion.

\subsection{Spectral clipping}

For \(c\in\mathbb R\) and \(R>0\), let
\[
 f_R(s):=\max\{-R,\min\{s,R\}\}
\]
and define by functional calculus
\begin{align}
 B_{c,R}
 &:=
 cI+f_R(A-cI),                                    \tag{58.2}\\
 G_{c,R}
 &:=
 A-B_{c,R}.                                       \tag{58.3}
\end{align}
Then \(B_{c,R}-cI\) is bounded by \(R\), and
\begin{equation}
 |G_{c,R}|=(|A-cI|-R)_+.                          \tag{58.4}
\end{equation}
Define the two-state thermal tail
\begin{equation}
 \tau_c(R)
 :=
 \operatorname{Tr}\{\rho(|A-cI|-R)_+\}
 +
 \operatorname{Tr}\{\sigma(|A-cI|-R)_+\}.         \tag{58.5}
\end{equation}
The traces in (58.5) are understood as positive quadratic-form
expectations and may initially take the value \(+\infty\).

\begin{theorem}[Clipped entropy stability]
\label{thm:v12v58-clipped}
Let \(A\) be an entropy-admissible difference and suppose
\(\tau_c(R)<\infty\).  Then
\begin{equation}
 \boxed{
 \|\rho-\sigma\|_1
 \leq
 {\beta R+
 \sqrt{\beta^2R^2+4\beta\tau_c(R)}
 \over2}.}                                        \tag{58.6}
\end{equation}
The right side may be minimized over \(c\in\mathbb R\) and \(R>0\),
and may be replaced by its minimum with \(2\).
\end{theorem}

\begin{proof}
Set \(X=\rho-\sigma\) and
\(\delta=\|X\|_1\).  The two Pinsker inequalities used in (57.7),
together with entropy admissibility, give
\begin{equation}
 \delta^2\leq\beta\operatorname{Tr}(AX).           \tag{58.7}
\end{equation}
Because \(\operatorname{Tr}X=0\), equations (58.2)--(58.3) imply
\[
 \operatorname{Tr}(AX)
 =
 \operatorname{Tr}\{f_R(A-cI)X\}
 +
 \operatorname{Tr}(G_{c,R}X).
\]
The first term is at most \(R\delta\) in absolute value.  For the
second, the form inequalities
\(-|G_{c,R}|\preceq G_{c,R}\preceq|G_{c,R}|\) give
\begin{align}
 |\operatorname{Tr}(G_{c,R}X)|
 &\leq
 |\operatorname{Tr}(G_{c,R}\rho)|
 +
 |\operatorname{Tr}(G_{c,R}\sigma)|              \nonumber\\
 &\leq\tau_c(R).                                  \tag{58.8}
\end{align}
Thus
\begin{equation}
                         \delta^2
 \leq\beta R\delta+\beta\tau_c(R).                \tag{58.9}
\end{equation}
The nonnegative solutions of this quadratic inequality lie below
its positive root, which is (58.6).
\end{proof}

If \(A-cI\) is bounded by \(R\), then \(\tau_c(R)=0\), and (58.6)
reduces to
\(\|\rho-\sigma\|_1\leq\beta R\).  Minimizing over \(c\) recovers
the sharp V57 theorem.

\subsection{A thermal-moment corollary}

For \(p>1\), define
\begin{equation}
 M_p(c)
 :=
 \operatorname{Tr}\{\rho|A-cI|^p\}
 +
 \operatorname{Tr}\{\sigma|A-cI|^p\}.             \tag{58.10}
\end{equation}

\begin{corollary}[Moment form of the unbounded estimate]
\label{cor:v12v58-moment}
If \(M_p(c)<\infty\), then
\begin{equation}
 \boxed{
 \|\rho-\sigma\|_1
 \leq
 \left\{
 \beta C_p M_p(c)^{1/p}
 \right\}^{p/(p+1)},\qquad
 C_p:=p(p-1)^{-(p-1)/p}.}                         \tag{58.11}
\end{equation}
Again the bound may be minimized over \(c\) and capped by \(2\).
\end{corollary}

\begin{proof}
The scalar inequality
\[
 (|s|-R)_+\leq {|s|^p\over R^{p-1}}
\]
and functional calculus give
\begin{equation}
                         \tau_c(R)
 \leq {M_p(c)\over R^{p-1}}.                      \tag{58.12}
\end{equation}
Insert this in (58.9):
\begin{equation}
 \delta^2
 \leq
 \beta\left\{R\delta+{M_p(c)\over R^{p-1}}\right\}.
                                                            \tag{58.13}
\end{equation}
For \(\delta>0\), minimize the braces over \(R>0\).  The minimizer
satisfies
\[
 R^p={(p-1)M_p(c)\over\delta},
\]
and the minimum is
\[
 C_p M_p(c)^{1/p}\delta^{(p-1)/p}.
\]
Consequently,
\[
 \delta^{(p+1)/p}
 \leq\beta C_p M_p(c)^{1/p},
\]
which is (58.11).  The case \(\delta=0\) is immediate.
\end{proof}

The exponent in (58.11) reflects the information loss in replacing
an unbounded perturbation by one moment.  It is not claimed to be
optimal for a particular elliptic family.

\subsection{Scaled principal perturbations}

Suppose \(A=\varepsilon V\) with \(V\) self-adjoint and
\[
 {\cal M}_p(d)
 :=
 \operatorname{Tr}\{\rho|V-dI|^p\}
 +
 \operatorname{Tr}\{\sigma|V-dI|^p\}<\infty.
\]
Taking \(c=\varepsilon d\) in (58.10) gives
\[
                         M_p(c)
 =
 |\varepsilon|^p{\cal M}_p(d).
\]
Therefore (58.11) becomes
\begin{equation}
 \|\rho-\sigma\|_1
 \leq
 (\beta C_p)^{p/(p+1)}
 |\varepsilon|^{p/(p+1)}
 {\cal M}_p(d)^{1/(p+1)}.                         \tag{58.14}
\end{equation}
Uniform thermal \(p\)-moments with \(p\) large make the exponent
\(p/(p+1)\) approach the bounded-perturbation exponent one.  This is
a possible route for a first-order or principal-symbol perturbation,
provided the entropy identity and uniform moment stock are proved
for the physical fibre.

\subsection{Retained maps and cofinal budgets}

For a positive trace-preserving retained-fibre map \({\cal R}\),
trace-norm contraction passes either (58.6) or (58.11) directly to
\({\cal R}\rho-{\cal R}\sigma\).  At level \(n\), define
\begin{equation}
 \varepsilon_n^{\rm tail}
 :=
 \inf_{c\in\mathbb R,\ R>0}
 {\beta_n R+
 \sqrt{\beta_n^2R^2+4\beta_n\tau_{n,c}(R)}
 \over2}.                                         \tag{58.15}
\end{equation}
Then
\[
 \sum_n\varepsilon_n^{\rm tail}<\infty
\]
is sufficient for the unbounded-Hamiltonian contribution to a
telescoping retained-state budget.  The infimum in (58.15) is a
mathematical quantity; an executable certificate must give explicit
\(c,R\) and an upper bound for the corresponding thermal tail.

\subsection{Acquisition required from an elliptic SM fibre}

To apply V58 to a metric- or connection-dependent physical
Hamiltonian, one must prove:
\begin{enumerate}
\item self-adjoint realizations of \(H\), \(K\), and a self-adjoint
representative \(A\) of their form difference;
\item trace-class heat kernels at the selected inverse temperature;
\item the entropy identity (58.1), including finiteness of both
relative entropies;
\item a scalar centre \(c\) and a uniform estimate of either
\(\tau_c(R)\) or \(M_p(c)\);
\item compatibility of the retained map with the physical legal
packet; and
\item summability of the resulting levelwise stock.
\end{enumerate}
Heat-kernel or elliptic estimates are expected to supply item 4.
The abstract entropy calculation does not supply them.

\begin{assessment}[Status after V58]
\label{ass:v12v58-status}
The bounded-difference obstruction in V57 has been separated into a
bounded scalar-quotient core and an explicit two-state thermal tail.
A finite thermal \(p\)-moment gives the quantitative unbounded
estimate (58.11), and retained positive trace-preserving maps
preserve it.  This creates a rigorous conditional route for
principal-part perturbations.
\end{assessment}

\begin{firewall}[Claim boundary after V58]
\label{fire:v12v58-claim}
Entropy admissibility is an additional hypothesis for an unbounded
form difference and has not been proved for the RPESM fibre.  No
uniform thermal moment, heat-kernel estimate, or self-adjoint
principal perturbation has been constructed.  The theorem does not
establish locality, Ward compatibility, reflection positivity,
determinant-line gluing, gravitational tightness, or the full
SM--Einstein continuum.
\end{firewall}

\subsection{Parent record}

This V58 note appends to the validated V57 file with record
\[
\begin{array}{c|c}
\text{lines}&16352\\
\text{bytes}&542144\\
\text{SHA--256}&
\texttt{8A544C1CF8AF2E84692EB05D22DA79154C849C6326A5DF3D3C1B204E5C05A9BA}.
\end{array}
\]

% =============================================================================
% END APPENDED TWELFTH-CYCLE V58 MATERIAL
% =============================================================================
% =============================================================================
% BEGIN APPENDED TWELFTH-CYCLE V59 MATERIAL
% =============================================================================

\section{Twelfth-cycle V59: a high-temperature certificate for
entropy admissibility}
\label{sec:v12v59}

V58 treated the symmetrized Gibbs identity for an unbounded
difference as an acquisition condition.  This note gives a concrete
sufficient criterion.  A small high-temperature trace margin makes
the own-Hamiltonian energies integrable, while a two-state thermal
first moment of the difference makes the cross energies integrable.
The relative-entropy identity then follows directly.

\subsection{Energy moments from a heat-trace margin}

\begin{lemma}[High-temperature margin]
\label{lem:v12v59-heat-margin}
Let \(L\) be self-adjoint and bounded below.  If
\[
 \operatorname{Tr}e^{-\beta_0L}<\infty
 \qquad\text{for some }0<\beta_0<\beta,
\]
then, for every real \(p>0\),
\begin{equation}
 \operatorname{Tr}\{|L|^p e^{-\beta L}\}<\infty.  \tag{59.1}
\end{equation}
In particular, the Gibbs state
\(\Gamma_\beta(L)\) has a finite absolute \(L\)-energy.
\end{lemma}

\begin{proof}
Let \(\ell_0\) be a lower bound for \(L\).  On
\([\ell_0,0]\), the function
\(|\lambda|^pe^{-(\beta-\beta_0)\lambda}\) is bounded.  On
\([0,\infty)\), it is also bounded because exponential decay
dominates the polynomial.  Hence there is a finite constant
\(C=C(p,\beta-\beta_0,\ell_0)\) such that
\[
 |\lambda|^pe^{-\beta\lambda}
 \leq C e^{-\beta_0\lambda},
 \qquad\lambda\geq\ell_0.
\]
Functional calculus and the positive trace give
\[
 \operatorname{Tr}\{|L|^pe^{-\beta L}\}
 \leq C\operatorname{Tr}e^{-\beta_0L}<\infty.
\]
\end{proof}

The margin \(\beta-\beta_0\) can be arbitrarily small but must be
strictly positive.  Trace class only at the single endpoint
\(\beta\) does not by itself control an inserted power of the
energy.

\subsection{A common-domain criterion}

Let \(H,K\) be self-adjoint, bounded below, and have the same
operator domain \({\cal D}\).  Suppose
\[
 (K-H)\psi=A\psi,\qquad\psi\in{\cal D},            \tag{59.2}
\]
where the difference is essentially self-adjoint on \({\cal D}\)
and \(A\) denotes its self-adjoint closure.  Assume that for some
\(0<\beta_0<\beta\),
\begin{equation}
 \operatorname{Tr}e^{-\beta_0H}
 +
 \operatorname{Tr}e^{-\beta_0K}<\infty.           \tag{59.3}
\end{equation}
Set
\[
 \rho=\Gamma_\beta(H),\qquad
 \sigma=\Gamma_\beta(K).
\]

\begin{theorem}[Thermal first-moment entropy certificate]
\label{thm:v12v59-certificate}
Under (59.2)--(59.3), suppose
\begin{equation}
 \operatorname{Tr}(\rho|A|)
 +
 \operatorname{Tr}(\sigma|A|)<\infty.             \tag{59.4}
\end{equation}
Then \(A\) is entropy admissible in the sense of V58.  More
explicitly,
\begin{align}
 D(\rho\|\sigma)
 &=
 \beta\operatorname{Tr}(\rho A)
 +\log{Z_K\over Z_H},                             \tag{59.5}\\
 D(\sigma\|\rho)
 &=
 -\beta\operatorname{Tr}(\sigma A)
 +\log{Z_H\over Z_K},                             \tag{59.6}\\
 D(\rho\|\sigma)+D(\sigma\|\rho)
 &=
 \beta\operatorname{Tr}\{A(\rho-\sigma)\}.        \tag{59.7}
\end{align}
All terms in (59.5)--(59.7) are finite.
\end{theorem}

\begin{proof}
Lemma \ref{lem:v12v59-heat-margin}, applied with \(p=1\), gives
\begin{equation}
 \operatorname{Tr}(\rho|H|)<\infty,
 \qquad
 \operatorname{Tr}(\sigma|K|)<\infty.             \tag{59.8}
\end{equation}
The spectral eigenvectors of the trace-class heat kernels lie in the
corresponding operator domains.  The common-domain identity (59.2),
together with (59.4) and (59.8), therefore makes the cross quadratic
expectations absolutely summable and gives
\begin{align}
 \operatorname{Tr}(\rho K)
 &=
 \operatorname{Tr}(\rho H)+\operatorname{Tr}(\rho A),
                                                        \tag{59.9}\\
 \operatorname{Tr}(\sigma H)
 &=
 \operatorname{Tr}(\sigma K)-\operatorname{Tr}(\sigma A).
                                                        \tag{59.10}
\end{align}
For example, in an \(H\)-eigenbasis the absolute sum of the
\(A\)-quadratic expectations is bounded by
\(\operatorname{Tr}(\rho|A|)\); the \(H\) part is controlled by
(59.8).  The \(K\)-eigenbasis proves the second identity.

Functional calculus gives
\[
 \log\rho=-\beta H-(\log Z_H)I,\qquad
 \log\sigma=-\beta K-(\log Z_K)I.
\]
The finiteness just proved permits these logarithms to be paired with
both states without an indeterminate difference of infinities.
Thus
\begin{align*}
 D(\rho\|\sigma)
 &=
 -\beta\operatorname{Tr}(\rho H)-\log Z_H
 +\beta\operatorname{Tr}(\rho K)+\log Z_K\\
 &=
 \beta\operatorname{Tr}(\rho A)+\log(Z_K/Z_H),
\end{align*}
which is (59.5).  The same calculation with the states reversed and
(59.10) proves (59.6).  Adding them proves (59.7).
\end{proof}

\subsection{The V58 moment stock already implies the first moment}

\begin{corollary}[Moment-to-admissibility bridge]
\label{cor:v12v59-moment-bridge}
Assume (59.2)--(59.3).  If for some \(p>1\) and \(c\in\mathbb R\)
the V58 stock
\[
 M_p(c)
 =
 \operatorname{Tr}\{\rho|A-cI|^p\}
 +
 \operatorname{Tr}\{\sigma|A-cI|^p\}
\]
is finite, then (59.4) holds.  Consequently \(A\) is entropy
admissible and the moment estimate (58.11) applies without assuming
(58.1) separately.
\end{corollary}

\begin{proof}
For either state \(\omega\), the spectral measure of
\(|A-cI|\) under \(\omega\) is a probability measure.  Classical
Hölder inequality on that measure gives
\[
 \operatorname{Tr}\{\omega|A-cI|\}
 \leq
 \left[
 \operatorname{Tr}\{\omega|A-cI|^p\}
 \right]^{1/p}.
\]
Since
\[
                         |A|
 \preceq |A-cI|+|c|I
\]
in quadratic-form expectation, the two absolute first moments of
\(A\) are finite.  Apply Theorem
\ref{thm:v12v59-certificate}.
\end{proof}

Combining the corollary with V58 gives the explicit conclusion
\begin{equation}
 \|\Gamma_\beta(H)-\Gamma_\beta(K)\|_1
 \leq
 \inf_{c\in\mathbb R}
 \left\{\beta C_pM_p(c)^{1/p}\right\}^{p/(p+1)}
                                                        \tag{59.11}
\end{equation}
under the common-domain and high-temperature hypotheses.

\subsection{A form-bounded first-moment shortcut}

Sometimes the full moment is not yet available, but the form
difference satisfies
\begin{equation}
 |\langle\psi,A\psi\rangle|
 \leq
 a\langle\psi,(H+bI)\psi\rangle,
 \qquad\psi\in{\cal D},                            \tag{59.12}
\end{equation}
after shifting \(H+bI\succeq0\).  Equation (59.12) controls
\(|\operatorname{Tr}(\rho A)|\), but it does not by itself imply
\(\operatorname{Tr}(\rho|A|)<\infty\), nor does it control the
\(\sigma\)-expectation.  A legal use of Theorem
\ref{thm:v12v59-certificate} therefore needs either operator-form
domination
\[
                         |A|\preceq a(H+bI)
                                                        \tag{59.13}
\]
and its \(K\)-side analogue, or a direct proof of (59.4).
This distinction prevents a signed form estimate from being
mistaken for the absolute thermal tail required in V58.

\subsection{Cofinal acquisition card}

At each refinement level \(n\), the unbounded route is now reduced
to four checkable analytic stocks:
\begin{enumerate}
\item a common self-adjoint domain and the operator identity
\(K_n-H_n=A_n\);
\item a high-temperature margin
\(\beta_{0,n}<\beta_n\) with both heat traces finite;
\item a scalar centre \(c_n\) and finite two-state thermal
\(p\)-moment \(M_{p,n}(c_n)\);
\item summability of the resulting bound in (59.11), after the
physical retained map is applied.
\end{enumerate}
Items 1--3 prove entropy admissibility and the tail estimate
together; no independent relative-entropy hypothesis remains.

\begin{assessment}[Status after V59]
\label{ass:v12v59-status}
A common-domain unbounded Gibbs perturbation with a
high-temperature heat-trace margin and a finite two-state thermal
first moment satisfies the exact symmetrized entropy identity.
The \(p>1\) moment already required by V58 supplies that first
moment, so (59.11) is now a self-contained conditional theorem rather
than an entropy-identity assumption.
\end{assessment}

\begin{firewall}[Claim boundary after V59]
\label{fire:v12v59-claim}
Metric-dependent principal operators may fail the common operator
domain hypothesis and instead share only form domains after unitary
transport.  No RPESM high-temperature trace margin, thermal moment,
or transported common-domain realization has been proved.  The
result does not establish locality, Ward compatibility, reflection
positivity, determinant-line gluing, gravitational tightness, or
the full SM--Einstein continuum.
\end{firewall}

\subsection{Parent record}

This V59 note appends to the validated V58 file with record
\[
\begin{array}{c|c}
\text{lines}&16653\\
\text{bytes}&551490\\
\text{SHA--256}&
\texttt{4AB481DEB3775DB00DA5EFE0CD36B362AB5876D5764FC13442FC060CF3DBE668}.
\end{array}
\]

% =============================================================================
% END APPENDED TWELFTH-CYCLE V59 MATERIAL
% =============================================================================
% =============================================================================
% BEGIN APPENDED TWELFTH-CYCLE V60 MATERIAL
% =============================================================================

\section{Twelfth-cycle V60: mandatory companion audit}
\label{sec:v12v60}

This is the required five-version review after V55.  It is performed
against the newest active companion files before addressing the
varying-Hilbert-space obstruction recorded in V59.

\subsection{Files inspected}

The active companion records are
\[
\begin{array}{c|r|r}
\text{file}&\text{lines}&\text{bytes}\\ \hline
\texttt{Renewal\_Yang\_Mills\_Mass\_Gap\_Research\_Notes\_V56.tex}
 &21223&723854\\
\texttt{Renewal\_NS\_Stability\_Research\_Notes\_V26.tex}
 &5319&278283.
\end{array}                                      \tag{60.1}
\]
Their SHA--256 digests at this audit are
\[
\begin{split}
 &\texttt{627247BCFC1A511C5158B1D3DED869FCA598501597EF3B8CF174C3A0BFEF54DC},\\
 &\texttt{82C887F8C2C69E4F28FAD7C3DC8DE5B9099CB5A4BF431EBA1FFF3E91935978AD}.
\end{split}                                      \tag{60.2}
\]

\subsection{Yang--Mills V56}

YM V56 proves that the parity-mask algebra used at the sixteen
\(M=2\) corners is a proper quotient of the full Laurent algebra.
It constructs an exact Laurent lift of one Wilson quartic response
matrix and verifies that reducing the lifted exponents modulo two
recovers every corner value.  The endpoint-completion and
cubic-bubble lifts remain open there.

The result changes one interpretation of SM V54.  Exact dyadic cell
averages and centre samples can approximate a known continuous
Gibbs-shape field in trace norm.  Finite samples alone do not
determine its Laurent/Fourier coefficients or its local Wilson
symbol.  Therefore the cell route may certify convergence of states,
but any assertion about a signed Poisson core, momentum locality, or
a Wilsonian symbol still requires an exponent-resolved lift or an
independent analytic representation.

This is consistent with V47--V48, where Fourier coefficients were
derived from the full function rather than inferred from grid
values.  No SM formula or numerical constant is imported from the
YM compiler.

\subsection{Navier--Stokes V26}

NS remains at V26.  Its pointwise rank-one improvement applies only
after the input is proved to lie in the rank-one cone.  The V55
transfer remains unchanged: a common-domain, block, tensor, or
transported-fibre reduction must be proved before a norm specialized
to that structure is used.  It supplies no operator-domain,
heat-kernel, entropy, or Standard Model estimate.

\subsection{Audit decision}

\begin{decision}[Direction after the V60 audit]
\label{dec:v12v60-direction}
Preserve the two approximation routes as distinct:
\begin{enumerate}
\item use V47--V52 only when the complete parameter dependence is
available in an exponent-resolved analytic or differentiable form;
\item use V54 and V57 for state convergence from exact cell
averages or certified samples, without inferring a Wilson symbol
from those samples;
\item next construct a coherent unitary identification of
metric-dependent Hilbert fibres, transport their Gibbs states and
Hamiltonians to a common fibre, and quantify both Hamiltonian and
transport-cocycle defects.
\end{enumerate}
\end{decision}

\begin{assessment}[Status after V60]
\label{ass:v12v60-status}
The mandatory audit is complete.  YM V56 exposes the information
lost by a finite sampling quotient; NS V26 again enforces proof of
structured membership.  The next upstream task is the
metric-dependent Hilbert-fibre identification required before the
common-domain entropy theorem V59 can be applied.
\end{assessment}

\begin{firewall}[Claim boundary after V60]
\label{fire:v12v60-claim}
No YM Laurent coefficient or NS kernel constant is used as an SM
input.  The audit does not construct a spinor-bundle
trivialization, prove a common elliptic domain, or identify the
physical Wilson symbol.  The full SM--Einstein continuum remains
open.
\end{firewall}

\subsection{Parent record}

This V60 note appends to the validated V59 file with record
\[
\begin{array}{c|c}
\text{lines}&16928\\
\text{bytes}&560561\\
\text{SHA--256}&
\texttt{63E9AF4A7EC3C089879CD8CE00E7FFFD4A3E5CFAE36BC206D8A619FF81501271}.
\end{array}
\]

% =============================================================================
% END APPENDED TWELFTH-CYCLE V60 MATERIAL
% =============================================================================
% =============================================================================
% BEGIN APPENDED TWELFTH-CYCLE V61 MATERIAL
% =============================================================================

\section{Twelfth-cycle V61: coherent unitary transport of
metric-dependent Gibbs fibres}
\label{sec:v12v61}

V59 requires two Hamiltonians to act on one Hilbert space with a
common domain.  In the Einstein sector, the natural \(L^2\) space
and spinor inner product depend on the metric.  Comparing operators
before identifying those Hilbert fibres is meaningless.  This note
introduces the required unitary transport, proves its interaction
with Gibbs normalization and entropy bounds, quantifies a failure
of the transport cocycle, and gives a trace-class limit theorem for
a cofinal chain of transported fibre states.

\subsection{Pairwise transport}

For each parameter \(x\), let \({\cal H}_x\) be a Hilbert space,
\(H_x\) a self-adjoint operator on it, and
\[
 \rho_x:=\Gamma_{\beta_x}(H_x)
\]
a well-defined Gibbs state.  Let
\[
 T_{x\leftarrow y}:{\cal H}_y\longrightarrow{\cal H}_x
\]
be unitary.  Define
\begin{align}
 H_{y\to x}
 &:=
 T_{x\leftarrow y}H_yT_{x\leftarrow y}^*,        \tag{61.1}\\
 \rho_{y\to x}
 &:=
 T_{x\leftarrow y}\rho_yT_{x\leftarrow y}^*.     \tag{61.2}
\end{align}

\begin{lemma}[Functional calculus commutes with transport]
\label{lem:v12v61-functional}
For every bounded Borel function \(f\),
\[
 f(H_{y\to x})
 =
 T_{x\leftarrow y}f(H_y)T_{x\leftarrow y}^*.
\]
In particular, if \(\beta_x=\beta_y=\beta\), then
\begin{equation}
 \rho_{y\to x}=\Gamma_\beta(H_{y\to x}),\qquad
 \operatorname{Tr}e^{-\beta H_{y\to x}}
 =
 \operatorname{Tr}e^{-\beta H_y}.                 \tag{61.3}
\end{equation}
\end{lemma}

\begin{proof}
The spectral measure of \(H_{y\to x}\) is
\(T_{x\leftarrow y}E_{H_y}(\,\cdot\,)
T_{x\leftarrow y}^*\).  Integration against this measure proves the
functional-calculus identity.  Apply it to the heat exponential and
use unitary invariance of the trace.
\end{proof}

Thus the comparison
\begin{equation}
 d_T(x,y):=
 \|\rho_x-\rho_{y\to x}\|_1                       \tag{61.4}
\end{equation}
is a trace distance between operators on one Hilbert space.

\subsection{Transported Hamiltonian estimates}

Assume first that \(\beta_x=\beta_y=\beta\), that
\(H_x\) and \(H_{y\to x}\) have a common domain, and that
\begin{equation}
 A_{x,y}:=H_{y\to x}-H_x                          \tag{61.5}
\end{equation}
is bounded and self-adjoint.

\begin{theorem}[Bounded transported-fibre estimate]
\label{thm:v12v61-bounded}
Under these hypotheses,
\begin{equation}
 d_T(x,y)
 \leq
 \beta\|[A_{x,y}]\|_{\rm q}.                      \tag{61.6}
\end{equation}
\end{theorem}

\begin{proof}
By Lemma \ref{lem:v12v61-functional}, both states in (61.4) are
Gibbs states on \({\cal H}_x\) for the Hamiltonians \(H_x\) and
\(H_{y\to x}\).  Apply Theorem
\ref{thm:v12v57-sharp-Gibbs}.
\end{proof}

If \(A_{x,y}\) is unbounded but satisfies the common-domain,
high-temperature, and thermal-moment hypotheses of V59 after
transport, then V58--V59 instead give
\begin{equation}
 d_T(x,y)
 \leq
 \inf_{c\in\mathbb R}
 \left\{\beta C_pM_{p;x,y}(c)^{1/p}\right\}^{p/(p+1)},        \tag{61.7}
\end{equation}
where
\begin{align}
 M_{p;x,y}(c)
 :={}&
 \operatorname{Tr}\{\rho_x|A_{x,y}-cI|^p\}
 \nonumber\\
 &+
 \operatorname{Tr}\{\rho_{y\to x}|A_{x,y}-cI|^p\}.           \tag{61.8}
\end{align}
All domain and moment assertions in (61.7) are made after the
unitary identification.  Fibrewise formulas before transport do not
supply them.

\subsection{Coherent transports and cocycle defects}

Exact pairwise transports should obey
\begin{equation}
 T_{x\leftarrow x}=I,\qquad
 T_{x\leftarrow y}^*=T_{y\leftarrow x},\qquad
 T_{x\leftarrow z}
 =
 T_{x\leftarrow y}T_{y\leftarrow z}.              \tag{61.9}
\end{equation}
For a possibly approximate system define
\begin{equation}
 \kappa(x,y,z)
 :=
 \|T_{x\leftarrow z}
   -T_{x\leftarrow y}T_{y\leftarrow z}\|.          \tag{61.10}
\end{equation}

\begin{lemma}[Unitary conjugation defect]
\label{lem:v12v61-unitary-defect}
If \(U,V:{\cal K}\to{\cal H}\) are unitary and \(\tau\) is a density
operator on \({\cal K}\), then
\begin{equation}
 \|U\tau U^*-V\tau V^*\|_1
 \leq2\|U-V\|.                                    \tag{61.11}
\end{equation}
\end{lemma}

\begin{proof}
Write
\[
 U\tau U^*-V\tau V^*
 =
 (U-V)\tau U^*+V\tau(U^*-V^*).
\]
The trace ideal inequality and
\(\|\tau\|_1=\|U\|=\|V\|=1\) bound each term by
\(\|U-V\|\).
\end{proof}

\begin{proposition}[Transported triangle inequality]
\label{prop:v12v61-triangle}
For any three fibre states,
\begin{equation}
 d_T(x,z)
 \leq
 d_T(x,y)+d_T(y,z)+2\kappa(x,y,z).                \tag{61.12}
\end{equation}
\end{proposition}

\begin{proof}
Insert, on \({\cal H}_x\), the two intermediate states
\[
 T_{x\leftarrow y}\rho_yT_{x\leftarrow y}^*
\]
and
\[
 T_{x\leftarrow y}T_{y\leftarrow z}\rho_z
 T_{y\leftarrow z}^*T_{x\leftarrow y}^*.
\]
The first difference is \(d_T(x,y)\).  Unitary invariance of the
trace norm makes the second difference \(d_T(y,z)\).  Lemma
\ref{lem:v12v61-unitary-defect} bounds the last difference by
\(2\kappa(x,y,z)\).
\end{proof}

A nonzero \(\kappa\) is a transport error, separate from the
Hamiltonian Gibbs error.  It cannot be absorbed into a scalar energy
shift.

\subsection{Exact trivializations}

Let \({\cal H}_*\) be a reference Hilbert space and suppose unitaries
\[
 U_x:{\cal H}_x\longrightarrow{\cal H}_*
\]
are chosen.  Then
\begin{equation}
 T_{x\leftarrow y}:=U_x^*U_y                     \tag{61.13}
\end{equation}
satisfies the exact cocycle (61.9).  The transported reference-space
operators are
\begin{equation}
 \widehat H_x:=U_xH_xU_x^*,\qquad
 \widehat\rho_x:=U_x\rho_xU_x^*.                  \tag{61.14}
\end{equation}
Moreover,
\begin{equation}
 d_T(x,y)=\|\widehat\rho_x-\widehat\rho_y\|_1,
 \qquad
 \|[A_{x,y}]\|_{\rm q}
 =
 \|[\widehat H_y-\widehat H_x]\|_{\rm q}.          \tag{61.15}
\end{equation}

\begin{proof}
Equation (61.9) follows by cancellation of \(U_yU_y^*=I\).
Conjugating both terms in (61.4) by \(U_x\) gives the first identity
in (61.15).  Conjugating (61.5) by \(U_x\) gives
\(\widehat H_y-\widehat H_x\), and the quotient norm is unitarily
invariant.
\end{proof}

If a second trivialization is
\(\widetilde U_x=WU_x\) with one fixed unitary \(W\), all quantities
in (61.15) are unchanged.  An \(x\)-dependent \(W_x\) changes the
identification between distinct fibres and contributes a connection
term; it is not a harmless simultaneous conjugation.

\subsection{The metric-density unitary}

Let \(M\) be compact, let \(d\mu_g=J_g\,d\mu_0\), and suppose
\(B_g:E_g\to E_0\) is a measurable fibrewise unitary identification
of the relevant finite-rank Hermitian bundles.  Define
\begin{equation}
 U_g:L^2(M,E_g,d\mu_g)\longrightarrow
     L^2(M,E_0,d\mu_0),
 \qquad
 (U_g\psi)(x):=J_g(x)^{1/2}B_g(x)\psi(x).          \tag{61.16}
\end{equation}

\begin{proposition}[Metric \(L^2\) trivialization]
\label{prop:v12v61-metric-unitary}
The map \(U_g\) in (61.16) is unitary.  The transports
\(T_{g\leftarrow h}=U_g^*U_h\) obey the exact cocycle.
\end{proposition}

\begin{proof}
Fibrewise unitarity of \(B_g\) gives
\[
 \|U_g\psi\|_{L^2(d\mu_0)}^2
 =
 \int_M J_g|\psi|_{E_g}^2\,d\mu_0
 =
 \int_M|\psi|_{E_g}^2\,d\mu_g.
\]
The inverse is \(B_g^*J_g^{-1/2}\), so \(U_g\) is onto.  Equation
(61.13) proves the cocycle.
\end{proof}

This proposition resolves the measure-space mismatch once the bundle
identifications \(B_g\) have been chosen.  It does not prove that
the transported differential operators have a common domain.

For a differentiable fixed-space realization
\(\widehat H_s=U_sH_sU_s^*\), formal differentiation gives
\begin{equation}
 {d\widehat H_s\over ds}
 =
 U_s\dot H_sU_s^*
 +
 [\Omega_s,\widehat H_s],
 \qquad
 \Omega_s:=\dot U_sU_s^*,\quad\Omega_s^*=-\Omega_s.           \tag{61.17}
\end{equation}
Equation (61.17) is rigorous on a common invariant core when both
terms extend there.  The commutator is the connection contribution
created by the moving trivialization; omitting it gives the wrong
transported perturbation.

\subsection{A cofinal trace-class limit}

Consider Hilbert spaces \({\cal H}_n\), density operators \(\rho_n\),
and adjacent unitaries
\[
 T_{n\leftarrow n+1}:{\cal H}_{n+1}\to{\cal H}_n.
\]
Define
\begin{equation}
 e_n:=
 \|\rho_n-
 T_{n\leftarrow n+1}\rho_{n+1}
 T_{n\leftarrow n+1}^*\|_1.                       \tag{61.18}
\end{equation}
Let
\[
 U_0=I,\qquad
 U_n:=T_{0\leftarrow1}T_{1\leftarrow2}\cdots
      T_{n-1\leftarrow n}:{\cal H}_n\to{\cal H}_0,
\]
and set \(\widehat\rho_n=U_n\rho_nU_n^*\).

\begin{theorem}[Transported fibre-state continuum]
\label{thm:v12v61-limit}
If
\begin{equation}
                         \sum_{n=0}^\infty e_n<\infty,        \tag{61.19}
\end{equation}
then \(\widehat\rho_n\) converges in trace norm to a positive
trace-one operator \(\rho_\infty\) on \({\cal H}_0\).  Quantitatively,
\begin{equation}
 \|\widehat\rho_m-\rho_\infty\|_1
 \leq\sum_{n=m}^\infty e_n.                       \tag{61.20}
\end{equation}
\end{theorem}

\begin{proof}
Since
\(U_{n+1}=U_nT_{n\leftarrow n+1}\), unitary invariance gives
\[
 \|\widehat\rho_{n+1}-\widehat\rho_n\|_1=e_n.
\]
Thus (61.19) makes \((\widehat\rho_n)\) Cauchy in the Banach space
\({\mathfrak S}_1({\cal H}_0)\), and summing the tail proves
(61.20).  The positive cone is trace-norm closed and the trace is
trace-norm continuous.  Since every \(\widehat\rho_n\) is positive
with trace one, the limit has the same properties.
\end{proof}

The adjacent errors \(e_n\) may be bounded by (61.6) in the bounded
case or by (61.7) in the thermal-tail case.  If independently
constructed direct transports are used instead of the composite
\(U_n\), their cocycle defects must also be summed through
(61.12).

\begin{assessment}[Status after V61]
\label{ass:v12v61-status}
Metric-dependent Gibbs states can now be compared only after an
explicit unitary fibre identification.  Hamiltonian errors and
transport-cocycle errors have separate quantitative bounds.  A
summable adjacent transported-state defect produces a genuine
trace-class continuum state on a reference Hilbert space, with an
explicit tail rate.
\end{assessment}

\begin{firewall}[Claim boundary after V61]
\label{fire:v12v61-claim}
The metric-density unitary does not construct the spinor-bundle
identifications \(B_g\), prove a common Sobolev domain for the
transported RPESM operators, or bound the commutator in (61.17).
The limit in Theorem \ref{thm:v12v61-limit} concerns one chain of
finite-fibre density operators; it is not the full field measure,
does not establish Ward or reflection-positive compatibility, and
does not prove gravitational tightness, determinant-line gluing, or
the full SM--Einstein continuum.
\end{firewall}

\subsection{Parent record}

This V61 note appends to the validated V60 file with record
\[
\begin{array}{c|c}
\text{lines}&17042\\
\text{bytes}&565048\\
\text{SHA--256}&
\texttt{607CA86F7A140A600E7DE1E9A71EFF1C1E3377A6C1D6E9B0E823BB9BD26C55B6}.
\end{array}
\]

% =============================================================================
% END APPENDED TWELFTH-CYCLE V61 MATERIAL
% =============================================================================
% =============================================================================
% BEGIN APPENDED TWELFTH-CYCLE V62 MATERIAL
% =============================================================================

\section{Twelfth-cycle V62: common Sobolev domains and heat margins
for transported elliptic fibres}
\label{sec:v12v62}

V61 places metric-dependent operators on one reference \(L^2\)
space, but unitary transport alone does not identify their operator
domains.  This note gives a uniform elliptic certificate for a
common \(H^2\) domain, compact resolvent, and heat trace at every
positive inverse temperature.  It then states the additional
relative-power estimate that turns an unbounded elliptic difference
into the V58--V59 thermal-moment stock.

\subsection{Reference geometry and coefficient class}

Let \(M\) be a closed \(d\)-dimensional manifold, let \(E\to M\) be
a finite-rank Hermitian bundle, and work on
\[
 {\cal H}_0:=L^2(M,E,d\mu_0).
\]
Fix a Hermitian connection \(\nabla\) and a positive reference
Laplace-type operator
\[
 \Lambda:=I+\nabla^*\nabla,
 \qquad
 {\cal D}(\Lambda)=H^2(M,E).
\]
Its form domain is \(H^1(M,E)\).  Write its eigenvalues as
\(\nu_j\), in nondecreasing order.  There are constants
\(c_0>0\) and \(c_1\geq0\) such that
\begin{equation}
                         \nu_j\geq c_0j^{2/d}-c_1.
                                                        \tag{62.1}
\end{equation}

Let \(G\) be a parameter set.  For \(x\in G\), consider a formally
self-adjoint, strongly elliptic, second-order differential expression
on \(C^\infty(M,E)\),
\begin{equation}
 L_xu
 =
 -\nabla_i(a_x^{ij}\nabla_j u)
 +b_x^j\nabla_j u+c_xu.                           \tag{62.2}
\end{equation}
The displayed form includes bundle-matrix coefficients and is
understood after absorbing the fixed reference density into the
lower-order terms.  Assume:
\begin{enumerate}
\item the principal coefficients are uniformly Legendre positive,
\begin{equation}
 \operatorname{Re}
 \langle a_x^{ij}(q)\xi_jv,\xi_iv\rangle
 \geq a_0|\xi|^2|v|^2                            \tag{62.3}
\end{equation}
for one \(a_0>0\);
\item the \(W^{1,\infty}\) norms of \(a_x^{ij}\) and the
\(L^\infty\) norms of \(b_x^j,c_x\) are bounded uniformly in \(x\);
and
\item the formal-symmetry relations for (62.2) hold with respect to
the fixed \(L^2\) inner product.
\end{enumerate}

\subsection{Uniform realization theorem}

\begin{theorem}[Common transported elliptic domain]
\label{thm:v12v62-common-domain}
Under (62.2)--(62.3), the Friedrichs realization, still denoted
\(L_x\), is self-adjoint and bounded below.  Its operator domain is
\begin{equation}
                         {\cal D}(L_x)=H^2(M,E)    \tag{62.4}
\end{equation}
for every \(x\in G\).  Moreover there are constants
\(C_{\rm gr}\geq1\) and \(B_{\rm gr}\geq0\), independent of \(x\),
such that
\begin{equation}
 C_{\rm gr}^{-1}\|u\|_{H^2}
 \leq
 \|L_xu\|_2+(B_{\rm gr}+1)\|u\|_2
 \leq
 C_{\rm gr}\|u\|_{H^2}.                           \tag{62.5}
\end{equation}
Each \(L_x\) has compact resolvent.
\end{theorem}

\begin{proof}
Uniform Legendre positivity, integration by parts, and absorption of
the first-order terms give a G{\aa}rding inequality
\begin{equation}
 \langle u,L_xu\rangle
 \geq
 a_1\|u\|_{H^1}^2-b_1\|u\|_2^2                  \tag{62.6}
\end{equation}
with \(a_1>0,b_1\geq0\) independent of \(x\).  The closed
semibounded form on \(H^1(M,E)\) therefore has a self-adjoint
Friedrichs realization.

For completeness, the regularity step is local.  In a coordinate
and bundle chart, apply first difference quotients to the weak
equation
\[
                         L_xu=f\in L^2.
\]
Testing against the corresponding difference quotient of a compactly
supported \(u\), the principal term controls all second weak
derivatives by (62.3).  Commuting a difference quotient through
\(a_x^{ij}\) costs only its \(W^{1,\infty}\) norm; the lower-order
terms are absorbed using (62.6).  A finite partition of unity on the
closed manifold yields, uniformly in \(x\),
\begin{equation}
 \|u\|_{H^2}
 \leq C\{\|L_xu\|_2+\|u\|_2\}.                   \tag{62.7}
\end{equation}
Thus every vector in the Friedrichs operator domain lies in \(H^2\).
Conversely, the coefficient bounds in (62.2) make
\(L_x:H^2\to L^2\) bounded, and density of smooth sections identifies
the operator action with the Friedrichs realization.  This proves
(62.4) and the left side of (62.5); the mapping bound proves the
right side after enlarging the constants.

Finally, (62.5) makes the resolvent a bounded map from \(L^2\) into
\(H^2\).  The Rellich embedding \(H^2(M,E)\hookrightarrow L^2(M,E)\)
is compact, so the resolvent is compact.
\end{proof}

The theorem is stated for a closed manifold.  A boundary would
require one fixed elliptic boundary condition and a separate
uniform complementing-condition estimate.

\subsection{Uniform heat traces}

The same coefficient assumptions give constants \(a_2>0,b_2\geq0\)
such that, as quadratic forms on \(H^1\),
\begin{equation}
                         L_x\ \succeq\ a_2\Lambda-b_2I.       \tag{62.8}
\end{equation}

\begin{corollary}[Heat trace at every positive temperature]
\label{cor:v12v62-heat}
For every \(\gamma>0\),
\begin{equation}
 \sup_{x\in G}\operatorname{Tr}e^{-\gamma L_x}
 \leq
 e^{\gamma b_2}
 \sum_{j=1}^\infty e^{-\gamma a_2\nu_j}
 <\infty.                                         \tag{62.9}
\end{equation}
Consequently, for any target \(\beta>0\), every
\(0<\beta_0<\beta\) supplies the high-temperature margin required
by V59, uniformly in \(x\).
\end{corollary}

\begin{proof}
The min--max principle applied to (62.8) gives
\[
 \lambda_j(L_x)\geq a_2\nu_j-b_2.
\]
Sum the exponentials to obtain the first inequality in (62.9).
Equation (62.1) makes the final series converge.
\end{proof}

The same comparison and Lemma \ref{lem:v12v59-heat-margin} give
uniform finite moments of each Hamiltonian in its own Gibbs state.
They do not yet give a moment of the difference of two Hamiltonians.

\subsection{Size of an elliptic difference}

For \(x,y\in G\), let
\[
                         A_{x,y}:=L_y-L_x
\]
on \(H^2(M,E)\).  Directly from (62.2),
\begin{equation}
 \|A_{x,y}u\|_2
 \leq
 C\,\Xi(x,y)\|u\|_{H^2},                          \tag{62.10}
\end{equation}
where one may take
\begin{align}
 \Xi(x,y):={}&
 \|a_y-a_x\|_{W^{1,\infty}}
 +\|b_y-b_x\|_{L^\infty}
 +\|c_y-c_x\|_{L^\infty}.                         \tag{62.11}
\end{align}
Thus \(A_{x,y}\) is bounded from the common graph domain to \(L^2\),
but it is generally unbounded on \(L^2\) when the principal
coefficients differ.

If \(a_y=a_x\) and \(b_y=b_x\), then
\(A_{x,y}=c_y-c_x\) is a bounded multiplication operator and V57
applies directly:
\begin{equation}
 \|\Gamma_\beta(L_y)-\Gamma_\beta(L_x)\|_1
 \leq
 \beta\|[c_y-c_x]\|_{\rm q}.                      \tag{62.12}
\end{equation}
Equality of only the principal coefficients leaves a first-order
difference and does not justify (62.12).

\subsection{Relative-power thermal certificate}

The following condition packages the additional elliptic or
pseudodifferential estimate needed for V58.  Suppose
\(A_{x,y}\) is essentially self-adjoint on \(H^2\), let \(A\) be its
closure, and choose lower bounds
\[
 L_x\geq\ell_xI,\qquad L_y\geq\ell_yI.
\]
For \(p>1\) and \(c\in\mathbb R\), assume the quadratic-form
inequalities
\begin{align}
 |A-cI|^p
 &\preceq
 C_x\{I+L_x-\ell_xI\}^p,                          \tag{62.13}\\
 |A-cI|^p
 &\preceq
 C_y\{I+L_y-\ell_yI\}^p.                          \tag{62.14}
\end{align}

\begin{theorem}[Elliptic moment handoff]
\label{thm:v12v62-moment-handoff}
Under (62.13)--(62.14), the two-state thermal moment satisfies
\begin{align}
 M_{p;x,y}(c)
 \leq{}&
 {C_x\over Z_x}
 \operatorname{Tr}\!\left[
 \{I+L_x-\ell_xI\}^p e^{-\beta L_x}\right]
 \nonumber\\
 &+
 {C_y\over Z_y}
 \operatorname{Tr}\!\left[
 \{I+L_y-\ell_yI\}^p e^{-\beta L_y}\right]
 <\infty.                                         \tag{62.15}
\end{align}
Consequently the entropy-admissibility theorem V59 and the
unbounded stability estimate V58 apply to the pair \(L_x,L_y\).
\end{theorem}

\begin{proof}
Take the expectation of (62.13) in
\(\Gamma_\beta(L_x)\) and of (62.14) in
\(\Gamma_\beta(L_y)\), then add.  This gives the displayed upper
bound.  Each trace is finite by Corollary
\ref{cor:v12v62-heat} and the scalar estimate used in Lemma
\ref{lem:v12v59-heat-margin}, with polynomial
\((1+\lambda-\ell)^p\).  Theorem
\ref{thm:v12v59-certificate} supplies the entropy identity, and
Corollary \ref{cor:v12v58-moment} gives the state estimate.
\end{proof}

The self-adjointness of the difference and the two inequalities
(62.13)--(62.14) are not consequences of the graph bound (62.10).
They are separate certificates.  This is the precise analytic work
still required when a metric variation changes the principal
symbol.

\subsection{Application after metric transport}

Let \(U_g\) be the metric-density unitary (61.16) and define
\[
                         L_g:=U_gH_gU_g^*
\]
on the fixed reference space.  Theorem
\ref{thm:v12v62-common-domain} applies if:
\begin{enumerate}
\item the transported expressions have the divergence form (62.2);
\item the transported metrics are uniformly elliptic relative to
the reference metric;
\item the metric, bundle-identification, connection, and potential
coefficients have the uniform regularity appearing above; and
\item formal symmetry is checked after including every density and
connection term generated by \(U_g\).
\end{enumerate}
Under these conditions the operator-domain and high-temperature
parts of V59 are no longer assumptions.  Only the self-adjoint
difference and relative-power moment certificates remain.

\begin{assessment}[Status after V62]
\label{ass:v12v62-status}
A uniformly elliptic transported family on a closed manifold has one
common \(H^2\) domain, uniform graph norms, compact resolvent, and a
uniform heat trace at every positive inverse temperature.  Bounded
potential changes fall directly under V57.  Principal-part changes
enter V58 once the explicit self-adjoint-difference and
relative-power inequalities (62.13)--(62.14) are supplied.
\end{assessment}

\begin{firewall}[Claim boundary after V62]
\label{fire:v12v62-claim}
The physical RPESM operator has not been written in the transported
coefficient form (62.2), and the required uniform metric class,
spinor identification, formal-symmetry check, and relative-power
inequalities have not been proved.  The note treats closed
manifolds and supplies no boundary condition.  It does not establish
Ward or reflection-positive compatibility, gravitational
tightness, determinant-line gluing, or the full SM--Einstein
continuum.
\end{firewall}

\subsection{Parent record}

This V62 note appends to the validated V61 file with record
\[
\begin{array}{c|c}
\text{lines}&17412\\
\text{bytes}&576342\\
\text{SHA--256}&
\texttt{205999C72BC0B071A4439E842B1E1D63ECD5413662E35FA5C032F82BA9C0C250}.
\end{array}
\]

% =============================================================================
% END APPENDED TWELFTH-CYCLE V62 MATERIAL
% =============================================================================
% =============================================================================
% BEGIN APPENDED TWELFTH-CYCLE V63 MATERIAL
% =============================================================================

\section{Twelfth-cycle V63: thermal control of squared Dirac
connection and Yukawa perturbations}
\label{sec:v12v63}

V62 reduces an unbounded elliptic comparison to a relative-power
thermal moment.  For a fixed metric and fixed Dirac principal
symbol, changing the gauge connection or a bounded Yukawa
endomorphism changes the Dirac operator by a zeroth-order operator.
After squaring, the Hamiltonian difference is first order.  This
note proves the required second-moment estimate explicitly.

\subsection{Fixed-principal-symbol hypotheses}

Let \(M\) be closed and let \(E\to M\) be a fixed Hermitian Dirac
bundle.  Let \(D_0,D_1\) be self-adjoint Dirac-type operators on
\(L^2(M,E)\) with the same principal symbol and common domain
\(H^1(M,E)\).  Assume
\begin{equation}
                         D_1=D_0+V,               \tag{63.1}
\end{equation}
where \(V=V^*\) is a bounded bundle endomorphism preserving
\(H^1\), and
\begin{equation}
 v:=\|V\|<\infty,\qquad
 c_1:=\|[D_0,V]\|<\infty.                         \tag{63.2}
\end{equation}
For a Dirac operator, the commutator in (63.2) is Clifford
multiplication by the covariant derivative of \(V\); a
\(W^{1,\infty}\) coefficient bound is sufficient.

Fix \(\mu^2\geq0\) and define positive elliptic Hamiltonians
\begin{equation}
 L_j:=D_j^2+\mu^2I,\qquad j=0,1.                  \tag{63.3}
\end{equation}
The uniform elliptic theorem V62 gives
\[
 {\cal D}(L_0)={\cal D}(L_1)=H^2(M,E)
\]
and trace-class heat kernels at every positive inverse temperature.

\subsection{The first-order square difference}

On \(H^2(M,E)\),
\begin{align}
 W
 &:=
 L_1-L_0                                             \nonumber\\
 &=
 D_0V+VD_0+V^2                                    \tag{63.4}\\
 &=
 2VD_0+[D_0,V]+V^2.                              \tag{63.5}
\end{align}
The same operator can be written relative to \(D_1\) as
\begin{equation}
 W=2VD_1+[D_1,V]-V^2,
 \qquad [D_1,V]=[D_0,V].                          \tag{63.6}
\end{equation}

\begin{lemma}[Self-adjoint first-order difference]
\label{lem:v12v63-selfadjoint}
Suppose the coefficients of \(D_0\) and \(V\) are smooth, or have
the standard Lipschitz regularity needed for first-order
commutator approximation.  Then the symmetric expression \(W\) on
\(C^\infty(M,E)\) is essentially self-adjoint.  Its closure contains
\(H^1(M,E)\) in its operator domain and is bounded from \(H^1\) to
\(L^2\).
\end{lemma}

\begin{proof}
Equations (63.4)--(63.5) show that \(W\) is a symmetric first-order
differential expression with bounded principal symbol on the compact
manifold.  Choose a Friedrichs smoothing family \(J_\varepsilon\), built from
local convolution mollifiers and a finite partition of unity.  It
maps \(L^2\) into smooth sections and converges strongly to the
identity.  The first-order Friedrichs commutator lemma for bounded
Lipschitz coefficients gives
\[
 \sup_{0<\varepsilon\leq1}\|[W,J_\varepsilon]\|<\infty,
 \qquad
 [W,J_\varepsilon]u\longrightarrow0
 \quad(u\in L^2).
\]
The second assertion follows first for smooth \(u\) by direct
differentiation of the mollifier and then for every \(L^2\) vector
from the uniform bound.

If \(u\in{\cal D}(W^*)\), distributional formal symmetry gives
\[
 WJ_\varepsilon u
 =
 J_\varepsilon W^*u+[W,J_\varepsilon]u.
\]
Thus \(J_\varepsilon u\to u\) and
\(WJ_\varepsilon u\to W^*u\), so smooth sections approximate every
adjoint-domain vector in the graph norm.  Hence
\(C^\infty\) is a core for \(W^*\), so the closure of \(W\) equals
\(W^*\).  Equations (63.5) and (63.2) make the differential expression
bounded from \(H^1\) to \(L^2\).  Approximation of an \(H^1\)
section by smooth sections therefore places \(H^1\) inside the
operator domain of the closure with the asserted mapping bound.
\end{proof}

The core argument is recorded because self-adjointness of a
difference is not implied merely by self-adjointness of its two
summands.

\subsection{Two-sided relative square bounds}

Define
\begin{equation}
 Q(v,c_1)
 :=
 \max\{8v^2,\ 2(c_1+v^2)^2\}.                    \tag{63.7}
\end{equation}

\begin{theorem}[Relative second-moment certificate]
\label{thm:v12v63-relative-square}
As quadratic forms on \(H^1(M,E)\),
\begin{align}
 |W|^2
 &\preceq
 Q(v,c_1)(I+D_0^2),                              \tag{63.8}\\
 |W|^2
 &\preceq
 Q(v,c_1)(I+D_1^2).                              \tag{63.9}
\end{align}
\end{theorem}

\begin{proof}
For \(u\in C^\infty(M,E)\), equation (63.5) gives
\[
 \|Wu\|_2
 \leq
 2v\|D_0u\|_2+(c_1+v^2)\|u\|_2.
\]
Using \((a+b)^2\leq2a^2+2b^2\),
\[
 \|Wu\|_2^2
 \leq
 8v^2\|D_0u\|_2^2
 +
 2(c_1+v^2)^2\|u\|_2^2
 \leq
 Q(v,c_1)\langle u,(I+D_0^2)u\rangle.
\]
Closure in the \(H^1\) form norm proves (63.8).  Equation (63.6)
gives the same estimate with \(D_1\), since
\(\|[D_1,V]\|=c_1\), proving (63.9).
\end{proof}

These are the V62 relative-power inequalities with \(p=2\), and are
stronger than the version with \((I+D_j^2)^2\) on the right.

\subsection{The thermal second moment}

Let
\[
 \rho_j:=\Gamma_\beta(L_j),\qquad
 {\cal E}_\beta
 :=
 2+\operatorname{Tr}(\rho_0D_0^2)
   +\operatorname{Tr}(\rho_1D_1^2).               \tag{63.10}
\]

\begin{corollary}[Finite two-state moment]
\label{cor:v12v63-moment}
The V58 moment with the common scalar centre \(c=0\) satisfies
\begin{equation}
 M_2(0)
 :=
 \operatorname{Tr}(\rho_0|W|^2)
 +
 \operatorname{Tr}(\rho_1|W|^2)
 \leq
 Q(v,c_1){\cal E}_\beta<\infty.                   \tag{63.11}
\end{equation}
\end{corollary}

\begin{proof}
Take the \(\rho_0\)-expectation of (63.8) and the
\(\rho_1\)-expectation of (63.9), then add.  Finiteness of the energy
terms follows from the high-temperature margin in Corollary
\ref{cor:v12v62-heat} and Lemma
\ref{lem:v12v59-heat-margin}.
\end{proof}

The energy stock can also be read from the partition functions:
\begin{equation}
 \operatorname{Tr}(\rho_jD_j^2)
 =
 -{d\over d\beta}\log
   \operatorname{Tr}e^{-\beta L_j}
 -\mu^2.                                          \tag{63.12}
\end{equation}
Differentiation is justified by the heat-trace margin.

\subsection{Gibbs-state comparison}

For \(p=2\), the V58 constant is
\[
                         C_2=2.
\]
The common domain, heat margin, self-adjoint difference, and moment
conditions are now all available.

\begin{theorem}[Dirac-square Gibbs stability]
\label{thm:v12v63-Gibbs}
Under the hypotheses of this note, including the coefficient
regularity in Lemma \ref{lem:v12v63-selfadjoint},
\begin{equation}
 \boxed{
 \|\Gamma_\beta(L_1)-\Gamma_\beta(L_0)\|_1
 \leq
 \min\left\{
 2,\
 \left(2\beta
 \sqrt{Q(v,c_1){\cal E}_\beta}\right)^{2/3}
 \right\}.}                                      \tag{63.13}
\end{equation}
Every positive trace-preserving retained-fibre map preserves the
same upper bound.
\end{theorem}

\begin{proof}
Lemma \ref{lem:v12v63-selfadjoint}, Corollary
\ref{cor:v12v62-heat}, and (63.11) verify the hypotheses of
Corollary \ref{cor:v12v59-moment-bridge}.  Apply (58.11) with
\(p=2\), \(C_2=2\), and the upper bound (63.11).  The cap by \(2\)
is the diameter of the state space in trace norm.  Contractivity
under a retained map is (53.10).
\end{proof}

The exponent \(2/3\) arises from the general clipping argument.  A
stronger model-specific covariance or relative-bound estimate could
improve it, but no such improvement is assumed here.

\subsection{Gauge and Yukawa interpretation}

At fixed metric and spin structure, write schematically
\[
 D_{A,\Phi}
 =
 D_{\rm spin}\otimes I
 +
 c(A)
 +
 Y(\Phi),
\]
where \(c(A)\) is Clifford contraction of the gauge connection and
\(Y(\Phi)\) is the finite internal Yukawa--Higgs endomorphism.  For
two backgrounds,
\begin{equation}
 V
 =
 c(A_1-A_0)+Y(\Phi_1)-Y(\Phi_0).                  \tag{63.14}
\end{equation}
On a compact manifold,
\begin{align}
 v
 &\leq
 C_{\rm cl}\|A_1-A_0\|_{L^\infty}
 +
 \|Y(\Phi_1)-Y(\Phi_0)\|_{L^\infty},              \tag{63.15}\\
 c_1
 &\leq
 C_{\rm der}
 \left\{
 \|A_1-A_0\|_{W^{1,\infty}}
 +
 \|Y(\Phi_1)-Y(\Phi_0)\|_{W^{1,\infty}}
 \right\},                                        \tag{63.16}
\end{align}
with constants fixed by the Clifford representation, reference
connection, and finite internal matrices.  Once those constants and
the full legal packet are specified, (63.13) is an explicit
background-stability estimate.

The quotient and thermal centre must be common to the complete
direct-sum packet, as proved in V56.  Species-dependent scalar
subtractions cannot be made before the packet Gibbs weights are
assembled.

\subsection{What changes when the metric moves}

If the metric changes, the transported Dirac operators generally
have different principal symbols.  Their difference is then
first-order rather than the bounded \(V\) in (63.1), and the square
difference can remain second-order.  Equations (63.5)--(63.16) do
not cover that case.  One must instead verify the V62
relative-power inequalities directly for the transported
metric-dependent pair, including the connection commutator
(61.17).

\begin{assessment}[Status after V63]
\label{ass:v12v63-status}
For a fixed metric, bounded gauge-connection and Yukawa changes give
a first-order self-adjoint difference of squared Dirac
Hamiltonians.  Its two-state thermal second moment is explicitly
controlled by the zeroth-order perturbation, its Dirac commutator,
and the two Gibbs energy expectations.  This closes the V58--V59
handoff for that sector and yields (63.13).
\end{assessment}

\begin{firewall}[Claim boundary after V63]
\label{fire:v12v63-claim}
The schematic RPESM operator, Clifford constants, Yukawa matrices,
legal chiral packet, and energy stocks have not been populated.
Metric variations and boundaries are not covered by (63.13).
The result does not establish determinant/Pfaffian phase gluing,
Ward compatibility, reflection positivity, gravitational
tightness, or the full SM--Einstein continuum.
\end{firewall}

\subsection{Parent record}

This V63 note appends to the validated V62 file with record
\[
\begin{array}{c|c}
\text{lines}&17729\\
\text{bytes}&587513\\
\text{SHA--256}&
\texttt{A2C31D0DB3C89AEF991F47A486A0119E310DBC126CDE658A6B66E370C1DFF8F1}.
\end{array}
\]

% =============================================================================
% END APPENDED TWELFTH-CYCLE V63 MATERIAL
% =============================================================================
% =============================================================================
% BEGIN APPENDED TWELFTH-CYCLE V64 MATERIAL
% =============================================================================

\section{Twelfth-cycle V64: quadratic-form Gibbs stability for
metric principal variations}
\label{sec:v12v64}

V63 handles a fixed principal symbol by representing the squared
Dirac difference as a self-adjoint first-order operator.  A metric
variation can change the second-order principal part, and the
difference of the two self-adjoint Hamiltonians need not itself have
a useful self-adjoint closure.  The correct upstream object is then
the difference of their closed quadratic forms.  This note derives a
trace-distance estimate directly from relative entropy on a common
form domain.

\subsection{Two-sided relative form oscillation}

Let \(H,K\) be self-adjoint and bounded below on one Hilbert space,
with a common closed form domain
\[
                         {\cal Q}:={\cal Q}(H)={\cal Q}(K).
\]
Let their closed forms be \(h\) and \(k\), and choose lower bounds
\[
                         H\geq\ell_HI,\qquad K\geq\ell_KI.
\]
Assume that for some \(0<\beta_0<\beta\),
\begin{equation}
 \operatorname{Tr}e^{-\beta_0H}
 +
 \operatorname{Tr}e^{-\beta_0K}<\infty.           \tag{64.1}
\end{equation}
Put
\[
 \rho:=\Gamma_\beta(H),\qquad
 \sigma:=\Gamma_\beta(K),
\]
and define the positive thermal energy stocks
\begin{align}
 E_H&:=
 \operatorname{Tr}\{\rho(I+H-\ell_HI)\},           \tag{64.2}\\
 E_K&:=
 \operatorname{Tr}\{\sigma(I+K-\ell_KI)\}.         \tag{64.3}
\end{align}
They are finite by Lemma \ref{lem:v12v59-heat-margin}.

For \(c\in\mathbb R\), define the scalar-centered form difference
\begin{equation}
 r_c[u]
 :=
 k[u]-h[u]-c\|u\|^2,\qquad u\in{\cal Q}.           \tag{64.4}
\end{equation}
Suppose there are finite \(\epsilon_H,\epsilon_K\geq0\) such that
\begin{align}
 |r_c[u]|
 &\leq
 \epsilon_H\{
 \|u\|^2+h[u]-\ell_H\|u\|^2\},                    \tag{64.5}\\
 |r_c[u]|
 &\leq
 \epsilon_K\{
 \|u\|^2+k[u]-\ell_K\|u\|^2\}.                    \tag{64.6}
\end{align}

\subsection{The form entropy identity}

If \(\omega\) is a Gibbs state whose Hamiltonian eigenvectors lie in
\({\cal Q}\), define the expectation of \(r_c\) by summing the form
over an eigenbasis of the density.  Bounds (64.5)--(64.6) make the
relevant sums absolutely convergent.

\begin{lemma}[Cross-energy and entropy identity]
\label{lem:v12v64-form-entropy}
Under (64.1)--(64.6), both relative entropies are finite and
\begin{align}
 D(\rho\|\sigma)
 &=
 \beta\{\langle r_c\rangle_\rho+c\}
 +\log{Z_K\over Z_H},                             \tag{64.7}\\
 D(\sigma\|\rho)
 &=
 -\beta\{\langle r_c\rangle_\sigma+c\}
 +\log{Z_H\over Z_K}.                             \tag{64.8}
\end{align}
Consequently,
\begin{equation}
 D(\rho\|\sigma)+D(\sigma\|\rho)
 =
 \beta\{
 \langle r_c\rangle_\rho-\langle r_c\rangle_\sigma\}.         \tag{64.9}
\end{equation}
Moreover,
\begin{equation}
 |\langle r_c\rangle_\rho|\leq\epsilon_HE_H,
 \qquad
 |\langle r_c\rangle_\sigma|\leq\epsilon_KE_K.    \tag{64.10}
\end{equation}
\end{lemma}

\begin{proof}
The heat margin (64.1) and Lemma
\ref{lem:v12v59-heat-margin} make the own energies
\(\operatorname{Tr}(\rho|H|)\) and
\(\operatorname{Tr}(\sigma|K|)\) finite.  Apply (64.5) to every
\(H\)-eigenvector and sum with its Gibbs weight.  This proves the
first inequality in (64.10) and absolute convergence of
\[
 \operatorname{Tr}(\rho K)
 =
 \operatorname{Tr}(\rho H)+c+\langle r_c\rangle_\rho.         \tag{64.11}
\]
Likewise, (64.6) gives the second inequality in (64.10) and
\[
 \operatorname{Tr}(\sigma H)
 =
 \operatorname{Tr}(\sigma K)-c-\langle r_c\rangle_\sigma.     \tag{64.12}
\]
Here each cross energy is understood through the common closed form;
the absolute estimates show that no difference of infinities occurs.

Since
\[
 \log\rho=-\beta H-(\log Z_H)I,\qquad
 \log\sigma=-\beta K-(\log Z_K)I,
\]
substitution of (64.11) in the definition of
\(D(\rho\|\sigma)\) proves (64.7), and (64.12) proves (64.8).
Adding cancels both the partition functions and the common scalar
\(c\), giving (64.9).
\end{proof}

The common scalar cancels, but a sector-dependent scalar does not;
the packet rule of V56 remains in force.

\subsection{Form-stability theorem}

\begin{theorem}[Gibbs stability from a common form domain]
\label{thm:v12v64-form-stability}
Under (64.1)--(64.6),
\begin{equation}
 \boxed{
 \|\Gamma_\beta(K)-\Gamma_\beta(H)\|_1
 \leq
 \min\left\{
 2,\
 \sqrt{\beta(\epsilon_HE_H+\epsilon_KE_K)}
 \right\}.}                                      \tag{64.13}
\end{equation}
The same estimate holds after every positive trace-preserving
retained-fibre map.
\end{theorem}

\begin{proof}
Put \(\delta=\|\rho-\sigma\|_1\).  Applying quantum Pinsker in both
directions as in (57.7), then using (64.9)--(64.10), gives
\[
 \delta^2
 \leq
 D(\rho\|\sigma)+D(\sigma\|\rho)
 \leq
 \beta(\epsilon_HE_H+\epsilon_KE_K).
\]
Take square roots and use \(\delta\leq2\).  Retained-map
contractivity is (53.10).
\end{proof}

No self-adjoint operator representing \(K-H\) appears in the
theorem.  The price is a square-root dependence on the relative form
oscillation.  V57 remains stronger for bounded differences, and V58
can be stronger when a self-adjoint unbounded difference has high
thermal moments.

\subsection{Coefficient realization}

For the uniformly elliptic family \(L_x\) of V62, let
\(\mathfrak l_x\) be its closed form on \(H^1(M,E)\).  The uniform
G{\aa}rding inequality gives constants \(C_{\rm form}\) such that
\begin{equation}
 \|u\|_{H^1}^2
 \leq
 C_{\rm form}
 \{\|u\|^2+\mathfrak l_x[u]-\ell_x\|u\|^2\}
                                                        \tag{64.14}
\end{equation}
for every \(x\) in the coefficient class.

For \(x,y\), choose one real scalar \(c\) and define
\begin{align}
 \Xi_c^{\rm form}(x,y)
 :={}&
 \|a_y-a_x\|_{L^\infty}
 +\|b_y-b_x\|_{L^\infty}                         \nonumber\\
 &+
 \|c_y-c_x-cI_E\|_{L^\infty}.                    \tag{64.15}
\end{align}
The first-order term is interpreted in the symmetric integrated
form.

\begin{proposition}[Coefficient-to-form handoff]
\label{prop:v12v64-coefficient}
There is a constant \(C_{\rm coef}\), depending only on the fixed
atlas, bundle connection, rank, and uniform coefficient class, such
that
\begin{equation}
 |\mathfrak l_y[u]-\mathfrak l_x[u]-c\|u\|^2|
 \leq
 C_{\rm coef}\Xi_c^{\rm form}(x,y)\|u\|_{H^1}^2.  \tag{64.16}
\end{equation}
Consequently (64.5)--(64.6) hold with
\begin{equation}
 \epsilon_H,\epsilon_K
 \leq
 C_{\rm coef}C_{\rm form}
 \Xi_c^{\rm form}(x,y).                           \tag{64.17}
\end{equation}
\end{proposition}

\begin{proof}
Subtract the two integrated forms.  The principal term is bounded
by
\(\|a_y-a_x\|_\infty\|\nabla u\|_2^2\).
Every symmetric first-order term is bounded by its coefficient
difference times \(\|\nabla u\|_2\|u\|_2\), and
\(2ab\leq a^2+b^2\) absorbs it into the \(H^1\) norm.  The centered
zeroth-order term is bounded by its \(L^\infty\) norm times
\(\|u\|_2^2\).  A finite partition of unity produces the uniform
constant in (64.16).  Apply (64.14) at \(x\) and at \(y\) to obtain
(64.17).
\end{proof}

If
\[
                         E_*:=\sup_{x\in G}E_{L_x}<\infty,
\]
Theorem \ref{thm:v12v64-form-stability} and (64.17) give
\begin{equation}
 \|\Gamma_\beta(L_y)-\Gamma_\beta(L_x)\|_1
 \leq
 \sqrt{
 2\beta C_{\rm coef}C_{\rm form}E_*
 \inf_{c\in\mathbb R}\Xi_c^{\rm form}(x,y)}.       \tag{64.18}
\end{equation}
To see the uniform energy stock, the lower spectral comparison
(62.8) bounds the polynomially weighted heat-trace numerator
uniformly.  The coefficient upper bounds give an upper form
comparison with \(\Lambda\); evaluation on one fixed normalized
smooth section yields \(\sup_x\lambda_1(L_x)<\infty\), and hence a
uniform partition lower bound
\(Z_x\geq e^{-\beta\sup_x\lambda_1(L_x)}>0\).  Their ratio gives
\(E_*<\infty\) on every compact inverse-temperature interval bounded
away from zero.

\subsection{Hölder cell martingales for moving metrics}

Suppose a parameterized transported elliptic family satisfies
\begin{equation}
 \inf_c\Xi_c^{\rm form}(x,y)
 \leq L_{\rm form}d_{\mathbb T}(x,y).             \tag{64.19}
\end{equation}
Define
\[
 C_{\beta,\rm form}
 :=
 \sqrt{2\beta C_{\rm coef}C_{\rm form}E_*L_{\rm form}}.
\]
Then (64.18) becomes the \(1/2\)-Hölder estimate
\begin{equation}
 \|\rho(y)-\rho(x)\|_1
 \leq C_{\beta,\rm form}d_{\mathbb T}(x,y)^{1/2}. \tag{64.20}
\end{equation}
Repeating the Bochner cell-average argument of V54 gives
\begin{align}
 \|\rho-E_n\rho\|_{L^1}
 &\leq
 C_{\beta,\rm form}\delta_n^{1/2}
 =
 C_{\beta,\rm form}D^{1/4}2^{-n/2},               \tag{64.21}\\
 \|E_n(S_{n+1}\rho)-S_n\rho\|_{L^1}
 &\leq
 C_{\beta,\rm form}
 \{\delta_n^{1/2}+\delta_{n+1}^{1/2}\}.           \tag{64.22}
\end{align}
The adjacent errors remain summable because
\(\sum_n2^{-n/2}<\infty\).  Thus principal metric variation weakens
the rate but does not destroy the cofinal cell construction.

\subsection{Application after the V61 metric unitary}

For transported metric Hamiltonians
\[
 L_g=U_gH_gU_g^*
\]
the theorem applies once:
\begin{enumerate}
\item their closed forms have the common domain \(H^1(M,E_0)\);
\item the transported coefficients satisfy the uniform V62 class;
\item the full density and connection terms from \(U_g\) are
included in (64.15);
\item the legal packet uses one common scalar \(c\); and
\item the coefficient variation and energy stock are uniform along
the cofinal metric set.
\end{enumerate}
No self-adjointness of \(L_h-L_g\) is needed.

\begin{assessment}[Status after V64]
\label{ass:v12v64-status}
A common transported form domain now suffices to compare Gibbs
states under a change of elliptic principal symbol.  The
trace-distance cost is the square root of a two-sided relative form
oscillation times finite Gibbs energy stocks.  Uniformly Lipschitz
metric coefficients yield a \(1/2\)-Hölder state field and a
geometrically summable dyadic cell error.
\end{assessment}

\begin{firewall}[Claim boundary after V64]
\label{fire:v12v64-claim}
The physical metric-dependent RPESM coefficients, their common
spinor identification, uniform form class, and energy stock have not
been populated.  The theorem assumes a closed manifold and a
positive Gibbs Hamiltonian.  It does not prove compatibility with
the chiral determinant/Pfaffian phase, Ward repair, reflection
positivity, gravitational tightness, or the full SM--Einstein
continuum.
\end{firewall}

\subsection{Parent record}

This V64 note appends to the validated V63 file with record
\[
\begin{array}{c|c}
\text{lines}&18061\\
\text{bytes}&598034\\
\text{SHA--256}&
\texttt{2712D6448333F54293BAE2F2C57A6A371A5A881953F7C021A9C02144B2091157}.
\end{array}
\]

% =============================================================================
% END APPENDED TWELFTH-CYCLE V64 MATERIAL
% =============================================================================
% =============================================================================
% BEGIN APPENDED TWELFTH-CYCLE V65 MATERIAL
% =============================================================================

\section{Twelfth-cycle V65: mandatory companion audit}
\label{sec:v12v65}

This is the required five-version review after V60 and before the
background-mixture step.

\subsection{Files inspected}

The newest active companion files remain
\[
\begin{array}{c|r|r}
\text{file}&\text{lines}&\text{bytes}\\ \hline
\texttt{Renewal\_Yang\_Mills\_Mass\_Gap\_Research\_Notes\_V56.tex}
 &21223&723854\\
\texttt{Renewal\_NS\_Stability\_Research\_Notes\_V26.tex}
 &5319&278283.
\end{array}                                      \tag{65.1}
\]
Their SHA--256 digests are
\[
\begin{split}
 &\texttt{627247BCFC1A511C5158B1D3DED869FCA598501597EF3B8CF174C3A0BFEF54DC},\\
 &\texttt{82C887F8C2C69E4F28FAD7C3DC8DE5B9099CB5A4BF431EBA1FFF3E91935978AD}.
\end{split}                                      \tag{65.2}
\]
Thus neither companion changed after the V60 audit.

\subsection{Effect on V61--V64}

YM V56 still requires exponent-resolved data before a finite
momentum sample can be interpreted as a Laurent or Wilson symbol.
The operator and form estimates V61--V64 compare genuine
Hamiltonians and their Gibbs states; they do not infer those
Hamiltonians from parity samples.  No change to their statements is
needed.

NS V26 still permits a reduced norm only after its structured input
class is proved.  V63 uses a fixed-principal-symbol identity
\(D_1-D_0=V\), and V64 uses a common \(H^1\) form domain.  These are
explicit structural hypotheses, so the transferred NS warning is
satisfied.  Neither hypothesis has yet been verified for the
physical RPESM packet.

\subsection{Audit decision}

\begin{decision}[Direction after the V65 audit]
\label{dec:v12v65-direction}
Retain V61--V64 unchanged.  The next missing bridge is probabilistic:
combine a convergent or tight sequence of Einstein-background
measures with the transported conditional SM Gibbs states.  Prove a
coupling estimate for their trace-class barycentres and a
bounded-Lipschitz estimate for the resulting classical--quantum
joint states.  Keep the gravitational convergence hypothesis
explicit.
\end{decision}

\begin{assessment}[Status after V65]
\label{ass:v12v65-status}
The mandatory audit is complete.  It supplies no new companion
constant and changes no prior SM theorem.  It confirms that the next
step must preserve the distinction between a certified conditional
fibre and the still-missing probability law of the Einstein
background.
\end{assessment}

\begin{firewall}[Claim boundary after V65]
\label{fire:v12v65-claim}
No gravitational measure, tightness estimate, physical Hamiltonian
field, or companion result has been imported.  The full
SM--Einstein continuum remains open.
\end{firewall}

\subsection{Parent record}

This V65 note appends to the validated V64 file with record
\[
\begin{array}{c|c}
\text{lines}&18403\\
\text{bytes}&609035\\
\text{SHA--256}&
\texttt{A018D6E2CF64D4FC356052C10039DEB1725ED476E3E80E767537C91BA653E54E}.
\end{array}
\]

% =============================================================================
% END APPENDED TWELFTH-CYCLE V65 MATERIAL
% =============================================================================
% =============================================================================
% BEGIN APPENDED TWELFTH-CYCLE V66 MATERIAL
% =============================================================================

\section{Twelfth-cycle V66: coupling conditional SM fibres to
Einstein-background measures}
\label{sec:v12v66}

V61--V64 construct and compare transported conditional Gibbs states
over a fixed background.  A joint SM--Einstein object must also
integrate those states against the probability law of the metric and
other classical backgrounds.  This note proves the required
classical--quantum coupling theorem.  It separates background
transport from conditional-state error and gives a summable
criterion that constructs the joint limit.

\subsection{Bounded transport metric}

Let \((X,d)\) be a complete separable metric space of background
configurations and put
\begin{equation}
 \bar d(x,y):=\min\{1,d(x,y)\}.                   \tag{66.1}
\end{equation}
For \(0<\alpha\leq1\), \(\bar d^\alpha\) is again a complete bounded
metric with the same topology.  For probability measures
\(\mu,\nu\) on \(X\), define
\begin{equation}
 {\mathsf W}_\alpha(\mu,\nu)
 :=
 \inf_{\pi\in\Pi(\mu,\nu)}
 \int_{X\times X}\bar d(x,y)^\alpha\,d\pi(x,y).   \tag{66.2}
\end{equation}
Because the cost is bounded and generates the topology,
\({\mathsf W}_\alpha\) metrizes weak convergence of probability
measures on \(X\).

Let \({\cal H}\) be the reference Hilbert space obtained after the
unitary transports of V61.  A conditional state field is a strongly
measurable map
\[
 \rho:X\longrightarrow{\mathfrak S}_1({\cal H}),
 \qquad
 \rho(x)\succeq0,\quad\operatorname{Tr}\rho(x)=1.
\]
Assume it is Hölder continuous:
\begin{equation}
 \|\rho(x)-\rho(y)\|_1
 \leq C_\rho\bar d(x,y)^\alpha.                  \tag{66.3}
\end{equation}
V57 supplies \(\alpha=1\) for bounded scalar-quotient
Hamiltonian differences; V64 supplies \(\alpha=1/2\) for principal
metric variations controlled at form level.

\subsection{Finite background spaces}

Let \(X_n\) be a finite or standard Borel background space,
\(\mu_n\) a probability measure on it, and
\[
 \iota_n:X_n\longrightarrow X
\]
a measurable reconstruction map.  Let
\[
 \rho_n:X_n\longrightarrow{\mathfrak S}_1({\cal H})
\]
be a measurable density-operator field.  Define the uniform
conditional error
\begin{equation}
 \epsilon_n
 :=
 \operatorname*{ess\,sup}_{s\in X_n}
 \|\rho_n(s)-\rho(\iota_n(s))\|_1.                \tag{66.4}
\end{equation}
The reconstructed background law is
\[
                         \nu_n:=(\iota_n)_*\mu_n.
\]

The trace-class barycentres are
\begin{equation}
 \overline\rho_n:=\int_{X_n}\rho_n(s)\,d\mu_n(s),
 \qquad
 \overline\rho:=\int_X\rho(x)\,d\mu(x),            \tag{66.5}
\end{equation}
where the integrals are Bochner integrals.

\begin{theorem}[Background--fibre barycentre bound]
\label{thm:v12v66-barycentre}
For every probability measure \(\mu\) on \(X\),
\begin{equation}
 \boxed{
 \|\overline\rho_n-\overline\rho\|_1
 \leq
 \epsilon_n+C_\rho{\mathsf W}_\alpha(\nu_n,\mu).} \tag{66.6}
\end{equation}
Both barycentres are positive and have trace one.
\end{theorem}

\begin{proof}
Choose a coupling \(\pi\in\Pi(\nu_n,\mu)\).  Equivalently, after
disintegrating over \(\iota_n\), use a coupling of \(s\sim\mu_n\)
and \(x\sim\mu\) whose reconstructed first coordinate has law
\(\nu_n\).  Bochner's inequality, (66.4), and (66.3) give
\begin{align}
 \|\overline\rho_n-\overline\rho\|_1
 &\leq
 \int\|\rho_n(s)-\rho(x)\|_1\,d\pi(s,x)           \nonumber\\
 &\leq
 \epsilon_n+
 C_\rho\int\bar d(\iota_n(s),x)^\alpha\,d\pi(s,x).
                                                        \tag{66.7}
\end{align}
Take the infimum over couplings.  Positivity follows because the
positive trace-class cone is closed and convex; trace one follows
from continuity of the trace under the Bochner integral.
\end{proof}

The barycentre records the unconditional quantum state but forgets
correlation with the classical background.  The next construction
retains that correlation.

\subsection{The joint classical--quantum state}

Let \({\cal A}\) be the unital \(C^*\)-algebra of bounded
norm-continuous fields
\[
 F:X\longrightarrow{\cal B}({\cal H})
\]
with pointwise operations and the supremum norm.  Define
\begin{align}
 \Omega_n(F)
 &:=
 \int_{X_n}
 \operatorname{Tr}\{\rho_n(s)F(\iota_n(s))\}
 \,d\mu_n(s),                                     \tag{66.8}\\
 \Omega(F)
 &:=
 \int_X\operatorname{Tr}\{\rho(x)F(x)\}\,d\mu(x). \tag{66.9}
\end{align}
These are positive normalized functionals wherever the displayed
measurability holds.

For a field \(F\), put
\begin{equation}
 [F]_\alpha
 :=
 \sup_{x\ne y}
 {\|F(x)-F(y)\|\over\bar d(x,y)^\alpha}.          \tag{66.10}
\end{equation}
Define the bounded-Hölder distance between joint states by
\begin{equation}
 d_{\rm BH,\alpha}(\Omega_1,\Omega_2)
 :=
 \sup_{\substack{\|F\|_\infty\leq1\\ {}[F]_\alpha\leq1}}
 |\Omega_1(F)-\Omega_2(F)|.                       \tag{66.11}
\end{equation}

\begin{theorem}[Joint-state coupling bound]
\label{thm:v12v66-joint}
Under (66.3)--(66.4),
\begin{equation}
 \boxed{
 d_{\rm BH,\alpha}(\Omega_n,\Omega)
 \leq
 \epsilon_n+(C_\rho+1)
 {\mathsf W}_\alpha(\nu_n,\mu).}                  \tag{66.12}
\end{equation}
More generally, for any \(F\) with finite Hölder seminorm,
\begin{equation}
 |\Omega_n(F)-\Omega(F)|
 \leq
 \|F\|_\infty\epsilon_n+
 \{C_\rho\|F\|_\infty+[F]_\alpha\}
 {\mathsf W}_\alpha(\nu_n,\mu).                   \tag{66.13}
\end{equation}
\end{theorem}

\begin{proof}
Use the same coupling as in (66.7).  Add and subtract
\(\operatorname{Tr}\{\rho(x)F(\iota_n(s))\}\).  Trace duality gives
\begin{align}
 &\left|
 \operatorname{Tr}\{\rho_n(s)F(\iota_n(s))\}
 -
 \operatorname{Tr}\{\rho(x)F(x)\}
 \right|                                          \nonumber\\
 &\quad\leq
 \|F\|_\infty\|\rho_n(s)-\rho(x)\|_1
 +
 \|F(\iota_n(s))-F(x)\|                           \nonumber\\
 &\quad\leq
 \|F\|_\infty\epsilon_n+
 \{C_\rho\|F\|_\infty+[F]_\alpha\}
 \bar d(\iota_n(s),x)^\alpha.                     \tag{66.14}
\end{align}
Integrate and minimize over the coupling to prove (66.13).
Taking the supremum in (66.11) gives (66.12).
\end{proof}

Scalar background observables correspond to
\(F(x)=f(x)I\); constant quantum observables correspond to
\(F(x)=A\).  Hence the joint estimate controls both sectors and
their correlations.

\subsection{A target-free Cauchy construction}

The preceding theorems use a stated limit \((\mu,\rho)\).  A
cofinal construction should obtain it from adjacent data.  Suppose
now that all reconstructed laws live on \(X\), and write them simply
as \(\mu_n\).  Let
\[
 \rho_n:X\to{\mathfrak S}_1({\cal H})
\]
be state fields satisfying one common Hölder bound
\begin{equation}
 \|\rho_n(x)-\rho_n(y)\|_1
 \leq C_\rho\bar d(x,y)^\alpha                    \tag{66.15}
\end{equation}
and adjacent uniform errors
\begin{equation}
 a_n:=
 \sup_x\|\rho_{n+1}(x)-\rho_n(x)\|_1.             \tag{66.16}
\end{equation}
Set
\[
                         w_n:={\mathsf W}_\alpha(\mu_{n+1},\mu_n).
\]

\begin{theorem}[Cofinal classical--quantum limit]
\label{thm:v12v66-cofinal}
If
\begin{equation}
 \sum_{n=0}^\infty(a_n+w_n)<\infty,               \tag{66.17}
\end{equation}
then:
\begin{enumerate}
\item \(\mu_n\) converges in \({\mathsf W}_\alpha\) to a probability
measure \(\mu_\infty\);
\item \(\rho_n\) converges uniformly in trace norm to a
\(C_\rho\)-Hölder state field \(\rho_\infty\);
\item the barycentres
\[
 \overline\rho_n=\int_X\rho_n\,d\mu_n
\]
converge in trace norm to
\(\overline\rho_\infty=\int_X\rho_\infty\,d\mu_\infty\); and
\item the joint states \(\Omega_n\) converge in
\(d_{\rm BH,\alpha}\) to the state defined by
\((\mu_\infty,\rho_\infty)\).
\end{enumerate}
Quantitatively,
\begin{align}
 \|\overline\rho_m-\overline\rho_\infty\|_1
 &\leq
 \sum_{n=m}^\infty(a_n+C_\rho w_n),               \tag{66.18}\\
 d_{\rm BH,\alpha}(\Omega_m,\Omega_\infty)
 &\leq
 \sum_{n=m}^\infty
 \{a_n+(C_\rho+1)w_n\}.                           \tag{66.19}
\end{align}
\end{theorem}

\begin{proof}
The space of probability measures on the complete separable space
\(X\), equipped with the Wasserstein distance for the complete
bounded metric \(\bar d^\alpha\), is complete.  The summability of
\(w_n\) therefore gives \(\mu_\infty\).  The summability of \(a_n\)
makes \(\rho_n\) uniformly Cauchy in the complete trace-class Banach
space.  Positivity and trace one pass to the uniform limit, and
(66.15) passes to the limit as well.

For adjacent barycentres, choose an almost optimal coupling of
\(\mu_n,\mu_{n+1}\).  Then
\[
 \|\rho_{n+1}(y)-\rho_n(x)\|_1
 \leq a_n+C_\rho\bar d(x,y)^\alpha.
\]
The proof of Theorem \ref{thm:v12v66-barycentre} gives
\[
 \|\overline\rho_{n+1}-\overline\rho_n\|_1
 \leq a_n+C_\rho w_n.
\]
Sum from \(m\) and pass to the limit to obtain (66.18).
The adjacent version of (66.12) gives
\[
 d_{\rm BH,\alpha}(\Omega_{n+1},\Omega_n)
 \leq a_n+(C_\rho+1)w_n,
\]
and summation proves (66.19).  Pointwise limits of positive
normalized functionals are positive and normalized, so the limit is
the asserted joint state.
\end{proof}

\subsection{Error-stock assignment}

For an RPESM cofinal sequence, the adjacent conditional stock
\(a_n\) may receive contributions from:
\begin{enumerate}
\item transported Hamiltonian errors from V57, V58, V63, or V64;
\item unitary cocycle errors from V61;
\item signed Fourier or cell approximation errors from V47--V54;
\item Ward repair and retained-map errors from V19--V34; and
\item any legal packet or visibility remainder.
\end{enumerate}
The background stock \(w_n\) is different: it must come from the
Einstein/gravitational measure programme.  Equation (66.17) requires
both stocks to be summable and does not allow one to substitute for
the other.

\begin{assessment}[Status after V66]
\label{ass:v12v66-status}
A convergent Einstein-background law and a uniformly controlled
transported SM conditional state now determine a joint
classical--quantum continuum state.  More strongly, summable
adjacent background Wasserstein errors and conditional trace-norm
errors construct both limits directly, with the quantitative tails
(66.18)--(66.19).
\end{assessment}

\begin{firewall}[Claim boundary after V66]
\label{fire:v12v66-claim}
No Einstein-background probability measures or Wasserstein stocks
\(w_n\) have been constructed, and no complete physical conditional
RPESM state field has been populated.  The joint state is a
classical--quantum mixture; additional work is required to identify
the local observable algebra, prove reflection positivity and Ward
identities on it, establish determinant-line gluing, and obtain the
full SM--Einstein continuum.
\end{firewall}

\subsection{Parent record}

This V66 note appends to the validated V65 file with record
\[
\begin{array}{c|c}
\text{lines}&18494\\
\text{bytes}&612364\\
\text{SHA--256}&
\texttt{80D4EE2AA5BFE24866867E36E6FC1F49873B467B1781E1A74E4626CC600071E2}.
\end{array}
\]

% =============================================================================
% END APPENDED TWELFTH-CYCLE V66 MATERIAL
% =============================================================================
% =============================================================================
% BEGIN APPENDED TWELFTH-CYCLE V67 MATERIAL
% =============================================================================

\section{Twelfth-cycle V67: Ward and reflection-positive closure of
the joint limit}
\label{sec:v12v67}

V66 constructs a joint classical--quantum state from summable
background and conditional-fibre errors.  Convergence of states
alone does not identify the limit as an Euclidean quantum field
theory.  This note proves the closedness statements needed to pass
exact or asymptotically exact Ward identities and
Osterwalder--Schrader reflection positivity to that joint limit.

\subsection{Reflected state cone}

Let \({\cal A}\) be a unital \(C^*\)-algebra, let
\({\cal A}_+\subset{\cal A}\) be a unital subalgebra, and let
\[
 \Theta:{\cal A}\longrightarrow{\cal A}
\]
be a norm-continuous conjugate-linear involution preserving products
and the star structure in the chosen Euclidean reflection
convention.  A state \(\omega\) is reflection positive on
\(({\cal A},\Theta,{\cal A}_+)\) when
\begin{equation}
 \omega(\Theta(F)F)\geq0,\qquad F\in{\cal A}_+.    \tag{67.1}
\end{equation}
The inequality includes the assertion that the displayed number is
real.

\begin{theorem}[Weak-star closedness of reflection positivity]
\label{thm:v12v67-OS-closed}
The set of states satisfying (67.1) is weak-star closed and convex.
Consequently, if \(\omega_n\) are reflection positive and
\[
                         \omega_n(A)\longrightarrow\omega(A)
 \qquad(A\in{\cal A}),                            \tag{67.2}
\]
then \(\omega\) is reflection positive.
\end{theorem}

\begin{proof}
For fixed \(F\in{\cal A}_+\), evaluation at the fixed algebra element
\(\Theta(F)F\) is weak-star continuous and affine.  The inverse image
of the closed half-line \([0,\infty)\) is therefore weak-star closed
and convex.  Intersect these sets over all \(F\in{\cal A}_+\).
Equation (67.2) then passes every inequality to the limit.
\end{proof}

It is enough to verify (67.1) on a norm-dense subalgebra
\({\cal C}_+\subset{\cal A}_+\).  Indeed
\(F\mapsto\Theta(F)F\) is norm continuous and every state has norm
one.

\subsection{A quantitative varying-reflection criterion}

At finite level the reflected algebra and its representatives may
vary.  Let \(\Theta_n\) be conjugate-linear maps on a common
normed core, and let \(F_n\) approximate \(F\).

\begin{proposition}[Approximate OS passage]
\label{prop:v12v67-OS-approx}
Suppose, for each \(F\) in a dense positive-time core, there are
\(F_n\) such that
\begin{align}
 F_n&\longrightarrow F,                           \tag{67.3}\\
 \Theta_n(F_n)&\longrightarrow\Theta(F)            \tag{67.4}
\end{align}
in norm.  Assume
\begin{equation}
 \omega_n(\Theta_n(F_n)F_n)\geq-r_n(F),
 \qquad r_n(F)\longrightarrow0,                   \tag{67.5}
\end{equation}
and that
\begin{equation}
 \omega_n(\Theta(F)F)\longrightarrow
 \omega(\Theta(F)F).                              \tag{67.6}
\end{equation}
Then \(\omega\) is reflection positive.
\end{proposition}

\begin{proof}
State norm one and the product decomposition give
\begin{align}
 &\left|
 \omega_n(\Theta_n(F_n)F_n)
 -
 \omega_n(\Theta(F)F)
 \right|                                          \nonumber\\
 &\quad\leq
 \|\Theta_n(F_n)-\Theta(F)\|\,\|F_n\|
 +
 \|\Theta(F)\|\,\|F_n-F\|.                        \tag{67.7}
\end{align}
The right side tends to zero by (67.3)--(67.4).  Equations
(67.5)--(67.6) therefore give
\(\omega(\Theta(F)F)\geq0\) on the core.  Extend by density.
\end{proof}

The convergence of states must be tested on the OS products
\(\Theta(F)F\).  Convergence only on linear fields \(F\) is
insufficient.

\subsection{Ward closure}

Let \({\cal C}\subset{\cal A}\) be a common local test algebra and
let
\[
 {\mathsf s}:{\cal C}\to{\cal A}
\]
be the limiting Ward or BRST operator on that core.  At level \(n\),
let \({\mathsf s}_n\) be its finite approximation.

\begin{theorem}[Asymptotic Ward identity]
\label{thm:v12v67-Ward}
Suppose for every \(F\in{\cal C}\),
\begin{align}
 \|{\mathsf s}_nF-{\mathsf s}F\|
 &\longrightarrow0,                              \tag{67.8}\\
 |\omega_n({\mathsf s}_nF)|
 &\leq r_n^{\rm W}(F)\longrightarrow0,            \tag{67.9}\\
 \omega_n({\mathsf s}F)
 &\longrightarrow\omega({\mathsf s}F).            \tag{67.10}
\end{align}
Then
\begin{equation}
                         \omega({\mathsf s}F)=0,
 \qquad F\in{\cal C}.                             \tag{67.11}
\end{equation}
\end{theorem}

\begin{proof}
Since every state has norm one,
\begin{align}
 |\omega({\mathsf s}F)|
 &\leq
 |\omega({\mathsf s}F)-\omega_n({\mathsf s}F)|
 +\|{\mathsf s}F-{\mathsf s}_nF\|
 +|\omega_n({\mathsf s}_nF)|.
\end{align}
All three terms tend to zero by (67.8)--(67.10).
\end{proof}

If \({\mathsf s}^2=0\) on \({\cal C}\), the limit identity is a Ward
identity for the same complex.  The theorem does not infer
nilpotence from (67.11); nilpotence must be supplied by the finite
BRST certificate and its operator convergence, as separated in
V17--V19.

\subsection{Symmetry covariance}

Let \(\alpha\) be a norm-continuous automorphism of \({\cal A}\).
If \(\omega_n\circ\alpha=\omega_n\) and
\(\omega_n\to\omega\) pointwise, then
\begin{equation}
                         \omega\circ\alpha=\omega.             \tag{67.12}
\end{equation}

\begin{proof}
For every \(A\in{\cal A}\),
\[
 \omega(\alpha(A))
 =
 \lim_n\omega_n(\alpha(A))
 =
 \lim_n\omega_n(A)
 =
 \omega(A).
\]
\end{proof}

The same argument permits approximate covariance with a residual
that tends to zero uniformly on the chosen local core.

\subsection{Mixtures with a decomposable reflection}

Let \(x\mapsto\omega_x\) be a measurable family of states on one
algebra \({\cal B}\).  Suppose the reflection and positive-time
algebra are fixed and each \(\omega_x\) is reflection positive.
Then for every probability measure \(\mu\),
\begin{equation}
                         \omega_\mu(B)
 :=
 \int_X\omega_x(B)\,d\mu(x)                       \tag{67.13}
\end{equation}
is reflection positive.

More generally, for an operator-field algebra, suppose
\[
 (\Theta F)(x)=\Theta_x(F(x))
\]
acts fibrewise, \(F(x)\) belongs to the positive-time fibre algebra,
and every conditional state is reflection positive for
\(\Theta_x\).  Then the V66 joint state satisfies
\begin{equation}
 \Omega(\Theta(F)F)
 =
 \int_X
 \omega_x(\Theta_x(F(x))F(x))\,d\mu(x)\geq0.       \tag{67.14}
\end{equation}

\begin{proof}
The integrand in (67.13) or (67.14) is nonnegative almost
everywhere.  Integration preserves the inequality.
\end{proof}

Equation (67.14) applies only to a decomposable reflection.  Physical
time reflection generally also sends the background
\(x\) to a reflected background \(\vartheta x\).  In that case
fibrewise positivity is not sufficient; the finite joint state must
satisfy the coupled OS inequality, and Theorem
\ref{thm:v12v67-OS-closed} or Proposition
\ref{prop:v12v67-OS-approx} must be used.

\subsection{Insertion into the V66 convergence metric}

For a bounded operator field \(G\), define
\[
 \|G\|_{{\rm BH},\alpha}
 :=
 \|G\|_\infty+[G]_\alpha.
\]
The proof of V66 gives, after homogeneity,
\begin{equation}
 |\Omega_m(G)-\Omega_\infty(G)|
 \leq
 \|G\|_{{\rm BH},\alpha}
 \sum_{n=m}^\infty
 \{a_n+(C_\rho+1)w_n\},                           \tag{67.15}
\end{equation}
up to replacing the displayed norm by the equivalent maximum norm
used in (66.11).

Therefore (67.6) is automatic when
\begin{equation}
 \|\Theta(F)F\|_{{\rm BH},\alpha}<\infty.          \tag{67.16}
\end{equation}
Likewise, (67.10) is automatic when
\begin{equation}
 \|{\mathsf s}F\|_{{\rm BH},\alpha}<\infty.        \tag{67.17}
\end{equation}
The finite-level residuals \(r_n(F)\) and
\(r_n^{\rm W}(F)\), and the operator convergence in
(67.3)--(67.4) and (67.8), remain separate stocks.

\subsection{Cofinal structural theorem}

\begin{theorem}[Joint-limit structure]
\label{thm:v12v67-joint-structure}
Assume the cofinal hypotheses (66.15)--(66.17), so that
\(\Omega_n\to\Omega_\infty\).  On dense local cores, also assume:
\begin{enumerate}
\item the approximate OS conditions (67.3)--(67.5) and the
bounded-Hölder test condition (67.16);
\item the approximate Ward conditions (67.8)--(67.9) and the
bounded-Hölder test condition (67.17);
\item exact or asymptotically exact covariance for every named
compact symmetry.
\end{enumerate}
Then \(\Omega_\infty\) is a normalized positive joint state that is
reflection positive, satisfies the limiting Ward identities, and is
covariant under those symmetries.
\end{theorem}

\begin{proof}
V66 gives the normalized positive state.  Equation (67.15) supplies
the required convergence on OS products and Ward images.
Proposition \ref{prop:v12v67-OS-approx}, Theorem
\ref{thm:v12v67-Ward}, and (67.12) give the three structural
conclusions.
\end{proof}

\subsection{Acquisition ledger}

A physical application must now populate four logically independent
rows:
\begin{enumerate}
\item the V66 summable background and conditional convergence
stocks \(w_n,a_n\);
\item the finite joint OS residual and convergence of the reflection
maps on the positive-time core;
\item the finite Ward residual and convergence of the BRST/Ward
operators on the local core; and
\item the common local observable algebra on which all products and
limits are defined.
\end{enumerate}
Positivity of each conditional density matrix supplies none of rows
2--4 automatically.

\begin{assessment}[Status after V67]
\label{ass:v12v67-status}
Reflection positivity is a weak-star closed cone of joint states,
and Ward identities are closed under convergence once their
operators and residuals converge on a common local core.  The V66
bounded-Hölder error controls the required OS products and Ward
images when they remain in that test class.  Thus a joint
SM--Einstein limit retains these structures under explicit,
separate finite-level certificates.
\end{assessment}

\begin{firewall}[Claim boundary after V67]
\label{fire:v12v67-claim}
No physical positive-time algebra, coupled background reflection,
finite joint OS certificate, limiting Ward operator, or local
observable completion has been constructed.  Fibrewise positivity
is not asserted to imply the coupled gravitational OS inequality.
Determinant-line gluing, the gravitational measures, and the full
SM--Einstein continuum remain open.
\end{firewall}

\subsection{Parent record}

This V67 note appends to the validated V66 file with record
\[
\begin{array}{c|c}
\text{lines}&18838\\
\text{bytes}&623560\\
\text{SHA--256}&
\texttt{4C2143912FCC3C30EFE3F8F3C6EB5FB493654EB35339804D190DC4D791D10E1C}.
\end{array}
\]

% =============================================================================
% END APPENDED TWELFTH-CYCLE V67 MATERIAL
% =============================================================================
% =============================================================================
% BEGIN APPENDED TWELFTH-CYCLE V68 MATERIAL
% =============================================================================

\section{Twelfth-cycle V68: exhausted determinant-line gluing and
paired complex-state transport}
\label{sec:v12v68}

The attached sector-patched fermion-line programme V15 already proves
that summable transition, reflection, holonomy, current, and
Pfaffian-sign defects converge on one fixed finite line screen.
That theorem is not repeated here.  The present note addresses two
later gaps: passage from fixed screens to a locally finite atlas by
compact exhaustion, and coupling of the resulting line-valued
fermion section to the V66 joint background--fibre limit.  The
second step requires an independent dual section; a Hermitian norm
does not provide complex-linear scalarization.

\subsection{A locally finite physical atlas}

Let \(X\) be a paracompact, sigma-compact Hausdorff space with a
locally finite cover
\[
                         {\cal U}=(U_i)_{i\in\mathbb N}.
\]
Choose a compact exhaustion
\[
 K_1\subset\operatorname{int}K_2\subset\cdots,
 \qquad
 \bigcup_{r\geq1}K_r=X,                            \tag{68.1}
\]
such that each \(K_r\) meets only finitely many \(U_i\).  At cutoff
\(n\), suppose the physically pulled-back determinant line is
presented on this same atlas by continuous unitary transitions
\begin{equation}
 g_{n,ij}:U_i\cap U_j\longrightarrow{\rm U}(1),   \tag{68.2}
\end{equation}
with
\begin{equation}
 g_{n,ii}=1,\qquad
 g_{n,ij}g_{n,jk}=g_{n,ik}                        \tag{68.3}
\end{equation}
on every triple overlap.  No cutoff-dependent rephasing is made
after the physical pullback.

For a compact screen \(K_r\), define
\begin{equation}
 d_{n,r}^{\rm tr}
 :=
 \max_{\substack{i,j\\K_r\cap U_i\cap U_j\ne\varnothing}}
 \|g_{n+1,ij}-g_{n,ij}\|_
 {C^0(K_r\cap U_i\cap U_j)}.                      \tag{68.4}
\end{equation}

\begin{theorem}[Exhausted Hermitian line limit]
\label{thm:v12v68-line-limit}
If, for every \(r\),
\begin{equation}
                         \sum_{n=1}^\infty
 d_{n,r}^{\rm tr}<\infty,                         \tag{68.5}
\end{equation}
then:
\begin{enumerate}
\item every \(g_{n,ij}\) converges uniformly on compact subsets of
\(U_i\cap U_j\) to a continuous \(g_{ij}:U_i\cap U_j\to{\rm U}(1)\);
\item the limits obey (68.3) and define a Hermitian line bundle
\(\mathscr L_\infty\to X\);
\item the restriction of \(\mathscr L_\infty\) to each \(K_r\) is
the unique limit, up to line isomorphism, of the corresponding
finite-screen data.
\end{enumerate}
\end{theorem}

\begin{proof}
Fix \(r\).  Only finitely many overlaps meet \(K_r\).  The telescopic
estimate from (68.4)--(68.5) makes every transition sequence
uniformly Cauchy on its compact overlap.  The unit circle is complete
and closed in \(\mathbb C\), so the limit is continuous and
unitary.  The limits obtained on \(K_r\) and \(K_{r+1}\) agree by
uniqueness of uniform limits.  They therefore define compact-open
limits on all overlaps.

The normalization and cocycle equations in (68.3) are closed under
uniform convergence, so they hold for \(g_{ij}\).  The standard
quotient of
\(\coprod_i U_i\times\mathbb C\) by
\[
 (x,z)_j\sim(x,g_{ij}(x)z)_i
\]
is consequently a complex line bundle.  Since the transitions are
unitary, the local standard metrics glue to a Hermitian metric.
Finite-screen uniqueness is precisely the finite descent statement,
and compatibility through the exhaustion gives the global
isomorphism class.
\end{proof}

Condition (68.5) is local in \(r\); it does not demand a finite
global cover when the physical background space is noncompact.

\subsection{Reflected real structures and Pfaffian signs}

Let \(\vartheta:X\to X\) be an involution and suppose the cover is
involuted, \(U_i\mapsto U_{\bar i}\).  At cutoff \(n\), write a local
anti-linear lift as
\[
 z_i\longmapsto
 \eta_{n,i}(x)\overline{z_i}
 \quad\text{in the chart }U_{\bar i}\text{ at }\vartheta x,
 \qquad \eta_{n,i}:U_i\to{\rm U}(1).
\]
Compatibility means
\begin{align}
 \eta_{n,i}(x)\overline{g_{n,ij}(x)}
 &=
 g_{n,\bar i\bar j}(\vartheta x)\eta_{n,j}(x),     \tag{68.6}\\
 \eta_{n,\bar i}(\vartheta x)\overline{\eta_{n,i}(x)}
 &=1.                                             \tag{68.7}
\end{align}
Define \(d_{n,r}^{\rm ref}\) by the maximum \(C^0(K_r)\) difference
of the finitely many \(\eta_{n,i}\) meeting \(K_r\).

\begin{corollary}[Reflected and real/Pfaffian limit]
\label{cor:v12v68-real-limit}
If
\[
 \sum_n(d_{n,r}^{\rm tr}+d_{n,r}^{\rm ref})<\infty
 \qquad\text{for every }r,                        \tag{68.8}
\]
then the limiting line carries an anti-linear involutive lift
\(\jmath_\infty\) of \(\vartheta\).  On a reflection-fixed real
Pfaffian stratum, suppose the physically transported orientation
signs take values in \(\{+1,-1\}\) and their screenwise adjacent
supremum differences are summable.  Then the signs are eventually
constant on every fixed compact component and define a limiting
Pfaffian orientation there.
\end{corollary}

\begin{proof}
The same compact-screen Cauchy argument gives local limits
\(\eta_i\).  Equations (68.6)--(68.7) are closed algebraic
relations, so they pass to the limit and glue to
\(\jmath_\infty\).  Distinct signs have distance two.  A sequence
whose adjacent sign distances are summable can therefore jump only
finitely many times on each fixed compact component.
\end{proof}

This extends the attached fixed-screen result through the
exhaustion; it does not choose an orientation when no transported
finite-cutoff signs were supplied.

\subsection{Connections and sections}

If the transitions converge in \(C^1\) on compact overlaps and local
unitary connection forms \({\cal A}_{n,i}\) converge in \(C^0\) on
compact subsets while satisfying
\begin{equation}
 {\cal A}_{n,i}
 =
 {\cal A}_{n,j}
 +{\rm i}\,g_{n,ij}^{-1}dg_{n,ij},                \tag{68.9}
\end{equation}
then the limits define a unitary connection with continuous local
coefficients, and every holonomy along a fixed finite
piecewise-smooth loop passes to the limit.  If the connection forms
converge in \(C^1\) and the transitions in \(C^2\), their curvature
two-forms also converge locally uniformly and obey the limiting
curvature identity.

Likewise, suppose local coefficient functions \(s_{n,i}\) satisfy
\begin{equation}
                         s_{n,i}=g_{n,ij}s_{n,j}  \tag{68.10}
\end{equation}
and are locally uniformly Cauchy in the same physically transported
frames.  Then they converge to a section \(s_\infty\) of
\(\mathscr L_\infty\).  Its divisor is not determined merely by
local uniform convergence: zeros may appear in the limit unless a
local lower bound is also supplied.

\begin{proof}
Compact-open convergence permits passage to (68.9)--(68.10).
The connection transformation rule glues the local limits; line
integrals and exponentiation give loop holonomy convergence.
Equation (68.10) glues the limiting section.  The final warning
follows from the scalar example \(s_n\equiv1/n\).
\end{proof}

\subsection{Why a dual is still required}

A Hermitian metric gives an anti-linear Riesz map
\(\mathscr L_\infty\to\mathscr L_\infty^*\).  Pairing a section with
its Riesz dual gives \(\|s\|^2\), which contains no complex phase.
To obtain the one-determinant amplitude, assume an independently
specified complex-linear dual section
\begin{equation}
 t_n\in\Gamma(\mathscr L_n^*),                    \tag{68.11}
\end{equation}
transported on the same atlas.  Then
\begin{equation}
                         a_n(x):=t_n(x)(s_n(x))   \tag{68.12}
\end{equation}
is a gauge-invariant scalar.  If both local section packets converge
compact-open, then \(a_n\to a_\infty=t_\infty(s_\infty)\)
compact-open.  A real/Pfaffian branch instead uses its separately
transported orientation and direct Pfaffian sign.

\subsection{Coupling the paired line to the V66 state}

Use the bounded metric and exponent \(\alpha\) of V66.  Let
\(\mu_n,\mu\) be background laws on \(X\), let
\(\rho_n,\rho\) be the transported conditional density fields, and
let \(a_n,a:X\to\mathbb C\) be the gauge-invariant pairings
(68.12).  Assume
\begin{align}
 |a_n(x)|,\ |a(x)|&\leq A,                        \tag{68.13}\\
 |a(x)-a(y)|&\leq C_a\bar d(x,y)^\alpha,          \tag{68.14}\\
 \|\rho(x)-\rho(y)\|_1
 &\leq C_\rho\bar d(x,y)^\alpha,                  \tag{68.15}\\
 \epsilon_n^a&:=\sup_x|a_n(x)-a(x)|\to0,           \tag{68.16}\\
 \epsilon_n^\rho&:=\sup_x\|\rho_n(x)-\rho(x)\|_1\to0.          \tag{68.17}
\end{align}
Put \(w_n={\mathsf W}_\alpha(\mu_n,\mu)\) and define the unnormalized
paired trace-class amplitudes
\begin{equation}
 T_n:=\int_X a_n(x)\rho_n(x)\,d\mu_n(x),
 \qquad
 T:=\int_X a(x)\rho(x)\,d\mu(x).                   \tag{68.18}
\end{equation}
Their scalar partitions are
\begin{equation}
                         Z_n:=\operatorname{Tr}T_n,\qquad
 Z:=\operatorname{Tr}T.                           \tag{68.19}
\end{equation}

\begin{theorem}[Paired line-amplitude convergence]
\label{thm:v12v68-paired}
Under (68.13)--(68.17),
\begin{align}
 \|T_n-T\|_1
 &\leq
 \delta_n^T
 :=
 \epsilon_n^a+A\epsilon_n^\rho
 +(C_a+AC_\rho)w_n,                               \tag{68.20}\\
 |Z_n-Z|
 &\leq
 \delta_n^Z
 :=
 \epsilon_n^a+C_aw_n.                             \tag{68.21}
\end{align}
If
\begin{equation}
                         |Z|\geq\mathfrak v>0
 \quad\text{and}\quad
 \delta_n^Z\leq\mathfrak v/2,                     \tag{68.22}
\end{equation}
then the normalized complex trace-class functionals satisfy
\begin{equation}
 \boxed{
 \left\|{T_n\over Z_n}-{T\over Z}\right\|_1
 \leq
 {2\delta_n^T\over\mathfrak v}
 +
 {2A\delta_n^Z\over\mathfrak v^2}.}               \tag{68.23}
\end{equation}
\end{theorem}

\begin{proof}
Choose a coupling \(\pi_n\) of \(\mu_n,\mu\).  At two coupled points
\(x,y\),
\begin{align}
 \|a_n(x)\rho_n(x)-a(y)\rho(y)\|_1
 \leq{}&
 |a_n(x)-a(x)|
 +|a(x)|\|\rho_n(x)-\rho(x)\|_1                  \nonumber\\
 &+|a(x)-a(y)|
 +|a(y)|\|\rho(x)-\rho(y)\|_1                    \nonumber\\
 \leq{}&
 \epsilon_n^a+A\epsilon_n^\rho
 +(C_a+AC_\rho)\bar d(x,y)^\alpha.                \tag{68.24}
\end{align}
Integrate, use Bochner's inequality, and minimize over couplings to
obtain (68.20).  Since
\(Z_n=\int a_n\,d\mu_n\) and \(Z=\int a\,d\mu\), the scalar version
gives (68.21).

Equation (68.22) implies \(|Z_n|\geq\mathfrak v/2\).  Moreover
\(\|T\|_1\leq\int|a|\,d\mu\leq A\).  Therefore
\begin{align*}
 \left\|{T_n\over Z_n}-{T\over Z}\right\|_1
 &\leq
 {\|T_n-T\|_1\over|Z_n|}
 +
 {\|T\|_1|Z_n-Z|\over|Z_n||Z|}\\
 &\leq
 {2\delta_n^T\over\mathfrak v}
 +
 {2A\delta_n^Z\over\mathfrak v^2}.
\end{align*}
\end{proof}

The normalized object in (68.23) is generally a complex functional,
not a positive state.  It becomes the V66 positive state only on a
constant-phase positive branch, or after an independently justified
conjugate-doubled or real/Pfaffian construction.

\subsection{Cofinal version}

If adjacent line pairings, conditional states, and background laws
satisfy
\begin{equation}
 \sum_n
 \left\{
 \|a_{n+1}-a_n\|_\infty
 +A\|\rho_{n+1}-\rho_n\|_{L^\infty({\mathfrak S}_1)}
 +(C_a+AC_\rho)
 {\mathsf W}_\alpha(\mu_{n+1},\mu_n)
 \right\}
 <\infty,                                         \tag{68.25}
\end{equation}
then \(T_n\) is trace-norm Cauchy.  If, in addition, the scalar
partitions have a cofinal visibility floor
\begin{equation}
                         \inf_{n\geq n_0}|Z_n|>0, \tag{68.26}
\end{equation}
the normalized paired functionals are trace-norm Cauchy as well.
The proof is the adjacent form of (68.20)--(68.23).

The line-atlas summability (68.5), the section/dual convergence, the
background Wasserstein stock, and the visibility floor are four
separate rows.  None follows from anomaly cancellation.

\begin{assessment}[Status after V68]
\label{ass:v12v68-status}
The attached fixed-screen determinant-line theorem now extends to a
locally finite physical atlas through compact exhaustion.  Reflected
real structures, Pfaffian signs, connections, and sections pass
under their stated local summability.  Once an independent dual
section supplies a scalar amplitude, the V66 coupling estimate gives
explicit unnormalized and normalized complex-state convergence with
an average-phase visibility floor.
\end{assessment}

\begin{firewall}[Claim boundary after V68]
\label{fire:v12v68-claim}
No physical RPESM determinant atlas, dual section, Berry connection,
Pfaffian orientation, divisor control, or visibility floor has been
populated.  The theorem does not turn a Hermitian norm into a chiral
phase and does not infer line triviality from anomaly cancellation.
A nonconstant-phase normalized functional is not positive, so the
joint crossing OS gate of the attached programme remains required.
The full SM--Einstein continuum is not yet constructed.
\end{firewall}

\subsection{Parent record}

This V68 note appends to the validated V67 file with record
\[
\begin{array}{c|c}
\text{lines}&19165\\
\text{bytes}&634442\\
\text{SHA--256}&
\texttt{49C4747E66A9C637DBDBEC49607DDFF576B9735BFFA14D1FE7C748138DE8C3ED}.
\end{array}
\]

% =============================================================================
% END APPENDED TWELFTH-CYCLE V68 MATERIAL
% =============================================================================
% =============================================================================
% BEGIN APPENDED TWELFTH-CYCLE V69 MATERIAL
% =============================================================================

\section{Twelfth-cycle V69: gravitational compact containment and a
summable background subsequence}
\label{sec:v12v69}

The attached SM--Einstein V34 note already proves deterministic
\(C^{2,\alpha}\) convergence when one scheduler series has a
summable tail.  It does not construct probability laws for random
Einstein backgrounds.  V66 requires such laws and a bounded
Wasserstein convergence stock.  This note gives a distinct
probability theorem: uniform coercive moments on a gauge-fixed,
uniformly nondegenerate metric class imply tightness in a weaker
two-derivative topology, and every convergent subsequence contains a
further subsequence with a summable V66 transport budget.

\subsection{The admissible metric space}

Let \(M\) be compact without boundary, fix a smooth positive
reference metric \(g_0\), and choose
\[
                         0<\beta<\alpha<1.
\]
Let
\[
 X:=C^{2,\beta}(M;\operatorname{Sym}^2T^*M)
\]
with its norm metric, replaced by the bounded metric
\(\bar d=\min\{1,d_X\}\) as in V66.  Fix
\(0<\lambda<1\) and define the uniformly nondegenerate gauge-fixed
class
\begin{equation}
 {\cal G}_\lambda
 :=
 \{g\in C^{2,\alpha}:
   \lambda g_0\preceq g\preceq\lambda^{-1}g_0\}.  \tag{69.1}
\end{equation}
The coordinate or gauge condition selecting the representatives in
\({\cal G}_\lambda\) is treated as part of the definition of the
physical chart.

For \(R>0\), put
\begin{equation}
 K_R:=
 \{g\in{\cal G}_\lambda:\|g\|_{C^{2,\alpha}}\leq R\}.          \tag{69.2}
\end{equation}

\begin{lemma}[Metric compact screen]
\label{lem:v12v69-compact-screen}
The closure of \(K_R\) in \(X=C^{2,\beta}\) is compact and consists
of metrics satisfying the same nondegeneracy inequalities (69.1).
\end{lemma}

\begin{proof}
The compact Hölder embedding
\[
 C^{2,\alpha}(M)\hookrightarrow C^{2,\beta}(M)
\]
makes every bounded \(C^{2,\alpha}\) sequence precompact in
\(C^{2,\beta}\).  The tensor rank is finite, so the statement applies
componentwise in a finite atlas and patches by a partition of unity.
The pointwise inequalities in (69.1) are closed under uniform
\(C^0\) convergence, which is implied by \(C^{2,\beta}\)
convergence.
\end{proof}

\subsection{Moment tightness}

Let \(\mu_n\) be Borel probability measures on \(X\), supported on
\({\cal G}_\lambda\).

\begin{theorem}[Uniform metric moment implies tightness]
\label{thm:v12v69-moment-tightness}
If for some \(p>0\),
\begin{equation}
 \sup_n\int_X\|g\|_{C^{2,\alpha}}^p\,d\mu_n(g)
 \leq C_p<\infty,                                 \tag{69.3}
\end{equation}
then \((\mu_n)\) is tight on \(X\).  More precisely,
\begin{equation}
                         \mu_n(K_R^c)\leq C_pR^{-p}.           \tag{69.4}
\end{equation}
Hence a subsequence converges weakly to a probability measure
\(\mu_\infty\) supported on the \(C^{2,\beta}\)-closure of
\({\cal G}_\lambda\).
\end{theorem}

\begin{proof}
Markov's inequality gives (69.4).  Lemma
\ref{lem:v12v69-compact-screen} supplies a compact set carrying at
least \(1-C_pR^{-p}\) mass, uniformly in \(n\).  This is tightness.
The space \(X\) is Polish, so Prokhorov's theorem gives a weakly
convergent subsequence.  Since every compact screen closure obeys
(69.1), so does the support of the limit.
\end{proof}

This is compact containment of gauge-fixed metric representatives.
It does not prove tightness modulo diffeomorphisms unless a
measurable gauge section or a separate quotient-space argument is
supplied.

\subsection{A finite-action certificate for the moment}

Suppose the cutoff law has the positive form
\begin{equation}
 d\mu_n(g)
 =
 Z_n^{-1}e^{-S_n(g)}\,d\lambda_n(g),              \tag{69.5}
\end{equation}
where \(\lambda_n\) is a Borel probability measure on
\({\cal G}_\lambda\).  Assume there are constants
\(a>0,b,B<\infty\) and \(m>0\), independent of \(n\), together with
sets \(A_n\), such that
\begin{align}
 S_n(g)&\geq
 a\|g\|_{C^{2,\alpha}}^p-b
 \quad\text{for }\lambda_n\text{-almost every }g, \tag{69.6}\\
 \lambda_n(A_n)&\geq m,\qquad
 \sup_{g\in A_n}S_n(g)\leq B.                     \tag{69.7}
\end{align}

\begin{proposition}[Coercive action-to-tightness handoff]
\label{prop:v12v69-action}
Equations (69.6)--(69.7) imply
\begin{align}
 Z_n&\geq me^{-B},                                \tag{69.8}\\
 \int\|g\|_{C^{2,\alpha}}^p\,d\mu_n(g)
 &\leq
 {e^{b+B}\over ma e}.                             \tag{69.9}
\end{align}
Consequently the laws (69.5) are tight in \(C^{2,\beta}\).
\end{proposition}

\begin{proof}
Integrating over \(A_n\) gives (69.8).  From (69.6),
\[
 \|g\|_{C^{2,\alpha}}^pe^{-S_n(g)}
 \leq
 e^b r^pe^{-ar^p},
 \qquad r=\|g\|_{C^{2,\alpha}}.
\]
The scalar maximum of \(r^pe^{-ar^p}\) on \(r\geq0\) is
\(1/(ae)\).  Integrate against the probability \(\lambda_n\) and
divide by (69.8) to obtain (69.9).  Apply Theorem
\ref{thm:v12v69-moment-tightness}.
\end{proof}

The Euclidean Einstein--Hilbert action is not asserted to satisfy
(69.6); its conformal direction is a known obstruction to such a
naive coercive bound.  A physical application must obtain (69.6)
from the actual stabilized, constrained, gauge-fixed action or
replace it with another compact-containment Lyapunov estimate.

\subsection{From weak convergence to the V66 metric}

For \(0<\gamma\leq1\), use the bounded transport distance
\[
 {\mathsf W}_\gamma(\mu,\nu)
 =
 \inf_{\pi\in\Pi(\mu,\nu)}
 \int\bar d(g,h)^\gamma\,d\pi(g,h).
\]
On the Polish space \(X\), this distance metrizes weak convergence.

\begin{theorem}[Summable Wasserstein extraction]
\label{thm:v12v69-summable-subsequence}
If \(\mu_{n_k}\) converges weakly to \(\mu_\infty\), then it has a
further subsequence, relabeled \(\mu_{m_j}\), such that
\begin{align}
 {\mathsf W}_\gamma(\mu_{m_j},\mu_\infty)
 &\leq2^{-j},                                     \tag{69.10}\\
 \sum_{j=1}^\infty
 {\mathsf W}_\gamma(\mu_{m_{j+1}},\mu_{m_j})
 &\leq {3\over2}.                                 \tag{69.11}
\end{align}
More generally, \(2^{-j}\) may be replaced by any prescribed
summable positive sequence tending to zero.
\end{theorem}

\begin{proof}
Weak convergence is convergence in
\({\mathsf W}_\gamma\).  Choose successively \(m_j\) far enough along
the convergent subsequence that (69.10) holds.  The triangle
inequality gives
\[
 {\mathsf W}_\gamma(\mu_{m_{j+1}},\mu_{m_j})
 \leq2^{-(j+1)}+2^{-j}.
\]
Summing over \(j\geq1\) gives \(3/2\).
\end{proof}

Thus tightness supplies a summable background transport stock after
subsequence extraction.  It does not supply summability on the
original sequence.

\subsection{Einstein writers under the compact topology}

On the uniformly nondegenerate class (69.1), the maps
\begin{align}
 g&\longmapsto g^{-1},\quad
 |\det g|^{1/2},\quad
 \Gamma(g),                                      \tag{69.12}\\
 g&\longmapsto
 \operatorname{Rm}(g),\quad
 \mathscr E_\Lambda(g)                            \tag{69.13}
\end{align}
are continuous from \(C^{2,\beta}\) into their natural
\(C^{0,\beta}\) or \(C^{1,\beta}\) spaces.  Therefore weak
convergence of \(\mu_{m_j}\) passes every bounded continuous
cylinder built from these writers.

\begin{proof}
Each map is a finite combination of matrix inversion, determinant,
at most two derivatives, products, and contractions.  Uniform
nondegeneracy makes inversion smooth on the coefficient range, and
Hölder spaces are Banach algebras.  The standard coordinate formulas
used in the attached deterministic V34 theorem are therefore
continuous in the stated topology.  The continuous-mapping theorem
passes bounded continuous cylinders.
\end{proof}

Unbounded Einstein or stress insertions require uniform
integrability in addition to tightness.  The moment (69.3) supplies
it only for insertions dominated by a strictly smaller power of the
same \(C^{2,\alpha}\) norm, or after a separate de la
Vall\'ee--Poussin estimate.

\subsection{Handoff to the joint SM--Einstein limit}

Let \(\rho_n(g)\) be transported conditional SM states on the
reference Hilbert space.  Suppose the V64 form estimate gives one
uniform \(1/2\)-Hölder constant on the metric class and, along the
subsequence selected above,
\begin{equation}
 \sum_j
 \sup_g\|\rho_{m_{j+1}}(g)-\rho_{m_j}(g)\|_1
 <\infty.                                         \tag{69.14}
\end{equation}
Use \(\gamma=1/2\) in Theorem
\ref{thm:v12v69-summable-subsequence}.  Then
\begin{equation}
 \sum_j
 \left\{
 \sup_g\|\rho_{m_{j+1}}(g)-\rho_{m_j}(g)\|_1
 +
 {\mathsf W}_{1/2}(\mu_{m_{j+1}},\mu_{m_j})
 \right\}<\infty.                                 \tag{69.15}
\end{equation}

\begin{corollary}[Conditional gravitational handoff]
\label{cor:v12v69-joint}
Under (69.3), the V64 uniform form hypotheses, and (69.14), the
selected subsequence satisfies the cofinal criterion (66.17).
It therefore has a joint classical--quantum state limit with the
barycentre and bounded-Hölder convergence rates of V66.  If the
finite OS, Ward, symmetry, and determinant-line hypotheses of
V67--V68 also hold with summable residuals on the same subsequence,
those structures pass to the joint limit.
\end{corollary}

\begin{proof}
Tightness and Theorem
\ref{thm:v12v69-summable-subsequence} give the summable background
part of (69.15).  Equation (69.14) gives the conditional part.
Apply Theorem \ref{thm:v12v66-cofinal}, then Theorem
\ref{thm:v12v67-joint-structure} and the V68 paired-line bounds for
the additional conditional assertions.
\end{proof}

\subsection{Nontriviality and uniqueness are separate}

The tight limit may be a delta measure or may depend on the selected
subsequence.  To prevent collapse, one needs a uniform lower bound
on the mass of a physically nontrivial open cylinder, or an
equivalent variance/response floor.  To obtain a unique full-sequence
limit, one needs a uniqueness or contraction theorem, or a direct
full-sequence Cauchy estimate.  Neither follows from (69.3).

\begin{assessment}[Status after V69]
\label{ass:v12v69-status}
A uniform \(C^{2,\alpha}\) moment on uniformly nondegenerate
gauge-fixed metrics gives gravitational tightness in
\(C^{2,\beta}\).  An explicit coercive action and partition-floor
certificate yields that moment.  Every weakly convergent subsequence
can be thinned to have a summable bounded-Wasserstein budget, which
feeds directly into the V64/V66 joint-state construction.
\end{assessment}

\begin{firewall}[Claim boundary after V69]
\label{fire:v12v69-claim}
No physical RPESM gravitational action has been proved coercive, and
the conformal instability is not resolved.  No gauge section,
partition floor, nontriviality estimate, uniqueness theorem,
conditional-state stock (69.14), or uniform-integrability bound for
unbounded stress insertions has been populated.  The theorem is a
conditional tightness and subsequence handoff, not a construction of
the full SM--Einstein continuum.
\end{firewall}

\subsection{Parent record}

This V69 note appends to the validated V68 file with record
\[
\begin{array}{c|c}
\text{lines}&19533\\
\text{bytes}&647798\\
\text{SHA--256}&
\texttt{72B81F47AE8B7616C350E8A4208B62BDC3CB520AA3FA71D003740ACDDD35EFBF}.
\end{array}
\]

% =============================================================================
% END APPENDED TWELFTH-CYCLE V69 MATERIAL
% =============================================================================
% =============================================================================
% BEGIN APPENDED TWELFTH-CYCLE V70 MATERIAL
% =============================================================================

\section{Twelfth-cycle V70: mandatory companion audit and
fibre-first gravitational compact containment}
\label{sec:v12v70}

This multiple-of-five checkpoint reads the current Yang--Mills and
Navier--Stokes companions before changing the gravitational
probability route of V69.  The audit supplies a proof-order rule.
It does not import either programme's numerical constants.

\subsection{Mandatory V70 companion audit}

The live companion records inspected read-only are
\[
\begin{array}{c|r|r|c}
\text{record}&\text{lines}&\text{bytes}&\text{SHA--256}\\ \hline
\text{Yang--Mills V56}&21223&723854&
\texttt{627247BCFC1A511C5158B1D3DED869FCA598501597EF3B8CF174C3A0BFEF54DC}\\
\text{Navier--Stokes V26}&5319&278283&
\texttt{82C887F8C2C69E4F28FAD7C3DC8DE5B9099CB5A4BF431EBA1FFF3E91935978AD}.
\end{array}                                                   \tag{70.1}
\]
These are unchanged from the V65 audit.

Yang--Mills V55 proves that an absolute Fourier-tail estimate is
legal only after all cards in the physical shift packet have been
traced and their signed congruence core retained.  V56 proves the
sharper information-loss statement that values on all parity corners
do not determine the arbitrary-momentum Laurent observable.  The
corresponding gravitational rule is that one must first integrate or
constrain the complete conformal fibre.  A pointwise lower estimate
on selected modes, or a finite collection of Hessian values, cannot
stand in for an estimate on the resulting marginal weight.

Navier--Stokes V26 proves that its smaller rank-one tensor norm is
usable only after the actual input is shown to remain in the
rank-one cone.  Here the analogue is that a coercive estimate on a
gauge slice is usable only after the regulated law is proved to be
supported on that slice, or after the quotient law is constructed
directly.

\begin{decision}[Effect of the V70 audit]
\label{dec:v12v70-audit}
Retain V61--V69 unchanged.  Replace the attempted pointwise attack on
the indefinite conformal action by the following order:
\[
 \boxed{
 \text{complete conformal fibre}
 \longrightarrow
 \text{positive marginal weight}
 \longrightarrow
 \text{effective coercivity and fibre tails}
 \longrightarrow
 \text{compact containment}.}                              \tag{70.2}
\]
No cancellation-sensitive or slice-sensitive estimate may be applied
before the object on its left has been constructed.
\end{decision}

\subsection{A fibre-first probability theorem}

Let \(B\) and \(F\) be Polish spaces.  A point of the regulated
configuration chart is denoted \((b,\varphi)\), where \(b\) contains
the nonconformal metric variables and all other fields, while
\(\varphi\) is the complete regulated conformal fibre.  For every
\(n\), let \(\eta_n\) be a probability measure on \(B\), and let
\(b\mapsto\kappa_n^b\) be a Borel probability kernel on \(F\).
Let \(S_n:B\times F\to(-\infty,\infty]\) be Borel and set
\begin{align}
 q_n(b)&:=\int_F e^{-S_n(b,\varphi)}\,d\kappa_n^b(\varphi),
                                                               \tag{70.3}\\
 Z_n&:=\int_B q_n(b)\,d\eta_n(b).                              \tag{70.4}
\end{align}
Assume \(0<q_n(b)<\infty\) for \(\eta_n\)-almost every \(b\) and
\(0<Z_n<\infty\).  The full law and its base marginal are
\begin{align}
 d{\mathbb P}_n(b,\varphi)
 &=Z_n^{-1}e^{-S_n(b,\varphi)}
   d\kappa_n^b(\varphi)d\eta_n(b),                            \tag{70.5}\\
 d\overline\mu_n(b)&=Z_n^{-1}q_n(b)\,d\eta_n(b).              \tag{70.6}
\end{align}
The conditional conformal law is
\[
 d\nu_n^b(\varphi)
 :=q_n(b)^{-1}e^{-S_n(b,\varphi)}d\kappa_n^b(\varphi).        \tag{70.7}
\]

Let \(U:B\to[0,\infty]\) and
\(V:B\times F\to[0,\infty]\) be Borel.  Assume that for all
\(R,T<\infty\), the closure in \(B\times F\) of
\begin{equation}
 {\cal C}_{R,T}
 :=\{(b,\varphi):U(b)\leq R,\ V(b,\varphi)\leq T\}            \tag{70.8}
\end{equation}
is compact.

\begin{theorem}[Fibre-first compact-containment criterion]
\label{thm:v12v70-fibre-first}
Suppose constants
\[
 A,a,C,c,z_0>0
\]
independent of \(n\) satisfy
\begin{align}
 Z_n&\geq z_0,                                                \tag{70.9}\\
 q_n(b)&\leq A e^{-aU(b)}
 \quad\text{for }\eta_n\text{-almost every }b,               \tag{70.10}\\
 \nu_n^b\{V(b,\varphi)>T\}&\leq C e^{-cT}
 \quad\text{for }\overline\mu_n\text{-almost every }b
 \text{ and every }T\geq0.                                  \tag{70.11}
\end{align}
Then
\begin{equation}
 {\mathbb P}_n({\cal C}_{R,T}^{\,c})
 \leq {A\over z_0}e^{-aR}+Ce^{-cT}.                          \tag{70.12}
\end{equation}
Consequently \(({\mathbb P}_n)\) is tight on \(B\times F\), and
\((\overline\mu_n)\) is tight on \(B\).  In addition,
\begin{align}
 \sup_n\int_B U\,d\overline\mu_n
 &\leq {A\over a z_0},                                      \tag{70.13}\\
 \sup_n\int_{B\times F}V\,d{\mathbb P}_n
 &\leq {C\over c}.                                          \tag{70.14}
\end{align}
\end{theorem}

\begin{proof}
By (70.6), (70.9), and (70.10),
\[
 \overline\mu_n\{U>R\}
 ={1\over Z_n}\int_{\{U>R\}}q_n\,d\eta_n
 \leq {A\over z_0}e^{-aR},                                  \tag{70.15}
\]
because \(\eta_n\) is a probability measure.  Disintegration
(70.5)--(70.7) and (70.11) give
\[
 {\mathbb P}_n\{V>T\}
 =\int_B\nu_n^b\{V(b,\cdot)>T\}\,d\overline\mu_n(b)
 \leq Ce^{-cT}.                                              \tag{70.16}
\]
The complement of \({\cal C}_{R,T}\) is contained in
\(\{U>R\}\cup\{V>T\}\).  The union bound proves (70.12).
The compact-closure hypothesis in (70.8), followed by arbitrarily
large \(R,T\), proves tightness.  Finally, the layer-cake formula and
(70.15)--(70.16) give
\[
 \int U\,d\overline\mu_n
 =\int_0^\infty\overline\mu_n\{U>R\}\,dR
 \leq {A\over az_0},
\]
and the same calculation gives (70.14).
\end{proof}

If the effective conformal integral is written
\begin{equation}
                         S_n^{\rm eff}(b):=-\log q_n(b),      \tag{70.17}
\end{equation}
then (70.10) is exactly the lower bound
\begin{equation}
 S_n^{\rm eff}(b)\geq aU(b)-\log A.                          \tag{70.18}
\end{equation}
It is a statement about the completed fibre integral.  The original
action \(S_n(b,\varphi)\) need not have a pointwise lower bound in
\(\varphi\).

\begin{proposition}[A partition floor after fibre integration]
\label{prop:v12v70-floor}
Suppose there are Borel sets \(D_n\subset B\) and constants
\(m,B_0>0\), independent of \(n\), such that
\begin{equation}
 \eta_n(D_n)\geq m,\qquad q_n(b)\geq e^{-B_0}
 \quad(b\in D_n).                                           \tag{70.19}
\end{equation}
Then
\begin{equation}
                              Z_n\geq me^{-B_0}.              \tag{70.20}
\end{equation}
Thus (70.19), (70.10), and (70.11) imply Theorem
\ref{thm:v12v70-fibre-first} with \(z_0=me^{-B_0}\).
\end{proposition}

\begin{proof}
Integrate the second inequality in (70.19) over \(D_n\) with respect
to \(\eta_n\):
\[
 Z_n=\int_Bq_n\,d\eta_n
 \geq\int_{D_n}q_n\,d\eta_n
 \geq me^{-B_0}.
\]
\end{proof}

\subsection{Relation to V69}

Take \(B\times F\) to be a regulated metric-and-matter chart and
choose \(U,V\) so that (70.8) implies a uniform
\(C^{2,\alpha}\) bound and the nondegeneracy inequalities (69.1).
The compact Hölder embedding of Lemma
\ref{lem:v12v69-compact-screen} then verifies the compact-closure
hypothesis.  Theorem \ref{thm:v12v70-fibre-first} supplies the
moment/tightness input of V69 without requiring the pointwise
coercivity (69.6) of the full bare action.  Prokhorov extraction and
Theorem \ref{thm:v12v69-summable-subsequence} then give a summable
bounded-Wasserstein background subsequence.

The theorem separates two estimates that cannot replace each other.
Effective base coercivity (70.10) gives no conformal-fibre
containment.  The conditional tail (70.11) gives no base
containment.  Both are needed for the full metric law.

\begin{warning}[Complex contours require a different gate]
\label{warn:v12v70-complex}
The proof uses the positive conditional probability (70.7).  If the
conformal fibre is an imaginary or more general complex contour,
\(q_n\) may be complex and cancellation may make its logarithm
multivalued or its value zero.  In that branch, (70.9)--(70.11) must
be replaced by total-variation domination, relative-homology control,
and a nonzero denominator/phase anchor of the kind developed in
V1--V12.  Taking absolute values inside (70.3) would erase precisely
the cancellation that the V70 audit requires one to preserve.
\end{warning}

\begin{assessment}[Status after V70]
\label{ass:v12v70-status}
The mandatory companion audit is complete.  It changes the
gravitational acquisition order from pointwise coercivity to a
fibre-first marginal estimate.  The new theorem proves tightness of
the full positive law from effective base coercivity, a partition
floor, and a uniform conditional conformal tail, with explicit
exponential compact-containment and first-moment bounds.
\end{assessment}

\begin{firewall}[Claim boundary after V70]
\label{fire:v12v70-claim}
No physical conformal integral \(q_n\), positivity theorem,
effective-action bound (70.10), conditional tail (70.11), partition
floor, or compact gauge quotient has been constructed for the full
RPESM action.  Finite conformal Hessian data do not establish these
nonlinear uniform estimates.  The complex-contour branch remains
subject to its independent denominator and reflection-positivity
gates.  The full SM--Einstein continuum is not yet constructed.
\end{firewall}

\subsection{Parent record}

This V70 note appends to the validated V69 file with record
\[
\begin{array}{c|c}
\text{lines}&19849\\
\text{bytes}&659268\\
\text{SHA--256}&
\texttt{B56051BC6FE2C3081B63BA82C84197A0FDBD1DF138C57E95044590F0889C0E35}.
\end{array}
\]

% =============================================================================
% END APPENDED TWELFTH-CYCLE V70 MATERIAL
% =============================================================================

% =============================================================================
% BEGIN APPENDED TWELFTH-CYCLE V71 MATERIAL
% =============================================================================

\section{Twelfth-cycle V71: compact containment directly on the
diffeomorphism quotient}
\label{sec:v12v71}

V69 obtained tightness for gauge-fixed metric representatives and
correctly left a measurable global gauge section open.  V70 shows
that support on such a slice must be established before slice-specific
coercivity can be used.  This note gives a separate route: when the
orbit space has a Polish metric compatible with the quotient
topology, invariant compact containment and bounded-Wasserstein
convergence descend directly to that orbit space.  A global section
is not part of the theorem.

\subsection{Metric quotient hypotheses}

Let \((X,d)\) be a Polish space with \(d\leq1\), and let a group
\({\cal D}\) act continuously on \(X\).  Write
\[
                         \pi:X\longrightarrow Y:=X/{\cal D}.
\]
Assume \(Y\) is Hausdorff and carries a complete separable metric
\(\delta\leq1\) that induces the quotient topology and satisfies
\begin{equation}
                 \delta(\pi x,\pi y)\leq d(x,y).              \tag{71.1}
\end{equation}
For an isometric action, the natural candidate is
\[
 \delta(\pi x,\pi y)
 :=\inf_{\psi\in{\cal D}}d(x,\psi y),                         \tag{71.2}
\]
provided it separates orbits and its metric topology is the quotient
topology.  These provisos are essential for a noncompact
diffeomorphism group.

\begin{lemma}[A compact-group sufficient condition]
\label{lem:v12v71-compact-group}
If \({\cal D}\) is compact and acts continuously by isometries on
\(X\), then (71.2) is a metric on \(Y\), it induces the quotient
topology, \(Y\) is Polish, and (71.1) holds.
\end{lemma}

\begin{proof}
The triangle inequality follows by composing two group elements and
using invariance of \(d\).  Compactness of \({\cal D}\) makes the
infimum in (71.2) attained.  If the distance of two orbits is zero,
a minimizing group element maps one point to the other, so the
orbits agree.  The inequality (71.1) follows by using the identity
element.

The quotient map is open: the saturation of an open set is the
union of its translates.  Hence the metric and quotient topologies
agree.  The image of a countable dense subset of \(X\) is dense in
\(Y\), so \(Y\) is separable.  If \((\pi x_n)\) is
\(\delta\)-Cauchy, choose a subsequence \(n_j\) and group elements
\(\psi_j\) recursively so that
\[
 d(\psi_jx_{n_j},\psi_{j+1}x_{n_{j+1}})<2^{-j}.
\]
The adjusted subsequence converges in \(X\), hence its orbit classes
converge in \(Y\).  The original Cauchy sequence has the same limit.
Thus \(Y\) is complete.
\end{proof}

The compact-group lemma is a sufficient model, not a claim that the
full diffeomorphism group is compact.  The later theorems use only
the stated metric quotient hypotheses.

\subsection{Tightness modulo gauge without representative tightness}

Let \(L:X\to[0,\infty]\) be a Borel, orbit-invariant function:
\[
                         L(\psi x)=L(x).                      \tag{71.3}
\]
It therefore induces \(\overline L:Y\to[0,\infty]\) by
\(\overline L(\pi x)=L(x)\).

\begin{theorem}[Quotient compact-containment criterion]
\label{thm:v12v71-quotient-tight}
Let \((\mu_n)\) be Borel probability measures on \(X\).  Suppose
\begin{align}
 \sup_n\mu_n\{L>R\}&\leq\varepsilon(R),
 &\varepsilon(R)&\longrightarrow0,                           \tag{71.4}\\
 \overline{\{\pi x:L(x)\leq R\}}^{\,Y}
 &\text{ is compact for every }R<\infty.                     \tag{71.5}
\end{align}
Then the orbit laws
\[
                         \overline\mu_n:=\pi_\#\mu_n          \tag{71.6}
\]
are tight on \(Y\).  In particular, they have a weakly convergent
subsequence.  No tightness of \((\mu_n)\) on \(X\) is required.
\end{theorem}

\begin{proof}
Let
\[
 K_R:=\overline{\{\pi x:L(x)\leq R\}}^{\,Y}.
\]
It is compact by (71.5).  Orbit invariance gives
\[
 \overline\mu_n(K_R^c)
 \leq\overline\mu_n\{\overline L>R\}
 =\mu_n\{L>R\}
 \leq\varepsilon(R).                                        \tag{71.7}
\]
Letting \(R\) increase proves tightness.  Since \(Y\) is Polish,
Prokhorov's theorem gives the subsequence.
\end{proof}

This theorem permits arbitrary drift along diffeomorphism orbits.
It asks instead for a proper physical Lyapunov function on the
quotient.  Thus a family may fail every fixed-coordinate
\(C^{2,\alpha}\) compact screen in \(X\) and still be tight in
physical geometry.

\begin{corollary}[Fibre-first quotient tightness]
\label{cor:v12v71-fibre-quotient}
In Theorem \ref{thm:v12v70-fibre-first}, suppose the full
configuration space \(B\times F\) carries an action satisfying the
metric quotient hypotheses, and suppose the compact-control
functions \(U\) and \(V\) combine to an orbit-invariant function
\[
                         L(b,\varphi):=U(b)+V(b,\varphi).     \tag{71.8}
\]
If the quotient closures of its sublevel sets are compact, then the
orbit laws \(\pi_\#{\mathbb P}_n\) are tight.  Explicitly,
\begin{equation}
 (\pi_\#{\mathbb P}_n)\{\overline L>R\}
 \leq {A\over z_0}e^{-aR/2}+Ce^{-cR/2}.                     \tag{71.9}
\end{equation}
\end{corollary}

\begin{proof}
If \(U+V>R\), then \(U>R/2\) or \(V>R/2\).  Apply
(70.15)--(70.16), followed by Theorem
\ref{thm:v12v71-quotient-tight}.
\end{proof}

\subsection{Transport distances descend}

For \(0<\gamma\leq1\), write
\[
 {\mathsf W}_{d,\gamma}(\mu,\nu)
 :=\inf_{\Pi\in\Pi(\mu,\nu)}
   \int d(x,x')^\gamma\,d\Pi(x,x')
\]
and define \({\mathsf W}_{\delta,\gamma}\) analogously on \(Y\).

\begin{theorem}[Wasserstein contraction under the quotient]
\label{thm:v12v71-Wasserstein}
For all probability measures \(\mu,\nu\) on \(X\),
\begin{equation}
 {\mathsf W}_{\delta,\gamma}(\pi_\#\mu,\pi_\#\nu)
 \leq{\mathsf W}_{d,\gamma}(\mu,\nu).                        \tag{71.10}
\end{equation}
Consequently
\begin{equation}
 \sum_n{\mathsf W}_{d,\gamma}(\mu_{n+1},\mu_n)<\infty
 \quad\Longrightarrow\quad
 \sum_n{\mathsf W}_{\delta,\gamma}
     (\pi_\#\mu_{n+1},\pi_\#\mu_n)<\infty.                   \tag{71.11}
\end{equation}
If only quotient tightness is known, every weakly convergent quotient
subsequence has a further subsequence with a summable
\({\mathsf W}_{\delta,\gamma}\) budget, as in V69.
\end{theorem}

\begin{proof}
For any coupling \(\Pi\) of \(\mu,\nu\), the pushforward
\((\pi,\pi)_\#\Pi\) couples \(\pi_\#\mu\) and \(\pi_\#\nu\).
Equation (71.1) gives
\[
 \int\delta(\pi x,\pi x')^\gamma\,d\Pi
 \leq\int d(x,x')^\gamma\,d\Pi.
\]
Take the infimum over \(\Pi\) to obtain (71.10), then sum to get
(71.11).  On a bounded Polish metric space the corresponding
Wasserstein distance metrizes weak convergence; the thinning proof
of Theorem \ref{thm:v12v69-summable-subsequence} applies on \(Y\).
\end{proof}

\subsection{Physical observables and conditional states}

\begin{proposition}[Descent of invariant observables]
\label{prop:v12v71-observable-descent}
If \(f:X\to E\) is a continuous orbit-invariant map into a Hausdorff
space \(E\), there is a unique continuous
\(\overline f:Y\to E\) such that
\[
                              f=\overline f\circ\pi.          \tag{71.12}
\]
For every probability measure \(\mu\) and every bounded scalar
\(f\),
\begin{equation}
                  \int_X f\,d\mu
                  =\int_Y\overline f\,d(\pi_\#\mu).          \tag{71.13}
\end{equation}
\end{proposition}

\begin{proof}
Set \(\overline f(\pi x)=f(x)\).  Orbit invariance makes this
well-defined and unique.  A map out of a quotient is continuous
exactly when its pullback by the quotient map is continuous, proving
(71.12).  Equation (71.13) is the definition of pushforward.
\end{proof}

Let \({\cal T}_1({\cal H}_*)\) be the trace-class operators on the V61
reference Hilbert space.  Suppose the transported conditional state
field is physically invariant,
\[
                         \rho(\psi x)=\rho(x),                \tag{71.14}
\]
and obeys the physical Hölder estimate
\begin{equation}
 \|\rho(x)-\rho(x')\|_1
 \leq C_\rho\delta(\pi x,\pi x')^\gamma.                     \tag{71.15}
\end{equation}
Proposition \ref{prop:v12v71-observable-descent} gives a unique
Hölder field \(\overline\rho:Y\to{\cal T}_1({\cal H}_*)\).

\begin{corollary}[V66 construction on orbit space]
\label{cor:v12v71-joint-quotient}
Replace the background space of V66 by \(Y\), its bounded metric by
\(\delta\), and the conditional state by
\(\overline\rho_n\).  If
\begin{equation}
 \sum_n\left\{
 \sup_{y\in Y}\|\overline\rho_{n+1}(y)-\overline\rho_n(y)\|_1
 +{\mathsf W}_{\delta,\gamma}
       (\overline\mu_{n+1},\overline\mu_n)\right\}<\infty,    \tag{71.16}
\end{equation}
then V66 constructs the joint classical--quantum continuum state
directly on the physical orbit space.  The result is independent of
any choice of representatives.
\end{corollary}

\begin{proof}
All hypotheses of Theorem \ref{thm:v12v66-cofinal} hold on the
bounded Polish space \((Y,\delta)\).  Apply that theorem.
Representative independence follows from (71.12)--(71.14).
\end{proof}

\subsection{The remaining quotient gate}

For the actual diffeomorphism action, the hard inputs have now been
isolated:
\begin{enumerate}
\item a Hausdorff Polish physical orbit space, or a specified regular
      stratum on which the metric quotient hypotheses hold;
\item compact quotient sublevels of a physical Lyapunov function;
\item equivariance of the finite laws and of the V61 transports;
\item treatment of stabilisers, Gribov multiplicities, and singular
      strata in the measure, not merely in the set-theoretic quotient.
\end{enumerate}
A local Faddeev--Popov gap from the frozen notes helps construct local
slices, but does not by itself prove these global properties.

\begin{assessment}[Status after V71]
\label{ass:v12v71-status}
Gravitational compact containment no longer logically depends on a
global measurable gauge section.  Under explicit Polish metric
quotient and compact physical-sublevel hypotheses, invariant tails
give quotient tightness, bounded Wasserstein budgets contract under
the quotient, and the V66 joint state is constructed on orbit space.
\end{assessment}

\begin{firewall}[Claim boundary after V71]
\label{fire:v12v71-claim}
The full diffeomorphism orbit space has not been proved Hausdorff or
Polish, and no compact physical Lyapunov sublevels, equivariant
finite-cutoff laws, stabiliser weights, or Gribov corrections have
been populated.  The compact-group lemma is only a sufficient model
and is not being applied to the noncompact diffeomorphism group.
The full SM--Einstein continuum is not yet constructed.
\end{firewall}

\subsection{Parent record}

This V71 note appends to the validated V70 file with record
\[
\begin{array}{c|c}
\text{lines}&20110\\
\text{bytes}&669450\\
\text{SHA--256}&
\texttt{FB40A4E16AF8E9A5676F18BE032853A978070A20BDE8B91CD5F9386113960B63}.
\end{array}
\]

% =============================================================================
% END APPENDED TWELFTH-CYCLE V71 MATERIAL
% =============================================================================

% =============================================================================
% BEGIN APPENDED TWELF-CYCLE V72 MATERIAL
% =============================================================================

\section{Twelfth-cycle V72: a best-representative Hölder Lyapunov
function and compact orbit screens}
\label{sec:v12v72}

V71 reduces quotient tightness to compact sublevels of an invariant
Lyapunov function.  This note constructs such a function from the
same strong Hölder screens used in V69.  Its infimum is taken over
each diffeomorphism orbit.  The compactness proof chooses a
representative separately for each member of a sequence; it does not
construct, or require, a measurable global gauge section.

\subsection{The orbit-infimum control}

Let \(M\) be compact, fix \(0<\beta<\alpha<1\), a smooth reference
metric \(g_0\), and \(0<\lambda<1\).  Put
\[
 X:=\operatorname{Met}^{2,\beta}(M),
\]
viewed as an open subset of
\(C^{2,\beta}(M;\operatorname{Sym}^2T^*M)\).  Let
\({\cal D}\) be a specified group of diffeomorphisms whose pullback
action preserves \(X\), and suppose its orbit space \(Y=X/{\cal D}\)
satisfies the metric quotient hypotheses of V71.

Define the extended strong norm
\begin{equation}
 {\cal N}_{\lambda,\alpha}(h)
 :=
 \begin{cases}
  \|h\|_{C^{2,\alpha}},&
  \lambda g_0\preceq h\preceq\lambda^{-1}g_0,\\
  +\infty,&\text{otherwise},
 \end{cases}                                                \tag{72.1}
\end{equation}
and its orbit infimum
\begin{equation}
 {\cal L}_{\lambda,\alpha}(g)
 :=
 \inf_{\psi\in{\cal D}}
 {\cal N}_{\lambda,\alpha}(\psi^*g).                         \tag{72.2}
\end{equation}
For probability statements below, assume
\({\cal L}_{\lambda,\alpha}\) is Borel.  This measurability is
separate from compactness and must be checked for the chosen
diffeomorphism class and chart.

\begin{lemma}[Orbit invariance]
\label{lem:v12v72-invariance}
For every \(\chi\in{\cal D}\),
\begin{equation}
 {\cal L}_{\lambda,\alpha}(\chi^*g)
 ={\cal L}_{\lambda,\alpha}(g).                              \tag{72.3}
\end{equation}
Hence (72.2) defines a function
\(\overline{\cal L}_{\lambda,\alpha}:Y\to[0,\infty]\).
\end{lemma}

\begin{proof}
Since \({\cal D}\) is a group,
\(\{\psi\circ\chi:\psi\in{\cal D}\}={\cal D}\).  Therefore
\[
 \inf_{\psi\in{\cal D}}
 {\cal N}_{\lambda,\alpha}(\psi^*\chi^*g)
 =
 \inf_{\theta\in{\cal D}}
 {\cal N}_{\lambda,\alpha}(\theta^*g).
\]
This is (72.3), up to the immaterial convention for the order of
pullbacks.
\end{proof}

\subsection{Compactness of the physical sublevels}

\begin{theorem}[Best-representative compact orbit screen]
\label{thm:v12v72-compact-orbit}
For every \(R<\infty\), the closure in \(Y\) of
\begin{equation}
 {\cal K}_R^Y
 :=
 \{\pi g:{\cal L}_{\lambda,\alpha}(g)\leq R\}                \tag{72.4}
\end{equation}
is compact.
\end{theorem}

\begin{proof}
It suffices to prove that every sequence in (72.4) has a convergent
subsequence.  Choose orbit classes \(y_j=\pi g_j\) in (72.4).
By the definition of the infimum, there is
\(\psi_j\in{\cal D}\) such that
\begin{equation}
 h_j:=\psi_j^*g_j,\qquad
 {\cal N}_{\lambda,\alpha}(h_j)\leq R+j^{-1}.                \tag{72.5}
\end{equation}
Thus \((h_j)\) is bounded in \(C^{2,\alpha}\) and obeys
\[
                   \lambda g_0\preceq h_j
                   \preceq\lambda^{-1}g_0.                  \tag{72.6}
\]
The compact Hölder embedding
\[
 C^{2,\alpha}\hookrightarrow C^{2,\beta}
\]
gives a subsequence, not relabeled, and an \(h\) such that
\(h_j\to h\) in \(C^{2,\beta}\).  Uniform convergence and (72.6)
show that \(h\) is a metric satisfying the same nondegeneracy
bounds.  In particular \(h\in X\).

Because \(\pi\) is continuous,
\[
                         y_j=\pi h_j\longrightarrow\pi h
                         \quad\text{in }Y.                   \tag{72.7}
\]
Thus (72.4) is relatively sequentially compact.  Its closure is
sequentially compact, and a sequentially compact subset of the
metric space \(Y\) is compact.
\end{proof}

The representatives \(\psi_j^*g_j\) in (72.5) may depend on the
sequence and on \(j\).  The proof makes no consistency, continuity,
or measurability assertion about \(g\mapsto\psi_g\).

\begin{corollary}[Invariant moment implies quotient tightness]
\label{cor:v12v72-moment}
Let \(\mu_n\) be Borel probability measures on \(X\).  If for some
\(p>0\),
\begin{equation}
 \sup_n\int_X
 {\cal L}_{\lambda,\alpha}(g)^p\,d\mu_n(g)
 \leq C_p<\infty,                                           \tag{72.8}
\end{equation}
then the orbit laws \(\pi_\#\mu_n\) are tight on \(Y\), with
\begin{equation}
 (\pi_\#\mu_n)\{
 \overline{\cal L}_{\lambda,\alpha}>R\}
 \leq C_pR^{-p}.                                            \tag{72.9}
\end{equation}
They therefore have a weakly convergent subsequence.
\end{corollary}

\begin{proof}
Markov's inequality gives (72.9), Lemma
\ref{lem:v12v72-invariance} identifies the event on \(Y\), and
Theorem \ref{thm:v12v72-compact-orbit} supplies the compact screens.
Apply Theorem \ref{thm:v12v71-quotient-tight}.
\end{proof}

\subsection{Effective action criterion on orbit space}

Suppose the fibre-integrated positive marginal of V70 is a law
\[
 d\mu_n(g)=Z_n^{-1}q_n(g)\,d\eta_n(g)                        \tag{72.10}
\]
on \(X\), where \(\eta_n\) is a probability measure.  This formula
may already include the matter, gauge, ghost, and conformal
conditional integrals, provided the result is a positive Borel
weight.

\begin{proposition}[Physical effective coercivity]
\label{prop:v12v72-effective}
Assume constants \(A,a,z_0>0\), independent of \(n\), satisfy
\begin{align}
 Z_n&\geq z_0,                                               \tag{72.11}\\
 q_n(g)&\leq
 A\exp\{-a{\cal L}_{\lambda,\alpha}(g)^p\}
 \quad\text{for }\eta_n\text{-almost every }g.              \tag{72.12}
\end{align}
Then for every \(R>0\),
\begin{equation}
 (\pi_\#\mu_n)\{
 \overline{\cal L}_{\lambda,\alpha}>R\}
 \leq {A\over z_0}e^{-aR^p}.                                \tag{72.13}
\end{equation}
In particular the quotient laws are tight.
\end{proposition}

\begin{proof}
On the event
\({\cal L}_{\lambda,\alpha}>R\), equation (72.12) bounds
\(q_n\) by \(Ae^{-aR^p}\).  Integrating against the probability
\(\eta_n\), dividing by (72.11), and using orbit invariance gives
(72.13).  Theorem \ref{thm:v12v72-compact-orbit} then proves
tightness.
\end{proof}

Equivalently, for
\(S_n^{\rm eff}:=-\log q_n\), condition (72.12) is
\begin{equation}
 S_n^{\rm eff}(g)
 \geq a{\cal L}_{\lambda,\alpha}(g)^p-\log A.                \tag{72.14}
\end{equation}
This is a coercive lower bound in the best representative of each
physical orbit.  It does not demand coercivity in every coordinate
representative.

\subsection{Comparison with a gauge-fixed screen}

If a genuine gauge section
\(\sigma:Y\to X\) exists and satisfies
\[
 {\cal N}_{\lambda,\alpha}(\sigma(\pi g))
 \leq C\{1+{\cal L}_{\lambda,\alpha}(g)\},                   \tag{72.15}
\]
then an orbit moment controls the section representatives.  But
(72.15) is not used in Theorem
\ref{thm:v12v72-compact-orbit}: the latter is strictly a
compactness-modulo-diffeomorphism statement.

Conversely, a moment bound on arbitrary representatives does imply
an orbit moment because
\[
 {\cal L}_{\lambda,\alpha}(g)
 \leq{\cal N}_{\lambda,\alpha}(g),                           \tag{72.16}
\]
whenever the identity diffeomorphism is admissible.  Thus V69 is a
stronger sufficient route when its gauge-fixed moment is available,
while V72 can succeed when only physical geometry is controlled.

\subsection{What has and has not moved}

The abstract compact-sublevel input (71.5) is now populated by the
explicit orbit control (72.2).  The remaining analytic acquisition
is a uniform tail such as (72.8) or the fibre-integrated effective
bound (72.12).  Because (72.2) is defined after optimizing over the
whole orbit, proving that tail from a Faddeev--Popov chart requires
overlap and multiplicity control; a local slice estimate alone still
does not provide it.

\begin{assessment}[Status after V72]
\label{ass:v12v72-status}
The strong \(C^{2,\alpha}\) compact screen of V69 has been converted
into an invariant best-representative Lyapunov function.  Its
sublevels are compact in the physical \(C^{2,\beta}\) orbit space,
and either a uniform moment or a fibre-integrated effective-action
bound yields quotient tightness with an explicit tail.  No global
choice of representatives enters the proof.
\end{assessment}

\begin{firewall}[Claim boundary after V72]
\label{fire:v12v72-claim}
The metric quotient hypotheses, Borel measurability of (72.2), and a
uniform physical tail (72.8) or (72.12) have not been proved for the
full diffeomorphism group and RPESM action.  The orbit infimum does
not supply Faddeev--Popov determinants, Gribov multiplicities, or
stabiliser weights.  It is a compactness device, not a construction
of the gravitational measure.  The full SM--Einstein continuum is
not yet constructed.
\end{firewall}

\subsection{Parent record}

This V72 note appends to the validated V71 file with record
\[
\begin{array}{c|c}
\text{lines}&20406\\
\text{bytes}&680604\\
\text{SHA--256}&
\texttt{B8EBD5C46B8F234389531B93F5B5DFACE190D2A27FCCC8E36E5A78857BE8DE08}.
\end{array}
\]

% =============================================================================
% END APPENDED TWELFTH-CYCLE V72 MATERIAL
% =============================================================================

% =============================================================================
% BEGIN APPENDED TWELFTH-CYCLE V73 MATERIAL
% =============================================================================

\section{Twelfth-cycle V73: measurable orbit infima through a
separable Sobolev screen}
\label{sec:v12v73}

V72 assumes that its best-representative Hölder control is Borel.
The usual big Hölder spaces are nonseparable and therefore are not the cleanest setting
for proving that assertion.  This note replaces the strong screen by
a separable Sobolev space compactly embedded in the same weak metric
topology.  The resulting extended norm is Borel.  For a Polish
diffeomorphism group with a Borel pullback action, the orbit infimum
is universally measurable; under a natural orbitwise continuity
condition it is Borel.

\subsection{Correction of the ambient Polish topology}

In V69--V72, every use of a Polish two-derivative Hölder background
space is to be read in the little Hölder space
\[
 c^{2,\beta}
 :=\overline{C^\infty}^{\,C^{2,\beta}},
                                                               \tag{73.1}
\]
or in another explicitly specified separable Polish topology with
the same compact embeddings.  The ordinary big Hölder space
\(C^{2,\beta}\) is generally nonseparable, so the unqualified claim
in V69 that it is Polish is superseded here.  All compactness proofs
in V69--V72 remain valid after this replacement because bounded
\(C^{2,\alpha}\) sets with \(\alpha>\beta\) are relatively compact
in \(c^{2,\beta}\).

\subsection{A Borel strong control}

Let \(d=\dim M\), fix \(0<\beta<1\), and choose
\begin{equation}
                         s>{d\over2}+2+\beta.                 \tag{73.2}
\end{equation}
Let
\[
 E:=c^{2,\beta}(M;\operatorname{Sym}^2T^*M),\qquad
 H:=H^s(M;\operatorname{Sym}^2T^*M).
\]
The Sobolev embedding \(J:H\to E\) is continuous, injective, and
compact.  Both \(H\) and \(E\) are separable Banach spaces.

Fix \(0<\lambda<1\).  On
\(X=\operatorname{Met}^{2,\beta}(M)\subset E\), define
\begin{equation}
 {\cal N}_{\lambda,s}(h)
 :=
 \begin{cases}
  \|J^{-1}h\|_{H^s},&
  h\in J(H),\
  \lambda g_0\preceq h\preceq\lambda^{-1}g_0,\\
  +\infty,&\text{otherwise}.
 \end{cases}                                                \tag{73.3}
\end{equation}

\begin{lemma}[Borel Sobolev screen]
\label{lem:v12v73-Borel-screen}
The function \({\cal N}_{\lambda,s}:X\to[0,\infty]\) is Borel.
For every \(R<\infty\), its sublevel set is relatively compact in
\(X\) with the \(c^{2,\beta}\) topology.
\end{lemma}

\begin{proof}
The space \(H\) is Polish and \(J:H\to E\) is a continuous
injection into a Polish space.  The Lusin--Souslin theorem implies
that \(J(H)\) is Borel in \(E\) and that
\(J^{-1}:J(H)\to H\) is Borel.  Hence
\(h\mapsto\|J^{-1}h\|_{H^s}\) is Borel on \(J(H)\).

The pointwise inequalities
\(\lambda g_0\preceq h\preceq\lambda^{-1}g_0\) define a closed
subset of \(E\), because \(c^{2,\beta}\) convergence implies uniform
convergence.  Extending the Sobolev norm by \(+\infty\) outside the
intersection therefore gives a Borel function, proving the first
claim.

A sublevel of (73.3) is bounded in \(H^s\).  Compactness of \(J\)
gives a convergent subsequence in \(E\) from every sequence in that
sublevel.  The uniform nondegeneracy inequalities are closed, so the
limit remains a metric in \(X\).  This proves relative compactness.
\end{proof}

\subsection{Projection measurability of the orbit infimum}

Let \({\cal D}\) be a Polish group and assume the pullback action
\[
                         {\cal D}\times X\longrightarrow X,
 \qquad(\psi,g)\longmapsto\psi^*g                            \tag{73.4}
\]
is Borel.  Define
\begin{equation}
 {\cal L}_{\lambda,s}(g)
 :=\inf_{\psi\in{\cal D}}
 {\cal N}_{\lambda,s}(\psi^*g).                              \tag{73.5}
\end{equation}

Recall that a real-valued function on a Polish space is upper
semianalytic when every strict sublevel set is analytic.  Such a
function is measurable for the completion of every Borel
probability measure.

\begin{theorem}[Universal measurability of best-representative
control]
\label{thm:v12v73-universal}
Under (73.4), \({\cal L}_{\lambda,s}\) is orbit invariant and upper
semianalytic.  In particular it is universally measurable.  If
\(Y=X/{\cal D}\) is a Polish metric quotient as in V71, the induced
function
\[
 \overline{\cal L}_{\lambda,s}(\pi g)
 :={\cal L}_{\lambda,s}(g)                                  \tag{73.6}
\]
is also upper semianalytic and universally measurable on \(Y\).
\end{theorem}

\begin{proof}
Orbit invariance is the change-of-variable argument of Lemma
\ref{lem:v12v72-invariance}.  By Lemma
\ref{lem:v12v73-Borel-screen} and the Borel action, the function
\[
 F(\psi,g):={\cal N}_{\lambda,s}(\psi^*g)                    \tag{73.7}
\]
is Borel on \({\cal D}\times X\).  For \(a\in\mathbb R\),
\begin{equation}
 \{g:{\cal L}_{\lambda,s}(g)<a\}
 =
 \operatorname{proj}_X\{(\psi,g):F(\psi,g)<a\}.              \tag{73.8}
\end{equation}
The set on the right before projection is Borel, and its projection
is analytic.  Thus \({\cal L}_{\lambda,s}\) is upper semianalytic
and hence universally measurable.

The strict sublevel of (73.5) is the image under the continuous
quotient map \(\pi\) of the analytic set (73.8).  A continuous image
of an analytic set is analytic.  This proves the quotient assertion.
\end{proof}

Universal measurability is enough for moment and tail estimates:
each finite-cutoff probability law is used with its completion.
No choice function \(g\mapsto\psi_g\) is being asserted.

\begin{proposition}[A Borel upgrade]
\label{prop:v12v73-Borel-upgrade}
Let \(\{\psi_k:k\geq1\}\) be dense in \({\cal D}\).  Suppose that for
every \(g\in X\), the function
\[
                  \psi\longmapsto
                  {\cal N}_{\lambda,s}(\psi^*g)              \tag{73.9}
\]
has open strict sublevel sets in \({\cal D}\); in particular, this
holds if (73.9) is upper semicontinuous.  Then
\begin{equation}
 {\cal L}_{\lambda,s}(g)
 =\inf_{k\geq1}{\cal N}_{\lambda,s}(\psi_k^*g),              \tag{73.10}
\end{equation}
and \({\cal L}_{\lambda,s}\) is Borel on \(X\).
\end{proposition}

\begin{proof}
The right side of (73.10) is at least the infimum over all of
\({\cal D}\).  Conversely, if
\({\cal N}_{\lambda,s}(\psi^*g)<a\), the strict sublevel containing
\(\psi\) is open and therefore contains some \(\psi_k\).  Hence the
countable infimum is below every \(a\) strictly above the full
infimum, proving equality.  Every term in the countable infimum is
Borel by (73.4) and Lemma \ref{lem:v12v73-Borel-screen}; a countable
infimum of Borel functions is Borel.
\end{proof}

The condition (73.9) is a sufficient upgrade, not a hidden
regularity fact about the full diffeomorphism action.  It must be
verified in the chosen Sobolev diffeomorphism category.

\subsection{Compact quotient sublevels}

\begin{theorem}[Measurable Sobolev orbit screen]
\label{thm:v12v73-compact-screen}
Assume the Polish metric quotient hypotheses of V71.  Then for every
\(R<\infty\), the closure in \(Y\) of
\begin{equation}
 \{\pi g:{\cal L}_{\lambda,s}(g)\leq R\}                    \tag{73.11}
\end{equation}
is compact.
\end{theorem}

\begin{proof}
Choose a sequence \(y_j=\pi g_j\) in (73.11).  By the definition of
the infimum, select \(\psi_j\in{\cal D}\) with
\[
 h_j:=\psi_j^*g_j,\qquad
 {\cal N}_{\lambda,s}(h_j)\leq R+j^{-1}.                    \tag{73.12}
\]
Lemma \ref{lem:v12v73-Borel-screen} gives a subsequence converging in
\(c^{2,\beta}\) to a uniformly nondegenerate metric \(h\).
Continuity of the quotient map gives
\[
                             y_j=\pi h_j\longrightarrow\pi h.
\]
Thus (73.11) is relatively sequentially compact in the metric space
\(Y\), and its closure is compact.
\end{proof}

\begin{corollary}[Completed-law moment criterion]
\label{cor:v12v73-moment}
Let \(\mu_n\) be Borel probability measures on \(X\), completed so
that \({\cal L}_{\lambda,s}\) is measurable.  If
\begin{equation}
 \sup_n\int_X{\cal L}_{\lambda,s}(g)^p\,d\mu_n(g)
 \leq C_p<\infty                                             \tag{73.13}
\end{equation}
for some \(p>0\), then \((\pi_\#\mu_n)\) is tight on \(Y\) and
\begin{equation}
 (\pi_\#\mu_n)
 \{\overline{\cal L}_{\lambda,s}>R\}
 \leq C_pR^{-p}.                                            \tag{73.14}
\end{equation}
The fibre-first effective-action criterion of V70 has the identical
conclusion if its \(U\) is replaced by
\({\cal L}_{\lambda,s}^p\).
\end{corollary}

\begin{proof}
Universal measurability from Theorem
\ref{thm:v12v73-universal} makes the integrals and events measurable
in the completed laws.  Markov's inequality gives (73.14).
Theorem \ref{thm:v12v73-compact-screen} and the quotient tightness
criterion of V71 finish the proof.  The V70 assertion follows by
repeating (70.15) with \(U={\cal L}_{\lambda,s}^p\).
\end{proof}

\subsection{Why the measurability distinction matters}

Theorem \ref{thm:v12v73-universal} is sufficient for existence and
tightness under each completed cutoff law.  Borel measurability is
preferable when defining a common Markov kernel, a disintegration
simultaneously in external parameters, or a Borel conditional state
field.  Proposition \ref{prop:v12v73-Borel-upgrade} states one exact
route to that stronger property.  Neither form supplies the tail
estimate (73.12); it only makes that estimate a well-posed
acquisition target.

\begin{assessment}[Status after V73]
\label{ass:v12v73-status}
The measurability assumption left in V72 has been reduced to standard
descriptive-set hypotheses.  A separable high-Sobolev screen is
Borel and compact in the weak \(c^{2,\beta}\) topology.  Its
best-orbit infimum is universally measurable for every Polish Borel
action, has compact quotient sublevels, and is Borel under the
explicit orbitwise open-sublevel condition (73.9).
\end{assessment}

\begin{firewall}[Claim boundary after V73]
\label{fire:v12v73-claim}
No specific Sobolev diffeomorphism group, global Polish quotient,
orbitwise continuity property (73.9), or uniform moment (73.12) has
been established for the full RPESM regulator.  Universal
measurability does not resolve Gribov multiplicities, stabilisers,
partition floors, conformal positivity, or reflection positivity.
The full SM--Einstein continuum is not yet constructed.
\end{firewall}

\subsection{Parent record}

This V73 note appends to the validated V72 file with record
\[
\begin{array}{c|c}
\text{lines}&20674\\
\text{bytes}&690055\\
\text{SHA--256}&
\texttt{CF9205DDC01F59FF0DA70A4B62A1E716455E9EA2FC5F33FC1C71A97985BD2664}.
\end{array}
\]

% =============================================================================
% END APPENDED TWELFTH-CYCLE V73 MATERIAL
% =============================================================================



% =============================================================================
% BEGIN APPENDED TWELFTH-CYCLE V74 MATERIAL
% =============================================================================

\section{Twelfth-cycle V74: a compact gauge-invariant observable
continuum without a global orbit space}
\label{sec:v12v74}

V71--V73 construct a geometric route through a Polish
diffeomorphism quotient.  A global quotient can fail to be Hausdorff
or can contain singular orbit strata.  This note gives a weaker but
unconditional compactness route for positive laws: retain a
countable family of bounded gauge-invariant observables, embed their
joint values in the Hilbert cube, and take the closure.  This
constructs a subsequential continuum for those cylinders without
claiming that every limiting point is a classical smooth geometry.

\subsection{The physical observable compactification}

Let \(X\) be the finite-regularity configuration space at all
cutoffs, embedded in one measurable ambient set, and let
\({\cal D}\) denote the physical gauge and diffeomorphism action.
Choose bounded continuous invariant observables
\[
 f_k:X\longrightarrow[-1,1],\qquad
 f_k(\psi x)=f_k(x),\qquad k\geq1.                            \tag{74.1}
\]
Examples of candidates are bounded transforms of smeared curvature
writers, matter bilinears, and Wilson observables with data drawn
from countable dense test families.  Their continuity and invariance
must be verified in the selected regulator.

Let
\begin{align}
 {\mathbb Q}&:=[-1,1]^{\mathbb N},                           \tag{74.2}\\
 r(u,v)&:=\sum_{k=1}^{\infty}2^{-k-1}|u_k-v_k|,              \tag{74.3}\\
 \Phi(x)&:=(f_1(x),f_2(x),\ldots),                           \tag{74.4}\\
 K_{\rm phys}&:=\overline{\Phi(X)}^{\,{\mathbb Q}}.          \tag{74.5}
\end{align}
The metric satisfies \(r\leq1\).

\begin{lemma}[Compact physical cylinder space]
\label{lem:v12v74-compact}
The metric \(r\) induces the product topology on \(\mathbb Q\).
The spaces \(\mathbb Q\) and \(K_{\rm phys}\) are compact Polish
spaces, and \(\Phi:X\to K_{\rm phys}\) is continuous and invariant.
\end{lemma}

\begin{proof}
If \(r(u^{(n)},u)\to0\), then for every fixed \(k\),
\[
 |u_k^{(n)}-u_k|
 \leq2^{k+1}r(u^{(n)},u)\longrightarrow0,
\]
so metric convergence implies coordinatewise convergence.
Conversely, if all coordinates converge, choose \(N\) so that
\(\sum_{k>N}2^{-k}<\varepsilon/2\), and then make the first \(N\)
weighted coordinate differences sum to less than
\(\varepsilon/2\).  Thus coordinatewise convergence implies metric
convergence.

The countable product of compact intervals is compact in the product
topology, hence compact under \(r\).  It is separable because
sequences that are rational and eventually zero form a countable
dense set.  Compact metric spaces are complete, so \(\mathbb Q\) is
Polish.  The closed subset \(K_{\rm phys}\) has the same properties.
Continuity of all coordinates \(f_k\), together with the finite-tail
argument above, proves continuity of \(\Phi\).  Invariance follows
coordinatewise from (74.1).
\end{proof}

\subsection{Automatic subsequential cylinder convergence}

Let \(\mu_n\) be arbitrary Borel probability laws on \(X\) and put
\begin{equation}
                         \nu_n:=\Phi_\#\mu_n.                 \tag{74.6}
\end{equation}

\begin{theorem}[Positive observable subsequence]
\label{thm:v12v74-subsequence}
Every sequence \((\nu_n)\) has a weakly convergent subsequence
\(\nu_{n_j}\Rightarrow\nu_\infty\) on \(K_{\rm phys}\).  It may be
thinned further so that
\begin{align}
 {\mathsf W}_r(\nu_{n_j},\nu_\infty)&\leq2^{-j},              \tag{74.7}\\
 \sum_{j=1}^{\infty}
 {\mathsf W}_r(\nu_{n_{j+1}},\nu_{n_j})&\leq{3\over2}.        \tag{74.8}
\end{align}
Here \({\mathsf W}_r\) is the order-one transport distance with
bounded cost \(r\).
\end{theorem}

\begin{proof}
All probability measures on the compact metric space
\(K_{\rm phys}\) form a weakly sequentially compact family.
Therefore some subsequence converges weakly.  For a bounded metric,
\({\mathsf W}_r\) metrizes weak convergence.  Choose successive
indices so that (74.7) holds.  The triangle inequality gives
\[
 {\mathsf W}_r(\nu_{n_{j+1}},\nu_{n_j})
 \leq2^{-(j+1)}+2^{-j},
\]
whose sum is \(3/2\).
\end{proof}

\begin{corollary}[Continuum of bounded physical cylinders]
\label{cor:v12v74-cylinders}
For every \(F\in C(K_{\rm phys})\),
\begin{equation}
 \int_X F(\Phi(x))\,d\mu_{n_j}(x)
 \longrightarrow
 \int_{K_{\rm phys}}F(u)\,d\nu_\infty(u).                   \tag{74.9}
\end{equation}
In particular all finite polynomials in the coordinates \(f_k\)
converge along the same subsequence.
\end{corollary}

\begin{proof}
Equation (74.9) is the definition of weak convergence, and every
finite coordinate polynomial is continuous on \(K_{\rm phys}\).
\end{proof}

This result requires neither representative compactness nor a
Hausdorff set-theoretic orbit space.  It pays for that robustness by
allowing ideal points in
\(K_{\rm phys}\setminus\Phi(X)\).

\subsection{When the compactification represents physical orbits}

\begin{proposition}[Separation and absence of ideal points]
\label{prop:v12v74-separation}
Suppose the observables \(f_k\) separate \({\cal D}\)-orbits:
\begin{equation}
 \Phi(x)=\Phi(y)
 \quad\Longrightarrow\quad
 y=\psi x\text{ for some }\psi\in{\cal D}.                  \tag{74.10}
\end{equation}
Then the induced map from the set-theoretic orbit space into
\(K_{\rm phys}\) is injective.  If, in addition,
\begin{equation}
                              \Phi(X)\text{ is closed in }\mathbb Q,
                                                                    \tag{74.11}
\end{equation}
then \(K_{\rm phys}=\Phi(X)\), so every limiting point represents an
actual orbit in \(X\).
\end{proposition}

\begin{proof}
Invariant observables make \(\Phi\) constant on each orbit, while
(74.10) says two distinct orbits cannot have the same value.  This
proves injectivity.  Equation (74.11) and the definition (74.5) give
\(K_{\rm phys}=\Phi(X)\).
\end{proof}

Condition (74.11) is a no-escape statement and is generally as hard
as geometric compact containment.  Without it, V74 constructs a
continuum state on the chosen observable algebra, not a probability
measure on smooth metrics and fields.

\subsection{Closed positivity and identity gates}

Assume there is a continuous reflection
\(\Theta:K_{\rm phys}\to K_{\rm phys}\) and a subalgebra
\({\cal A}_+\subset C(K_{\rm phys})\) of positive-time cylinders.

\begin{theorem}[Observable reflection positivity is weakly closed]
\label{thm:v12v74-OS}
Suppose for every \(F\in{\cal A}_+\),
\begin{equation}
 \int_{K_{\rm phys}}
 \overline{F(\Theta u)}F(u)\,d\nu_{n_j}(u)
 \geq-\epsilon_j(F),\qquad
 \epsilon_j(F)\longrightarrow0.                             \tag{74.12}
\end{equation}
Then
\begin{equation}
 \int_{K_{\rm phys}}
 \overline{F(\Theta u)}F(u)\,d\nu_\infty(u)\geq0             \tag{74.13}
\end{equation}
for every \(F\in{\cal A}_+\).
\end{theorem}

\begin{proof}
The integrand in (74.12) is bounded and continuous because \(F\)
and \(\Theta\) are continuous.  Weak convergence passes its integral
to the limit.  Let \(j\to\infty\) in (74.12).
\end{proof}

\begin{proposition}[Bounded Ward residuals pass]
\label{prop:v12v74-Ward}
Let \(D_m:C(K_{\rm phys})\to C(K_{\rm phys})\) be specified Ward
writers.  If for every cylinder \(F\) in their common domain,
\begin{equation}
 \left|\int D_mF\,d\nu_{n_j}\right|
 \leq\varepsilon_{j,m}(F),\qquad
 \varepsilon_{j,m}(F)\longrightarrow0,                      \tag{74.14}
\end{equation}
then
\begin{equation}
                         \int D_mF\,d\nu_\infty=0.            \tag{74.15}
\end{equation}
\end{proposition}

\begin{proof}
The function \(D_mF\) is bounded and continuous by hypothesis.
Weak convergence passes the integral, and the residual tends to
zero.
\end{proof}

These are observable versions of V67.  They do not prove that the
required reflection or Ward writers descend from the full regulated
SM--Einstein system.

\subsection{Handoff to the conditional quantum state}

Suppose transported conditional states factor through the physical
coordinates:
\[
                    \rho_n(x)=\overline\rho_n(\Phi(x)),       \tag{74.16}
\]
where
\(\overline\rho_n:K_{\rm phys}\to{\cal T}_1({\cal H}_*)\)
has one uniform Hölder constant with respect to \(r\).  If, along the
subsequence in Theorem \ref{thm:v12v74-subsequence},
\begin{equation}
 \sum_j\sup_{u\in K_{\rm phys}}
 \|\overline\rho_{n_{j+1}}(u)
   -\overline\rho_{n_j}(u)\|_1<\infty,                       \tag{74.17}
\end{equation}
then (74.8), (74.17), and V66 construct the joint
classical--quantum state on \(K_{\rm phys}\).  V67 and Theorem
\ref{thm:v12v74-OS} pass any populated reflection, Ward, and symmetry
gates.

\begin{assessment}[Status after V74]
\label{ass:v12v74-status}
Every positive cutoff family now has a subsequential continuum on
the compact space generated by any countable set of bounded
gauge-invariant observables, with a summable transport thinning.
Reflection positivity and bounded Ward identities are closed under
that limit.  This route bypasses a global gauge section and a
Hausdorff geometric orbit space at the cylinder level.
\end{assessment}

\begin{firewall}[Claim boundary after V74]
\label{fire:v12v74-claim}
No countable family has been proved to separate all physical
SM--Einstein orbits, no-escape condition (74.11) is unproved, and the
limit may live on ideal generalized configurations or collapse to a
trivial law.  Unbounded stress tensors require uniform
integrability.  Complex determinant amplitudes require the
independent denominator and phase controls of V68 and cannot be
treated as the positive laws (74.6).  The reflection, Ward, and
conditional-state hypotheses have not been populated for the full
model.  The full SM--Einstein continuum is not yet constructed.
\end{firewall}

\subsection{Parent record}

This V74 note appends to the validated V73 file with record
\[
\begin{array}{c|c}
\text{lines}&20959\\
\text{bytes}&700878\\
\text{SHA--256}&
\texttt{694C3D0AD0FC6F66E0B9349282BE978C157144A0D07DF89468F10FC570872BC1}.
\end{array}
\]

% =============================================================================
% END APPENDED TWELFTH-CYCLE V74 MATERIAL
% =============================================================================

% =============================================================================
% BEGIN APPENDED TWELFTH-CYCLE V75 MATERIAL
% =============================================================================

\section{Twelfth-cycle V75: mandatory companion audit}
\label{sec:v12v75}

This multiple-of-five checkpoint inspects the current Yang--Mills and
Navier--Stokes companions before continuing the observable-space
construction of V74.

\subsection{Live companion records}

The files inspected read-only are
\[
\begin{array}{c|r|r|c}
\text{record}&\text{lines}&\text{bytes}&\text{SHA--256}\\ \hline
\text{Yang--Mills V59}&22075&750408&
\texttt{AACEF6B54245542726D0D597400DC0C0A42EA6025F04C66244BDEE606831399A}\\
\text{Navier--Stokes V26}&5319&278283&
\texttt{82C887F8C2C69E4F28FAD7C3DC8DE5B9099CB5A4BF431EBA1FFF3E91935978AD}.
\end{array}                                                   \tag{75.1}
\]
The Yang--Mills note has advanced from V56 at the V70 audit to V59.
The Navier--Stokes note is unchanged.

\subsection{New Yang--Mills information}

Yang--Mills V57 replaces one monolithic endpoint calculation by
\(384\) exact ownerwise Laurent lifts with restartable certificates.
V58 combines the completion and endpoint Wilson parts into the full
endpoint quartic arbitrary-momentum oracle.  V59 constructs exact
Laurent cubic vertices through external order two, proves their
orientation parity coefficientwise, and thereby completes all
one-loop numerator oracles.  The inverse-Hessian jet contraction and
the signed Fourier core remain open there.

The relevant transferable statement is informational.  Equality on
all sixteen parity corners did not determine the arbitrary-momentum
Laurent function; exponent data had to be retained before
evaluation.  Therefore a finite or merely convenient list of
SM--Einstein cylinders cannot be declared physically complete
because it agrees on the cutoffs already sampled.  Separation of
physical configurations and determination of probability laws must
be proved for the entire chosen observable family.

The ownerwise streaming mechanism is computational rather than a new
SM theorem.  It nevertheless confirms the sound acquisition order:
construct exact observables on the full regulated domain, prove their
invariance and continuity, and only then use their values as
coordinates of the V74 compactification.

\subsection{Navier--Stokes comparison}

Navier--Stokes V26 still shows that a reduced norm is legal only when
the input lies in the exact structured class on which that norm was
computed.  No new metric, measure, reflection, or gauge result is
available to import.  Applied as a proof rule, it agrees with the
Yang--Mills conclusion: an observable compactification is physically
faithful only on the class it is proved to separate.

\begin{decision}[Direction after the V75 audit]
\label{dec:v12v75-next}
Retain V69--V74, including the little-Hölder correction of V73.
At V76, assume the physical orbit space of V71 is Polish and use a
countable dense set in its bounded metric to construct distance
coordinates.  Prove that these coordinates:
\begin{enumerate}
\item separate physical orbits and recover the quotient topology;
\item generate its Borel sigma-algebra and determine probability
      laws;
\item make the V74 polynomial cylinders dense on the compact
      observable closure;
\item carry no limiting mass on ideal boundary points whenever the
      V71--V73 quotient-tightness stock is available.
\end{enumerate}
This will identify precisely how the geometric and observable routes
fit together.
\end{decision}

\begin{assessment}[Status after V75]
\label{ass:v12v75-status}
The mandatory audit is complete.  YM V57--V59 supplies no
gravitational estimate, but it rules out treating finite sampled
values as a complete observable.  The next SM step is therefore a
proof of topological and measure separation by a countable invariant
family, followed by a no-ideal-mass handoff from geometric tightness.
\end{assessment}

\begin{firewall}[Claim boundary after V75]
\label{fire:v12v75-claim}
No Yang--Mills Laurent coefficient is imported into the gravitational
sector, and no finite set of observables is claimed to separate
SM--Einstein configurations.  The physical orbit space, its metric,
and the required exact invariant writers remain conditional inputs.
The full SM--Einstein continuum is not yet constructed.
\end{firewall}

\subsection{Parent record}

This V75 note appends to the validated V74 file with record
\[
\begin{array}{c|c}
\text{lines}&21232\\
\text{bytes}&711243\\
\text{SHA--256}&
\texttt{8A25C8C76C6AAC807F1F3A7F7AA4EAE2D99FAB2B06834DDBD22DB74ED23700D2}.
\end{array}
\]

% =============================================================================
% END APPENDED TWELFTH-CYCLE V75 MATERIAL
% =============================================================================

% =============================================================================
% BEGIN APPENDED TWELFTH-CYCLE V76 MATERIAL
% =============================================================================

\section{Twelfth-cycle V76: a countable orbit-separating algebra and
the no-ideal-mass theorem}
\label{sec:v12v76}

V74 constructs a compact continuum for any countable bounded
invariant observable family, but leaves separation and ideal boundary
points open.  Under the Polish physical-quotient hypothesis of V71,
this note supplies a canonical separating family from the quotient
metric itself.  It also proves that the quotient-tightness estimates
of V71--V73 prevent every subsequential compactified limit from
charging ideal points.

\subsection{Distance coordinates reconstruct the quotient metric}

Let \((Y,\delta)\) be the bounded Polish physical orbit space of V71,
with \(\delta\leq1\), and choose a countable dense set
\[
                         \{a_k:k\geq1\}\subset Y.
\]
Define
\begin{align}
 f_k(y)&:=\delta(y,a_k)\in[0,1],                              \tag{76.1}\\
 \iota(y)&:=(f_1(y),f_2(y),\ldots)\in[0,1]^{\mathbb N},       \tag{76.2}\\
 K_\delta&:=\overline{\iota(Y)}^{\,[0,1]^{\mathbb N}}.       \tag{76.3}
\end{align}
Equip the product with the V74 metric
\[
 r(u,v):=\sum_{k=1}^{\infty}2^{-k-1}|u_k-v_k|.               \tag{76.4}
\]

\begin{theorem}[Exact reconstruction by countably many distances]
\label{thm:v12v76-reconstruction}
For all \(y,z\in Y\),
\begin{equation}
 \boxed{\displaystyle
 \delta(y,z)=\sup_{k\geq1}
 |f_k(y)-f_k(z)|.}                                          \tag{76.5}
\end{equation}
Consequently \(\iota\) is injective.  It is a homeomorphism from
\(Y\) onto \(\iota(Y)\) with the subspace topology inherited from
\(K_\delta\).
\end{theorem}

\begin{proof}
The triangle inequality gives
\[
 |f_k(y)-f_k(z)|
 =|\delta(y,a_k)-\delta(z,a_k)|
 \leq\delta(y,z),
\]
so the right side of (76.5) is at most the left side.  Choose
\(a_{k_j}\to y\).  Then
\[
 f_{k_j}(y)\longrightarrow0,\qquad
 f_{k_j}(z)\longrightarrow\delta(z,y)
\]
by continuity of the distance.  The corresponding absolute
differences tend to \(\delta(y,z)\), proving (76.5) and injectivity.

If \(y_n\to y\), every \(f_k(y_n)\to f_k(y)\), so
\(\iota(y_n)\to\iota(y)\) in the product topology and hence under
\(r\).  Conversely, suppose
\(\iota(y_n)\to\iota(y)\).  Given \(\varepsilon>0\), choose \(k\)
with \(\delta(y,a_k)<\varepsilon/4\).  For all sufficiently large
\(n\),
\[
 |\delta(y_n,a_k)-\delta(y,a_k)|<\varepsilon/2.
\]
Therefore
\[
 \delta(y_n,y)
 \leq\delta(y_n,a_k)+\delta(a_k,y)
 <2\delta(y,a_k)+{\varepsilon\over2}<\varepsilon.
\]
Thus \(y_n\to y\), proving continuity of the inverse on
\(\iota(Y)\).
\end{proof}

The embedding is also Lipschitz in the forward direction:
\begin{equation}
 r(\iota(y),\iota(z))
 \leq {1\over2}\delta(y,z),                                 \tag{76.6}
\end{equation}
because every coordinate difference is at most \(\delta(y,z)\) and
\(\sum_{k\geq1}2^{-k-1}=1/2\).

\subsection{Pullback to gauge-invariant observables}

Let \(\pi:X\to Y\) be the physical quotient map.  Put
\begin{equation}
                         F_k(x):=f_k(\pi x).                 \tag{76.7}
\end{equation}
Each \(F_k\) is bounded, continuous, and gauge invariant.  Moreover,
\begin{equation}
 \pi x=\pi x'
 \quad\Longleftrightarrow\quad
 F_k(x)=F_k(x')\ \text{for every }k.                        \tag{76.8}
\end{equation}
Thus the countable family \((F_k)\) populates the separation
hypothesis (74.10), conditional only on the existence of
\((Y,\delta)\).

\begin{corollary}[Countable recovery of physical convergence]
\label{cor:v12v76-convergence}
For orbit classes \(y_n,y\in Y\),
\begin{equation}
 y_n\longrightarrow y
 \quad\Longleftrightarrow\quad
 f_k(y_n)\longrightarrow f_k(y)
 \text{ for every }k.                                      \tag{76.9}
\end{equation}
Equivalently, convergence of all the invariant writers \(F_k\)
recovers the physical quotient topology.
\end{corollary}

\begin{proof}
This is the homeomorphism assertion in Theorem
\ref{thm:v12v76-reconstruction}.
\end{proof}

\subsection{Borel and probability determination}

\begin{theorem}[The distance algebra determines physical laws]
\label{thm:v12v76-law-determination}
The image \(\iota(Y)\) is Borel in \(K_\delta\), and
\(\iota^{-1}:\iota(Y)\to Y\) is Borel.  Hence
\begin{equation}
 {\cal B}(Y)=\sigma(f_1,f_2,\ldots).                         \tag{76.10}
\end{equation}
If two Borel probability measures \(\mu,\nu\) on \(Y\) satisfy
\begin{equation}
 \int_Y P(f_{k_1},\ldots,f_{k_m})\,d\mu
 =
 \int_Y P(f_{k_1},\ldots,f_{k_m})\,d\nu                     \tag{76.11}
\end{equation}
for every finite list of indices and every real polynomial \(P\),
then \(\mu=\nu\).
\end{theorem}

\begin{proof}
The map \(\iota\) is a continuous injection from the Polish space
\(Y\) into the Polish space \(K_\delta\).  The Lusin--Souslin theorem
makes \(\iota(Y)\) Borel and the inverse Borel.  Since the product
Borel sigma-algebra is generated by its coordinates, pulling it back
gives (76.10).

Let \({\cal A}\) be the unital real algebra on \(K_\delta\) generated
by the coordinate projections.  It separates points because two
different sequences differ in some coordinate.  The Stone--Weierstrass
theorem makes \({\cal A}\) uniformly dense in
\(C(K_\delta)\).  Equation (76.11) says that
\(\iota_\#\mu\) and \(\iota_\#\nu\) agree on \({\cal A}\); uniform
approximation shows that they agree on every continuous function,
hence are equal.  Applying the Borel inverse of \(\iota\) gives
\(\mu=\nu\).
\end{proof}

Thus one common subsequence of all finite distance-coordinate
polynomials determines at most one physical probability law.  This
is stronger than agreement on any fixed finite sampling screen.

\subsection{Geometric tightness excludes ideal boundary mass}

The ideal boundary of the observable compactification is
\begin{equation}
                         \partial_{\rm id}K_\delta
 :=K_\delta\setminus\iota(Y).                               \tag{76.12}
\end{equation}
It is Borel because \(\iota(Y)\) is Borel, but it need not be empty.

\begin{theorem}[No ideal mass under quotient tightness]
\label{thm:v12v76-no-ideal}
Let \((\mu_n)\) be a uniformly tight sequence of Borel probability
measures on \(Y\).  Suppose along a subsequence
\begin{equation}
                         \iota_\#\mu_n\Rightarrow\eta
                         \quad\text{on }K_\delta.            \tag{76.13}
\end{equation}
Then
\begin{equation}
                  \eta(\iota(Y))=1,\qquad
                  \eta(\partial_{\rm id}K_\delta)=0.         \tag{76.14}
\end{equation}
There is a unique probability measure \(\mu\) on \(Y\) with
\(\iota_\#\mu=\eta\), and in fact
\begin{equation}
                              \mu_n\Rightarrow\mu
                              \quad\text{on }Y.              \tag{76.15}
\end{equation}
\end{theorem}

\begin{proof}
Given \(\varepsilon>0\), uniform tightness gives a compact
\(C_\varepsilon\subset Y\) with
\[
                         \inf_n\mu_n(C_\varepsilon)
                         \geq1-\varepsilon.
\]
The set \(\iota(C_\varepsilon)\) is compact and therefore closed in
\(K_\delta\).  By the closed-set half of the Portmanteau theorem,
\[
 \eta(\iota(C_\varepsilon))
 \geq\limsup_n
 (\iota_\#\mu_n)(\iota(C_\varepsilon))
 \geq1-\varepsilon.                                        \tag{76.16}
\]
Since \(\iota(C_\varepsilon)\subset\iota(Y)\) and
\(\iota(Y)\) is Borel, let \(\varepsilon\downarrow0\) to obtain
(76.14).

Define
\[
 \mu:=(\iota^{-1})_\#\eta
\]
on the full-measure Borel set \(\iota(Y)\).  This is the unique law
with pushforward \(\eta\), by Theorem
\ref{thm:v12v76-law-determination}.

To prove (76.15), take an arbitrary subsequence of \((\mu_n)\).
Uniform tightness and Prokhorov's theorem give a further subsequence
converging weakly on \(Y\), say to \(\mu'\).  Continuity of \(\iota\)
makes its pushforwards converge to \(\iota_\#\mu'\).  They also
converge to \(\eta\) by (76.13), so uniqueness of weak limits on
\(K_\delta\) gives
\(\iota_\#\mu'=\eta=\iota_\#\mu\).  Theorem
\ref{thm:v12v76-law-determination} gives \(\mu'=\mu\).
Every subsequence has a further subsequence with this same limit;
therefore the full sequence in (76.15) converges to \(\mu\).
\end{proof}

\begin{corollary}[Ideal mass diagnoses geometric escape]
\label{cor:v12v76-escape}
If a compactified subsequential limit in (76.13) satisfies
\(\eta(\partial_{\rm id}K_\delta)>0\), then the corresponding
physical laws are not uniformly tight on \(Y\).
\end{corollary}

\begin{proof}
This is the contrapositive of Theorem
\ref{thm:v12v76-no-ideal}.
\end{proof}

\subsection{Transport compatibility}

For probability measures \(\mu,\nu\) on \(Y\), equation (76.6)
implies
\begin{equation}
 {\mathsf W}_r(\iota_\#\mu,\iota_\#\nu)
 \leq{1\over2}{\mathsf W}_\delta(\mu,\nu).                   \tag{76.17}
\end{equation}

\begin{proof}
Push any coupling of \(\mu,\nu\) forward by
\((\iota,\iota)\), apply (76.6) pointwise, and take the infimum.
\end{proof}

Hence the summable quotient-Wasserstein budgets of V69 and V71 pass
automatically to the canonical observable compactification.  The
reverse estimate need not hold for the weighted metric \(r\);
topological recovery is supplied instead by (76.5) and tightness.

\subsection{Combined geometric and observable handoff}

If either the physical moment criterion (73.12) or the fibre-first
effective tail of V70 holds, then V71--V73 give uniform tightness on
\(Y\).  The V74 compactified subsequence constructed from the
canonical coordinates (76.7) therefore has no ideal mass by Theorem
\ref{thm:v12v76-no-ideal}.  Its bounded cylinder correlations
determine a unique physical law on \(Y\) by Theorem
\ref{thm:v12v76-law-determination}.  Any populated OS and Ward gates
then pass by V74, while the conditional quantum-state handoff remains
the separate V66 summability requirement.

\begin{assessment}[Status after V76]
\label{ass:v12v76-status}
Conditional on a bounded Polish physical quotient, a countable
canonical family of bounded invariant observables now reconstructs
the quotient metric, topology, Borel structure, and probability
laws.  The geometric tightness route and the automatic observable
compactification have been joined: tightness forces every
compactified limit to represent an actual physical probability law,
while ideal boundary mass is a rigorous certificate of geometric
escape.
\end{assessment}

\begin{firewall}[Claim boundary after V76]
\label{fire:v12v76-claim}
The bounded Polish quotient metric remains unconstructed for the
full diffeomorphism action, so the distance writers (76.7) are not
yet explicit RPESM observables.  No effective conformal tail,
partition floor, nontriviality estimate, full-sequence uniqueness,
conditional-state stock, determinant phase, or reflection-positive
SM--Einstein kernel has been populated.  The full SM--Einstein
continuum is not yet constructed.
\end{firewall}

\subsection{Parent record}

This V76 note appends to the validated V75 file with record
\[
\begin{array}{c|c}
\text{lines}&21346\\
\text{bytes}&716124\\
\text{SHA--256}&
\texttt{84F0A7B16F419BC6AC3E2B92975858CB875D80AC6C135785E73F6F95303A3AB3}.
\end{array}
\]

% =============================================================================
% END APPENDED TWELFTH-CYCLE V76 MATERIAL
% =============================================================================

% =============================================================================
% BEGIN APPENDED TWELFTH-CYCLE V77 MATERIAL
% =============================================================================

\section{Twelfth-cycle V77: full-sequence convergence from a
countable observable ledger}
\label{sec:v12v77}

V74 supplies compactness and V76 supplies a countable
measure-determining physical algebra, but compactness alone selects
only subsequences.  V66 gives a stronger joint-state theorem when a
summable background Wasserstein stock is already known.  The result
below is complementary: on the compact observable space, convergence
or summable adjacent control of one countable ledger of expectations
constructs the full probability sequence, without presupposing a
coupling or a target law.

\subsection{A complete metric made from expectation values}

Let \(K\) be a compact metric space; in the application,
\(K=K_\delta\) from V76 or \(K=K_{\rm phys}\) from V74.  The Banach
space \(C(K;\mathbb R)\) is separable.  Choose a sequence
\[
 \{h_m:m\geq1\}
 \quad\text{dense in}\quad
 \{h\in C(K;\mathbb R):\|h\|_\infty\leq1\}.                  \tag{77.1}
\]
For Borel probability measures \(\mu,\nu\) on \(K\), define
\begin{equation}
 d_{\rm obs}(\mu,\nu)
 :=
 \sum_{m=1}^{\infty}2^{-m-1}
 \left|\int_Kh_m\,d\mu-\int_Kh_m\,d\nu\right|.               \tag{77.2}
\end{equation}

\begin{theorem}[Complete observable metric]
\label{thm:v12v77-metric}
The function \(d_{\rm obs}\) is a metric on
\({\cal P}(K)\), takes values in \([0,1]\), and induces the topology
of weak convergence.  The metric space
\(({\cal P}(K),d_{\rm obs})\) is compact and complete.
\end{theorem}

\begin{proof}
Each summand in (77.2) is nonnegative, symmetric, and satisfies the
triangle inequality.  The difference of two expectations is at most
\(2\), while
\(\sum_{m\geq1}2^{-m-1}=1/2\), so \(d_{\rm obs}\leq1\).

If \(d_{\rm obs}(\mu,\nu)=0\), the two measures agree on every
\(h_m\).  Given a real \(h\in C(K)\), after rescaling it into the
unit ball choose \(h_{m_j}\to h/\max\{1,\|h\|_\infty\}\)
uniformly.  The expectation functionals have norm one, so uniform
approximation shows that they agree on \(h\).  The Riesz
representation theorem gives \(\mu=\nu\).  Thus (77.2) is a metric.

If \(\mu_n\Rightarrow\mu\), every fixed expectation in (77.2)
converges.  A finite-head and uniformly bounded-tail argument gives
\(d_{\rm obs}(\mu_n,\mu)\to0\).  Conversely, convergence in
\(d_{\rm obs}\) gives convergence of every \(h_m\) expectation.
Uniform approximation as above gives convergence for every
\(h\in C(K)\), which is weak convergence.

The probability measures on compact \(K\) are weakly compact.
Since \(d_{\rm obs}\) induces that topology, they are compact under
\(d_{\rm obs}\).  Every compact metric space is complete.
\end{proof}

\subsection{A target-free full-sequence theorem}

Let \(\nu_n\in{\cal P}(K)\) be the pushforward of the \(n\)-th
regulated SM--Einstein law under a common observable map.  Put
\begin{equation}
                         e_n:=
 d_{\rm obs}(\nu_{n+1},\nu_n).                               \tag{77.3}
\end{equation}

\begin{theorem}[Summable observable-ledger construction]
\label{thm:v12v77-summable}
If
\begin{equation}
                              \sum_{n=1}^{\infty}e_n<\infty,  \tag{77.4}
\end{equation}
then there is a unique \(\nu_\infty\in{\cal P}(K)\) such that
\[
                         \nu_n\Rightarrow\nu_\infty,
\]
and
\begin{equation}
 d_{\rm obs}(\nu_N,\nu_\infty)
 \leq\sum_{n=N}^{\infty}e_n.                                \tag{77.5}
\end{equation}
For every \(m\),
\begin{equation}
 \left|\int h_m\,d\nu_N-\int h_m\,d\nu_\infty\right|
 \leq 2^{m+1}\sum_{n=N}^{\infty}e_n.                        \tag{77.6}
\end{equation}
\end{theorem}

\begin{proof}
The triangle inequality and (77.4) make \((\nu_n)\) Cauchy in
\(d_{\rm obs}\).  Completeness from Theorem
\ref{thm:v12v77-metric} gives a limit \(\nu_\infty\).
Letting the upper endpoint tend to infinity in
\[
 d_{\rm obs}(\nu_M,\nu_N)\leq
 \sum_{n=N}^{M-1}e_n
\]
proves (77.5).  The \(m\)-th nonnegative summand of (77.2) is at most
the whole metric, giving (77.6).  Uniqueness follows because
\(d_{\rm obs}\) is a metric.
\end{proof}

The stock \(e_n\) is observable and coupling-free.  It can be
estimated from a finite head plus a universal tail:
\begin{equation}
 e_n\leq
 \sum_{m=1}^{M}2^{-m-1}
 \left|\nu_{n+1}(h_m)-\nu_n(h_m)\right|
 +2^{-M}.                                                   \tag{77.7}
\end{equation}
The tail follows because each expectation difference is at most two.

\begin{warning}[A fixed ledger truncation is insufficient]
\label{warn:v12v77-truncation}
Using one fixed \(M\) in (77.7) leaves a constant error \(2^{-M}\)
at every refinement step, which is not summable.  A valid adjacent
certificate must choose \(M=M_n\to\infty\) with
\(\sum_n2^{-M_n}<\infty\), while also making the finite-head
expectation defects summable.  This is the observable analogue of
the exponent-retention lesson in the V75 Yang--Mills audit.
\end{warning}

\subsection{Convergence without an adjacent summability rate}

\begin{theorem}[Dense-ledger convergence criterion]
\label{thm:v12v77-dense}
Suppose that for every \(m\), the scalar sequence
\[
                              \int_Kh_m\,d\nu_n              \tag{77.8}
\]
converges.  Then the full sequence \((\nu_n)\) converges weakly to a
unique probability measure on \(K\).
\end{theorem}

\begin{proof}
Compactness of \({\cal P}(K)\) gives a weakly convergent subsequence,
say \(\nu_{n_j}\Rightarrow\nu\).  For each \(m\), its limiting
expectation is the prescribed full-sequence limit in (77.8).
Any other weakly convergent subsequence has a limit \(\nu'\) with the
same \(h_m\) expectations.  Density as in Theorem
\ref{thm:v12v77-metric} gives \(\nu'=\nu\).

If the full sequence did not converge to \(\nu\), some subsequence
would remain outside a fixed weak neighbourhood of \(\nu\).
Compactness would give that subsequence a further convergent
subsequence, whose limit must equal \(\nu\), a contradiction.
\end{proof}

For the canonical distance algebra of V76, one may take a countable
dense subalgebra of finite rational polynomials in the coordinate
functions, truncated and rescaled into the unit ball.  Theorem
\ref{thm:v12v76-law-determination} then gives the same uniqueness
directly.

\subsection{Return from the compactification to physical geometry}

Assume \(K=K_\delta\), with embedding
\(\iota:Y\to K_\delta\) from V76, and suppose the original physical
laws \(\mu_n\) on \(Y\) satisfy
\[
                              \nu_n=\iota_\#\mu_n.            \tag{77.9}
\]
If the quotient laws are uniformly tight by V71--V73, then the limit
\(\nu_\infty\) furnished by either Theorem
\ref{thm:v12v77-summable} or
\ref{thm:v12v77-dense} obeys
\[
                              \nu_\infty(\iota(Y))=1.         \tag{77.10}
\]
The no-ideal-mass theorem V76 therefore gives a unique
\(\mu_\infty\in{\cal P}(Y)\) and
\[
                              \mu_n\Rightarrow\mu_\infty.    \tag{77.11}
\]
Thus the observable ledger supplies full-sequence uniqueness, while
the geometric Lyapunov estimate supplies absence of escape.  Neither
one substitutes for the other.

\subsection{Handoff to OS, Ward, and conditional quantum data}

Once (77.11) holds, every bounded continuous physical cylinder has a
full-sequence limit.  The closed OS and bounded Ward theorems of V74
therefore apply without further subsequence extraction.  To upgrade
this scalar law to the joint state of V66, one still needs convergence
of the transported conditional density fields.  The observable stock
(77.4) controls the background law only; it does not bound the
trace-norm fibre stock \(a_n\) in (66.16).

\begin{assessment}[Status after V77]
\label{ass:v12v77-status}
A countable ledger of bounded physical expectations now gives two
full-sequence criteria on the compact observable space.  Summable
adjacent ledger distance yields an explicit convergence tail, while
mere convergence of every dense test expectation gives qualitative
uniqueness.  Combined with quotient tightness, the resulting limit
is a unique physical law rather than an ideal-boundary measure.
\end{assessment}

\begin{firewall}[Claim boundary after V77]
\label{fire:v12v77-claim}
No RPESM estimate has yet populated the adjacent stock (77.4) or
proved convergence of every ledger entry (77.8).  Quotient tightness,
nontriviality, conditional-state convergence, determinant phase,
unbounded stress insertions, and reflection positivity remain
separate gates.  The theorem is a full-sequence criterion, not a
proof that the physical regulator satisfies it.  The full
SM--Einstein continuum is not yet constructed.
\end{firewall}

\subsection{Parent record}

This V77 note appends to the validated V76 file with record
\[
\begin{array}{c|c}
\text{lines}&21662\\
\text{bytes}&727598\\
\text{SHA--256}&
\texttt{64DB8AEA615378B705B5F6AD64F4D4C7EAA66B817457E5EE8CE50676B1366EBE}.
\end{array}
\]

% =============================================================================
% END APPENDED TWELFTH-CYCLE V77 MATERIAL
% =============================================================================

% =============================================================================
% BEGIN APPENDED TWELFTH-CYCLE V78 MATERIAL
% =============================================================================

\section{Twelfth-cycle V78: quantitative noncollapse of the physical
probability law}
\label{sec:v12v78}

V69 notes that tightness may converge to a delta law, and V77 now
provides full-sequence uniqueness once an observable ledger
converges.  This note supplies the missing abstract noncollapse gate.
A connected two-point floor for one bounded physical observable, or
positive mass on two disjoint compact physical regions, passes to
the continuum and rules out a Dirac limit.  This is a probabilistic
nontriviality statement, separate from the determinant visibility
floor of V68.

\subsection{A variance floor survives weak convergence}

Let \(K\) be the compact physical observable space of V74 or V76,
and let \(\nu_n\Rightarrow\nu_\infty\) be the full sequence obtained
in V77.  For real \(h\in C(K)\), write
\[
 \operatorname{Var}_\nu(h)
 :=
 \int_Kh^2\,d\nu-\left(\int_Kh\,d\nu\right)^2.               \tag{78.1}
\]

\begin{theorem}[Connected-observable noncollapse]
\label{thm:v12v78-variance}
If there are a real \(h\in C(K)\) and \(v>0\) such that
\begin{equation}
                 \liminf_{n\to\infty}
                 \operatorname{Var}_{\nu_n}(h)\geq v,        \tag{78.2}
\end{equation}
then
\begin{equation}
                 \operatorname{Var}_{\nu_\infty}(h)\geq v.  \tag{78.3}
\end{equation}
In particular, \(\nu_\infty\) is not a Dirac measure.
\end{theorem}

\begin{proof}
Both \(h\) and \(h^2\) are bounded continuous functions on compact
\(K\).  Weak convergence therefore gives
\[
 \int h\,d\nu_n\longrightarrow\int h\,d\nu_\infty,\qquad
 \int h^2\,d\nu_n\longrightarrow\int h^2\,d\nu_\infty.
\]
Substitution in (78.1) shows that the variances converge.  Equation
(78.2) then gives (78.3).  Every observable has zero variance under
a Dirac measure, proving the last assertion.
\end{proof}

If \(h\) is \(L\)-Lipschitz with respect to a metric \(r\) on \(K\),
Popoviciu's inequality gives
\begin{equation}
 \operatorname{Var}_{\nu_\infty}(h)
 \leq {1\over4}
 \operatorname{osc}_{\operatorname{supp}\nu_\infty}(h)^2
 \leq {L^2\over4}
 \operatorname{diam}_r(\operatorname{supp}\nu_\infty)^2.    \tag{78.4}
\end{equation}
Consequently (78.3) implies
\begin{equation}
 \operatorname{diam}_r(\operatorname{supp}\nu_\infty)
 \geq {2\sqrt v\over L}.                                    \tag{78.5}
\end{equation}

\begin{proof}
A real random variable supported in an interval of length \(D\) has
variance at most \(D^2/4\).  Apply this to \(h\), then use the
Lipschitz bound on its oscillation.
\end{proof}

\subsection{Varying cutoff writers}

At finite cutoff the physical observable itself may depend on \(n\).
The next estimate makes the replacement error explicit.

\begin{lemma}[Stability of variance under uniform writer error]
\label{lem:v12v78-writer}
Let \(h,h_n:K\to\mathbb R\) satisfy
\[
 \|h\|_\infty\leq1,\qquad
 \|h_n\|_\infty\leq1,\qquad
 \|h_n-h\|_\infty\leq\epsilon_n.                            \tag{78.6}
\]
Then for every probability measure \(\nu\),
\begin{equation}
 \left|
 \operatorname{Var}_\nu(h_n)-\operatorname{Var}_\nu(h)
 \right|\leq4\epsilon_n.                                    \tag{78.7}
\end{equation}
\end{lemma}

\begin{proof}
Since both functions have norm at most one,
\[
 |h_n^2-h^2|\leq2|h_n-h|,
\]
so their second moments differ by at most \(2\epsilon_n\).  Their
means differ by at most \(\epsilon_n\), and two numbers in
\([-1,1]\) have squared values differing by at most twice their
difference.  Thus the squared means differ by at most
\(2\epsilon_n\).  Add the two bounds.
\end{proof}

\begin{corollary}[Cutoff response floor]
\label{cor:v12v78-cutoff}
Suppose \(\epsilon_n\to0\) in (78.6) and
\begin{equation}
                  \liminf_n
                  \operatorname{Var}_{\nu_n}(h_n)\geq v>0.  \tag{78.8}
\end{equation}
Then (78.3) holds for \(h\), and the continuum law is non-Dirac.
\end{corollary}

\begin{proof}
Lemma \ref{lem:v12v78-writer} transfers (78.8) to (78.2), and
Theorem \ref{thm:v12v78-variance} applies.
\end{proof}

For a bounded source \(t h_n\), define
\[
 Z_n(t):=\int_K e^{t h_n}\,d\nu_n.
\]
Direct differentiation gives
\begin{equation}
 \left.{d^2\over dt^2}\log Z_n(t)\right|_{t=0}
 =\operatorname{Var}_{\nu_n}(h_n).                          \tag{78.9}
\end{equation}
Thus (78.8) can be acquired as a uniform connected susceptibility
floor.  Equation (78.9) is only an identity; it does not produce the
floor.

\subsection{Two compact physical regions}

\begin{theorem}[Separated-region noncollapse]
\label{thm:v12v78-two-regions}
Let \(A,B\subset K\) be disjoint compact sets.  If constants
\(m_A,m_B>0\) satisfy
\begin{equation}
 \limsup_{n\to\infty}\nu_n(A)\geq m_A,\qquad
 \limsup_{n\to\infty}\nu_n(B)\geq m_B,                       \tag{78.10}
\end{equation}
then
\begin{equation}
 \nu_\infty(A)\geq m_A,\qquad
 \nu_\infty(B)\geq m_B.                                     \tag{78.11}
\end{equation}
In particular \(\nu_\infty\) is not a Dirac measure.
\end{theorem}

\begin{proof}
Compact subsets of the metric space \(K\) are closed.  The
closed-set Portmanteau inequality gives
\[
 \limsup_n\nu_n(A)\leq\nu_\infty(A),\qquad
 \limsup_n\nu_n(B)\leq\nu_\infty(B),
\]
which proves (78.11).  A Dirac measure cannot assign positive mass
to two disjoint closed sets.
\end{proof}

The use of fixed compact regions is essential.  Lower bounds on
moving sets can drift into the same limiting point, and lower bounds
on arbitrary open sets do not pass in the direction used in
(78.11).

\subsection{Return to the physical quotient}

In the canonical V76 setting, let
\(\iota:Y\to K_\delta\) and assume the quotient-tightness hypotheses
needed for the no-ideal-mass theorem.  The V77 limit is then
\(\nu_\infty=\iota_\#\mu_\infty\) for a unique law
\(\mu_\infty\) on \(Y\).  A variance floor for
\(h=f_k\), or for a finite polynomial in the distance coordinates,
shows that \(\mu_\infty\) is non-Dirac on physical orbit space.
Likewise, compact \(A,B\subset Y\) have compact disjoint images under
\(\iota\), so the two-region theorem applies to their images.

Neither criterion requires a choice of gauge representative.
They do require the observable or regions to be physically
invariant before the floor is estimated.

\subsection{Separation from other nonvanishing gates}

The conclusions above concern a positive probability law.  They do
not imply:
\begin{enumerate}
\item a nonzero complex determinant denominator or average-phase
      visibility floor;
\item nonzero connected correlations at arbitrarily large physical
      separation;
\item a particle pole, mass gap, scattering state, or interacting
      renormalised coupling;
\item uniqueness with respect to a change of regulator or
      universality class.
\end{enumerate}
Those statements require their own estimates.  Conversely, a
nonzero determinant phase anchor does not prevent the normalized
positive background law from becoming a delta measure.

\begin{assessment}[Status after V78]
\label{ass:v12v78-status}
The probabilistic nontriviality row left open in V69 is now an exact
gate.  A uniform connected variance floor for one convergent bounded
physical writer, with explicit writer-error tolerance, or positive
mass in two fixed disjoint compact physical regions passes to the
V77 full-sequence limit and rules out collapse to one orbit.
\end{assessment}

\begin{firewall}[Claim boundary after V78]
\label{fire:v12v78-claim}
No RPESM susceptibility floor, separated-region mass floor, or
uniformly convergent nonconstant physical writer has been populated.
Non-Dirac probability is weaker than interacting quantum dynamics,
clustering, a spectral particle interpretation, or universality.
The determinant and reflection-positive joint-state gates remain
separate.  The full SM--Einstein continuum is not yet constructed.
\end{firewall}

\subsection{Parent record}

This V78 note appends to the validated V77 file with record
\[
\begin{array}{c|c}
\text{lines}&21898\\
\text{bytes}&736819\\
\text{SHA--256}&
\texttt{C97C62760358C3433BF14B93BE7B660B65E82714448FB3763020A6C0C3890D25}.
\end{array}
\]

% =============================================================================
% END APPENDED TWELFTH-CYCLE V78 MATERIAL
% =============================================================================

% =============================================================================
% BEGIN APPENDED TWELFTH-CYCLE V79 MATERIAL
% =============================================================================

\section{Twelfth-cycle V79: unbounded Einstein writers and
uniformly integrable correlation limits}
\label{sec:v12v79}

V69 passes bounded continuous Einstein cylinders and leaves
unbounded curvature and stress insertions to a separate
uniform-integrability estimate.  V77 gives a full physical law once
its bounded ledger converges.  This note proves the missing handoff:
local uniform convergence of cutoff writers plus one
superlinear moment passes their expectations and, with the correct
Hölder exponents, their finite products.

\subsection{Varying continuous writers}

Let \(Y\) be the Polish physical space of V71 and suppose
\[
                              \mu_n\Rightarrow\mu             \tag{79.1}
\]
as obtained from V76--V77.  Let
\(T_n,T:Y\to\mathbb R\) be continuous and assume
\begin{equation}
 T_n\longrightarrow T
 \quad\text{uniformly on every compact subset of }Y.         \tag{79.2}
\end{equation}

\begin{lemma}[Varying bounded continuous mapping]
\label{lem:v12v79-bounded}
If \(\sup_n\|T_n\|_\infty\leq B<\infty\), then
\begin{equation}
                         \int_YT_n\,d\mu_n
                         \longrightarrow\int_YT\,d\mu.       \tag{79.3}
\end{equation}
\end{lemma}

\begin{proof}
Weak convergence makes \((\mu_n)\) uniformly tight.  Given
\(\varepsilon>0\), choose a compact \(C\subset Y\) such that
\[
 \mu(C)\geq1-\varepsilon,\qquad
 \inf_{n\geq1}\mu_n(C)\geq1-\varepsilon
\]
Extend the continuous restriction
\(T|_C\) to a bounded continuous function
\(\widetilde T\) on \(Y\), with
\(\|\widetilde T\|_\infty\leq B\), by the Tietze extension theorem.
Then
\[
 \begin{split}
 \left|\int T_n\,d\mu_n-\int T\,d\mu\right|
 &\leq
 \sup_C|T_n-T|
 +\left|\int\widetilde T\,d\mu_n
              -\int\widetilde T\,d\mu\right|\\
 &\quad+4B\varepsilon
 \end{split}                                                \tag{79.4}
\]
for all large \(n\).  The first term tends to zero by (79.2), and
the second by weak convergence.  Let \(\varepsilon\downarrow0\).
\end{proof}

The common bound \(B\) also bounds \(T\), because every singleton is
compact and (79.2) gives pointwise convergence.

\begin{theorem}[Superlinear moment handoff]
\label{thm:v12v79-moment}
Assume (79.1)--(79.2) and, for some \(p>1\),
\begin{equation}
                         \sup_n\int_Y|T_n|^p\,d\mu_n
                         \leq C_p<\infty.                    \tag{79.5}
\end{equation}
Then \(T\in L^p(\mu)\),
\begin{equation}
                         \int_Y|T|^p\,d\mu\leq C_p,           \tag{79.6}
\end{equation}
and
\begin{equation}
                         \int_YT_n\,d\mu_n
                         \longrightarrow\int_YT\,d\mu.       \tag{79.7}
\end{equation}
For every \(R>0\), both cutoff and limiting tails obey
\begin{align}
 \sup_n\int_{\{|T_n|>R\}}|T_n|\,d\mu_n
 &\leq C_pR^{1-p},                                          \tag{79.8}\\
 \int_{\{|T|>R\}}|T|\,d\mu
 &\leq C_pR^{1-p}.                                          \tag{79.9}
\end{align}
\end{theorem}

\begin{proof}
For a bounded continuous function
\(\chi:\mathbb R\to\mathbb R\), the functions
\(\chi\circ T_n\) are uniformly bounded and converge uniformly on
compact subsets to \(\chi\circ T\).  Lemma
\ref{lem:v12v79-bounded} gives
\[
                    (T_n)_\#\mu_n\Rightarrow T_\#\mu
                    \quad\text{on }\mathbb R.                \tag{79.10}
\]
Portmanteau applied to the nonnegative lower-semicontinuous function
\(|t|^p\) yields
\[
 \int|T|^p\,d\mu
 \leq\liminf_n\int|T_n|^p\,d\mu_n
 \leq C_p,
\]
which proves (79.6).

On \(\{|T_n|>R\}\),
\(|T_n|\leq R^{1-p}|T_n|^p\), giving (79.8); the same argument with
(79.6) gives (79.9).  Let
\[
 \tau_R(t):=\max\{-R,\min\{R,t\}\}.
\]
Lemma \ref{lem:v12v79-bounded} applied to
\(\tau_R\circ T_n\) gives convergence of the truncated expectations.
The difference between a writer and its truncation is bounded by its
tail first moment.  Equations (79.8)--(79.9), followed by
\(R\to\infty\), prove (79.7).
\end{proof}

\subsection{Finite products and connected correlations}

Let \(T_{i,n},T_i:Y\to\mathbb R\), \(1\leq i\leq m\), be continuous
and converge locally uniformly.  Suppose exponents \(p_i>1\) obey
\begin{equation}
 \sup_n\int|T_{i,n}|^{p_i}\,d\mu_n\leq C_i,\qquad
                         \theta:=\sum_{i=1}^m{1\over p_i}<1. \tag{79.11}
\end{equation}
Put \(r:=1/\theta>1\).

\begin{theorem}[Unbounded correlation handoff]
\label{thm:v12v79-products}
Under (79.11),
\begin{equation}
 \int_Y\prod_{i=1}^mT_{i,n}\,d\mu_n
 \longrightarrow
 \int_Y\prod_{i=1}^mT_i\,d\mu.                              \tag{79.12}
\end{equation}
The products have the uniform superlinear bound
\begin{equation}
 \sup_n\int_Y
 \left|\prod_{i=1}^mT_{i,n}\right|^r\,d\mu_n
 \leq\prod_{i=1}^m C_i^{\,r/p_i}.                           \tag{79.13}
\end{equation}
\end{theorem}

\begin{proof}
The functions
\[
                         P_n:=\prod_{i=1}^mT_{i,n},
 \qquad P:=\prod_{i=1}^mT_i
\]
are continuous, and finite products preserve uniform convergence on
each compact set.  Since
\(\sum_i r/p_i=1\), generalized Hölder gives
\[
 \int|P_n|^r\,d\mu_n
 \leq\prod_i
 \left(\int|T_{i,n}|^{p_i}\,d\mu_n\right)^{r/p_i},
\]
proving (79.13).  Apply Theorem
\ref{thm:v12v79-moment} to \(P_n,P\) with exponent \(r>1\).
\end{proof}

Every connected correlation is a finite linear combination of
ordinary products and products of lower moments.  Therefore the same
hypotheses pass connected correlations of the writers
\(T_{i,n}\).

The strict inequality \(\theta<1\) is the superlinear margin.  At
\(\theta=1\), Hölder gives only an \(L^1\) bound, which does not by
itself imply uniform integrability.

\subsection{Growth controlled by the physical Lyapunov function}

Let \({\cal L}:Y\to[0,\infty]\) be the physical Sobolev orbit control
of V73, and suppose for some \(P>0\),
\begin{equation}
                         \sup_n\int_Y{\cal L}^P\,d\mu_n
                         \leq C_L.                           \tag{79.14}
\end{equation}
Assume a writer has the uniform growth bound
\begin{equation}
                         |T_n(y)|
                         \leq A\{1+{\cal L}(y)^q\}            \tag{79.15}
\end{equation}
with \(0<q<P\).  Choose any \(p\) with
\[
                              1<p\leq {P\over q}.             \tag{79.16}
\]

\begin{corollary}[Lyapunov-to-writer uniform integrability]
\label{cor:v12v79-growth}
Equations (79.14)--(79.16) imply a bound of the form (79.5), uniform
in \(n\).  Hence every locally uniformly convergent writer satisfying
(79.15) passes its expectation to the continuum.
\end{corollary}

\begin{proof}
For \(p\geq1\),
\[
 |T_n|^p
 \leq 2^{p-1}A^p\{1+{\cal L}^{qp}\}.
\]
Since \(qp\leq P\),
\(t^{qp}\leq1+t^P\) for \(t\geq0\).  Integrate and use (79.14).
Then apply Theorem \ref{thm:v12v79-moment}.
\end{proof}

For an \(m\)-point correlation with growth orders \(q_i\), it is
enough to choose \(p_i\leq P/q_i\) with
\(\sum_i1/p_i<1\).  A sufficient condition is
\begin{equation}
                              \sum_{i=1}^m q_i<P.             \tag{79.17}
\end{equation}
Indeed, choose \(p_i\) slightly below \(P/q_i\) while retaining the
strict reciprocal-sum margin.

\subsection{Einstein and stress-tensor use}

On a uniformly nondegenerate Sobolev screen, the curvature and
Einstein writers are continuous in the weak topology described in
V69 and corrected in V73, after smearing against fixed smooth test
tensors.  If their cutoff reconstructions converge locally uniformly
and their growth is bounded by (79.15), Corollary
\ref{cor:v12v79-growth} passes their one-point functions.
Theorem \ref{thm:v12v79-products} passes any finite collection whose
total growth satisfies (79.17).

A distribution-valued stress tensor is handled only after smearing.
Point evaluation of a distribution and a supremum over all test
functions are not consequences of these scalar theorems.  Ward
identities containing such insertions require the same common moment
stock.

\begin{assessment}[Status after V79]
\label{ass:v12v79-status}
The bounded-cylinder continuum of V77 now extends to unbounded
locally convergent Einstein and stress writers whenever a uniform
superlinear moment is available.  The tail rate is explicit, finite
products pass under the sharp Hölder margin
\(\sum_i1/p_i<1\), and the physical Lyapunov moment of V73 supplies
those writer moments whenever total growth remains below its
available power.
\end{assessment}

\begin{firewall}[Claim boundary after V79]
\label{fire:v12v79-claim}
No RPESM physical Lyapunov moment, writer growth exponent, local
uniform reconstruction, or stress-tensor moment has been populated.
The theorem does not define unsmeared distributional products,
renormalise coincident insertions, or prove the Einstein equation in
expectation.  Reflection positivity, nontriviality, determinant
phase, conditional quantum convergence, and universality remain
separate gates.  The full SM--Einstein continuum is not yet
constructed.
\end{firewall}

\subsection{Parent record}

This V79 note appends to the validated V78 file with record
\[
\begin{array}{c|c}
\text{lines}&22137\\
\text{bytes}&745294\\
\text{SHA--256}&
\texttt{8BE6FDE433967A48215A09E111C1907C9CFE740A45D62858ADFF34144D794E24}.
\end{array}
\]

% =============================================================================
% END APPENDED TWELFTH-CYCLE V79 MATERIAL
% =============================================================================



% =============================================================================
% BEGIN APPENDED TWELFTH-CYCLE V80 MATERIAL
% =============================================================================

\section{Twelfth-cycle V80: mandatory companion audit}
\label{sec:v12v80}

This multiple-of-five checkpoint inspects the live companion notes
before the joint-state construction is continued beyond V79.

\subsection{Companion records}

The files inspected read-only are
\[
\begin{array}{c|r|r|c}
\text{record}&\text{lines}&\text{bytes}&\text{SHA--256}\\ \hline
\text{Yang--Mills V59}&22075&750408&
\texttt{AACEF6B54245542726D0D597400DC0C0A42EA6025F04C66244BDEE606831399A}\\
\text{Navier--Stokes V26}&5319&278283&
\texttt{82C887F8C2C69E4F28FAD7C3DC8DE5B9099CB5A4BF431EBA1FFF3E91935978AD}.
\end{array}                                                   \tag{80.1}
\]
Both records are unchanged from the V75 audit.

Yang--Mills V57--V59 still requires exact arbitrary-momentum
constituents to be assembled before scalar response extraction.
The corresponding SM rule has already been enforced in V74--V77:
the joint observable must be defined on the physical algebra before
its countable ledger is compressed.  Navier--Stokes V26 still
requires proof that an input lies in a structured class before a
smaller structured norm is used.  The SM analogue is that a
trace-class conditional-density estimate may be used only where
normality and the energy-tail hypotheses have been established.

No companion constant or theorem supplies the missing gravitational
measure, conformal tail, or conditional SM state.

\begin{decision}[Direction after the V80 audit]
\label{dec:v12v80-next}
Retain V69--V79.  At V81, construct the joint limit directly in the
state space of the separable algebra
\[
 C(K)\mathbin{\widehat\otimes}
 \widetilde{\mathcal K}({\cal H}_*),
\]
where \(K\) is the compact physical observable space and
\(\widetilde{\mathcal K}\) is the unitised compact-operator algebra
on the V61 reference Hilbert space.  Prove:
\begin{enumerate}
\item a countable state ledger metrises weak-star convergence and
      gives a full-sequence Cauchy theorem;
\item cutoff background--density pairs define positive normalized
      states on that algebra without requiring a pointwise limiting
      density field;
\item a uniform compact-resolvent energy moment excludes quantum mass
      at infinity in the limit;
\item the scalar physical marginal agrees with the V77 law.
\end{enumerate}
This weakens the V66 conditional convergence input while keeping
normality as an explicit energy gate.
\end{decision}

\begin{assessment}[Status after V80]
\label{ass:v12v80-status}
The mandatory audit is complete and changes no prior theorem.  It
confirms that the next bottleneck is joint-state convergence, not
another scalar compactness estimate.  V81 will work at the level of
the complete physical state rather than demand uniform pointwise
convergence of every conditional density operator.
\end{assessment}

\begin{firewall}[Claim boundary after V80]
\label{fire:v12v80-claim}
No companion result is imported as an SM estimate.  No joint-state
ledger, energy tail, normal limiting state, or physical local
operator algebra has yet been populated.  The full SM--Einstein
continuum is not yet constructed.
\end{firewall}

\subsection{Parent record}

This V80 note appends to the validated V79 file with record
\[
\begin{array}{c|c}
\text{lines}&22417\\
\text{bytes}&754912\\
\text{SHA--256}&
\texttt{CEBFDB7F5DC953BB2CAE3908DE5B279D317DAF042AF996E83DF49CB7316DDEC9}.
\end{array}
\]

% =============================================================================
% END APPENDED TWELFTH-CYCLE V80 MATERIAL
% =============================================================================

% =============================================================================
% BEGIN APPENDED TWELFTH-CYCLE V81 MATERIAL
% =============================================================================

\section{Twelfth-cycle V81: direct compactness of the joint
classical--quantum state and exclusion of quantum escape}
\label{sec:v12v81}

V66 constructs a joint limit from uniform convergence of conditional
density fields.  That hypothesis is stronger than convergence of the
physical state itself.  This note works directly on a separable
classical--quantum \(C^*\)-algebra.  Its state space is compact and
metrizable, a countable ledger gives a full-sequence criterion, and a
uniform compact-resolvent energy moment prevents the limiting state
from losing normal quantum mass at arbitrarily high modes.

\subsection{The separable joint algebra}

Let \(K\) be the compact physical observable space of V74--V77 and
let \({\cal H}_*\) be the separable reference Hilbert space supplied
by the V61 transports.  Write
\[
 {\cal K}_*:={\cal K}({\cal H}_*),\qquad
 \widetilde{\cal K}_*:={\cal K}_*+\mathbb C I              \tag{81.1}
\]
for the compact operators and their unitisation.  Define
\begin{equation}
 {\mathfrak A}
 :=C(K)\otimes_{\min}\widetilde{\cal K}_*
 \cong C(K;\widetilde{\cal K}_*).                            \tag{81.2}
\end{equation}
This is a separable unital \(C^*\)-algebra.  Its closed ideal
\begin{equation}
 {\mathfrak I}
 :=C(K)\otimes_{\min}{\cal K}_*
 \cong C(K;{\cal K}_*)                                      \tag{81.3}
\end{equation}
contains the quantum-normal observables.

Let \({\cal S}({\mathfrak A})\) be the positive normalized
functionals on \({\mathfrak A}\).  Choose a sequence
\(\{A_m:m\geq1\}\) dense in the self-adjoint unit ball of
\({\mathfrak A}\), and set
\begin{equation}
 d_{\mathfrak A}(\Omega,\Psi)
 :=
 \sum_{m=1}^{\infty}2^{-m-1}
 |\Omega(A_m)-\Psi(A_m)|.                                   \tag{81.4}
\end{equation}

\begin{theorem}[Compact metrizable joint-state space]
\label{thm:v12v81-state-space}
The metric \(d_{\mathfrak A}\) induces the weak-star topology on
\({\cal S}({\mathfrak A})\).  This state space is compact and
complete under \(d_{\mathfrak A}\).
\end{theorem}

\begin{proof}
States have norm one, so (81.4) converges and is at most one.
Agreement on the dense self-adjoint family implies agreement on the
entire self-adjoint part by norm approximation, and then on all of
\({\mathfrak A}\) by complex linearity.  Thus (81.4) is a metric.

Weak-star convergence gives convergence on every \(A_m\); a
finite-head and uniformly bounded-tail argument then gives
\(d_{\mathfrak A}\)-convergence.  Conversely,
\(d_{\mathfrak A}\)-convergence gives convergence on the dense
family, and norm approximation gives convergence on every element of
\({\mathfrak A}\).  Hence the two topologies agree.

The unit ball of \({\mathfrak A}^*\) is weak-star compact by
Banach--Alaoglu.  Positivity and the equation \(\Omega(I)=1\) are
weak-star closed, so the state space is compact.  Separability of
\({\mathfrak A}\) makes this compact topology metrizable by (81.4).
A compact metric space is complete.
\end{proof}

\subsection{Cutoff density fields define joint states}

Let \(\nu_n\) be a probability law on \(K\), and let
\[
 \rho_n:K\longrightarrow{\mathfrak S}_1({\cal H}_*)
\]
be strongly measurable with
\(\rho_n(u)\succeq0\) and \(\operatorname{Tr}\rho_n(u)=1\)
for \(\nu_n\)-almost every \(u\).  Define
\begin{equation}
 \Omega_n(F)
 :=
 \int_K\operatorname{Tr}\{\rho_n(u)F(u)\}\,d\nu_n(u),
 \qquad F\in{\mathfrak A}.                                  \tag{81.5}
\end{equation}

\begin{lemma}[Finite joint states]
\label{lem:v12v81-finite-state}
Equation (81.5) defines a state on \({\mathfrak A}\), and its
classical marginal is
\begin{equation}
                         \Omega_n(fI)=\int_Kf\,d\nu_n
                         \qquad(f\in C(K)).                  \tag{81.6}
\end{equation}
\end{lemma}

\begin{proof}
Trace duality gives
\[
 |\operatorname{Tr}\{\rho_n(u)F(u)\}|
 \leq\|F\|_\infty,
\]
so the integral exists under the stated measurability.  Pointwise
positivity of \(F\) gives positivity of the integral, and
\(\Omega_n(I)=1\).  Equation (81.6) follows from
\(\operatorname{Tr}\rho_n(u)=1\).
\end{proof}

No limiting conditional field is needed to regard a weak-star limit
of the \(\Omega_n\) as a joint physical state.

\subsection{Full-sequence state-ledger construction}

Put
\begin{equation}
                         b_n:=
 d_{\mathfrak A}(\Omega_{n+1},\Omega_n).                     \tag{81.7}
\end{equation}

\begin{theorem}[Target-free joint-state limit]
\label{thm:v12v81-cofinal}
If
\begin{equation}
                              \sum_{n=1}^{\infty}b_n<\infty,  \tag{81.8}
\end{equation}
then there is a unique
\(\Omega_\infty\in{\cal S}({\mathfrak A})\) such that
\[
                         \Omega_n\longrightarrow\Omega_\infty
                         \quad\text{weak-star},
\]
and
\begin{equation}
 d_{\mathfrak A}(\Omega_N,\Omega_\infty)
 \leq\sum_{n=N}^{\infty}b_n.                                \tag{81.9}
\end{equation}
It is enough, without a quantitative rate, that every scalar
sequence \(\Omega_n(A_m)\) converge.
\end{theorem}

\begin{proof}
Under (81.8), the triangle inequality makes \((\Omega_n)\) Cauchy in
the complete metric of Theorem \ref{thm:v12v81-state-space}.
Letting the upper endpoint tend to infinity gives (81.9).

For the qualitative statement, compactness gives a convergent
subsequence.  Every subsequential limit has the prescribed values on
all \(A_m\), hence is the same state by density.  If the full
sequence failed to converge to that state, compactness would produce
a subsequence with a different cluster point, a contradiction.
\end{proof}

As in V77, a finite ledger head gives
\begin{equation}
 b_n\leq
 \sum_{m=1}^{M}2^{-m-1}
 |\Omega_{n+1}(A_m)-\Omega_n(A_m)|
 +2^{-M}.                                                   \tag{81.10}
\end{equation}
A summable certificate therefore requires \(M=M_n\) with
\(\sum_n2^{-M_n}<\infty\), together with summable finite-head
defects.

\begin{corollary}[Compatibility with the V77 marginal]
\label{cor:v12v81-marginal}
If the classical laws satisfy
\(\nu_n\Rightarrow\nu_\infty\), then every weak-star limit in
Theorem \ref{thm:v12v81-cofinal} obeys
\begin{equation}
                 \Omega_\infty(fI)=\int_Kf\,d\nu_\infty
                 \qquad(f\in C(K)).                         \tag{81.11}
\end{equation}
\end{corollary}

\begin{proof}
Use (81.6), weak convergence of \(\nu_n\), and weak-star convergence
of \(\Omega_n\) on the fixed algebra element \(fI\).
\end{proof}

\subsection{The quantum escape channel}

A state on the unitised algebra may have part of its norm supported
on the scalar quotient rather than on the compact-operator ideal.
The next condition excludes this loss.

Let \(H_0\geq0\) be self-adjoint on \({\cal H}_*\), with compact
resolvent.  Enumerate its eigenvalues with multiplicity,
\[
 0\leq E_1\leq E_2\leq\cdots,\qquad E_j\longrightarrow\infty,
\]
and let \(P_J\) project onto the first \(J\) eigenvectors.  Put
\(Q_J=I-P_J\).

Assume the cutoff density fields have the uniform energy moment
\begin{equation}
 \sup_n\int_K
 \operatorname{Tr}\{\rho_n(u)H_0\}\,d\nu_n(u)
 \leq{\cal E}<\infty.                                       \tag{81.12}
\end{equation}

\begin{lemma}[Uniform quantum compact containment]
\label{lem:v12v81-energy-tail}
For every \(J\) with \(E_{J+1}>0\),
\begin{equation}
 \Omega_n(1\otimes Q_J)
 \leq{{\cal E}\over E_{J+1}}.                               \tag{81.13}
\end{equation}
\end{lemma}

\begin{proof}
The spectral theorem gives the quadratic-form inequality
\[
                              H_0\geq E_{J+1}Q_J.
\]
Trace it against each positive \(\rho_n(u)\), integrate, and use
(81.12).
\end{proof}

\begin{theorem}[No quantum mass at infinity]
\label{thm:v12v81-no-escape}
Let \(\Omega_n\to\Omega_\infty\) weak-star and assume (81.12).
Then
\begin{align}
 \Omega_\infty(1\otimes Q_J)
 &\leq{{\cal E}\over E_{J+1}},                              \tag{81.14}\\
 \lim_{J\to\infty}\Omega_\infty(1\otimes P_J)&=1,            \tag{81.15}\\
 \|\Omega_\infty|_{\mathfrak I}\|&=1.                       \tag{81.16}
\end{align}
Thus the limiting state has no norm defect in the
compact-operator ideal.  Its constant-quantum restriction is a
normal state: there is a unique
\(\overline\rho_\infty\in{\mathfrak S}_1({\cal H}_*)\) with
\[
 \overline\rho_\infty\succeq0,\qquad
 \operatorname{Tr}\overline\rho_\infty=1,\qquad
 \Omega_\infty(1\otimes A)
 =\operatorname{Tr}(\overline\rho_\infty A)
 \quad(A\in{\cal K}_*).                                     \tag{81.17}
\]
\end{theorem}

\begin{proof}
The element \(1\otimes Q_J\) belongs to \({\mathfrak A}\), so
weak-star convergence passes (81.13) to (81.14).  Since
\(E_{J+1}\to\infty\) and
\(P_J+Q_J=I\), equation (81.15) follows.

The elements \(e_J:=1\otimes P_J\) form a contractive approximate
identity for \({\mathfrak I}\).  Indeed, \(P_JA\to A\) and
\(AP_J\to A\) in norm for each compact operator \(A\); the convergence
is uniform on the compact range of every
\(F\in C(K;{\cal K}_*)\).  For a positive functional, the norm of its
restriction to an ideal is the supremum of its values on a positive
contractive approximate identity.  Equation (81.15) therefore gives
(81.16).

Restrict \(\Omega_\infty\) further to the constant copy of
\({\cal K}_*\).  It is a positive functional of norm one by
(81.15).  Since
\({\cal K}({\cal H}_*)^*={\mathfrak S}_1({\cal H}_*)\), it has the
unique trace-class representation (81.17); positivity and norm one
give the stated properties.
\end{proof}

The theorem proves normality of the unconditional quantum marginal
and full norm on the joint compact-operator ideal.  It does not
assert the existence of a pointwise trace-class disintegration
\(u\mapsto\rho_\infty(u)\).  Such a disintegration is an additional
Radon--Nikodym/measurability question and is unnecessary for defining
the state on \({\mathfrak I}\).

\subsection{Closed structural gates}

If a reflection operation and positive-time subalgebra act
continuously on \({\mathfrak A}\), every inequality
\[
 \Omega_n((\Theta A)^*A)\geq-\varepsilon_n(A),
 \qquad\varepsilon_n(A)\to0
\]
passes to \(\Omega_\infty\) by weak-star convergence, as in V67 and
V74.  Bounded Ward residuals pass for the same reason.  Unbounded
stress or Hamiltonian insertions require the V79 uniform-integrability
stock in addition to the energy containment (81.12).

\begin{assessment}[Status after V81]
\label{ass:v12v81-status}
The SM--Einstein joint state can now be constructed directly from a
countable full-state ledger on a compact physical observable space,
without uniform pointwise convergence of conditional density fields.
A uniform compact-resolvent energy moment gives an explicit
high-mode tail and proves that the limiting state retains norm one
on the joint compact-operator ideal, with a normal trace-class
unconditional quantum marginal.
\end{assessment}

\begin{firewall}[Claim boundary after V81]
\label{fire:v12v81-claim}
No RPESM joint-state ledger, common physical compact space, or
uniform energy moment (81.12) has been populated.  The algebra
\(\widetilde{\cal K}_*\) is a controlled separable core and has not
been proved to contain the full local bounded-observable algebra.
Pointwise conditional disintegration, determinant phases,
reflection positivity, nontrivial dynamics, clustering, and
universality remain separate gates.  The full SM--Einstein continuum
is not yet constructed.
\end{firewall}

\subsection{Parent record}

This V81 note appends to the validated V80 file with record
\[
\begin{array}{c|c}
\text{lines}&22512\\
\text{bytes}&758697\\
\text{SHA--256}&
\texttt{8F7DF59F30EDC298DD6D9149CE370136100937B4435C371E11C78281B54EE3CD}.
\end{array}
\]

% =============================================================================
% END APPENDED TWELFTH-CYCLE V81 MATERIAL
% =============================================================================

% =============================================================================
% BEGIN APPENDED TWELFTH-CYCLE V82 MATERIAL
% =============================================================================

\section{Twelfth-cycle V82: disintegration of the no-escape joint
state into conditional density operators}
\label{sec:v12v82}

V81 constructs a joint state on
\(C(K;\widetilde{\cal K}_*)\) and deliberately does not infer a
pointwise limiting density field.  The no-quantum-escape conclusion
is exactly the additional input needed for such a disintegration.
This note proves it by taking Radon--Nikodym derivatives of the
matrix-unit measures.  Thus the direct state-space route recovers a
conditional trace-class field after the limit, without requiring
uniform pointwise convergence before the limit.

\subsection{Statement of the disintegration theorem}

Let \(K\) be compact metric, let \({\cal H}_*\) be separable with
orthonormal basis \((e_j)_{j\geq1}\), and retain
\[
 {\mathfrak A}=C(K;\widetilde{\cal K}_*),\qquad
 {\mathfrak I}=C(K;{\cal K}_*).
\]
Let \(P_J\) be projection onto
\(\operatorname{span}\{e_1,\ldots,e_J\}\) and \(Q_J=I-P_J\).
Let \(\Omega\) be a state on \({\mathfrak A}\).  Its central
classical marginal is the unique probability measure \(\mu\) such
that
\begin{equation}
                         \Omega(fI)=\int_Kf\,d\mu
                         \qquad(f\in C(K)).                  \tag{82.1}
\end{equation}

\begin{theorem}[No-escape conditional-density disintegration]
\label{thm:v12v82-disintegration}
Assume
\begin{equation}
                         \Omega(1\otimes Q_J)\longrightarrow0.
                                                                    \tag{82.2}
\end{equation}
Then there is a strongly \(\mu\)-measurable field
\[
 \rho:K\longrightarrow{\mathfrak S}_1({\cal H}_*)
\]
such that, for \(\mu\)-almost every \(u\),
\begin{equation}
                         \rho(u)\succeq0,\qquad
                         \operatorname{Tr}\rho(u)=1,         \tag{82.3}
\end{equation}
and
\begin{equation}
 \boxed{\displaystyle
 \Omega(F)=\int_K
 \operatorname{Tr}\{\rho(u)F(u)\}\,d\mu(u)}
 \qquad(F\in{\mathfrak A}).                                 \tag{82.4}
\end{equation}
The field is unique up to a \(\mu\)-null set.
\end{theorem}

\subsection{Matrix-unit measures}

Let \(E_{ij}=|e_i\rangle\langle e_j|\).  For each \(i,j\), the map
\[
 f\longmapsto\Omega(f\otimes E_{ji})
\]
is a bounded linear functional on \(C(K)\).  By the complex Riesz
representation theorem there is a finite complex Radon measure
\(m_{ij}\) with
\begin{equation}
                         \Omega(f\otimes E_{ji})
                         =\int_Kf\,dm_{ij}.                  \tag{82.5}
\end{equation}

\begin{lemma}[Absolute continuity and positive matrix densities]
\label{lem:v12v82-matrix-density}
Every \(m_{ij}\) is absolutely continuous with respect to \(\mu\).
Writing
\[
                         r_{ij}:={dm_{ij}\over d\mu},         \tag{82.6}
\]
one may choose versions on one common full-measure set such that,
for every \(J\),
\begin{equation}
                         R_J(u):=[r_{ij}(u)]_{i,j\leq J}
                         \succeq0,                           \tag{82.7}
\end{equation}
the matrices are consistent under compression, and
\begin{equation}
                  0\leq t_J(u):=\operatorname{Tr}R_J(u)
                  \leq1.                                    \tag{82.8}
\end{equation}
\end{lemma}

\begin{proof}
For \(f\geq0\) and a finite vector
\(v=\sum_i c_ie_i\), positivity of \(\Omega\) gives
\begin{equation}
 \sum_{i,j}\overline c_i c_j
 \int_Kf\,dm_{ij}
 =
 \Omega(f\otimes|v\rangle\langle v|)\geq0.                  \tag{82.9}
\end{equation}
In particular the diagonal measures are positive.  Since
\(0\leq E_{ii}\leq I\),
\[
                         0\leq m_{ii}\leq\mu.                \tag{82.10}
\]
The \(2\times2\) positivity in (82.9) gives, for every Borel set
\(B\),
\[
 |m_{ij}(B)|^2\leq m_{ii}(B)m_{jj}(B)
\]
first on continuity sets and then on all Borel sets by regular
approximation.  Hence \(\mu(B)=0\) implies \(m_{ij}(B)=0\), proving
absolute continuity.

There are countably many rational finite vectors.  Apply (82.9) to
them and use Radon--Nikodym differentiation; outside one common null
set, every corresponding quadratic form in the densities is
nonnegative.  Density of rational vectors gives (82.7), and the
definitions make the finite matrices compression-consistent.

Finally \(0\leq P_J\leq I\) gives, for \(f\geq0\),
\[
 \int_K f\,t_J\,d\mu
 =\Omega(f\otimes P_J)
 \leq\Omega(fI)=\int_Kf\,d\mu.
\]
Thus \(0\leq t_J\leq1\) almost everywhere.
\end{proof}

\subsection{Construction of the trace-class field}

\begin{lemma}[Unit trace from no escape]
\label{lem:v12v82-unit-trace}
Under (82.2),
\begin{equation}
                              t_J(u)\uparrow1
                              \quad\text{for }\mu\text{-almost every }u.
                                                                    \tag{82.11}
\end{equation}
For almost every \(u\), the consistent positive matrices
\(R_J(u)\) are the compressions of a unique positive trace-class
operator \(\rho(u)\) satisfying (82.3).
\end{lemma}

\begin{proof}
Positivity and compression consistency make \(t_J\) increasing.
Moreover,
\[
 \int_Kt_J\,d\mu
 =\Omega(1\otimes P_J)
 =1-\Omega(1\otimes Q_J)\longrightarrow1.                  \tag{82.12}
\]
By (82.8), the monotone limit \(t_\infty\) is at most one.
Monotone convergence and (82.12) give
\(\int t_\infty\,d\mu=1\), hence \(t_\infty=1\) almost
everywhere.

Fix such a point \(u\).  The compatible positive forms defined by
the \(R_J(u)\) have uniformly bounded trace one.  They therefore
define a positive bounded operator \(\rho(u)\): for finitely
supported \(v\),
\[
 0\leq\langle v,\rho(u)v\rangle
 \leq\|v\|^2\sup_J\operatorname{Tr}R_J(u)\leq\|v\|^2.
\]
Its diagonal trace is
\(\sum_i r_{ii}(u)=1\).  A positive operator with finite diagonal
sum in one orthonormal basis is trace class and has trace equal to
that sum.  Its finite compressions are exactly \(R_J(u)\), proving
existence and uniqueness.
\end{proof}

\begin{lemma}[Strong trace-class measurability]
\label{lem:v12v82-measurable}
The field \(u\mapsto\rho(u)\) can be chosen strongly measurable as a
map into \({\mathfrak S}_1({\cal H}_*)\).
\end{lemma}

\begin{proof}
For each \(J\),
\[
                         \rho_J(u):=P_J\rho(u)P_J
\]
has the finite matrix of measurable entries
\([r_{ij}(u)]_{i,j\leq J}\), so it is strongly measurable in the
finite-dimensional trace-class space.  At every point where
\(\rho(u)\) is positive trace class,
\[
                         \|\rho_J(u)-\rho(u)\|_1\longrightarrow0.
\]
Thus \(\rho\) is the pointwise trace-norm limit of strongly
measurable functions into the separable Banach space
\({\mathfrak S}_1\), and is strongly measurable.
\end{proof}

\subsection{Proof of the disintegration formula}

\begin{proof}[Proof of Theorem \ref{thm:v12v82-disintegration}]
Equation (82.5) and the definition of \(r_{ij}\) give (82.4) for
fields \(f\otimes E_{ji}\).  By linearity it holds for
finite-rank-valued continuous fields.

Such fields are norm dense in
\(C(K;{\cal K}_*)\).  Indeed, if \(F:K\to{\cal K}_*\) is continuous,
its image is compact in the compact operators, and
\[
                         \sup_{u\in K}
 \|P_JF(u)P_J-F(u)\|\longrightarrow0.                       \tag{82.13}
\]
Both sides of (82.4) are norm-continuous functionals, so the identity
extends to \({\mathfrak I}\).

For \(f\in C(K)\), equation (82.3) gives
\[
 \int_K\operatorname{Tr}\{\rho(u)f(u)I\}\,d\mu
 =\int_Kf\,d\mu=\Omega(fI).
\]
Every element of \({\mathfrak A}\) is uniquely a sum of a
compact-valued field and a scalar field times \(I\), so (82.4)
follows on all of \({\mathfrak A}\).

If \(\rho'\) is another such field, test (82.4) on
\(f\otimes E_{ji}\).  Uniqueness of Radon--Nikodym derivatives gives
\(\langle e_i,\rho e_j\rangle
 =\langle e_i,\rho'e_j\rangle\) almost everywhere for every
\(i,j\).  Outside the union of these countably many null sets, the
two trace-class operators agree.
\end{proof}

\subsection{Energy passes to the conditional field}

Let \(H_0\) and its spectral projections \(P_J\) be as in V81, with
the basis chosen to be an eigenbasis.  Put
\[
                         H_{0,J}:=P_JH_0P_J\in{\cal K}_*.
\]

\begin{corollary}[Limit energy moment]
\label{cor:v12v82-energy}
Suppose \(\Omega_n\to\Omega\) weak-star and the cutoff fields obey
the uniform energy bound (81.12).  Then the disintegrated field in
Theorem \ref{thm:v12v82-disintegration} satisfies
\begin{equation}
 \int_K\operatorname{Tr}\{\rho(u)H_0\}\,d\mu(u)
 \leq{\cal E}.                                               \tag{82.14}
\end{equation}
\end{corollary}

\begin{proof}
For fixed \(J\),
\[
 \Omega(1\otimes H_{0,J})
 =\lim_n\Omega_n(1\otimes H_{0,J})
 \leq{\cal E}.                                               \tag{82.15}
\]
Formula (82.4) identifies the left side with the integral of
\(\operatorname{Tr}(\rho H_{0,J})\).  These nonnegative functions
increase pointwise to
\(\operatorname{Tr}(\rho H_0)\).  Monotone convergence and
(82.15) prove (82.14).
\end{proof}

Consequently, for \(\mu\)-almost every \(u\), the conditional state
has finite energy, except possibly on a null set; the average bound
alone does not give a uniform pointwise energy ceiling.

\subsection{Relation to V66}

V66 remains the quantitative route when a uniformly Hölder
conditional field and summable pointwise trace-norm errors are
available.  V81--V82 give a weaker compactness route:
\[
 \text{joint state ledger}
 \longrightarrow
 \text{weak-star joint limit}
 \longrightarrow
 \text{energy no escape}
 \longrightarrow
 \text{measurable conditional density}.                    \tag{82.16}
\]
This route supplies no uniform Hölder modulus for the limiting field
and no trace-norm convergence \(\rho_n(u)\to\rho(u)\) at fixed
\(u\).  It does construct the same kind of joint integral
representation after the limit.

\begin{assessment}[Status after V82]
\label{ass:v12v82-status}
The no-escape joint state of V81 now has a unique almost-everywhere
disintegration into positive trace-one conditional density
operators over its physical classical marginal.  The construction
uses only countably many matrix-unit Radon--Nikodym derivatives, and
the uniform cutoff energy moment passes to the disintegrated field.
Uniform pointwise conditional convergence is therefore sufficient
but no longer necessary for existence of a conditional continuum
state.
\end{assessment}

\begin{firewall}[Claim boundary after V82]
\label{fire:v12v82-claim}
No physical RPESM joint ledger, energy estimate, or no-escape state
has been populated.  The disintegrated field is only measurable; no
continuity, locality, KMS property, spectral gap, or pointwise
Hamiltonian representation follows.  The full local observable
algebra, determinant phase, reflection positivity, nontrivial
dynamics, clustering, and universality remain open.  The full
SM--Einstein continuum is not yet constructed.
\end{firewall}

\subsection{Parent record}

This V82 note appends to the validated V81 file with record
\[
\begin{array}{c|c}
\text{lines}&22842\\
\text{bytes}&770751\\
\text{SHA--256}&
\texttt{C00A1AE94202DAED560CD0EE2DA85B7D0076E799D936CA72BF3596EC306E80A4}.
\end{array}
\]

% =============================================================================
% END APPENDED TWELFTH-CYCLE V82 MATERIAL
% =============================================================================

% =============================================================================
% BEGIN APPENDED TWELFTH-CYCLE V83 MATERIAL
% =============================================================================

\section{Twelfth-cycle V83: compatible local restrictions and the
quasi-local continuum state}
\label{sec:v12v83}

V81--V82 construct and disintegrate a joint state on a controlled
compact-operator core.  The full continuum also needs one state on
the completed algebra generated by all bounded local observables.
This note gives the exact inductive-limit handoff.  It uses
convergence of every fixed local restriction and does not require a
cutoff state to be defined on observables that have not yet appeared
at that cutoff.

\subsection{The nested local system}

Let
\[
 {\mathfrak A}_1\subset{\mathfrak A}_2\subset\cdots          \tag{83.1}
\]
be an increasing sequence of separable unital \(C^*\)-algebras with
the same identity and isometric unital inclusions.  Define
\begin{equation}
 {\mathfrak A}_{\rm loc}:=\bigcup_{m\geq1}{\mathfrak A}_m,
 \qquad
 {\mathfrak A}_{\rm ql}
 :=\overline{{\mathfrak A}_{\rm loc}}^{\,\|\cdot\|}.         \tag{83.2}
\end{equation}
The notation permits \({\mathfrak A}_m\) to encode a bounded
spacetime region, a finite observable resolution, and a finite
energy window simultaneously.  Exact inclusions mean that a retained
observable has the same algebraic meaning at every later cutoff.

For each \(n\), let \(\omega_n\) be a state on
\({\mathfrak A}_n\).  If \(m\leq n\), write
\[
                         \omega_n^{(m)}
 :=\omega_n|_{{\mathfrak A}_m}.                              \tag{83.3}
\]

\begin{theorem}[Local-restriction continuum state]
\label{thm:v12v83-local-limit}
Suppose for every fixed \(m\) there is a state
\(\omega_\infty^{(m)}\) on \({\mathfrak A}_m\) such that
\begin{equation}
                         \omega_n^{(m)}
 \longrightarrow\omega_\infty^{(m)}
 \quad\text{weak-star as }n\to\infty.                       \tag{83.4}
\end{equation}
Then the limiting states are compatible:
\begin{equation}
 \omega_\infty^{(m+1)}|_{{\mathfrak A}_m}
 =\omega_\infty^{(m)}.                                      \tag{83.5}
\end{equation}
There is a unique state \(\omega_\infty\) on
\({\mathfrak A}_{\rm ql}\) whose restriction to every
\({\mathfrak A}_m\) is \(\omega_\infty^{(m)}\).  Moreover, for every
local observable \(A\in{\mathfrak A}_{\rm loc}\),
\begin{equation}
                         \omega_n(A)\longrightarrow
                         \omega_\infty(A)                    \tag{83.6}
\end{equation}
once \(n\) is large enough that \(A\in{\mathfrak A}_n\).
\end{theorem}

\begin{proof}
For \(A\in{\mathfrak A}_m\), the number
\(\omega_n(A)\) is the same whether \(A\) is regarded as an element
of \({\mathfrak A}_m\) or \({\mathfrak A}_{m+1}\).  Taking the limit
in (83.4) proves (83.5).

Define \(\omega_0\) on the algebraic union by
\[
                         \omega_0(A):=\omega_\infty^{(m)}(A)
 \quad\text{if }A\in{\mathfrak A}_m.                        \tag{83.7}
\]
Compatibility makes this well-defined.  It is linear and positive
because every finite collection of its arguments lies in one
\({\mathfrak A}_m\).  It satisfies
\[
                         |\omega_0(A)|\leq\|A\|,\qquad
                         \omega_0(I)=1.                      \tag{83.8}
\]
Thus it extends uniquely by norm continuity to a positive
norm-one functional \(\omega_\infty\) on the completion
\({\mathfrak A}_{\rm ql}\).  This is a state.  Density of the union
gives uniqueness, and (83.6) is (83.4) in any local algebra
containing \(A\).
\end{proof}

\subsection{One countable cofinal ledger}

For each \(m\), choose a sequence
\(\{A_{m,k}:k\geq1\}\) dense in the self-adjoint unit ball of
\({\mathfrak A}_m\), and define on its state space
\begin{equation}
 d_m(\varphi,\psi)
 :=
 \sum_{k=1}^{\infty}2^{-k-1}
 |\varphi(A_{m,k})-\psi(A_{m,k})|.                           \tag{83.9}
\end{equation}
As in V81, \(d_m\) metrizes the weak-star topology and makes the
state space compact and complete.

For \(n\geq1\), put
\begin{equation}
 c_n:=
 \sum_{m=1}^{n}2^{-m}
 d_m(\omega_{n+1}^{(m)},\omega_n^{(m)}).                    \tag{83.10}
\end{equation}

\begin{theorem}[Summable quasi-local state ledger]
\label{thm:v12v83-ledger}
If
\begin{equation}
                              \sum_{n=1}^{\infty}c_n<\infty, \tag{83.11}
\end{equation}
then the hypotheses of Theorem
\ref{thm:v12v83-local-limit} hold.  Quantitatively, for every fixed
\(m\) and \(N\geq m\),
\begin{equation}
 d_m(\omega_N^{(m)},\omega_\infty^{(m)})
 \leq 2^m\sum_{n=N}^{\infty}c_n.                            \tag{83.12}
\end{equation}
\end{theorem}

\begin{proof}
For \(n\geq m\), the \(m\)-th nonnegative term in (83.10) gives
\[
 d_m(\omega_{n+1}^{(m)},\omega_n^{(m)})
 \leq2^mc_n.                                                 \tag{83.13}
\]
Equation (83.11) therefore makes
\((\omega_n^{(m)})_{n\geq m}\) Cauchy in the complete state space of
\({\mathfrak A}_m\).  Let its limit be
\(\omega_\infty^{(m)}\).  Summing (83.13) from \(N\) to infinity
proves (83.12).  Apply Theorem
\ref{thm:v12v83-local-limit}.
\end{proof}

The ledger (83.10) contains every fixed local scale for all
sufficiently late \(n\).  It does not require uniform control over
all \(m\) at one cutoff.  The geometric weights make the infinitely
many local scales summable.

\begin{theorem}[Qualitative diagonal criterion]
\label{thm:v12v83-qualitative}
It is enough that, for every \(m,k\), the scalar sequence
\begin{equation}
                              \omega_n(A_{m,k})               \tag{83.14}
\end{equation}
converge as \(n\to\infty\), \(n\geq m\).  Then a unique
quasi-local state satisfying (83.6) exists.
\end{theorem}

\begin{proof}
For a fixed \(m\), compactness of the state space gives cluster
points of \(\omega_n^{(m)}\).  Every cluster point has the prescribed
values on the dense family \(A_{m,k}\), so density makes it unique.
A compact sequence with one cluster point converges, giving (83.4).
Theorem \ref{thm:v12v83-local-limit} applies.
\end{proof}

\subsection{Embedding the V81 core}

Suppose there are compatible unital embeddings
\begin{equation}
 \jmath_m:
 C(K_m;\widetilde{\cal K}({\cal H}_m))
 \longrightarrow{\mathfrak A}_m                              \tag{83.15}
\end{equation}
of the V81 controlled cores, and the restriction of \(\omega_n\) to
each retained core obeys its energy no-escape bound.  Then V81--V82
give a normal density disintegration on each fixed core, while
Theorem \ref{thm:v12v83-local-limit} places all those compatible
restrictions in one quasi-local state.

This does not assert that the quasi-local state is normal in a
single global representation.  Local normality and global
quasi-equivalence are distinct representation-theoretic conditions.

\subsection{Closed local OS and Ward gates}

Let \(\Theta\) be a reflection operation preserving a nested
positive-time system
\({\mathfrak A}_{m,+}\subset{\mathfrak A}_m\).  Suppose for every
fixed \(m\) and \(A\in{\mathfrak A}_{m,+}\),
\begin{equation}
 \omega_n((\Theta A)^*A)\geq-\varepsilon_{n,m}(A),
 \qquad\varepsilon_{n,m}(A)\longrightarrow0.                \tag{83.16}
\end{equation}
Then
\begin{equation}
                         \omega_\infty((\Theta A)^*A)\geq0.  \tag{83.17}
\end{equation}
Likewise, if a bounded local Ward writer \(D_mA\in{\mathfrak A}_m\)
satisfies
\begin{equation}
                         |\omega_n(D_mA)|
 \leq\eta_{n,m}(A)\longrightarrow0,                         \tag{83.18}
\end{equation}
then \(\omega_\infty(D_mA)=0\).

\begin{proof}
For a fixed local element, equation (83.6) passes the left sides of
(83.16) and (83.18) to the limit.  Let the residuals tend to zero.
\end{proof}

Unbounded Ward insertions first require the uniform-integrability
handoff of V79.

\subsection{What remains to populate}

The theorem turns a compatible local family into one completed
state, but the physical regulator must still supply:
\begin{enumerate}
\item exact or controlled retained embeddings between cutoff local
      algebras;
\item convergence of the countable local ledger (83.11) or (83.14);
\item consistency of reflections, gauge quotients, and spacetime
      localization under those embeddings;
\item local energy tails and determinant-line phases on the same
      cofinal schedule.
\end{enumerate}

\begin{assessment}[Status after V83]
\label{ass:v12v83-status}
The joint compact core of V81--V82 now has a rigorous route into a
single separable quasi-local continuum algebra.  Compatible
weak-star limits on every fixed local algebra define a unique state
on the norm completion.  One summable weighted ledger controls all
local scales and gives explicit restriction tails, while local OS
and bounded Ward gates pass to that state.
\end{assessment}

\begin{firewall}[Claim boundary after V83]
\label{fire:v12v83-claim}
No physical RPESM net
\(({\mathfrak A}_m)\), retained embeddings, local ledger, or
consistent reflection/localization system has been populated.
Local normality on compact-operator cores does not prove global
normality, Haag--Kastler axioms, covariance, clustering, a spectrum
condition, or universality.  The full SM--Einstein continuum is not
yet constructed.
\end{firewall}

\subsection{Parent record}

This V83 note appends to the validated V82 file with record
\[
\begin{array}{c|c}
\text{lines}&23170\\
\text{bytes}&782272\\
\text{SHA--256}&
\texttt{E4C56972FF1ADDED3F35513528A00C0B5EF8658F1129B57DC11368A8A045CF23}.
\end{array}
\]

% =============================================================================
% END APPENDED TWELFTH-CYCLE V83 MATERIAL
% =============================================================================

% =============================================================================
% BEGIN APPENDED TWELFTH-CYCLE V84 MATERIAL
% =============================================================================

\section{Twelfth-cycle V84: passage of clustering to the quasi-local
state and channel spectral thresholds}
\label{sec:v12v84}

V83 constructs the quasi-local continuum state but does not yet pass
long-distance correlation estimates.  This note proves that a
uniform cutoff connected-correlation envelope survives when the
translated cutoff writer converges in norm to the continuum
translation.  It also records the exact OS spectral consequence of
an exponential Euclidean-time envelope in one channel.  No global
mass gap is expected for a theory containing photon and graviton
sectors.

\subsection{Translation compatibility}

Let \({\mathfrak A}_{\rm ql}\) and \(\omega_\infty\) be as in V83.
Assume a group of isometric automorphisms
\[
                         \tau_x:{\mathfrak A}_{\rm ql}
                         \longrightarrow{\mathfrak A}_{\rm ql},
 \qquad x\in\mathbb R^d,                                    \tag{84.1}
\]
which is strongly continuous in \(x\).  Let \(A,B\) be fixed bounded
local observables.  For each fixed \(x\), suppose
\(A,\tau_x(B)\), and \(A\tau_x(B)\) belong to some fixed local
algebra once the local index is enlarged.

At cutoff \(n\), let \(B_{n,x}\in{\mathfrak A}_n\) be the retained
translate of \(B\), defined for all sufficiently large \(n\), and
assume
\begin{equation}
                         \|B_{n,x}-\tau_x(B)\|
                         \leq\eta_n(B,x),
 \qquad\eta_n(B,x)\longrightarrow0                          \tag{84.2}
\end{equation}
for every fixed \(x\).

\begin{lemma}[Translated products pass]
\label{lem:v12v84-products}
Under the local convergence (83.6) and (84.2),
\begin{align}
 \omega_n(B_{n,x})&\longrightarrow
 \omega_\infty(\tau_x(B)),                                  \tag{84.3}\\
 \omega_n(AB_{n,x})&\longrightarrow
 \omega_\infty(A\tau_x(B))                                  \tag{84.4}
\end{align}
for every fixed \(x\).
\end{lemma}

\begin{proof}
For all sufficiently large \(n\),
\[
 |\omega_n(B_{n,x})-\omega_n(\tau_x(B))|
 \leq\eta_n(B,x),
\]
and
\[
 |\omega_n(AB_{n,x})-\omega_n(A\tau_x(B))|
 \leq\|A\|\eta_n(B,x).
\]
The fixed elements \(\tau_x(B)\) and \(A\tau_x(B)\) lie in one local
algebra, so (83.6) passes their expectations.  Let \(n\to\infty\).
\end{proof}

\subsection{Clustering envelopes}

Define the connected correlation
\begin{equation}
 {\cal C}_\omega(A,B;x)
 :=
 \omega(A\tau_x(B))-\omega(A)\omega(\tau_x(B)).              \tag{84.5}
\end{equation}
Let \(\Phi:[0,\infty)\to[0,\infty)\) satisfy
\[
                              \Phi(r)\longrightarrow0
                              \quad(r\to\infty).              \tag{84.6}
\]

\begin{theorem}[Uniform cutoff clustering passes]
\label{thm:v12v84-clustering}
Suppose constants \(C_{A,B}<\infty\) and residuals
\(\varepsilon_n(A,B,x)\) satisfy, for every fixed \(x\) and all
sufficiently large \(n\),
\begin{equation}
 \left|
 \omega_n(AB_{n,x})-\omega_n(A)\omega_n(B_{n,x})
 \right|
 \leq C_{A,B}\Phi(|x|)
      +\varepsilon_n(A,B,x),                                \tag{84.7}
\end{equation}
where
\begin{equation}
                         \varepsilon_n(A,B,x)\longrightarrow0
                         \quad\text{for each fixed }x.       \tag{84.8}
\end{equation}
Then
\begin{equation}
 \boxed{\displaystyle
 |{\cal C}_{\omega_\infty}(A,B;x)|
 \leq C_{A,B}\Phi(|x|)}                                     \tag{84.9}
\end{equation}
for every \(x\).  Consequently
\[
                         {\cal C}_{\omega_\infty}(A,B;x)
                         \longrightarrow0
                         \quad(|x|\to\infty).                \tag{84.10}
\]
\end{theorem}

\begin{proof}
By (83.6), \(\omega_n(A)\to\omega_\infty(A)\).
Lemma \ref{lem:v12v84-products} passes the other two expectations in
(84.7).  Products of convergent scalar sequences converge, so the
left side of (84.7) tends to the left side of (84.9).  Use (84.8)
and take \(n\to\infty\).  Then (84.6) gives (84.10).
\end{proof}

The residual need not be uniform in \(x\) to prove (84.9), because
the cutoff limit is taken at each fixed \(x\) before the
long-distance limit.  The envelope and its constant must be uniform
in the cutoff.

\begin{corollary}[Power-law and exponential channels]
\label{cor:v12v84-rates}
If (84.7) holds with
\[
 \Phi(r)=(1+r)^{-q}
 \quad\text{or}\quad
 \Phi(r)=e^{-mr},
\]
then the same power-law or exponential envelope holds in the
continuum with the identical \(q\) or \(m\).
\end{corollary}

\begin{proof}
Apply Theorem \ref{thm:v12v84-clustering}.
\end{proof}

\subsection{Translation invariance is a separate residual}

Suppose the cutoff retained translates also satisfy
\begin{equation}
 |\omega_n(B_{n,x})-\omega_n(B)|
 \leq\zeta_n(B,x),\qquad
 \zeta_n(B,x)\longrightarrow0                               \tag{84.11}
\end{equation}
for every local \(B\) and fixed \(x\).

\begin{proposition}[Continuum translation invariance]
\label{prop:v12v84-invariance}
Under (84.2), (84.11), and (83.6),
\begin{equation}
                         \omega_\infty(\tau_x(B))
                         =\omega_\infty(B)                   \tag{84.12}
\end{equation}
for every local \(B\).  By norm density it holds for every
\(B\in{\mathfrak A}_{\rm ql}\).
\end{proposition}

\begin{proof}
Lemma \ref{lem:v12v84-products} passes the first expectation in
(84.11), while (83.6) passes the second.  Let \(n\to\infty\).
The equality extends from the local union to its norm completion
because states and automorphisms are norm continuous.
\end{proof}

Clustering does not imply invariance without (84.11), and invariance
does not imply clustering.

\subsection{OS time reconstruction in one channel}

Assume the populated reflection-positive gate reconstructs a Hilbert
space \({\cal H}_{\rm OS}\), a unit vacuum vector \(\Omega\), and a
positive self-adjoint Euclidean-time generator \(H\) with
\[
                              H\Omega=0.
\]
Let \({\cal V}\subset\Omega^\perp\) be a subspace generated densely
in a closed channel \({\cal H}_{\rm ch}\subset\Omega^\perp\) by
zero-mean reflected local observables.  Suppose that for some \(m>0\)
and every \(v\in{\cal V}\), there is \(C_v<\infty\) such that
\begin{equation}
                         0\leq
 \langle v,e^{-tH}v\rangle
 \leq C_ve^{-mt}
 \qquad(t\geq0).                                            \tag{84.13}
\end{equation}

\begin{theorem}[Exponential channel decay implies a spectral
threshold]
\label{thm:v12v84-spectral}
Under (84.13),
\begin{equation}
 {\bf1}_{[0,m)}(H)v=0\qquad(v\in{\cal V}).                   \tag{84.14}
\end{equation}
If \({\cal V}\) is dense in the invariant channel
\({\cal H}_{\rm ch}\), then
\begin{equation}
 \operatorname{spec}(H|_{{\cal H}_{\rm ch}})
 \subset[m,\infty).                                         \tag{84.15}
\end{equation}
\end{theorem}

\begin{proof}
Let \(\sigma_v\) be the positive spectral measure of \(H\) in the
vector \(v\).  Then
\[
 \langle v,e^{-tH}v\rangle
 =\int_{[0,\infty)}e^{-t\lambda}\,d\sigma_v(\lambda).        \tag{84.16}
\]
Fix \(0\leq a<m\).  If
\(\sigma_v([0,a])>0\), then
\[
 \langle v,e^{-tH}v\rangle
 \geq e^{-at}\sigma_v([0,a]),
\]
which contradicts (84.13) after multiplication by \(e^{mt}\) and
letting \(t\to\infty\).  Thus
\(\sigma_v([0,a])=0\) for every \(a<m\), proving (84.14).

The spectral projection \({\bf1}_{[0,m)}(H)\) is bounded.  If it
vanishes on the dense subspace \({\cal V}\), it vanishes on the
closed channel.  This is equivalent to (84.15), assuming the channel
is invariant under the spectral calculus of \(H\).
\end{proof}

This theorem gives a gap only in the channel for which (84.13) is
proved.  A photon, graviton, or other massless channel is expected to
have a power-law or zero-threshold spectrum and is not covered by an
exponential assertion.

\subsection{Unbounded correlations}

If \(A\) or \(B\) is an unbounded smeared Einstein or stress writer,
Theorem \ref{thm:v12v84-clustering} cannot be applied directly in a
\(C^*\)-algebra.  First truncate the writers, pass the bounded
clustering estimate, and then remove the truncation using the
superlinear product moments of V79.  Uniform \(L^1\) bounds alone do
not justify that removal.

\begin{assessment}[Status after V84]
\label{ass:v12v84-status}
Uniform cutoff connected-correlation envelopes now pass to the V83
quasi-local state under explicit retained-translation convergence.
Power-law and exponential rates are preserved.  If reflection
positivity reconstructs Euclidean time, exponential diagonal decay
on a dense specified channel gives a spectral threshold in that
channel, without imposing a false global gap on massless SM--Einstein
sectors.
\end{assessment}

\begin{firewall}[Claim boundary after V84]
\label{fire:v12v84-claim}
No physical translation or relational-separation action,
retained-translate estimate (84.2), cutoff clustering envelope,
translation-invariance residual, or OS reconstruction has been
populated for the full RPESM system.  The spectral theorem is
conditional and channel-specific.  It proves neither a global mass
gap nor the required photon and graviton infrared structure.
The full SM--Einstein continuum is not yet constructed.
\end{firewall}

\subsection{Parent record}

This V84 note appends to the validated V83 file with record
\[
\begin{array}{c|c}
\text{lines}&23434\\
\text{bytes}&791992\\
\text{SHA--256}&
\texttt{00891D41EDCDB7F64CCD23457C706E9C7284552164112C0985E0D22FE03A9603}.
\end{array}
\]

% =============================================================================
% END APPENDED TWELFTH-CYCLE V84 MATERIAL
% =============================================================================

% =============================================================================
% BEGIN APPENDED TWELFTH-CYCLE V85 MATERIAL
% =============================================================================

\section{Twelfth-cycle V85: mandatory companion audit}
\label{sec:v12v85}

This multiple-of-five checkpoint inspects the live companion files
before the regulator-independence step.

\subsection{Live records and the duplicate YM alias}

The active-looking companion files are
\[
\begin{array}{c|r|r|c}
\text{record}&\text{lines}&\text{bytes}&\text{SHA--256}\\ \hline
\text{Yang--Mills V60}&22478&765479&
\texttt{6F145D0871A5EC33C78DE20B410FB1F80458CDB31A7D85053C84EFCCCFAF5CAC}\\
\text{Yang--Mills V61}&22478&765479&
\texttt{6F145D0871A5EC33C78DE20B410FB1F80458CDB31A7D85053C84EFCCCFAF5CAC}\\
\text{Navier--Stokes V26}&5319&278283&
\texttt{82C887F8C2C69E4F28FAD7C3DC8DE5B9099CB5A4BF431EBA1FFF3E91935978AD}.
\end{array}                                                   \tag{85.1}
\]
The two Yang--Mills files are byte-identical and both end with the
V60 section and V60 claim boundary.  Accordingly V60 is the latest
substantive companion result; the V61 filename supplies no additional
mathematical content.  Neither file is altered here.  Navier--Stokes
is unchanged.

\subsection{Substantive YM V60 result}

YM V60 completes a coefficient-exact arbitrary-real-momentum
one-loop response interface.  It combines exact Laurent Hessian
jets, quartic banks, both orientations of the cubic banks, inverse
jets, and tadpole--bubble contractions.  Its exact corner
specialization reproduces the earlier rational cards.  General
momenta still use floating matrix inversion, and no interval
Brillouin integral, continuum response sign, nonperturbative
remainder, continuum measure, or mass gap is claimed.

The transferable lesson is again one of comparison scope.  A
regulator cannot be declared equivalent to another because finitely
many sampled responses agree.  Equality of continuum states must be
tested on a measure-determining or norm-dense physical observable
ledger after both schemes have been mapped into the same retained
local algebra.

NS V26 supplies no new result.  Its structured-norm warning remains
consistent with this rule: a comparison estimate is valid only on
the common physical subspace for which both reconstruction maps have
been proved.

\begin{decision}[Direction after the V85 audit]
\label{dec:v12v85-next}
Retain V69--V84.  At V86, compare two arbitrary regulator schemes
after their cutoff observables have been reconstructed in the common
nested algebras of V83.  Prove that:
\begin{enumerate}
\item vanishing cross-scheme discrepancy on one dense local ledger
      transfers convergence from either scheme to the other;
\item if both limits exist, the same discrepancy proves equality of
      the continuum states;
\item a weighted summable cross ledger gives quantitative local
      bounds;
\item parameter matching and reconstruction errors remain explicit
      inputs and cannot be inferred from finite samples.
\end{enumerate}
This is a regulator-independence criterion, not a claim that the
physical RPESM schemes satisfy it.
\end{decision}

\begin{assessment}[Status after V85]
\label{ass:v12v85-status}
The mandatory audit is complete.  The duplicate YM V61 alias has
been identified without modifying the shared programme.  YM V60
strengthens the reason to formulate universality on the full dense
physical ledger.  No companion numerical result enters the
SM--Einstein estimates.
\end{assessment}

\begin{firewall}[Claim boundary after V85]
\label{fire:v12v85-claim}
No cross-regulator reconstruction, counterterm matching, parameter
matching, or universality estimate has been populated.  The YM
arbitrary-momentum response is not a gravitational or full-SM
continuum result.  The full SM--Einstein continuum is not yet
constructed.
\end{firewall}

\subsection{Parent record}

This V85 note appends to the validated V84 file with record
\[
\begin{array}{c|c}
\text{lines}&23716\\
\text{bytes}&801787\\
\text{SHA--256}&
\texttt{B0183456BA29C078C44FBD5A4E25AA59978B88ED01CB1AD83962B98D33035B37}.
\end{array}
\]

% =============================================================================
% END APPENDED TWELFTH-CYCLE V85 MATERIAL
% =============================================================================

% =============================================================================
% BEGIN APPENDED TWELFTH-CYCLE V86 MATERIAL
% =============================================================================

\section{Twelfth-cycle V86: cross-regulator observable matching and
universality of the continuum state}
\label{sec:v12v86}

V83 constructs a continuum state for one regulator from its local
ledger.  Universality requires more: two different regulators,
blocking rules, or contour implementations must define the same
state after their observables and renormalised parameters have been
matched.  This note proves the comparison theorem on the common
physical local net.  It does not infer the matching maps or their
error bounds from finite samples.

\subsection{Two asynchronous cofinal schedules}

Let
\[
 {\mathfrak A}_1\subset{\mathfrak A}_2\subset\cdots,\qquad
 {\mathfrak A}_{\rm ql}
 =\overline{\bigcup_m{\mathfrak A}_m}
\]
be the common physical net of V83, and let \(d_m\) be the
weak-star state metric (83.9).  Consider two regulator schemes,
denoted \(S\) and \(T\).  After applying their proposed physical
reconstruction maps, suppose their matched cofinal schedules give
states
\[
 \omega_j^S\text{ on }{\mathfrak A}_{p_j},\qquad
 \omega_j^T\text{ on }{\mathfrak A}_{q_j},                  \tag{86.1}
\]
where \(p_j,q_j\to\infty\).  For \(m\leq p_j\) or \(q_j\), write
\(\omega_j^{S,(m)}\) and \(\omega_j^{T,(m)}\) for the restrictions
to \({\mathfrak A}_m\).

Choose integers
\[
 M_j\leq\min\{p_j,q_j\},\qquad M_j\longrightarrow\infty,     \tag{86.2}
\]
and define the matched cross-ledger discrepancy
\begin{equation}
 X_j:=
 2^{-M_j}
 +\sum_{m=1}^{M_j}2^{-m}
 d_m(\omega_j^{S,(m)},\omega_j^{T,(m)}).                    \tag{86.3}
\end{equation}

\begin{lemma}[Every fixed local discrepancy is visible]
\label{lem:v12v86-visible}
For every fixed \(m\) and every \(j\) with \(M_j\geq m\),
\begin{equation}
 d_m(\omega_j^{S,(m)},\omega_j^{T,(m)})
 \leq2^mX_j.                                                 \tag{86.4}
\end{equation}
Consequently \(X_j\to0\) implies vanishing cross-scheme discrepancy
on every fixed local algebra.
\end{lemma}

\begin{proof}
The \(m\)-th summand in (86.3) is nonnegative and is at most the
whole sum.  Multiply by \(2^m\).  Equation (86.2) ensures that every
fixed \(m\) eventually occurs.
\end{proof}

The tail \(2^{-M_j}\) prevents a comparison on one permanently fixed
finite head from being labeled a full physical comparison.

\subsection{Transfer of convergence}

\begin{theorem}[One-scheme convergence transfers to the other]
\label{thm:v12v86-transfer}
Assume \(X_j\to0\).  Suppose scheme \(S\) has the quasi-local
continuum state \(\omega_\infty^S\), meaning
\begin{equation}
 \omega_j^{S,(m)}
 \longrightarrow\omega_\infty^S|_{{\mathfrak A}_m}
 \quad\text{weak-star for every fixed }m.                   \tag{86.5}
\end{equation}
Then scheme \(T\) also converges on every fixed local algebra and
its V83 continuum state exists.  Moreover,
\begin{equation}
                              \boxed{\omega_\infty^T
                              =\omega_\infty^S}
                              \quad\text{on }{\mathfrak A}_{\rm ql}.
                                                                    \tag{86.6}
\end{equation}
Quantitatively, for \(M_j\geq m\),
\begin{equation}
 d_m(\omega_j^{T,(m)},
          \omega_\infty^S|_{{\mathfrak A}_m})
 \leq
 d_m(\omega_j^{S,(m)},
          \omega_\infty^S|_{{\mathfrak A}_m})
 +2^mX_j.                                                    \tag{86.7}
\end{equation}
\end{theorem}

\begin{proof}
The triangle inequality and Lemma
\ref{lem:v12v86-visible} give (86.7).  Both terms on its right tend
to zero by (86.5) and \(X_j\to0\).  Hence the \(T\)-restrictions
converge to the \(S\)-restriction for every fixed \(m\).  Theorem
\ref{thm:v12v83-local-limit} constructs \(\omega_\infty^T\).
The two states agree on the dense local union, so norm continuity
gives (86.6).
\end{proof}

\begin{corollary}[Equality of independently constructed limits]
\label{cor:v12v86-equality}
If both schemes already satisfy either V83 convergence criterion and
\(X_j\to0\), their continuum states are equal.
\end{corollary}

\begin{proof}
Apply Theorem \ref{thm:v12v86-transfer} with either scheme as the
reference.
\end{proof}

If \(\sum_jX_j<\infty\), then \(X_j\to0\), and (86.7) gives a
summable matched comparison term at each fixed local scale.  More
generally, any certified bound \(X_j\leq R_j\to0\) gives the explicit
local error \(2^mR_j\).

\subsection{A countable scalar version}

Let \(A_{m,k}\) be the dense self-adjoint unit-ball ledger of V83.
It is enough that for every fixed \(m,k\),
\begin{equation}
 |\omega_j^S(A_{m,k})-\omega_j^T(A_{m,k})|
 \longrightarrow0.                                         \tag{86.8}
\end{equation}
If one scheme converges, then the other has the same continuum state.

\begin{proof}
Convergence of scheme \(S\) and (86.8) make every scalar
\(T\)-ledger entry converge to the corresponding \(S\)-limit.
The qualitative diagonal criterion
\ref{thm:v12v83-qualitative} constructs the \(T\)-state.  Equality
on the dense ledger gives equality on each local algebra and then on
the quasi-local completion.
\end{proof}

This criterion contains countably many observables.  Agreement on
any fixed finite subset does not imply (86.8).

\subsection{Separating reconstruction and parameter matching}

A cross-regulator estimate normally contains two different errors.
Let \(\lambda_j^S,\lambda_j^T\in\mathbb R^r\) be the renormalised
relevant-parameter vectors.  For each fixed \(m\), suppose there is a
comparison family of states
\(\chi_{j,m}(\lambda)\) on \({\mathfrak A}_m\) such that
\begin{align}
 d_m(\omega_j^{S,(m)},\chi_{j,m}(\lambda_j^S))
 &\leq r_{j,m}^S,                                           \tag{86.9}\\
 d_m(\omega_j^{T,(m)},\chi_{j,m}(\lambda_j^T))
 &\leq r_{j,m}^T,                                           \tag{86.10}\\
 d_m(\chi_{j,m}(\lambda),\chi_{j,m}(\lambda'))
 &\leq L_{j,m}\|\lambda-\lambda'\|.                         \tag{86.11}
\end{align}

\begin{proposition}[Reconstruction--parameter decomposition]
\label{prop:v12v86-parameters}
Under (86.9)--(86.11),
\begin{equation}
 d_m(\omega_j^{S,(m)},\omega_j^{T,(m)})
 \leq
 r_{j,m}^S+r_{j,m}^T
 +L_{j,m}\|\lambda_j^S-\lambda_j^T\|.                       \tag{86.12}
\end{equation}
Thus the cross ledger vanishes if the weighted reconstruction errors
and the weighted parameter mismatch vanish on a growing local head.
\end{proposition}

\begin{proof}
Insert
\(\chi_{j,m}(\lambda_j^S)\) and
\(\chi_{j,m}(\lambda_j^T)\) between the two states, then apply the
triangle inequality and (86.9)--(86.11).
\end{proof}

The decomposition prevents a counterterm or reconstruction error
from being hidden inside a claimed parameter match.  It also shows
that bare-parameter equality is irrelevant unless it implies
convergence of the renormalised physical vectors appearing in
(86.12).

\subsection{Consequences of equality of the state}

If the two schemes use the same identified quasi-local algebra and
(86.6) holds, then:
\begin{enumerate}
\item every bounded local correlation and connected correlation
      agrees;
\item every common OS quadratic form and bounded Ward functional
      agrees;
\item their GNS triples are unitarily equivalent by the uniqueness
      theorem for the GNS representation of one state;
\item any common translation automorphisms have the same clustering
      and channel spectral consequences of V84.
\end{enumerate}
Unbounded Einstein and stress writers agree only after their two
schemes share the V79 uniform-integrability stock and their
reconstructed writers converge to the same affiliated observable.

\subsection{Scope of universality}

Theorem \ref{thm:v12v86-transfer} is a uniqueness theorem after
physical identification.  It does not construct the common
observable net or prove that two bare actions flow to it.  A complete
RPESM universality proof must populate:
\begin{enumerate}
\item reconstruction maps into the same
      \({\mathfrak A}_m\);
\item tuning of all relevant couplings, masses, gauge parameters, and
      gravitational constants;
\item vanishing irrelevant-operator and contour discrepancies;
\item a growing ledger estimate such as (86.3), not finitely many
      matching conditions.
\end{enumerate}

\begin{assessment}[Status after V86]
\label{ass:v12v86-status}
Regulator independence now has an exact cofinal criterion.  A growing
weighted cross ledger transfers local convergence from one scheme to
another and proves equality of their quasi-local continuum states,
with explicit local error bounds.  Reconstruction errors and
renormalised-parameter mismatch enter separately, so finite response
agreement cannot be mistaken for universality.
\end{assessment}

\begin{firewall}[Claim boundary after V86]
\label{fire:v12v86-claim}
No pair of physical RPESM regulators has been reconstructed into a
common local net, and no bounds in (86.9)--(86.12) have been
populated.  Equality of continuum states is conditional on the full
growing ledger and does not follow from the YM V60 finite response
oracle or from equality of bare parameters.  Existence, positivity,
nontriviality, clustering, and parameter matching remain physical
acquisition tasks.  The full SM--Einstein continuum is not yet
constructed.
\end{firewall}

\subsection{Parent record}

This V86 note appends to the validated V85 file with record
\[
\begin{array}{c|c}
\text{lines}&23822\\
\text{bytes}&806153\\
\text{SHA--256}&
\texttt{9AD96A1FF95A355A59682450BF9305A785C33F91CDCCBE8072C1528B6732EE87}.
\end{array}
\]

% =============================================================================
% END APPENDED TWELFTH-CYCLE V86 MATERIAL
% =============================================================================

% =============================================================================
% BEGIN APPENDED TWELFTH-CYCLE V87 MATERIAL
% =============================================================================

\section{Twelfth-cycle V87: passage of Einstein--matter residuals and
Schwinger--Dyson identities}
\label{sec:v12v87}

V79 passes unbounded smeared writers and V83 supplies one quasi-local
state, but neither result by itself states the Einstein--matter
equation.  This note formulates three distinct continuum gates:
vanishing of the mean residual, vanishing of the residual inside all
admissible correlations, and concentration of the law on
distributional solutions.  The implications are proved separately
so that a mean equation is not mistaken for a pathwise constraint.

\subsection{Smeared residual writers}

Let \(Y\) be the physical configuration space with
\(\mu_n\Rightarrow\mu\).  Let \({\cal T}\) be a separable locally
convex space of smooth compactly supported symmetric test tensors
and choose a countable dense set
\[
                         \{\varphi_j:j\geq1\}\subset{\cal T}. \tag{87.1}
\]
For each cutoff, let
\[
                         {\cal R}_n(\varphi):Y\to\mathbb R
\]
be the reconstructed smeared Einstein--matter residual.  Formally
its intended limit is
\begin{equation}
 {\cal R}(\varphi)
 =
 \left\langle
 G(g)+\Lambda g-8\pi G_{\!N}\,T_{\rm ren},
 \varphi\right\rangle,                                     \tag{87.2}
\end{equation}
with all gauge fixing, conditional quantum expectation, and
renormalisation already included in the definition of the writer.
Assume for each fixed \(\varphi\) that
\begin{equation}
 {\cal R}_n(\varphi)\longrightarrow{\cal R}(\varphi)
 \quad\text{uniformly on every physical compact subset of }Y.       \tag{87.3}
\end{equation}

\subsection{The mean Einstein equation}

\begin{theorem}[Mean residual passage]
\label{thm:v12v87-mean}
Fix \(\varphi\in{\cal T}\).  Suppose (87.3) holds and for some
\(p>1\),
\begin{equation}
 \sup_n\int_Y|{\cal R}_n(\varphi)|^p\,d\mu_n<\infty.         \tag{87.4}
\end{equation}
If
\begin{equation}
 \left|\int_Y{\cal R}_n(\varphi)\,d\mu_n\right|
 \leq\varepsilon_n(\varphi),
 \qquad\varepsilon_n(\varphi)\longrightarrow0,              \tag{87.5}
\end{equation}
then
\begin{equation}
                         \int_Y{\cal R}(\varphi)\,d\mu=0.     \tag{87.6}
\end{equation}
\end{theorem}

\begin{proof}
The superlinear moment theorem
\ref{thm:v12v79-moment}, applied to the writers
\({\cal R}_n(\varphi)\), gives
\[
 \int{\cal R}_n(\varphi)\,d\mu_n
 \longrightarrow
 \int{\cal R}(\varphi)\,d\mu.
\]
Equation (87.5) makes the left limit zero.
\end{proof}

If (87.4)--(87.5) hold on the dense tests (87.1) and the limiting
mean functional
\[
 \varphi\longmapsto\int{\cal R}(\varphi)\,d\mu              \tag{87.7}
\]
is continuous on \({\cal T}\), then (87.6) extends to every
\(\varphi\in{\cal T}\).

\begin{proof}
A continuous linear functional vanishing on a dense subset vanishes
everywhere.
\end{proof}

Equation (87.6) is only an equation in expectation.  Cancellation
between configurations with opposite residuals can make it true even
when no individual configuration solves (87.2).

\subsection{Schwinger--Dyson identities}

Let \(F_n,F:Y\to\mathbb C\) be reconstructed cylinder insertions,
with \(F_n\to F\) locally uniformly.  Let
\({\cal D}_{n,\varphi}F_n\) be the cutoff contact or functional-
derivative term and suppose
\[
 {\cal D}_{n,\varphi}F_n
 \longrightarrow{\cal D}_{\varphi}F
\quad\text{locally uniformly}.                              \tag{87.8}
\]
Assume moment exponents \(p,q>1\) satisfy
\begin{align}
 \sup_n\int|{\cal R}_n(\varphi)|^p\,d\mu_n&<\infty,          \tag{87.9}\\
 \sup_n\int|F_n|^q\,d\mu_n&<\infty,
 &{1\over p}+{1\over q}&<1,                                \tag{87.10}
\end{align}
and assume a superlinear moment for
\({\cal D}_{n,\varphi}F_n\) sufficient for V79.

\begin{theorem}[Correlated residual passage]
\label{thm:v12v87-SD}
If the cutoff Schwinger--Dyson residual obeys
\begin{equation}
 \left|
 \int_Y\left\{
 {\cal R}_n(\varphi)F_n
 -{\cal D}_{n,\varphi}F_n
 \right\}\,d\mu_n
 \right|
 \leq\epsilon_n(\varphi,F),
 \qquad\epsilon_n(\varphi,F)\longrightarrow0,               \tag{87.11}
\end{equation}
then
\begin{equation}
 \boxed{\displaystyle
 \int_Y{\cal R}(\varphi)F\,d\mu
 =
 \int_Y{\cal D}_{\varphi}F\,d\mu.}                          \tag{87.12}
\end{equation}
\end{theorem}

\begin{proof}
The product theorem \ref{thm:v12v79-products} passes the first
integral in (87.11), using the strict reciprocal-exponent margin in
(87.10).  The superlinear single-writer theorem passes the contact
term.  Let \(n\to\infty\) in (87.11).
\end{proof}

The contact term in (87.12) is essential.  For an action-derived
quantum measure, the correct equation inside correlations is
generally Schwinger--Dyson, not
\(\int{\cal R}(\varphi)F\,d\mu=0\) for arbitrary \(F\).

If (87.12) is first proved on countable dense test and cylinder
families, it extends to their completed domains only after the two
sides are shown continuous in the corresponding test norms.

\subsection{Mean-square concentration on distributional solutions}

The strongest of the three gates is support on zero residual.

\begin{theorem}[Mean-square residual closure]
\label{thm:v12v87-L2}
Fix \(\varphi\in{\cal T}\) and assume (87.3).  If
\begin{equation}
                         \int_Y
 |{\cal R}_n(\varphi)|^2\,d\mu_n
 \longrightarrow0,                                         \tag{87.13}
\end{equation}
then
\begin{equation}
                         {\cal R}(\varphi)=0
                         \quad\mu\text{-almost surely}.      \tag{87.14}
\end{equation}
\end{theorem}

\begin{proof}
For every bounded continuous
\(\chi:\mathbb R\to\mathbb R\), the functions
\(\chi({\cal R}_n(\varphi))\) converge locally uniformly to
\(\chi({\cal R}(\varphi))\).  The varying bounded mapping lemma
\ref{lem:v12v79-bounded} therefore gives
\[
 ({\cal R}_n(\varphi))_\#\mu_n
 \Rightarrow
 ({\cal R}(\varphi))_\#\mu.                                 \tag{87.15}
\]
Portmanteau for the nonnegative lower-semicontinuous function
\(t\mapsto t^2\) yields
\[
 \int|{\cal R}(\varphi)|^2\,d\mu
 \leq\liminf_n\int|{\cal R}_n(\varphi)|^2\,d\mu_n=0.
\]
This proves (87.14).
\end{proof}

\begin{corollary}[Almost-sure distributional constraint]
\label{cor:v12v87-distribution}
Suppose (87.13) holds for every \(\varphi_j\) in (87.1).
Assume that for \(\mu\)-almost every \(y\), the map
\begin{equation}
                         \varphi\longmapsto
                         {\cal R}(y;\varphi)                 \tag{87.16}
\end{equation}
is a continuous linear functional on \({\cal T}\).  Then, outside
one \(\mu\)-null set,
\begin{equation}
                         {\cal R}(y;\varphi)=0
                         \quad\text{for every }\varphi\in{\cal T}.
                                                                    \tag{87.17}
\end{equation}
\end{corollary}

\begin{proof}
For each \(j\), Theorem \ref{thm:v12v87-L2} gives a full-measure set
on which \({\cal R}(\varphi_j)=0\).  Intersect these countably many
sets and the full-measure set on which (87.16) is continuous.
At any point in the intersection, the continuous functional
(87.16) vanishes on the dense set (87.1), hence everywhere.
\end{proof}

\subsection{Quantum meaning of the three gates}

For a coupled quantum gravitational law, an almost-sure classical
Einstein equation may be the wrong target: quantum fluctuations
usually satisfy Schwinger--Dyson identities with contact terms.
Accordingly:
\begin{enumerate}
\item (87.6) is the weakest mean semiclassical equation;
\item (87.12) is the correlation-level quantum equation;
\item (87.17) is a support theorem appropriate only if the intended
      construction truly concentrates on classical constrained
      solutions.
\end{enumerate}
None implies either of the others without additional hypotheses.

The stress writer in (87.2) must be the renormalised writer on the
same regulator schedule as the metric writer.  V79 supplies passage
after smearing and moment control; it does not manufacture the
counterterms or prove conservation.

\subsection{Compatibility with universality}

If two regulators satisfy V86 and their limiting state is the same,
then any common limiting residual and contact writers satisfying the
V79 moment hypotheses give the same identities (87.6) and (87.12).
For unbounded writers, equality of the bounded quasi-local state
alone is insufficient unless both schemes also have the same
affiliated-writer reconstruction and uniform-integrability stock.

\begin{assessment}[Status after V87]
\label{ass:v12v87-status}
The Einstein--matter equation is now separated into three rigorous
continuum gates.  Superlinear moments pass mean residuals and
Schwinger--Dyson identities with their contact terms; vanishing
cutoff mean-square residuals concentrate the limit on
distributional solutions through a countable dense test family.
The proof prevents a vanishing mean from being misreported as an
almost-sure field equation.
\end{assessment}

\begin{firewall}[Claim boundary after V87]
\label{fire:v12v87-claim}
No renormalised RPESM stress tensor, residual writer, contact term,
moment estimate, or vanishing residual stock has been populated.
The almost-sure theorem is not asserted to be the correct quantum
gravity equation.  Conservation, anomalies, reflection positivity,
local covariance, nontriviality, clustering, and parameter matching
remain separate physical gates.  The full SM--Einstein continuum is
not yet constructed.
\end{firewall}

\subsection{Parent record}

This V87 note appends to the validated V86 file with record
\[
\begin{array}{c|c}
\text{lines}&24083\\
\text{bytes}&816007\\
\text{SHA--256}&
\texttt{F5459AE4C1553C7AE93FA6864B85F6386355AD2C8BDD335F29ABE950333B3F4D}.
\end{array}
\]

% =============================================================================
% END APPENDED TWELFTH-CYCLE V87 MATERIAL
% =============================================================================

% =============================================================================
% BEGIN APPENDED TWELFTH-CYCLE V88 MATERIAL
% =============================================================================

\section{Twelfth-cycle V88: deriving the V81 quantum energy tail from
the V62 elliptic heat stock}
\label{sec:v12v88}

V81 assumes a uniform compact-resolvent reference-energy moment.
V62 already proves uniform ellipticity and heat traces for transported
second-order conditional Hamiltonians, but does not write the
specific implication needed by V81.  This note supplies it.  A fixed
Rayleigh test gives a uniform partition floor, the heat trace at half
the inverse temperature controls the Gibbs energy, and the V62 form
comparison transfers that energy to the reference operator.

\subsection{Positive shifted Hamiltonians}

Retain the V62 setting.  Let
\[
                         K_x:=L_x+b_2I,                      \tag{88.1}
\]
where \(b_2\) is the constant in (62.8).  Then
\begin{equation}
                         K_x\succeq a_2\Lambda\succeq0.      \tag{88.2}
\end{equation}
The normalized Gibbs state is unchanged by this scalar shift:
\begin{equation}
 \Gamma_\beta(K_x)
 ={e^{-\beta K_x}\over\operatorname{Tr}e^{-\beta K_x}}
 =
 {e^{-\beta L_x}\over\operatorname{Tr}e^{-\beta L_x}}
 =\Gamma_\beta(L_x).                                       \tag{88.3}
\end{equation}

Choose one fixed smooth unit vector \(u_0\in{\cal H}_0\).
The uniform coefficient bounds in (62.2) imply
\begin{equation}
                         \langle u_0,K_xu_0\rangle
                         \leq\kappa_0<\infty                 \tag{88.4}
\end{equation}
with \(\kappa_0\) independent of \(x\).

\begin{lemma}[Uniform partition floor]
\label{lem:v12v88-floor}
For every \(\beta>0\),
\begin{equation}
 Z_x(\beta):=\operatorname{Tr}e^{-\beta K_x}
 \geq e^{-\beta\kappa_0}.                                  \tag{88.5}
\end{equation}
\end{lemma}

\begin{proof}
The min--max principle and (88.4) give
\[
                         \lambda_1(K_x)\leq\kappa_0.
\]
The first eigenvalue alone contributes
\(e^{-\beta\lambda_1(K_x)}\) to the trace, proving (88.5).
\end{proof}

This floor uses an upper Rayleigh bound, whereas coercivity supplies
a lower spectral bound.  The two inequalities play different roles.

\subsection{Half-temperature energy estimate}

\begin{lemma}[Scalar half-temperature bound]
\label{lem:v12v88-scalar}
For every \(t\geq0\) and \(\beta>0\),
\begin{equation}
                         t e^{-\beta t}
 \leq {2\over e\beta}e^{-\beta t/2}.                        \tag{88.6}
\end{equation}
\end{lemma}

\begin{proof}
Write
\[
 t e^{-\beta t}
 =\left(t e^{-\beta t/2}\right)e^{-\beta t/2}.
\]
The maximum of \(t e^{-\beta t/2}\) on \([0,\infty)\) occurs at
\(t=2/\beta\) and equals \(2/(e\beta)\).
\end{proof}

Let
\begin{equation}
 C_{\beta/2}
 :=\sup_x\operatorname{Tr}e^{-(\beta/2)K_x}.                \tag{88.7}
\end{equation}
It is finite by (62.9), since shifting by \(b_2I\) only multiplies
the heat trace by \(e^{-(\beta/2)b_2}\).

\begin{theorem}[Uniform conditional Gibbs energy]
\label{thm:v12v88-Gibbs-energy}
For every \(x\),
\begin{equation}
 \operatorname{Tr}\{\Gamma_\beta(K_x)K_x\}
 \leq
 {2e^{\beta\kappa_0}\over e\beta}\,C_{\beta/2}.              \tag{88.8}
\end{equation}
Consequently
\begin{equation}
 \operatorname{Tr}\{\Gamma_\beta(L_x)\Lambda\}
 \leq
 {\cal E}_\Lambda
 :=
 {2e^{\beta\kappa_0}\over
  e\beta a_2}\,C_{\beta/2}.                                 \tag{88.9}
\end{equation}
Both bounds are uniform in the background \(x\).
\end{theorem}

\begin{proof}
Functional calculus and Lemma \ref{lem:v12v88-scalar} give
\[
 \operatorname{Tr}(K_xe^{-\beta K_x})
 \leq {2\over e\beta}
 \operatorname{Tr}e^{-(\beta/2)K_x}
 \leq {2\over e\beta}C_{\beta/2}.                           \tag{88.10}
\]
Divide by the partition floor (88.5) to obtain (88.8).
The form inequality (88.2) gives
\(\Lambda\preceq a_2^{-1}K_x\).  Take its expectation in the
positive Gibbs state and use (88.3) and (88.8), proving (88.9).
\end{proof}

\subsection{Averaging over Einstein backgrounds}

Let \(\nu_n\) be any probability law on the admissible background
class, and let the conditional density field be
\[
                         \rho_n(x):=\Gamma_\beta(L_{n,x}),   \tag{88.11}
\]
where all cutoff families obey the same constants in
(62.3), (62.8), (88.4), and (88.7).  After the V61 unitary
identifications, take the V81 reference energy to be
\[
                              H_0:=\Lambda.                  \tag{88.12}
\]

\begin{corollary}[The V81 energy hypothesis]
\label{cor:v12v88-V81}
Under these uniform elliptic hypotheses,
\begin{equation}
 \sup_n\int
 \operatorname{Tr}\{\rho_n(x)H_0\}\,d\nu_n(x)
 \leq{\cal E}_\Lambda.                                      \tag{88.13}
\end{equation}
Hence the high-mode tail (81.13), no-escape theorem V81, and
conditional-density disintegration V82 apply to every joint-state
limit constructed from these Gibbs fibres.
\end{corollary}

\begin{proof}
Equation (88.9) bounds the integrand pointwise and uniformly in
\(n,x\).  Integrate against the probability \(\nu_n\).  The cited
V81--V82 results then apply.
\end{proof}

This implication does not require convergence of the background law
or pointwise convergence of the conditional Gibbs states.  Those
inputs enter existence and identification of the joint limit, not
the energy containment.

\subsection{Higher reference moments}

For \(r>0\), the scalar maximum
\[
 \sup_{t\geq0}t^re^{-\beta t/2}
 =\left({2r\over e\beta}\right)^r                            \tag{88.14}
\]
gives
\begin{equation}
 \operatorname{Tr}\{\Gamma_\beta(K_x)K_x^r\}
 \leq
 e^{\beta\kappa_0}
 \left({2r\over e\beta}\right)^r C_{\beta/2}.                \tag{88.15}
\end{equation}
When operator monotonicity or an explicit relative-power inequality
justifies
\begin{equation}
                         \Lambda^r\preceq C_r(I+K_x^r),      \tag{88.16}
\end{equation}
the same calculation gives a uniform \(r\)-th reference-energy
moment.  For \(r>1\), (88.16) is not a consequence of (88.2) alone,
because the map \(t\mapsto t^r\) is not operator monotone.  It must be
proved separately, as in the relative-power philosophy of V62.

\begin{proof}
Apply
\[
 t^re^{-\beta t}
 =\left(t^re^{-\beta t/2}\right)e^{-\beta t/2}
 \leq\left({2r\over e\beta}\right)^re^{-\beta t/2}
\]
by functional calculus, trace, divide by (88.5), and then use
(88.16) when available.
\end{proof}

\subsection{Physical scope}

The V62 operator \(L_x\) is a finite-rank-bundle, second-order
elliptic Hamiltonian on a closed manifold.  A many-body Fock
Hamiltonian, an indefinite ghost sector, a chiral determinant
amplitude, or a metric contour is not automatically of this form.
For a tensor product of independently controlled positive fibres,
energy moments add under the corresponding sum Hamiltonian; an
interacting Fock theory needs its own relative-form and stability
bounds.

The partition floor (88.5) is uniform only while the coefficient
bounds and the fixed Rayleigh test survive the regulator.  If the
rank, volume, or number of independent degrees of freedom grows,
the normalized local energy density rather than the extensive total
energy may be the quantity with a uniform limit.  V81 requires a
specified compact-resolvent reference whose tail estimate is
dimensionally compatible with that growth.

\begin{assessment}[Status after V88]
\label{ass:v12v88-status}
The abstract V81 quantum no-escape hypothesis is now derived from the
V62 elliptic stock for one transported conditional Gibbs fibre.
Uniform coefficient bounds give a partition floor, the
half-temperature heat trace gives an explicit Gibbs-energy ceiling,
and elliptic form comparison transfers it to the compact-resolvent
reference operator before averaging over gravitational backgrounds.
\end{assessment}

\begin{firewall}[Claim boundary after V88]
\label{fire:v12v88-claim}
The physical full RPESM conditional Hamiltonian has not been proved
to satisfy the uniform V62 coefficient class, fixed-rank heat bound,
or Rayleigh ceiling across its regulator.  Ghost, determinant,
many-body, and contour sectors are not covered by the positive
single-particle Gibbs argument.  Extensivity and local energy-density
renormalisation remain open.  Thus (88.13) is a populated logical
handoff under explicit elliptic hypotheses, not yet a full physical
energy estimate.  The full SM--Einstein continuum is not
constructed.
\end{firewall}

\subsection{Parent record}

This V88 note appends to the validated V87 file with record
\[
\begin{array}{c|c}
\text{lines}&24368\\
\text{bytes}&826062\\
\text{SHA--256}&
\texttt{545C029D4BC5C5BBB28AF19DE78174229BCFF21B89ACC2E9604621BD7C6A5548}.
\end{array}
\]

% =============================================================================
% END APPENDED TWELFTH-CYCLE V88 MATERIAL
% =============================================================================

% =============================================================================
% BEGIN APPENDED TWELFTH-CYCLE V89 MATERIAL
% =============================================================================

\section{Twelfth-cycle V89: local energy compactness and projective
normality in the infinite-volume limit}
\label{sec:v12v89}

V88 derives a uniform reference-energy moment for one elliptic Gibbs
fibre, but warns that total energy is extensive when the physical
volume or number of degrees of freedom grows.  The correct
quasi-local requirement is weaker: every fixed local algebra has its
own cutoff-uniform compact-resolvent moment.  This note proves that
such regionwise bounds give normality on every local quantum core,
with no uniform bound in the region index and no claim of a global
density matrix.

\subsection{Local compact-operator cores}

Retain the nested system
\[
 {\mathfrak A}_1\subset{\mathfrak A}_2\subset\cdots
\]
and its V83 limit state \(\omega_\infty\).  For each region index
\(m\), suppose there is a compact physical background space \(K_m\),
a separable Hilbert space \({\cal H}_m\), and a unital embedding
\begin{equation}
 \jmath_m:
 {\mathfrak B}_m:=
 C(K_m;\widetilde{\cal K}({\cal H}_m))
 \longrightarrow{\mathfrak A}_m.                            \tag{89.1}
\end{equation}
Let \(H_m\geq0\) be self-adjoint on \({\cal H}_m\) with compact
resolvent.  Write its eigenvalues and spectral projections as
\[
 E_{m,1}\leq E_{m,2}\leq\cdots\longrightarrow\infty,\qquad
 P_{m,J},\qquad Q_{m,J}:=I-P_{m,J}.                          \tag{89.2}
\]

For \(n\geq m\), let
\[
                         \Omega_{n,m}
 :=\omega_n|_{\jmath_m({\mathfrak B}_m)}
\]
after identifying the embedded core with \({\mathfrak B}_m\).
Assume it has the cutoff density representation
\begin{equation}
 \Omega_{n,m}(F)
 =
 \int_{K_m}\operatorname{Tr}\{
 \rho_{n,m}(u)F(u)\}\,d\nu_{n,m}(u).                        \tag{89.3}
\end{equation}

\subsection{Regionwise energy is sufficient}

\begin{theorem}[Local energy no escape]
\label{thm:v12v89-local-normal}
Suppose that for every fixed \(m\) there is a finite constant
\({\cal E}_m\), independent of \(n\geq m\), such that
\begin{equation}
 \sup_{n\geq m}\int_{K_m}
 \operatorname{Tr}\{\rho_{n,m}(u)H_m\}\,d\nu_{n,m}(u)
 \leq{\cal E}_m.                                            \tag{89.4}
\end{equation}
Then the limiting restriction
\[
                         \Omega_{\infty,m}
 :=\omega_\infty|_{\jmath_m({\mathfrak B}_m)}
\]
satisfies
\begin{align}
 \Omega_{\infty,m}(1\otimes Q_{m,J})
 &\leq{{\cal E}_m\over E_{m,J+1}},                          \tag{89.5}\\
 \|\Omega_{\infty,m}|_{C(K_m;{\cal K}({\cal H}_m))}\|
 &=1.                                                       \tag{89.6}
\end{align}
Consequently there are a probability measure \(\nu_{\infty,m}\) on
\(K_m\) and a strongly measurable field
\[
 \rho_{\infty,m}:K_m\to{\mathfrak S}_1({\cal H}_m)
\]
such that
\begin{align}
 \rho_{\infty,m}(u)&\succeq0,\qquad
 \operatorname{Tr}\rho_{\infty,m}(u)=1
 \quad\text{almost everywhere},                            \tag{89.7}\\
 \Omega_{\infty,m}(F)
 &=\int_{K_m}\operatorname{Tr}\{
 \rho_{\infty,m}(u)F(u)\}\,d\nu_{\infty,m}(u).              \tag{89.8}
\end{align}
No bound on \(\sup_m{\cal E}_m\) is required.
\end{theorem}

\begin{proof}
For \(n\geq m\), the spectral inequality
\[
                         H_m\geq E_{m,J+1}Q_{m,J}
\]
and (89.4) give
\[
 \Omega_{n,m}(1\otimes Q_{m,J})
 \leq{{\cal E}_m\over E_{m,J+1}}.                           \tag{89.9}
\]
The V83 local convergence passes each fixed bounded element
\(1\otimes Q_{m,J}\) to \(\Omega_{\infty,m}\), proving (89.5).
Let \(J\to\infty\).  The V81 no-escape theorem gives (89.6), and the
V82 disintegration theorem gives (89.7)--(89.8).
Every step fixes \(m\) before taking \(n\to\infty\); the constants
may therefore grow with \(m\).
\end{proof}

\begin{corollary}[Local V88 acquisition]
\label{cor:v12v89-V88}
For a fixed region \(m\), suppose its transported conditional
Hamiltonians obey the V62 elliptic hypotheses with constants allowed
to depend on \(m\) but uniform in \(n\geq m\) and in the local
background.  Then V88 supplies a finite
\({\cal E}_m\) in (89.4).  Hence that local core is normal in the
continuum.
\end{corollary}

\begin{proof}
Apply Corollary \ref{cor:v12v88-V81} in the fixed region and then
Theorem \ref{thm:v12v89-local-normal}.
\end{proof}

This use of V88 avoids a false demand that one extensive total-energy
constant remain bounded as \(m\to\infty\).

\subsection{Projective consistency of local density fields}

The state restrictions are automatically compatible:
\begin{equation}
 \Omega_{\infty,\ell}|_{{\mathfrak B}_m}
 =\Omega_{\infty,m}\qquad(m<\ell),                           \tag{89.10}
\end{equation}
whenever the core embeddings commute with the inclusions in the
physical net.  The next proposition spells out what this means for
density fields in a standard tensor-product localization.

Assume for \(m<\ell\):
\begin{enumerate}
\item a continuous projection \(p_{\ell m}:K_\ell\to K_m\);
\item a factorization
\({\cal H}_\ell={\cal H}_m\otimes{\cal H}_{\ell\setminus m}\) with \(\dim{\cal H}_{\ell\setminus m}<\infty\);
\item the embedding
\[
 F\longmapsto
 \bigl[v\mapsto
 F(p_{\ell m}v)\otimes I_{\ell\setminus m}\bigr];            \tag{89.11}
\]
\item \((p_{\ell m})_\#\nu_{\infty,\ell}=\nu_{\infty,m}\).
\end{enumerate}
Disintegrate
\begin{equation}
 d\nu_{\infty,\ell}(v)
 =d\kappa_{\ell m}^{\,u}(v)\,d\nu_{\infty,m}(u),
 \qquad p_{\ell m}(v)=u.                                    \tag{89.12}
\end{equation}

\begin{proposition}[Conditional partial-trace consistency]
\label{prop:v12v89-partial-trace}
Under (89.10)--(89.12),
\begin{equation}
 \boxed{\displaystyle
 \rho_{\infty,m}(u)
 =
 \int_{p_{\ell m}^{-1}(u)}
 \operatorname{Tr}_{\ell\setminus m}
 \rho_{\infty,\ell}(v)\,
 d\kappa_{\ell m}^{\,u}(v)}
                                                               \tag{89.13}
\end{equation}
for \(\nu_{\infty,m}\)-almost every \(u\).  The integral is a
Bochner integral in trace class.
\end{proposition}

\begin{proof}
The partial trace is a positive trace-norm contraction and maps a
density operator to a density operator.  Strong measurability of the
large-region field and standard properties of disintegration make
the right side of (89.13) a strongly measurable positive trace-one
field.

For \(f\in C(K_m)\) and \(A\in{\cal K}({\cal H}_m)\), equations
(89.8), (89.11), and (89.12) give
\[
 \begin{split}
 \Omega_{\infty,\ell}
 (f\circ p_{\ell m}\otimes A\otimes I)
 &=
 \int_{K_m} f(u)
 \operatorname{Tr}\left[
 \left\{\int
 \operatorname{Tr}_{\ell\setminus m}
 \rho_{\infty,\ell}(v)\,
 d\kappa_{\ell m}^{\,u}(v)\right\}A\right]
 d\nu_{\infty,m}(u).
 \end{split}                                                \tag{89.14}
\]
Compatibility (89.10) identifies the left side with
\[
 \int_{K_m}f(u)
 \operatorname{Tr}\{\rho_{\infty,m}(u)A\}
 d\nu_{\infty,m}(u).
\]
Use a countable dense family of compact operators \(A\) and
continuous functions \(f\), followed by uniqueness of
Radon--Nikodym densities, to obtain (89.13) outside one null set.
\end{proof}

The conditional field at a smaller region is generally an average
over the external classical configurations in the fibre of
\(p_{\ell m}\), followed by quantum partial trace.  It is not the
partial trace at one arbitrarily chosen extension.

\subsection{No global density-matrix claim}

Even though every fixed local restriction is normal, the
quasi-local GNS representation may be disjoint from a chosen
infinite tensor-product reference representation.  Infinite-volume
KMS states and phases commonly have no global trace-class density
matrix.  The physically relevant conclusion of V89 is local
normality and projective consistency, not a global trace.

Local energy density bounds can be added by choosing positive local
Hamiltonians \(H_m\) with
\({\cal E}_m\leq C\,\operatorname{Vol}(m)\).  This controls energy per
volume but is not needed for the fixed-\(m\) normality proof.

\begin{assessment}[Status after V89]
\label{ass:v12v89-status}
The extensivity caveat of V88 is resolved at the logical level.
A cutoff-uniform energy moment for every fixed region gives quantum
no escape and a trace-class conditional disintegration on that
region, even when the constants grow with volume.  Compatible local
states yield the correct conditional partial-trace relation after
averaging over external classical backgrounds.
\end{assessment}

\begin{firewall}[Claim boundary after V89]
\label{fire:v12v89-claim}
No physical RPESM local net, tensor factorization, background
projection, local elliptic Hamiltonian, or regionwise energy bound
has been populated.  Gauge theories may have edge modes and
nonfactorizing local Hilbert spaces, in which case (89.11)--(89.13)
must be replaced by the appropriate inclusion and conditional
expectation.  Local normality does not prove a global density matrix,
clustering, phase uniqueness, or universality.  The full
SM--Einstein continuum is not yet constructed.
\end{firewall}

\subsection{Parent record}

This V89 note appends to the validated V88 file with record
\[
\begin{array}{c|c}
\text{lines}&24626\\
\text{bytes}&835001\\
\text{SHA--256}&
\texttt{F09EFE2987B0E718596B0980E4FAC1CD5F76A099E2B07E3230B1879DBF1E37B9}.
\end{array}
\]

% =============================================================================
% END APPENDED TWELFTH-CYCLE V89 MATERIAL
% =============================================================================


% =============================================================================
% BEGIN APPENDED TWELFTH-CYCLE V90 MATERIAL
% =============================================================================

\section{V90 audit: locality before global majorants}
\label{sec:v12v90}

This is the mandatory five-version audit following V85.  It compares
the present Standard Model--Einstein line with the newest active
Yang--Mills and Navier--Stokes notes in the shared directory.  Its
purpose is mathematical: a useful companion result is imported only
after its hypotheses and claim boundary have been separated from its
numerical evidence.

\subsection{Files examined and new information}

The active companions examined for this audit were
\[
\begin{array}{c|c|c|c}
\text{programme}&\text{version}&\text{bytes}&\text{SHA--256}\\ \hline
\mathrm{YM}&62&787188&
\texttt{635C7B5BB3ED2F5E64C22C568390387D2F68F5F561F4F9449BD24BBB5B444EB1}\\
\mathrm{NS}&26&278283&
\texttt{82C887F8C2C69E4F28FAD7C3DC8DE5B9099CB5A4BF431EBA1FFF3E91935978AD}.
\end{array}
\]
The Navier--Stokes companion is unchanged from the preceding audit,
so no new assertion is imported from it.

The Yang--Mills line has advanced through two relevant steps.  Its
V61 derives a conjugate-reflection law for the response and evaluates
uncompressed \(M=3\) and \(M=4\) grids.  Those computations are
diagnostic evidence rather than an interval proof of a continuum
sign.  Its V62 proves a rational real-domain derivative majorant and
a valid \(M^{-2}\) quadrature estimate.  The resulting numerical
constant is about thirty orders of magnitude too large for the
desired sign conclusion.  Thus the estimate is correct but
nonclosing.  The stated next route is cellwise recursion or a local
Hessian floor which retains cancellation and reflection structure.

\subsection{Transfer to the SM--Einstein programme}

The transferable lesson is a proof-design rule.  If a theory is
assembled from local or sectorial data, replacing all of that data by
one global ambient norm can erase precisely the structure needed to
close the estimate.  In V89 the analogous danger is to impose a
global tensor factorization merely to obtain a partial trace.  Gauge
edge modes, Gauss-law centres, and superselection sectors need not
admit that factorization.

Accordingly, the next step will retain the local algebraic
inclusions themselves.  Let
\[
 {\cal M}_m\xrightarrow{\ \iota_{m\ell}\ }{\cal M}_\ell,
 \qquad m\leq\ell,
\]
be normal unital inclusions of local von Neumann algebras.  A normal
large-region state has a canonical small-region restriction through
the preadjoint
\[
 \iota_{m\ell *}:({\cal M}_\ell)_*\longrightarrow({\cal M}_m)_*,
 \qquad
 \iota_{m\ell *}\omega=\omega\circ\iota_{m\ell}.
\]
This operation requires neither a Hilbert-space tensor product nor a
conditional expectation.  It reduces to the partial trace in the
type-I tensor-product case treated in V89.

There are two points to keep distinct in the next proof.

\begin{enumerate}
\item Restriction along a normal inclusion is canonical and is enough
for projective consistency of locally normal states.
\item A normal conditional expectation from the large algebra onto
the embedded small algebra is additional structure.  Asking the
large state to be recovered through such an expectation is a Markov
or recovery property, not a consequence of consistency.
\end{enumerate}

The distinction prevents a circular argument: the continuum local
state will not be assumed to possess a factorization or recovery map
which the construction itself is supposed to justify.

\begin{assessment}[Decision after the V90 audit]
\label{ass:v12v90-decision}
V91 will replace the finite-dimensional tensor-shell caveat in V89
by a theorem for directed systems of sigma-finite von Neumann
algebras with separable preduals.  It will prove positivity, norm
contraction, projective consistency, stability of approximate
compatibility, and the exact relation with the type-I partial trace.
It will also state precisely what may be concluded about centre and
sector weights.
\end{assessment}

\begin{firewall}[Claim boundary after V90]
\label{fire:v12v90-claim}
The Yang--Mills calculations do not establish a continuum response
sign or a Yang--Mills mass gap.  The audit does not populate the
local SM--Einstein algebras, prove local normality, construct edge
sectors, or supply the physical energy bounds required in V89.
Those inputs remain conditional.  The full SM--Einstein continuum is
not yet constructed.
\end{firewall}

\subsection{Parent record}

This V90 note appends to the validated V89 file with record
\[
\begin{array}{c|c}
\text{lines}&24895\\
\text{bytes}&844545\\
\text{SHA--256}&
\texttt{C46CCE11D70A04BF18E7F1F941F156DD502412B1378997ADCF36F74C0ECDCB8D}.
\end{array}
\]

% =============================================================================
% END APPENDED TWELFTH-CYCLE V90 MATERIAL
% =============================================================================

% =============================================================================
% BEGIN APPENDED TWELFTH-CYCLE V91 MATERIAL
% =============================================================================

\section{V91: local normal inclusions, edge sectors, and projective
predual compactness}
\label{sec:v12v91}

V89 obtained local density fields and, for a finite-dimensional
tensor shell, their partial-trace consistency.  The finite-shell
hypothesis was used only because \(A\otimes I\) need not be compact
when the exterior Hilbert space is infinite-dimensional.  It is not
a physical locality principle.  This note moves to the natural
von Neumann local algebras, where restriction is defined by the
preadjoint of a normal inclusion.  This covers arbitrary separable
type-I shells and also local algebras with centres or edge sectors.

\subsection{The directed local normal system}

Let the regions be indexed by \(\mathbb N\).  For \(m\leq\ell\), let
\({\cal M}_m\) and \({\cal M}_\ell\) be sigma-finite von Neumann
algebras with separable preduals, and let
\[
 \iota_{m\ell}:{\cal M}_m\longrightarrow{\cal M}_\ell             \tag{91.1}
\]
be a normal unital injective star homomorphism.  Assume
\[
 \iota_{mm}=\operatorname{id},\qquad
 \iota_{\ell r}\circ\iota_{m\ell}=\iota_{mr}
 \quad(m\leq\ell\leq r).                                          \tag{91.2}
\]
No tensor-product representation is included in these hypotheses.

Recall that \(({\cal M}_m)_*\) denotes the unique predual, whose
positive norm-one elements are precisely the normal states.  Define
\[
 R_{m\ell}:({\cal M}_\ell)_*\longrightarrow({\cal M}_m)_*,
 \qquad R_{m\ell}\omega=\omega\circ\iota_{m\ell}.                  \tag{91.3}
\]

\begin{proposition}[Canonical normal restriction]
\label{prop:v12v91-restriction}
The map \(R_{m\ell}\) is positive and contractive.  It maps normal
states to normal states and satisfies
\[
 R_{m\ell}R_{\ell r}=R_{mr}.                                      \tag{91.4}
\]
For every \(\omega,\eta\in({\cal M}_\ell)_*\),
\[
 \|R_{m\ell}\omega-R_{m\ell}\eta\|
 \leq \|\omega-\eta\|.                                           \tag{91.5}
\]
\end{proposition}

\begin{proof}
Normality of \(\iota_{m\ell}\) means that composition with it sends
normal functionals to normal functionals, so (91.3) is the
preadjoint map.  Positivity is immediate.  If \(\omega\geq0\), then
a positive functional has norm equal to its value at the unit, and
unitality gives
\[
 \|R_{m\ell}\omega\|
 =(R_{m\ell}\omega)(1)
 =\omega(\iota_{m\ell}(1))
 =\omega(1)=\|\omega\|.                                          \tag{91.6}
\]
Thus normal states remain normal states.  For a general predual
functional, the operator norm of composition is bounded by the norm
of \(\iota_{m\ell}\), which is one; this proves (91.5).  Finally,
for \(A\in{\cal M}_m\),
\[
 (R_{m\ell}R_{\ell r}\omega)(A)
 =\omega(\iota_{\ell r}(\iota_{m\ell}(A)))
 =\omega(\iota_{mr}(A))
 =(R_{mr}\omega)(A),
\]
which proves (91.4).
\end{proof}

Consequently, projective consistency of local normal states is the
intrinsic identity
\[
 \omega_m=R_{m\ell}\omega_\ell,\qquad m\leq\ell.                  \tag{91.7}
\]
It does not require a chosen extension of a classical background, a
split inclusion, or an exterior vacuum.

\subsection{A diagonal theorem with approximate cutoff consistency}

At a finite regulator \(n\), the states at two regions may not be
exactly compatible because their ultraviolet cutoffs or boundary
collars differ.  The following result isolates the convergence
statement actually needed.

\begin{theorem}[Local predual compactness closes the projective
limit]
\label{thm:v12v91-projective}
For every \(n,m\in\mathbb N\), let
\(\omega_m^{(n)}\) be a normal state on \({\cal M}_m\).  Assume:

\begin{enumerate}
\item for each fixed \(m\), the set
\(\{\omega_m^{(n)}:n\in\mathbb N\}\) is relatively compact in the
predual norm of \(({\cal M}_m)_*\);
\item for each fixed \(m\leq\ell\),
\[
 \delta_{m\ell}^{(n)}
 :=\|\omega_m^{(n)}-R_{m\ell}\omega_\ell^{(n)}\|
 \longrightarrow0.                                               \tag{91.8}
\]
\end{enumerate}

Then there are a subsequence \(n_j\) and normal states
\(\omega_m\in({\cal M}_m)_*\) such that, for every fixed \(m\),
\[
 \|\omega_m^{(n_j)}-\omega_m\|\longrightarrow0,                  \tag{91.9}
\]
and the limiting family obeys (91.7).  It therefore determines a
unique state on the norm closure of the directed union of the local
C-star algebras.
\end{theorem}

\begin{proof}
Relative norm compactness at \(m=1\) gives a subsequence converging in
\(({\cal M}_1)_*\).  From it select a further subsequence converging
at \(m=2\), and continue.  The diagonal subsequence converges at
every fixed \(m\).  The positive cone is norm closed, while
evaluation at the unit is norm continuous, so every limit
\(\omega_m\) is a normal positive functional of norm one.

Fix \(m\leq\ell\).  Contractivity (91.5) and (91.8) give
\[
\begin{split}
 \|\omega_m-R_{m\ell}\omega_\ell\|
 &\leq
 \|\omega_m-\omega_m^{(n_j)}\|
 +\delta_{m\ell}^{(n_j)}
 +\|R_{m\ell}(\omega_\ell^{(n_j)}-\omega_\ell)\|\\
 &\leq
 \|\omega_m-\omega_m^{(n_j)}\|
 +\delta_{m\ell}^{(n_j)}
 +\|\omega_\ell^{(n_j)}-\omega_\ell\|
 \longrightarrow0.                                               \tag{91.10}
\end{split}
\]
Hence (91.7) holds.

For an element represented by \(A\in{\cal M}_m\), set
\(\omega_{\mathrm{loc}}(A)=\omega_m(A)\).  Equation (91.7) makes
this definition independent of the representing region.  It is
positive and bounded with norm one on the algebraic directed union,
so it extends uniquely by continuity to its C-star norm closure.
\end{proof}

The theorem also has a pointwise version.  If
\(\omega_m^{(n_j)}(A)\to\omega_m(A)\) for every \(A\in{\cal M}_m\),
if every resulting \(\omega_m\) is known independently to be normal,
and if (91.8) holds, then (91.10) applied to each observable proves
compatibility.  The independent normality clause is essential:
pointwise limits of normal states can be singular.  The regionwise
energy no-escape mechanism of V89 is one concrete way to supply that
clause in a type-I realization.

\subsection{Recovery of partial trace without a finite shell}

\begin{corollary}[The type-I formula]
\label{cor:v12v91-typei}
Let
\[
 {\cal M}_m={\cal B}({\cal H}_m),\qquad
 {\cal M}_\ell={\cal B}({\cal H}_m\otimes{\cal H}_{\ell\setminus m}),
 \qquad
 \iota_{m\ell}(A)=A\otimes I,                                    \tag{91.11}
\]
where both Hilbert spaces are separable and the exterior space may
be infinite-dimensional.  If
\(\omega_\ell(X)=\operatorname{Tr}(\rho_\ell X)\) for a positive
trace-class \(\rho_\ell\) of trace one, then
\[
 R_{m\ell}\omega_\ell(A)
 =\operatorname{Tr}\!\left[
   (\operatorname{Tr}_{\ell\setminus m}\rho_\ell)A\right].
                                                                    \tag{91.12}
\]
Thus the predual restriction is represented by the usual partial
trace, with no finite-dimensional shell assumption.
\end{corollary}

\begin{proof}
The partial trace of a trace-class operator over an arbitrary
separable Hilbert space is the unique trace-class operator satisfying
\[
 \operatorname{Tr}\!\left[
  (\operatorname{Tr}_{\ell\setminus m}\rho_\ell)A\right]
 =\operatorname{Tr}[\rho_\ell(A\otimes I)]
\]
for all \(A\in{\cal B}({\cal H}_m)\).  The right side is
\((R_{m\ell}\omega_\ell)(A)\), proving (91.12).
\end{proof}

This corollary removes the technical finite-dimensional restriction
in V89 once the locally normal state has been extended from the
compact-operator core to the local von Neumann algebra.  For a
non-type-I gauge algebra, (91.3), rather than a matrix partial trace,
is the definition that survives.

\subsection{Centres, sectors, and quantitative stability}

Let \(z\in Z({\cal M}_m)\) be a central projection.  Compatibility
implies
\[
 \omega_m(z)=\omega_\ell(\iota_{m\ell}(z)).                       \tag{91.13}
\]
The projection \(\iota_{m\ell}(z)\) need not be central in
\({\cal M}_\ell\).  Therefore (91.13) always preserves the
probability of the embedded event, but it describes a nested
superselection sector only under the additional centre-preserving
condition
\[
 \iota_{m\ell}(Z({\cal M}_m))\subseteq Z({\cal M}_\ell).           \tag{91.14}
\]
If \(\{z_a\}_{a\in{\cal A}}\) is a finite or countable orthogonal
central partition of the unit in \({\cal M}_m\), normality yields
\[
 p_{m,a}:=\omega_m(z_a)\geq0,\qquad
 \sum_{a\in{\cal A}}p_{m,a}=1.                                   \tag{91.15}
\]
For approximately compatible cutoff states, (91.8) gives the sharp
observable-independent estimate
\[
 \left|\omega_m^{(n)}(z_a)
 -\omega_\ell^{(n)}(\iota_{m\ell}(z_a))\right|
 \leq\delta_{m\ell}^{(n)},                                       \tag{91.16}
\]
because a nonzero projection has norm one.  Hence sector weights
cannot jump in the limit once the predual defect closes.

This is the algebraic place for Gauss-law boundary charge and edge
sector data.  Their existence and tightness are separate analytic
inputs; the theorem proves that such data are not discarded by the
passage from a large region to a small one.

\subsection{Conditional expectations are additional structure}

Suppose a normal conditional expectation
\[
 E_{\ell m}:{\cal M}_\ell\longrightarrow
 \iota_{m\ell}({\cal M}_m)                                      \tag{91.17}
\]
has been chosen.  A small-region state then has the particular
extension
\[
 \widehat\omega_\ell
 =\omega_m\circ\iota_{m\ell}^{-1}\circ E_{\ell m}.                \tag{91.18}
\]
Equation (91.18) implies (91.7), but the converse is false.

Indeed, in finite dimensions take
\({\cal M}_\ell={\cal B}({\cal H}_m\otimes{\cal H}_e)\), choose a
state \(\tau_e\), and set
\[
 E_{\ell m}(X)
 =\bigl((\operatorname{id}\otimes\tau_e)(X)\bigr)\otimes I_e.
                                                                    \tag{91.19}
\]
The extension (91.18) is the product state
\(\omega_m\otimes\tau_e\).  A correlated or entangled state can have
the same marginal \(\omega_m\), hence obey (91.7), while differing
from that product state.  Requiring (91.18) is consequently a
Markov, recovery, or specified-environment condition.  It cannot be
inserted merely to prove local consistency.

\begin{assessment}[Status after V91]
\label{ass:v12v91-status}
The finite tensor-shell caveat in V89 is now removed at the
algebraic level.  Normal restriction along local inclusions is
positive, contractive, and projectively consistent.  Local predual
relative compactness plus vanishing compatibility defects produces
a compatible continuum family by a diagonal argument.  The
construction preserves embedded edge events and, for
centre-preserving inclusions, their superselection weights.  In a
type-I representation it reproduces partial trace for arbitrary
separable exterior Hilbert space.
\end{assessment}

\begin{firewall}[Claim boundary after V91]
\label{fire:v12v91-claim}
No concrete SM--Einstein local von Neumann net, normal inclusion,
centre-preserving edge system, or predual compactness estimate has
been constructed here.  V89 supplies a conditional type-I
no-escape mechanism, not those physical inputs in full generality.
A conditional expectation or Markov recovery law is not obtained.
The quasi-local state is not asserted to be normal in any preferred
global representation.  Reflection positivity, clustering,
universality, and the full SM--Einstein continuum remain open.
\end{firewall}

\subsection{Next bottleneck}

The remaining analytic gap in this branch is now precise.  One needs
a usable criterion that produces the relative predual compactness in
Theorem~\ref{thm:v12v91-projective} for local gauge algebras with
centres, preferably from the V62/V88 elliptic energy control without
assuming a type-I tensor factorization.  The next note should
formulate a noncommutative uniform-integrability criterion using
finite-weight spectral projections or a faithful semifinite weight,
and should identify exactly which modular and sector-tail estimates
must be supplied by the RPESM construction.

\subsection{Parent record}

This V91 note appends to the validated V90 file with record
\[
\begin{array}{c|c}
\text{lines}&25015\\
\text{bytes}&849537\\
\text{SHA--256}&
\texttt{A961499A1029C1B717E7A43884FF4AE4B41EB276487762338B3F44FA5937F989}.
\end{array}
\]

% =============================================================================
% END APPENDED TWELFTH-CYCLE V91 MATERIAL
% =============================================================================

% =============================================================================
% BEGIN APPENDED TWELFTH-CYCLE V92 MATERIAL
% =============================================================================

\section{V92: a noncommutative compactness criterion with sector
tails}
\label{sec:v12v92}

The next bottleneck identified in V91 is not solved by asking merely
for finite trace or finite weight of an energy sublevel.  Such a
sublevel can still carry infinitely many mutually separated normal
states.  This note gives a sufficient compactness theorem with two
independent ingredients: tightness across central sectors and
finite-rank approximation inside the retained sectors.  It then
derives the familiar compact-resolvent energy criterion as a
type-I corollary.

\subsection{Uniform finite-rank approximation}

Let \({\cal M}\) be a von Neumann algebra and
\({\cal K}\subset({\cal M})_*\) a norm-bounded family of normal
functionals.

\begin{proposition}[Finite-rank predual approximation]
\label{prop:v12v92-finite-rank}
Suppose that for every \(\varepsilon>0\) there is a normal unital
positive finite-rank map
\[
 T_\varepsilon:{\cal M}\longrightarrow{\cal M}                  \tag{92.1}
\]
such that
\[
 \sup_{\omega\in{\cal K}}
 \|\omega-T_{\varepsilon *}\omega\|<\varepsilon.                 \tag{92.2}
\]
Then \({\cal K}\) is relatively compact in the predual norm.
\end{proposition}

\begin{proof}
A unital positive map has operator norm one, so its preadjoint is a
contraction.  Because \(T_\varepsilon\) has finite-dimensional
range, \(T_{\varepsilon *}\) also has finite-dimensional range.
The image of the bounded set \({\cal K}\) is therefore totally
bounded.  Choose a finite \(\varepsilon\)-net for
\(T_{\varepsilon *}{\cal K}\).  Equation (92.2) makes the same set a
\(2\varepsilon\)-net for \({\cal K}\).  Thus \({\cal K}\) is totally
bounded.  Its closure in the Banach space \(({\cal M})_*\) is
complete and totally bounded, hence compact.
\end{proof}

Complete positivity may be imposed in (92.1), and is natural for
physical coarse graining, but positivity is enough for this proof.
Pointwise ultraweak approximation of the identity by finite-rank
maps is insufficient: the approximation in (92.2) must be uniform
over the regulated state family under consideration.

\subsection{Central sector assembly}

Let \(\{z_a:a\in{\cal A}\}\) be a finite or countable family of
mutually orthogonal central projections with strong sum one.  For a
normal positive functional \(\omega\), define
\[
 \omega_a(A)=\omega(z_aA),\qquad
 z_F=\sum_{a\in F}z_a,\qquad
 \omega_F(A)=\omega(z_FA)                                       \tag{92.3}
\]
for finite \(F\subset{\cal A}\).  Normality gives the norm identity
\[
 \|\omega-\omega_F\|=\omega(1-z_F).                              \tag{92.4}
\]

\begin{theorem}[Sector tail plus compact cores]
\label{thm:v12v92-sector}
Let \({\cal K}\) be a family of normal states on \({\cal M}\).
Assume:

\begin{enumerate}
\item for every \(\varepsilon>0\) there is a finite
\(F\subset{\cal A}\) such that
\[
 \sup_{\omega\in{\cal K}}\omega(1-z_F)<\varepsilon;              \tag{92.5}
\]
\item for every finite \(F\), the truncated family
\[
 {\cal K}_F=\{\omega_F:\omega\in{\cal K}\}
 \subset(z_F{\cal M})_*                                         \tag{92.6}
\]
is relatively compact in norm.
\end{enumerate}

Then \({\cal K}\) is relatively compact in \(({\cal M})_*\).
Condition \(2\) follows, in particular, if the truncated family has
the uniform finite-rank approximation property of
Proposition~\ref{prop:v12v92-finite-rank}.
\end{theorem}

\begin{proof}
Fix \(\varepsilon>0\), and choose \(F\) from (92.5).
Relative compactness of \({\cal K}_F\) provides a finite
\(\varepsilon\)-net \(\{\eta_1,\ldots,\eta_N\}\) for that family.
For each \(\omega\in{\cal K}\), choose \(j\) with
\(\|\omega_F-\eta_j\|<\varepsilon\).  By (92.4),
\[
 \|\omega-\eta_j\|
 \leq\|\omega-\omega_F\|+\|\omega_F-\eta_j\|
 <2\varepsilon.                                                  \tag{92.7}
\]
Thus \({\cal K}\) is totally bounded, and the conclusion follows as
in Proposition~\ref{prop:v12v92-finite-rank}.
\end{proof}

A central coercive observable gives (92.5).  More precisely, choose
numbers \(q_a\geq0\) such that
\[
 \#\{a:q_a\leq R\}<\infty\quad\text{for every }R<\infty,          \tag{92.8}
\]
and let \(Q=\sum_a q_a z_a\) be the affiliated positive central
operator.  If
\[
 \sup_{\omega\in{\cal K}}\omega(Q)\leq C,                         \tag{92.9}
\]
then for \(F_R=\{a:q_a\leq R\}\),
\[
 \sup_{\omega\in{\cal K}}\omega(1-z_{F_R})
 \leq \frac{C}{R}.                                               \tag{92.10}
\]
This is the sector form of Markov's inequality.  It controls escape
to high boundary charge or high superselection label, but it says
nothing by itself about compactness within any fixed sector.

\subsection{The type-I compact-resolvent corollary}

The preceding separation becomes explicit for a countable direct
sum
\[
 {\cal M}=\bigoplus_{a\in{\cal A}}{\cal B}({\cal H}_a),\qquad
 {\cal M}_*=\bigoplus_{a\in{\cal A}}^{\ell^1}
 {\cal T}_1({\cal H}_a).                                        \tag{92.11}
\]
Let \(H_a\geq0\) be self-adjoint with compact resolvent on
\({\cal H}_a\), and let the sector penalties \(q_a\) satisfy
(92.8).  Define the affiliated operator
\[
 G=\bigoplus_{a\in{\cal A}}(q_a I_a+H_a).                         \tag{92.12}
\]

\begin{corollary}[Energy compactness with countable sectors]
\label{cor:v12v92-energy}
For \(C<\infty\), the family of normal states
\[
 {\cal K}_C=
 \{\omega_\rho:\rho\geq0,\ \operatorname{Tr}\rho=1,\
                  \operatorname{Tr}(\rho G)\leq C\}              \tag{92.13}
\]
is relatively compact in predual norm.
\end{corollary}

\begin{proof}
For \(R>0\), let \(P_R=\mathbf 1_{[0,R]}(G)\).  Only finitely many
sectors have \(q_a\leq R\), and compact resolvent makes the
corresponding spectral projection of each \(H_a\) finite rank.
Hence \(P_R\) is finite rank.  The energy bound gives
\[
 \delta_R(\rho):=\operatorname{Tr}(\rho(1-P_R))
 \leq\frac{C}{R}.                                                \tag{92.14}
\]
Writing \(P=P_R\) and \(P^\perp=1-P\), the off-diagonal trace-norm
Cauchy--Schwarz inequality yields
\[
 \|P\rho P^\perp\|_1
 \leq
 \{\operatorname{Tr}(P\rho P)
   \operatorname{Tr}(P^\perp\rho P^\perp)\}^{1/2}
 \leq\sqrt{\delta_R(\rho)}.                                      \tag{92.15}
\]
The same bound holds for \(P^\perp\rho P\), while
\(\|P^\perp\rho P^\perp\|_1=\delta_R(\rho)\).  Therefore
\[
 \|\rho-P_R\rho P_R\|_1
 \leq2\sqrt{C/R}+C/R.                                            \tag{92.16}
\]
The compressed positive operators lie in the unit ball of the
finite-dimensional trace-class space
\(P_R{\cal T}_1P_R\), which is compact.  Equation (92.16) gives a
uniform finite-dimensional approximation, proving total boundedness
of (92.13).
\end{proof}

For a single sector this recovers the local energy compactness used
in V81 and V89.  The new content is that the coercive sector penalty
and the internal compact-resolvent Hamiltonians perform different
jobs and can be verified separately.

\subsection{Why finite weight is not finite-dimensional
compactness}

A tempting replacement for the finite-rank projection \(P_R\) is a
projection of finite value under a faithful semifinite trace or
weight.  That replacement is invalid without an additional
uniform-integrability or finite-rank approximation theorem.

For example, take the diffuse finite von Neumann algebra
\[
 {\cal M}=L^\infty([0,1]),\qquad
 \tau(f)=\int_0^1 f(x)\,dx.                                     \tag{92.17}
\]
Its unit has finite trace.  Choose pairwise disjoint measurable
sets \(E_n\) of positive measure and define the normal states
\[
 \omega_n(f)=\frac{1}{|E_n|}\int_{E_n}f(x)\,dx.                  \tag{92.18}
\]
For \(n\neq k\), the corresponding \(L^1\) densities have disjoint
support, so
\[
 \|\omega_n-\omega_k\|=2.                                       \tag{92.19}
\]
Thus even a family supported by one finite-trace projection need
not be predual-norm precompact.  A fortiori, finite weight of each
energy sublevel cannot replace the compact-core hypothesis in
Theorem~\ref{thm:v12v92-sector}.

For type-III gauge algebras the correct replacement may be a split
or nuclearity estimate, modular compactness, or an explicit family
of normal finite-rank completely positive coarse grainings which
satisfies (92.2).  V92 does not assume which mechanism the physical
construction will provide; it identifies the exact uniform
approximation statement required by V91.

\begin{assessment}[Status after V92]
\label{ass:v12v92-status}
The local compactness input of V91 has been reduced to two
quantitative tasks.  A coercive central observable controls escape
among countable edge or superselection sectors.  Uniform finite-rank
predual approximation controls the remaining degrees of freedom
inside finitely many retained sectors.  In countable type-I sectors,
a compact-resolvent Hamiltonian plus a proper sector penalty proves
both requirements and hence predual-norm compactness.
\end{assessment}

\begin{firewall}[Claim boundary after V92]
\label{fire:v12v92-claim}
A finite trace, finite weight, or bounded first moment alone does not
prove norm compactness in a diffuse or type-III local algebra.  No
RPESM nuclearity, split, modular compactness, sector penalty, or
finite-rank coarse-graining estimate has yet been populated.
Consequently the hypotheses of V91 are still conditional outside
the type-I compact-resolvent realization.  The full SM--Einstein
continuum remains open.
\end{firewall}

\subsection{Next bottleneck}

The next upstream task is to connect the physical heat-kernel control
already isolated in V62 and V88 to the uniform finite-rank
approximation (92.2) on a fixed local algebra.  One should avoid
assuming a global tensor factorization.  A viable intermediate
target is a local phase-space or nuclearity map whose approximation
numbers are summable uniformly over the admitted background core;
such a bound would produce explicit finite-rank normal maps and
would remain meaningful for non-type-I local algebras.

\subsection{Parent record}

This V92 note appends to the validated V91 file with record
\[
\begin{array}{c|c}
\text{lines}&25340\\
\text{bytes}&862071\\
\text{SHA--256}&
\texttt{9C3904E0011130A8F339400721C05CCF580C4D8C29DBCEE47C97D6F47C755AAF}.
\end{array}
\]

% =============================================================================
% END APPENDED TWELFTH-CYCLE V92 MATERIAL
% =============================================================================

% =============================================================================
% BEGIN APPENDED TWELFTH-CYCLE V93 MATERIAL
% =============================================================================

\section{V93: fixed-reference elliptic coercivity produces the
compact core}
\label{sec:v12v93}

V92 separated sector tightness from compactness inside a retained
sector.  This note supplies the latter in the type-I local
realization by one fixed finite-rank coarse graining.  The important
uniformity is that the projection is taken from a reference
elliptic operator and therefore does not move with the classical
background.

\subsection{A finite-rank quantum channel}

Let \({\cal H}\) be separable and let
\(\Lambda\geq0\) be self-adjoint with compact resolvent.  Write
\[
 P_R=\mathbf 1_{[0,R]}(\Lambda),\qquad Q_R=I-P_R.                 \tag{93.1}
\]
For \(R\) above the lowest eigenvalue, choose a density matrix
\(\sigma_R\) supported in \(P_R{\cal H}\), and define
\[
 {\cal T}_R(A)
 =P_RAP_R+\operatorname{Tr}(\sigma_RA)Q_R,
 \qquad A\in{\cal B}({\cal H}).                                  \tag{93.2}
\]
This is a normal unital completely positive finite-rank map.  Its
preadjoint is
\[
 {\cal T}_{R*}(\rho)
 =P_R\rho P_R+\operatorname{Tr}(Q_R\rho)\sigma_R.                 \tag{93.3}
\]

\begin{proposition}[Reference-energy channel bound]
\label{prop:v12v93-channel}
If \(\rho\geq0\), \(\operatorname{Tr}\rho=1\), and
\[
 \operatorname{Tr}(\rho\Lambda)\leq L,                            \tag{93.4}
\]
then
\[
 \|\rho-{\cal T}_{R*}(\rho)\|_1
 \leq2\sqrt{L/R}+2L/R.                                           \tag{93.5}
\]
Consequently, every family satisfying the uniform bound (93.4) has
the uniform finite-rank approximation property (92.2), and is
relatively compact in trace norm.
\end{proposition}

\begin{proof}
Spectral calculus and Markov's inequality give
\[
 \delta_R:=\operatorname{Tr}(\rho Q_R)\leq L/R.                  \tag{93.6}
\]
As in (92.15),
\[
 \|P_R\rho Q_R\|_1+\|Q_R\rho P_R\|_1
 \leq2\sqrt{\delta_R},
 \qquad
 \|Q_R\rho Q_R\|_1=\delta_R.                                    \tag{93.7}
\]
Thus
\[
 \|\rho-P_R\rho P_R\|_1
 \leq2\sqrt{\delta_R}+\delta_R.                                  \tag{93.8}
\]
The second term in (93.3) has trace norm
\(\delta_R\|\sigma_R\|_1=\delta_R\).  The triangle inequality,
(93.6), and (93.8) prove (93.5).  Compact resolvent makes \(P_R\)
finite rank, so Proposition~\ref{prop:v12v92-finite-rank} applies.
\end{proof}

The channel retains the low reference modes and moves all discarded
probability to one fixed low-mode state.  This restores trace and
unitality of the Heisenberg map.  Plain compression \(A\mapsto
P_RAP_R\) would not be unital and would lose the tail mass.

\subsection{Uniform form comparison over a background core}

The reference moment (93.4) follows from a background-dependent
energy moment if all local quadratic forms have one common
coercivity floor.

\begin{lemma}[Common G{\aa}rding comparison]
\label{lem:v12v93-garding}
Let \(D\) be a compact manifold or a bounded domain with fixed local
boundary condition, and let \({\cal H}=L^2(D;E)\) for a
finite-rank bundle \(E\).  Suppose a family of closed
semibounded Laplace-type forms \(k_x\) has common form domain
\({\cal V}\subset H^1(D;E)\), and that there are constants
\(c_0>0\) and \(c_1<\infty\), independent of \(x\), for which
\[
 k_x[u]\geq c_0\|u\|_{H^1}^2-c_1\|u\|_2^2,
 \qquad u\in{\cal V}.                                            \tag{93.9}
\]
Let \(\Lambda\geq0\) be a fixed elliptic reference operator with
compact resolvent and form domain \({\cal V}\), and suppose
\[
 \langle u,\Lambda u\rangle
 \leq c_2\|u\|_{H^1}^2+c_3\|u\|_2^2.                            \tag{93.10}
\]
If \(K_x\) is the operator associated with \(k_x\), then there are
\(a>0\) and \(b<\infty\), independent of \(x\), such that
\[
 K_x\geq a\Lambda-bI                                             \tag{93.11}
\]
in quadratic-form sense.
\end{lemma}

\begin{proof}
Equation (93.9) implies
\[
 \|u\|_{H^1}^2
 \leq c_0^{-1}k_x[u]+c_0^{-1}c_1\|u\|_2^2.                      \tag{93.12}
\]
Insert this into (93.10):
\[
 \langle u,\Lambda u\rangle
 \leq(c_2/c_0)k_x[u]
   +(c_2c_1/c_0+c_3)\|u\|_2^2.                                  \tag{93.13}
\]
Rearranging gives (93.11) with
\[
 a=c_0/c_2,\qquad b=c_1+c_0c_3/c_2,                              \tag{93.14}
\]
after increasing \(b\) harmlessly if needed.  Compactness of the
embedding \({\cal V}\hookrightarrow{\cal H}\) gives compact
resolvent for the standard reference choices.
\end{proof}

For metrics in a compact uniformly elliptic \(C^{2,\alpha}\) core,
one first identifies the varying \(L^2\) spaces with a fixed
reference space by the volume-density unitary.  Uniform coefficient
bounds and a fixed boundary condition then imply (93.9).  This last
statement is a route for populating the lemma; it still requires the
actual RPESM kinetic operators, gauge fixing, and boundary domains
to be checked.

\begin{proposition}[Moving energy to fixed reference energy]
\label{prop:v12v93-transfer}
Assume (93.11), let \(\rho_x\) be normal states, and suppose
\[
 \sup_x\operatorname{Tr}(\rho_xK_x)\leq C.                       \tag{93.15}
\]
Then
\[
 \sup_x\operatorname{Tr}(\rho_x\Lambda)
 \leq\frac{C+b}{a}.                                              \tag{93.16}
\]
The common channel \({\cal T}_R\) therefore approximates every
\(\rho_x\) with error bounded by
\[
 2\sqrt{\frac{C+b}{aR}}+
 2\frac{C+b}{aR}.                                                \tag{93.17}
\]
\end{proposition}

\begin{proof}
Apply the positive normal functional defined by \(\rho_x\) to the
form inequality (93.11), using increasing bounded spectral
truncations when an expectation is initially infinite.  This gives
\[
 C\geq\operatorname{Tr}(\rho_xK_x)
 \geq a\operatorname{Tr}(\rho_x\Lambda)-b,
\]
and hence (93.16).  Proposition~\ref{prop:v12v93-channel} gives
(93.17).
\end{proof}

\subsection{The Gibbs input from a heat trace}

The remaining energy bound can be obtained from the heat trace and
the partition floor already isolated in V62 and V88.

\begin{corollary}[Uniform Gibbs compactness]
\label{cor:v12v93-gibbs}
Let \(K_x\geq0\) satisfy (93.11), fix \(\beta>0\), and suppose
\[
 Z_x:=\operatorname{Tr}(e^{-\beta K_x})\geq z_0>0,               \tag{93.18}
\]
while
\[
 \operatorname{Tr}(e^{-\beta a\Lambda/2})<\infty.                \tag{93.19}
\]
For
\[
 \rho_x=Z_x^{-1}e^{-\beta K_x},                                  \tag{93.20}
\]
one has
\[
 \sup_x\operatorname{Tr}(\rho_xK_x)
 \leq
 C_\beta:=
 \frac{2e^{\beta b/2}}{e\beta z_0}
 \operatorname{Tr}(e^{-\beta a\Lambda/2}).                       \tag{93.21}
\]
Consequently (93.17) holds with \(C=C_\beta\), and the Gibbs family
is relatively compact in trace norm.
\end{corollary}

\begin{proof}
For \(t\geq0\),
\[
 te^{-\beta t}\leq\frac{2}{e\beta}e^{-\beta t/2}.                 \tag{93.22}
\]
Functional calculus therefore gives
\[
 \operatorname{Tr}(K_xe^{-\beta K_x})
 \leq\frac{2}{e\beta}
       \operatorname{Tr}(e^{-\beta K_x/2}).                      \tag{93.23}
\]
The min--max principle applied to (93.11) gives
\[
 \operatorname{Tr}(e^{-\beta K_x/2})
 \leq e^{\beta b/2}
       \operatorname{Tr}(e^{-\beta a\Lambda/2}).                 \tag{93.24}
\]
Divide (93.23) by the partition lower bound (93.18) to obtain
(93.21), then use
Proposition~\ref{prop:v12v93-transfer}.
\end{proof}

The partition floor cannot be discarded.  It prevents normalization
from amplifying an otherwise small numerator.  As in V88, it may
come from one fixed trial vector whose \(K_x\)-energy is uniformly
bounded over the admitted background core.

\subsection{What this supplies to the projective theorem}

For every fixed region \(m\), suppose its background-dependent local
Hamiltonians can be transported to one Hilbert space
\({\cal H}_m\), and suppose (93.11), (93.18), and (93.19) hold with
constants depending on \(m\) but not on the regulator or admitted
background.  Then the local Gibbs densities form a relatively
trace-norm compact family.  Theorem~\ref{thm:v12v91-projective} can
therefore be applied region by region; the constants may grow with
\(m\), because its diagonal argument fixes \(m\) before taking the
cutoff limit.

Together with V92, a countable collection of central edge sectors
can also be included, provided a proper sector penalty controls
their tail and the same fixed-reference estimate holds in each
finite set of retained sectors.

\begin{assessment}[Status after V93]
\label{ass:v12v93-status}
A concrete sufficient route from local elliptic analysis to the
predual compactness required in V91 is now proved.  Uniform
G{\aa}rding comparison transfers moving-background energy to one
compact-resolvent reference operator.  A fixed finite-rank normal
quantum channel then approximates all admitted local states in
trace norm with the explicit error (93.17).  Uniform heat trace and
a partition floor supply the required Gibbs energy moment.
\end{assessment}

\begin{firewall}[Claim boundary after V93]
\label{fire:v12v93-claim}
The proof is conditional on a common Hilbert-space trivialization,
common form domain, uniform G{\aa}rding constants, a partition
floor, and the stated heat trace.  Gauge constraints, ghosts,
fermionic signs, interacting non-Gaussian states, edge-sector tails,
and dynamical gravity have not been shown to satisfy these
conditions.  The finite-rank channel acts in a type-I local
realization and does not by itself establish nuclearity for a
type-III observable algebra.  The full SM--Einstein continuum
remains open.
\end{firewall}

\subsection{Next bottleneck}

The next step should separate gauge constraints from elliptic
coercivity.  The fixed-reference projection need not preserve the
Gauss-law subspace.  One needs a quantitative statement showing that
group averaging or BRST reduction commutes with the finite-rank
approximation up to a controlled defect, and that the defect
vanishes uniformly on the gravitational compact core.  This is
upstream of any claim that the compact local states belong to the
physical observable algebra.

\subsection{Parent record}

This V93 note appends to the validated V92 file with record
\[
\begin{array}{c|c}
\text{lines}&25617\\
\text{bytes}&872715\\
\text{SHA--256}&
\texttt{E3226A9A169401B440E70D2E2C4597BE70C346FD4D9830D0AA8797A9F4252848}.
\end{array}
\]

% =============================================================================
% END APPENDED TWELFTH-CYCLE V93 MATERIAL
% =============================================================================

% =============================================================================
% BEGIN APPENDED TWELFTH-CYCLE V94 MATERIAL
% =============================================================================

\section{V94: equivariant compact cores and compact-group Gauss
reduction}
\label{sec:v12v94}

V93 produced one fixed finite-rank channel from a reference elliptic
operator, but did not show that the channel respects a gauge
constraint.  This note proves exact descent when the regulated gauge
symmetry is a compact group and the reference operator is invariant.
It also distinguishes invariant observables from the stronger
singlet-support condition on vectors.

\subsection{Gauge averaging and the fixed-point algebra}

Let \(G\) be a second-countable compact group with normalized Haar
measure \(dg\), and let \(U:G\to{\cal U}({\cal H})\) be a strongly
continuous unitary representation on a separable Hilbert space.
Set
\[
 \alpha_g(A)=U_gAU_g^*,\qquad
 {\cal M}^G=\{A\in{\cal B}({\cal H}):\alpha_g(A)=A
                    \text{ for all }g\}.                         \tag{94.1}
\]
The Haar average
\[
 {\cal E}_G(A)=\int_G\alpha_g(A)\,dg                              \tag{94.2}
\]
is a normal unital completely positive conditional expectation from
\({\cal B}({\cal H})\) onto \({\cal M}^G\).  If
\(\varphi\) is a normal state on \({\cal M}^G\), its canonical
normal invariant extension is
\[
 \widetilde\varphi=\varphi\circ{\cal E}_G.                        \tag{94.3}
\]
Conversely, restriction of an invariant normal state to
\({\cal M}^G\) loses no expectation of an invariant observable.

Assume the reference operator \(\Lambda\) from V93 strongly commutes
with \(U_g\) for every \(g\).  Then \(P_R\) and \(Q_R\) in (93.1)
are invariant.  Choose
\[
 \sigma_R=\frac{P_R}{\operatorname{rank}P_R}                     \tag{94.4}
\]
whenever \(P_R\neq0\), and use it in the channel (93.2).

\begin{proposition}[Equivariant finite-rank channel]
\label{prop:v12v94-equivariant}
The channel \({\cal T}_R\) is \(G\)-equivariant:
\[
 {\cal T}_R\circ\alpha_g=\alpha_g\circ{\cal T}_R
 \quad(g\in G).                                                   \tag{94.5}
\]
It therefore restricts to a normal unital completely positive
finite-rank map
\[
 {\cal T}_R^G:{\cal M}^G\longrightarrow{\cal M}^G.               \tag{94.6}
\]
Its preadjoint maps invariant density matrices to invariant density
matrices.
\end{proposition}

\begin{proof}
Strong commutation gives
\(U_gP_R=P_RU_g\), \(U_gQ_R=Q_RU_g\), and
\(U_g\sigma_RU_g^*=\sigma_R\).  Trace cyclicity then yields
\[
 \operatorname{Tr}(\sigma_R\alpha_g(A))
 =\operatorname{Tr}(\sigma_RA).
\]
Substitution in (93.2) proves (94.5).  If \(A\) is fixed by every
\(\alpha_g\), so is \({\cal T}_R(A)\), proving (94.6).  Its range is
contained in the finite-dimensional space
\[
 P_R{\cal B}({\cal H})P_R+\mathbb C Q_R.                         \tag{94.7}
\]
The predual assertion follows by dualizing (94.5).
\end{proof}

\begin{theorem}[Predual compactness after compact-group reduction]
\label{thm:v12v94-fixed}
Let \({\cal K}^G_L\) be a family of normal states \(\varphi\) on
\({\cal M}^G\) whose canonical extensions (94.3), represented by
density matrices \(\widetilde\rho_\varphi\), satisfy
\[
 \operatorname{Tr}(\widetilde\rho_\varphi\Lambda)\leq L.          \tag{94.8}
\]
Then
\[
 \sup_{\varphi\in{\cal K}^G_L}
 \|\varphi-{\cal T}_{R*}^G\varphi\|
 \leq2\sqrt{L/R}+2L/R,                                           \tag{94.9}
\]
and \({\cal K}^G_L\) is relatively compact in
\(({\cal M}^G)_*\).
\end{theorem}

\begin{proof}
Proposition~\ref{prop:v12v93-channel} applied to the canonical
extension gives
\[
 \|\widetilde\varphi-
       {\cal T}_{R*}\widetilde\varphi\|
 \leq2\sqrt{L/R}+2L/R.                                           \tag{94.10}
\]
Restriction of normal functionals from \({\cal B}({\cal H})\) to
\({\cal M}^G\) is contractive by
Proposition~\ref{prop:v12v91-restriction}.  Equivariance identifies
the restricted approximant with
\({\cal T}_{R*}^G\varphi\), proving (94.9).  The map in (94.6) has
finite rank, so
Proposition~\ref{prop:v12v92-finite-rank} proves relative
compactness.
\end{proof}

The same conclusion holds for any density matrix \(\rho\), invariant
or not, when only invariant observables are retained.  Its Haar
average
\[
 \overline\rho=\int_GU_g\rho U_g^*\,dg                           \tag{94.11}
\]
has the same restriction to \({\cal M}^G\).  Because \(\Lambda\)
commutes with \(G\),
\[
 \operatorname{Tr}(\overline\rho\Lambda)
 =\operatorname{Tr}(\rho\Lambda).                                \tag{94.12}
\]
Thus averaging does not worsen the reference-energy compactness
bound.

\subsection{Charge sectors emerge from the representation}

Peter--Weyl decomposition gives
\[
 {\cal H}\cong
 \bigoplus_{\pi\in\widehat G}
 V_\pi\otimes{\cal N}_\pi,\qquad
 U_g\cong\bigoplus_{\pi\in\widehat G}
 \pi(g)\otimes I_{{\cal N}_\pi},                                 \tag{94.13}
\]
where the irreducible spaces \(V_\pi\) are finite-dimensional.
Consequently,
\[
 {\cal M}^G=U(G)'
 \cong\bigoplus_{\pi\in\widehat G}
 I_{V_\pi}\otimes{\cal B}({\cal N}_\pi).                         \tag{94.14}
\]
The isotypic projections \(z_\pi\) are central in
\({\cal M}^G\).  A strongly \(G\)-invariant reference operator
decomposes as
\[
 \Lambda\cong\bigoplus_{\pi\in\widehat G}
 I_{V_\pi}\otimes\Lambda_\pi.                                   \tag{94.15}
\]
If \(\Lambda\) has compact resolvent, then for every \(R\) only
finitely many \(\pi\) contribute spectrum below \(R\), and each
contributing low-energy multiplicity space is finite-dimensional.
Hence the single spectral tail estimate in V93 simultaneously
provides the sector tail (92.5) and the internal compact core
(92.6).  This is the representation-theoretic realization of the
two conditions separated in V92.

\begin{corollary}[Quantitative charge-sector tail]
\label{cor:v12v94-sector}
Define the bottom energy in sector \(\pi\) by
\[
 q_\pi=\inf\operatorname{spec}\Lambda_\pi.                       \tag{94.16}
\]
For every normal state \(\varphi\) satisfying (94.8),
\[
 \sum_{\pi:q_\pi>R}\varphi(z_\pi)\leq\frac{L}{R}.                 \tag{94.17}
\]
\end{corollary}

\begin{proof}
The central projection
\(z_{>R}=\sum_{q_\pi>R}z_\pi\) obeys
\(\Lambda\geq Rz_{>R}\).  Apply the canonical extension
\(\widetilde\varphi\) and use normality to interchange the countable
central sum with the state.  Equation (94.8) then gives
\(L\geq R\varphi(z_{>R})\), which is (94.17).
\end{proof}

\subsection{Invariant observables versus singlet vectors}

Let
\[
 P_{\mathrm{sing}}=\int_GU_g\,dg                                \tag{94.18}
\]
be the orthogonal projection onto invariant vectors.  Requiring a
density matrix to satisfy
\[
 \rho=P_{\mathrm{sing}}\rho P_{\mathrm{sing}}                    \tag{94.19}
\]
is stronger than requiring \(\rho\) to commute with \(U_g\), and is
also stronger than merely restricting expectations to
\({\cal M}^G\).

If \(P_{\mathrm{sing}}P_R\neq0\), choose \(\sigma_R\) supported in
that intersection.  Since both projections commute, (93.3) then
maps every state satisfying (94.19) to another such state.  The
bound (93.5) is unchanged.  Thus the finite-rank approximation
respects an exact singlet Gauss constraint whenever the invariant
reference operator has a nonempty retained singlet subspace.

The distinction matters for boundary regions.  Charged edge sectors
may be present in \({\cal M}^G\) even when a globally closed system
is a singlet.  Erasing them by imposing (94.19) separately in every
subregion would contradict the sector bookkeeping of V91--V92.

\subsection{Compatibility with the local net}

Suppose each regulated inclusion
\(\iota_{m\ell}:{\cal M}_m\to{\cal M}_\ell\) intertwines compact
gauge actions:
\[
 \iota_{m\ell}\circ\alpha_g^{(m)}
 =\alpha_g^{(\ell)}\circ\iota_{m\ell}.                            \tag{94.20}
\]
It then maps fixed observables into fixed observables, and the
restriction maps of V91 descend to
\[
 R_{m\ell}^G:({\cal M}_\ell^{G})_*
 \longrightarrow({\cal M}_m^{G})_*.                             \tag{94.21}
\]
All contraction and diagonal consistency estimates in V91 remain
valid after this descent.  If the fixed reference channels are also
compatible up to a defect tending to zero, that defect simply joins
\(\delta_{m\ell}^{(n)}\) in (91.8).

\begin{assessment}[Status after V94]
\label{ass:v12v94-status}
For a compact regulated gauge group, an invariant
compact-resolvent reference operator yields gauge-equivariant
finite-rank channels.  The V93 trace-norm bound descends unchanged
to the fixed observable algebra, and compact resolvent controls
both the countable charge-sector tail and the internal sector
modes.  Exact singlet support is preserved when requested and when
the retained reference subspace contains a singlet.
\end{assessment}

\begin{firewall}[Claim boundary after V94]
\label{fire:v12v94-claim}
Continuum local gauge groups are generally infinite-dimensional and
do not possess the normalized Haar measure used here.  BRST
cohomology, Gribov copies, anomalies, chiral fermion measure,
noncompact diffeomorphism reduction, and gravitational constraints
are not handled by compact-group averaging.  The intertwining
condition (94.20) and an invariant RPESM reference operator have not
been populated.  These results therefore close the compact
regulated constraint step only conditionally; the full
SM--Einstein continuum remains open.
\end{firewall}

\subsection{Next bottleneck}

The next mandatory audit should compare this exact compact-group
descent with the newest Yang--Mills boundary-sector work.  After
that audit, the next substantive note should formulate an
anomaly-sensitive cochain criterion for passing finite-regulator
Ward or BRST identities to the local projective limit.  The aim is
to replace unavailable continuum Haar averaging by a closed,
nilpotent constraint identity while keeping the edge-sector
compatibility established above.

\subsection{Parent record}

This V94 note appends to the validated V93 file with record
\[
\begin{array}{c|c}
\text{lines}&25915\\
\text{bytes}&883317\\
\text{SHA--256}&
\texttt{ECDD3B7A6555D744C4975B1EB4BC4C6FC9DF3185FF52C28B3E348823C37135D9}.
\end{array}
\]

% =============================================================================
% END APPENDED TWELFTH-CYCLE V94 MATERIAL
% =============================================================================

% =============================================================================
% BEGIN APPENDED TWELFTH-CYCLE V95 MATERIAL
% =============================================================================

\section{V95 audit: unchanged companions and the next constraint
closure}
\label{sec:v12v95}

This is the mandatory five-version audit following V90.  The newest
companion files in the shared directory are still
\[
\begin{array}{c|c|c|c}
\text{programme}&\text{version}&\text{bytes}&\text{SHA--256}\\ \hline
\mathrm{YM}&62&787188&
\texttt{635C7B5BB3ED2F5E64C22C568390387D2F68F5F561F4F9449BD24BBB5B444EB1}\\
\mathrm{NS}&26&278283&
\texttt{82C887F8C2C69E4F28FAD7C3DC8DE5B9099CB5A4BF431EBA1FFF3E91935978AD}.
\end{array}
\]
Both hashes agree with the files examined in V90.  There is
therefore no new companion theorem to import, and their earlier
results are not repeated here.

The direction chosen in V90 remains consistent with the unchanged
Yang--Mills evidence: global absolute-value propagation is valid but
destroys the local matrix structure needed for a useful conclusion.
V91--V94 implemented the analogous repair on the SM--Einstein side
by retaining normal local inclusions, central sectors, fixed
elliptic compact cores, and exact compact-group equivariance.

A search of the accumulated SM notes also prevents duplication.
V6--V23 already develop finite cochain anomaly tests, refinement
transport, counterterm repair, and Ward-compatible complex shell
martingales.  V67 already proves pointwise passage of a bounded-core
asymptotic Ward identity through a convergent joint state.  What is
still missing after V91--V94 is a single theorem joining these
results to the locally normal projective system when the BRST
operator is unbounded.

\begin{assessment}[Decision after the V95 audit]
\label{ass:v12v95-decision}
V96 will work with the graph norm of a closed local BRST derivation.
It will prove uniform passage of a possibly nonzero anomaly
functional under predual-norm state convergence, establish its
compatibility under local restriction, and quantify the commutator
defect of the V93 finite-rank channel.  This strengthens V67 rather
than restating it: the identity will extend from a chosen test core
to the completed graph domain and will retain, rather than suppress,
a surviving anomaly.
\end{assessment}

\begin{firewall}[Claim boundary after V95]
\label{fire:v12v95-claim}
The unchanged companion files supply no new continuum result.
Neither V94 nor this audit constructs a continuum BRST differential,
proves its closedness or nilpotence, cancels the retained anomaly,
or establishes compatibility with the physical finite-rank
channels.  The full SM--Einstein continuum remains open.
\end{firewall}

\subsection{Parent record}

This V95 note appends to the validated V94 file with record
\[
\begin{array}{c|c}
\text{lines}&26198\\
\text{bytes}&893778\\
\text{SHA--256}&
\texttt{8581C2B47306809459F328A0537EEB0FEC9EB638E6FD42FAEA31C801F7454D90}.
\end{array}
\]

% =============================================================================
% END APPENDED TWELFTH-CYCLE V95 MATERIAL
% =============================================================================

% =============================================================================
% BEGIN APPENDED TWELFTH-CYCLE V96 MATERIAL
% =============================================================================

\section{V96: graph-norm BRST closure and natural anomaly
restriction}
\label{sec:v12v96}

V67 passes a Ward identity pointwise on a prescribed bounded core.
V91 instead constructs locally normal limits in predual norm, while
V93--V94 approximate those states by finite-rank channels.  The
missing bridge is that a BRST derivation is generally unbounded.
This note works in its graph norm, retains a possibly nonzero anomaly
as part of the conclusion, and quantifies compatibility with local
restriction and finite-rank coarse graining.

\subsection{Closed local BRST data}

Let \({\cal M}_m\) be a graded von Neumann algebra and let
\[
 s_m:{\cal D}_m\subset{\cal M}_m\longrightarrow{\cal M}_m        \tag{96.1}
\]
be a closed densely defined odd graded derivation.  Assume
\({\cal D}_m\) is a graded star algebra,
\(s_m{\cal D}_m\subset{\cal D}_m\), and
\[
 s_m^2F=0,\qquad F\in{\cal D}_m.                                 \tag{96.2}
\]
Equip the domain with the graph norm
\[
 \|F\|_{s_m}:=\|F\|+\|s_mF\|.                                   \tag{96.3}
\]
Closedness makes \(({\cal D}_m,\|\cdot\|_{s_m})\) a Banach space.

For a normal state \(\omega_m\), define its Ward functional on the
graph domain by
\[
 {\cal W}_{\omega_m}(F)=\omega_m(s_mF).                          \tag{96.4}
\]
It is automatically continuous and
\[
 \|{\cal W}_{\omega_m}\|_{{\cal D}_m^*}\leq1.                    \tag{96.5}
\]
A declared anomaly is a continuous functional
\(a_m\in{\cal D}_m^*\) of the corresponding ghost degree.  The
anomalous Ward identity is
\[
 {\cal W}_{\omega_m}=a_m.                                       \tag{96.6}
\]
Anomaly freedom is the additional statement \(a_m=0\); it is not
built into the limit theorem.

\subsection{Uniform passage from regulated derivations}

Let \(\omega_m^{(n)}\) be normal states on \({\cal M}_m\), and let
\(s_m^{(n)}:{\cal D}_m\to{\cal M}_m\) be regulated odd maps.  Define
\[
 \begin{split}
 \tau_m^{(n)}
 &:=\|\omega_m^{(n)}-\omega_m\|,\\
 \kappa_m^{(n)}
 &:=\sup_{\|F\|_{s_m}\leq1}
       \|(s_m^{(n)}-s_m)F\|,\\
 \epsilon_m^{(n)}
 &:=\|\omega_m^{(n)}\circ s_m^{(n)}
                   -a_m^{(n)}\|_{{\cal D}_m^*}.                  \tag{96.7}
 \end{split}
\]

\begin{theorem}[Graph-norm anomalous Ward limit]
\label{thm:v12v96-graph}
Suppose
\[
 \tau_m^{(n)}\to0,\qquad
 \kappa_m^{(n)}\to0,\qquad
 \epsilon_m^{(n)}\to0,\qquad
 \|a_m^{(n)}-a_m\|_{{\cal D}_m^*}\to0.                           \tag{96.8}
\]
Then (96.6) holds on the entire completed graph domain.  More
precisely,
\[
 \|\omega_m\circ s_m-a_m\|_{{\cal D}_m^*}
 \leq
 \tau_m^{(n)}+\kappa_m^{(n)}
 +\epsilon_m^{(n)}
 +\|a_m^{(n)}-a_m\|_{{\cal D}_m^*}.                              \tag{96.9}
\]
If \(s_m^2=0\), the limiting anomaly obeys the
Wess--Zumino consistency identity
\[
 a_m(s_mG)=0,\qquad G\in{\cal D}_m.                              \tag{96.10}
\]
\end{theorem}

\begin{proof}
For \(\|F\|_{s_m}\leq1\),
\[
\begin{split}
 |\omega_m(s_mF)-a_m(F)|
 &\leq
 |(\omega_m-\omega_m^{(n)})(s_mF)|\\
 &\quad+
 |\omega_m^{(n)}((s_m-s_m^{(n)})F)|\\
 &\quad+
 |\omega_m^{(n)}(s_m^{(n)}F)-a_m^{(n)}(F)|\\
 &\quad+
 |(a_m^{(n)}-a_m)(F)|.
                                                                    \tag{96.11}
\end{split}
\]
The four terms are bounded respectively by the four quantities on
the right of (96.9), because
\(\|s_mF\|\leq\|F\|_{s_m}\) and every state has norm one.  Take the
supremum and then the limit.  Finally, (96.6), stability of the
domain under \(s_m\), and (96.2) give
\[
 a_m(s_mG)=\omega_m(s_m^2G)=0.
\]
\end{proof}

Unlike the pointwise theorem in V67, (96.9) is uniform on the graph
unit ball and therefore controls the identity after graph
completion.  It also prevents a surviving \(a_m\) from being
silently relabelled as a vanishing regulator error.

\subsection{Restriction through the local net}

Let the normal inclusions \(\iota_{m\ell}\) of V91 satisfy
\[
 \iota_{m\ell}({\cal D}_m)\subset{\cal D}_\ell,\qquad
 s_\ell\iota_{m\ell}F=\iota_{m\ell}s_mF.                         \tag{96.12}
\]
Assume they are graph bounded:
\[
 \|\iota_{m\ell}F\|_{s_\ell}
 \leq C_{m\ell}\|F\|_{s_m}.                                     \tag{96.13}
\]
For a functional \(a_\ell\in{\cal D}_\ell^*\), write
\[
 \iota_{m\ell}^{\sharp}a_\ell(F)
 :=a_\ell(\iota_{m\ell}F).                                      \tag{96.14}
\]

\begin{proposition}[The anomaly is a projective natural
functional]
\label{prop:v12v96-natural}
If \(\omega_m=R_{m\ell}\omega_\ell\) and both states satisfy
(96.6), then
\[
 a_m=\iota_{m\ell}^{\sharp}a_\ell.                              \tag{96.15}
\]
Thus a nonzero anomaly cannot disappear merely by restricting from
a large region to a small one; it must restrict to zero as a
functional or be removed by the cohomological mechanism of
V19--V23.
\end{proposition}

\begin{proof}
For \(F\in{\cal D}_m\), compatibility, (96.12), and the two Ward
identities give
\[
 \begin{split}
 a_m(F)
 &=\omega_m(s_mF)
 =\omega_\ell(\iota_{m\ell}s_mF)\\
 &=\omega_\ell(s_\ell\iota_{m\ell}F)
 =a_\ell(\iota_{m\ell}F).
 \end{split}                                                     \tag{96.16}
\]
This is (96.15).
\end{proof}

There is also a finite-cutoff defect estimate.  Define
\[
 \begin{split}
 \delta_{m\ell}^{(n)}
 &:=\|\omega_m^{(n)}
       -R_{m\ell}\omega_\ell^{(n)}\|,\\
 \chi_{m\ell}^{(n)}
 &:=\sup_{\|F\|_{s_m}\leq1}
 \|s_\ell^{(n)}\iota_{m\ell}F
       -\iota_{m\ell}s_m^{(n)}F\|.                              \tag{96.17}
 \end{split}
\]
Let \(\epsilon_m^{(n)}\) be as in (96.7), and use the analogous
definition at \(\ell\).  Then
\[
\begin{split}
 \|a_m^{(n)}-\iota_{m\ell}^{\sharp}a_\ell^{(n)}\|_{{\cal D}_m^*}
 \leq{}&
 \epsilon_m^{(n)}
 +C_{m\ell}\epsilon_\ell^{(n)}
 +\delta_{m\ell}^{(n)}(1+\kappa_m^{(n)})
 +\chi_{m\ell}^{(n)}.                                           \tag{96.18}
\end{split}
\]

To prove (96.18), insert the two regulated Ward functionals.  The
small-region Ward residual contributes \(\epsilon_m^{(n)}\), the
large-region residual contributes at most
\(C_{m\ell}\epsilon_\ell^{(n)}\), state incompatibility contributes
\[
 \delta_{m\ell}^{(n)}\|s_m^{(n)}F\|
 \leq\delta_{m\ell}^{(n)}(1+\kappa_m^{(n)})
\]
on the graph unit ball, and failure of the chain-map identity
contributes \(\chi_{m\ell}^{(n)}\).  This proves the stated sum.

Combining (96.18) with Theorem
\ref{thm:v12v91-projective} shows that anomaly compatibility follows
from four independently measurable stocks: Ward residual,
predual restriction defect, BRST operator convergence, and
intertwining defect.

\subsection{The finite-rank channel commutator}

Let \({\cal T}:{\cal M}_m\to{\cal M}_m\) be a normal unital positive
map satisfying
\[
 {\cal T}({\cal D}_m)\subset{\cal D}_m.                           \tag{96.19}
\]
Define its BRST commutator on the graph domain by
\[
 [\,{\cal T},s_m\,]F
 :={\cal T}(s_mF)-s_m({\cal T}F).                                \tag{96.20}
\]

\begin{proposition}[Constraint defect of a compactness channel]
\label{prop:v12v96-channel}
If \(\omega_m\circ s_m=a_m\), then the coarse-grained state
\(\omega_m^{\cal T}:={\cal T}_*\omega_m=\omega_m\circ{\cal T}\)
satisfies
\[
 \omega_m^{\cal T}(s_mF)
 =a_m({\cal T}F)+
   \omega_m([\,{\cal T},s_m\,]F),                                \tag{96.21}
\]
and hence
\[
 \left|
 \omega_m^{\cal T}(s_mF)-a_m({\cal T}F)
 \right|
 \leq\|[\,{\cal T},s_m\,]F\|.                                   \tag{96.22}
\]
If \({\cal T}\) commutes with \(s_m\), an anomaly-free state remains
anomaly free.  More generally, if a family \({\cal T}_R\) has
\[
 \gamma_R:=
 \sup_{\|F\|_{s_m}\leq1}
 \|[\,{\cal T}_R,s_m\,]F\|\longrightarrow0,                      \tag{96.23}
\]
then its induced Ward defect vanishes uniformly on the graph unit
ball.
\end{proposition}

\begin{proof}
Use (96.20) to write
\[
 {\cal T}(s_mF)=s_m({\cal T}F)+[\,{\cal T},s_m\,]F.
\]
Apply \(\omega_m\), use the anomalous Ward identity on
\({\cal T}F\), and bound the last expectation by the operator norm.
\end{proof}

For the spectral channel \({\cal T}_R\) of V93, the trace-norm
compactness estimate (93.5) does not control \(\gamma_R\).  One must
separately prove that \(P_R\), \(Q_R\), and the replacement state
\(\sigma_R\) respect the BRST graph domain and that the commutator
in (96.23) vanishes.  Compact-group equivariance in V94 gives this
for infinitesimal compact symmetry generators on a common analytic
core, but it does not give the nonlinear BRST identity with ghosts
or diffeomorphisms.

\subsection{Descent to degree-zero cohomology}

If \(a_m=0\), then
\[
 \omega_m(s_mG)=0                                               \tag{96.24}
\]
for every degree-minus-one \(G\in{\cal D}_m\).  Hence on
degree-zero closed elements, the expectation depends only on the
BRST cohomology class:
\[
 F_1-F_2=s_mG
 \quad\Longrightarrow\quad
 \omega_m(F_1)=\omega_m(F_2).                                   \tag{96.25}
\]
If \(a_m\neq0\), equation (96.6) records exactly the failure of
(96.25).

This is a linear cohomological descent.  Positivity on a physical
C-star quotient requires an independently constructed positive
observable algebra or BRST reconstruction.  It is not inferred from
an indefinite ghost extension.

\begin{assessment}[Status after V96]
\label{ass:v12v96-status}
Predual-norm local convergence now carries a regulated anomalous
Ward identity to the complete graph domain of a closed BRST
derivation with the explicit bound (96.9).  The anomaly forms a
natural projective family under BRST-intertwining local inclusions,
and (96.18) separates every finite-cutoff compatibility defect.
Finite-rank compactness channels preserve the constraint precisely
when their graph commutator is controlled.
\end{assessment}

\begin{firewall}[Claim boundary after V96]
\label{fire:v12v96-claim}
No full SM--Einstein closed nilpotent BRST derivation or common graph
domain has been constructed.  The anomaly functionals have not been
shown to vanish, and the bounds
\(\kappa_m^{(n)},\epsilon_m^{(n)},\chi_{m\ell}^{(n)}\), or
\(\gamma_R\) have not been populated for the RPESM regulator.
The theorem does not solve Gribov, global-anomaly, chiral-measure, or
diffeomorphism-cohomology problems.  The full SM--Einstein continuum
remains open.
\end{firewall}

\subsection{Next bottleneck}

The next analytic target is the spectral-channel commutator
\(\gamma_R\).  A sufficient route is a relative BRST-energy bound
of the form
\[
 \|s_mF\|_{\Lambda,\mathrm{tail}}
 \leq C_m\|F\|_{s_m}
\]
together with off-diagonal decay of \(s_m\) between low and high
reference-energy sectors.  The next note should formulate and prove
an explicit block-matrix estimate converting such off-diagonal
decay into (96.23), while keeping any zero-mode or anomaly block
outside the estimate.

\subsection{Parent record}

This V96 note appends to the validated V95 file with record
\[
\begin{array}{c|c}
\text{lines}&26274\\
\text{bytes}&896978\\
\text{SHA--256}&
\texttt{BBADF21954DF01F1E04D9DED25B4C337456FF6E7260E9E084C4BE3CBE3EBB21F}.
\end{array}
\]

% =============================================================================
% END APPENDED TWELFTH-CYCLE V96 MATERIAL
% =============================================================================

% =============================================================================
% BEGIN APPENDED TWELFTH-CYCLE V97 MATERIAL
% =============================================================================

\section{V97: a BRST-Hodge reference cutoff eliminates the spectral
commutator}
\label{sec:v12v97}

V96 leaves the graph commutator of a generic reference-energy
channel as a new stock.  Estimating a sharp moving cutoff directly
is usually wasteful and may fail in operator norm.  The upstream
repair is to choose the reference operator from the BRST complex
itself.  Spectral projections of its Hodge Laplacian reduce the BRST
charge, so the off-diagonal cutoff commutator vanishes identically.

\subsection{The closed Hilbert complex}

Let \({\cal H}={\cal H}^{\bar0}\oplus{\cal H}^{\bar1}\) be a graded
separable Hilbert space.  Let
\[
 {\mathbb Q}:\operatorname{Dom}{\mathbb Q}\subset{\cal H}
 \longrightarrow{\cal H}                                      \tag{97.1}
\]
be a closed densely defined odd operator satisfying
\[
 \operatorname{Ran}{\mathbb Q}\subset\operatorname{Dom}{\mathbb Q},
 \qquad {\mathbb Q}^2=0.                                       \tag{97.2}
\]
The closed nonnegative form
\[
 h[\psi]=\|{\mathbb Q}\psi\|^2+\|{\mathbb Q}^*\psi\|^2,
 \qquad
 \operatorname{Dom}h=
 \operatorname{Dom}{\mathbb Q}\cap
 \operatorname{Dom}{\mathbb Q}^*                               \tag{97.3}
\]
defines the BRST-Hodge Laplacian
\[
 \Delta_{\mathbb Q}
 ={\mathbb Q}^*{\mathbb Q}+{\mathbb Q}{\mathbb Q}^*.             \tag{97.4}
\]

\begin{lemma}[Hodge spectral reduction]
\label{lem:v12v97-hodge}
For every bounded Borel function \(f\),
\(f(\Delta_{\mathbb Q})\) preserves
\(\operatorname{Dom}{\mathbb Q}\) and
\[
 {\mathbb Q}f(\Delta_{\mathbb Q})\psi
 =
 f(\Delta_{\mathbb Q}){\mathbb Q}\psi,
 \qquad \psi\in\operatorname{Dom}{\mathbb Q}.                   \tag{97.5}
\]
In particular, every spectral projection
\[
 P_R=\mathbf1_{[0,R]}(\Delta_{\mathbb Q})                        \tag{97.6}
\]
reduces \({\mathbb Q}\), and
\[
 \ker\Delta_{\mathbb Q}
 =\ker{\mathbb Q}\cap\ker{\mathbb Q}^*.                         \tag{97.7}
\]
\end{lemma}

\begin{proof}
On the common operator core of the two products in (97.4),
nilpotence gives
\[
 {\mathbb Q}\Delta_{\mathbb Q}
 =
 {\mathbb Q}{\mathbb Q}^*{\mathbb Q}
 =
 \Delta_{\mathbb Q}{\mathbb Q}.                                \tag{97.8}
\]
Apply (97.8) to the resolvent equation and use closedness of
\({\mathbb Q}\) to obtain
\[
 {\mathbb Q}(\Delta_{\mathbb Q}+\lambda)^{-1}
 =
 (\Delta_{\mathbb Q}+\lambda)^{-1}{\mathbb Q},
 \qquad\lambda>0,                                               \tag{97.9}
\]
on \(\operatorname{Dom}{\mathbb Q}\).  The bounded Borel functional
calculus, first for resolvent-generated functions and then by
monotone spectral approximation, gives (97.5) and (97.6).
Finally,
\[
 \langle\psi,\Delta_{\mathbb Q}\psi\rangle
 =\|{\mathbb Q}\psi\|^2+\|{\mathbb Q}^*\psi\|^2
\]
on the form domain, which proves (97.7).
\end{proof}

The lemma is the infinite-dimensional closed-complex counterpart of
the finite Ward-Laplacian decomposition used in V7.  The new role is
analytic: when \(\Delta_{\mathbb Q}\) has compact resolvent, its
spectral projections are simultaneously finite rank and exactly
BRST reducing.

\subsection{The bounded commutator domain}

For a homogeneous bounded operator \(A\), define
\[
 s(A)=[{\mathbb Q},A]_{\rm gr}
 ={\mathbb Q}A-(-1)^{|A|}A{\mathbb Q}                           \tag{97.10}
\]
whenever \(A\) preserves \(\operatorname{Dom}{\mathbb Q}\) and the
displayed operator extends boundedly.  Let \({\cal D}_{\mathbb Q}\)
be a graded algebraic core on which (97.10) is defined and stable.
Lemma~\ref{lem:v12v97-hodge} gives
\[
 s(P_R)=0,\qquad s(I-P_R)=0.                                    \tag{97.11}
\]

Let \(\sigma_R\) be any density matrix supported in
\(P_R{\cal H}\), and form the normal finite-rank channel
\[
 {\cal T}_R(A)
 =
 P_RAP_R+\operatorname{Tr}(\sigma_RA)(I-P_R).                   \tag{97.12}
\]
For homogeneous \(A\in{\cal D}_{\mathbb Q}\), the exact commutator
formula is
\[
 [\,{\cal T}_R,s\,]A
 =
 \operatorname{Tr}(\sigma_Rs(A))(I-P_R).                        \tag{97.13}
\]

\begin{proof}
Use the graded derivation rule and (97.11):
\[
 s(P_RAP_R)=P_Rs(A)P_R,
 \qquad
 s\!\left(\operatorname{Tr}(\sigma_RA)(I-P_R)\right)=0.
\]
On the other hand,
\[
 {\cal T}_R(s(A))
 =
 P_Rs(A)P_R+
 \operatorname{Tr}(\sigma_Rs(A))(I-P_R).
\]
Subtract the two expressions.
\end{proof}

Thus the entire cutoff commutator is concentrated in the Ward
functional of the one replacement state.  There are no uncontrolled
off-diagonal blocks at the spectral boundary.

\subsection{Exact preservation from a harmonic replacement state}

Assume
\[
 {\cal H}_{\rm harm}:=\ker\Delta_{\mathbb Q}\neq\{0\}.            \tag{97.14}
\]
Choose a positive trace-one \(\sigma_0\) supported in a
finite-dimensional subspace of \({\cal H}_{\rm harm}\).  Then
\[
 {\mathbb Q}\sigma_0=0,\qquad
 \sigma_0{\mathbb Q}=0,                                        \tag{97.15}
\]
on the natural domains, and consequently
\[
 \operatorname{Tr}(\sigma_0s(A))=0,
 \qquad A\in{\cal D}_{\mathbb Q}.                               \tag{97.16}
\]

\begin{theorem}[BRST-exact compactness channel]
\label{thm:v12v97-exact}
Suppose \(\Delta_{\mathbb Q}\) has compact resolvent and use
\(\sigma_0\) in (97.12).  Then \({\cal T}_R\) is normal, unital,
completely positive, finite rank, and
\[
 {\cal T}_R s=s{\cal T}_R
 \quad\text{on }{\cal D}_{\mathbb Q}.                            \tag{97.17}
\]
If a normal state \(\omega\) has anomaly \(a=\omega\circ s\), then
\[
 ({\cal T}_{R*}\omega)\circ s=a\circ{\cal T}_R.                  \tag{97.18}
\]
In particular, an anomaly-free state remains anomaly free under
every cutoff channel.
\end{theorem}

\begin{proof}
The channel properties are those of (93.2).  Equations
(97.13) and (97.16) prove (97.17).  For \(A\) in the core,
\[
 ({\cal T}_{R*}\omega)(sA)
 =\omega({\cal T}_R(sA))
 =\omega(s({\cal T}_RA))
 =a({\cal T}_RA),
\]
which is (97.18).
\end{proof}

The channel transports an anomaly rather than cancelling it.
Vanishing still requires the retained cohomological tests of
V19--V23.

\subsection{Tail-suppressed replacement defect}

If no harmonic density \(\sigma_0\) is available, retain an arbitrary
\(\sigma_R\) supported below \(R\).  Let \(\omega\) be anomaly free
and assume the BRST-Hodge energy moment
\[
 \omega(\Delta_{\mathbb Q})\leq L.                              \tag{97.19}
\]
Then (97.13) yields a defect smaller than the operator-norm estimate
suggested in V96.

\begin{proposition}[Energy-tail Ward defect]
\label{prop:v12v97-tail}
For every \(A\in{\cal D}_{\mathbb Q}\),
\[
 |({\cal T}_{R*}\omega)(sA)|
 \leq\frac{L}{R}\,
       |\operatorname{Tr}(\sigma_Rs(A))|
 \leq\frac{L}{R}\|s(A)\|.                                      \tag{97.20}
\]
Hence
\[
 \|({\cal T}_{R*}\omega)\circ s\|_{
       ({\cal D}_{\mathbb Q},\|\cdot\|_s)^*}
 \leq L/R.                                                      \tag{97.21}
\]
\end{proposition}

\begin{proof}
Because \(\omega\circ s=0\),
\[
 \omega(s({\cal T}_RA))=0.
\]
Use (97.13) in
\[
 ({\cal T}_{R*}\omega)(sA)
 =\omega({\cal T}_R(sA))
 =\omega(s({\cal T}_RA))
  +\operatorname{Tr}(\sigma_Rs(A))\omega(I-P_R).                 \tag{97.22}
\]
Spectral Markov gives
\[
 \omega(I-P_R)\leq L/R.                                        \tag{97.23}
\]
The state bound
\(|\operatorname{Tr}(\sigma_Rs(A))|\leq\|s(A)\|\)
proves (97.20), and the graph unit ball proves (97.21).
\end{proof}

This replaces the strong operator commutator demand (96.23) by a
state-weighted defect whenever the reference cutoff is Hodge
adapted.  It is enough for Theorem~\ref{thm:v12v96-graph}.

\subsection{Simultaneous compactness and constraint closure}

Combining V93 with the Hodge choice gives the following conditional
package.  Suppose a family of local normal states satisfies
\[
 \sup_n\omega_m^{(n)}(\Delta_{{\mathbb Q},m})\leq L_m,            \tag{97.24}
\]
and \(\Delta_{{\mathbb Q},m}\) has compact resolvent.  Then the
channels (97.12) give the trace-norm compactness estimate
\[
 \sup_n\|\rho_m^{(n)}
       -{\cal T}_{R*}\rho_m^{(n)}\|_1
 \leq2\sqrt{L_m/R}+2L_m/R.                                     \tag{97.25}
\]
If a harmonic replacement state exists, the Ward defect is exactly
zero.  Otherwise its graph-dual norm is at most \(L_m/R\).
Thus one and the same cutoff proves local predual compactness and
constraint closure.

The price is a new physical comparison.  The local Hamiltonian used
in V93 must control the Hodge energy:
\[
 K_{m,x}\geq a_m\Delta_{{\mathbb Q},m}-b_mI                     \tag{97.26}
\]
uniformly over the admitted background core.  Equation (97.26) is
stronger than ordinary elliptic coercivity if the BRST charge
contains ghost, gauge, or diffeomorphism directions not seen by
\(K_{m,x}\).  Those zero and longitudinal blocks must be retained
explicitly rather than hidden in \(b_m\).

\begin{assessment}[Status after V97]
\label{ass:v12v97-status}
A Hodge-adapted compact-resolvent reference operator eliminates the
spectral BRST commutator exactly.  With a harmonic replacement state,
the finite-rank compactness channel preserves the anomalous Ward
identity exactly as in (97.18).  Without one, an anomaly-free state
acquires only the explicit energy-tail defect \(L/R\).  This joins
local predual compactness and graph-norm constraint closure through
one spectral construction.
\end{assessment}

\begin{firewall}[Claim boundary after V97]
\label{fire:v12v97-claim}
No RPESM continuum BRST charge has been constructed as a closed
Hilbert-complex differential, and compact resolvent of its Hodge
Laplacian is not known.  The uniform comparison (97.26), existence
of a positive harmonic replacement density, and compatibility with
local inclusions have not been proved.  Indefinite ghost inner
products may require a Hilbert majorant before this argument
applies.  Global anomalies and Gribov obstructions remain separate.
The full SM--Einstein continuum is not yet constructed.
\end{firewall}

\subsection{Next bottleneck}

The remaining upstream question is whether the BRST-Hodge operator
has compact resolvent after removal of gauge and diffeomorphism zero
modes.  A direct global claim would be circular.  The next note
should split a finite-dimensional retained cohomology block from an
orthogonal positive block, prove compactness and moment bounds only
on the latter, and carry the retained block through the projective
anomaly ledger without assigning it a false spectral gap.

\subsection{Parent record}

This V97 note appends to the validated V96 file with record
\[
\begin{array}{c|c}
\text{lines}&26623\\
\text{bytes}&908395\\
\text{SHA--256}&
\texttt{BB240421CE256CC4AED92674C6578EAD33DE2E9C0DDE7914D2EB95D2CEBA0BD4}.
\end{array}
\]

% =============================================================================
% END APPENDED TWELFTH-CYCLE V97 MATERIAL
% =============================================================================

% =============================================================================
% BEGIN APPENDED TWELFTH-CYCLE V98 MATERIAL
% =============================================================================

\section{V98: zero-mode-safe local BRST compactness and continuum
closure}
\label{sec:v12v98}

V97 uses a compact-resolvent BRST-Hodge operator.  A direct gap
claim at zero would be wrong: harmonic vectors represent retained
cohomology and must not be penalized merely to make an estimate
coercive.  This note keeps the harmonic block exactly and imposes
compactness only on the positive block.  It then combines the
result with V91 and V96 into one local continuum-closure theorem.

\subsection{Retained harmonic block}

For each region \(m\), let
\[
 \Delta_m=\Delta_{{\mathbb Q},m},\qquad
 P_m^0=\mathbf1_{\{0\}}(\Delta_m),\qquad
 {\cal H}_m^0=P_m^0{\cal H}_m.                                  \tag{98.1}
\]
Assume
\[
 d_m:=\dim{\cal H}_m^0<\infty                                   \tag{98.2}
\]
and that the restriction
\[
 \Delta_m^+:=\Delta_m\big|_{(I-P_m^0){\cal H}_m}                 \tag{98.3}
\]
has compact resolvent.  No lower bound uniform in \(m\) or in the
ultraviolet regulator is assumed.

Define the zero-mode-retaining projection
\[
 P_{m,R}
 =
 P_m^0+\mathbf1_{(0,R]}(\Delta_m).                              \tag{98.4}
\]
It is finite rank.  For any density matrix \(\rho\) with
\[
 \operatorname{Tr}(\rho\Delta_m)\leq L_m,                        \tag{98.5}
\]
spectral Markov gives
\[
 \operatorname{Tr}(\rho(I-P_{m,R}))\leq L_m/R.                  \tag{98.6}
\]
All mass and all coherences inside the harmonic block are retained;
only positive energy above \(R\) is discarded.

Choose a density matrix \(\sigma_m^0\) supported in
\({\cal H}_m^0\) and set
\[
 {\cal T}_{m,R}(A)
 =
 P_{m,R}AP_{m,R}
 +\operatorname{Tr}(\sigma_m^0A)(I-P_{m,R}).                    \tag{98.7}
\]
By V97, this channel commutes with the BRST derivation on its bounded
commutator core.

\begin{proposition}[Zero-mode-safe compact core]
\label{prop:v12v98-core}
For every density matrix satisfying (98.5),
\[
 \|\rho-{\cal T}_{m,R*}\rho\|_1
 \leq2\sqrt{L_m/R}+2L_m/R.                                     \tag{98.8}
\]
The family of all such states is relatively compact in trace norm,
and the approximating channel is exactly BRST equivariant.
\end{proposition}

\begin{proof}
Equation (98.6) replaces (93.6), and the block estimate
(93.7)--(93.8) gives (98.8).  The range of the channel is finite
dimensional because \(P_{m,R}\) is finite rank, so
Proposition~\ref{prop:v12v92-finite-rank} proves compactness.
Since \(\sigma_m^0\) is harmonic,
Theorem~\ref{thm:v12v97-exact} gives exact BRST equivariance.
\end{proof}

This proof neither inverts \(\Delta_m\) on its kernel nor uses the
smallest positive eigenvalue.  It therefore remains valid when the
positive spectral gap tends to zero along a cofinal sequence, as
long as the fixed-region compactness and moment bounds are uniform
in the ultraviolet regulator.

\subsection{Integrated local closure theorem}

Let \(\iota_{m\ell}:{\cal B}({\cal H}_m)\to
{\cal B}({\cal H}_\ell)\) be the normal local inclusions from V91.
For every regulator \(n\), let \(\rho_m^{(n)}\) be a local density
matrix and write \(\omega_m^{(n)}\) for its state.

\begin{theorem}[Zero-mode-safe BRST-normal projective limit]
\label{thm:v12v98-integrated}
Assume, for every fixed region \(m\):

\begin{enumerate}
\item the closed nilpotent charge \({\mathbb Q}_m\) satisfies
(98.2)--(98.3);
\item the uniform positive-block moment holds:
\[
 \sup_n\operatorname{Tr}(\rho_m^{(n)}\Delta_m)\leq L_m<\infty;
                                                                    \tag{98.9}
\]
\item for \(m\leq\ell\), the predual compatibility defect obeys
\[
 \|\omega_m^{(n)}
 -R_{m\ell}\omega_\ell^{(n)}\|\longrightarrow0;                 \tag{98.10}
\]
\item the inclusions are graph bounded and intertwine the limiting
BRST derivations as in (96.12)--(96.13);
\item the regulated Ward data satisfy the four graph-norm limits in
(96.8).
\end{enumerate}

Then there is one subsequence \(n_j\) and normal local states
\(\omega_m\) such that
\[
 \|\omega_m^{(n_j)}-\omega_m\|\longrightarrow0
 \quad\text{for every fixed }m,                                 \tag{98.11}
\]
the family is projectively compatible,
\[
 \omega_m=R_{m\ell}\omega_\ell,                                 \tag{98.12}
\]
and it satisfies on the completed graph domain
\[
 \omega_m\circ s_m=a_m,\qquad
 a_m=\iota_{m\ell}^{\sharp}a_\ell,\qquad
 a_m\circ s_m=0.                                                \tag{98.13}
\]
It determines a unique quasi-local state on the norm closure of the
directed local union.  If \(a_m=0\) for every \(m\), its local
expectations descend to degree-zero BRST cohomology.
\end{theorem}

\begin{proof}
Proposition~\ref{prop:v12v98-core} makes
\(\{\omega_m^{(n)}\}_n\) relatively compact in predual norm for each
fixed \(m\).  The diagonal theorem
\ref{thm:v12v91-projective}, using (98.10), gives one subsequence,
(98.11), (98.12), and the quasi-local extension.

The graph-norm convergence and Ward residual assumptions allow
Theorem~\ref{thm:v12v96-graph} at every fixed \(m\), giving the first
and third identities in (98.13).  Proposition
\ref{prop:v12v96-natural}, together with (98.12) and the chain-map
property of the inclusions, gives the middle identity.  Finally,
(96.24)--(96.25) give cohomological descent when every \(a_m\)
vanishes.
\end{proof}

The theorem permits \(L_m\), \(d_m\), and every finite-region
spectral scale to grow with \(m\).  It is a local-normal theorem and
does not assert a trace-class density matrix in infinite volume.

\subsection{Two harmonic spaces that must not be identified}

There are now two finite-dimensional objects with different roles:
\[
 \begin{array}{c|c|c}
 \text{object}&\text{definition}&\text{role}\\ \hline
 {\cal H}_m^0&
 \ker\Delta_{{\mathbb Q},m}\text{ in the state Hilbert space}&
 \text{zero modes retained by }P_{m,R}\\
 {\mathfrak H}_m^1&
 \text{harmonic degree-one regulator cochains of V19}&
 \text{retained anomaly class}.
 \end{array}                                                     \tag{98.14}
\]
Finite dimensionality of \({\cal H}_m^0\) proves compactness of the
state zero-mode block.  It does not imply that the anomaly vector in
\({\mathfrak H}_m^1\) vanishes.  Conversely, cancellation of the
degree-one anomaly does not remove physical states from
\({\cal H}_m^0\).

Let \(H_m^{\rm an}\) denote the harmonic projection on the anomaly
cochain complex and define
\[
 A_m^{\rm an}:=H_m^{\rm an}a_m.                                \tag{98.15}
\]
The projective identity in (98.13) transports this vector through
the cohomology maps only when the refinement hypotheses of V8 and
V19 hold.  Compactness merely gives convergent subsequences of the
finite-dimensional vectors \(A_m^{\rm an}\); it does not force their
limit to be zero.  Exact species cancellation or the lattice
systole criterion of V21 is still required.

\subsection{No mass-gap inference}

For each fixed \(m\), compact resolvent on the positive block implies
that the first positive eigenvalue is nonzero.  This is a
finite-region fact.  Neither (98.9) nor the proof of Theorem
\ref{thm:v12v98-integrated} supplies a lower bound uniform in
\(m\).  Therefore the construction makes no assertion about an
infinite-volume particle or transfer-Hamiltonian mass gap.  This is
especially necessary in the presence of the photon and graviton
channels already isolated in V84.

\begin{assessment}[Status after V98]
\label{ass:v12v98-status}
Finite-dimensional BRST harmonic state modes can be retained
exactly while energy compactness is imposed only on the positive
Hodge block.  Under local moment, compatibility, and graph-Ward
stocks, one subsequence now yields a projectively consistent,
locally normal quasi-local state with a natural Wess--Zumino
anomaly family.  If that family vanishes by the separate anomaly
ledger, expectations descend to local degree-zero BRST cohomology.
\end{assessment}

\begin{firewall}[Claim boundary after V98]
\label{fire:v12v98-claim}
The hypotheses of the integrated theorem have not been populated
for a full RPESM regulator.  Finite-dimensional state cohomology,
positive-block compact resolvent, uniform Hodge moments, graph
intertwining, and anomaly cancellation all remain separate physical
tasks.  No uniform positive spectral gap, mass gap, global
trace-class state, or positive BRST quotient has been proved.  The
full SM--Einstein continuum remains open.
\end{firewall}

\subsection{Next bottleneck}

The next note should address the energy comparison without charging
the retained zero modes.  The required statement is a uniform form
inequality controlling only
\((I-P_m^0)\Delta_m(I-P_m^0)\) by the gauge-fixed local RPESM
Hamiltonian.  Its proof must expose the principal-symbol coercivity,
boundary conditions, and lower-order constants, and must leave the
finite harmonic block untouched.

\subsection{Parent record}

This V98 note appends to the validated V97 file with record
\[
\begin{array}{c|c}
\text{lines}&26956\\
\text{bytes}&919462\\
\text{SHA--256}&
\texttt{65D7EE644C9C6246858B4F3F0E24C9A33EB3BD518AB1F819C1505E9BA5DB78D9}.
\end{array}
\]

% =============================================================================
% END APPENDED TWELFTH-CYCLE V98 MATERIAL
% =============================================================================

% =============================================================================
% BEGIN APPENDED TWELFTH-CYCLE V99 MATERIAL
% =============================================================================

\section{V99: projected G{\aa}rding coercivity without charging BRST
zero modes}
\label{sec:v12v99}

V98 requires a uniform moment of the positive BRST-Hodge operator.
The physical Gibbs estimate, however, is naturally stated for a
background-dependent gauge-fixed local Hamiltonian.  This note gives
a form comparison which leaves the finite harmonic block at zero
energy and controls all cross terms explicitly.

\subsection{Common form domain and harmonic splitting}

Fix one region and suppress its index.  Let
\[
 {\cal H}={\cal H}^0\oplus{\cal H}^+,\qquad
 P_0:{\cal H}\to{\cal H}^0,\qquad P_+=I-P_0,                    \tag{99.1}
\]
where \({\cal H}^0=\ker\Delta_{\mathbb Q}\) is finite-dimensional.
Let \({\cal V}\) be the common form domain of
\(\Delta_{\mathbb Q}\) and a family of closed semibounded forms
\(k_x\).  Assume
\[
 P_0{\cal V}\subset{\cal V},\qquad P_+{\cal V}\subset{\cal V}.    \tag{99.2}
\]
Write \(u=h+v\), with \(h=P_0u\) and \(v=P_+u\).

Choose a Hilbert norm \(\|\cdot\|_{\cal V}\) on the form domain.
Assume there are constants \(c>0\), \(b_+<\infty\),
\(b_0<\infty\), and \(b_\times<\infty\), independent of the
background \(x\), such that
\[
 \begin{split}
 k_x[v]&\geq c\|v\|_{\cal V}^2-b_+\|v\|^2,
 &&v\in P_+{\cal V},\\
 k_x[h]&\geq-b_0\|h\|^2,
 &&h\in{\cal H}^0,\\
 2|\operatorname{Re}k_x[v,h]|
 &\leq\frac c2\|v\|_{\cal V}^2+b_\times\|h\|^2.
                                                                    \tag{99.3}
 \end{split}
\]
The last estimate includes the lower-order and boundary cross
terms.  It is automatic with a uniform constant once those terms
are form bounded on \(P_+{\cal V}\times{\cal H}^0\), because
\({\cal H}^0\) is finite-dimensional, but it must still be verified
for the selected boundary realization.

Assume also that the positive Hodge form has the upper comparison
\[
 \langle v,\Delta_{\mathbb Q}v\rangle
 \leq c_\Delta\|v\|_{\cal V}^2,
 \qquad v\in P_+{\cal V}.                                      \tag{99.4}
\]

\begin{theorem}[Zero-mode-safe form comparison]
\label{thm:v12v99-comparison}
Under (99.1)--(99.4),
\[
 k_x[u]\geq
 \frac{c}{2c_\Delta}
 \langle u,\Delta_{\mathbb Q}u\rangle
 -b\|u\|^2,\qquad u\in{\cal V},                                \tag{99.5}
\]
uniformly in \(x\), where one may take
\[
 b=b_++b_0+b_\times.                                            \tag{99.6}
\]
Equivalently, for the self-adjoint operator \(K_x\) associated with
\(k_x\),
\[
 K_x+bI\geq a\Delta_{\mathbb Q}
 \quad\text{in form sense},\qquad
 a=\frac{c}{2c_\Delta}.                                         \tag{99.7}
\]
The right side vanishes on \({\cal H}^0\); no energy penalty or
inverse Hodge operator is introduced there.
\end{theorem}

\begin{proof}
Sesquilinearity and (99.3) give
\[
\begin{split}
 k_x[u]
 &=k_x[v]+2\operatorname{Re}k_x[v,h]+k_x[h]\\
 &\geq
 \frac c2\|v\|_{\cal V}^2
 -b_+\|v\|^2-(b_0+b_\times)\|h\|^2\\
 &\geq
 \frac c2\|v\|_{\cal V}^2-b\|u\|^2.                            \tag{99.8}
\end{split}
\]
Since \(\Delta_{\mathbb Q}P_0=0\) and its spectral projection
\(P_0\) reduces the operator,
\[
 \langle u,\Delta_{\mathbb Q}u\rangle
 =\langle v,\Delta_{\mathbb Q}v\rangle
 \leq c_\Delta\|v\|_{\cal V}^2.                                \tag{99.9}
\]
Insert (99.9) into (99.8) to obtain (99.5).  The representation
theorem for closed semibounded forms gives (99.7).
\end{proof}

\subsection{How the projected estimates arise from an elliptic
boundary problem}

Let \(D\) be a compact manifold with boundary, \(E\to D\) a
finite-rank graded bundle, and choose a fixed strongly elliptic local
boundary condition.  After transporting every admitted metric and
bundle density to one reference \(L^2(D;E)\), suppose:

\begin{enumerate}
\item the boundary condition gives one background-independent
closed subspace \({\cal V}\subset H^1(D;E)\);
\item the positive-block principal symbol of \(K_x\) has the uniform
Legendre bound
\[
 \operatorname{Re}
 \langle\sigma_2(K_x)(y,\xi)w,w\rangle
 \geq\kappa|\xi|^2|w|^2                                       \tag{99.10}
\]
for \(w\) in the constrained positive bundle directions, with
\(\kappa>0\) independent of \(x,y,\xi\);
\item all first-order, zeroth-order, and boundary coefficients are
uniformly form bounded on the compact gravitational background
core;
\item the Hodge form has uniformly bounded second-order
coefficients relative to the same \(H^1\) norm.
\end{enumerate}

The standard parameter-uniform G{\aa}rding inequality then gives
the first line of (99.3) on \(P_+{\cal V}\).  For completeness, the
lower-order terms obey, for every \(\eta>0\),
\[
 |\ell_x[v]|
 \leq\eta\|\nabla v\|_2^2+C_\eta\|v\|_2^2,                     \tag{99.11}
\]
with \(C_\eta\) uniform on the coefficient core.  Choose
\(\eta<\kappa\) and absorb the gradient term into (99.10).  The
trace theorem gives the same form estimate for admissible boundary
terms.  Finite dimensionality of \({\cal H}^0\) makes every form
norm equivalent to the \(L^2\) norm on that block; applying Young's
inequality to \(k_x[v,h]\) gives the last two lines of (99.3).
The coefficient upper bound for the Hodge form gives (99.4).

This argument identifies three gates which cannot be replaced by the
word elliptic: the positive constrained symbol, the common boundary
form domain, and uniform lower-order form bounds.

\subsection{Transfer of the Gibbs energy moment}

\begin{corollary}[Physical energy controls positive Hodge energy]
\label{cor:v12v99-moment}
Let \(\rho_x\) be normal states with
\[
 \sup_x\operatorname{Tr}(\rho_xK_x)\leq C.                      \tag{99.12}
\]
Then
\[
 \sup_x\operatorname{Tr}(\rho_x\Delta_{\mathbb Q})
 \leq \frac{C+b}{a}
 =\frac{2c_\Delta(C+b)}{c}.                                    \tag{99.13}
\]
Consequently the zero-mode-retaining channels of V98 satisfy
\[
 \sup_x\|\rho_x-{\cal T}_{R*}\rho_x\|_1
 \leq
 2\sqrt{\frac{C+b}{aR}}+
 2\frac{C+b}{aR}.                                               \tag{99.14}
\]
\end{corollary}

\begin{proof}
Apply the normal functional represented by \(\rho_x\) to increasing
bounded spectral truncations of the form inequality (99.7), then
use monotone convergence.  This gives
\[
 C+b\geq
 \operatorname{Tr}(\rho_x(K_x+bI))
 \geq a\operatorname{Tr}(\rho_x\Delta_{\mathbb Q}),
\]
which is (99.13).  Proposition~\ref{prop:v12v98-core} gives
(99.14).
\end{proof}

If \(\rho_x=Z_x^{-1}e^{-\beta K_x}\), the uniform energy constant
\(C\) follows from the partition floor and heat-trace argument in
Corollary~\ref{cor:v12v93-gibbs}.  Inserting its explicit
\(C_\beta\) in (99.13) supplies the V98 moment \(L\) without a
separate assumption.

\subsection{Background-dependent zero modes}

The proof above uses one reference projection \(P_0\).  If the
kernel of the BRST-Hodge operator varies with \(x\), pointwise
finite dimensionality is insufficient.  A valid replacement
requires:

\begin{enumerate}
\item constant kernel dimension on the admitted background core;
\item a norm-continuous spectral projection \(x\mapsto P_0(x)\);
\item a uniform positive contour separating zero from the rest of
the local spectrum; and
\item a finite-rank bundle trivialization transporting \(P_0(x)\)
to one reference harmonic space.
\end{enumerate}

Under these hypotheses, Kato spectral transport gives bounded
unitaries \(U_x\) with
\[
 U_xP_0(x)U_x^*=P_0(x_0),                                      \tag{99.15}
\]
and the preceding proof applies after conjugation, with uniformly
equivalent form norms.  If the kernel dimension jumps, no such
trivialization exists and the cohomology change must enter the V8
refinement obstruction.  It cannot be absorbed into the constant
\(b\).

\subsection{Relation to Einstein backgrounds}

For an Einstein-background law, the coefficient uniformity must hold
on the compact core selected before the fibre analysis, as in
V69--V70.  Outside that core one needs a separate probability-tail
bound.  The logical order is
\[
 \begin{split}
 &\text{compact gravitational coefficient core}\\
 &\quad\Longrightarrow
 \text{uniform projected G{\aa}rding constants}\\
 &\quad\Longrightarrow
 \text{uniform positive Hodge moment}\\
 &\quad\Longrightarrow
 \text{local predual compactness and BRST closure},
 \end{split}                                                     \tag{99.16}
\]
followed by removal of the background-core tail.  Reversing the
order and averaging nonuniform G{\aa}rding constants can produce an
infinite bound.

\begin{assessment}[Status after V99]
\label{ass:v12v99-status}
The Hodge moment required in V98 now follows from a projected,
parameter-uniform G{\aa}rding comparison.  The proof retains the
finite BRST harmonic block at zero energy, controls its cross terms
by finite dimensionality, and exposes the principal-symbol,
boundary-domain, and lower-order hypotheses.  Combined with the V93
Gibbs estimate, it gives the explicit compactness error (99.14).
\end{assessment}

\begin{firewall}[Claim boundary after V99]
\label{fire:v12v99-claim}
No concrete coupled SM--Einstein gauge-fixed operator has been shown
to obey the positive-symbol bound (99.10), a common strongly
elliptic boundary condition, or the required uniform coefficient
bounds.  Kernel-dimension constancy and Kato transport have not been
proved on the physical gravitational core.  Ghost indefiniteness,
Gribov horizons, chiral determinants, and diffeomorphism zero modes
remain outside the theorem.  The full SM--Einstein continuum is not
yet constructed.
\end{firewall}

\subsection{Rollover instruction}

V99 is the final active note before the mandatory V100 rollover.
At the next iteration, V100 must first perform the companion audit.
The completed V100 file must then be frozen with the next suffix
\(\texttt{\_f12}\), and the active sequence restarted at V1 with a
compact context and major-results summary.  At least one substantive
post-rollover version must follow that summary.

\subsection{Parent record}

This V99 note appends to the validated V98 file with record
\[
\begin{array}{c|c}
\text{lines}&27206\\
\text{bytes}&928937\\
\text{SHA--256}&
\texttt{695BAB518603C908F8005E4A26FCCB34967C60DEC00F56D990968A0348806EA5}.
\end{array}
\]

% =============================================================================
% END APPENDED TWELFTH-CYCLE V99 MATERIAL
% =============================================================================

% =============================================================================
% BEGIN APPENDED TWELFTH-CYCLE V100 MATERIAL
% =============================================================================

\section{V100 audit: exact structural floors before background
perturbation}
\label{sec:v12v100}

This is the mandatory five-version audit following V95 and the
terminal audit of the twelfth compact cycle.  The newest companion
files examined were
\[
\begin{array}{c|c|c|c}
\text{programme}&\text{version}&\text{bytes}&\text{SHA--256}\\ \hline
\mathrm{YM}&63&801223&
\texttt{D91787B1DCC4E770D518E96F1CDC9E90F808D1979B2C956369CE48A1BA2769AA}\\
\mathrm{NS}&26&278283&
\texttt{82C887F8C2C69E4F28FAD7C3DC8DE5B9099CB5A4BF431EBA1FFF3E91935978AD}.
\end{array}
\]
The Navier--Stokes file is unchanged, so its earlier results are not
repeated.

\subsection{New Yang--Mills input}

Yang--Mills V63 supplies a genuinely new upstream result.  Its
positive owner endpoint has the exact factorization
\[
 H_{\rm rr}=K_R^*K_R,\qquad K_R=ED_R,                            \tag{100.1}
\]
and a sixteen-alias constraint permits the global real-momentum
Loewner estimate
\[
 H_{\rm rr}(k)\succeq2D_R(k)^*D_R(k).                            \tag{100.2}
\]
Since \(D_R^*D_R=3H_{\rm W}\) in that construction, the direct
endpoint satisfies
\[
 H_*(k)=H_{\rm W}(k)+H_{\rm rr}(k)
 \succeq7H_{\rm W}(k)
 \succeq\frac{343}{1250}I.                                     \tag{100.3}
\]
This is an exact all-momentum result, not a grid inference.  It
improves the endpoint inverse floor by a factor of seven relative to
the inherited Wilson value.  Its claim boundary remains flat,
quadratic, reduced-fibre, and real; it does not establish a
continuum response sign or a mass gap.

\subsection{Transfer to the SM--Einstein line}

V99 currently obtains the positive Hodge comparison from a generic
projected G{\aa}rding inequality.  V63 shows why a sharper order of
proof is preferable when exact structure is available:
\[
 \text{factorization and constraint}
 \ \Longrightarrow\ \text{Loewner floor}
 \ \Longrightarrow\ \text{inverse and heat bounds},             \tag{100.4}
\]
rather than repeatedly propagating entrywise absolute values.

The SM analogue is conditional but precise.  On the positive
BRST-Hodge block, seek a factorization
\[
 K_{x_0}^+={\cal C}_{x_0}^*{\cal C}_{x_0},\qquad
 \Delta_{\mathbb Q}^+={\cal D}^*{\cal D},                        \tag{100.5}
\]
at one reference background \(x_0\), together with a constrained
estimate
\[
 \|{\cal C}_{x_0}v\|^2\geq a_0\|{\cal D}v\|^2.                  \tag{100.6}
\]
If the background perturbation form has relative size
\(\eta_x<a_0\), then
\[
 K_x^+\succeq(a_0-\eta_x)\Delta_{\mathbb Q}^+
             -b_xI.                                            \tag{100.7}
\]
The exact reference floor \(a_0\) should be proved before estimating
\(\eta_x\); otherwise the perturbation budget may consume an
artificially weak inherited constant.

\begin{assessment}[Decision after the V100 audit]
\label{ass:v12v100-decision}
The twelfth cycle ends with the local-normal BRST closure reduced to
a positive-block structural coercivity problem.  After rollover, the
first substantive note will prove an abstract
factorization-to-Loewner theorem and a sharp relative-form
perturbation lemma which preserve the harmonic kernel.  This imports
the proof architecture of Yang--Mills V63 without importing its
flat numerical constant into the curved SM--Einstein problem.
\end{assessment}

\begin{firewall}[Claim boundary after V100]
\label{fire:v12v100-claim}
No factorization (100.5), constrained comparison (100.6), or
background perturbation stock has been constructed for the coupled
RPESM Hessian.  The Yang--Mills factor seven is specific to its
sixteen-alias endpoint and cannot be reused as an SM--Einstein
constant.  No continuum measure, anomaly cancellation, reflection
positivity, clustering, or full SM--Einstein continuum follows from
this audit.
\end{firewall}

\subsection{Rollover record}

This V100 note completes the twelfth active cycle.  It will be frozen
as
\[
 \texttt{Renewal\_SM\_Einstein\_Continuum\_Research\_Notes\_V100\_f12.tex}.
                                                                    \tag{100.8}
\]
A new thirteenth-cycle V1 will contain context, a compact statement
of the major acquired results, the open gates, and the reference
record of this archive.  Subsequent versions will append to that new
file.

\subsection{Parent record}

This V100 note appends to the validated V99 file with record
\[
\begin{array}{c|c}
\text{lines}&27495\\
\text{bytes}&939485\\
\text{SHA--256}&
\texttt{24A82C3CB690DC32E2EBB739C47E645E7E83A7CA2FC5E9A8F95BE0596EFDDFC4}.
\end{array}
\]

% =============================================================================
% END APPENDED TWELFTH-CYCLE V100 MATERIAL
% =============================================================================

\end{document}
